\documentclass[11pt,a4paper]{article}
\usepackage[utf8]{inputenc}
\usepackage[T1]{fontenc}
\usepackage{geometry}
\usepackage{amsmath}
\usepackage{amssymb}
\usepackage{listings}
\usepackage{xcolor}
\usepackage{hyperref}
\usepackage{fancyhdr}
\usepackage{titlesec}

\geometry{margin=2.5cm}

% Code listing style for Witness
\lstdefinestyle{witnessstyle}{
    language=C++,
    basicstyle=\ttfamily\small,
    keywordstyle=\color{blue}\bfseries,
    commentstyle=\color{gray}\itshape,
    stringstyle=\color{red},
    numbers=left,
    numberstyle=\tiny\color{gray},
    stepnumber=1,
    numbersep=5pt,
    frame=single,
    breaklines=true,
    breakatwhitespace=true,
    tabsize=2,
    showstringspaces=false
}

\lstdefinelanguage{Witness}{
    keywords={asset, clause, subject, action, object, service, oblig, claim, not, AND, OR, IMPLIES, EQUIV, global},
    sensitive=true,
    comment=[l]{//},
    morecomment=[s]{/*}{*/},
    morestring=[b]"
}

\title{Legal Theses Formalization Documentation}
\author{Computational Analysis of Judicial Consistency}
\date{\today}

\begin{document}
\maketitle
\tableofcontents
\newpage

\section{Thesis 161518: PRESCRIPCIÓN POSITIVA. PARA ADQUIRIR UN BIEN INMUEBLE A TRAVÉS DE ESTA FIGURA, SIN NECESIDAD DE TÍTULO, ES MENESTER QUE SE DEMUESTRE QUE SE}
\label{thesis:161518}

\subsection{Original Thesis Text}

\begin{quote}
                                                   Semanario Judicial de la Federación

                                                  Tesis

 Registro digital: 161518

 Instancia: Tribunales             Novena Época                           Materia(s): Civil
 Colegiados de Circuito

 Tesis: XXXI. J/5                  Fuente: Semanario Judicial de la  Tipo: Jurisprudencia
                                   Federación y su Gaceta.
                                   Tomo XXXIV, Julio de 2011, página
                                   1880


PRESCRIPCIÓN POSITIVA. PARA ADQUIRIR UN BIEN INMUEBLE A TRAVÉS DE ESTA
FIGURA, SIN NECESIDAD DE TÍTULO, ES MENESTER QUE SE DEMUESTRE QUE SE
ADQUIRIÓ LA POSESIÓN EN CONCEPTO DE DUEÑO O DE PROPIETARIO, Y NO EN FORMA
DERIVADA NI PRECARIA (LEGISLACIÓN DEL ESTADO DE CAMPECHE).


El numeral 1157 del Código Civil del Estado establece: "La posesión necesaria para prescribir debe
ser: I. En concepto de propietario; II. Pacífica; III. Continua; IV. Pública.". Por su parte, el artículo
1158 del mismo ordenamiento legal contempla diversas hipótesis en las que procede la prescripción
de bienes inmuebles, distinguiendo aquellos casos en los que la posesión se ejerce con justo título,
ya sea de buena fe o no, de los que la posesión es sin título. En este último supuesto, la fracción II
de este numeral, señala que prescriben en quince años los bienes inmuebles, cuando son poseídos
sin título, pero siempre y cuando dicha posesión sea en concepto de propietario y de manera
pacífica, continua y pública. Por otro lado, el artículo 1159 establece: "Se entiende por justo título el
que es traslativo de dominio.". De la interpretación sistemática de estos artículos se advierte que
cuando se carece de título, no toda posesión es apta para prescribir el bien inmueble, sino sólo
aquella que cumple con los requisitos previstos en el artículo 1157 en cita, pues en la fracción II del
numeral 1158 sólo se liberó el requisito de demostrar únicamente mediante prueba documental tal
circunstancia, pero no de evidenciar que la posesión se tiene en concepto de propietario, esto es,
con pleno dominio del inmueble en cuestión, lo cual debe demostrarse en el juicio, aun con otro
medio probatorio. Por lo tanto, para que prospere una declaración en el sentido de que se adquirió
la posesión en concepto de dueño o de propietario, es menester que se demuestre la causa que le
dio ese carácter, aun cuando sea con medios distintos a la prueba documental, pues sólo la
posesión que se adquiere y disfruta en concepto de dueño de la cosa poseída puede producir la
prescripción, no así la posesión derivada o precaria.

TRIBUNAL COLEGIADO DEL TRIGÉSIMO PRIMER CIRCUITO.


Amparo directo 81/2009. 15 de julio de 2009. Unanimidad de votos. Ponente: Mayra González Solís.
Secretario: Aarón Alberto Pereira Lizama.

Amparo directo 103/2009. Wilberth Pérez Carrillo y otro. 15 de octubre de 2009. Unanimidad de
votos. Ponente: José Rubén Ruiz Ramírez, secretario de tribunal autorizado para desempeñar las
funciones de Magistrado, en términos del artículo 81, fracción XXII, de la Ley Orgánica del Poder
Judicial de la Federación, en relación con el artículo 52, fracción V, del Acuerdo General del Pleno
del Consejo de la Judicatura Federal, que reglamenta la organización y funcionamiento del propio


                                              Pág. 1 de 2              Fecha de impresión 01/10/2025
  https://sjf2.scjn.gob.mx/detalle/tesis/161518
                                                  Semanario Judicial de la Federación

consejo. Secretario: Aarón Alberto Pereira Lizama.

Amparo directo 885/2009. René Leal Botello. 6 de mayo de 2010. Unanimidad de votos. Ponente:
Mayra González Solís. Secretario: Aarón Alberto Pereira Lizama.

Amparo directo 895/2009. Wendy Mariana Concha Uc. 6 de mayo de 2010. Unanimidad de votos.
Ponente: Mayra González Solís. Secretario: Aarón Alberto Pereira Lizama.

Amparo directo 261/2011. María Concepción Segovia Núñez. 25 de mayo de 2011. Unanimidad de
votos. Ponente: David Alberto Barredo Villanueva. Secretario: Carlos David González Vargas.




                                             Pág. 2 de 2       Fecha de impresión 01/10/2025
  https://sjf2.scjn.gob.mx/detalle/tesis/161518

\end{quote}

\subsection{Formalization}

\subsubsection{Clause: tesis161518\_principio\_fundamental}

\textit{Faith exclusivity: Possession is either good faith XOR bad faith (mutually exclusive) Thesis-specific principle: Unregistered contracts cannot establish adverse possession against registered owners}

\textbf{Intermediate Representation:}

\begin{quote}
registered ownership is obligated and unregistered contract is obligated
\end{quote}

\textbf{Witness Clause:}

\begin{lstlisting}[style=witnessstyle, language=Witness]
clause tesis161518_principio_fundamental = oblig(titularidad_registral) AND oblig(contrato_no_inscrito) IMPLIES not(claim(demanda_prescripcion));
\end{lstlisting}

\subsubsection{Clause: tesis161518\_proteccion\_registral}

\textit{Supporting legal logic}

\textbf{Intermediate Representation:}

\begin{quote}
registered ownership is obligated
\end{quote}

\textbf{Witness Clause:}

\begin{lstlisting}[style=witnessstyle, language=Witness]
clause tesis161518_proteccion_registral = oblig(titularidad_registral) IMPLIES claim(oponibilidad_terceros);
\end{lstlisting}

\subsubsection{Clause: tesis161518\_limitacion\_no\_inscripcion}

\textbf{Intermediate Representation:}

\begin{quote}
unregistered contract is obligated
\end{quote}

\textbf{Witness Clause:}

\begin{lstlisting}[style=witnessstyle, language=Witness]
clause tesis161518_limitacion_no_inscripcion = oblig(contrato_no_inscrito) IMPLIES not(oblig(oponibilidad_terceros));
\end{lstlisting}

\subsubsection{Clause: tesis161518\_consecuencia\_logica}

\textbf{Intermediate Representation:}

\begin{quote}
not (opposability to third parties is obligated)
\end{quote}

\textbf{Witness Clause:}

\begin{lstlisting}[style=witnessstyle, language=Witness]
clause tesis161518_consecuencia_logica = not(oblig(oponibilidad_terceros)) IMPLIES not(claim(demanda_prescripcion));
\end{lstlisting}

\newpage

\section{Thesis 162032: PRESCRIPCIÓN POSITIVA. REQUISITOS QUE DEBEN ACREDITARSE PARA SU PROCEDENCIA (LEGISLACIÓN DEL ESTADO DE SONORA).}
\label{thesis:162032}

\subsection{Original Thesis Text}

\begin{quote}
                                                   Semanario Judicial de la Federación

                                                  Tesis

 Registro digital: 162032

 Instancia: Primera Sala          Novena Época                         Materia(s): Civil

 Tesis: 1a./J. 125/2010           Fuente: Semanario Judicial de la  Tipo: Jurisprudencia
                                  Federación y su Gaceta.
                                  Tomo XXXIII, Mayo de 2011, página
                                  101


PRESCRIPCIÓN POSITIVA. REQUISITOS QUE DEBEN ACREDITARSE PARA SU
PROCEDENCIA (LEGISLACIÓN DEL ESTADO DE SONORA).


La prescripción positiva o adquisitiva es un medio de adquirir el dominio mediante la posesión
pacífica, continua, pública, cierta y en concepto de dueño, por el tiempo que establezca la
normatividad aplicable, según se desprende de los artículos 998, 1307, párrafo primero, y 1323 del
Código Civil para el Estado de Sonora. El concepto de dueño no proviene del fuero interno del
poseedor, sino que le es aplicable precisamente a quien entró a poseer la cosa mediante un acto o
hecho que le permite ostentarse como tal, siempre que sea poseedor originario, dado que en el
ordenamiento de referencia, es el único que puede usucapir. Es relevante señalar que la posesión
originaria puede ser justa o de hecho. Por ello, además de que el poseedor deberá probar el tiempo
por el que ininterrumpidamente poseyó (cinco o diez años según el caso, atendiendo al citado
artículo 1323 del Código Civil para el Estado de Sonora), siempre deberá probar la causa
generadora de la posesión. Consecuentemente, si pretende que se declare su adquisición por
usucapión, por haber detentado la cosa durante cinco años en su calidad de poseedor originario,
jurídico y de buena fe, debe exigírsele que demuestre el justo título, en el que basa su pretensión.
Así mismo, si pretende que se declare su adquisición, por haber detentado la cosa durante cinco
años en su calidad de poseedor originario, de hecho y de buena fe, debe exigírsele que pruebe el
hecho generador de la posesión, al igual que si pretende que se declare su adquisición por haber
detentado la cosa durante diez años en su calidad de poseedor originario, de hecho, aunque de
mala fe.


Contradicción de tesis 175/2010. Entre las sustentadas por los Tribunales Colegiados Primero y
Segundo, ambos en Materias Civil y de Trabajo del Quinto Circuito. 17 de noviembre de 2010.
Unanimidad de cuatro votos. Ponente: Juan N. Silva Meza. Secretario: Roberto Ávila Ornelas.

Tesis de jurisprudencia 125/2010. Aprobada por la Primera Sala de este Alto Tribunal, en sesión de
fecha veinticuatro de noviembre de dos mil diez.




                                             Pág. 1 de 2            Fecha de impresión 01/10/2025
  https://sjf2.scjn.gob.mx/detalle/tesis/162032
                                                Semanario Judicial de la Federación




                                           Pág. 2 de 2       Fecha de impresión 01/10/2025
https://sjf2.scjn.gob.mx/detalle/tesis/162032

\end{quote}

\subsection{Formalization}

\subsubsection{Clause: tesis162032\_principio\_fundamental}

\textit{Thesis-specific principle: Possession must be "in concept of owner" (peaceful, continuous, public) for both good and bad faith}

\textbf{Intermediate Representation:}

\begin{quote}
possession as owner is obligated and (oblig(ejercicio\_buena\_fe) OR oblig(ejercicio\_mala\_fe)) IMPLIES oblig(ejercer\_prescripcion\_positiva)
\end{quote}

\textbf{Witness Clause:}

\begin{lstlisting}[style=witnessstyle, language=Witness]
clause tesis162032_principio_fundamental = oblig(posesion_como_propietario) AND (oblig(ejercicio_buena_fe) OR oblig(ejercicio_mala_fe)) IMPLIES oblig(ejercer_prescripcion_positiva);
\end{lstlisting}

\subsubsection{Clause: tesis162032\_posesion\_ambas\_fe}

\textit{Supporting logic: Possession as owner applies to both good faith and bad faith cases}

\textbf{Intermediate Representation:}

\begin{quote}
If possession as owner is obligated, then (good faith exercise is obligated or bad faith exercise is obligated)
\end{quote}

\textbf{Witness Clause:}

\begin{lstlisting}[style=witnessstyle, language=Witness]
clause tesis162032_posesion_ambas_fe = oblig(posesion_como_propietario) IMPLIES (oblig(ejercicio_buena_fe) OR oblig(ejercicio_mala_fe));
\end{lstlisting}

\subsubsection{Clause: tesis162032\_dominio\_economico}

\textit{Economic dominion: Must be dominator of the thing, commanding and enjoying it as owner}

\textbf{Intermediate Representation:}

\begin{quote}
possession as owner is obligated
\end{quote}

\textbf{Witness Clause:}

\begin{lstlisting}[style=witnessstyle, language=Witness]
clause tesis162032_dominio_economico = oblig(posesion_como_propietario) IMPLIES oblig(causa_generadora_probada);
\end{lstlisting}

\subsubsection{Clause: tesis162032\_posesion\_originaria}

\textit{Original possession: Must begin with cause different from derived possession}

\textbf{Intermediate Representation:}

\begin{quote}
possession as owner is obligated and not (temporal possession is obligated)
\end{quote}

\textbf{Witness Clause:}

\begin{lstlisting}[style=witnessstyle, language=Witness]
clause tesis162032_posesion_originaria = oblig(posesion_como_propietario) AND not(oblig(posesion_temporal)) IMPLIES oblig(causa_generadora_probada);
\end{lstlisting}

\newpage

\section{Thesis 162513: PRESCRIPCIÓN ADQUISITIVA. NO SE REQUIERE QUE EL ACTO TRASLATIVO DE DOMINIO EN QUE SE FUNDE SEA DE FECHA CIERTA, SI QUIENES LO SUSCRIBIERON SON PARTE EN}
\label{thesis:162513}

\subsection{Original Thesis Text}

\begin{quote}
                                                   Semanario Judicial de la Federación

                                                  Tesis

 Registro digital: 162513

 Instancia: Tribunales             Novena Época                           Materia(s): Civil
 Colegiados de Circuito

 Tesis: I.7o.C.160 C               Fuente: Semanario Judicial de la       Tipo: Aislada
                                   Federación y su Gaceta.
                                   Tomo XXXIII, Marzo de 2011,
                                   página 2395


PRESCRIPCIÓN ADQUISITIVA. NO SE REQUIERE QUE EL ACTO TRASLATIVO DE DOMINIO
EN QUE SE FUNDE SEA DE FECHA CIERTA, SI QUIENES LO SUSCRIBIERON SON PARTE EN
EL JUICIO RESPECTIVO.


La jurisprudencia 1a./J. 9/2008, emitida por la Primera Sala de la Suprema Corte de Justicia de la
Nación, publicada en el Semanario Judicial de la Federación y su Gaceta, Novena Época, Tomo
XXVII, abril de 2008, página 315, de rubro: "PRESCRIPCIÓN ADQUISITIVA. EL CONTRATO
PRIVADO DE COMPRAVENTA QUE SE EXHIBE PARA ACREDITAR EL JUSTO TÍTULO O LA
CAUSA GENERADORA DE LA POSESIÓN, DEBE SER DE FECHA CIERTA (LEGISLACIÓN DEL
ESTADO DE NUEVO LEÓN).", establece que en tratándose de la acción de prescripción positiva,
para que un contrato traslativo de dominio pueda tener valor probatorio debe ser de fecha cierta; sin
embargo, de la ejecutoria respectiva se desprende que las consideraciones de la Corte en el sentido
de la exigencia del justo título de fecha cierta son para que: 1) Tenga plena eficacia probatoria; y, 2)
Pueda surtir consecuencias contra terceros. Todo ello para proporcionar certidumbre respecto de la
buena fe del acto ahí contenido y evitar actos fraudulentos o dolosos cometidos por las partes en
ese acto jurídico y en perjuicio de terceros. En esas condiciones, el supuesto de la fecha cierta a
que se refiere tal criterio no es aplicable cuando el contrato traslativo de dominio que se exhibe en el
juicio de prescripción adquisitiva para acreditar la causa generadora de la posesión no proviene de
un tercero extraño al juicio, sino de los propios contendientes pues, en esos casos, la fecha y
demás elementos del contrato privado deben considerarse verdaderos, mientras no se demuestre
su falsedad, ya que con el llamamiento a juicio del vendedor que haya suscrito el documento se
permite cumplir la razón subyacente tras la jurisprudencia en cuestión, habida cuenta que el
demandado tendrá la oportunidad de objetar o impugnar cualquier simulación o acto doloso, con lo
que se evitará la comisión de fraudes a través de ese acto jurídico.

SÉPTIMO TRIBUNAL COLEGIADO EN MATERIA CIVIL DEL PRIMER CIRCUITO.


Amparo directo 58/2011. Rocío Hernández Ibarra. 17 de febrero de 2011. Unanimidad de votos.
Ponente: Julio César Vázquez-Mellado García. Secretaria: Alicia Ramírez Ricárdez.

Nota:

Por ejecutoria del 16 de noviembre de 2011, la Primera Sala declaró inexistente la contradicción de
tesis 221/2011, derivada de la denuncia de la que fue objeto el criterio contenido en esta tesis, al
estimarse que no son discrepantes los criterios materia de la denuncia respectiva.


                                              Pág. 1 de 2             Fecha de impresión 01/10/2025
  https://sjf2.scjn.gob.mx/detalle/tesis/162513
                                                  Semanario Judicial de la Federación

Por ejecutoria del 10 de junio de 2015, la Primera Sala declaró inexistente la contradicción de tesis
201/2014, derivada de la denuncia de la que fue objeto el criterio contenido en esta tesis, al
estimarse que no son discrepantes los criterios materia de la denuncia respectiva.

Por ejecutoria del 22 de abril de 2015, la Primera Sala declaró sin materia la contradicción de tesis
251/2014 derivada de la denuncia de la que fue objeto el criterio contenido en esta tesis, al
estimarse que la tesis sobre la cual se genera la contradicción, dejó de ser un criterio obligatorio al
haber sido interrumpido al resolverse una diversa contradicción de tesis.




                                              Pág. 2 de 2             Fecha de impresión 01/10/2025
  https://sjf2.scjn.gob.mx/detalle/tesis/162513

\end{quote}

\subsection{Formalization}

\subsubsection{Clause: tesis162513\_principio\_fundamental}

\textit{Thesis-specific principle: Adverse possession can be invoked by action or exception, but requires same proof elements}

\textbf{Intermediate Representation:}

\begin{quote}
possession as owner is obligated and (good faith exercise is obligated or bad faith exercise is obligated) and (oblig(posesion\_tiempo\_5\_anos) OR oblig(posesion\_tiempo\_10\_anos)) IMPLIES oblig(ejercer\_prescripcion\_positiva)
\end{quote}

\textbf{Witness Clause:}

\begin{lstlisting}[style=witnessstyle, language=Witness]
clause tesis162513_principio_fundamental = oblig(posesion_como_propietario) AND (oblig(ejercicio_buena_fe) OR oblig(ejercicio_mala_fe)) AND (oblig(posesion_tiempo_5_anos) OR oblig(posesion_tiempo_10_anos)) IMPLIES oblig(ejercer_prescripcion_positiva);
\end{lstlisting}

\subsubsection{Clause: tesis162513\_igualdad\_prueba}

\textit{Supporting logic: Both action and exception require same proof elements regardless of procedural route}

\textbf{Intermediate Representation:}

\begin{quote}
ejercer prescripcion positiva is obligated
\end{quote}

\textbf{Witness Clause:}

\begin{lstlisting}[style=witnessstyle, language=Witness]
clause tesis162513_igualdad_prueba = oblig(ejercer_prescripcion_positiva) IMPLIES oblig(posesion_como_propietario);
\end{lstlisting}

\subsubsection{Clause: tesis162513\_plazos\_temporales}

\textit{Time requirements: 5 years for good faith, 10 years for bad faith}

\textbf{Intermediate Representation:}

\begin{quote}
good faith exercise is obligated
\end{quote}

\textbf{Witness Clause:}

\begin{lstlisting}[style=witnessstyle, language=Witness]
clause tesis162513_plazos_temporales = oblig(ejercicio_buena_fe) IMPLIES oblig(posesion_tiempo_5_anos);
\end{lstlisting}

\subsubsection{Clause: tesis162513\_plazos\_mala\_fe}

\textbf{Intermediate Representation:}

\begin{quote}
bad faith exercise is obligated
\end{quote}

\textbf{Witness Clause:}

\begin{lstlisting}[style=witnessstyle, language=Witness]
clause tesis162513_plazos_mala_fe = oblig(ejercicio_mala_fe) IMPLIES oblig(posesion_tiempo_10_anos);
\end{lstlisting}

\subsubsection{Clause: tesis162513\_objetivo\_registral}

\textit{Target requirement: Action must be directed against registered owner}

\textbf{Intermediate Representation:}

\begin{quote}
ejercer prescripcion positiva is obligated
\end{quote}

\textbf{Witness Clause:}

\begin{lstlisting}[style=witnessstyle, language=Witness]
clause tesis162513_objetivo_registral = oblig(ejercer_prescripcion_positiva) IMPLIES oblig(titularidad_registral);
\end{lstlisting}

\newpage

\section{Thesis 164273: PRESCRIPCIÓN ADQUISITIVA. CUANDO SE INVOCA UN CONTRATO VERBAL DE COMPRAVENTA COMO CAUSA GENERADORA DE LA POSESIÓN, CUYA EXISTENCIA}
\label{thesis:164273}

\subsection{Original Thesis Text}

\begin{quote}
                                                   Semanario Judicial de la Federación

                                                  Tesis

 Registro digital: 164273

 Instancia: Tribunales            Novena Época                        Materia(s): Civil
 Colegiados de Circuito

 Tesis: IV.3o.C.40 C              Fuente: Semanario Judicial de la    Tipo: Aislada
                                  Federación y su Gaceta.
                                  Tomo XXXII, Julio de 2010, página
                                  2043


PRESCRIPCIÓN ADQUISITIVA. CUANDO SE INVOCA UN CONTRATO VERBAL DE
COMPRAVENTA COMO CAUSA GENERADORA DE LA POSESIÓN, CUYA EXISTENCIA
PRETENDE ACREDITARSE CON PRUEBA TESTIMONIAL, LOS TESTIGOS DEBEN
MANIFESTAR LA FECHA EXACTA EN QUE AQUÉL SE CELEBRÓ (LEGISLACIÓN DEL ESTADO
DE NUEVO LEÓN).


De una interpretación armónica de los artículos 806 y 1176 del Código Civil para el Estado de
Nuevo León se concluye que para la procedencia de la acción de prescripción adquisitiva es
necesario acreditar la fecha exacta en que se entró a poseer el bien en disputa, ya que es a partir
de entonces cuando empieza a correr el término de cinco años para que opere la prescripción de
buena fe; pues si bien es cierto que el numeral 1173 de la propia codificación, establece que la
prescripción se cuenta por años y no de momento a momento, ello constituye la regla general que
no aplica a la prescripción adquisitiva, porque el mismo precepto en su segunda parte señala: "...
excepto en los casos en que así lo determine la ley expresamente."; y si en este tipo de acción, la
posesión se cuenta a partir de que la persona entra a poseer el bien, es inconcuso que debe
mencionarse la fecha exacta de ese acontecimiento, entendida como tal, el día, mes y año; por
ende, el actor debe precisarla en los hechos de la demanda, ya que éstos son los que se
encuentran sujetos a prueba en términos del diverso artículo 223 del Código de Procedimientos
Civiles del Estado. Por tanto, si se invoca como causa generadora de la posesión un contrato verbal
de compraventa, cuya existencia se pretende acreditar mediante la prueba testimonial, es menester
que los declarantes manifiesten las circunstancias de tiempo (día, mes y año), modo y lugar en que
se concertó ese acuerdo de voluntades.

TERCER TRIBUNAL COLEGIADO EN MATERIA CIVIL DEL CUARTO CIRCUITO.


Amparo directo 498/2009. Isidro Mata Cortez. 24 de marzo de 2010. Unanimidad de votos. Ponente:
Pedro Pablo Hernández Lobato. Secretaria: María Guadalupe Campa Molina.

Amparo directo 452/2009. Florentino Maldonado Salazar, su sucesión. 7 de abril de 2010.
Unanimidad de votos. Ponente: Eduardo Ochoa Torres. Secretaria: Daniela Judith Sáenz Treviño.




                                             Pág. 1 de 2           Fecha de impresión 01/10/2025
  https://sjf2.scjn.gob.mx/detalle/tesis/164273
                                                Semanario Judicial de la Federación




                                           Pág. 2 de 2       Fecha de impresión 01/10/2025
https://sjf2.scjn.gob.mx/detalle/tesis/164273

\end{quote}

\subsection{Formalization}

\subsubsection{Clause: tesis164273\_principio\_fundamental}

\textit{Thesis-specific principle: Adverse possession can be invoked by action or exception, but requires same proof elements}

\textbf{Intermediate Representation:}

\begin{quote}
possession as owner is obligated and (good faith exercise is obligated or bad faith exercise is obligated) and (oblig(posesion\_tiempo\_5\_anos) OR oblig(posesion\_tiempo\_10\_anos)) IMPLIES oblig(ejercer\_prescripcion\_positiva)
\end{quote}

\textbf{Witness Clause:}

\begin{lstlisting}[style=witnessstyle, language=Witness]
clause tesis164273_principio_fundamental = oblig(posesion_como_propietario) AND (oblig(ejercicio_buena_fe) OR oblig(ejercicio_mala_fe)) AND (oblig(posesion_tiempo_5_anos) OR oblig(posesion_tiempo_10_anos)) IMPLIES oblig(ejercer_prescripcion_positiva);
\end{lstlisting}

\subsubsection{Clause: tesis164273\_igualdad\_prueba}

\textit{Supporting logic: Both action and exception require same proof elements regardless of procedural route}

\textbf{Intermediate Representation:}

\begin{quote}
ejercer prescripcion positiva is obligated
\end{quote}

\textbf{Witness Clause:}

\begin{lstlisting}[style=witnessstyle, language=Witness]
clause tesis164273_igualdad_prueba = oblig(ejercer_prescripcion_positiva) IMPLIES oblig(posesion_como_propietario);
\end{lstlisting}

\subsubsection{Clause: tesis164273\_plazos\_temporales}

\textit{Time requirements: 5 years for good faith, 10 years for bad faith}

\textbf{Intermediate Representation:}

\begin{quote}
good faith exercise is obligated
\end{quote}

\textbf{Witness Clause:}

\begin{lstlisting}[style=witnessstyle, language=Witness]
clause tesis164273_plazos_temporales = oblig(ejercicio_buena_fe) IMPLIES oblig(posesion_tiempo_5_anos);
\end{lstlisting}

\subsubsection{Clause: tesis164273\_plazos\_mala\_fe}

\textbf{Intermediate Representation:}

\begin{quote}
bad faith exercise is obligated
\end{quote}

\textbf{Witness Clause:}

\begin{lstlisting}[style=witnessstyle, language=Witness]
clause tesis164273_plazos_mala_fe = oblig(ejercicio_mala_fe) IMPLIES oblig(posesion_tiempo_10_anos);
\end{lstlisting}

\subsubsection{Clause: tesis164273\_objetivo\_registral}

\textit{Target requirement: Action must be directed against registered owner}

\textbf{Intermediate Representation:}

\begin{quote}
ejercer prescripcion positiva is obligated
\end{quote}

\textbf{Witness Clause:}

\begin{lstlisting}[style=witnessstyle, language=Witness]
clause tesis164273_objetivo_registral = oblig(ejercer_prescripcion_positiva) IMPLIES oblig(titularidad_registral);
\end{lstlisting}

\newpage

\section{Thesis 164722: PRESCRIPCIÓN ADQUISITIVA. A LA PREVISTA EN EL CÓDIGO CIVIL DEL ESTADO DE MÉXICO ABROGADO LE ES APLICABLE LA JURISPRUDENCIA 1a./J. 9/2008, AL}
\label{thesis:164722}

\subsection{Original Thesis Text}

\begin{quote}
                                                   Semanario Judicial de la Federación

                                                  Tesis

 Registro digital: 164722

 Instancia: Tribunales            Novena Época                         Materia(s): Civil
 Colegiados de Circuito

 Tesis: II.4o.C. J/2              Fuente: Semanario Judicial de la     Tipo: Jurisprudencia
                                  Federación y su Gaceta.
                                  Tomo XXXI, Abril de 2010, página
                                  2412


PRESCRIPCIÓN ADQUISITIVA. A LA PREVISTA EN EL CÓDIGO CIVIL DEL ESTADO DE
MÉXICO ABROGADO LE ES APLICABLE LA JURISPRUDENCIA 1a./J. 9/2008, AL
INTERPRETAR, SOBRE ESA FIGURA, PRECEPTOS DE LA LEGISLACIÓN DEL ESTADO DE
NUEVO LEÓN DE SIMILAR CONTENIDO A AQUEL ORDENAMIENTO.


La Primera Sala de la Suprema Corte de Justicia de la Nación, en la jurisprudencia número 1a./J.
9/2008, publicada en el Semanario Judicial de la Federación y su Gaceta, Novena Época, Tomo
XXVII, abril de 2008, página 315, de rubro: "PRESCRIPCIÓN ADQUISITIVA. EL CONTRATO
PRIVADO DE COMPRAVENTA QUE SE EXHIBE PARA ACREDITAR EL JUSTO TÍTULO O LA
CAUSA GENERADORA DE LA POSESIÓN, DEBE SER DE FECHA CIERTA (LEGISLACIÓN DEL
ESTADO DE NUEVO LEÓN).", estableció que para que opere la prescripción adquisitiva, es
indispensable que el bien a usucapir, se posea en concepto de propietario, sin que sea suficiente
que se revele la causa generadora de la posesión para tener por acreditado ese requisito, pues es
necesario comprobar el acto jurídico o hecho que justifique ese carácter; por consiguiente, si el
dominio tiene su origen en un instrumento traslativo consistente en un contrato privado, para
acreditar el justo título o la causa generadora de la posesión, es indispensable que sea de fecha
cierta, pues ese dato proporciona certidumbre respecto de la buena fe del acto contenido en el
referido documento y otorga eficacia probatoria a la fecha que consta en él; por lo que la citada
jurisprudencia es aplicable a la prescripción adquisitiva prevista en el código sustantivo del Estado
de México abrogado, al ser interpretada esta figura en los artículos 806, 826, 1136, 1148, 1149,
1151 y 1152 del Código Civil del Estado de Nuevo León, que son de similar contenido a los artículos
781, 801, 911, 912, 914, 915 y 919 del Código Civil del Estado de México abrogado;
consecuentemente, para acreditar el justo título o la causa generadora de la posesión, como
elemento constitutivo de la acción de usucapión, es indispensable que éste sea de fecha cierta.

CUARTO TRIBUNAL COLEGIADO EN MATERIA CIVIL DEL SEGUNDO CIRCUITO.


Amparo directo 845/2008. Luis Manuel Mejía Landa. 29 de enero de 2009. Unanimidad de votos.
Ponente: José Martínez Guzmán. Secretaria: Mónica Saloma Palacios.

Amparo directo 707/2009. Marco Antonio Villegas Granados. 16 de octubre de 2009. Unanimidad de
votos. Ponente: José Martínez Guzmán. Secretario: Francisco Peñaloza Heras.

Amparo directo 860/2009. Enrique Hernández Robles. 16 de octubre de 2009. Unanimidad de votos.
Ponente: Jorge Mario Pardo Rebolledo. Secretaria: Mercedes Verónica Sánchez Miguez.


                                             Pág. 1 de 2             Fecha de impresión 01/10/2025
  https://sjf2.scjn.gob.mx/detalle/tesis/164722
                                                  Semanario Judicial de la Federación

Amparo directo 1009/2009. César de Jesús Urbina. 27 de noviembre de 2009. Unanimidad de
votos. Ponente: Jorge Mario Pardo Rebolledo. Secretaria: Mariana Olúa Aguilar.

Amparo directo 78/2010. **********. 18 de febrero de 2010. Unanimidad de votos. Ponente: Javier
Cardoso Chávez. Secretario: Antonio Salazar López.




                                             Pág. 2 de 2        Fecha de impresión 01/10/2025
  https://sjf2.scjn.gob.mx/detalle/tesis/164722

\end{quote}

\subsection{Formalization}

\subsubsection{Clause: tesis164722\_principio\_fundamental}

\textit{Thesis-specific principle: Just title not required for adverse possession, but must reveal origin of possession and affirm possession as owner}

\textbf{Intermediate Representation:}

\begin{quote}
ejercer prescripcion positiva is obligated and possession as owner is obligated
\end{quote}

\textbf{Witness Clause:}

\begin{lstlisting}[style=witnessstyle, language=Witness]
clause tesis164722_principio_fundamental = oblig(ejercer_prescripcion_positiva) IMPLIES oblig(causa_generadora_revelada) AND oblig(posesion_como_propietario);
\end{lstlisting}

\subsubsection{Clause: tesis164722\_determinacion\_judicial}

\textit{Supporting logic: Judge needs generating fact/act to determine possession quality and timing}

\textbf{Intermediate Representation:}

\begin{quote}
If causa generadora revelada is obligated, then (possession as owner is obligated and (good faith exercise is obligated and not (bad faith exercise is obligated)))
\end{quote}

\textbf{Witness Clause:}

\begin{lstlisting}[style=witnessstyle, language=Witness]
clause tesis164722_determinacion_judicial = oblig(causa_generadora_revelada) IMPLIES (oblig(posesion_como_propietario) AND (oblig(ejercicio_buena_fe) AND not(oblig(ejercicio_mala_fe))));
\end{lstlisting}

\subsubsection{Clause: tesis164722\_tipo\_posesion}

\textit{Possession type determination: Can determine if original or derived possession}

\textbf{Intermediate Representation:}

\begin{quote}
If causa generadora revelada is obligated, then (possession as owner is obligated and not (temporal possession is obligated))
\end{quote}

\textbf{Witness Clause:}

\begin{lstlisting}[style=witnessstyle, language=Witness]
clause tesis164722_tipo_posesion = oblig(causa_generadora_revelada) IMPLIES (oblig(posesion_como_propietario) AND not(oblig(posesion_temporal)));
\end{lstlisting}

\subsubsection{Clause: tesis164722\_determinacion\_temporal}

\textit{Timing requirement: Must establish when prescription period starts counting}

\textbf{Intermediate Representation:}

\begin{quote}
If causa generadora revelada is obligated, then (posesion tiempo 5 anos is obligated or posesion tiempo 10 anos is obligated)
\end{quote}

\textbf{Witness Clause:}

\begin{lstlisting}[style=witnessstyle, language=Witness]
clause tesis164722_determinacion_temporal = oblig(causa_generadora_revelada) IMPLIES (oblig(posesion_tiempo_5_anos) OR oblig(posesion_tiempo_10_anos));
\end{lstlisting}

\newpage

\section{Thesis 166602: PRESCRIPCIÓN POSITIVA. PARA ADQUIRIR UN BIEN INMUEBLE A TRAVÉS DE ESTA FIGURA, SIN NECESIDAD DE TÍTULO, ES MENESTER QUE SE DEMUESTRE QUE SE OBTUVO}
\label{thesis:166602}

\subsection{Original Thesis Text}

\begin{quote}
                                                   Semanario Judicial de la Federación

                                                  Tesis

 Registro digital: 166602

 Instancia: Tribunales             Novena Época                           Materia(s): Civil
 Colegiados de Circuito

 Tesis: XXXI.3 C                   Fuente: Semanario Judicial de la Tipo: Aislada
                                   Federación y su Gaceta.
                                   Tomo XXX, Agosto de 2009, página
                                   1675


PRESCRIPCIÓN POSITIVA. PARA ADQUIRIR UN BIEN INMUEBLE A TRAVÉS DE ESTA
FIGURA, SIN NECESIDAD DE TÍTULO, ES MENESTER QUE SE DEMUESTRE QUE SE OBTUVO
LA POSESIÓN EN CONCEPTO DE DUEÑO O DE PROPIETARIO, Y NO EN FORMA DERIVADA
NI PRECARIA (LEGISLACIÓN DEL ESTADO DE CAMPECHE).


El numeral 1157 del Código Civil del Estado establece: "La posesión necesaria para prescribir debe
ser: I. En concepto de propietario; II. Pacífica; III. Continua; IV. Pública.". Por su parte, el artículo
1158 del mismo ordenamiento legal contempla diversas hipótesis en las que procede la prescripción
de bienes inmuebles, distinguiendo aquellos casos en los que la posesión se ejerce con justo título,
ya sea de buena fe o no, de los que la posesión es sin título. En este último supuesto, la fracción II
del artículo 1158 del Código Civil del Estado de Campeche, señala que prescriben en quince años
los bienes inmuebles, cuando son poseídos sin título, pero siempre y cuando dicha posesión sea en
concepto de propietario y de manera pacífica, continua y pública. Por otro lado, el artículo 1159
establece que "Se entiende por justo título el que es traslativo de dominio.". De la interpretación
sistemática de estos artículos se desprende que cuando se carece de título, no toda posesión es
apta para prescribir el bien inmueble, sino sólo aquella que cumple con los requisitos previstos en el
artículo 1157 en cita, pues en la fracción II del numeral 1158 sólo se liberó el requisito de demostrar
únicamente mediante prueba documental tal circunstancia, pero no de evidenciar que la posesión
se tiene en concepto de propietario, esto es, con pleno dominio del inmueble en cuestión, lo cual
debe demostrarse en el juicio, aun con otro medio probatorio. Por lo tanto, es evidente que para que
prospere una revelación en el sentido de que se adquirió la posesión en concepto de dueño o de
propietario, es menester que se demuestre la causa que le dio ese carácter, aun cuando sea con
medios distintos a la prueba documental, pues sólo la posesión que se adquiere y disfruta en
concepto de dueño de la cosa poseída puede producir la prescripción, no así la posesión derivada o
precaria.

TRIBUNAL COLEGIADO DEL TRIGÉSIMO PRIMER CIRCUITO.


Amparo directo 81/2009. Martha Patricia Garma Blanquet. 15 de julio de 2009. Unanimidad de
votos. Ponente: Mayra González Solís. Secretario: Aarón Alberto Pereira Lizama.




                                              Pág. 1 de 2              Fecha de impresión 01/10/2025
  https://sjf2.scjn.gob.mx/detalle/tesis/166602
                                                Semanario Judicial de la Federación




                                           Pág. 2 de 2       Fecha de impresión 01/10/2025
https://sjf2.scjn.gob.mx/detalle/tesis/166602

\end{quote}

\subsection{Formalization}

\subsubsection{Clause: tesis166602\_principio\_fundamental}

\textit{Thesis-specific principle: Just title not required for adverse possession, but must reveal origin of possession and affirm possession as owner}

\textbf{Intermediate Representation:}

\begin{quote}
ejercer prescripcion positiva is obligated and possession as owner is obligated
\end{quote}

\textbf{Witness Clause:}

\begin{lstlisting}[style=witnessstyle, language=Witness]
clause tesis166602_principio_fundamental = oblig(ejercer_prescripcion_positiva) IMPLIES oblig(causa_generadora_revelada) AND oblig(posesion_como_propietario);
\end{lstlisting}

\subsubsection{Clause: tesis166602\_determinacion\_judicial}

\textit{Supporting logic: Judge needs generating fact/act to determine possession quality and timing}

\textbf{Intermediate Representation:}

\begin{quote}
If causa generadora revelada is obligated, then (possession as owner is obligated and (good faith exercise is obligated and not (bad faith exercise is obligated)))
\end{quote}

\textbf{Witness Clause:}

\begin{lstlisting}[style=witnessstyle, language=Witness]
clause tesis166602_determinacion_judicial = oblig(causa_generadora_revelada) IMPLIES (oblig(posesion_como_propietario) AND (oblig(ejercicio_buena_fe) AND not(oblig(ejercicio_mala_fe))));
\end{lstlisting}

\subsubsection{Clause: tesis166602\_tipo\_posesion}

\textit{Possession type determination: Can determine if original or derived possession}

\textbf{Intermediate Representation:}

\begin{quote}
If causa generadora revelada is obligated, then (possession as owner is obligated and not (temporal possession is obligated))
\end{quote}

\textbf{Witness Clause:}

\begin{lstlisting}[style=witnessstyle, language=Witness]
clause tesis166602_tipo_posesion = oblig(causa_generadora_revelada) IMPLIES (oblig(posesion_como_propietario) AND not(oblig(posesion_temporal)));
\end{lstlisting}

\subsubsection{Clause: tesis166602\_determinacion\_temporal}

\textit{Timing requirement: Must establish when prescription period starts counting}

\textbf{Intermediate Representation:}

\begin{quote}
If causa generadora revelada is obligated, then (posesion tiempo 5 anos is obligated or posesion tiempo 10 anos is obligated)
\end{quote}

\textbf{Witness Clause:}

\begin{lstlisting}[style=witnessstyle, language=Witness]
clause tesis166602_determinacion_temporal = oblig(causa_generadora_revelada) IMPLIES (oblig(posesion_tiempo_5_anos) OR oblig(posesion_tiempo_10_anos));
\end{lstlisting}

\newpage

\section{Thesis 168238: ACCIÓN DE USUCAPIÓN. A LA PREVISTA EN EL CÓDIGO CIVIL DEL ESTADO DE MÉXICO ABROGADO, LE ES APLICABLE LA JURISPRUDENCIA 1a./J. 9/2008, AL ANALIZAR EN LA}
\label{thesis:168238}

\subsection{Original Thesis Text}

\begin{quote}
                                                   Semanario Judicial de la Federación

                                                  Tesis

 Registro digital: 168238

 Instancia: Tribunales            Novena Época                         Materia(s): Civil
 Colegiados de Circuito

 Tesis: II.4o.C.33 C              Fuente: Semanario Judicial de la Tipo: Aislada
                                  Federación y su Gaceta.
                                  Tomo XXIX, Enero de 2009, página
                                  2628


ACCIÓN DE USUCAPIÓN. A LA PREVISTA EN EL CÓDIGO CIVIL DEL ESTADO DE MÉXICO
ABROGADO, LE ES APLICABLE LA JURISPRUDENCIA 1a./J. 9/2008, AL ANALIZAR EN LA
EJECUTORIA QUE LE DIO ORIGEN, PRECEPTOS DE LA LEGISLACIÓN DEL ESTADO DE
NUEVO LEÓN QUE CONCUERDAN CON AQUEL ORDENAMIENTO, SOBRE ESA FIGURA.


A la acción de usucapión o prescripción adquisitiva sustentada en la legislación civil del Estado de
México abrogada, le es aplicable la jurisprudencia 1a./J. 9/2008, publicada en el Semanario Judicial
de la Federación y su Gaceta, Novena Época, Tomo XXVII, abril de 2008, página 315, de rubro:
"PRESCRIPCIÓN ADQUISITIVA. EL CONTRATO PRIVADO DE COMPRAVENTA QUE SE
EXHIBE PARA ACREDITAR EL JUSTO TÍTULO O LA CAUSA GENERADORA DE LA POSESIÓN,
DEBE SER DE FECHA CIERTA (LEGISLACIÓN DEL ESTADO DE NUEVO LEÓN).", porque en
ésta se interpreta lo que sobre el tema establecen los artículos 806, 826 y 1148 del Código Civil del
Estado de Nuevo León, es decir, los elementos referentes a la calidad de la posesión (buena o mala
fe), al título o a su causa generadora y los requisitos que ésta debe reunir para prosperar, tales
como que sea en concepto de propietario, pacífica, pública y continua, los cuales concuerdan con
los numerales 781, 801 y 911 del Código Civil del Estado de México abrogado.

CUARTO TRIBUNAL COLEGIADO EN MATERIA CIVIL DEL SEGUNDO CIRCUITO.


Amparo directo 810/2008. José Samuel Martínez Real y otros. 23 de octubre de 2008. Unanimidad
de votos. Ponente: Jorge Mario Pardo Rebolledo. Secretaria: Mercedes Verónica Sánchez Miguez.

Nota: Por ejecutoria del 2 de diciembre de 2014, el Pleno en Materia Civil del Segundo Circuito
declaró sin materia la contradicción de tesis 4/2014 derivada de la denuncia de la que fue objeto el
criterio contenido en esta tesis, al existir la jurisprudencia 1a./J. 82/2014 (10a.) que resuelve el
mismo problema jurídico.




                                             Pág. 1 de 2            Fecha de impresión 01/10/2025
  https://sjf2.scjn.gob.mx/detalle/tesis/168238
                                                Semanario Judicial de la Federación




                                           Pág. 2 de 2       Fecha de impresión 01/10/2025
https://sjf2.scjn.gob.mx/detalle/tesis/168238

\end{quote}

\subsection{Formalization}

\subsubsection{Clause: tesis168238\_principio\_fundamental}

\textit{Faith exclusivity: Possession is either good faith XOR bad faith (mutually exclusive) Thesis-specific principle: Unregistered contracts cannot establish adverse possession against registered owners}

\textbf{Intermediate Representation:}

\begin{quote}
registered ownership is obligated and unregistered contract is obligated
\end{quote}

\textbf{Witness Clause:}

\begin{lstlisting}[style=witnessstyle, language=Witness]
clause tesis168238_principio_fundamental = oblig(titularidad_registral) AND oblig(contrato_no_inscrito) IMPLIES not(claim(demanda_prescripcion));
\end{lstlisting}

\subsubsection{Clause: tesis168238\_proteccion\_registral}

\textit{Supporting legal logic}

\textbf{Intermediate Representation:}

\begin{quote}
registered ownership is obligated
\end{quote}

\textbf{Witness Clause:}

\begin{lstlisting}[style=witnessstyle, language=Witness]
clause tesis168238_proteccion_registral = oblig(titularidad_registral) IMPLIES claim(oponibilidad_terceros);
\end{lstlisting}

\subsubsection{Clause: tesis168238\_limitacion\_no\_inscripcion}

\textbf{Intermediate Representation:}

\begin{quote}
unregistered contract is obligated
\end{quote}

\textbf{Witness Clause:}

\begin{lstlisting}[style=witnessstyle, language=Witness]
clause tesis168238_limitacion_no_inscripcion = oblig(contrato_no_inscrito) IMPLIES not(oblig(oponibilidad_terceros));
\end{lstlisting}

\subsubsection{Clause: tesis168238\_consecuencia\_logica}

\textbf{Intermediate Representation:}

\begin{quote}
not (opposability to third parties is obligated)
\end{quote}

\textbf{Witness Clause:}

\begin{lstlisting}[style=witnessstyle, language=Witness]
clause tesis168238_consecuencia_logica = not(oblig(oponibilidad_terceros)) IMPLIES not(claim(demanda_prescripcion));
\end{lstlisting}

\newpage

\section{Thesis 168592: PRESCRIPCIÓN ADQUISITIVA EN MATERIA AGRARIA. PARA QUE OPERE ES NECESARIO QUE EL ACTOR REVELE Y DEMUESTRE LA CAUSA GENERADORA DE LA POSESIÓN, QUE}
\label{thesis:168592}

\subsection{Original Thesis Text}

\begin{quote}
                                                   Semanario Judicial de la Federación

                                                  Tesis

 Registro digital: 168592

 Instancia: Tribunales            Novena Época                         Materia(s): Administrativa
 Colegiados de Circuito

 Tesis: XI.2o. J/37               Fuente: Semanario Judicial de la     Tipo: Jurisprudencia
                                  Federación y su Gaceta.
                                  Tomo XXVIII, Octubre de 2008,
                                  página 2270


PRESCRIPCIÓN ADQUISITIVA EN MATERIA AGRARIA. PARA QUE OPERE ES NECESARIO
QUE EL ACTOR REVELE Y DEMUESTRE LA CAUSA GENERADORA DE LA POSESIÓN, QUE
DEBE SER DE NATURALEZA ORIGINARIA Y NO PRECARIA O DERIVADA.


El artículo 48 de la Ley Agraria dispone, entre otras cosas, que quien haya poseído tierras ejidales
"en concepto de titular de derechos de ejidatario", que no sean las destinadas al asentamiento
humano ni se trate de bosques o selvas, de manera pacífica, continua y pública durante un periodo
de cinco años si la posesión es de buena fe, o de diez si fuera de mala fe, adquirirá sobre dichas
tierras los mismos derechos que cualquier ejidatario sobre su parcela. Luego, para que se entienda
satisfecha la posesión en los términos descritos y opere la prescripción adquisitiva en materia
agraria, es necesario que el actor revele y demuestre la causa generadora de aquélla,
entendiéndose por tal, cualquier acto que fundadamente se considere bastante para transferirle el
dominio sobre la unidad de dotación de que se trate, a fin de estar en aptitud de dilucidar
primeramente, si en realidad poseyó "en concepto de titular de derechos de ejidatario", para en
seguida determinar lo referente a los plazos en que habrá de operar la prescripción, ya que la
posesión debe ser de naturaleza originaria y no precaria o derivada.

SEGUNDO TRIBUNAL COLEGIADO DEL DÉCIMO PRIMER CIRCUITO.


Amparo directo 466/95. Amalia Fernández Badajosa. 10 de agosto de 1995. Unanimidad de votos.
Ponente: Juan Díaz Ponce de León. Secretario: Gilberto Díaz Ortiz.

Amparo directo 773/2003. Manuel Vázquez Cuevas. 12 de noviembre de 2003. Unanimidad de
votos. Ponente: Raúl Murillo Delgado. Secretaria: Libertad Rodríguez Verduzco.

Amparo directo 892/2004. Alejandro Mejía Pérez y otro. 26 de enero de 2005. Unanimidad de votos.
Ponente: Hugo Sahuer Hernández. Secretaria: Ma. de la Cruz Estrada Flores.

Amparo directo 1032/2005. Octaviano Duarte Acosta. 19 de abril de 2006. Unanimidad de votos.
Ponente: Raúl Murillo Delgado. Secretario: Pedro Garibay García.

Amparo directo 500/2008. 4 de septiembre de 2008. Unanimidad de votos. Ponente: Raúl Murillo
Delgado. Secretario: Juan Gabriel Sánchez Iriarte.




                                             Pág. 1 de 2             Fecha de impresión 01/10/2025
  https://sjf2.scjn.gob.mx/detalle/tesis/168592
                                                Semanario Judicial de la Federación




                                           Pág. 2 de 2       Fecha de impresión 01/10/2025
https://sjf2.scjn.gob.mx/detalle/tesis/168592

\end{quote}

\subsection{Formalization}

\subsubsection{Clause: tesis168592\_principio\_fundamental}

\textit{Thesis-specific principle: Possession must be "in concept of owner" (peaceful, continuous, public) for both good and bad faith}

\textbf{Intermediate Representation:}

\begin{quote}
possession as owner is obligated and (oblig(ejercicio\_buena\_fe) OR oblig(ejercicio\_mala\_fe)) IMPLIES oblig(ejercer\_prescripcion\_positiva)
\end{quote}

\textbf{Witness Clause:}

\begin{lstlisting}[style=witnessstyle, language=Witness]
clause tesis168592_principio_fundamental = oblig(posesion_como_propietario) AND (oblig(ejercicio_buena_fe) OR oblig(ejercicio_mala_fe)) IMPLIES oblig(ejercer_prescripcion_positiva);
\end{lstlisting}

\subsubsection{Clause: tesis168592\_posesion\_ambas\_fe}

\textit{Supporting logic: Possession as owner applies to both good faith and bad faith cases}

\textbf{Intermediate Representation:}

\begin{quote}
If possession as owner is obligated, then (good faith exercise is obligated or bad faith exercise is obligated)
\end{quote}

\textbf{Witness Clause:}

\begin{lstlisting}[style=witnessstyle, language=Witness]
clause tesis168592_posesion_ambas_fe = oblig(posesion_como_propietario) IMPLIES (oblig(ejercicio_buena_fe) OR oblig(ejercicio_mala_fe));
\end{lstlisting}

\subsubsection{Clause: tesis168592\_dominio\_economico}

\textit{Economic dominion: Must be dominator of the thing, commanding and enjoying it as owner}

\textbf{Intermediate Representation:}

\begin{quote}
possession as owner is obligated
\end{quote}

\textbf{Witness Clause:}

\begin{lstlisting}[style=witnessstyle, language=Witness]
clause tesis168592_dominio_economico = oblig(posesion_como_propietario) IMPLIES oblig(causa_generadora_probada);
\end{lstlisting}

\subsubsection{Clause: tesis168592\_posesion\_originaria}

\textit{Original possession: Must begin with cause different from derived possession}

\textbf{Intermediate Representation:}

\begin{quote}
possession as owner is obligated and not (temporal possession is obligated)
\end{quote}

\textbf{Witness Clause:}

\begin{lstlisting}[style=witnessstyle, language=Witness]
clause tesis168592_posesion_originaria = oblig(posesion_como_propietario) AND not(oblig(posesion_temporal)) IMPLIES oblig(causa_generadora_probada);
\end{lstlisting}

\newpage

\section{Thesis 169022: USUCAPIÓN. EL CONTRATO PRIVADO DE COMPRAVENTA QUE SE EXHIBE PARA ACREDITAR EL JUSTO TÍTULO O LA CAUSA GENERADORA DE LA POSESIÓN, DEBE SER DE}
\label{thesis:169022}

\subsection{Original Thesis Text}

\begin{quote}
                                                   Semanario Judicial de la Federación

                                                  Tesis

 Registro digital: 169022

 Instancia: Tribunales             Novena Época                          Materia(s): Civil
 Colegiados de Circuito

 Tesis: VI.2o.C.619 C              Fuente: Semanario Judicial de la      Tipo: Aislada
                                   Federación y su Gaceta.
                                   Tomo XXVIII, Agosto de 2008,
                                   página 1215


USUCAPIÓN. EL CONTRATO PRIVADO DE COMPRAVENTA QUE SE EXHIBE PARA
ACREDITAR EL JUSTO TÍTULO O LA CAUSA GENERADORA DE LA POSESIÓN, DEBE SER DE
FECHA CIERTA (LEGISLACIÓN DEL ESTADO DE PUEBLA).


Al analizar los elementos de la acción de prescripción adquisitiva, conforme a las reglas contenidas
en los artículos 806, 826, 1136, 1148, 1149, 1151 y 1152 del Código Civil para el Estado de Nuevo
León, la Primera Sala de la Suprema Corte de Justicia de la Nación concluyó que para la
demostración del justo título como un elemento de esa pretensión, es necesario que el contrato
traslativo de dominio en que se sustenta el poder que se ejerce en calidad de dueño sobre el bien
poseído, además de revelar la causa generadora de la posesión, debe gozar de fecha cierta, ya que
ese dato propicia certidumbre respecto de la buena fe del acto jurídico contenido en el documento
en cuestión, evitando así actos fraudulentos o dolosos. Dicho criterio se encuentra plasmado en la
jurisprudencia 1a./J. 9/2008, de rubro: "PRESCRIPCIÓN ADQUISITIVA. EL CONTRATO PRIVADO
DE COMPRAVENTA QUE SE EXHIBE PARA ACREDITAR EL JUSTO TÍTULO O LA CAUSA
GENERADORA DE LA POSESIÓN, DEBE SER DE FECHA CIERTA (LEGISLACIÓN DEL ESTADO
DE NUEVO LEÓN).", publicada en la página 315, Tomo XXVII, abril de 2008, Novena Época del
Semanario Judicial de la Federación y su Gaceta. Pues bien, los preceptos analizados en la
ejecutoria que dio origen al referido criterio jurisprudencial, son semejantes en su contenido a los
artículos 1367, 1372, 1394, 1397, 1401, 1404 y 1407 del Código Civil para el Estado de Puebla y,
por ende, el criterio definido resulta aplicable para los casos juzgados con sustento en la legislación
para esta entidad federativa, de manera que para que un contrato privado de compraventa sea
susceptible de acreditar el justo título o la causa generadora de la posesión en un juicio de
usucapión, es necesario que goce de fecha cierta.

SEGUNDO TRIBUNAL COLEGIADO EN MATERIA CIVIL DEL SEXTO CIRCUITO.


Amparo directo 63/2008. Wulfrano Sánchez Hoyos y otro. 26 de junio de 2008. Unanimidad de
votos. Ponente: Humberto Schettino Reyna, secretario de tribunal en funciones de Magistrado por
ministerio de ley, en términos del artículo 81, fracción XXII, de la Ley Orgánica del Poder Judicial de
la Federación. Secretario: Juan Carlos Cortés Salgado.




                                              Pág. 1 de 2             Fecha de impresión 01/10/2025
  https://sjf2.scjn.gob.mx/detalle/tesis/169022
                                                Semanario Judicial de la Federación




                                           Pág. 2 de 2       Fecha de impresión 01/10/2025
https://sjf2.scjn.gob.mx/detalle/tesis/169022

\end{quote}

\subsection{Formalization}

\subsubsection{Clause: tesis169022\_principio\_fundamental}

\textit{Faith exclusivity: Possession is either good faith XOR bad faith (mutually exclusive) Thesis-specific principle: Unregistered contracts cannot establish adverse possession against registered owners}

\textbf{Intermediate Representation:}

\begin{quote}
registered ownership is obligated and unregistered contract is obligated
\end{quote}

\textbf{Witness Clause:}

\begin{lstlisting}[style=witnessstyle, language=Witness]
clause tesis169022_principio_fundamental = oblig(titularidad_registral) AND oblig(contrato_no_inscrito) IMPLIES not(claim(demanda_prescripcion));
\end{lstlisting}

\subsubsection{Clause: tesis169022\_proteccion\_registral}

\textit{Supporting legal logic}

\textbf{Intermediate Representation:}

\begin{quote}
registered ownership is obligated
\end{quote}

\textbf{Witness Clause:}

\begin{lstlisting}[style=witnessstyle, language=Witness]
clause tesis169022_proteccion_registral = oblig(titularidad_registral) IMPLIES claim(oponibilidad_terceros);
\end{lstlisting}

\subsubsection{Clause: tesis169022\_limitacion\_no\_inscripcion}

\textbf{Intermediate Representation:}

\begin{quote}
unregistered contract is obligated
\end{quote}

\textbf{Witness Clause:}

\begin{lstlisting}[style=witnessstyle, language=Witness]
clause tesis169022_limitacion_no_inscripcion = oblig(contrato_no_inscrito) IMPLIES not(oblig(oponibilidad_terceros));
\end{lstlisting}

\subsubsection{Clause: tesis169022\_consecuencia\_logica}

\textbf{Intermediate Representation:}

\begin{quote}
not (opposability to third parties is obligated)
\end{quote}

\textbf{Witness Clause:}

\begin{lstlisting}[style=witnessstyle, language=Witness]
clause tesis169022_consecuencia_logica = not(oblig(oponibilidad_terceros)) IMPLIES not(claim(demanda_prescripcion));
\end{lstlisting}

\newpage

\section{Thesis 169244: PRESCRIPCIÓN ADQUISITIVA EN MATERIA AGRARIA. EL CERTIFICADO DE DERECHOS EXPEDIDO POR AUTORIDAD COMPETENTE ES APTO PARA ACREDITAR LA POSESIÓN EN}
\label{thesis:169244}

\subsection{Original Thesis Text}

\begin{quote}
                                                   Semanario Judicial de la Federación

                                                  Tesis

 Registro digital: 169244

 Instancia: Tribunales             Novena Época                          Materia(s): Administrativa
 Colegiados de Circuito

 Tesis: XXII.2o.17 A               Fuente: Semanario Judicial de la   Tipo: Aislada
                                   Federación y su Gaceta.
                                   Tomo XXVIII, Julio de 2008, página
                                   1830


PRESCRIPCIÓN ADQUISITIVA EN MATERIA AGRARIA. EL CERTIFICADO DE DERECHOS
EXPEDIDO POR AUTORIDAD COMPETENTE ES APTO PARA ACREDITAR LA POSESIÓN EN
CONCEPTO DE TITULAR DE DERECHOS DE EJIDATARIO, AUN CUANDO CON
POSTERIORIDAD A LA CONSUMACIÓN DE AQUÉLLA HAYA SIDO DECLARADO NULO POR
VICIOS DEL PROCEDIMIENTO EN QUE SE SUSTENTÓ SU EXPEDICIÓN.


Si al hacer valer la prescripción adquisitiva se invoca como causa generadora de la posesión un
certificado de derechos expedido por la autoridad competente, dicho documento es apto para
acreditar que la posesión fue adquirida en concepto de titular de derechos de ejidatario, tal como lo
exige el artículo 48 de la Ley Agraria. Lo anterior, independientemente de que con posterioridad a la
consumación de la prescripción, el referido título haya sido declarado nulo por la existencia de vicios
en el procedimiento en que se sustentó su expedición, ya que esa circunstancia no afecta el hecho
de que mientras transcurría el plazo para usucapir, la posesión fue ejercida de buena fe, al amparo
de un acto que generó la creencia fundada de haber obtenido el título suficiente para poseer en
concepto de titular de derechos de ejidatario.

SEGUNDO TRIBUNAL COLEGIADO DEL VIGÉSIMO SEGUNDO CIRCUITO.


Amparo directo 953/2005. Manuel Martínez López. 3 de octubre de 2006. Unanimidad de votos.
Ponente: Germán Tena Campero. Secretario: Carlos Ernesto Farías Flores.




                                              Pág. 1 de 1             Fecha de impresión 01/10/2025
  https://sjf2.scjn.gob.mx/detalle/tesis/169244

\end{quote}

\subsection{Formalization}

\subsubsection{Clause: tesis169244\_principio\_fundamental}

\textit{Thesis-specific principle: Good faith requires sufficient title/generating cause, ignorance of title defects, and founded belief of ownership}

\textbf{Intermediate Representation:}

\begin{quote}
good faith exercise is obligated and not (conocimiento vicios is obligated)
\end{quote}

\textbf{Witness Clause:}

\begin{lstlisting}[style=witnessstyle, language=Witness]
clause tesis169244_principio_fundamental = oblig(ejercicio_buena_fe) IMPLIES oblig(causa_generadora_probada) AND not(oblig(conocimiento_vicios));
\end{lstlisting}

\subsubsection{Clause: tesis169244\_mala\_fe}

\textit{Supporting logic: Bad faith possessor has no title/generating cause OR knows title defects that prevent rightful possession}

\textbf{Intermediate Representation:}

\begin{quote}
If bad faith exercise is obligated, then (not (proven generating cause is obligated) or conocimiento vicios is obligated)
\end{quote}

\textbf{Witness Clause:}

\begin{lstlisting}[style=witnessstyle, language=Witness]
clause tesis169244_mala_fe = oblig(ejercicio_mala_fe) IMPLIES (not(oblig(causa_generadora_probada)) OR oblig(conocimiento_vicios));
\end{lstlisting}

\subsubsection{Clause: tesis169244\_requisitos\_buena\_fe}

\textit{Good faith requirements: Must have sufficient title and ignorance of defects}

\textbf{Intermediate Representation:}

\begin{quote}
good faith exercise is obligated
\end{quote}

\textbf{Witness Clause:}

\begin{lstlisting}[style=witnessstyle, language=Witness]
clause tesis169244_requisitos_buena_fe = oblig(ejercicio_buena_fe) IMPLIES oblig(causa_generadora_probada);
\end{lstlisting}

\subsubsection{Clause: tesis169244\_condiciones\_mala\_fe}

\textit{Bad faith conditions: No title or knowledge of defects}

\textbf{Intermediate Representation:}

\begin{quote}
(not(oblig(causa\_generadora\_probada)) OR oblig(conocimiento\_vicios)) IMPLIES oblig(ejercicio\_mala\_fe)
\end{quote}

\textbf{Witness Clause:}

\begin{lstlisting}[style=witnessstyle, language=Witness]
clause tesis169244_condiciones_mala_fe = (not(oblig(causa_generadora_probada)) OR oblig(conocimiento_vicios)) IMPLIES oblig(ejercicio_mala_fe);
\end{lstlisting}

\newpage

\section{Thesis 169650: PRESCRIPCIÓN POSITIVA EN MATERIA AGRARIA. PARA QUE PROCEDA, LA CAUSA GENERADORA DE LA POSESIÓN PUEDE ACREDITARSE MEDIANTE UN CONTRATO}
\label{thesis:169650}

\subsection{Original Thesis Text}

\begin{quote}
                                                   Semanario Judicial de la Federación

                                                  Tesis

 Registro digital: 169650

 Instancia: Tribunales              Novena Época                           Materia(s): Administrativa
 Colegiados de Circuito

 Tesis: III.1o.A.144 A              Fuente: Semanario Judicial de la Tipo: Aislada
                                    Federación y su Gaceta.
                                    Tomo XXVII, Mayo de 2008, página
                                    1115


PRESCRIPCIÓN POSITIVA EN MATERIA AGRARIA. PARA QUE PROCEDA, LA CAUSA
GENERADORA DE LA POSESIÓN PUEDE ACREDITARSE MEDIANTE UN CONTRATO
PRIVADO DE CESIÓN DE DERECHOS CELEBRADO BAJO LA VIGENCIA DE LA DEROGADA
LEY FEDERAL DE REFORMA AGRARIA.


El artículo 48 de la Ley Agraria dispone, entre otras cosas, que quien haya poseído tierras ejidales
"en concepto de titular de derechos de ejidatario", que no sean las destinadas al asentamiento
humano ni se trate de bosques o selvas, de manera pacífica, continua y pública durante un periodo
de cinco años si la posesión es de buena fe o de diez si fuera de mala fe, adquirirá sobre ellas los
mismos derechos que cualquier ejidatario sobre su parcela, precisando así los requisitos para que
proceda la prescripción positiva de derechos agrarios. Luego, para que se entienda satisfecha la
aludida posesión en las condiciones apuntadas, es menester que se demuestre su causa
generadora, no obstante que ésta se acredite mediante un contrato privado de cesión de derechos
celebrado bajo la vigencia de la derogada Ley Federal de Reforma Agraria, que prohibía este tipo
de actos, pues tal documento surte efectos a partir de que comenzó a regir el ordenamiento
primeramente citado, aunado a que no se requiere de justo título para que proceda la mencionada
prescripción y, por tanto, tal criterio no implica la aplicación retroactiva de la ley en perjuicio de los
ejidatarios.

PRIMER TRIBUNAL COLEGIADO EN MATERIA ADMINISTRATIVA DEL TERCER CIRCUITO.


Amparo directo 216/2007. Javier Almaraz Martínez. 6 de febrero de 2008. Unanimidad de votos.
Ponente: Rogelio Camarena Cortés. Secretario: José Vega Cortez.

Nota: Por ejecutoria del 12 de julio de 2017, la Segunda Sala declaró inexistente la contradicción de
tesis 45/2017 derivada de la denuncia de la que fue objeto el criterio contenido en esta tesis, al
estimarse que no son discrepantes los criterios materia de la denuncia respectiva.




                                               Pág. 1 de 2              Fecha de impresión 01/10/2025
  https://sjf2.scjn.gob.mx/detalle/tesis/169650
                                                Semanario Judicial de la Federación




                                           Pág. 2 de 2       Fecha de impresión 01/10/2025
https://sjf2.scjn.gob.mx/detalle/tesis/169650

\end{quote}

\subsection{Formalization}

\subsubsection{Clause: tesis169650\_principio\_fundamental}

\textit{Thesis-specific principle: Just title not required for adverse possession, but must reveal origin of possession and affirm possession as owner}

\textbf{Intermediate Representation:}

\begin{quote}
ejercer prescripcion positiva is obligated and possession as owner is obligated
\end{quote}

\textbf{Witness Clause:}

\begin{lstlisting}[style=witnessstyle, language=Witness]
clause tesis169650_principio_fundamental = oblig(ejercer_prescripcion_positiva) IMPLIES oblig(causa_generadora_revelada) AND oblig(posesion_como_propietario);
\end{lstlisting}

\subsubsection{Clause: tesis169650\_determinacion\_judicial}

\textit{Supporting logic: Judge needs generating fact/act to determine possession quality and timing}

\textbf{Intermediate Representation:}

\begin{quote}
If causa generadora revelada is obligated, then (possession as owner is obligated and (good faith exercise is obligated and not (bad faith exercise is obligated)))
\end{quote}

\textbf{Witness Clause:}

\begin{lstlisting}[style=witnessstyle, language=Witness]
clause tesis169650_determinacion_judicial = oblig(causa_generadora_revelada) IMPLIES (oblig(posesion_como_propietario) AND (oblig(ejercicio_buena_fe) AND not(oblig(ejercicio_mala_fe))));
\end{lstlisting}

\subsubsection{Clause: tesis169650\_tipo\_posesion}

\textit{Possession type determination: Can determine if original or derived possession}

\textbf{Intermediate Representation:}

\begin{quote}
If causa generadora revelada is obligated, then (possession as owner is obligated and not (temporal possession is obligated))
\end{quote}

\textbf{Witness Clause:}

\begin{lstlisting}[style=witnessstyle, language=Witness]
clause tesis169650_tipo_posesion = oblig(causa_generadora_revelada) IMPLIES (oblig(posesion_como_propietario) AND not(oblig(posesion_temporal)));
\end{lstlisting}

\subsubsection{Clause: tesis169650\_determinacion\_temporal}

\textit{Timing requirement: Must establish when prescription period starts counting}

\textbf{Intermediate Representation:}

\begin{quote}
If causa generadora revelada is obligated, then (posesion tiempo 5 anos is obligated or posesion tiempo 10 anos is obligated)
\end{quote}

\textbf{Witness Clause:}

\begin{lstlisting}[style=witnessstyle, language=Witness]
clause tesis169650_determinacion_temporal = oblig(causa_generadora_revelada) IMPLIES (oblig(posesion_tiempo_5_anos) OR oblig(posesion_tiempo_10_anos));
\end{lstlisting}

\newpage

\section{Thesis 169830: PRESCRIPCIÓN ADQUISITIVA. EL CONTRATO PRIVADO DE COMPRAVENTA QUE SE EXHIBE PARA ACREDITAR EL JUSTO TÍTULO O LA CAUSA GENERADORA DE LA POSESIÓN, DEBE}
\label{thesis:169830}

\subsection{Original Thesis Text}

\begin{quote}
                                                   Semanario Judicial de la Federación

                                                  Tesis

 Registro digital: 169830

 Instancia: Primera Sala          Novena Época                          Materia(s): Civil

 Tesis: 1a./J. 9/2008             Fuente: Semanario Judicial de la      Tipo: Jurisprudencia
                                  Federación y su Gaceta.
                                  Tomo XXVII, Abril de 2008, página
                                  315


PRESCRIPCIÓN ADQUISITIVA. EL CONTRATO PRIVADO DE COMPRAVENTA QUE SE EXHIBE
PARA ACREDITAR EL JUSTO TÍTULO O LA CAUSA GENERADORA DE LA POSESIÓN, DEBE
SER DE FECHA CIERTA (LEGISLACIÓN DEL ESTADO DE NUEVO LEÓN).


De los artículos 806, 826, 1136, 1148, 1149, 1151 y 1152 del Código Civil del Estado de Nuevo
León se advierte que son poseedores de buena fe tanto el que entra en la posesión en virtud de un
título suficiente para darle derecho de poseer como quien ignora los vicios de su título que le
impiden poseer con derecho; que la posesión necesaria para prescribir debe ser en concepto de
propietario y con justo título, pacífica, continua y pública; y que sólo la posesión que se adquiere y
disfruta en concepto de dueño de la cosa poseída puede producir la prescripción. De manera que si
para que opere la prescripción adquisitiva es indispensable que el bien a usucapir se posea en
concepto de propietario, no basta con revelar la causa generadora de la posesión para tener por
acreditado ese requisito, sino que es necesario comprobar el acto jurídico o hecho que justifique ese
carácter, esto es, el justo título, entendiéndose por tal el que es o fundadamente se cree bastante
para transferir el dominio. Ahora bien, los documentos privados adquieren certeza de su contenido a
partir del día en que se inscriben en un registro público de la propiedad, se presentan ante un
fedatario público o muere alguno de los firmantes, pues si no se actualiza uno de esos supuestos no
puede otorgarse valor probatorio frente a terceros. Así, se concluye que si el dominio tiene su origen
en un instrumento traslativo consistente en un contrato privado de compraventa, para acreditar el
justo título o la causa generadora de la posesión es indispensable que sea de fecha cierta, pues ese
dato proporciona certidumbre respecto de la buena fe del acto contenido en el referido documento y
otorga eficacia probatoria a la fecha que consta en él, para evitar actos fraudulentos o dolosos, ya
que la exhibición del contrato tiene como finalidad la acreditación del derecho que le asiste a una
persona y que la legitima para promover un juicio de usucapión; de ahí que la autoridad debe contar
con elementos de convicción idóneos para fijar la calidad de la posesión y computar su término.


Contradicción de tesis 27/2007-PS. Entre las sustentadas por los Tribunales Colegiados Segundo y
Tercero, ambos en Materia Civil del Cuarto Circuito. 9 de enero de 2008. Mayoría de cuatro votos.
Disidente: Juan N. Silva Meza. Ponente: Olga Sánchez Cordero de García Villegas. Secretaria:
Beatriz Joaquina Jaimes Ramos.

Tesis de jurisprudencia 9/2008. Aprobada por la Primera Sala de este Alto Tribunal, en sesión
privada de dieciséis de enero de dos mil ocho.

Nota: Este criterio fue interrumpido por la tesis 1a./J. 82/2014 (10a.), de título y subtítulo:
"PRESCRIPCIÓN ADQUISITIVA. AUNQUE LA LEGISLACIÓN APLICABLE NO EXIJA QUE EL

                                             Pág. 1 de 2             Fecha de impresión 01/10/2025
  https://sjf2.scjn.gob.mx/detalle/tesis/169830
                                                  Semanario Judicial de la Federación

JUSTO TÍTULO O ACTO TRASLATIVO DE DOMINIO QUE CONSTITUYE LA CAUSA
GENERADORA DE LA POSESIÓN DE BUENA FE, SEA DE FECHA CIERTA, LA CERTEZA DE LA
FECHA DEL ACTO JURÍDICO DEBE PROBARSE EN FORMA FEHACIENTE POR SER UN
ELEMENTO DEL JUSTO TÍTULO (INTERRUPCIÓN DE LA JURISPRUDENCIA 1a./J. 9/2008).",
publicada el viernes 5 de diciembre de 2014, a las 10:05 horas en el Semanario Judicial de la
Federación y en la Gaceta del Semanario Judicial de la Federación, Décima Época, Libro 13, Tomo
I, diciembre de 2014, página 200.




                                             Pág. 2 de 2        Fecha de impresión 01/10/2025
  https://sjf2.scjn.gob.mx/detalle/tesis/169830

\end{quote}

\subsection{Formalization}

\subsubsection{Clause: tesis169830\_principio\_fundamental}

\textit{Faith exclusivity: Possession is either good faith XOR bad faith (mutually exclusive) Thesis-specific principle: Unregistered contracts cannot establish adverse possession against registered owners}

\textbf{Intermediate Representation:}

\begin{quote}
registered ownership is obligated and unregistered contract is obligated
\end{quote}

\textbf{Witness Clause:}

\begin{lstlisting}[style=witnessstyle, language=Witness]
clause tesis169830_principio_fundamental = oblig(titularidad_registral) AND oblig(contrato_no_inscrito) IMPLIES not(claim(demanda_prescripcion));
\end{lstlisting}

\subsubsection{Clause: tesis169830\_proteccion\_registral}

\textit{Supporting legal logic}

\textbf{Intermediate Representation:}

\begin{quote}
registered ownership is obligated
\end{quote}

\textbf{Witness Clause:}

\begin{lstlisting}[style=witnessstyle, language=Witness]
clause tesis169830_proteccion_registral = oblig(titularidad_registral) IMPLIES claim(oponibilidad_terceros);
\end{lstlisting}

\subsubsection{Clause: tesis169830\_limitacion\_no\_inscripcion}

\textbf{Intermediate Representation:}

\begin{quote}
unregistered contract is obligated
\end{quote}

\textbf{Witness Clause:}

\begin{lstlisting}[style=witnessstyle, language=Witness]
clause tesis169830_limitacion_no_inscripcion = oblig(contrato_no_inscrito) IMPLIES not(oblig(oponibilidad_terceros));
\end{lstlisting}

\subsubsection{Clause: tesis169830\_consecuencia\_logica}

\textbf{Intermediate Representation:}

\begin{quote}
not (opposability to third parties is obligated)
\end{quote}

\textbf{Witness Clause:}

\begin{lstlisting}[style=witnessstyle, language=Witness]
clause tesis169830_consecuencia_logica = not(oblig(oponibilidad_terceros)) IMPLIES not(claim(demanda_prescripcion));
\end{lstlisting}

\newpage

\section{Thesis 172709: PRESCRIPCIÓN ADQUISITIVA DE INMUEBLES CUYA POSESIÓN SEA POR MÁS DE VEINTE AÑOS, PARA QUE PROCEDA, NO ES NECESARIO ACREDITAR UN JUSTO TÍTULO NI LA}
\label{thesis:172709}

\subsection{Original Thesis Text}

\begin{quote}
                                                   Semanario Judicial de la Federación

                                                  Tesis

 Registro digital: 172709

 Instancia: Primera Sala          Novena Época                          Materia(s): Civil

 Tesis: 1a./J. 19/2007            Fuente: Semanario Judicial de la      Tipo: Jurisprudencia
                                  Federación y su Gaceta.
                                  Tomo XXV, Abril de 2007, página
                                  312


PRESCRIPCIÓN ADQUISITIVA DE INMUEBLES CUYA POSESIÓN SEA POR MÁS DE VEINTE
AÑOS, PARA QUE PROCEDA, NO ES NECESARIO ACREDITAR UN JUSTO TÍTULO NI LA
CAUSA GENERADORA DE LA POSESIÓN (LEGISLACIÓN DEL ESTADO DE GUANAJUATO).


De los artículos 1037, 1039, 1044, 1054, 1055, 1074, 1246, 1248, 1250 y 1251 del Código Civil para
el Estado de Guanajuato, se desprende que la regla general para que opere la prescripción de
bienes inmuebles en la entidad, se requiere que la posesión sea civil, pacífica, continua y pública, y
con base en estos elementos, prescribirán, en cinco años cuando se posean con justo título y con
buena fe y en diez años cuando se posean con justo título y de mala fe, y toda vez que la posesión
civil es aquella que se tiene a título de propietario, debe acreditarse la causa generadora o el justo
título que legitime esa posesión. Sin embargo, esta regla no aplica en su totalidad, tratándose de
aquellos inmuebles cuya posesión sea por más de veinte años, porque en estos casos, el Código
Civil Estatal, no exige que deba demostrarse o acreditarse un justo título ni la causa generadora de
la posesión; por lo que para usucapir este tipo de bienes, basta con que la posesión se ejerza a
nombre propio, como dueño, esto es, que no se posea precariamente o a nombre de otro, mediante
la exteriorización del dominio sobre el inmueble a través de la ejecución de actos que revelen que el
poseedor es el dominador, el que manda en él y lo disfruta para sí, como dueño en sentido
económico, para hacer suya la propiedad desde el punto de vista de los hechos, para tener por
acreditado que dicho bien se posee de manera civil, sin que sea menester acreditar un justo título ni
la causa generadora de la posesión, además de demostrarse que ha sido pacífica, continua y
pública.


Contradicción de tesis 115/2005-PS. Entre las sustentadas por el Primer Tribunal Colegiado en
Materia Civil, Tribunal Colegiado en Materia Penal y Segundo Tribunal Colegiado en Materia Civil,
todos del Décimo Sexto Circuito (antes Segundo, Tercero y Cuarto Tribunales Colegiados del citado
circuito). 18 de octubre de 2006. Cinco votos. Ponente: Juan N. Silva Meza. Secretario: Martín
Adolfo Santos Pérez.

Tesis de jurisprudencia 19/2007. Aprobada por la Primera Sala de este Alto Tribunal, en sesión de
fecha treinta y uno de enero de dos mil siete.

Nota:

Esta tesis fue objeto de la denuncia relativa a la contradicción de tesis 183/2020 en la Suprema
Corte de Justicia de la Nación, desechada por notoriamente improcedente por lo que hace a la
denuncia de los criterios sostenidos por la Primera y Segunda Salas mediante acuerdo de

                                             Pág. 1 de 2             Fecha de impresión 01/10/2025
  https://sjf2.scjn.gob.mx/detalle/tesis/172709
                                                  Semanario Judicial de la Federación

presidencia de 22 de septiembre de 2020 y la admitió a trámite para resolver por el Pleno de la
Suprema Corte de Justicia de la Nación, por lo que hace a los criterios sustentados por los
Tribunales Colegiados de Circuito en Materias Civil y Administrativa.

Por ejecutoria del 10 de febrero de 2021, la Segunda Sala declaró inexistente la contradicción de
tesis 183/2020, derivada de la denuncia de la que fue objeto el criterio contenido en esta tesis.




                                             Pág. 2 de 2          Fecha de impresión 01/10/2025
  https://sjf2.scjn.gob.mx/detalle/tesis/172709

\end{quote}

\subsection{Formalization}

\subsubsection{Clause: tesis172709\_principio\_fundamental}

\textit{Thesis-specific principle: When defendant admits possession without title for more than 10 years with required qualities, is bad faith possessor and property is eligible for prescription}

\textbf{Intermediate Representation:}

\begin{quote}
If (posesion\_como\_propietario AND posesion\_pacifico\_continuo\_publico AND posesion\_tiempo\_10\_anos AND not(causa\_generadora\_probada)), then (ejercicio\_mala\_fe AND oblig(ejercer\_prescripcion\_positiva))
\end{quote}

\textbf{Witness Clause:}

\begin{lstlisting}[style=witnessstyle, language=Witness]
clause tesis172709_principio_fundamental = (posesion_como_propietario AND posesion_pacifico_continuo_publico AND posesion_tiempo_10_anos AND not(causa_generadora_probada)) IMPLIES (ejercicio_mala_fe AND oblig(ejercer_prescripcion_positiva));
\end{lstlisting}

\subsubsection{Clause: tesis172709\_no\_prueba\_causa}

\textit{Supporting logic: Does not need to prove generating cause of possession in trial when no title exists}

\textbf{Intermediate Representation:}

\begin{quote}
(not(causa\_generadora\_probada) AND posesion\_tiempo\_10\_anos) IMPLIES oblig(ejercer\_prescripcion\_positiva)
\end{quote}

\textbf{Witness Clause:}

\begin{lstlisting}[style=witnessstyle, language=Witness]
clause tesis172709_no_prueba_causa = (not(causa_generadora_probada) AND posesion_tiempo_10_anos) IMPLIES oblig(ejercer_prescripcion_positiva);
\end{lstlisting}

\subsubsection{Clause: tesis172709\_definicion\_titulo}

\textit{Title definition: In state law, title is understood as the generating cause of possession}

\textbf{Intermediate Representation:}

\begin{quote}
not (causa\_generadora\_probada) IMPLIES not(titularidad\_registral)
\end{quote}

\textbf{Witness Clause:}

\begin{lstlisting}[style=witnessstyle, language=Witness]
clause tesis172709_definicion_titulo = not(causa_generadora_probada) IMPLIES not(titularidad_registral);
\end{lstlisting}

\subsubsection{Clause: tesis172709\_apto\_prescripcion}

\textit{Prescription eligibility: If 10 years have passed and other elements are met, property is eligible for prescription}

\textbf{Intermediate Representation:}

\begin{quote}
(posesion\_tiempo\_10\_anos AND posesion\_como\_propietario AND posesion\_pacifico\_continuo\_publico) IMPLIES oblig(ejercer\_prescripcion\_positiva)
\end{quote}

\textbf{Witness Clause:}

\begin{lstlisting}[style=witnessstyle, language=Witness]
clause tesis172709_apto_prescripcion = (posesion_tiempo_10_anos AND posesion_como_propietario AND posesion_pacifico_continuo_publico) IMPLIES oblig(ejercer_prescripcion_positiva);
\end{lstlisting}

\newpage

\section{Thesis 172955: PRESCRIPCIÓN ADQUISITIVA. ES IMPROCEDENTE LA ACCIÓN EJERCITADA POR EL COMPRADOR DE UNA COSA CIERTA Y DETERMINADA CONTRA SU VENDEDOR, SI LO QUE}
\label{thesis:172955}

\subsection{Original Thesis Text}

\begin{quote}
                                                   Semanario Judicial de la Federación

                                                  Tesis

 Registro digital: 172955

 Instancia: Tribunales            Novena Época                          Materia(s): Civil
 Colegiados de Circuito

 Tesis: XVI.2o.C.28 C             Fuente: Semanario Judicial de la      Tipo: Aislada
                                  Federación y su Gaceta.
                                  Tomo XXV, Marzo de 2007, página
                                  1746


PRESCRIPCIÓN ADQUISITIVA. ES IMPROCEDENTE LA ACCIÓN EJERCITADA POR EL
COMPRADOR DE UNA COSA CIERTA Y DETERMINADA CONTRA SU VENDEDOR, SI LO QUE
PRETENDE ES LA OBTENCIÓN DE LA ESCRITURA PÚBLICA QUE AMPARE SU DERECHO DE
PROPIEDAD (LEGISLACIÓN DEL ESTADO DE GUANAJUATO).


La prescripción positiva es una acción que tiene como contenido la pretensión de adquirir el derecho
de propiedad sobre un bien que se ha poseído a título de dueño, pacífica, continuamente y por un
tiempo determinado por la circunstancia de si existe buena o mala fe y si se cuenta con justo título o
no; por su parte, la acción pro forma comprende la pretensión del comprador de que le sea otorgada
escritura pública para acreditar frente a todo mundo que es propietario del bien. En ese orden de
ideas, si de conformidad con en el artículo 1742 del Código Civil del Estado de Guanajuato, la
compraventa es un contrato consensual que cuando versa sobre cosas ciertas y determinadas es
válido desde que las partes se ponen de acuerdo en precio y cosa, aun cuando no se haya
entregado el bien ni satisfecho el precio; entonces es claro que dándose ese supuesto, el
comprador adquiere la propiedad de la cosa desde el momento en que se verifica el asentimiento de
voluntades, a pesar de que la cosa no le sea entregada o que no se haya cubierto el precio; por ello,
la naturaleza de la acción de prescripción adquisitiva es improcedente cuando el comprador de una
cosa cierta y determinada la ejercita contra su vendedor, pues la naturaleza de aquella acción es
incompatible con la situación jurídica creada por los hechos, ya que conforme a su postura, el
comprador ya cuenta con la propiedad del bien y en realidad lo que persigue es la obtención de la
escritura pública que ampare aquel derecho; reclamación que debe intentarse al través de la acción
personal pro forma y no de la acción de prescripción adquisitiva, pues no puede alcanzarse a través
de ésta lo que ya se tiene, es decir, nadie puede prescribir contra su propio título.

SEGUNDO TRIBUNAL COLEGIADO EN MATERIA CIVIL DEL DÉCIMO SEXTO CIRCUITO.


Amparo directo 1006/2006. José Ignacio Moreno, su sucesión. 19 de enero de 2006. Unanimidad de
votos. Ponente: Juan Manuel Arredondo Elías. Secretario: Salvador Álvarez Villanueva.

Nota:

Por ejecutoria de fecha 4 de junio de 2008, la Primera Sala declaró inexistente la contradicción de
tesis 157/2007-PS en que participó el presente criterio.

Por ejecutoria de fecha 27 de mayo de 2009, la Segunda Sala declaró inexistente la contradicción


                                             Pág. 1 de 2             Fecha de impresión 01/10/2025
  https://sjf2.scjn.gob.mx/detalle/tesis/172955
                                                  Semanario Judicial de la Federación

de tesis 141/2009 en que participó el presente criterio.




                                              Pág. 2 de 2      Fecha de impresión 01/10/2025
  https://sjf2.scjn.gob.mx/detalle/tesis/172955

\end{quote}

\subsection{Formalization}

\subsubsection{Clause: tesis172955\_principio\_fundamental}

\textit{Thesis-specific principle: Good faith possessors require just title as generating cause of possession for adverse possession to operate}

\textbf{Intermediate Representation:}

\begin{quote}
good faith exercise is obligated and possession as owner is obligated
\end{quote}

\textbf{Witness Clause:}

\begin{lstlisting}[style=witnessstyle, language=Witness]
clause tesis172955_principio_fundamental = oblig(ejercicio_buena_fe) AND oblig(posesion_como_propietario) IMPLIES oblig(causa_generadora_probada);
\end{lstlisting}

\subsubsection{Clause: tesis172955\_revelar\_causa}

\textit{Supporting logic: Must reveal and fully prove the generating cause of possession}

\textbf{Intermediate Representation:}

\begin{quote}
ejercer prescripcion positiva is obligated
\end{quote}

\textbf{Witness Clause:}

\begin{lstlisting}[style=witnessstyle, language=Witness]
clause tesis172955_revelar_causa = oblig(ejercer_prescripcion_positiva) IMPLIES oblig(causa_generadora_revelada);
\end{lstlisting}

\subsubsection{Clause: tesis172955\_determinacion\_judicial}

\textit{Judicial determination: Judge needs generating fact/act to determine possession quality and timing}

\textbf{Intermediate Representation:}

\begin{quote}
If causa generadora revelada is obligated, then (possession as owner is obligated and (good faith exercise is obligated and not (bad faith exercise is obligated)))
\end{quote}

\textbf{Witness Clause:}

\begin{lstlisting}[style=witnessstyle, language=Witness]
clause tesis172955_determinacion_judicial = oblig(causa_generadora_revelada) IMPLIES (oblig(posesion_como_propietario) AND (oblig(ejercicio_buena_fe) AND not(oblig(ejercicio_mala_fe))));
\end{lstlisting}

\subsubsection{Clause: tesis172955\_determinacion\_temporal}

\textit{Timing requirement: Must establish when prescription period begins}

\textbf{Intermediate Representation:}

\begin{quote}
If causa generadora revelada is obligated, then (posesion tiempo 5 anos is obligated or posesion tiempo 10 anos is obligated)
\end{quote}

\textbf{Witness Clause:}

\begin{lstlisting}[style=witnessstyle, language=Witness]
clause tesis172955_determinacion_temporal = oblig(causa_generadora_revelada) IMPLIES (oblig(posesion_tiempo_5_anos) OR oblig(posesion_tiempo_10_anos));
\end{lstlisting}

\newpage

\section{Thesis 174862: PRESCRIPCIÓN POSITIVA. EL ADMINISTRADOR ÚNICO DE UNA EMPRESA NO PUEDE PRESCRIBIR PARA SÍ, EL INMUEBLE QUE DETENTA SU REPRESENTADA,}
\label{thesis:174862}

\subsection{Original Thesis Text}

\begin{quote}
                                                   Semanario Judicial de la Federación

                                                  Tesis

 Registro digital: 174862

 Instancia: Tribunales            Novena Época                          Materia(s): Civil
 Colegiados de Circuito

 Tesis: I.3o.C.549 C              Fuente: Semanario Judicial de la      Tipo: Aislada
                                  Federación y su Gaceta.
                                  Tomo XXIII, Junio de 2006, página
                                  1199


PRESCRIPCIÓN POSITIVA. EL ADMINISTRADOR ÚNICO DE UNA EMPRESA NO PUEDE
PRESCRIBIR PARA SÍ, EL INMUEBLE QUE DETENTA SU REPRESENTADA,
INDEPENDIENTEMENTE DE CUÁL SEA LA CAUSA GENERADORA DE ESA POSESIÓN.


De conformidad con lo dispuesto en los artículos 1135 y 1136 del Código Civil para el Distrito
Federal, la figura jurídica de la prescripción es un medio de adquirir bienes o de librarse de
obligaciones, mediante el transcurso de cierto tiempo y bajo las condiciones establecidas en la ley.
La adquisición de bienes en virtud de la posesión, se denomina prescripción positiva; en tanto que
la liberación de obligaciones, por no exigirse su cumplimiento, se llama prescripción negativa. La
institución de la prescripción positiva, también denominada como usucapión, consiste en adquirir la
propiedad de un bien a través de una posesión calificada, por el término que determine la ley, de tal
forma que esa institución no actúa instantáneamente, sino que es el resultado de la permanencia en
un estado posesorio, el que con el transcurso del tiempo genera derechos que son regulados y
protegidos por la legislación. Al respecto, el artículo 1151 del Código Civil para el Distrito Federal
dispone que la posesión apta para prescribir debe ejercerse en concepto de propietario, en forma
pacífica, continua, pública y por el tiempo que determina la ley; lo que pone de manifiesto que no
toda posesión es apta para prescribir, pues para que la acción relativa prospere es menester que se
goce de la posesión originaria y no de la derivada; es decir, debe poseerse a nombre propio y no en
nombre de otro, pues la operancia de la prescripción adquisitiva excluye los conceptos que por su
definición y naturaleza no revisten el ánimo de poseer para sí y, por consecuencia "en concepto de
propietario". En ese tenor, dadas las funciones que desarrolla el administrador único de una
persona moral, e incluso un apoderado, que consisten, generalmente, en representar a la sociedad
y realizar todas las operaciones inherentes a su objeto social, de conformidad con lo dispuesto en el
artículo 10 de la Ley General de Sociedades Mercantiles, no se encuentran facultados para
prescribir a su favor, como personas físicas, determinado bien que se encuentre en posesión de su
representada, sin importar cuál sea el origen de esa detentación (posesión originaria o derivada),
pues no pueden desconocer el vínculo jurídico que los une a ella, habida cuenta que la posesión
que en todo caso tuvieren, con motivo del cargo que ostentan, no puede considerarse apta para
prescribir de buena, ni de mala fe, ya que esa posesión no sería originaria, sino derivada de la
posesión de su representada, por lo que no se configuraría uno de los elementos fundamentales
para la actualización de la figura jurídica de mérito, dado que no se poseería el bien "en concepto
de propietario".

TERCER TRIBUNAL COLEGIADO EN MATERIA CIVIL DEL PRIMER CIRCUITO.




                                             Pág. 1 de 2             Fecha de impresión 01/10/2025
  https://sjf2.scjn.gob.mx/detalle/tesis/174862
                                                  Semanario Judicial de la Federación

Amparo directo 854/2005. Román Sandoval Belmont y/o Román Sandoval XX. 16 de marzo de
2006. Unanimidad de votos. Ponente: Raúl Alfaro Telpalo, secretario de tribunal autorizado por el
Consejo de la Judicatura Federal para desempeñar las funciones de Magistrado. Secretario: Gabriel
Regis López.




                                             Pág. 2 de 2          Fecha de impresión 01/10/2025
  https://sjf2.scjn.gob.mx/detalle/tesis/174862

\end{quote}

\subsection{Formalization}

\subsubsection{Clause: tesis174862\_principio\_fundamental}

\textit{Faith exclusivity: Possession is either good faith XOR bad faith (mutually exclusive) Thesis-specific principle: Unregistered contracts cannot establish adverse possession against registered owners}

\textbf{Intermediate Representation:}

\begin{quote}
registered ownership is obligated and unregistered contract is obligated
\end{quote}

\textbf{Witness Clause:}

\begin{lstlisting}[style=witnessstyle, language=Witness]
clause tesis174862_principio_fundamental = oblig(titularidad_registral) AND oblig(contrato_no_inscrito) IMPLIES not(claim(demanda_prescripcion));
\end{lstlisting}

\subsubsection{Clause: tesis174862\_proteccion\_registral}

\textit{Supporting legal logic}

\textbf{Intermediate Representation:}

\begin{quote}
registered ownership is obligated
\end{quote}

\textbf{Witness Clause:}

\begin{lstlisting}[style=witnessstyle, language=Witness]
clause tesis174862_proteccion_registral = oblig(titularidad_registral) IMPLIES claim(oponibilidad_terceros);
\end{lstlisting}

\subsubsection{Clause: tesis174862\_limitacion\_no\_inscripcion}

\textbf{Intermediate Representation:}

\begin{quote}
unregistered contract is obligated
\end{quote}

\textbf{Witness Clause:}

\begin{lstlisting}[style=witnessstyle, language=Witness]
clause tesis174862_limitacion_no_inscripcion = oblig(contrato_no_inscrito) IMPLIES not(oblig(oponibilidad_terceros));
\end{lstlisting}

\subsubsection{Clause: tesis174862\_consecuencia\_logica}

\textbf{Intermediate Representation:}

\begin{quote}
not (opposability to third parties is obligated)
\end{quote}

\textbf{Witness Clause:}

\begin{lstlisting}[style=witnessstyle, language=Witness]
clause tesis174862_consecuencia_logica = not(oblig(oponibilidad_terceros)) IMPLIES not(claim(demanda_prescripcion));
\end{lstlisting}

\newpage

\section{Thesis 175034: PRESCRIPCIÓN POSITIVA EN MATERIA AGRARIA. PARA QUE PROCEDA DEBE INTENTARSE SOBRE TIERRAS FORMALMENTE PARCELADAS.}
\label{thesis:175034}

\subsection{Original Thesis Text}

\begin{quote}
                                                   Semanario Judicial de la Federación

                                                  Tesis

 Registro digital: 175034

 Instancia: Tribunales            Novena Época                          Materia(s): Administrativa
 Colegiados de Circuito

 Tesis: II.1o.A.117 A             Fuente: Semanario Judicial de la      Tipo: Aislada
                                  Federación y su Gaceta.
                                  Tomo XXIII, Mayo de 2006, página
                                  1838


PRESCRIPCIÓN POSITIVA EN MATERIA AGRARIA. PARA QUE PROCEDA DEBE INTENTARSE
SOBRE TIERRAS FORMALMENTE PARCELADAS.


Para la procedencia de la prescripción positiva en materia agraria, el artículo 48 de la ley de la
materia prevé cuatro requisitos de carácter subjetivo: que se posean tierras ejidales; en concepto de
titular de derechos de ejidatario; de manera pacífica, continua y pública, y durante cinco años, si la
posesión es de buena fe, o de diez si fuera de mala fe. Además, tal precepto establece un requisito
de carácter objetivo, que alude a las cualidades que debe cubrir el predio objeto de la prescripción,
y que consiste en que no se trate de tierras destinadas al asentamiento humano, ni de bosques o
selvas, lo que deriva de la protección especial que caracteriza a tales bienes del ejido, como se
advierte de los artículos 27 constitucional; 64 y 74 de la Ley Agraria. Así, aunque el citado artículo
48 no excluye de la prescripción positiva a las tierras de uso común que no sean bosques o selvas,
de una interpretación teleológica y sistemática de éste y de los numerales 23 y 56 de la referida ley,
se concluye que para que un predio ejidal pueda reclamarse por la vía de la prescripción positiva,
debe tratarse de tierras formalmente parceladas, ya que sólo en tal supuesto podría prosperar la
demanda de prescripción positiva, puesto que de lo contrario, correspondería a la asamblea del
núcleo de población decidir el destino del inmueble.

PRIMER TRIBUNAL COLEGIADO EN MATERIA ADMINISTRATIVA DEL SEGUNDO CIRCUITO.


Amparo directo 549/2005. Francisco Carmona Neri. 23 de febrero de 2006. Mayoría de votos.
Disidente: Angelina Hernández Hernández. Ponente: Salvador Mondragón Reyes. Secretaria: Sonia
Rojas Castro.

Nota: Por ejecutoria del 22 de enero de 2025, la Segunda Sala declaró improcedente
la contradicción de criterios 238/2024, derivada de la denuncia de la que fue objeto el criterio
contenido en esta tesis, en virtud de que el criterio sostenido por el Primer Tribunal Colegiado en
Materia Administrativa del Segundo Circuito, en el amparo directo 549/2005, fue superado ante lo
resuelto por el Pleno en Materia Administrativa del Segundo Circuito, al resolver la contradicción de
tesis 13/2019, que dio origen a la tesis de jurisprudencia PC.II.A. J/1 A (11a.), de rubro:
"PRESCRIPCIÓN ADQUISITIVA EN MATERIA AGRARIA. LA CAUSA GENERADORA DE LA
POSESIÓN NO PUEDE ACREDITARSE VÁLIDAMENTE EN RELACIÓN CON UNA PARCELA NO
ASIGNADA POR LA ASAMBLEA DE EJIDATARIOS."




                                             Pág. 1 de 2             Fecha de impresión 01/10/2025
  https://sjf2.scjn.gob.mx/detalle/tesis/175034
                                                Semanario Judicial de la Federación




                                           Pág. 2 de 2       Fecha de impresión 01/10/2025
https://sjf2.scjn.gob.mx/detalle/tesis/175034

\end{quote}

\subsection{Formalization}

\subsubsection{Clause: tesis175034\_principio\_fundamental}

\textit{Thesis-specific principle: Just title not required for adverse possession, but must reveal origin of possession and affirm possession as owner}

\textbf{Intermediate Representation:}

\begin{quote}
ejercer prescripcion positiva is obligated and possession as owner is obligated
\end{quote}

\textbf{Witness Clause:}

\begin{lstlisting}[style=witnessstyle, language=Witness]
clause tesis175034_principio_fundamental = oblig(ejercer_prescripcion_positiva) IMPLIES oblig(causa_generadora_revelada) AND oblig(posesion_como_propietario);
\end{lstlisting}

\subsubsection{Clause: tesis175034\_determinacion\_judicial}

\textit{Supporting logic: Judge needs generating fact/act to determine possession quality and timing}

\textbf{Intermediate Representation:}

\begin{quote}
If causa generadora revelada is obligated, then (possession as owner is obligated and (good faith exercise is obligated and not (bad faith exercise is obligated)))
\end{quote}

\textbf{Witness Clause:}

\begin{lstlisting}[style=witnessstyle, language=Witness]
clause tesis175034_determinacion_judicial = oblig(causa_generadora_revelada) IMPLIES (oblig(posesion_como_propietario) AND (oblig(ejercicio_buena_fe) AND not(oblig(ejercicio_mala_fe))));
\end{lstlisting}

\subsubsection{Clause: tesis175034\_tipo\_posesion}

\textit{Possession type determination: Can determine if original or derived possession}

\textbf{Intermediate Representation:}

\begin{quote}
If causa generadora revelada is obligated, then (possession as owner is obligated and not (temporal possession is obligated))
\end{quote}

\textbf{Witness Clause:}

\begin{lstlisting}[style=witnessstyle, language=Witness]
clause tesis175034_tipo_posesion = oblig(causa_generadora_revelada) IMPLIES (oblig(posesion_como_propietario) AND not(oblig(posesion_temporal)));
\end{lstlisting}

\subsubsection{Clause: tesis175034\_determinacion\_temporal}

\textit{Timing requirement: Must establish when prescription period starts counting}

\textbf{Intermediate Representation:}

\begin{quote}
If causa generadora revelada is obligated, then (posesion tiempo 5 anos is obligated or posesion tiempo 10 anos is obligated)
\end{quote}

\textbf{Witness Clause:}

\begin{lstlisting}[style=witnessstyle, language=Witness]
clause tesis175034_determinacion_temporal = oblig(causa_generadora_revelada) IMPLIES (oblig(posesion_tiempo_5_anos) OR oblig(posesion_tiempo_10_anos));
\end{lstlisting}

\newpage

\section{Thesis 175851: PRESCRIPCIÓN POSITIVA. SI LA ACCIÓN SE EJERCE CON BASE EN LA POSESIÓN DE BUENA FE, EL JUZGADOR SE ENCUENTRA IMPEDIDO PARA ANALIZAR DE OFICIO LA}
\label{thesis:175851}

\subsection{Original Thesis Text}

\begin{quote}
                                                   Semanario Judicial de la Federación

                                                  Tesis

 Registro digital: 175851

 Instancia: Primera Sala          Novena Época                          Materia(s): Civil

 Tesis: 1a./J. 200/2005           Fuente: Semanario Judicial de la      Tipo: Jurisprudencia
                                  Federación y su Gaceta.
                                  Tomo XXIII, Febrero de 2006,
                                  página 441


PRESCRIPCIÓN POSITIVA. SI LA ACCIÓN SE EJERCE CON BASE EN LA POSESIÓN DE
BUENA FE, EL JUZGADOR SE ENCUENTRA IMPEDIDO PARA ANALIZAR DE OFICIO LA
POSESIÓN DE MALA FE.


Si se atiende al principio de congruencia en las sentencias, conforme al cual el juzgador solamente
debe atender a las acciones y excepciones hechas valer por las partes en el juicio, sin introducir
cuestiones ajenas al debate, se concluye que cuando la prescripción se ejerce con base en una
posesión de buena fe, el Juez no puede analizar de oficio la existencia de una posesión de mala fe,
ya que ésta no fue planteada en la demanda, porque de lo contrario se dejaría en estado de
indefensión al demandado, en tanto que su defensa se endereza contra lo abordado en aquélla, por
lo que si la parte actora al hacer valer su acción de prescripción aduce ser poseedor de buena fe, en
caso de no acreditarse la posesión en esos términos, el juzgador está impedido para analizar si la
que ostenta el actor es de mala fe, pues ello no forma parte de la litis planteada.


Contradicción de tesis 18/2005-PS. Entre las sustentadas por el Noveno Tribunal Colegiado en
Materia Civil del Primer Circuito y el Tercer Tribunal Colegiado en Materia Civil del Tercer Circuito.
23 de noviembre de 2005. Cinco votos. Ponente: Olga Sánchez Cordero de García Villegas.
Secretaria: Rosaura Rivera Salcedo.

Tesis de jurisprudencia 200/2005. Aprobada por la Primera Sala de este Alto Tribunal, en sesión de
fecha treinta de noviembre de dos mil cinco.




                                             Pág. 1 de 1             Fecha de impresión 01/10/2025
  https://sjf2.scjn.gob.mx/detalle/tesis/175851

\end{quote}

\subsection{Formalization}

\subsubsection{Clause: tesis175851\_principio\_fundamental}

\textit{Thesis-specific principle: Good faith possessors require just title as generating cause of possession for adverse possession to operate}

\textbf{Intermediate Representation:}

\begin{quote}
good faith exercise is obligated and possession as owner is obligated
\end{quote}

\textbf{Witness Clause:}

\begin{lstlisting}[style=witnessstyle, language=Witness]
clause tesis175851_principio_fundamental = oblig(ejercicio_buena_fe) AND oblig(posesion_como_propietario) IMPLIES oblig(causa_generadora_probada);
\end{lstlisting}

\subsubsection{Clause: tesis175851\_revelar\_causa}

\textit{Supporting logic: Must reveal and fully prove the generating cause of possession}

\textbf{Intermediate Representation:}

\begin{quote}
ejercer prescripcion positiva is obligated
\end{quote}

\textbf{Witness Clause:}

\begin{lstlisting}[style=witnessstyle, language=Witness]
clause tesis175851_revelar_causa = oblig(ejercer_prescripcion_positiva) IMPLIES oblig(causa_generadora_revelada);
\end{lstlisting}

\subsubsection{Clause: tesis175851\_determinacion\_judicial}

\textit{Judicial determination: Judge needs generating fact/act to determine possession quality and timing}

\textbf{Intermediate Representation:}

\begin{quote}
If causa generadora revelada is obligated, then (possession as owner is obligated and (good faith exercise is obligated and not (bad faith exercise is obligated)))
\end{quote}

\textbf{Witness Clause:}

\begin{lstlisting}[style=witnessstyle, language=Witness]
clause tesis175851_determinacion_judicial = oblig(causa_generadora_revelada) IMPLIES (oblig(posesion_como_propietario) AND (oblig(ejercicio_buena_fe) AND not(oblig(ejercicio_mala_fe))));
\end{lstlisting}

\subsubsection{Clause: tesis175851\_determinacion\_temporal}

\textit{Timing requirement: Must establish when prescription period begins}

\textbf{Intermediate Representation:}

\begin{quote}
If causa generadora revelada is obligated, then (posesion tiempo 5 anos is obligated or posesion tiempo 10 anos is obligated)
\end{quote}

\textbf{Witness Clause:}

\begin{lstlisting}[style=witnessstyle, language=Witness]
clause tesis175851_determinacion_temporal = oblig(causa_generadora_revelada) IMPLIES (oblig(posesion_tiempo_5_anos) OR oblig(posesion_tiempo_10_anos));
\end{lstlisting}

\newpage

\section{Thesis 179504: PRESCRIPCIÓN ADQUISITIVA EN MATERIA AGRARIA. PARA SU PROCEDENCIA NO SE REQUIERE DE JUSTO TÍTULO.}
\label{thesis:179504}

\subsection{Original Thesis Text}

\begin{quote}
                                                   Semanario Judicial de la Federación

                                                  Tesis

 Registro digital: 179504

 Instancia: Segunda Sala          Novena Época                         Materia(s): Administrativa

 Tesis: 2a./J. 207/2004           Fuente: Semanario Judicial de la     Tipo: Jurisprudencia
                                  Federación y su Gaceta.
                                  Tomo XXI, Enero de 2005, página
                                  575


PRESCRIPCIÓN ADQUISITIVA EN MATERIA AGRARIA. PARA SU PROCEDENCIA NO SE
REQUIERE DE JUSTO TÍTULO.


Conforme al artículo 48 de la Ley Agraria, la prescripción adquisitiva de parcelas ejidales depende
de la actualización de ciertas condiciones ineludibles, consistentes en que la posesión se haya
ejercido por quien pueda adquirir la titularidad de derechos de ejidatario; que se haya ejercido
respecto de tierras ejidales que no sean de las destinadas al asentamiento humano, ni se trate de
bosques o selvas; que la posesión haya sido pacífica, pública y continua por los plazos que señala
la ley; y que se ejerza en concepto de titular de derechos de ejidatario. Como se advierte, la ley no
exige un "justo título" para poseer, entendido éste como el que es, o el que fundadamente se cree,
bastante para transferir derechos agrarios sobre la parcela o parte de ella, pues aun el poseedor de
mala fe, que es el que entra en posesión sin título alguno o el que conoce los vicios de su título,
puede adquirir la titularidad de tales derechos por prescripción. Sin embargo, lo anterior no implica
que no deba acreditarse la causa generadora de la posesión, pues ello es necesario para conocer la
calidad con la que se ejerce, es decir, si es originaria o derivada, si es de buena o mala fe, y la
fecha cierta a partir de la cual ha de computarse el término legal de la prescripción.


Contradicción de tesis 96/2004-SS. Entre las sustentadas por los Tribunales Colegiados Segundo y
Tercero, ambos en Materia Administrativa del Sexto Circuito y Primer Tribunal Colegiado en Materia
Administrativa del Segundo Circuito. 1o. de diciembre de 2004. Cinco votos. Ponente: Margarita
Beatriz Luna Ramos. Secretaria: Constanza Tort San Román.

Tesis de jurisprudencia 207/2004. Aprobada por la Segunda Sala de este Alto Tribunal, en sesión
privada del ocho de diciembre de dos mil cuatro.

Nota:

Esta tesis fue objeto de la denuncia relativa a la contradicción de tesis 183/2020 en la Suprema
Corte de Justicia de la Nación, desechada por notoriamente improcedente por lo que hace a la
denuncia de los criterios sostenidos por la Primera y Segunda Salas mediante acuerdo de
presidencia de 22 de septiembre de 2020 y la admitió a trámite para resolver por el Pleno de la
Suprema Corte de Justicia de la Nación, por lo que hace a los criterios sustentados por los
Tribunales Colegiados de Circuito en Materias Civil y Administrativa.

Por ejecutoria del 10 de febrero de 2021, la Segunda Sala declaró inexistente la contradicción de
tesis 183/2020, derivada de la denuncia de la que fue objeto el criterio contenido en esta tesis.

                                             Pág. 1 de 2             Fecha de impresión 01/10/2025
  https://sjf2.scjn.gob.mx/detalle/tesis/179504
                                                Semanario Judicial de la Federación




                                           Pág. 2 de 2       Fecha de impresión 01/10/2025
https://sjf2.scjn.gob.mx/detalle/tesis/179504

\end{quote}

\subsection{Formalization}

\subsubsection{Clause: tesis179504\_principio\_fundamental}

\textit{Thesis-specific principle: Possession must be "in concept of owner" (peaceful, continuous, public) for both good and bad faith}

\textbf{Intermediate Representation:}

\begin{quote}
possession as owner is obligated and (oblig(ejercicio\_buena\_fe) OR oblig(ejercicio\_mala\_fe)) IMPLIES oblig(ejercer\_prescripcion\_positiva)
\end{quote}

\textbf{Witness Clause:}

\begin{lstlisting}[style=witnessstyle, language=Witness]
clause tesis179504_principio_fundamental = oblig(posesion_como_propietario) AND (oblig(ejercicio_buena_fe) OR oblig(ejercicio_mala_fe)) IMPLIES oblig(ejercer_prescripcion_positiva);
\end{lstlisting}

\subsubsection{Clause: tesis179504\_posesion\_ambas\_fe}

\textit{Supporting logic: Possession as owner applies to both good faith and bad faith cases}

\textbf{Intermediate Representation:}

\begin{quote}
If possession as owner is obligated, then (good faith exercise is obligated or bad faith exercise is obligated)
\end{quote}

\textbf{Witness Clause:}

\begin{lstlisting}[style=witnessstyle, language=Witness]
clause tesis179504_posesion_ambas_fe = oblig(posesion_como_propietario) IMPLIES (oblig(ejercicio_buena_fe) OR oblig(ejercicio_mala_fe));
\end{lstlisting}

\subsubsection{Clause: tesis179504\_dominio\_economico}

\textit{Economic dominion: Must be dominator of the thing, commanding and enjoying it as owner}

\textbf{Intermediate Representation:}

\begin{quote}
possession as owner is obligated
\end{quote}

\textbf{Witness Clause:}

\begin{lstlisting}[style=witnessstyle, language=Witness]
clause tesis179504_dominio_economico = oblig(posesion_como_propietario) IMPLIES oblig(causa_generadora_probada);
\end{lstlisting}

\subsubsection{Clause: tesis179504\_posesion\_originaria}

\textit{Original possession: Must begin with cause different from derived possession}

\textbf{Intermediate Representation:}

\begin{quote}
possession as owner is obligated and not (temporal possession is obligated)
\end{quote}

\textbf{Witness Clause:}

\begin{lstlisting}[style=witnessstyle, language=Witness]
clause tesis179504_posesion_originaria = oblig(posesion_como_propietario) AND not(oblig(posesion_temporal)) IMPLIES oblig(causa_generadora_probada);
\end{lstlisting}

\newpage

\section{Thesis 180040: USUCAPIÓN, CAUSA GENERADORA DE LA POSESIÓN. NO LA CONSTITUYE UN CONTRATO PRIVADO DE COMPRAVENTA RESCINDIDO PREVIAMENTE.}
\label{thesis:180040}

\subsection{Original Thesis Text}

\begin{quote}
                                                   Semanario Judicial de la Federación

                                                  Tesis

 Registro digital: 180040

 Instancia: Tribunales             Novena Época                          Materia(s): Civil
 Colegiados de Circuito

 Tesis: II.2o.C.480 C              Fuente: Semanario Judicial de la      Tipo: Aislada
                                   Federación y su Gaceta.
                                   Tomo XX, Noviembre de 2004,
                                   página 2044


USUCAPIÓN, CAUSA GENERADORA DE LA POSESIÓN. NO LA CONSTITUYE UN CONTRATO
PRIVADO DE COMPRAVENTA RESCINDIDO PREVIAMENTE.


De reconocido derecho resulta que en la acción de usucapión no basta con revelar la causa
generadora de la posesión para que la misma prospere, sino que aquélla debe corroborarse en el
juicio natural, ante lo cual, es necesario que el promovente cuente con un justo título eficaz y
subsistente al intentarla, con el fin de que pueda demostrar que posee como dueño y de buena fe.
Por consiguiente, si el contrato de compraventa con el cual se pretende dicha usucapión fue
rescindido previamente al ejercicio de esa acción, mediante sentencia que quedó firme, no puede ya
considerarse como causa generadora de la posesión, al haber dejado de tener efectos en el orden
jurídico, de ahí que no pueda estimarse a dicha compraventa como título eficiente, por carecer de
trascendencia, formal y legalmente, dado que en definitiva se rescindió con antelación la
compraventa que permitiría a la posesión producir la usucapión o prescripción adquisitiva.

SEGUNDO TRIBUNAL COLEGIADO EN MATERIA CIVIL DEL SEGUNDO CIRCUITO.


Amparo directo 493/2004. José Vázquez Márquez. 31 de agosto de 2004. Unanimidad de votos.
Ponente: Virgilio A. Solorio Campos. Secretaria: Aimeé Michelle Delgado Martínez.

Nota: Esta tesis fue objeto de la denuncia relativa a la contradicción de criterios 378/2023 del índice
de la Suprema Corte de Justicia de la Nación, desechada por notoriamente improcedente, mediante
acuerdo de presidencia de 4 de diciembre de 2023.




                                              Pág. 1 de 2             Fecha de impresión 01/10/2025
  https://sjf2.scjn.gob.mx/detalle/tesis/180040
                                                Semanario Judicial de la Federación




                                           Pág. 2 de 2       Fecha de impresión 01/10/2025
https://sjf2.scjn.gob.mx/detalle/tesis/180040

\end{quote}

\subsection{Formalization}

\subsubsection{Clause: tesis180040\_principio\_fundamental}

\textit{Thesis-specific principle: Possession must be "in concept of owner" (peaceful, continuous, public) for both good and bad faith}

\textbf{Intermediate Representation:}

\begin{quote}
possession as owner is obligated and (oblig(ejercicio\_buena\_fe) OR oblig(ejercicio\_mala\_fe)) IMPLIES oblig(ejercer\_prescripcion\_positiva)
\end{quote}

\textbf{Witness Clause:}

\begin{lstlisting}[style=witnessstyle, language=Witness]
clause tesis180040_principio_fundamental = oblig(posesion_como_propietario) AND (oblig(ejercicio_buena_fe) OR oblig(ejercicio_mala_fe)) IMPLIES oblig(ejercer_prescripcion_positiva);
\end{lstlisting}

\subsubsection{Clause: tesis180040\_posesion\_ambas\_fe}

\textit{Supporting logic: Possession as owner applies to both good faith and bad faith cases}

\textbf{Intermediate Representation:}

\begin{quote}
If possession as owner is obligated, then (good faith exercise is obligated or bad faith exercise is obligated)
\end{quote}

\textbf{Witness Clause:}

\begin{lstlisting}[style=witnessstyle, language=Witness]
clause tesis180040_posesion_ambas_fe = oblig(posesion_como_propietario) IMPLIES (oblig(ejercicio_buena_fe) OR oblig(ejercicio_mala_fe));
\end{lstlisting}

\subsubsection{Clause: tesis180040\_dominio\_economico}

\textit{Economic dominion: Must be dominator of the thing, commanding and enjoying it as owner}

\textbf{Intermediate Representation:}

\begin{quote}
possession as owner is obligated
\end{quote}

\textbf{Witness Clause:}

\begin{lstlisting}[style=witnessstyle, language=Witness]
clause tesis180040_dominio_economico = oblig(posesion_como_propietario) IMPLIES oblig(causa_generadora_probada);
\end{lstlisting}

\subsubsection{Clause: tesis180040\_posesion\_originaria}

\textit{Original possession: Must begin with cause different from derived possession}

\textbf{Intermediate Representation:}

\begin{quote}
possession as owner is obligated and not (temporal possession is obligated)
\end{quote}

\textbf{Witness Clause:}

\begin{lstlisting}[style=witnessstyle, language=Witness]
clause tesis180040_posesion_originaria = oblig(posesion_como_propietario) AND not(oblig(posesion_temporal)) IMPLIES oblig(causa_generadora_probada);
\end{lstlisting}

\newpage

\section{Thesis 184251: USUCAPIÓN. CASO EN QUE NO PROCEDE ESA ACCIÓN. SEGUNDO TRIBUNAL COLEGIADO EN MATERIA CIVIL DEL SEXTO CIRCUITO.}
\label{thesis:184251}

\subsection{Original Thesis Text}

\begin{quote}
                                                   Semanario Judicial de la Federación

                                                  Tesis

 Registro digital: 184251

 Instancia: Tribunales             Novena Época                           Materia(s): Civil
 Colegiados de Circuito

 Tesis: VI.2o.C.315 C              Fuente: Semanario Judicial de la       Tipo: Aislada
                                   Federación y su Gaceta.
                                   Tomo XVII, Mayo de 2003, página
                                   1287


USUCAPIÓN. CASO EN QUE NO PROCEDE ESA ACCIÓN.


Para que se declare probada la acción de usucapión sobre un bien inmueble, es menester que la
parte promovente acredite los siguientes elementos: a) Que cuenta con justo título para poseer el
bien a usucapir, demostrando la existencia de la causa generadora de la posesión, lo que implica
revelar el acto que la originó, la fecha y el lugar exactos en que tuvo verificativo, los sujetos que
intervinieron y la materia del mismo; y, b) Las cualidades de su posesión, es decir, que ha ejercido
la posesión a nombre propio, de manera pública, pacífica, continua y por diez años si es de buena
fe o veinte si es de mala fe; de ahí que si sólo se justifica el elemento relativo a la causa generadora
de la posesión, no así el segundo de los elementos, ya sea porque la posesión no cumplió con
alguna de las cualidades exigidas legalmente, o bien, con ninguna, la acción intentada no puede
prosperar.

SEGUNDO TRIBUNAL COLEGIADO EN MATERIA CIVIL DEL SEXTO CIRCUITO.


Amparo directo 491/2002. Isabel Pérez García. 13 de febrero de 2003. Unanimidad de votos.
Ponente: Raúl Armando Pallares Valdez. Secretaria: Gloria Margarita Romero Velázquez.

Véase: Semanario Judicial de la Federación y su Gaceta, Novena Época, Tomo I, junio de 1995,
página 374, tesis VI.2o. J/6, de rubro: "USUCAPIÓN. CAUSA GENERADORA DE LA POSESIÓN.
DEBE SEÑALARSE PROPORCIONANDO TODOS AQUELLOS DATOS QUE REVELAN SU
EXISTENCIA (LEGISLACIÓN DEL ESTADO DE PUEBLA).".




                                              Pág. 1 de 2              Fecha de impresión 01/10/2025
  https://sjf2.scjn.gob.mx/detalle/tesis/184251
                                                Semanario Judicial de la Federación




                                           Pág. 2 de 2       Fecha de impresión 01/10/2025
https://sjf2.scjn.gob.mx/detalle/tesis/184251

\end{quote}

\subsection{Formalization}

\subsubsection{Clause: tesis184251\_principio\_fundamental}

\textit{Thesis-specific principle: Good faith possessors require just title as generating cause of possession for adverse possession to operate}

\textbf{Intermediate Representation:}

\begin{quote}
good faith exercise is obligated and possession as owner is obligated
\end{quote}

\textbf{Witness Clause:}

\begin{lstlisting}[style=witnessstyle, language=Witness]
clause tesis184251_principio_fundamental = oblig(ejercicio_buena_fe) AND oblig(posesion_como_propietario) IMPLIES oblig(causa_generadora_probada);
\end{lstlisting}

\subsubsection{Clause: tesis184251\_revelar\_causa}

\textit{Supporting logic: Must reveal and fully prove the generating cause of possession}

\textbf{Intermediate Representation:}

\begin{quote}
ejercer prescripcion positiva is obligated
\end{quote}

\textbf{Witness Clause:}

\begin{lstlisting}[style=witnessstyle, language=Witness]
clause tesis184251_revelar_causa = oblig(ejercer_prescripcion_positiva) IMPLIES oblig(causa_generadora_revelada);
\end{lstlisting}

\subsubsection{Clause: tesis184251\_determinacion\_judicial}

\textit{Judicial determination: Judge needs generating fact/act to determine possession quality and timing}

\textbf{Intermediate Representation:}

\begin{quote}
If causa generadora revelada is obligated, then (possession as owner is obligated and (good faith exercise is obligated and not (bad faith exercise is obligated)))
\end{quote}

\textbf{Witness Clause:}

\begin{lstlisting}[style=witnessstyle, language=Witness]
clause tesis184251_determinacion_judicial = oblig(causa_generadora_revelada) IMPLIES (oblig(posesion_como_propietario) AND (oblig(ejercicio_buena_fe) AND not(oblig(ejercicio_mala_fe))));
\end{lstlisting}

\subsubsection{Clause: tesis184251\_determinacion\_temporal}

\textit{Timing requirement: Must establish when prescription period begins}

\textbf{Intermediate Representation:}

\begin{quote}
If causa generadora revelada is obligated, then (posesion tiempo 5 anos is obligated or posesion tiempo 10 anos is obligated)
\end{quote}

\textbf{Witness Clause:}

\begin{lstlisting}[style=witnessstyle, language=Witness]
clause tesis184251_determinacion_temporal = oblig(causa_generadora_revelada) IMPLIES (oblig(posesion_tiempo_5_anos) OR oblig(posesion_tiempo_10_anos));
\end{lstlisting}

\newpage

\section{Thesis 187145: PRESCRIPCIÓN POSITIVA. EL ALLANAMIENTO A LA DEMANDA ES APTO PARA DEMOSTRAR LA POSESIÓN POR PARTE DEL ACTOR Y SUS DEMÁS ATRIBUTOS (LEGISLACIÓN DEL}
\label{thesis:187145}

\subsection{Original Thesis Text}

\begin{quote}
                                                   Semanario Judicial de la Federación

                                                  Tesis

 Registro digital: 187145

 Instancia: Tribunales            Novena Época                           Materia(s): Civil
 Colegiados de Circuito

 Tesis: XVII.1o.30 C              Fuente: Semanario Judicial de la       Tipo: Aislada
                                  Federación y su Gaceta.
                                  Tomo XV, Abril de 2002, página
                                  1313


PRESCRIPCIÓN POSITIVA. EL ALLANAMIENTO A LA DEMANDA ES APTO PARA DEMOSTRAR
LA POSESIÓN POR PARTE DEL ACTOR Y SUS DEMÁS ATRIBUTOS (LEGISLACIÓN DEL
ESTADO DE CHIHUAHUA).


En ejercicio de la acción de prescripción positiva, el demandado al producir su contestación a la
demanda puede allanarse a todas y cada una de las prestaciones reclamadas e, incluso, aceptar en
su totalidad los hechos en que se sustente tal acción, máxime si el inmueble que el enjuiciante
pretende prescribir adquisitivamente se encuentra inscrito a nombre de dicho demandado en el
Registro Público de la Propiedad; que entre aquél y éste existió el acto jurídico generador de la
posesión del bien en litigio y, asimismo, la fecha en que se verificó el mismo, es decir, por el que le
fue transmitido el dominio de ese bien, todo lo cual incuestionablemente revela la intención evidente
del enjuiciado y la aceptación de que el accionante ha poseído el inmueble desde la fecha de
realización del referido acto, con los atributos inherentes a que ese hecho sea público, pacífico,
continuo, ininterrumpido y de buena fe. Consecuentemente, atento tal confesión judicial expresa de
la parte demandada, derivada del allanamiento debidamente ratificado, deviene indiscutible y
patente que el promovente ya no está obligado a rendir más pruebas para justificar la posesión del
bien para que proceda dicha prescripción positiva en los términos a que se refieren los artículos
1153 y 1158 del Código Civil para el Estado de Chihuahua, en relación con el 260 del Código de
Procedimientos Civiles de la propia entidad federativa.

PRIMER TRIBUNAL COLEGIADO DEL DÉCIMO SÉPTIMO CIRCUITO.


Amparo directo 368/2001. Graciela Mcgrew Hernández de Escalante. 8 de febrero de 2002.
Unanimidad de votos. Ponente: Amador Muñoz Torres, secretario de tribunal autorizado por el
Pleno del Consejo de la Judicatura Federal para desempeñar las funciones de Magistrado.
Secretario: Juan Carlos Zamora Tejeda.

Véase: Semanario Judicial de la Federación y su Gaceta, Novena Época, Tomo XIII, enero de 2001,
página 1810, tesis II.2o.C.258 C, de rubro: "USUCAPIÓN. EL ALLANAMIENTO A LA DEMANDA ES
APTO PARA DEMOSTRAR LA POSESIÓN POR PARTE DEL ACTOR Y SUS DEMÁS ATRIBUTOS
(LEGISLACIÓN DEL ESTADO DE MÉXICO).".

Nota: Esta tesis fue objeto de la denuncia relativa a la contradicción de tesis 4/2017 del Pleno del
Décimo Séptimo Circuito de la que derivó la tesis jurisprudencial PC.XVII. J/17 C (10a.) de título y
subtítulo: "PRESCRIPCIÓN POSITIVA. EL ALLANAMIENTO A LA DEMANDA NO ES APTO PARA


                                             Pág. 1 de 2              Fecha de impresión 01/10/2025
  https://sjf2.scjn.gob.mx/detalle/tesis/187145
                                                  Semanario Judicial de la Federación

DEMOSTRAR LOS 'ATRIBUTOS DE LA POSESIÓN' (LEGISLACIÓN DEL ESTADO DE
CHIHUAHUA)."

Por ejecutoria del 3 de abril de 2019, la Primera Sala declaró sin materia la contradicción de tesis
179/2018 derivada de la denuncia de la que fue objeto el criterio contenido en esta tesis, al existir la
jurisprudencia PC.XVII. J/17 C (10a.) que resuelve el mismo problema jurídico.




                                              Pág. 2 de 2              Fecha de impresión 01/10/2025
  https://sjf2.scjn.gob.mx/detalle/tesis/187145

\end{quote}

\subsection{Formalization}

\subsubsection{Clause: tesis187145\_principio\_fundamental}

\textit{Thesis-specific principle: Adverse possession can be invoked by action or exception, but requires same proof elements}

\textbf{Intermediate Representation:}

\begin{quote}
possession as owner is obligated and (good faith exercise is obligated or bad faith exercise is obligated) and (oblig(posesion\_tiempo\_5\_anos) OR oblig(posesion\_tiempo\_10\_anos)) IMPLIES oblig(ejercer\_prescripcion\_positiva)
\end{quote}

\textbf{Witness Clause:}

\begin{lstlisting}[style=witnessstyle, language=Witness]
clause tesis187145_principio_fundamental = oblig(posesion_como_propietario) AND (oblig(ejercicio_buena_fe) OR oblig(ejercicio_mala_fe)) AND (oblig(posesion_tiempo_5_anos) OR oblig(posesion_tiempo_10_anos)) IMPLIES oblig(ejercer_prescripcion_positiva);
\end{lstlisting}

\subsubsection{Clause: tesis187145\_igualdad\_prueba}

\textit{Supporting logic: Both action and exception require same proof elements regardless of procedural route}

\textbf{Intermediate Representation:}

\begin{quote}
ejercer prescripcion positiva is obligated
\end{quote}

\textbf{Witness Clause:}

\begin{lstlisting}[style=witnessstyle, language=Witness]
clause tesis187145_igualdad_prueba = oblig(ejercer_prescripcion_positiva) IMPLIES oblig(posesion_como_propietario);
\end{lstlisting}

\subsubsection{Clause: tesis187145\_plazos\_temporales}

\textit{Time requirements: 5 years for good faith, 10 years for bad faith}

\textbf{Intermediate Representation:}

\begin{quote}
good faith exercise is obligated
\end{quote}

\textbf{Witness Clause:}

\begin{lstlisting}[style=witnessstyle, language=Witness]
clause tesis187145_plazos_temporales = oblig(ejercicio_buena_fe) IMPLIES oblig(posesion_tiempo_5_anos);
\end{lstlisting}

\subsubsection{Clause: tesis187145\_plazos\_mala\_fe}

\textbf{Intermediate Representation:}

\begin{quote}
bad faith exercise is obligated
\end{quote}

\textbf{Witness Clause:}

\begin{lstlisting}[style=witnessstyle, language=Witness]
clause tesis187145_plazos_mala_fe = oblig(ejercicio_mala_fe) IMPLIES oblig(posesion_tiempo_10_anos);
\end{lstlisting}

\subsubsection{Clause: tesis187145\_objetivo\_registral}

\textit{Target requirement: Action must be directed against registered owner}

\textbf{Intermediate Representation:}

\begin{quote}
ejercer prescripcion positiva is obligated
\end{quote}

\textbf{Witness Clause:}

\begin{lstlisting}[style=witnessstyle, language=Witness]
clause tesis187145_objetivo_registral = oblig(ejercer_prescripcion_positiva) IMPLIES oblig(titularidad_registral);
\end{lstlisting}

\newpage

\section{Thesis 188142: PRESCRIPCIÓN ADQUISITIVA, NO BASTA CON REVELAR LA CAUSA GENERADORA DE LA POSESIÓN, SINO QUE DEBE ACREDITARSE (LEGISLACIÓN DEL ESTADO DE MÉXICO).}
\label{thesis:188142}

\subsection{Original Thesis Text}

\begin{quote}
                                                   Semanario Judicial de la Federación

                                                  Tesis

 Registro digital: 188142

 Instancia: Tribunales            Novena Época                           Materia(s): Civil
 Colegiados de Circuito

 Tesis: II.3o.C. J/2              Fuente: Semanario Judicial de la       Tipo: Jurisprudencia
                                  Federación y su Gaceta.
                                  Tomo XIV, Diciembre de 2001,
                                  página 1581


PRESCRIPCIÓN ADQUISITIVA, NO BASTA CON REVELAR LA CAUSA GENERADORA DE LA
POSESIÓN, SINO QUE DEBE ACREDITARSE (LEGISLACIÓN DEL ESTADO DE MÉXICO).


El artículo 911 del Código Civil del Estado de México, establece que la posesión necesaria para
usucapir debe ser en concepto de propietario, pacífica, continua y pública. De ahí que uno de los
requisitos para que opere la prescripción adquisitiva, es el relativo a que el bien a usucapir se posea
con el carácter de propietario y tal calidad sólo puede ser calificada si se invoca la causa
generadora de la posesión, dado que si ésta no se expone, el juzgador está imposibilitado para
determinar si se cumple con tal elemento. Así, el precepto en comento, en cuanto a la condición
reseñada se complementa con lo dispuesto en el artículo 801 del ordenamiento citado, en cuanto a
que sólo la posesión que se adquiere y disfruta en concepto de dueño de la cosa poseída puede
producir la usucapión. De tal manera que, cuando se promueve un juicio de usucapión, es menester
que el actor revele dicha causa y puede ser: el hecho o acto jurídico que hace adquirir un derecho y
que entronca con la causa; el documento en que consta ese acto o hecho adquisitivo; el derecho
mismo que asiste a una persona y que la legitima activa o pasivamente, tanto para que la autoridad
esté en aptitud de fijar la calidad de la posesión, originaria o derivada, como para que se pueda
computar el término de ella, ya sea de buena o mala fe. Por lo cual, si alguna de las partes invoca
como origen generador de su posesión, un contrato verbal de compraventa, ello no significa que
haya cumplido con el requisito citado, pues la adquisición, desde el punto de vista jurídico, es la
incorporación de una cosa o derecho a la esfera patrimonial de una persona, en tanto que aquella
declaración solamente constituye una expresión genérica que se utiliza para poner de manifiesto
que un bien o un derecho ha ingresado al patrimonio de una persona, pero no indica, por sí misma,
el medio o forma en que se ingresó, como tampoco señala las cualidades específicas o los efectos
de la obtención, ni precisa si esa incorporación es plena o limitada, si es originaria o derivada.
Consecuentemente, en términos de los numerales aludidos así como de su interpretación armónica
y sistemática con los demás que se refieren al título tercero (De la posesión), título cuarto (De la
propiedad en general y de los medios para adquirirla) y capítulo quinto (De la usucapión), no basta
con revelar la causa generadora de la posesión, sino que debe acreditarse. Lo cual se corrobora
con la jurisprudencia de rubro: "PRESCRIPCIÓN ADQUISITIVA. PARA QUE SE ENTIENDA
SATISFECHO EL REQUISITO DE LA EXISTENCIA DE LA ‘POSESIÓN EN CONCEPTO DE
PROPIETARIO’ EXIGIDO POR EL CÓDIGO CIVIL PARA EL DISTRITO FEDERAL Y POR LAS
DIVERSAS LEGISLACIONES DE LOS ESTADOS DE LA REPÚBLICA QUE CONTIENEN
DISPOSICIONES IGUALES, ES NECESARIO DEMOSTRAR LA EXISTENCIA DE UN TÍTULO DEL
QUE SE DERIVE LA POSESIÓN.", en la que la entonces Tercera Sala de la Suprema Corte de
Justicia de la Nación llegó a la misma conclusión, al analizar los artículos 826, 1151, fracción I y
1152 del Código Civil para el Distrito Federal en Materia Común y para toda la República en Materia


                                             Pág. 1 de 2              Fecha de impresión 01/10/2025
  https://sjf2.scjn.gob.mx/detalle/tesis/188142
                                                  Semanario Judicial de la Federación

Federal, que contienen iguales disposiciones que los artículos 801, 911, fracción I y 912 del Código
Civil del Estado de México.

TERCER TRIBUNAL COLEGIADO EN MATERIA CIVIL DEL SEGUNDO CIRCUITO.


Amparo directo 555/99. María Asunción García Martínez. 15 de febrero de 2000. Unanimidad de
votos. Ponente: Ricardo Romero Vázquez. Secretario: José Fernando García Quiroz.

Amparo directo 365/2000. Antonio Álvarez Martínez. 18 de octubre de 2000. Unanimidad de votos.
Ponente: Ana María Serrano Oseguera de Torres. Secretario: José Antonio Franco Vera.

Amparo directo 747/2000. José Carmen Martínez Moreno. 22 de noviembre de 2000. Unanimidad
de votos. Ponente: Felipe Alfredo Fuente Barrera. Secretario: Guillermo Hindman Pozos.

Amparo directo 557/2000. Transportes y Montajes, Construcciones, S.A. de C.V. 16 de enero de
2001. Unanimidad de votos. Ponente: Ana María Serrano Oseguera de Torres. Secretario: José
Antonio Franco Vera.

Amparo directo 456/2001. Guadalupe Torres García, en su carácter de albacea de la sucesión
intestamentaria de Carlos Manuel Cedillo Arce. 9 de octubre de 2001. Unanimidad de votos.
Ponente: Felipe Alfredo Fuentes Barrera. Secretario: José del Carmen Gutiérrez Meneses.

Notas:

La tesis citada aparece publicada con el número 322, en el Apéndice al Semanario Judicial de la
Federación 1917-2000, Tomo IV, Materia Civil, página 271.

Por ejecutoria de fecha 12 de noviembre de 2003, la Primera Sala declaró improcedente la
contradicción de tesis 26/2003-PS en que participó el presente criterio, en virtud de que el Primer
Tribunal Colegiado en Materia Civil del Segundo Circuito, previamente a la denuncia de
contradicción de tesis, se apartó del criterio divergente y adoptó uno nuevo coincidente con el
Tercer Tribunal Colegiado de la misma materia y circuito.




                                             Pág. 2 de 2            Fecha de impresión 01/10/2025
  https://sjf2.scjn.gob.mx/detalle/tesis/188142

\end{quote}

\subsection{Formalization}

\subsubsection{Clause: tesis188142\_principio\_fundamental}

\textit{Faith exclusivity: Possession is either good faith XOR bad faith (mutually exclusive) Thesis-specific principle: Unregistered contracts cannot establish adverse possession against registered owners}

\textbf{Intermediate Representation:}

\begin{quote}
registered ownership is obligated and unregistered contract is obligated
\end{quote}

\textbf{Witness Clause:}

\begin{lstlisting}[style=witnessstyle, language=Witness]
clause tesis188142_principio_fundamental = oblig(titularidad_registral) AND oblig(contrato_no_inscrito) IMPLIES not(claim(demanda_prescripcion));
\end{lstlisting}

\subsubsection{Clause: tesis188142\_proteccion\_registral}

\textit{Supporting legal logic}

\textbf{Intermediate Representation:}

\begin{quote}
registered ownership is obligated
\end{quote}

\textbf{Witness Clause:}

\begin{lstlisting}[style=witnessstyle, language=Witness]
clause tesis188142_proteccion_registral = oblig(titularidad_registral) IMPLIES claim(oponibilidad_terceros);
\end{lstlisting}

\subsubsection{Clause: tesis188142\_limitacion\_no\_inscripcion}

\textbf{Intermediate Representation:}

\begin{quote}
unregistered contract is obligated
\end{quote}

\textbf{Witness Clause:}

\begin{lstlisting}[style=witnessstyle, language=Witness]
clause tesis188142_limitacion_no_inscripcion = oblig(contrato_no_inscrito) IMPLIES not(oblig(oponibilidad_terceros));
\end{lstlisting}

\subsubsection{Clause: tesis188142\_consecuencia\_logica}

\textbf{Intermediate Representation:}

\begin{quote}
not (opposability to third parties is obligated)
\end{quote}

\textbf{Witness Clause:}

\begin{lstlisting}[style=witnessstyle, language=Witness]
clause tesis188142_consecuencia_logica = not(oblig(oponibilidad_terceros)) IMPLIES not(claim(demanda_prescripcion));
\end{lstlisting}

\newpage

\section{Thesis 188144: PRESCRIPCIÓN ADQUISITIVA. EL TÍTULO GENERADOR DE LA POSESIÓN SÓLO SE REQUIERE PARA DETERMINAR SU ORIGEN.}
\label{thesis:188144}

\subsection{Original Thesis Text}

\begin{quote}
                                                   Semanario Judicial de la Federación

                                                  Tesis

 Registro digital: 188144

 Instancia: Tribunales            Novena Época                          Materia(s): Civil
 Colegiados de Circuito

 Tesis: III.1o.C.126 C            Fuente: Semanario Judicial de la      Tipo: Aislada
                                  Federación y su Gaceta.
                                  Tomo XIV, Diciembre de 2001,
                                  página 1778


PRESCRIPCIÓN ADQUISITIVA. EL TÍTULO GENERADOR DE LA POSESIÓN SÓLO SE
REQUIERE PARA DETERMINAR SU ORIGEN.


Si el acto por virtud del cual se cree propietario quien opone la prescripción es imperfecto o
deficiente, ello no es causa suficiente para desestimarla, dado que éste sólo tiene la obligación de
justificar la existencia del hecho que señaló como causa generadora de su posesión, como un acto
jurídico que le permite comportarse objetivamente como propietario, mediante la realización de
actos que revelen su dominio sobre el inmueble para hacerlo suyo. Es decir, lo que realmente
importa es que el juzgador conozca el acto generador de la posesión, para que pueda determinar si
la calidad de ésta es en concepto de propietario, originaria o derivada, de buena o mala fe, porque
el acto jurídico defectuoso no es el que constituye la fuente de adquisición de la propiedad, sino que
ésta se encuentra en la propia ley que prevé la institución de la usucapión, dado que aquel acto sólo
cumple la función de poner de manifiesto que la posesión no se disfruta en forma derivada, sino en
concepto de propietario, sobre la base de un título que aun cuando esté viciado, la ley le atribuye
efectos jurídicos.

PRIMER TRIBUNAL COLEGIADO EN MATERIA CIVIL DEL TERCER CIRCUITO.


Amparo directo 4231/2000. José Vázquez Huerta y otra. 29 de marzo de 2001. Unanimidad de
votos. Ponente: Carlos Arturo González Zárate. Secretario: Miguel Ivo Moreno Vidrio.

Véase: Semanario Judicial de la Federación, Octava Época, Tomo XIII, junio de 1994, página 627,
tesis I.5o.C.553 C, de rubro: "PRESCRIPCIÓN POSITIVA. DOCUMENTO TRASLATIVO DE
DOMINIO VICIADO INVOCADO COMO CAUSA DE LA POSESIÓN EN CONCEPTO DE
PROPIETARIO, HACE FACTIBLE QUE PROSPERE LA ACCIÓN, SI SE IGNORAN LOS VICIOS
DEL TÍTULO Y SE PRUEBAN LOS DEMÁS REQUISITOS LEGALES.".

Nota: Por ejecutoria del 11 de abril de 2018, la Primera Sala declaró inexistente la contradicción de
tesis 346/2017 derivada de la denuncia de la que fue objeto el criterio contenido en esta tesis, al
estimarse que no son discrepantes los criterios materia de la denuncia respectiva.




                                             Pág. 1 de 2             Fecha de impresión 01/10/2025
  https://sjf2.scjn.gob.mx/detalle/tesis/188144
                                                Semanario Judicial de la Federación




                                           Pág. 2 de 2       Fecha de impresión 01/10/2025
https://sjf2.scjn.gob.mx/detalle/tesis/188144

\end{quote}

\subsection{Formalization}

\subsubsection{Clause: tesis188144\_principio\_fundamental}

\textit{Thesis-specific principle: Possession must be "in concept of owner" (peaceful, continuous, public) for both good and bad faith}

\textbf{Intermediate Representation:}

\begin{quote}
possession as owner is obligated and (oblig(ejercicio\_buena\_fe) OR oblig(ejercicio\_mala\_fe)) IMPLIES oblig(ejercer\_prescripcion\_positiva)
\end{quote}

\textbf{Witness Clause:}

\begin{lstlisting}[style=witnessstyle, language=Witness]
clause tesis188144_principio_fundamental = oblig(posesion_como_propietario) AND (oblig(ejercicio_buena_fe) OR oblig(ejercicio_mala_fe)) IMPLIES oblig(ejercer_prescripcion_positiva);
\end{lstlisting}

\subsubsection{Clause: tesis188144\_posesion\_ambas\_fe}

\textit{Supporting logic: Possession as owner applies to both good faith and bad faith cases}

\textbf{Intermediate Representation:}

\begin{quote}
If possession as owner is obligated, then (good faith exercise is obligated or bad faith exercise is obligated)
\end{quote}

\textbf{Witness Clause:}

\begin{lstlisting}[style=witnessstyle, language=Witness]
clause tesis188144_posesion_ambas_fe = oblig(posesion_como_propietario) IMPLIES (oblig(ejercicio_buena_fe) OR oblig(ejercicio_mala_fe));
\end{lstlisting}

\subsubsection{Clause: tesis188144\_dominio\_economico}

\textit{Economic dominion: Must be dominator of the thing, commanding and enjoying it as owner}

\textbf{Intermediate Representation:}

\begin{quote}
possession as owner is obligated
\end{quote}

\textbf{Witness Clause:}

\begin{lstlisting}[style=witnessstyle, language=Witness]
clause tesis188144_dominio_economico = oblig(posesion_como_propietario) IMPLIES oblig(causa_generadora_probada);
\end{lstlisting}

\subsubsection{Clause: tesis188144\_posesion\_originaria}

\textit{Original possession: Must begin with cause different from derived possession}

\textbf{Intermediate Representation:}

\begin{quote}
possession as owner is obligated and not (temporal possession is obligated)
\end{quote}

\textbf{Witness Clause:}

\begin{lstlisting}[style=witnessstyle, language=Witness]
clause tesis188144_posesion_originaria = oblig(posesion_como_propietario) AND not(oblig(posesion_temporal)) IMPLIES oblig(causa_generadora_probada);
\end{lstlisting}

\newpage

\section{Thesis 188544: PRESCRIPCIÓN POSITIVA. LA ACCIÓN ES IMPROCEDENTE CUANDO NO SE DEMUESTRA LA CAUSA GENERADORA DE LA POSESIÓN (LEGISLACIÓN AGRARIA VIGENTE).}
\label{thesis:188544}

\subsection{Original Thesis Text}

\begin{quote}
                                                   Semanario Judicial de la Federación

                                                  Tesis

 Registro digital: 188544

 Instancia: Tribunales            Novena Época                         Materia(s): Administrativa
 Colegiados de Circuito

 Tesis: VI.2o.A.27 A              Fuente: Semanario Judicial de la  Tipo: Aislada
                                  Federación y su Gaceta.
                                  Tomo XIV, Octubre de 2001, página
                                  1163


PRESCRIPCIÓN POSITIVA. LA ACCIÓN ES IMPROCEDENTE CUANDO NO SE DEMUESTRA LA
CAUSA GENERADORA DE LA POSESIÓN (LEGISLACIÓN AGRARIA VIGENTE).


Una recta interpretación de la obligación impuesta en el artículo 48 de la Ley Agraria vigente,
permite establecer, para que opere la prescripción positiva, ciertos presupuestos básicos,
consistentes en: a) Que se posean las tierras en concepto de titular de derechos de ejidatario; b)
Que la posesión sea respecto de tierras ejidales, siempre y cuando no se trate de aquellas
destinadas al asentamiento humano ni de bosques o selvas; y c) Que esa posesión debe ser de
manera pacífica, continua y pública durante un periodo de cinco años si es de buena fe, o de diez si
es de mala fe. Sin embargo, aun cuando tal precepto legal no exija el "justo título" como elemento
de la prescripción, ello no significa que toda posesión es apta para prescribir, pues para que
prospere la manifestación de que se adquirió la posesión y se disfruta en concepto o con el carácter
de propietario, es menester demostrar la causa generadora de la misma (como sería, por ejemplo,
cualquier acto traslativo de dominio), para que el juzgador esté en condiciones de determinar si la
posesión es en concepto de propietario, originaria o derivada, de buena o mala fe y a partir de qué
momento se contará el plazo para usucapir; por lo que si el demandante no revela, ni acredita, la
forma en que entró a poseer, resulta evidente, además de que sólo se consideraría subjetivamente
como propietario, que la acción prescriptiva es improcedente.

SEGUNDO TRIBUNAL COLEGIADO EN MATERIA ADMINISTRATIVA DEL SEXTO CIRCUITO.


Amparo directo 32/2000. Jaime Pérez Rosas. 15 de noviembre de 2000. Unanimidad de votos.
Ponente: Omar Losson Ovando. Secretario: Jaime Silva Hernández.

Véase: Semanario Judicial de la Federación y su Gaceta, Novena Época, Tomo VII, marzo de 1998,
página 813, tesis XI.1o.4 A, de rubro: "PRESCRIPCIÓN ADQUISITIVA EN MATERIA AGRARIA.
DEBE INVOCARSE Y ACREDITARSE LA CAUSA GENERADORA DE LA POSESIÓN.".

Notas:

Por ejecutoria de fecha 4 de abril de 2003, la Segunda Sala declaró inexistente la contradicción de
tesis 179/2002-SS en que participó el presente criterio.

Esta tesis contendió en la contradicción 96/2004-SS resuelta por la Segunda Sala, de la que derivó
la tesis 2a./J. 207/2004, que aparece publicada en el Semanario Judicial de la Federación y su


                                             Pág. 1 de 2            Fecha de impresión 01/10/2025
  https://sjf2.scjn.gob.mx/detalle/tesis/188544
                                                  Semanario Judicial de la Federación

Gaceta, Novena Época, Tomo XXI, enero de 2005, página 575, con el rubro: "PRESCRIPCIÓN
ADQUISITIVA EN MATERIA AGRARIA. PARA SU PROCEDENCIA NO SE REQUIERE DE JUSTO
TÍTULO."




                                             Pág. 2 de 2       Fecha de impresión 01/10/2025
  https://sjf2.scjn.gob.mx/detalle/tesis/188544

\end{quote}

\subsection{Formalization}

\subsubsection{Clause: tesis188544\_principio\_fundamental}

\textit{Thesis-specific principle: Possession must be "in concept of owner" (peaceful, continuous, public) for both good and bad faith}

\textbf{Intermediate Representation:}

\begin{quote}
possession as owner is obligated and (oblig(ejercicio\_buena\_fe) OR oblig(ejercicio\_mala\_fe)) IMPLIES oblig(ejercer\_prescripcion\_positiva)
\end{quote}

\textbf{Witness Clause:}

\begin{lstlisting}[style=witnessstyle, language=Witness]
clause tesis188544_principio_fundamental = oblig(posesion_como_propietario) AND (oblig(ejercicio_buena_fe) OR oblig(ejercicio_mala_fe)) IMPLIES oblig(ejercer_prescripcion_positiva);
\end{lstlisting}

\subsubsection{Clause: tesis188544\_posesion\_ambas\_fe}

\textit{Supporting logic: Possession as owner applies to both good faith and bad faith cases}

\textbf{Intermediate Representation:}

\begin{quote}
If possession as owner is obligated, then (good faith exercise is obligated or bad faith exercise is obligated)
\end{quote}

\textbf{Witness Clause:}

\begin{lstlisting}[style=witnessstyle, language=Witness]
clause tesis188544_posesion_ambas_fe = oblig(posesion_como_propietario) IMPLIES (oblig(ejercicio_buena_fe) OR oblig(ejercicio_mala_fe));
\end{lstlisting}

\subsubsection{Clause: tesis188544\_dominio\_economico}

\textit{Economic dominion: Must be dominator of the thing, commanding and enjoying it as owner}

\textbf{Intermediate Representation:}

\begin{quote}
possession as owner is obligated
\end{quote}

\textbf{Witness Clause:}

\begin{lstlisting}[style=witnessstyle, language=Witness]
clause tesis188544_dominio_economico = oblig(posesion_como_propietario) IMPLIES oblig(causa_generadora_probada);
\end{lstlisting}

\subsubsection{Clause: tesis188544\_posesion\_originaria}

\textit{Original possession: Must begin with cause different from derived possession}

\textbf{Intermediate Representation:}

\begin{quote}
possession as owner is obligated and not (temporal possession is obligated)
\end{quote}

\textbf{Witness Clause:}

\begin{lstlisting}[style=witnessstyle, language=Witness]
clause tesis188544_posesion_originaria = oblig(posesion_como_propietario) AND not(oblig(posesion_temporal)) IMPLIES oblig(causa_generadora_probada);
\end{lstlisting}

\newpage

\section{Thesis 189628: PRESCRIPCIÓN POSITIVA. LA AUSENCIA DE FORMALIDADES EN EL CONTRATO DE DONACIÓN, EXHIBIDO COMO JUSTO TÍTULO EN AQUÉLLA, NO IMPIDE LA PROCEDENCIA}
\label{thesis:189628}

\subsection{Original Thesis Text}

\begin{quote}
                                                   Semanario Judicial de la Federación

                                                  Tesis

 Registro digital: 189628

 Instancia: Tribunales            Novena Época                           Materia(s): Civil
 Colegiados de Circuito

 Tesis: I.3o.C.219 C              Fuente: Semanario Judicial de la       Tipo: Aislada
                                  Federación y su Gaceta.
                                  Tomo XIII, Mayo de 2001, página
                                  1201


PRESCRIPCIÓN POSITIVA. LA AUSENCIA DE FORMALIDADES EN EL CONTRATO DE
DONACIÓN, EXHIBIDO COMO JUSTO TÍTULO EN AQUÉLLA, NO IMPIDE LA PROCEDENCIA
DE LA ACCIÓN SIEMPRE Y CUANDO SE ACREDITE LA EXISTENCIA DE DICHO CONTRATO.


La donación es un contrato por el que una persona transfiere a otra, gratuitamente, una parte o la
totalidad de sus bienes presentes, y para que ésta se perfeccione, es preciso que el donatario la
acepte y haga saber su aceptación al donante, de conformidad con lo dispuesto en el artículo 2340
del Código Civil del Distrito Federal. Además, para la validez de la donación, debe hacerse constar
en escritura pública, según lo dispone el artículo 2345 del código sustantivo invocado. Sin embargo,
cuando existe ausencia de formalidad en la donación, no impide la prescripción adquisitiva, porque
el haber adquirido y disfrutar la posesión en concepto de dueño, implica contar con un justo título
que legitime la detentación que tiene del inmueble, y que para los efectos de la prescripción, es el
hecho que sirve de causa a la posesión, ya que conforme al artículo 806 del Código Civil para el
Distrito Federal, es poseedor de buena fe el que entra en la posesión en virtud de un título suficiente
para darle derecho de poseer, y también se tiene como tal, al que ignora los vicios de su título que
le impiden poseer con derecho. Para que sea apto para la usucapión, ese título debe ser justo,
verdadero y válido. Por justo título debe entenderse el que legalmente basta para transferir el
dominio de la cosa de cuya prescripción se trate, es decir, el que produciría la transmisión y
adquisición del dominio, sin tomar en cuenta el vicio o defecto que precisamente a través de la
prescripción se subsanará. Por tanto, son eficaces para ello, la compraventa, la permuta, la
donación, la herencia, el legado y, en general, todos aquellos que transmiten el dominio. Título
verdadero es el de existencia real y no asimilado; y el requisito de la validez se debe interpretar en
el sentido de que no se puede exigir que el título sea perfectamente válido, esto es, que reúna todas
las condiciones necesarias para producir la transmisión del dominio, porque de lo contrario, no haría
falta la prescripción. De ahí que no era necesario que la donación alegada como causa de posesión,
constara en escritura pública y existiera la aceptación de la donataria, puesto que de haberse
consignado en esa forma, la presunta donataria habría adquirido desde entonces, plena e
indiscutiblemente, la propiedad. Pero no obstante que la donación invocada como causa de la
posesión, que no cumple con las formalidades requeridas por la ley, sí es apta para adquirir la
propiedad por prescripción, resulta necesario acreditar la existencia de esa donación, toda vez que
el hecho o el acto en que se afirme en qué consistió la causa generadora de la posesión, siempre
debe acreditarse, para justificar que no se trata de una mera tenencia o disfrute de la cosa que
obedezca a una relación de arrendamiento, comodato, depósito o prenda.

TERCER TRIBUNAL COLEGIADO EN MATERIA CIVIL DEL PRIMER CIRCUITO.



                                             Pág. 1 de 2              Fecha de impresión 01/10/2025
  https://sjf2.scjn.gob.mx/detalle/tesis/189628
                                                  Semanario Judicial de la Federación


Amparo directo 8993/2000. Graciela Arriaga Hernández. 8 de diciembre de 2000. Unanimidad de
votos. Ponente: Neófito López Ramos. Secretario: Rómulo Amadeo Figueroa Salmorán.




                                             Pág. 2 de 2       Fecha de impresión 01/10/2025
  https://sjf2.scjn.gob.mx/detalle/tesis/189628

\end{quote}

\subsection{Formalization}

\subsubsection{Clause: tesis189628\_principio\_fundamental}

\textit{Thesis-specific principle: Good faith requires sufficient title/generating cause, ignorance of title defects, and founded belief of ownership}

\textbf{Intermediate Representation:}

\begin{quote}
good faith exercise is obligated and not (conocimiento vicios is obligated)
\end{quote}

\textbf{Witness Clause:}

\begin{lstlisting}[style=witnessstyle, language=Witness]
clause tesis189628_principio_fundamental = oblig(ejercicio_buena_fe) IMPLIES oblig(causa_generadora_probada) AND not(oblig(conocimiento_vicios));
\end{lstlisting}

\subsubsection{Clause: tesis189628\_mala\_fe}

\textit{Supporting logic: Bad faith possessor has no title/generating cause OR knows title defects that prevent rightful possession}

\textbf{Intermediate Representation:}

\begin{quote}
If bad faith exercise is obligated, then (not (proven generating cause is obligated) or conocimiento vicios is obligated)
\end{quote}

\textbf{Witness Clause:}

\begin{lstlisting}[style=witnessstyle, language=Witness]
clause tesis189628_mala_fe = oblig(ejercicio_mala_fe) IMPLIES (not(oblig(causa_generadora_probada)) OR oblig(conocimiento_vicios));
\end{lstlisting}

\subsubsection{Clause: tesis189628\_requisitos\_buena\_fe}

\textit{Good faith requirements: Must have sufficient title and ignorance of defects}

\textbf{Intermediate Representation:}

\begin{quote}
good faith exercise is obligated
\end{quote}

\textbf{Witness Clause:}

\begin{lstlisting}[style=witnessstyle, language=Witness]
clause tesis189628_requisitos_buena_fe = oblig(ejercicio_buena_fe) IMPLIES oblig(causa_generadora_probada);
\end{lstlisting}

\subsubsection{Clause: tesis189628\_condiciones\_mala\_fe}

\textit{Bad faith conditions: No title or knowledge of defects}

\textbf{Intermediate Representation:}

\begin{quote}
(not(oblig(causa\_generadora\_probada)) OR oblig(conocimiento\_vicios)) IMPLIES oblig(ejercicio\_mala\_fe)
\end{quote}

\textbf{Witness Clause:}

\begin{lstlisting}[style=witnessstyle, language=Witness]
clause tesis189628_condiciones_mala_fe = (not(oblig(causa_generadora_probada)) OR oblig(conocimiento_vicios)) IMPLIES oblig(ejercicio_mala_fe);
\end{lstlisting}

\newpage

\section{Thesis 189906: PRESCRIPCIÓN POSITIVA EN MATERIA AGRARIA. NO OPERA SI SE DESVIRTÚA LA POSESIÓN CON EL CARÁCTER DE TITULAR DE UNA PARCELA.}
\label{thesis:189906}

\subsection{Original Thesis Text}

\begin{quote}
                                                   Semanario Judicial de la Federación

                                                  Tesis

 Registro digital: 189906

 Instancia: Tribunales             Novena Época                           Materia(s): Administrativa
 Colegiados de Circuito

 Tesis: III.1o.A.78 A              Fuente: Semanario Judicial de la       Tipo: Aislada
                                   Federación y su Gaceta.
                                   Tomo XIII, Abril de 2001, página
                                   1108


PRESCRIPCIÓN POSITIVA EN MATERIA AGRARIA. NO OPERA SI SE DESVIRTÚA LA
POSESIÓN CON EL CARÁCTER DE TITULAR DE UNA PARCELA.


El artículo 48 de la Ley Agraria establece una posesión calificada, entendiéndose por ésta que debe
ser con el carácter de titular de derechos ejidales, tratándose de la prescripción adquisitiva. Por
tanto, no es suficiente que se hubiera acreditado que la posesión se ha tenido por más de cinco
años o de diez, ya fuera de buena o de mala fe, de manera pública, pacífica y continua, porque si no
se hizo con aquel carácter, no opera la prescripción a que se refiere el citado numeral. Ahora bien,
si se instaura un procedimiento de privación de derechos agrarios en contra de un ejidatario, sobre
determinada parcela, y a la postre este juicio se resuelve favorable a sus intereses, confirmándose
así su titularidad de aquel terreno, el tiempo de la posesión ejercida durante el juicio sobre privación
de derechos agrarios, que invocó su contrario en la controversia agraria, no es apta para prescribir,
porque esa posesión, por más que haya provenido de que su ejercitante fue designado para
sustituir a quien se pretendía privar, se desnaturaliza cuando se resuelve el juicio privativo
declarando improcedente la acción de privación, quedando vigente, para todos los efectos legales,
la titularidad del ejidatario vencedor.

PRIMER TRIBUNAL COLEGIADO EN MATERIA ADMINISTRATIVA DEL TERCER CIRCUITO.


Amparo directo 425/2000. Consuelo García Obledo y otro. 22 de noviembre de 2000. Unanimidad
de votos. Ponente: Jorge Alfonso Álvarez Escoto. Secretaria: Gabriela Guadalupe Huízar Flores.

Véase: Semanario Judicial de la Federación y su Gaceta, Novena Época, Tomo III, junio de 1996,
tesis XI.2o.6 A, de rubro: "PRESCRIPCIÓN ADQUISITIVA EN MATERIA AGRARIA. POSESIÓN
‘EN CONCEPTO DE TITULAR DE DERECHOS DE EJIDATARIO’ PARA QUE PROSPERE LA.".




                                              Pág. 1 de 2             Fecha de impresión 01/10/2025
  https://sjf2.scjn.gob.mx/detalle/tesis/189906
                                                Semanario Judicial de la Federación




                                           Pág. 2 de 2       Fecha de impresión 01/10/2025
https://sjf2.scjn.gob.mx/detalle/tesis/189906

\end{quote}

\subsection{Formalization}

\subsubsection{Clause: tesis189906\_principio\_fundamental}

\textit{Thesis-specific principle: Just title not required for adverse possession, but must reveal origin of possession and affirm possession as owner}

\textbf{Intermediate Representation:}

\begin{quote}
ejercer prescripcion positiva is obligated and possession as owner is obligated
\end{quote}

\textbf{Witness Clause:}

\begin{lstlisting}[style=witnessstyle, language=Witness]
clause tesis189906_principio_fundamental = oblig(ejercer_prescripcion_positiva) IMPLIES oblig(causa_generadora_revelada) AND oblig(posesion_como_propietario);
\end{lstlisting}

\subsubsection{Clause: tesis189906\_determinacion\_judicial}

\textit{Supporting logic: Judge needs generating fact/act to determine possession quality and timing}

\textbf{Intermediate Representation:}

\begin{quote}
If causa generadora revelada is obligated, then (possession as owner is obligated and (good faith exercise is obligated and not (bad faith exercise is obligated)))
\end{quote}

\textbf{Witness Clause:}

\begin{lstlisting}[style=witnessstyle, language=Witness]
clause tesis189906_determinacion_judicial = oblig(causa_generadora_revelada) IMPLIES (oblig(posesion_como_propietario) AND (oblig(ejercicio_buena_fe) AND not(oblig(ejercicio_mala_fe))));
\end{lstlisting}

\subsubsection{Clause: tesis189906\_tipo\_posesion}

\textit{Possession type determination: Can determine if original or derived possession}

\textbf{Intermediate Representation:}

\begin{quote}
If causa generadora revelada is obligated, then (possession as owner is obligated and not (temporal possession is obligated))
\end{quote}

\textbf{Witness Clause:}

\begin{lstlisting}[style=witnessstyle, language=Witness]
clause tesis189906_tipo_posesion = oblig(causa_generadora_revelada) IMPLIES (oblig(posesion_como_propietario) AND not(oblig(posesion_temporal)));
\end{lstlisting}

\subsubsection{Clause: tesis189906\_determinacion\_temporal}

\textit{Timing requirement: Must establish when prescription period starts counting}

\textbf{Intermediate Representation:}

\begin{quote}
If causa generadora revelada is obligated, then (posesion tiempo 5 anos is obligated or posesion tiempo 10 anos is obligated)
\end{quote}

\textbf{Witness Clause:}

\begin{lstlisting}[style=witnessstyle, language=Witness]
clause tesis189906_determinacion_temporal = oblig(causa_generadora_revelada) IMPLIES (oblig(posesion_tiempo_5_anos) OR oblig(posesion_tiempo_10_anos));
\end{lstlisting}

\newpage

\section{Thesis 190437: USUCAPIÓN. EL ALLANAMIENTO A LA DEMANDA ES APTO PARA DEMOSTRAR LA POSESIÓN POR PARTE DEL ACTOR Y SUS DEMÁS ATRIBUTOS (LEGISLACIÓN DEL ESTADO}
\label{thesis:190437}

\subsection{Original Thesis Text}

\begin{quote}
                                                   Semanario Judicial de la Federación

                                                  Tesis

 Registro digital: 190437

 Instancia: Tribunales            Novena Época                          Materia(s): Civil
 Colegiados de Circuito

 Tesis: II.2o.C.258 C             Fuente: Semanario Judicial de la      Tipo: Aislada
                                  Federación y su Gaceta.
                                  Tomo XIII, Enero de 2001, página
                                  1810


USUCAPIÓN. EL ALLANAMIENTO A LA DEMANDA ES APTO PARA DEMOSTRAR LA
POSESIÓN POR PARTE DEL ACTOR Y SUS DEMÁS ATRIBUTOS (LEGISLACIÓN DEL ESTADO
DE MÉXICO).


En el ejercicio de la acción de usucapión, el demandado al producir su contestación a la demanda
puede allanarse a todas y cada una de las prestaciones reclamadas e incluso aceptar en su
totalidad los hechos en que se sustente tal acción, máxime si reconoce ser cierto que el inmueble
que el enjuiciante pretende prescribir adquisitivamente se encuentra inscrito a nombre de dicho
demandado en el Registro Público de la Propiedad; que entre aquél y éste existió el acto jurídico de
compraventa en orden con la fracción del inmueble en litigio, y asimismo, la fecha de celebración
del contrato, igual que el acto de entrega real y material de la posesión, es decir, que le fue
transmitido el dominio de ese predio, todo lo cual incuestionablemente revela la intención evidente
del enjuiciado y la aceptación de que el accionante ha poseído el inmueble desde la fecha del
referido contrato de compraventa, con los atributos inherentes a que ese hecho sea público,
pacífico, continuo, ininterrumpido y de buena fe. Consecuentemente, atento a tal confesión judicial
expresa de la parte demandada, derivada del allanamiento, debidamente ratificado, deviene
indiscutible y patente que el promovente ya no está obligado a rendir más pruebas para justificar la
posesión del inmueble para que proceda dicha usucapión en los términos a que se refieren los
artículos 801 y 932 del Código Civil para el Estado de México, en relación con el 620 del Código de
Procedimientos Civiles de la propia entidad federativa.

SEGUNDO TRIBUNAL COLEGIADO EN MATERIA CIVIL DEL SEGUNDO CIRCUITO.


Amparo directo 516/2000. Bartolomé Sánchez Vivar. 17 de octubre de 2000. Unanimidad de votos.
Ponente: Virgilio A. Solorio Campos. Secretario: Faustino García Astudillo.

Nota: Por ejecutoria del 21 de febrero de 2018, la Primera Sala declaró inexistente la contradicción
de tesis 245/2017 derivada de la denuncia de la que fue objeto el criterio contenido en esta tesis, al
estimarse que no son discrepantes los criterios materia de la denuncia respectiva, en virtud de que
el Segundo Tribunal Colegiado en Materia Civil del Segundo Circuito, al resolver los amparos
directos 551/2015, 142/2017 y 303/2017, superó el criterio sostenido en esta tesis.

Por ejecutoria del 3 de abril de 2019, la Primera Sala declaró inexistente la contradicción de tesis
179/2018 derivada de la denuncia de la que fue objeto el criterio contenido en esta tesis, al
estimarse que no son discrepantes los criterios materia de la denuncia respectiva.


                                             Pág. 1 de 2             Fecha de impresión 01/10/2025
  https://sjf2.scjn.gob.mx/detalle/tesis/190437
                                                  Semanario Judicial de la Federación

Este criterio fue interrumpido por la tesis II.2o.C.26 C (10a.), de título y subtítulo: “USUCAPIÓN. LA
RATIFICACIÓN DEL ALLANAMIENTO A LAS PRESTACIONES DEL ACTOR NO TIENE COMO
ALCANCE UN RECONOCIMIENTO PROPIO ACERCA DE QUE LA POSESIÓN SE HAYA
DESARROLLADO EN FORMA PACÍFICA, CONTINUA Y PÚBLICA –INTERRUPCIÓN DEL
CRITERIO SOSTENIDO EN LA TESIS AISLADA II.2o.C.258 C– (LEGISLACIÓN DEL ESTADO DE
MÉXICO).”, publicada en el Semanario Judicial de la Federación del viernes 10 de enero de 2020 a
las 10:11 horas y en la Gaceta del Semanario Judicial de la Federación, Décima Época, Libro 74,
Tomo III, enero de 2020, página 2730.




                                             Pág. 2 de 2             Fecha de impresión 01/10/2025
  https://sjf2.scjn.gob.mx/detalle/tesis/190437

\end{quote}

\subsection{Formalization}

\subsubsection{Clause: tesis190437\_principio\_fundamental}

\textit{Thesis-specific principle: When defendant admits possession without title for more than 10 years with required qualities, is bad faith possessor and property is eligible for prescription}

\textbf{Intermediate Representation:}

\begin{quote}
If (posesion\_como\_propietario AND posesion\_pacifico\_continuo\_publico AND posesion\_tiempo\_10\_anos AND not(causa\_generadora\_probada)), then (ejercicio\_mala\_fe AND oblig(ejercer\_prescripcion\_positiva))
\end{quote}

\textbf{Witness Clause:}

\begin{lstlisting}[style=witnessstyle, language=Witness]
clause tesis190437_principio_fundamental = (posesion_como_propietario AND posesion_pacifico_continuo_publico AND posesion_tiempo_10_anos AND not(causa_generadora_probada)) IMPLIES (ejercicio_mala_fe AND oblig(ejercer_prescripcion_positiva));
\end{lstlisting}

\subsubsection{Clause: tesis190437\_no\_prueba\_causa}

\textit{Supporting logic: Does not need to prove generating cause of possession in trial when no title exists}

\textbf{Intermediate Representation:}

\begin{quote}
(not(causa\_generadora\_probada) AND posesion\_tiempo\_10\_anos) IMPLIES oblig(ejercer\_prescripcion\_positiva)
\end{quote}

\textbf{Witness Clause:}

\begin{lstlisting}[style=witnessstyle, language=Witness]
clause tesis190437_no_prueba_causa = (not(causa_generadora_probada) AND posesion_tiempo_10_anos) IMPLIES oblig(ejercer_prescripcion_positiva);
\end{lstlisting}

\subsubsection{Clause: tesis190437\_definicion\_titulo}

\textit{Title definition: In state law, title is understood as the generating cause of possession}

\textbf{Intermediate Representation:}

\begin{quote}
not (causa\_generadora\_probada) IMPLIES not(titularidad\_registral)
\end{quote}

\textbf{Witness Clause:}

\begin{lstlisting}[style=witnessstyle, language=Witness]
clause tesis190437_definicion_titulo = not(causa_generadora_probada) IMPLIES not(titularidad_registral);
\end{lstlisting}

\subsubsection{Clause: tesis190437\_apto\_prescripcion}

\textit{Prescription eligibility: If 10 years have passed and other elements are met, property is eligible for prescription}

\textbf{Intermediate Representation:}

\begin{quote}
(posesion\_tiempo\_10\_anos AND posesion\_como\_propietario AND posesion\_pacifico\_continuo\_publico) IMPLIES oblig(ejercer\_prescripcion\_positiva)
\end{quote}

\textbf{Witness Clause:}

\begin{lstlisting}[style=witnessstyle, language=Witness]
clause tesis190437_apto_prescripcion = (posesion_tiempo_10_anos AND posesion_como_propietario AND posesion_pacifico_continuo_publico) IMPLIES oblig(ejercer_prescripcion_positiva);
\end{lstlisting}

\newpage

\section{Thesis 192067: PRESCRIPCIÓN ADQUISITIVA DE TIERRAS EJIDALES POR UN PARTICULAR. CUARTO TRIBUNAL COLEGIADO EN MATERIA ADMINISTRATIVA DEL PRIMER CIRCUITO.}
\label{thesis:192067}

\subsection{Original Thesis Text}

\begin{quote}
                                                   Semanario Judicial de la Federación

                                                  Tesis

 Registro digital: 192067

 Instancia: Tribunales            Novena Época                          Materia(s): Administrativa
 Colegiados de Circuito

 Tesis: I.4o.A.309 A              Fuente: Semanario Judicial de la   Tipo: Aislada
                                  Federación y su Gaceta.
                                  Tomo XI, Abril de 2000, página 982


PRESCRIPCIÓN ADQUISITIVA DE TIERRAS EJIDALES POR UN PARTICULAR.


Si el ahora quejoso acudió al Tribunal Superior Agrario invocando el derecho que tiene a prescribir a
su favor tierras pertenecientes a un núcleo ejidal, afirmando que reúne los requisitos de posesión a
que alude el artículo 48 de la vigente Ley Agraria, ya que ha poseído esas tierras en forma continua,
pacífica y pública por más de cinco años, en el caso no puede darse la institución que reclama,
porque el artículo en cuestión establece una posesión calificada para que opere la prescripción,
consistente en el hecho de que aquélla debe ser con el carácter de titular de derechos ejidales,
pues de esta forma se evita la segregación de las tierras pretendidas del núcleo ejidal; por tanto, no
es suficiente que se hubiera acreditado que la posesión se tuvo por más de cinco o de diez años,
según que ella fuera de buena o de mala fe, que hubiera sido en forma pública, pacífica y continua,
ya que si no lo hizo con aquel carácter la prescripción resulta improcedente.

CUARTO TRIBUNAL COLEGIADO EN MATERIA ADMINISTRATIVA DEL PRIMER CIRCUITO.


Amparo directo 1694/99. Martín Villanueva Franco. 11 de agosto de 1999. Unanimidad de votos.
Ponente: Hilario Bárcenas Chávez. Secretaria: Mariza Arellano Pompa.

Véase: Semanario Judicial de la Federación y su Gaceta, Novena Época, Tomo III, junio de 1996,
página 895, tesis XI.2o.6 A, de rubro: "PRESCRIPCIÓN ADQUISITIVA EN MATERIA AGRARIA.
POSESIÓN ‘EN CONCEPTO DE TITULAR DE DERECHOS DE EJIDATARIO’ PARA QUE
PROSPERE LA.".




                                             Pág. 1 de 1             Fecha de impresión 01/10/2025
  https://sjf2.scjn.gob.mx/detalle/tesis/192067

\end{quote}

\subsection{Formalization}

\subsubsection{Clause: tesis192067\_principio\_fundamental}

\textit{Thesis-specific principle: Good faith possessors require just title as generating cause of possession for adverse possession to operate}

\textbf{Intermediate Representation:}

\begin{quote}
good faith exercise is obligated and possession as owner is obligated
\end{quote}

\textbf{Witness Clause:}

\begin{lstlisting}[style=witnessstyle, language=Witness]
clause tesis192067_principio_fundamental = oblig(ejercicio_buena_fe) AND oblig(posesion_como_propietario) IMPLIES oblig(causa_generadora_probada);
\end{lstlisting}

\subsubsection{Clause: tesis192067\_revelar\_causa}

\textit{Supporting logic: Must reveal and fully prove the generating cause of possession}

\textbf{Intermediate Representation:}

\begin{quote}
ejercer prescripcion positiva is obligated
\end{quote}

\textbf{Witness Clause:}

\begin{lstlisting}[style=witnessstyle, language=Witness]
clause tesis192067_revelar_causa = oblig(ejercer_prescripcion_positiva) IMPLIES oblig(causa_generadora_revelada);
\end{lstlisting}

\subsubsection{Clause: tesis192067\_determinacion\_judicial}

\textit{Judicial determination: Judge needs generating fact/act to determine possession quality and timing}

\textbf{Intermediate Representation:}

\begin{quote}
If causa generadora revelada is obligated, then (possession as owner is obligated and (good faith exercise is obligated and not (bad faith exercise is obligated)))
\end{quote}

\textbf{Witness Clause:}

\begin{lstlisting}[style=witnessstyle, language=Witness]
clause tesis192067_determinacion_judicial = oblig(causa_generadora_revelada) IMPLIES (oblig(posesion_como_propietario) AND (oblig(ejercicio_buena_fe) AND not(oblig(ejercicio_mala_fe))));
\end{lstlisting}

\subsubsection{Clause: tesis192067\_determinacion\_temporal}

\textit{Timing requirement: Must establish when prescription period begins}

\textbf{Intermediate Representation:}

\begin{quote}
If causa generadora revelada is obligated, then (posesion tiempo 5 anos is obligated or posesion tiempo 10 anos is obligated)
\end{quote}

\textbf{Witness Clause:}

\begin{lstlisting}[style=witnessstyle, language=Witness]
clause tesis192067_determinacion_temporal = oblig(causa_generadora_revelada) IMPLIES (oblig(posesion_tiempo_5_anos) OR oblig(posesion_tiempo_10_anos));
\end{lstlisting}

\newpage

\section{Thesis 193673: PRESCRIPCIÓN POSITIVA. SE INICIA A PARTIR DE QUE SE POSEE EL INMUEBLE EN CONCEPTO DE DUEÑO Y SE CONSUMA EN EL MOMENTO EN QUE HA TRANSCURRIDO EL}
\label{thesis:193673}

\subsection{Original Thesis Text}

\begin{quote}
                                                   Semanario Judicial de la Federación

                                                  Tesis

 Registro digital: 193673

 Instancia: Tribunales            Novena Época                          Materia(s): Civil
 Colegiados de Circuito

 Tesis: I.5o.C.86 C               Fuente: Semanario Judicial de la      Tipo: Aislada
                                  Federación y su Gaceta.
                                  Tomo X, Julio de 1999, página 892


PRESCRIPCIÓN POSITIVA. SE INICIA A PARTIR DE QUE SE POSEE EL INMUEBLE EN
CONCEPTO DE DUEÑO Y SE CONSUMA EN EL MOMENTO EN QUE HA TRANSCURRIDO EL
TIEMPO NECESARIO EXIGIDO POR LA LEY; DE MODO QUE LA SENTENCIA QUE LA
DECLARA PROCEDENTE, SÓLO CONSOLIDA EN FORMA RETROACTIVA EL TÍTULO DE
PROPIEDAD.


De la interpretación armónica y sistemática de los artículos 1135 y 1156 del Código Civil para el
Distrito Federal, se obtiene que la prescripción adquisitiva o usucapión, es un medio de adquirir la
propiedad de un inmueble, por la posesión prolongada del mismo, durante el tiempo y bajo las
condiciones determinadas por la ley; de ahí también se arriba a la conclusión de que la prescripción
se inicia precisamente a partir de que el interesado entra a poseer el bien; de esa manera, si tal
prescripción no se interrumpe por las causas naturales o legales requeridas o, si no se le hace
cesar, entonces, se consuma al momento en que se ha cumplido el plazo de posesión exigido por el
ordenamiento jurídico, según el caso concreto (posesión de buena o mala fe). La necesidad de
promover la acción de prescripción o de oponerla como excepción en el juicio relativo, se hace
patente si se toma en cuenta que la usucapión, aun consumada, no surte efectos de pleno derecho,
pues para que esto sea así, es necesario que se ejerza vía acción o vía excepción; pero, debe
aclararse, la sentencia judicial que declara propietario por prescripción al poseedor de un bien, no
es la que consuma la usucapión, pues ésta se consuma por el solo transcurso del tiempo y, la
sentencia que así lo declara, sólo consolida el título de propiedad, al declarar procedente el derecho
prescrito a favor del interesado. Consecuentemente, la sentencia que determina que es propietario
por prescripción, el poseedor de un bien, surte efectos desde que la prescripción se inició, pues se
entiende que desde entonces se poseyó animus dominiis el bien prescrito. Eso es lo que algunos
tratadistas denominan "retroactividad de la prescripción".

QUINTO TRIBUNAL COLEGIADO EN MATERIA CIVIL DEL PRIMER CIRCUITO.


Amparo directo 750/99. José Carlos Méndez Solórzano. 20 de mayo de 1999. Unanimidad de votos.
Ponente: Arturo Ramírez Sánchez. Secretario: José Manuel Quistián Espericueta.




                                             Pág. 1 de 2             Fecha de impresión 01/10/2025
  https://sjf2.scjn.gob.mx/detalle/tesis/193673
                                                Semanario Judicial de la Federación




                                           Pág. 2 de 2       Fecha de impresión 01/10/2025
https://sjf2.scjn.gob.mx/detalle/tesis/193673

\end{quote}

\subsection{Formalization}

\subsubsection{Clause: tesis193673\_principio\_fundamental}

\textit{Thesis-specific principle: Adverse possession can be invoked by action or exception, but requires same proof elements}

\textbf{Intermediate Representation:}

\begin{quote}
possession as owner is obligated and (good faith exercise is obligated or bad faith exercise is obligated) and (oblig(posesion\_tiempo\_5\_anos) OR oblig(posesion\_tiempo\_10\_anos)) IMPLIES oblig(ejercer\_prescripcion\_positiva)
\end{quote}

\textbf{Witness Clause:}

\begin{lstlisting}[style=witnessstyle, language=Witness]
clause tesis193673_principio_fundamental = oblig(posesion_como_propietario) AND (oblig(ejercicio_buena_fe) OR oblig(ejercicio_mala_fe)) AND (oblig(posesion_tiempo_5_anos) OR oblig(posesion_tiempo_10_anos)) IMPLIES oblig(ejercer_prescripcion_positiva);
\end{lstlisting}

\subsubsection{Clause: tesis193673\_igualdad\_prueba}

\textit{Supporting logic: Both action and exception require same proof elements regardless of procedural route}

\textbf{Intermediate Representation:}

\begin{quote}
ejercer prescripcion positiva is obligated
\end{quote}

\textbf{Witness Clause:}

\begin{lstlisting}[style=witnessstyle, language=Witness]
clause tesis193673_igualdad_prueba = oblig(ejercer_prescripcion_positiva) IMPLIES oblig(posesion_como_propietario);
\end{lstlisting}

\subsubsection{Clause: tesis193673\_plazos\_temporales}

\textit{Time requirements: 5 years for good faith, 10 years for bad faith}

\textbf{Intermediate Representation:}

\begin{quote}
good faith exercise is obligated
\end{quote}

\textbf{Witness Clause:}

\begin{lstlisting}[style=witnessstyle, language=Witness]
clause tesis193673_plazos_temporales = oblig(ejercicio_buena_fe) IMPLIES oblig(posesion_tiempo_5_anos);
\end{lstlisting}

\subsubsection{Clause: tesis193673\_plazos\_mala\_fe}

\textbf{Intermediate Representation:}

\begin{quote}
bad faith exercise is obligated
\end{quote}

\textbf{Witness Clause:}

\begin{lstlisting}[style=witnessstyle, language=Witness]
clause tesis193673_plazos_mala_fe = oblig(ejercicio_mala_fe) IMPLIES oblig(posesion_tiempo_10_anos);
\end{lstlisting}

\subsubsection{Clause: tesis193673\_objetivo\_registral}

\textit{Target requirement: Action must be directed against registered owner}

\textbf{Intermediate Representation:}

\begin{quote}
ejercer prescripcion positiva is obligated
\end{quote}

\textbf{Witness Clause:}

\begin{lstlisting}[style=witnessstyle, language=Witness]
clause tesis193673_objetivo_registral = oblig(ejercer_prescripcion_positiva) IMPLIES oblig(titularidad_registral);
\end{lstlisting}

\newpage

\section{Thesis 196001: PRESCRIPCIÓN ADQUISITIVA Y REIVINDICACIÓN O RESTITUCIÓN RESPECTO DE UNA MISMA UNIDAD DE DOTACIÓN, SON ACCIONES CONTRADICTORIAS LAS DE.}
\label{thesis:196001}

\subsection{Original Thesis Text}

\begin{quote}
                                                   Semanario Judicial de la Federación

                                                  Tesis

 Registro digital: 196001

 Instancia: Tribunales            Novena Época                          Materia(s): Administrativa
 Colegiados de Circuito

 Tesis: XI.3o.9 A                 Fuente: Semanario Judicial de la      Tipo: Aislada
                                  Federación y su Gaceta.
                                  Tomo VII, Junio de 1998, página
                                  687


PRESCRIPCIÓN ADQUISITIVA Y REIVINDICACIÓN O RESTITUCIÓN RESPECTO DE UNA
MISMA UNIDAD DE DOTACIÓN, SON ACCIONES CONTRADICTORIAS LAS DE.


Las acciones de prescripción adquisitiva y reivindicación o restitución son incompatibles, pues de
acuerdo con el artículo 48 de la Ley Agraria, la primera compete al poseedor de una parcela ejidal y
parte de la base de que el actor o aspirante a ejidatario se encuentra en posesión de la tierra o
unidad de dotación en concepto de titular de ese derecho, en forma pacífica, continua y pública,
durante un periodo de cinco años si la posesión es de buena fe o diez si es de mala fe; y tiene la
finalidad de que quien ocupa la parcela en conflicto de simple poseedor se convierta en titular de los
derechos de la misma. En tanto que la acción de reivindicación o restitución corresponde al titular o
propietario de la parcela ejidal, que no está en posesión de la misma y tiene como propósito que se
declare que tiene dominio sobre la cosa que reclama, y el demandado se la entregue; por lo que el
actor debe probar, si es un núcleo de población, la propiedad de las tierras que reclama, y si es un
ejidatario la titularidad de la parcela que exige se le devuelva, la posesión por el demandado de la
cosa perseguida y la identidad de la misma, es decir, que no pueda dudarse cuál es la cosa que
pretende se le restituya y a la que se refieren los documentos fundatorios de la acción. Así pues, si
una misma persona demanda respecto de una misma unidad de dotación su prescripción y
reivindicación, ello implica que no la posee y pretende recuperarla, por lo que es claro que ejercita
acciones contradictorias que se destruyen entre sí, pues mientras la primera parte de la base de
que el actor no es el titular de la parcela, pero que la posee como tal, la segunda se basa en que sí
es titular y ha perdido la posesión, ya que no puede una persona al mismo tiempo ser titular de una
parcela y no poseerla y ocuparla y no ser titular de la misma.

TERCER TRIBUNAL COLEGIADO DEL DÉCIMO PRIMER CIRCUITO.


Amparo directo 710/97. Ignacio Gutiérrez Ceja. 12 de marzo de 1998. Unanimidad de votos.
Ponente: Moisés Duarte Aguíñiga. Secretario: Juan García Orozco.




                                             Pág. 1 de 2             Fecha de impresión 01/10/2025
  https://sjf2.scjn.gob.mx/detalle/tesis/196001
                                                Semanario Judicial de la Federación




                                           Pág. 2 de 2       Fecha de impresión 01/10/2025
https://sjf2.scjn.gob.mx/detalle/tesis/196001

\end{quote}

\subsection{Formalization}

\subsubsection{Clause: tesis196001\_principio\_fundamental}

\textit{Thesis-specific principle: Adverse possession can be invoked by action or exception, but requires same proof elements}

\textbf{Intermediate Representation:}

\begin{quote}
possession as owner is obligated and (good faith exercise is obligated or bad faith exercise is obligated) and (oblig(posesion\_tiempo\_5\_anos) OR oblig(posesion\_tiempo\_10\_anos)) IMPLIES oblig(ejercer\_prescripcion\_positiva)
\end{quote}

\textbf{Witness Clause:}

\begin{lstlisting}[style=witnessstyle, language=Witness]
clause tesis196001_principio_fundamental = oblig(posesion_como_propietario) AND (oblig(ejercicio_buena_fe) OR oblig(ejercicio_mala_fe)) AND (oblig(posesion_tiempo_5_anos) OR oblig(posesion_tiempo_10_anos)) IMPLIES oblig(ejercer_prescripcion_positiva);
\end{lstlisting}

\subsubsection{Clause: tesis196001\_igualdad\_prueba}

\textit{Supporting logic: Both action and exception require same proof elements regardless of procedural route}

\textbf{Intermediate Representation:}

\begin{quote}
ejercer prescripcion positiva is obligated
\end{quote}

\textbf{Witness Clause:}

\begin{lstlisting}[style=witnessstyle, language=Witness]
clause tesis196001_igualdad_prueba = oblig(ejercer_prescripcion_positiva) IMPLIES oblig(posesion_como_propietario);
\end{lstlisting}

\subsubsection{Clause: tesis196001\_plazos\_temporales}

\textit{Time requirements: 5 years for good faith, 10 years for bad faith}

\textbf{Intermediate Representation:}

\begin{quote}
good faith exercise is obligated
\end{quote}

\textbf{Witness Clause:}

\begin{lstlisting}[style=witnessstyle, language=Witness]
clause tesis196001_plazos_temporales = oblig(ejercicio_buena_fe) IMPLIES oblig(posesion_tiempo_5_anos);
\end{lstlisting}

\subsubsection{Clause: tesis196001\_plazos\_mala\_fe}

\textbf{Intermediate Representation:}

\begin{quote}
bad faith exercise is obligated
\end{quote}

\textbf{Witness Clause:}

\begin{lstlisting}[style=witnessstyle, language=Witness]
clause tesis196001_plazos_mala_fe = oblig(ejercicio_mala_fe) IMPLIES oblig(posesion_tiempo_10_anos);
\end{lstlisting}

\subsubsection{Clause: tesis196001\_objetivo\_registral}

\textit{Target requirement: Action must be directed against registered owner}

\textbf{Intermediate Representation:}

\begin{quote}
ejercer prescripcion positiva is obligated
\end{quote}

\textbf{Witness Clause:}

\begin{lstlisting}[style=witnessstyle, language=Witness]
clause tesis196001_objetivo_registral = oblig(ejercer_prescripcion_positiva) IMPLIES oblig(titularidad_registral);
\end{lstlisting}

\newpage

\section{Thesis 199535: PRESCRIPCION ADQUISITIVA. ES UNA INSTITUCION DE DERECHO CIVIL Y NO AGRARIA, POR LO QUE EL ACCIONANTE DEBE AJUSTARSE A LOS REQUISITOS PREVISTOS EN EL}
\label{thesis:199535}

\subsection{Original Thesis Text}

\begin{quote}
                                                   Semanario Judicial de la Federación

                                                  Tesis

 Registro digital: 199535

 Instancia: Tribunales             Novena Época                           Materia(s): Civil
 Colegiados de Circuito

 Tesis: XV.2o. J/1                 Fuente: Semanario Judicial de la  Tipo: Jurisprudencia
                                   Federación y su Gaceta.
                                   Tomo V, Enero de 1997, página 320


PRESCRIPCION ADQUISITIVA. ES UNA INSTITUCION DE DERECHO CIVIL Y NO AGRARIA,
POR LO QUE EL ACCIONANTE DEBE AJUSTARSE A LOS REQUISITOS PREVISTOS EN EL
CODIGO CIVIL Y A LA JURISPRUDENCIA DEFINIDA AL RESPECTO.


Si bien es cierto que no existe disposición expresa en el Código Civil para el Estado de Baja
California que obligue a manifestar solemnemente en la demanda de prescripción positiva, la frase
"causa generadora de la posesión", la Suprema Corte de Justicia en la jurisprudencia por
contradicción de tesis número 18/94, visible en la página 1235, del Tomo IV, Tercera Sala, Segunda
Parte, Octava Epoca, de la Jurisprudencia por Contradicción de Tesis, se pronunció al respecto bajo
el rubro: "PRESCRIPCION ADQUISITIVA. PARA QUE SE ENTIENDA SATISFECHO EL
REQUISITO DE LA EXISTENCIA DE LA 'POSESION EN CONCEPTO DE PROPIETARIO'
EXIGIDO POR EL CODIGO CIVIL PARA EL DISTRITO FEDERAL Y POR LAS DIVERSAS
LEGISLACIONES DE LOS ESTADOS DE LA REPUBLICA QUE CONTIENEN DISPOSICIONES
IGUALES, ES NECESARIO DEMOSTRAR LA EXISTENCIA DE UN TITULO DEL QUE SE DERIVE
LA POSESION". Siendo así, el hecho de que se hayan manifestado, la fecha a partir de la cual
entró a poseer, considerándose el accionante subjetivamente como propietario, y la fecha en que se
solicitó al gobernador del Estado la tenencia de la tierra, como lo exigía la Ley Federal de Reforma
Agraria, vigente en la época de los acontecimientos, resulta insuficiente atendiendo a que la
prescripción adquisitiva es una institución de derecho civil, por tanto, de derecho estricto, de ahí que
para la procedencia de la usucapión se hace necesaria la prueba objetiva del origen o causa
generadora de la posesión que se narre en la demanda, como sería la existencia de determinado
acto traslativo de dominio, verbal o escrito, que produzca consecuencias de derecho y que legitime
al poseedor para comportarse ostensible y objetivamente como propietario del inmueble sobre del
cual ha realizado actos que revelan dominio o mandato; la precisión y prueba de dicha causa
generadora en los términos apuntados, permite establecer si la posesión es en concepto de
propietario, originaria o derivada, de buena o mala fe y determina la calidad y naturaleza de la
posesión, por tanto, el momento en que debe empezar a contar el plazo de la prescripción positiva.

SEGUNDO TRIBUNAL COLEGIADO DEL DECIMO QUINTO CIRCUITO.


Amparo directo 91/96. Antonio Armenta Lomelí. 15 de mayo de 1996. Unanimidad de votos.
Ponente: Adán Gilberto Villarreal Castro. Secretaria: Magaly Herrera Olaiz.

Amparo directo 95/96. Antonio Armenta Lomelí y otro. 15 de mayo de 1996. Unanimidad de votos.
Ponente: Adán Gilberto Villarreal Castro. Secretaria: Edith Ríos Torres.



                                              Pág. 1 de 2             Fecha de impresión 01/10/2025
  https://sjf2.scjn.gob.mx/detalle/tesis/199535
                                                  Semanario Judicial de la Federación

Amparo directo 103/96. Grupo Solicitante de Tierras Jardines del Rincón. 15 de mayo de 1996.
Unanimidad de votos. Ponente: Sergio Javier Coss Ramos. Secretaria: Nora Laura Gómez
Castellanos.

Amparo directo 89/96. Grupo Solicitante de Tierras Jardines del Rincón. 16 de mayo de 1996.
Unanimidad de votos. Ponente: Carlos Humberto Trujillo Altamirano. Secretario: José Neals André
Nalda.

Amparo directo 293/96. Antonio Armenta Lomelí. 26 de junio de 1996. Unanimidad de votos.
Ponente: Sergio Javier Coss Ramos. Secretario: Joaquín Gallegos Flores.




                                             Pág. 2 de 2        Fecha de impresión 01/10/2025
  https://sjf2.scjn.gob.mx/detalle/tesis/199535

\end{quote}

\subsection{Formalization}

\subsubsection{Clause: tesis199535\_principio\_fundamental}

\textit{Additional actions for this thesis Additional assets for proof requirements Faith exclusivity: Adverse possession is either good faith XOR bad faith Thesis-specific principle: Good faith adverse possession requires proving authenticity, date, and diligence}

\textbf{Intermediate Representation:}

\begin{quote}
If good faith exercise is obligated and adverse possession claim is claimed, then (autenticidad titulo is obligated and fecha fehaciente is obligated and diligencia adquirente is obligated)
\end{quote}

\textbf{Witness Clause:}

\begin{lstlisting}[style=witnessstyle, language=Witness]
clause tesis199535_principio_fundamental = oblig(ejercicio_buena_fe) AND claim(demanda_prescripcion) IMPLIES (oblig(autenticidad_titulo) AND oblig(fecha_fehaciente) AND oblig(diligencia_adquirente));
\end{lstlisting}

\subsubsection{Clause: tesis199535\_prueba\_autenticidad}

\textit{Proof requirements: Must prove title authenticity and date reliably}

\textbf{Intermediate Representation:}

\begin{quote}
If proven generating cause is obligated, then (autenticidad titulo is obligated and fecha fehaciente is obligated)
\end{quote}

\textbf{Witness Clause:}

\begin{lstlisting}[style=witnessstyle, language=Witness]
clause tesis199535_prueba_autenticidad = oblig(causa_generadora_probada) IMPLIES (oblig(autenticidad_titulo) AND oblig(fecha_fehaciente));
\end{lstlisting}

\subsubsection{Clause: tesis199535\_diligencia\_buena\_fe}

\textit{Diligence requirement: Good faith requires demonstrating due diligence in acquisition}

\textbf{Intermediate Representation:}

\begin{quote}
good faith exercise is obligated
\end{quote}

\textbf{Witness Clause:}

\begin{lstlisting}[style=witnessstyle, language=Witness]
clause tesis199535_diligencia_buena_fe = oblig(ejercicio_buena_fe) IMPLIES oblig(diligencia_adquirente);
\end{lstlisting}

\subsubsection{Clause: tesis199535\_prueba\_pagos}

\textit{Payment proof: Onerous contracts require proving payments made}

\textbf{Intermediate Representation:}

\begin{quote}
proven generating cause is obligated and contrato compraventa is obligated
\end{quote}

\textbf{Witness Clause:}

\begin{lstlisting}[style=witnessstyle, language=Witness]
clause tesis199535_prueba_pagos = oblig(causa_generadora_probada) AND oblig(contrato_compraventa) IMPLIES oblig(pagos_precio);
\end{lstlisting}

\subsubsection{Clause: tesis199535\_carga\_prueba}

\textit{Burden of proof: Claimant must prove generating cause with reliable evidence}

\textbf{Intermediate Representation:}

\begin{quote}
adverse possession claim is claimed
\end{quote}

\textbf{Witness Clause:}

\begin{lstlisting}[style=witnessstyle, language=Witness]
clause tesis199535_carga_prueba = claim(demanda_prescripcion) IMPLIES oblig(causa_generadora_revelada);
\end{lstlisting}

\subsubsection{Clause: tesis199535\_falta\_diligencia}

\textit{Insufficient diligence consequence: Lack of diligence means bad faith (10-year period)}

\textbf{Intermediate Representation:}

\begin{quote}
not (diligencia adquirente is obligated)
\end{quote}

\textbf{Witness Clause:}

\begin{lstlisting}[style=witnessstyle, language=Witness]
clause tesis199535_falta_diligencia = not(oblig(diligencia_adquirente)) IMPLIES oblig(ejercicio_mala_fe);
\end{lstlisting}

\newpage

\section{Thesis 199950: PRESCRIPCION ADQUISITIVA. LA BUENA FE DEBE PROBARSE (LEGISLACION DEL ESTADO PRIMER TRIBUNAL COLEGIADO EN MATERIAS CIVIL Y DE TRABAJO DEL SEGUNDO}
\label{thesis:199950}

\subsection{Original Thesis Text}

\begin{quote}
                                                   Semanario Judicial de la Federación

                                                  Tesis

 Registro digital: 199950

 Instancia: Tribunales            Novena Época                       Materia(s): Civil
 Colegiados de Circuito

 Tesis: II.1o.C.T.96 C            Fuente: Semanario Judicial de la   Tipo: Aislada
                                  Federación y su Gaceta.
                                  Tomo IV, Diciembre de 1996, página
                                  436


PRESCRIPCION ADQUISITIVA. LA BUENA FE DEBE PROBARSE (LEGISLACION DEL ESTADO
DE MEXICO).


Si el actor en la usucapión se ostenta como poseedor de buena fe, al ser un elemento de la acción,
debe probarlo, como lo ordena el artículo 269 del Código de Procedimientos Civiles del Estado de
México, aun cuando no hubiere controvertido el demandado la posesión como de mala fe.

PRIMER TRIBUNAL COLEGIADO EN MATERIAS CIVIL Y DE TRABAJO DEL SEGUNDO
CIRCUITO.


Amparo directo 789/95. Carmen Culebro viuda de Hernández, en su carácter de albacea de la
sucesión intestamentaria a bienes de Ernesto Hernández Pantoja. 1o. de agosto de 1996.
Unanimidad de votos. Ponente: Salvador Bravo Gómez. Secretario: José Fernando García Quiroz.




                                             Pág. 1 de 1          Fecha de impresión 01/10/2025
  https://sjf2.scjn.gob.mx/detalle/tesis/199950

\end{quote}

\subsection{Formalization}

\subsubsection{Clause: tesis199950\_principio\_fundamental}

\textit{Faith exclusivity: Possession is either good faith XOR bad faith (mutually exclusive) Thesis-specific principle: Unregistered contracts cannot establish adverse possession against registered owners}

\textbf{Intermediate Representation:}

\begin{quote}
registered ownership is obligated and unregistered contract is obligated
\end{quote}

\textbf{Witness Clause:}

\begin{lstlisting}[style=witnessstyle, language=Witness]
clause tesis199950_principio_fundamental = oblig(titularidad_registral) AND oblig(contrato_no_inscrito) IMPLIES not(claim(demanda_prescripcion));
\end{lstlisting}

\subsubsection{Clause: tesis199950\_proteccion\_registral}

\textit{Supporting legal logic}

\textbf{Intermediate Representation:}

\begin{quote}
registered ownership is obligated
\end{quote}

\textbf{Witness Clause:}

\begin{lstlisting}[style=witnessstyle, language=Witness]
clause tesis199950_proteccion_registral = oblig(titularidad_registral) IMPLIES claim(oponibilidad_terceros);
\end{lstlisting}

\subsubsection{Clause: tesis199950\_limitacion\_no\_inscripcion}

\textbf{Intermediate Representation:}

\begin{quote}
unregistered contract is obligated
\end{quote}

\textbf{Witness Clause:}

\begin{lstlisting}[style=witnessstyle, language=Witness]
clause tesis199950_limitacion_no_inscripcion = oblig(contrato_no_inscrito) IMPLIES not(oblig(oponibilidad_terceros));
\end{lstlisting}

\subsubsection{Clause: tesis199950\_consecuencia\_logica}

\textbf{Intermediate Representation:}

\begin{quote}
not (opposability to third parties is obligated)
\end{quote}

\textbf{Witness Clause:}

\begin{lstlisting}[style=witnessstyle, language=Witness]
clause tesis199950_consecuencia_logica = not(oblig(oponibilidad_terceros)) IMPLIES not(claim(demanda_prescripcion));
\end{lstlisting}

\newpage

\section{Thesis 202177: AGRARIAN ADVERSE POSSESSION REQUIREMENTS}
\label{thesis:202177}

\subsection{Original Thesis Text}

\begin{quote}
                                                   Semanario Judicial de la Federación

                                                  Tesis

 Registro digital: 202177

 Instancia: Tribunales            Novena Época                          Materia(s): Administrativa
 Colegiados de Circuito

 Tesis: XI.2o.6 A                 Fuente: Semanario Judicial de la    Tipo: Aislada
                                  Federación y su Gaceta.
                                  Tomo III, Junio de 1996, página 895


PRESCRIPCION ADQUISITIVA EN MATERIA AGRARIA. POSESION "EN CONCEPTO DE
TITULAR DE DERECHOS DE EJIDATARIO" PARA QUE PROSPERE LA.


El artículo 48 de la Ley Agraria en vigor dispone, entre otras cosas, que quien haya poseído tierras
ejidales "en concepto de titular de derechos de ejidatario", que no sean las destinadas al
asentamiento humano ni se trate de bosques o selvas, de manera pacífica, continua y pública
durante un período de cinco años si la posesión es de buena fe o de diez si fuera de mala fe,
adquirirá sobre dichas tierras los mismos derechos que cualquier ejidatario sobre su parcela. Luego,
para que se entienda satisfecha la posesión "en concepto de titular de derechos de ejidatario", que
como requisito para la usucapión contempla el apuntado normativo, es menester que el actor revele
(y demuestre) la causa generadora de su posesión; entendiéndose por tal, cualquier acto que
fundadamente se considere bastante para transferirle el dominio sobre la unidad de dotación de que
se trate, a fin de estar en aptitud de dilucidar primeramente si en realidad su posesión es o no "en
concepto de titular de derechos de ejidatario", para en seguida determinar lo referente a los plazos
legales en que habrá de operar la prescripción, según sea el caso de posesión de buena o mala fe.
Lo anterior, porque una recta interpretación del precepto en comento, conlleva a establecer que
dicha institución se encuentra reservada para aquellos cuya posesión de tierras ejidales sea de
naturaleza originaria (en eso se traduce la connotación de "en concepto de titular de derechos de
ejidatario"), y no para los que las detentan de manera precaria o derivada.

SEGUNDO TRIBUNAL COLEGIADO DEL DECIMO PRIMER CIRCUITO.


Amparo directo 466/95. Amalia Fernández Badajosa. 10 de agosto de 1995. Unanimidad de votos.
Ponente: Juan Díaz Ponce de León. Secretario: Gilberto Díaz Ortiz.

Nota: Este criterio ha integrado la jurisprudencia XI.2o. J/37, publicada en el Semanario Judicial de
la Federación y su Gaceta, Novena Época, Tomo XXVIII, octubre de 2008, página 2270, con el
rubro: "PRESCRIPCIÓN ADQUISITIVA EN MATERIA AGRARIA. PARA QUE OPERE ES
NECESARIO QUE EL ACTOR REVELE Y DEMUESTRE LA CAUSA GENERADORA DE LA
POSESIÓN, QUE DEBE SER DE NATURALEZA ORIGINARIA Y NO PRECARIA O DERIVADA."




                                             Pág. 1 de 2             Fecha de impresión 01/10/2025
  https://sjf2.scjn.gob.mx/detalle/tesis/202177
                                                Semanario Judicial de la Federación




                                           Pág. 2 de 2       Fecha de impresión 01/10/2025
https://sjf2.scjn.gob.mx/detalle/tesis/202177

\end{quote}

\subsection{Formalization}

\subsubsection{Clause: tesis202177\_principio\_fundamental}

\textit{Thesis-specific principle: Agrarian adverse possession requires original possession "in concept of ejidatario rights holder" with proven generating cause}

\textbf{Intermediate Representation:}

\begin{quote}
possession as owner is obligated and proven generating cause is obligated and posesion agraria is obligated
\end{quote}

\textbf{Witness Clause:}

\begin{lstlisting}[style=witnessstyle, language=Witness]
clause tesis202177_principio_fundamental = oblig(posesion_como_propietario) AND oblig(causa_generadora_probada) AND oblig(posesion_agraria) IMPLIES oblig(ejercer_prescripcion_positiva);
\end{lstlisting}

\subsubsection{Clause: tesis202177\_posesion\_originaria}

\textit{Supporting legal logic: Original possession excludes precarious or derived possession}

\textbf{Intermediate Representation:}

\begin{quote}
possession as owner is obligated and not (temporal possession is obligated)
\end{quote}

\textbf{Witness Clause:}

\begin{lstlisting}[style=witnessstyle, language=Witness]
clause tesis202177_posesion_originaria = oblig(posesion_como_propietario) AND not(oblig(posesion_temporal)) IMPLIES oblig(causa_generadora_probada);
\end{lstlisting}

\subsubsection{Clause: tesis202177\_plazos\_temporales}

\textit{Time periods: 5 years for good faith, 10 years for bad faith}

\textbf{Intermediate Representation:}

\begin{quote}
good faith exercise is obligated
\end{quote}

\textbf{Witness Clause:}

\begin{lstlisting}[style=witnessstyle, language=Witness]
clause tesis202177_plazos_temporales = oblig(ejercicio_buena_fe) IMPLIES oblig(posesion_tiempo_5_anos);
\end{lstlisting}

\subsubsection{Clause: tesis202177\_plazos\_mala\_fe}

\textbf{Intermediate Representation:}

\begin{quote}
bad faith exercise is obligated
\end{quote}

\textbf{Witness Clause:}

\begin{lstlisting}[style=witnessstyle, language=Witness]
clause tesis202177_plazos_mala_fe = oblig(ejercicio_mala_fe) IMPLIES oblig(posesion_tiempo_10_anos);
\end{lstlisting}

\newpage

\section{Thesis 202535: BAD FAITH POSSESSOR IN INMATRICULATION PROCEDURES}
\label{thesis:202535}

\subsection{Original Thesis Text}

\begin{quote}
                                                   Semanario Judicial de la Federación

                                                  Tesis

 Registro digital: 202535

 Instancia: Tribunales            Novena Época                        Materia(s): Civil
 Colegiados de Circuito

 Tesis: II.1o.C.T.44 C            Fuente: Semanario Judicial de la   Tipo: Aislada
                                  Federación y su Gaceta.
                                  Tomo III, Mayo de 1996, página 715


USUCAPION; ES TENEDOR DE MALA FE QUIEN PROMUEVE PROCEDIMIENTO DE
INMATRICULACION, A SABIENDAS DE QUE EL PROPIETARIO ES PERSONA DISTINTA DE
QUIEN LE TRANSMITIO EL DOMINIO DEL INMUEBLE.


No es factible considerar tenedor de buena fe, en el juicio de usucapión, a quien para acreditarla,
promueve un procedimiento de inmatriculación, conociendo que el propietario es una persona
distinta de quien le transmitió el dominio del inmueble.

PRIMER TRIBUNAL COLEGIADO EN MATERIAS CIVIL Y DE TRABAJO DEL SEGUNDO
CIRCUITO.


Amparo directo 992/95. María Dolores Bañuelos Sánchez. 31 de octubre de 1995. Unanimidad de
votos. Ponente: Salvador Bravo Gómez. Secretario: José Fernando García Quiroz.




                                             Pág. 1 de 1           Fecha de impresión 01/10/2025
  https://sjf2.scjn.gob.mx/detalle/tesis/202535

\end{quote}

\subsection{Formalization}

\subsubsection{Clause: tesis202535\_principio\_fundamental}

\textit{Thesis-specific asset: Inmatriculation procedure for this specific case Thesis-specific principle: Person promoting inmatriculation knowing true owner is different from transferor cannot be good faith possessor}

\textbf{Intermediate Representation:}

\begin{quote}
conocimiento vicios is obligated and procedimiento inmatriculacion 202535 is obligated
\end{quote}

\textbf{Witness Clause:}

\begin{lstlisting}[style=witnessstyle, language=Witness]
clause tesis202535_principio_fundamental = oblig(conocimiento_vicios) AND oblig(procedimiento_inmatriculacion_202535) IMPLIES oblig(ejercicio_mala_fe);
\end{lstlisting}

\subsubsection{Clause: tesis202535\_conocimiento\_defectos}

\textit{Supporting legal logic: Knowledge of true ownership defects precludes good faith}

\textbf{Intermediate Representation:}

\begin{quote}
conocimiento vicios is obligated
\end{quote}

\textbf{Witness Clause:}

\begin{lstlisting}[style=witnessstyle, language=Witness]
clause tesis202535_conocimiento_defectos = oblig(conocimiento_vicios) IMPLIES not(oblig(ejercicio_buena_fe));
\end{lstlisting}

\newpage

\section{Thesis 204300: AGRARIAN ADVERSE POSSESSION TIMING REQUIREMENTS}
\label{thesis:204300}

\subsection{Original Thesis Text}

\begin{quote}
                                                   Semanario Judicial de la Federación

                                                  Tesis

 Registro digital: 204300

 Instancia: Tribunales            Novena Época                         Materia(s): Administrativa
 Colegiados de Circuito

 Tesis: VI.2o.17 A                Fuente: Semanario Judicial de la     Tipo: Aislada
                                  Federación y su Gaceta.
                                  Tomo II, Septiembre de 1995,
                                  página 589


PRESCRIPCION ADQUISITIVA AGRARIA. TERMINO PARA QUE OPERE, DEBE TOMARSE EN
CUENTA LA POSESION QUE SE DETENTA A PARTIR DE LA VIGENCIA DE LA LEY AGRARIA
CUANDO EXISTEN EJIDATARIOS O COMUNEROS CON DERECHOS AGRARIOS.


De lo establecido en el artículo 48 de la Ley Agraria, se desprende que para que la acción de
prescripción adquisitiva opere, los términos aptos para ello (cinco años cuando la posesión ha sido
de buena fe y diez años cuando ha sido de mala fe), deben empezar a contar a partir de la vigencia
de dicho ordenamiento legal, y no de la posesión que se tenía con anterioridad a ella, toda vez que
la abrogada Ley Federal de Reforma Agraria no contemplaba esta figura jurídica como medio para
adquirir derechos agrarios reconocidos en favor de ejidatarios o comuneros.

SEGUNDO TRIBUNAL COLEGIADO DEL SEXTO CIRCUITO.


Amparo directo 299/95. Pedro Tlacotzi Ayapantécatl. 16 de agosto de 1995. Unanimidad de votos.
Ponente: María Eugenia Estela Martínez Cardiel. Secretario: Enrique Baigts Muñoz.




                                             Pág. 1 de 1             Fecha de impresión 01/10/2025
  https://sjf2.scjn.gob.mx/detalle/tesis/204300

\end{quote}

\subsection{Formalization}

\subsubsection{Clause: tesis204300\_principio\_fundamental}

\textit{Thesis-specific principle: Agrarian adverse possession periods start from Agrarian Law effective date, not from previous possession}

\textbf{Intermediate Representation:}

\begin{quote}
posesion agraria is obligated and posesion tiempo 5 anos is obligated and posesion tiempo 10 anos is obligated
\end{quote}

\textbf{Witness Clause:}

\begin{lstlisting}[style=witnessstyle, language=Witness]
clause tesis204300_principio_fundamental = oblig(posesion_agraria) AND oblig(posesion_tiempo_5_anos) AND oblig(posesion_tiempo_10_anos) IMPLIES oblig(ejercer_prescripcion_positiva);
\end{lstlisting}

\subsubsection{Clause: tesis204300\_posesion\_anterior}

\textit{Supporting legal logic: Previous possession under old law doesn't count toward prescription periods}

\textbf{Intermediate Representation:}

\begin{quote}
temporal possession is obligated and not (posesion agraria is obligated)
\end{quote}

\textbf{Witness Clause:}

\begin{lstlisting}[style=witnessstyle, language=Witness]
clause tesis204300_posesion_anterior = oblig(posesion_temporal) AND not(oblig(posesion_agraria)) IMPLIES not(oblig(posesion_tiempo_5_anos));
\end{lstlisting}

\subsubsection{Clause: tesis204300\_marco\_legal\_nuevo}

\textit{Legal framework change: New law recognizes adverse possession as means to acquire agrarian rights}

\textbf{Intermediate Representation:}

\begin{quote}
posesion agraria is obligated
\end{quote}

\textbf{Witness Clause:}

\begin{lstlisting}[style=witnessstyle, language=Witness]
clause tesis204300_marco_legal_nuevo = oblig(posesion_agraria) IMPLIES oblig(ejercer_prescripcion_positiva);
\end{lstlisting}

\newpage

\section{Thesis 204375: JUST TITLE REQUIREMENT FOR ADVERSE POSSESSION}
\label{thesis:204375}

\subsection{Original Thesis Text}

\begin{quote}
                                                   Semanario Judicial de la Federación

                                                  Tesis

 Registro digital: 204375

 Instancia: Tribunales             Novena Época                          Materia(s): Civil
 Colegiados de Circuito

 Tesis: III.1o.C. J/6              Fuente: Semanario Judicial de la      Tipo: Jurisprudencia
                                   Federación y su Gaceta.
                                   Tomo II, Septiembre de 1995,
                                   página 475


PRESCRIPCION ADQUISITIVA. JUSTO TITULO.


El justo título, aun cuando no en todos los casos es absolutamente necesario para prescribir, no ha
sido desterrado del Código Civil del Estado de Jalisco, pues a él corresponden las nociones de título
objetiva o subjetivamente válido a que se hace referencia en el artículo 849, en la medida en que
previene, en lo que interesa, que es poseedor de buena fe el que entra en la posesión en virtud de
un título suficiente para darle derecho a poseer y el que ignora los vicios de su título que le impiden
poseer con derecho. Por tanto, cuando se invoca como causa de la posesión, por tratarse de un
supuesto privilegiado para usucapir, es necesario acreditarlo, y no solamente revelar el origen de la
posesión y afirmar que se posee a título de dueño. De no ser así, el Juez estaría imposibilitado para
establecer si es en concepto de propietario, originaria o derivada, de buena o de mala fe y a partir
de qué momento debe contarse el plazo para usucapir.

PRIMER TRIBUNAL COLEGIADO EN MATERIA CIVIL DEL TERCER CIRCUITO.


Amparo directo 840/88. Carlota Reynoso Castillo. 3 de febrero de 1989. Unanimidad de votos.
Ponente: Carlos Arturo González Zárate. Secretario: Juan Bonilla Pizano.

Amparo directo 227/91. José Pagua Montaño y otra. 28 de junio de 1991. Unanimidad de votos.
Ponente: Carlos Arturo González Zárate. Secretario: Juan Bonilla Pizano.

Amparo directo 531/91. Margarito Ramos Lomelí. 17 de octubre de 1991. Unanimidad de votos.
Ponente: José de Jesús Gudiño Pelayo. Secretario: Simón Daniel Canales Aguiar.

Amparo directo 107/94. José Guadalupe Bejarano Casillas. 13 de mayo de 1994. Unanimidad de
votos. Ponente: Carlos Arturo González Zárate. Secretario: José Luis Fernández Jaramillo.

Amparo directo 44/95. José de Jesús Hernández Olivares. 18 de mayo de 1995. Unanimidad de
votos. Ponente: Héctor Soto Gallardo. Secretario: Simón Daniel Canales Aguiar.




                                              Pág. 1 de 2             Fecha de impresión 01/10/2025
  https://sjf2.scjn.gob.mx/detalle/tesis/204375
                                                Semanario Judicial de la Federación




                                           Pág. 2 de 2       Fecha de impresión 01/10/2025
https://sjf2.scjn.gob.mx/detalle/tesis/204375

\end{quote}

\subsection{Formalization}

\subsubsection{Clause: tesis204375\_principio\_fundamental}

\textit{Thesis-specific principle: Just title must be proven when invoked as cause of possession, not just revealed}

\textbf{Intermediate Representation:}

\begin{quote}
proven generating cause is obligated and conocimiento vicios is obligated
\end{quote}

\textbf{Witness Clause:}

\begin{lstlisting}[style=witnessstyle, language=Witness]
clause tesis204375_principio_fundamental = oblig(causa_generadora_probada) AND oblig(conocimiento_vicios) IMPLIES oblig(ejercicio_buena_fe);
\end{lstlisting}

\subsubsection{Clause: tesis204375\_titulo\_suficiente}

\textit{Supporting legal logic: Good faith possessor enters possession by virtue of sufficient title}

\textbf{Intermediate Representation:}

\begin{quote}
possession as owner is obligated and proven generating cause is obligated
\end{quote}

\textbf{Witness Clause:}

\begin{lstlisting}[style=witnessstyle, language=Witness]
clause tesis204375_titulo_suficiente = oblig(posesion_como_propietario) AND oblig(causa_generadora_probada) IMPLIES oblig(ejercicio_buena_fe);
\end{lstlisting}

\subsubsection{Clause: tesis204375\_validez\_titulo}

\textit{Title validity requirement: Title must be objectively or subjectively valid}

\textbf{Intermediate Representation:}

\begin{quote}
proven generating cause is obligated and not (conocimiento vicios is obligated)
\end{quote}

\textbf{Witness Clause:}

\begin{lstlisting}[style=witnessstyle, language=Witness]
clause tesis204375_validez_titulo = oblig(causa_generadora_probada) AND not(oblig(conocimiento_vicios)) IMPLIES oblig(ejercicio_buena_fe);
\end{lstlisting}

\subsubsection{Clause: tesis204375\_determinacion\_judicial}

\textit{Judicial determination: Judge must establish possession type and timing}

\textbf{Intermediate Representation:}

\begin{quote}
possession as owner is obligated
\end{quote}

\textbf{Witness Clause:}

\begin{lstlisting}[style=witnessstyle, language=Witness]
clause tesis204375_determinacion_judicial = oblig(posesion_como_propietario) IMPLIES oblig(causa_generadora_probada);
\end{lstlisting}

\newpage

\section{Thesis 206322: AGRARIAN ADVERSE POSSESSION TIMING FROM LAW EFFECTIVE DATE}
\label{thesis:206322}

\subsection{Original Thesis Text}

\begin{quote}
                                                   Semanario Judicial de la Federación

                                                  Tesis

 Registro digital: 206322

 Instancia: Segunda Sala           Octava Época                           Materia(s): Administrativa

 Tesis: 2a./J. 24/94               Fuente: Gaceta del Semanario       Tipo: Jurisprudencia
                                   Judicial de la Federación.
                                   Núm. 84, Diciembre de 1994, página
                                   20


PRESCRIPCION POSITIVA EN MATERIA AGRARIA, PARA COMPUTAR LOS PLAZOS
PREVISTOS EN EL ARTICULO 48 DE LA LEY AGRARIA, NO DEBE TOMARSE EN CUENTA EL
TIEMPO DE POSESION ANTERIOR A LA ENTRADA EN VIGOR DE LA LEY ACTUAL. (27 DE
FEBRERO DE 1992).


Para el cómputo de los plazos de cinco y diez años, que para efectos de la prescripción señala la
ley, según sea la posesión de buena o mala fe, se debe tomar en cuenta únicamente el tiempo que
se ha poseído la tierra ejidal, a partir de la entrada en vigor de la Ley Agraria, y no así el tiempo de
posesión anterior a la vigencia de ésta, toda vez que, la Ley Federal de la Reforma Agraria, no
contemplaba esta figura jurídica como medio para adquirir derechos agrarios, sino que en su
artículo 75 negaba la posibilidad de prescripción adquisitiva de derechos agrarios, por lo que éstos
no eran susceptibles de adquirirse por prescripción. Además, el tercer párrafo del artículo 48 de la
Ley Agraria establece las circunstancias para interrumpir la prescripción positiva, de lo que se
advierte que el término para que opere corre junto con la posibilidad de interrumpirlo, estimar que el
tiempo de posesión anterior a la entrada en vigor de la ley, es computable para el plazo de
prescripción positiva, dejaría al ejidatario que perdió la posesión sin la posibilidad de interrumpir ese
plazo en forma alguna; amén de que constituiría una aplicación retroactiva de la ley en su perjuicio,
lo cual es contrario a lo establecido en el artículo 14 constitucional.


Contradicción de tesis 16/94. Entre las sustentadas por el Tribunal Colegiado del Vigésimo Tercer
Circuito y el Segundo Tribunal Colegiado del Segundo Circuito. 18 de noviembre de 1994.
Unanimidad de cuatro votos. Ausente: Carlos de Silva Nava. Ponente: Fausta Moreno Flores.
Secretaria: Guadalupe Robles Denetro.

Tesis de Jurisprudencia 24/94. Aprobada por la Segunda Sala de este alto Tribunal, en sesión
pública de dieciocho de noviembre de mil novecientos noventa y cuatro, por unanimidad de cuatro
votos de los señores Ministros: Presidente: Atanasio González Martínez, José Manuel Villagordoa
Lozano, Fausta Moreno Flores y Noé Castañón León. Ausente: Carlos de Silva Nava.




                                              Pág. 1 de 2              Fecha de impresión 01/10/2025
  https://sjf2.scjn.gob.mx/detalle/tesis/206322
                                                Semanario Judicial de la Federación




                                           Pág. 2 de 2       Fecha de impresión 01/10/2025
https://sjf2.scjn.gob.mx/detalle/tesis/206322

\end{quote}

\subsection{Formalization}

\subsubsection{Clause: tesis206322\_principio\_fundamental}

\textit{Thesis-specific principle: Agrarian adverse possession periods count only from Agrarian Law effective date (Feb 27, 1992), not from previous possession}

\textbf{Intermediate Representation:}

\begin{quote}
posesion agraria is obligated and posesion tiempo 5 anos is obligated and posesion tiempo 10 anos is obligated
\end{quote}

\textbf{Witness Clause:}

\begin{lstlisting}[style=witnessstyle, language=Witness]
clause tesis206322_principio_fundamental = oblig(posesion_agraria) AND oblig(posesion_tiempo_5_anos) AND oblig(posesion_tiempo_10_anos) IMPLIES oblig(ejercer_prescripcion_positiva);
\end{lstlisting}

\subsubsection{Clause: tesis206322\_posesion\_anterior}

\textit{Supporting legal logic: Previous possession before law effective date doesn't count toward prescription periods}

\textbf{Intermediate Representation:}

\begin{quote}
temporal possession is obligated and not (posesion agraria is obligated)
\end{quote}

\textbf{Witness Clause:}

\begin{lstlisting}[style=witnessstyle, language=Witness]
clause tesis206322_posesion_anterior = oblig(posesion_temporal) AND not(oblig(posesion_agraria)) IMPLIES not(oblig(posesion_tiempo_5_anos));
\end{lstlisting}

\subsubsection{Clause: tesis206322\_derecho\_interrupcion}

\textit{Interruption rights: Ejidatarios who lost possession must be able to interrupt prescription}

\textbf{Intermediate Representation:}

\begin{quote}
interrupcion prescripcion is obligated
\end{quote}

\textbf{Witness Clause:}

\begin{lstlisting}[style=witnessstyle, language=Witness]
clause tesis206322_derecho_interrupcion = oblig(interrupcion_prescripcion) IMPLIES oblig(posesion_agraria);
\end{lstlisting}

\subsubsection{Clause: tesis206322\_proteccion\_constitucional}

\textit{Constitutional protection: Retroactive application would violate constitutional rights}

\textbf{Intermediate Representation:}

\begin{quote}
posesion agraria is obligated
\end{quote}

\textbf{Witness Clause:}

\begin{lstlisting}[style=witnessstyle, language=Witness]
clause tesis206322_proteccion_constitucional = oblig(posesion_agraria) IMPLIES oblig(proteccion_legal);
\end{lstlisting}

\newpage

\section{Thesis 210115: STATE/MUNICIPAL PROPERTY ADVERSE POSSESSION REQUIREMENTS}
\label{thesis:210115}

\subsection{Original Thesis Text}

\begin{quote}
                                                   Semanario Judicial de la Federación

                                                  Tesis

 Registro digital: 210115

 Instancia: Tribunales            Octava Época                          Materia(s): Civil
 Colegiados de Circuito

 Tesis: II. 2o. C.T. 6 C          Fuente: Semanario Judicial de la      Tipo: Aislada
                                  Federación.
                                  Tomo XIV, Noviembre de 1994,
                                  página 545


USUCAPION DE BIENES PROPIOS DEL ESTADO O MUNICIPIOS. DEBEN SER OBJETO DE
POSESION APTA PARA PRESCRIBIR POR EL DOBLE DEL TIEMPO FIJADO A INMUEBLES DE
PROPIEDAD PARTICULAR. (ESTADO DE MEXICO).


De acuerdo con lo previsto por el artículo 918 del Código Civil, los bienes pertenecientes al Estado o
municipios sólo pueden ser adquiridos mediante prescripción positiva cuando se han tenido en
posesión apta para usucapir por el doble del tiempo que la misma ley exige para las demás
personas, por eso, cuando llega a demostrarse que la tenencia material de un inmueble inscrito en
favor del Gobierno del Estado es de buena fe y se ha prolongado por más de diez años, con ello
prospera la usucapión, porque de conformidad con lo dispuesto por el artículo 912, fracción I del
Código Civil, los bienes inmuebles se prescriben en cinco años, cuando la posesión reúne las
características de ser continua, pública, de buena fe y en concepto de propietario; en tal virtud,
tratándose de bienes pertenecientes al Estado o Municipios, el plazo para la prescripción positiva
fijado por el artículo 918 de la legislación en consulta se consuma en diez años. Y sólo aquellos
casos en que la posesión de esos bienes sea de mala fe requerirá el lapso de veinte años.

SEGUNDO TRIBUNAL COLEGIADO EN MATERIAS CIVIL Y DE TRABAJO DEL SEGUNDO
CIRCUITO .


Amparo directo 38/94. Gobierno del Estado de México. 31 de agosto de 1994. Unanimidad de votos.
Ponente: Raúl Solís Solís. Secretario: Joel Alfonso Sierra Palacios.




                                             Pág. 1 de 1             Fecha de impresión 01/10/2025
  https://sjf2.scjn.gob.mx/detalle/tesis/210115

\end{quote}

\subsection{Formalization}

\subsubsection{Clause: tesis210115\_principio\_fundamental}

\textit{Thesis-specific principle: State/municipal property requires double the time period for adverse possession}

\textbf{Intermediate Representation:}

\begin{quote}
possession as owner is obligated and good faith exercise is obligated and posesion tiempo 10 anos estatal is obligated
\end{quote}

\textbf{Witness Clause:}

\begin{lstlisting}[style=witnessstyle, language=Witness]
clause tesis210115_principio_fundamental = oblig(posesion_como_propietario) AND oblig(ejercicio_buena_fe) AND oblig(posesion_tiempo_10_anos_estatal) IMPLIES oblig(ejercer_prescripcion_positiva);
\end{lstlisting}

\subsubsection{Clause: tesis210115\_buena\_fe\_estatal}

\textit{Supporting legal logic: Good faith state property possession requires 10 years (double of private property 5 years)}

\textbf{Intermediate Representation:}

\begin{quote}
good faith exercise is obligated and posesion tiempo 10 anos estatal is obligated
\end{quote}

\textbf{Witness Clause:}

\begin{lstlisting}[style=witnessstyle, language=Witness]
clause tesis210115_buena_fe_estatal = oblig(ejercicio_buena_fe) AND oblig(posesion_tiempo_10_anos_estatal) IMPLIES oblig(ejercer_prescripcion_positiva);
\end{lstlisting}

\subsubsection{Clause: tesis210115\_mala\_fe\_estatal}

\textit{Bad faith state property possession requires 20 years (double of private property 10 years)}

\textbf{Intermediate Representation:}

\begin{quote}
bad faith exercise is obligated and posesion tiempo 20 anos estatal is obligated
\end{quote}

\textbf{Witness Clause:}

\begin{lstlisting}[style=witnessstyle, language=Witness]
clause tesis210115_mala_fe_estatal = oblig(ejercicio_mala_fe) AND oblig(posesion_tiempo_20_anos_estatal) IMPLIES oblig(ejercer_prescripcion_positiva);
\end{lstlisting}

\subsubsection{Clause: tesis210115\_requisitos\_posesion}

\textit{Possession requirements: Must be continuous, public, good faith, and in concept of owner}

\textbf{Intermediate Representation:}

\begin{quote}
possession as owner is obligated and good faith exercise is obligated
\end{quote}

\textbf{Witness Clause:}

\begin{lstlisting}[style=witnessstyle, language=Witness]
clause tesis210115_requisitos_posesion = oblig(posesion_como_propietario) AND oblig(ejercicio_buena_fe) IMPLIES oblig(posesion_tiempo_10_anos_estatal);
\end{lstlisting}

\newpage

\section{Thesis 210588: AGRARIAN ADVERSE POSSESSION 5-YEAR WAITING PERIOD}
\label{thesis:210588}

\subsection{Original Thesis Text}

\begin{quote}
                                                   Semanario Judicial de la Federación

                                                  Tesis

 Registro digital: 210588

 Instancia: Tribunales             Octava Época                           Materia(s): Administrativa
 Colegiados de Circuito

 Tesis: XVII. 2o. 11 A             Fuente: Semanario Judicial de la       Tipo: Aislada
                                   Federación.
                                   Tomo XIV, Septiembre de 1994,
                                   página 393


PRESCRIPCION ADQUISITIVA DE BUENA FE SOBRE TIERRAS EJIDALES, FIGURA PREVISTA
POR LA LEY AGRARIA VIGENTE. DEBE TRANSCURRIR EL TERMINO DE CINCO AÑOS, A
PARTIR DE SU VIGENCIA, PARA QUE SE ACTUALICE.


Si bien es cierto que la Ley Agraria vigente, en su artículo 48, prevé la prescripción adquisitiva como
un medio de adquirir tierras ejidales, cuando se poseen las mismas por el tiempo y con los
requisitos que establece, no menos cierto es que la derogada Ley Federal de Reforma Agraria,
establecía en su artículo 52 que: "Los derechos que sobre bienes agrarios adquieren los núcleos de
población serán inalienables, imprescriptibles, inembargables e intransmisibles y, por tanto, no
podrán en ningún caso ni en forma alguna enajenarse, cederse, transmitirse, arrendarse,
hipotecarse o gravarse, en todo o en parte. Serán inexistentes las operaciones, actos o contratos
que se hayan ejecutado o que se pretenden llevar a cabo en contravención de este precepto. Las
tierras cultivables que de acuerdo con la ley pueden ser objeto de adjudicación individual entre los
miembros del ejido, en ningún momento dejarán de ser propiedad del núcleo de población ejidal. El
aprovechamiento individual, cuando exista, terminará al resolverse, de acuerdo con la ley, que la
explotación debe ser colectiva en beneficio de todos los integrantes del ejido y renacerá cuando
ésta termine. Las unidades de dotación y solares que hayan pertenecido a ejidatarios y que resulten
vacantes por ausencia de heredero o sucesor legal, quedarán a disposición del núcleo de población
correspondiente. Este artículo es aplicable a los bienes que pertenecen a los núcleos de población
que de hecho o por derecho guarden el estado comunal"; ahora bien, es conveniente evidenciar que
la prescripción adquisitiva de tierras ejidales, prevista en el precepto legal primeramente citado,
adquirió vigencia a partir del día siguiente al en que se publicó en el Diario Oficial de la Federación
la ley que lo contiene, conforme lo dispone su artículo primero transitorio, que textualmente dice: "La
presente ley entrará en vigor al día siguiente de su publicación en el Diario Oficial de la Federación",
lo que ocurrió el miércoles veintiséis de febrero de mil novecientos noventa y dos, por cuyo motivo
es inconcuso que el jueves veintisiete de febrero del citado año inició su vigencia;
consecuentemente, es claro concluir que el artículo 48 de la Ley Agraria en vigor sólo es aplicable a
partir del veintisiete de febrero de mil novecientos noventa y dos, y que, por ende, no puede
comprender situaciones que se dieron estando en vigor la derogada Ley Federal de Reforma
Agraria, puesto que, conforme a ésta, las tierras ejidales eran imprescriptibles. En tales condiciones,
si el cuatro de febrero de mil novecientos noventa y cuatro, ante el Tribunal Agrario, es presentada
en la vía de jurisdicción voluntaria una demanda, intentando la prescripción adquisitiva de tierras
ejidales, fundándose para dicho efecto en el artículo 48 de la Ley Agraria en vigor porque se poseen
las mismas de buena fe y se cumple con los demás requisitos previstos en el mismo, el
desechamiento de plano de esa demanda es legal, ya que aún no transcurría el término de cinco
años que para prescribir bienes ejidales establece el precitado artículo 48 de la Ley Agraria en vigor,


                                              Pág. 1 de 2             Fecha de impresión 01/10/2025
  https://sjf2.scjn.gob.mx/detalle/tesis/210588
                                                  Semanario Judicial de la Federación

además de que admitirla, implicaría transgredir la garantía de irretroactividad de la ley, establecida
en el artículo 14 constitucional.

SEGUNDO TRIBUNAL COLEGIADO DEL DECIMO SEPTIMO CIRCUITO.


Amparo directo 161/94. Jesús Manuel Sierra Herrera. 8 de junio de 1994. Unanimidad de votos.
Ponente: Angel Gregorio Vázquez González. Secretaria: Blanca Estela Quezada Rojas.




                                             Pág. 2 de 2             Fecha de impresión 01/10/2025
  https://sjf2.scjn.gob.mx/detalle/tesis/210588

\end{quote}

\subsection{Formalization}

\subsubsection{Clause: tesis210588\_principio\_fundamental}

\textit{Thesis-specific principle: Agrarian adverse possession requires 5 years from Agrarian Law effective date (Feb 27, 1992) before it can be invoked}

\textbf{Intermediate Representation:}

\begin{quote}
posesion agraria is obligated and posesion tiempo 5 anos is obligated and good faith exercise is obligated
\end{quote}

\textbf{Witness Clause:}

\begin{lstlisting}[style=witnessstyle, language=Witness]
clause tesis210588_principio_fundamental = oblig(posesion_agraria) AND oblig(posesion_tiempo_5_anos) AND oblig(ejercicio_buena_fe) IMPLIES oblig(ejercer_prescripcion_positiva);
\end{lstlisting}

\subsubsection{Clause: tesis210588\_posesion\_anterior}

\textit{Supporting legal logic: Previous possession under old law (imprescriptible) doesn't count toward 5-year requirement}

\textbf{Intermediate Representation:}

\begin{quote}
temporal possession is obligated and not (posesion agraria is obligated)
\end{quote}

\textbf{Witness Clause:}

\begin{lstlisting}[style=witnessstyle, language=Witness]
clause tesis210588_posesion_anterior = oblig(posesion_temporal) AND not(oblig(posesion_agraria)) IMPLIES not(oblig(posesion_tiempo_5_anos));
\end{lstlisting}

\subsubsection{Clause: tesis210588\_marco\_legal\_nuevo}

\textit{Legal framework change: New law recognizes adverse possession, old law made ejidal lands imprescriptible}

\textbf{Intermediate Representation:}

\begin{quote}
posesion agraria is obligated
\end{quote}

\textbf{Witness Clause:}

\begin{lstlisting}[style=witnessstyle, language=Witness]
clause tesis210588_marco_legal_nuevo = oblig(posesion_agraria) IMPLIES oblig(ejercer_prescripcion_positiva);
\end{lstlisting}

\subsubsection{Clause: tesis210588\_proteccion\_constitucional}

\textit{Constitutional protection: Retroactive application would violate constitutional guarantee}

\textbf{Intermediate Representation:}

\begin{quote}
posesion agraria is obligated
\end{quote}

\textbf{Witness Clause:}

\begin{lstlisting}[style=witnessstyle, language=Witness]
clause tesis210588_proteccion_constitucional = oblig(posesion_agraria) IMPLIES oblig(proteccion_legal);
\end{lstlisting}

\newpage

\section{Thesis 210589: MUTUAL EXCLUSIVITY OF GOOD FAITH AND BAD FAITH POSSESSION}
\label{thesis:210589}

\subsection{Original Thesis Text}

\begin{quote}
                                                   Semanario Judicial de la Federación

                                                  Tesis

 Registro digital: 210589

 Instancia: Tribunales            Octava Época                          Materia(s): Civil
 Colegiados de Circuito

 Tesis: VII. 2o. C. 28 C          Fuente: Semanario Judicial de la      Tipo: Aislada
                                  Federación.
                                  Tomo XIV, Septiembre de 1994,
                                  página 394


PRESCRIPCION ADQUISITIVA DE INMUEBLES, ACCION DE. NO PUEDE EJERCERSE BASADA
SIMULTANEAMENTE EN LOS SUPUESTOS DE BUENA Y DE MALA FE (LEGISLACION DEL
ESTADO DE VERACRUZ).


La posesión no puede ser al mismo tiempo de buena y de mala fe, lo que encuentra claro sustento
en el artículo 842 del Código Civil local, dado que no puede ocurrir simultáneamente que se cuente
con un título (causa generadora de la posesión) que ampare el derecho de posesión y a la vez no
se cuente con él, o bien que, contando con él, se conozcan y desconozcan a un tiempo los vicios de
que adolece, por lo cual, si en un caso se planteó en los hechos de la demanda que se contaba con
un justo título, pero no que se careciera de él, o bien que, teniéndolo, éste adoleciera de vicios que
impidieran poseer con derecho, al tenor de ese dispositivo legal, resulta irrelevante que en la
demanda nominalmente se dijera, como de modo secundario, que también se hacía valer la
prescripción basada en la posesión de mala fe, contradiciendo los hechos de la demanda, posesión
que no puede coexistir con la de buena fe. Por ello, no es cierto que debía estudiarse este otro
fundamento, diverso del que en los hechos se planteó, pues de acuerdo con el artículo 2o. del
Código de Procedimientos Civiles veracruzano deben tomarse en cuenta los hechos narrados en la
demanda como fundamento de la misma para establecer, además de la clase de prestación exigida,
"el título o causa de la acción", con independencia de cómo el deducente de la misma la califique,
pues son tales hechos los que establecen el planteamiento del litigio y el fundamento de la
demanda.

SEGUNDO TRIBUNAL COLEGIADO EN MATERIA CIVIL DEL SEPTIMO CIRCUITO.


Amparo directo 94/94. Bulmaro Méndez Valiente. 25 de marzo de 1994. Unanimidad de votos.
Ponente: Héctor Soto Gallardo. Secretario: Felipe Alfredo Fuentes Barrera.




                                             Pág. 1 de 2             Fecha de impresión 01/10/2025
  https://sjf2.scjn.gob.mx/detalle/tesis/210589
                                                Semanario Judicial de la Federación




                                           Pág. 2 de 2       Fecha de impresión 01/10/2025
https://sjf2.scjn.gob.mx/detalle/tesis/210589

\end{quote}

\subsection{Formalization}

\subsubsection{Clause: tesis210589\_principio\_fundamental}

\textit{Thesis-specific principle: Possession cannot be simultaneously good faith and bad faith - they are mutually exclusive}

\textbf{Intermediate Representation:}

\begin{quote}
good faith exercise is obligated and not (bad faith exercise is obligated)
\end{quote}

\textbf{Witness Clause:}

\begin{lstlisting}[style=witnessstyle, language=Witness]
clause tesis210589_principio_fundamental = oblig(ejercicio_buena_fe) AND not(oblig(ejercicio_mala_fe));
\end{lstlisting}

\subsubsection{Clause: tesis210589\_titulo\_contradiccion}

\textit{Supporting legal logic: Cannot have title (generating cause) and not have it simultaneously}

\textbf{Intermediate Representation:}

\begin{quote}
proven generating cause is obligated and not (not (proven generating cause is obligated))
\end{quote}

\textbf{Witness Clause:}

\begin{lstlisting}[style=witnessstyle, language=Witness]
clause tesis210589_titulo_contradiccion = oblig(causa_generadora_probada) AND not(not(oblig(causa_generadora_probada)));
\end{lstlisting}

\subsubsection{Clause: tesis210589\_conocimiento\_contradiccion}

\textit{Knowledge contradiction: Cannot simultaneously know and not know title defects}

\textbf{Intermediate Representation:}

\begin{quote}
conocimiento vicios is obligated and not (not (conocimiento vicios is obligated))
\end{quote}

\textbf{Witness Clause:}

\begin{lstlisting}[style=witnessstyle, language=Witness]
clause tesis210589_conocimiento_contradiccion = oblig(conocimiento_vicios) AND not(not(oblig(conocimiento_vicios)));
\end{lstlisting}

\subsubsection{Clause: tesis210589\_consistencia\_procesal}

\textit{Procedural requirement: Action must be based on consistent facts, not contradictory claims}

\textbf{Intermediate Representation:}

\begin{quote}
If possession as owner is obligated, then (good faith exercise is obligated and not (bad faith exercise is obligated))
\end{quote}

\textbf{Witness Clause:}

\begin{lstlisting}[style=witnessstyle, language=Witness]
clause tesis210589_consistencia_procesal = oblig(posesion_como_propietario) IMPLIES (oblig(ejercicio_buena_fe) AND not(oblig(ejercicio_mala_fe)));
\end{lstlisting}

\newpage

\section{Thesis 210888: POSSESSION AS OWNER REQUIREMENT FOR ADVERSE POSSESSION}
\label{thesis:210888}

\subsection{Original Thesis Text}

\begin{quote}
                                                   Semanario Judicial de la Federación

                                                  Tesis

 Registro digital: 210888

 Instancia: Tribunales            Octava Época                         Materia(s): Civil
 Colegiados de Circuito

 Tesis: II. 3o. 257 C             Fuente: Semanario Judicial de la     Tipo: Aislada
                                  Federación.
                                  Tomo XIV, Agosto de 1994, página
                                  645


PRESCRIPCION ADQUISITIVA. POSESION EN CONCEPTO DE PROPIETARIO. (LEGISLACION
DEL ESTADO DE MEXICO).


El Código Civil para el Estado de México, en su artículo 911 dispone que los requisitos para
usucapir consisten en que la posesión sea en concepto de propietario, en forma pacífica, continua y
pública; pero no prevé como elemento de la acción prescriptiva el que se acredite el llamado justo
título. Así, la exigencia de poseer en concepto de propietario para poder adquirir por prescripción,
comprende no sólo los casos de buena fe, sino también al caso de la posesión de mala fe, por
ende, para usucapir únicamente se requiere para evidenciar que se posee en concepto de
propietario, que el poseedor es el dominador de la cosa, el que manda en ella y la disfruta para sí
como dueño en sentido económico, aun cuando carezca de un título legítimo, frente a todo el
mundo y siempre y cuando haya comenzado a poseer en virtud de una causa diversa de la que
origina la posesión derivada.

TERCER TRIBUNAL COLEGIADO DEL SEGUNDO CIRCUITO.


Amparo directo 349/94. Virginia Alanís Vda. de Valdéz. 24 de mayo de 1994. Unanimidad de votos.
Ponente: Fernando Narváez Barker. Secretario: Alejandro García Gómez.

Reitera criterio de la tesis de jurisprudencia número 1379, páginas 2227 y 2228, Segunda Parte,
Apéndice al Semanario Judicial de la Federación 1917-1988.




                                             Pág. 1 de 1            Fecha de impresión 01/10/2025
  https://sjf2.scjn.gob.mx/detalle/tesis/210888

\end{quote}

\subsection{Formalization}

\subsubsection{Clause: tesis210888\_principio\_fundamental}

\textit{Thesis-specific principle: Possession must be "in concept of owner" (peaceful, continuous, public) for both good and bad faith}

\textbf{Intermediate Representation:}

\begin{quote}
possession as owner is obligated and (oblig(ejercicio\_buena\_fe) OR oblig(ejercicio\_mala\_fe)) IMPLIES oblig(ejercer\_prescripcion\_positiva)
\end{quote}

\textbf{Witness Clause:}

\begin{lstlisting}[style=witnessstyle, language=Witness]
clause tesis210888_principio_fundamental = oblig(posesion_como_propietario) AND (oblig(ejercicio_buena_fe) OR oblig(ejercicio_mala_fe)) IMPLIES oblig(ejercer_prescripcion_positiva);
\end{lstlisting}

\subsubsection{Clause: tesis210888\_posesion\_ambas\_fe}

\textit{Supporting logic: Possession as owner applies to both good faith and bad faith cases}

\textbf{Intermediate Representation:}

\begin{quote}
If possession as owner is obligated, then (good faith exercise is obligated or bad faith exercise is obligated)
\end{quote}

\textbf{Witness Clause:}

\begin{lstlisting}[style=witnessstyle, language=Witness]
clause tesis210888_posesion_ambas_fe = oblig(posesion_como_propietario) IMPLIES (oblig(ejercicio_buena_fe) OR oblig(ejercicio_mala_fe));
\end{lstlisting}

\subsubsection{Clause: tesis210888\_dominio\_economico}

\textit{Economic dominion: Must be dominator of the thing, commanding and enjoying it as owner}

\textbf{Intermediate Representation:}

\begin{quote}
possession as owner is obligated
\end{quote}

\textbf{Witness Clause:}

\begin{lstlisting}[style=witnessstyle, language=Witness]
clause tesis210888_dominio_economico = oblig(posesion_como_propietario) IMPLIES oblig(causa_generadora_probada);
\end{lstlisting}

\subsubsection{Clause: tesis210888\_posesion\_originaria}

\textit{Original possession: Must begin with cause different from derived possession}

\textbf{Intermediate Representation:}

\begin{quote}
possession as owner is obligated and not (temporal possession is obligated)
\end{quote}

\textbf{Witness Clause:}

\begin{lstlisting}[style=witnessstyle, language=Witness]
clause tesis210888_posesion_originaria = oblig(posesion_como_propietario) AND not(oblig(posesion_temporal)) IMPLIES oblig(causa_generadora_probada);
\end{lstlisting}

\newpage

\section{Thesis 210890: ADVERSE POSSESSION BY ACTION OR EXCEPTION}
\label{thesis:210890}

\subsection{Original Thesis Text}

\begin{quote}
                                                   Semanario Judicial de la Federación

                                                  Tesis

 Registro digital: 210890

 Instancia: Tribunales            Octava Época                          Materia(s): Civil
 Colegiados de Circuito

 Tesis: XXI. 2o. 42 C             Fuente: Semanario Judicial de la      Tipo: Aislada
                                  Federación.
                                  Tomo XIV, Agosto de 1994, página
                                  646


PRESCRIPCION POSITIVA, PROCEDE HACERLA VALER POR VIA DE ACCION O EXCEPCION
(LEGISLACION DEL ESTADO DE GUERRERO).


Atento a lo dispuesto en los artículos 1135 y 1136 del Código Civil del Estado de Guerrero vigente
hasta el dos de septiembre de mil novecientos noventa y tres, la prescripción positiva o adquisitiva,
es un medio de adquirir bienes mediante la posesión por el transcurso de cierto tiempo y bajo las
condiciones establecidas por la ley. Dicha figura jurídica, puede hacerse valer por vía de acción o de
excepción, dependiendo de la situación procesal en que se encuentre quien la hace valer. Sin
embargo, es inexacto que cuando se hace valer en vía de excepción, y no obstante que el actor no
acredite los elementos de su acción, la parte demandada no requiera la demostración
pormenorizada de las cualidades y circunstancias exigidas por la ley para la usucapión, tales como
poseer en calidad de dueño, por el tiempo que la ley exige, de buena o de mala fe, en forma
pacífica, continua y pública, y que sólo deba demostrarse el hecho que justifique el cómputo del
término prescriptivo. Esto es así, porque en ambos casos, la finalidad perseguida es la misma, a
saber, la ya precisada de adquirir bienes en virtud de la posesión por determinado tiempo y bajo las
condiciones establecidas por la ley, y por consiguiente, los medios para conseguirla, también deben
ser iguales, es decir, en ambas hipótesis, el interesado debe probar todos los elementos necesarios
para que opere la usucapión, concretamente, aquéllos exigidos en los artículos 1151, 1152 y 1156
del Código Civil antes citado (artículo 752, 755 y 768 del Código Civil del Estado de Guerrero en
vigor). En ese orden de ideas, tratándose de inmuebles, los elementos a demostrar son los
siguientes: Que el bien que se pretende prescribir, haya sido poseído en concepto de propietario, en
forma pacífica, continua y pública, por un término de cinco años cuando el poseedor es de buena fe,
o por un término de diez años cuando la posesión es de mala fe; además que, la acción respectiva
debe enderezarse en contra de quien aparezca como propietario del bien en el Registro Público.

SEGUNDO TRIBUNAL COLEGIADO DEL VIGESIMO PRIMER CIRCUITO.


Amparo directo 201/94. María Luisa Vázquez López. 16 de junio de 1994. Unanimidad de votos.
Ponente: Martiniano Bautista Espinosa. Secretario: Javier Cardoso Chávez.




                                             Pág. 1 de 2             Fecha de impresión 01/10/2025
  https://sjf2.scjn.gob.mx/detalle/tesis/210890
                                                Semanario Judicial de la Federación




                                           Pág. 2 de 2       Fecha de impresión 01/10/2025
https://sjf2.scjn.gob.mx/detalle/tesis/210890

\end{quote}

\subsection{Formalization}

\subsubsection{Clause: tesis210890\_principio\_fundamental}

\textit{Thesis-specific principle: Adverse possession can be invoked by action or exception, but requires same proof elements}

\textbf{Intermediate Representation:}

\begin{quote}
possession as owner is obligated and (good faith exercise is obligated or bad faith exercise is obligated) and (oblig(posesion\_tiempo\_5\_anos) OR oblig(posesion\_tiempo\_10\_anos)) IMPLIES oblig(ejercer\_prescripcion\_positiva)
\end{quote}

\textbf{Witness Clause:}

\begin{lstlisting}[style=witnessstyle, language=Witness]
clause tesis210890_principio_fundamental = oblig(posesion_como_propietario) AND (oblig(ejercicio_buena_fe) OR oblig(ejercicio_mala_fe)) AND (oblig(posesion_tiempo_5_anos) OR oblig(posesion_tiempo_10_anos)) IMPLIES oblig(ejercer_prescripcion_positiva);
\end{lstlisting}

\subsubsection{Clause: tesis210890\_igualdad\_prueba}

\textit{Supporting logic: Both action and exception require same proof elements regardless of procedural route}

\textbf{Intermediate Representation:}

\begin{quote}
ejercer prescripcion positiva is obligated
\end{quote}

\textbf{Witness Clause:}

\begin{lstlisting}[style=witnessstyle, language=Witness]
clause tesis210890_igualdad_prueba = oblig(ejercer_prescripcion_positiva) IMPLIES oblig(posesion_como_propietario);
\end{lstlisting}

\subsubsection{Clause: tesis210890\_plazos\_temporales}

\textit{Time requirements: 5 years for good faith, 10 years for bad faith}

\textbf{Intermediate Representation:}

\begin{quote}
good faith exercise is obligated
\end{quote}

\textbf{Witness Clause:}

\begin{lstlisting}[style=witnessstyle, language=Witness]
clause tesis210890_plazos_temporales = oblig(ejercicio_buena_fe) IMPLIES oblig(posesion_tiempo_5_anos);
\end{lstlisting}

\subsubsection{Clause: tesis210890\_plazos\_mala\_fe}

\textbf{Intermediate Representation:}

\begin{quote}
bad faith exercise is obligated
\end{quote}

\textbf{Witness Clause:}

\begin{lstlisting}[style=witnessstyle, language=Witness]
clause tesis210890_plazos_mala_fe = oblig(ejercicio_mala_fe) IMPLIES oblig(posesion_tiempo_10_anos);
\end{lstlisting}

\subsubsection{Clause: tesis210890\_objetivo\_registral}

\textit{Target requirement: Action must be directed against registered owner}

\textbf{Intermediate Representation:}

\begin{quote}
ejercer prescripcion positiva is obligated
\end{quote}

\textbf{Witness Clause:}

\begin{lstlisting}[style=witnessstyle, language=Witness]
clause tesis210890_objetivo_registral = oblig(ejercer_prescripcion_positiva) IMPLIES oblig(titularidad_registral);
\end{lstlisting}

\newpage

\section{Thesis 211744: REVEALING POSSESSION CAUSE REQUIREMENT}
\label{thesis:211744}

\subsection{Original Thesis Text}

\begin{quote}
                                                   Semanario Judicial de la Federación

                                                  Tesis

 Registro digital: 211744

 Instancia: Tribunales            Octava Época                         Materia(s): Civil
 Colegiados de Circuito

                                  Fuente: Semanario Judicial de la     Tipo: Aislada
                                  Federación.
                                  Tomo XIV, Julio de 1994, página
                                  719


PRESCRIPCION ADQUISITIVA. NECESIDAD DE REVELAR LA CAUSA DE LA POSESION.


El actor en un juicio de prescripción positiva, debe revelar la causa de su posesión, aun en el caso
de poseedor de mala fe, porque es necesario que el juzgador conozca el hecho o acto generador de
la misma, para poder determinar la calidad de la posesión, si es en concepto de propietario,
originaria o derivada, de buena o mala fe y para precisar el momento en que debe empezar a contar
el plazo de la prescripción.

SEGUNDO TRIBUNAL COLEGIADO DEL SEXTO CIRCUITO.


Amparo directo 22/92. Concepción Sánchez Nava. 29 de enero de 1992. Unanimidad de votos.
Ponente: Gustavo Calvillo Rangel. Secretario: Jorge Alberto González Alvarez.

Amparo directo 365/89. Roquelia Díaz Pineda. 24 de noviembre de 1989. Unanimidad de votos.
Ponente: José Galván Rojas. Secretario: Vicente Martínez Sánchez.

Véase: Apéndice de Jurisprudencia 1917-1985. Cuarta Parte. Tesis 218, pág. 631.




                                             Pág. 1 de 1             Fecha de impresión 01/10/2025
  https://sjf2.scjn.gob.mx/detalle/tesis/211744

\end{quote}

\subsection{Formalization}

\subsubsection{Clause: tesis211744\_principio\_fundamental}

\textit{Thesis-specific principle: Actor must reveal the cause of possession, even for bad faith possessors, for judicial determination}

\textbf{Intermediate Representation:}

\begin{quote}
ejercer prescripcion positiva is obligated
\end{quote}

\textbf{Witness Clause:}

\begin{lstlisting}[style=witnessstyle, language=Witness]
clause tesis211744_principio_fundamental = oblig(ejercer_prescripcion_positiva) IMPLIES oblig(causa_generadora_revelada);
\end{lstlisting}

\subsubsection{Clause: tesis211744\_determinacion\_judicial}

\textit{Supporting logic: Judge needs to know generating cause to determine possession quality and timing}

\textbf{Intermediate Representation:}

\begin{quote}
If causa generadora revelada is obligated, then (possession as owner is obligated and (good faith exercise is obligated and not (bad faith exercise is obligated)))
\end{quote}

\textbf{Witness Clause:}

\begin{lstlisting}[style=witnessstyle, language=Witness]
clause tesis211744_determinacion_judicial = oblig(causa_generadora_revelada) IMPLIES (oblig(posesion_como_propietario) AND (oblig(ejercicio_buena_fe) AND not(oblig(ejercicio_mala_fe))));
\end{lstlisting}

\subsubsection{Clause: tesis211744\_requisito\_universal}

\textit{Universal requirement: Applies to both good faith and bad faith possessors}

\textbf{Intermediate Representation:}

\begin{quote}
(oblig(ejercicio\_buena\_fe) OR oblig(ejercicio\_mala\_fe)) IMPLIES oblig(causa\_generadora\_revelada)
\end{quote}

\textbf{Witness Clause:}

\begin{lstlisting}[style=witnessstyle, language=Witness]
clause tesis211744_requisito_universal = (oblig(ejercicio_buena_fe) OR oblig(ejercicio_mala_fe)) IMPLIES oblig(causa_generadora_revelada);
\end{lstlisting}

\subsubsection{Clause: tesis211744\_determinacion\_temporal}

\textit{Timing determination: Must establish when prescription period begins to count}

\textbf{Intermediate Representation:}

\begin{quote}
If causa generadora revelada is obligated, then (posesion tiempo 5 anos is obligated or posesion tiempo 10 anos is obligated)
\end{quote}

\textbf{Witness Clause:}

\begin{lstlisting}[style=witnessstyle, language=Witness]
clause tesis211744_determinacion_temporal = oblig(causa_generadora_revelada) IMPLIES (oblig(posesion_tiempo_5_anos) OR oblig(posesion_tiempo_10_anos));
\end{lstlisting}

\newpage

\section{Thesis 212497: AGRARIAN ADVERSE POSSESSION PROCEDURAL REQUIREMENTS}
\label{thesis:212497}

\subsection{Original Thesis Text}

\begin{quote}
                                                   Semanario Judicial de la Federación

                                                  Tesis

 Registro digital: 212497

 Instancia: Tribunales            Octava Época                           Materia(s): Administrativa
 Colegiados de Circuito

 Tesis: IV.2o.63 A                Fuente: Semanario Judicial de la       Tipo: Aislada
                                  Federación.
                                  Tomo XIII, Mayo de 1994, página
                                  394


AGRARIO. RECONOCIMIENTO DE DERECHOS EJIDALES VIA PRESCRIPCION ADQUISITIVA,
DESECHAMIENTO ILEGAL DE LA SOLICITUD O DEMANDA DE (NUEVA LEY AGRARIA).


Conforme a lo dispuesto por el artículo 48 de la nueva Ley Agraria, el reconocimiento de derechos
ejidales vía prescripción adquisitiva, procede en el evento de que el interesado acredite que la
posesión de la tierra parcelada es en concepto de titular de ese derecho, en forma pacífica, continua
y pública, durante un período de cinco años si la posesión es de buena fe o de diez si es de mala fe.
El propio numeral preceptúa que el Tribunal Unitario Agrario, previa audiencia de los interesados,
del comisariado ejidal y los colindantes, en la vía de jurisdicción voluntaria o mediante el desahogo
del juicio correspondiente, emitirá la resolución sobre la pretendida adquisición de los derechos
anotados. Lo anterior lleva a la conclusión de que no es lógico ni jurídico desechar la solicitud de la
jurisdicción voluntaria o la demanda del juicio agrario, bajo el argumento de que no ha transcurrido
el término previsto por el artículo 48 de la Ley Agraria para que opere la prescripción adquisitiva,
pues por una parte no se da oportunidad al promovente de ser oído en el procedimiento relativo
conforme a las formalidades esenciales que al efecto exige dicha legislación y por otro lado, el
examen de los términos y condiciones necesarios para establecer si procede o no el reconocimiento
de derechos ejidales, es una cuestión propia de la resolución definitiva y no constituye, en
consecuencia, la resolución definitiva que justifique el desechamiento apriorístico de la solicitud de
jurisdicción voluntaria o de la demanda del juicio agrario, según sea el caso, pues la legislación
agraria vigente no contiene disposición legal alguna que establezca el supuesto anotado como
causa de improcedencia.

SEGUNDO TRIBUNAL COLEGIADO DEL CUARTO CIRCUITO.


Amparo directo 100/94. Bertha Silva Moreno. 2 de marzo de 1994. Unanimidad de votos. Ponente:
Arturo Barocio Villalobos. Secretario: Carlos Rafael Domínguez Avilán.

Nota: Esta tesis contendió en la contradicción 14/96 resuelta por la Segunda Sala, de la que
derivaron las tesis 2a./J. 6/98 y 2a./J. 7/98, que aparecen publicadas en el Semanario Judicial de la
Federación y su Gaceta, Novena Época, Tomo VII, febrero de 1998, páginas 113 y 130, con los
rubros: "ACCIÓN DE PRESCRIPCIÓN AGRARIA. EL AUTO INICIAL NO DEBE CALIFICAR SU
PROCEDENCIA." y "ACCIÓN DE PRESCRIPCIÓN AGRARIA, PROCEDENCIA DE LA. EL
ANÁLISIS DE SUS ELEMENTOS MATERIA DE LA RESOLUCIÓN DE FONDO, POR LO QUE SE
IMPONE EL TRÁMITE DEL PROCEDIMIENTO INSTAURADO.", respectivamente.



                                             Pág. 1 de 2              Fecha de impresión 01/10/2025
  https://sjf2.scjn.gob.mx/detalle/tesis/212497
                                                Semanario Judicial de la Federación




                                           Pág. 2 de 2       Fecha de impresión 01/10/2025
https://sjf2.scjn.gob.mx/detalle/tesis/212497

\end{quote}

\subsection{Formalization}

\subsubsection{Clause: tesis212497\_principio\_fundamental}

\textit{Thesis-specific principle: Agrarian adverse possession requires possession as titular of ejidal rights, peacefully, continuously, and publicly}

\textbf{Intermediate Representation:}

\begin{quote}
posesion agraria is obligated and (good faith exercise is obligated or bad faith exercise is obligated) and (oblig(posesion\_tiempo\_5\_anos) OR oblig(posesion\_tiempo\_10\_anos)) IMPLIES oblig(ejercer\_prescripcion\_positiva)
\end{quote}

\textbf{Witness Clause:}

\begin{lstlisting}[style=witnessstyle, language=Witness]
clause tesis212497_principio_fundamental = oblig(posesion_agraria) AND (oblig(ejercicio_buena_fe) OR oblig(ejercicio_mala_fe)) AND (oblig(posesion_tiempo_5_anos) OR oblig(posesion_tiempo_10_anos)) IMPLIES oblig(ejercer_prescripcion_positiva);
\end{lstlisting}

\subsubsection{Clause: tesis212497\_analisis\_temporal}

\textit{Supporting logic: Time period analysis is matter for final resolution, not for dismissal}

\textbf{Intermediate Representation:}

\begin{quote}
ejercer prescripcion positiva is obligated
\end{quote}

\textbf{Witness Clause:}

\begin{lstlisting}[style=witnessstyle, language=Witness]
clause tesis212497_analisis_temporal = oblig(ejercer_prescripcion_positiva) IMPLIES oblig(posesion_agraria);
\end{lstlisting}

\subsubsection{Clause: tesis212497\_debido\_proceso}

\textit{Due process requirement: Must give opportunity to be heard according to essential formalities}

\textbf{Intermediate Representation:}

\begin{quote}
ejercer prescripcion positiva is obligated
\end{quote}

\textbf{Witness Clause:}

\begin{lstlisting}[style=witnessstyle, language=Witness]
clause tesis212497_debido_proceso = oblig(ejercer_prescripcion_positiva) IMPLIES oblig(proteccion_legal);
\end{lstlisting}

\subsubsection{Clause: tesis212497\_flexibilidad\_procesal}

\textit{Procedural flexibility: Can be voluntary jurisdiction or adversarial proceeding}

\textbf{Intermediate Representation:}

\begin{quote}
If posesion agraria is obligated, then (ejercer prescripcion positiva is obligated or congruencia procesal is obligated)
\end{quote}

\textbf{Witness Clause:}

\begin{lstlisting}[style=witnessstyle, language=Witness]
clause tesis212497_flexibilidad_procesal = oblig(posesion_agraria) IMPLIES (oblig(ejercer_prescripcion_positiva) OR oblig(congruencia_procesal));
\end{lstlisting}

\newpage

\section{Thesis 213224: REVEALING POSSESSION ORIGIN REQUIREMENT FOR ADVERSE POSSESSION}
\label{thesis:213224}

\subsection{Original Thesis Text}

\begin{quote}
                                                   Semanario Judicial de la Federación

                                                  Tesis

 Registro digital: 213224

 Instancia: Tribunales            Octava Época                         Materia(s): Civil
 Colegiados de Circuito

 Tesis: II.2o.180 C               Fuente: Semanario Judicial de la     Tipo: Aislada
                                  Federación.
                                  Tomo XIII, Marzo de 1994, página
                                  428


PRESCRIPCION ADQUISITIVA. BASTA CON REVELAR EL ORIGEN DE LA POSESION Y
AFIRMAR QUE SE POSEE A TITULO DE DUEÑO. (LEGISLACION DEL ESTADO DE MEXICO).


De conformidad con lo dispuesto en el artículo 911 fracción I del Código Civil del Estado de México,
no se exige como requisito para prescribir adquisitivamente, que la posesión esté fundada en justo
título, sino basta con revelar el origen de la posesión y afirmar que se posee a título de dueño, y
esto es así, con el objeto de que el juzgador conozca el hecho o acto generador de la misma, para
poder determinar la calidad de la posesión, si es en concepto de propietario, originaria o derivada,
de buena o de mala fe, para precisar el momento en que debe empezar a contar el plazo de la
prescripción.

SEGUNDO TRIBUNAL COLEGIADO DEL SEGUNDO CIRCUITO.


Amparo directo 14/94. Agustín Juárez Sánchez. 9 de febrero de 1994. Unanimidad de votos.
Ponente: Raúl Solís Solís. Secretario: Pablo Rabanal Arroyo.

Nota: Por ejecutoria de fecha 12 de noviembre de 2003, la Primera Sala declaró improcedente la
contradicción de tesis 26/2003-PS en que participó el presente criterio.




                                             Pág. 1 de 1             Fecha de impresión 01/10/2025
  https://sjf2.scjn.gob.mx/detalle/tesis/213224

\end{quote}

\subsection{Formalization}

\subsubsection{Clause: tesis213224\_principio\_fundamental}

\textit{Thesis-specific principle: Just title not required for adverse possession, but must reveal origin of possession and affirm possession as owner}

\textbf{Intermediate Representation:}

\begin{quote}
ejercer prescripcion positiva is obligated and possession as owner is obligated
\end{quote}

\textbf{Witness Clause:}

\begin{lstlisting}[style=witnessstyle, language=Witness]
clause tesis213224_principio_fundamental = oblig(ejercer_prescripcion_positiva) IMPLIES oblig(causa_generadora_revelada) AND oblig(posesion_como_propietario);
\end{lstlisting}

\subsubsection{Clause: tesis213224\_determinacion\_judicial}

\textit{Supporting logic: Judge needs generating fact/act to determine possession quality and timing}

\textbf{Intermediate Representation:}

\begin{quote}
If causa generadora revelada is obligated, then (possession as owner is obligated and (good faith exercise is obligated and not (bad faith exercise is obligated)))
\end{quote}

\textbf{Witness Clause:}

\begin{lstlisting}[style=witnessstyle, language=Witness]
clause tesis213224_determinacion_judicial = oblig(causa_generadora_revelada) IMPLIES (oblig(posesion_como_propietario) AND (oblig(ejercicio_buena_fe) AND not(oblig(ejercicio_mala_fe))));
\end{lstlisting}

\subsubsection{Clause: tesis213224\_tipo\_posesion}

\textit{Possession type determination: Can determine if original or derived possession}

\textbf{Intermediate Representation:}

\begin{quote}
If causa generadora revelada is obligated, then (possession as owner is obligated and not (temporal possession is obligated))
\end{quote}

\textbf{Witness Clause:}

\begin{lstlisting}[style=witnessstyle, language=Witness]
clause tesis213224_tipo_posesion = oblig(causa_generadora_revelada) IMPLIES (oblig(posesion_como_propietario) AND not(oblig(posesion_temporal)));
\end{lstlisting}

\subsubsection{Clause: tesis213224\_determinacion\_temporal}

\textit{Timing requirement: Must establish when prescription period starts counting}

\textbf{Intermediate Representation:}

\begin{quote}
If causa generadora revelada is obligated, then (posesion tiempo 5 anos is obligated or posesion tiempo 10 anos is obligated)
\end{quote}

\textbf{Witness Clause:}

\begin{lstlisting}[style=witnessstyle, language=Witness]
clause tesis213224_determinacion_temporal = oblig(causa_generadora_revelada) IMPLIES (oblig(posesion_tiempo_5_anos) OR oblig(posesion_tiempo_10_anos));
\end{lstlisting}

\newpage

\section{Thesis 214456: AGRARIAN ADVERSE POSSESSION PREVIOUS POSSESSION TIME COUNTING}
\label{thesis:214456}

\subsection{Original Thesis Text}

\begin{quote}
                                                   Semanario Judicial de la Federación

                                                  Tesis

 Registro digital: 214456

 Instancia: Tribunales            Octava Época                         Materia(s): Administrativa
 Colegiados de Circuito

                                  Fuente: Semanario Judicial de la     Tipo: Aislada
                                  Federación.
                                  Tomo XII, Noviembre de 1993,
                                  página 398


PRESCRIPCION ADQUISITIVA EN MATERIA AGRARIA, SI DEBE CONTARSE EL TIEMPO DE
POSESION ANTERIOR A LA VIGENCIA DE LA NUEVA LEY AGRARIA PARA EL COMPUTO DEL
TERMINO DE LA.


La Ley Agraria actualmente en vigor confirma a los núcleos de población ejidal como propietarios de
las tierras que les han sido dotados o de las que hubieren adquirido por cualquier otro título, según
puede constatarse de lo dispuesto por el artículo 9 de la misma. Sin embargo, a diferencia de la
legislación anterior, se opera una transformación en el régimen de propiedad, pues permite que se
cambie de ejidal a dominio pleno, de acuerdo con lo dispuesto por los artículos 29, 81, 82 y 83,
párrafo primero. Pero no obstante ello, mientras se continúe con el régimen de explotación ejidal se
sigue el sistema de la legislación precedente en cuanto toca a que los ejidatarios únicamente tienen
derecho al uso, aprovechamiento y usufructo de sus parcelas. Asimismo, se trastoca la manera de
adquirir los derechos ejidales aludidos mediante la posesión de las tierras de cultivo, pues ahora el
artículo 48 fija plazos prescriptivos de cinco y diez años para las ocupaciones de buena y mala fe,
respectivamente. De lo anterior puede ya concluirse que la prescripción adquisitiva que prevé la Ley
Agraria, se refiere únicamente a la obtención de los derechos para usufructuar los predios de labor,
pero de ninguna manera trae como consecuencia la apropiación del inmueble mismo, que como se
ha visto, seguirá perteneciendo al núcleo de población ejidal en tanto no cambie el régimen de
propiedad por alguno de los medios instituidos por la ley para tal efecto. En suma, las dos
instituciones que se analizan convergen en un fin común, que es el de otorgar derechos ejidales a
quienes han cumplido con los términos y condiciones legales para obtenerlos; y por tanto, de ello es
posible determinar que no son de naturaleza jurídica distinta, sino que sólo se presenta una
modificación legislativa a los requisitos que deben acatarse para que se produzca la prescripción en
materia agraria; es decir, no se trata de una nueva figura jurídica, sino la continuación de la ya
existente con diferentes requisitos para que se realice. Así pues, todo el tiempo que se hubieren
poseído parcelas ejidales antes de la entrada en vigor de la Ley Agraria, puede y debe ser tomado
en cuenta como parte del plazo requerido para la prescripción, la cual llegará a consumarse si
además se reúne el resto de las exigencias que para tales casos establece el artículo 48 del
ordenamiento legal citado; sin que ello implique una aplicación retroactiva de la ley, porque para
estimarlo así, es necesario que ésta vuelva al pasado para apreciar las condiciones de legalidad de
un acto, o para modificar o suprimir los efectos de un derecho ya realizado; lo cual no ocurre en
casos como el presente en el que sólo se pide el reconocimiento del hecho posesorio como
generador de derechos agrarios, y en tales condiciones, la aplicación del artículo 48 de la Ley
Agraria sobre la posesión material de tierras ejidales de cultivo, anterior a su entrada en vigor, no
constituye una aplicación retroactiva de la misma.



                                             Pág. 1 de 2             Fecha de impresión 01/10/2025
  https://sjf2.scjn.gob.mx/detalle/tesis/214456
                                                  Semanario Judicial de la Federación

SEGUNDO TRIBUNAL COLEGIADO DEL SEGUNDO CIRCUITO.


Amparo directo 604/93. Espiridión Santos Torrijos. 25 de agosto de 1993. Unanimidad de votos.
Ponente: Raúl Solís Solís. Secretario: Joel A. Sierra Palacios.

Nota: Esta tesis contendió en la contradicción 16/94 resuelta por la Segunda Sala, de la que derivó
la tesis 2a./J. 24/94, que aparece publicada en la Gaceta del Semanario Judicial de la Federación,
Octava Época, número 84, diciembre de 1994, página 20, con el rubro: "PRESCRIPCION
POSITIVA EN MATERIA AGRARIA, PARA COMPUTAR LOS PLAZOS PREVISTOS EN EL
ARTICULO 48 DE LA LEY AGRARIA, NO DEBE TOMARSE EN CUENTA EL TIEMPO DE
POSESION ANTERIOR A LA ENTRADA EN VIGOR DE LA LEY ACTUAL. (27 DE FEBRERO DE
1992)."




                                             Pág. 2 de 2           Fecha de impresión 01/10/2025
  https://sjf2.scjn.gob.mx/detalle/tesis/214456

\end{quote}

\subsection{Formalization}

\subsubsection{Clause: tesis214456\_principio\_fundamental}

\textit{Thesis-specific principle: Time of possession before Agrarian Law effective date (Feb 27, 1992) can and should be counted for prescription}

\textbf{Intermediate Representation:}

\begin{quote}
posesion agraria is obligated and temporal possession is obligated and (oblig(posesion\_tiempo\_5\_anos) OR oblig(posesion\_tiempo\_10\_anos)) IMPLIES oblig(ejercer\_prescripcion\_positiva)
\end{quote}

\textbf{Witness Clause:}

\begin{lstlisting}[style=witnessstyle, language=Witness]
clause tesis214456_principio_fundamental = oblig(posesion_agraria) AND oblig(posesion_temporal) AND (oblig(posesion_tiempo_5_anos) OR oblig(posesion_tiempo_10_anos)) IMPLIES oblig(ejercer_prescripcion_positiva);
\end{lstlisting}

\subsubsection{Clause: tesis214456\_no\_retroactividad}

\textit{Supporting logic: Not retroactive application because it only recognizes existing possessory fact}

\textbf{Intermediate Representation:}

\begin{quote}
temporal possession is obligated and posesion agraria is obligated
\end{quote}

\textbf{Witness Clause:}

\begin{lstlisting}[style=witnessstyle, language=Witness]
clause tesis214456_no_retroactividad = oblig(posesion_temporal) AND oblig(posesion_agraria) IMPLIES oblig(proteccion_legal);
\end{lstlisting}

\subsubsection{Clause: tesis214456\_continuidad\_legal}

\textit{Legal continuity: Continuation of existing legal figure with different requirements, not new figure}

\textbf{Intermediate Representation:}

\begin{quote}
posesion agraria is obligated
\end{quote}

\textbf{Witness Clause:}

\begin{lstlisting}[style=witnessstyle, language=Witness]
clause tesis214456_continuidad_legal = oblig(posesion_agraria) IMPLIES oblig(ejercer_prescripcion_positiva);
\end{lstlisting}

\subsubsection{Clause: tesis214456\_plazos\_temporales}

\textit{Time requirements: 5 years for good faith, 10 years for bad faith}

\textbf{Intermediate Representation:}

\begin{quote}
good faith exercise is obligated
\end{quote}

\textbf{Witness Clause:}

\begin{lstlisting}[style=witnessstyle, language=Witness]
clause tesis214456_plazos_temporales = oblig(ejercicio_buena_fe) IMPLIES oblig(posesion_tiempo_5_anos);
\end{lstlisting}

\subsubsection{Clause: tesis214456\_plazos\_mala\_fe}

\textbf{Intermediate Representation:}

\begin{quote}
bad faith exercise is obligated
\end{quote}

\textbf{Witness Clause:}

\begin{lstlisting}[style=witnessstyle, language=Witness]
clause tesis214456_plazos_mala_fe = oblig(ejercicio_mala_fe) IMPLIES oblig(posesion_tiempo_10_anos);
\end{lstlisting}

\newpage

\section{Thesis 215033: ADVERSE POSSESSION ACTION TARGET REQUIREMENTS}
\label{thesis:215033}

\subsection{Original Thesis Text}

\begin{quote}
                                                   Semanario Judicial de la Federación

                                                  Tesis

 Registro digital: 215033

 Instancia: Tribunales             Octava Época                          Materia(s): Civil
 Colegiados de Circuito

                                   Fuente: Semanario Judicial de la      Tipo: Aislada
                                   Federación.
                                   Tomo XII, Septiembre de 1993,
                                   página 281


PRESCRIPCION ADQUISITIVA. CONTRA QUIEN DEBE ENTABLARSE LA DEMANDA DE.


Si el inmueble que se pretende adquirir mediante la acción de prescripción positiva pertenece a la
sociedad legal que se constituyó con motivo del matrimonio celebrado entre el demandado y la
quejosa, aun cuando dicha sociedad no se encuentre inscrita en el Registro Público de la Propiedad
y del Comercio, sino que el inmueble perteneciente a la misma, se encuentra registrado únicamente
a nombre del marido demandado y no de ambos cónyuges, lo que sería necesario para proteger a
los terceros adquirentes de buena fe que contrataren con los cónyuges a fin de evitar de que sean
defraudados por ocultaciones o modificaciones a las capitulaciones matrimoniales que sólo aquéllos
conocen, tal inscripción del inmueble a nombre de la sociedad conyugal no se hace necesaria
cuando los actores son hijos del demandado y la quejosa, porque es de estimarse fundadamente
que conocían el régimen matrimonial que regían sus bienes, por lo que también debieron demandar
a la quejosa, que resultó ser su madre, y no únicamente a su padre que era quien aparecía como
propietario del bien en cuestión y al no haberlo hecho así, es evidente que se violaron en perjuicio
de aquélla las garantías contenidas en los artículos 14 y 16 constitucionales, por no habérsele dado
la oportunidad de ser oída y vencida en juicio.

PRIMER TRIBUNAL COLEGIADO DEL DECIMO QUINTO CIRCUITO.


Amparo en revisión 179/92. María Isabel Márquez de García. 8 de julio de 1992. Unanimidad de
votos. Ponente: Pedro F. Reyes Colín. Secretario: José de Jesús Bernal Juárez.

Amparo en revisión 160/92. María Isabel Márquez de García. 2 de julio de 1992. Unanimidad de
votos. Ponente: Raúl Molina Torres. Secretario: Oralia Barba Ramírez.

Amparo en revisión 170/92. María Isabel Márquez de García. 1o. de julio de 1992. Unanimidad de
votos. Ponente: Raúl Molina Torres. Secretaria: Oralia Barba Ramírez.

Amparo en revisión 81/92. Adela García Márquez. 8 de abril de 1992. Unanimidad de votos.
Ponente: Raúl Molina Torres. Secretaria: Oralia Barba Ramírez.

Nota: El criterio contenido en esta tesis contendió en la contradicción de tesis 33/93, resuelta por la
Tercera Sala, de la que derivó la tesis 3a./J. 24/94, que aparece publicada en la Gaceta del
Semanario Judicial de la Federación, Octava Época, número 80, agosto de 1994, página 22, con el
rubro "SOCIEDAD CONYUGAL, FALTA DE INSCRIPCION. LA DEMANDA DE PRESCRIPCION


                                              Pág. 1 de 2             Fecha de impresión 01/10/2025
  https://sjf2.scjn.gob.mx/detalle/tesis/215033
                                                 Semanario Judicial de la Federación

POSITIVA DEBE ENTABLARSE SOLO EN CONTRA DEL CONYUGE QUE APAREZCA COMO
PROPIETARIO EN EL REGISTRO PUBLICO DE LA PROPIEDAD, Y NO DE AMBOS, AUNQUE
QUIEN EJERCITE LA ACCION CORRESPONDIENTE SEA HIJO DE LOS QUE LA
CONSTITUYERON. (LEGISLACION DEL ESTADO DE BAJA CALIFORNIA Y DISPOSICIONES
AFINES DEL DISTRITO FEDERAL)."




                                            Pág. 2 de 2       Fecha de impresión 01/10/2025
 https://sjf2.scjn.gob.mx/detalle/tesis/215033

\end{quote}

\subsection{Formalization}

\subsubsection{Clause: tesis215033\_principio\_fundamental}

\textit{Thesis-specific principle: Adverse possession action must be directed against all property owners, not just registered owner, to ensure due process}

\textbf{Intermediate Representation:}

\begin{quote}
ejercer prescripcion positiva is obligated and registered ownership is obligated
\end{quote}

\textbf{Witness Clause:}

\begin{lstlisting}[style=witnessstyle, language=Witness]
clause tesis215033_principio_fundamental = oblig(ejercer_prescripcion_positiva) AND oblig(titularidad_registral) IMPLIES oblig(proteccion_legal);
\end{lstlisting}

\subsubsection{Clause: tesis215033\_debido\_proceso}

\textit{Supporting logic: Must give opportunity to be heard and defeated in trial (constitutional guarantee)}

\textbf{Intermediate Representation:}

\begin{quote}
ejercer prescripcion positiva is obligated
\end{quote}

\textbf{Witness Clause:}

\begin{lstlisting}[style=witnessstyle, language=Witness]
clause tesis215033_debido_proceso = oblig(ejercer_prescripcion_positiva) IMPLIES oblig(proteccion_legal);
\end{lstlisting}

\subsubsection{Clause: tesis215033\_conocimiento\_regimen}

\textit{Knowledge requirement: When actors know marital regime, must demand all property owners}

\textbf{Intermediate Representation:}

\begin{quote}
conocimiento vicios is obligated and ejercer prescripcion positiva is obligated
\end{quote}

\textbf{Witness Clause:}

\begin{lstlisting}[style=witnessstyle, language=Witness]
clause tesis215033_conocimiento_regimen = oblig(conocimiento_vicios) AND oblig(ejercer_prescripcion_positiva) IMPLIES oblig(proteccion_legal);
\end{lstlisting}

\subsubsection{Clause: tesis215033\_requisito\_procesal}

\textit{Procedural requirement: Action must be directed against all owners to avoid due process violations}

\textbf{Intermediate Representation:}

\begin{quote}
ejercer prescripcion positiva is obligated
\end{quote}

\textbf{Witness Clause:}

\begin{lstlisting}[style=witnessstyle, language=Witness]
clause tesis215033_requisito_procesal = oblig(ejercer_prescripcion_positiva) IMPLIES oblig(congruencia_procesal);
\end{lstlisting}

\newpage

\section{Thesis 216431: JUST TITLE REQUIREMENT FOR GOOD FAITH ADVERSE POSSESSION}
\label{thesis:216431}

\subsection{Original Thesis Text}

\begin{quote}
                                                   Semanario Judicial de la Federación

                                                  Tesis

 Registro digital: 216431

 Instancia: Tribunales            Octava Época                       Materia(s): Civil
 Colegiados de Circuito

                                  Fuente: Semanario Judicial de la  Tipo: Aislada
                                  Federación.
                                  Tomo XI, Mayo de 1993, página 374


PRESCRIPCION POSITIVA. ES NECESARIO REVELAR LA CAUSA GENERADORA DE LA
POSESION. (LEGISLACION DEL ESTADO DE NAYARIT).


Conforme a la interpretación armónica de lo dispuesto por los artículos 794, 814 y 1135 del Código
Civil para el Estado de Nayarit, se llega a la conclusión, de que los poseedores de buena fe
requieren de justo título como causa generadora de la posesión para que opere la prescripción
positiva; de ahí que, si la quejosa no reveló ni acreditó en forma plena la causa generadora de la
posesión que detenta a título de dueña, no pueda producirse a su favor la prescripción que
demanda. En razón de que el juzgador no estuvo en aptitud de conocer el hecho o acto generador
de dicha posesión, a fin de determinar la calidad de la misma y precisar el momento en que
comienza el plazo de la prescripción.

SEGUNDO TRIBUNAL COLEGIADO DEL DECIMO SEGUNDO CIRCUITO.


Amparo directo 272/92. Juana Betancourt Alvarez. 5 de marzo de 1993. Unanimidad de votos.
Ponente: Nicolás Názar Sevilla. Secretario: Samuel Eduardo Corona Leurette.




                                             Pág. 1 de 1          Fecha de impresión 01/10/2025
  https://sjf2.scjn.gob.mx/detalle/tesis/216431

\end{quote}

\subsection{Formalization}

\subsubsection{Clause: tesis216431\_principio\_fundamental}

\textit{Thesis-specific principle: Good faith possessors require just title as generating cause of possession for adverse possession to operate}

\textbf{Intermediate Representation:}

\begin{quote}
good faith exercise is obligated and possession as owner is obligated
\end{quote}

\textbf{Witness Clause:}

\begin{lstlisting}[style=witnessstyle, language=Witness]
clause tesis216431_principio_fundamental = oblig(ejercicio_buena_fe) AND oblig(posesion_como_propietario) IMPLIES oblig(causa_generadora_probada);
\end{lstlisting}

\subsubsection{Clause: tesis216431\_revelar\_causa}

\textit{Supporting logic: Must reveal and fully prove the generating cause of possession}

\textbf{Intermediate Representation:}

\begin{quote}
ejercer prescripcion positiva is obligated
\end{quote}

\textbf{Witness Clause:}

\begin{lstlisting}[style=witnessstyle, language=Witness]
clause tesis216431_revelar_causa = oblig(ejercer_prescripcion_positiva) IMPLIES oblig(causa_generadora_revelada);
\end{lstlisting}

\subsubsection{Clause: tesis216431\_determinacion\_judicial}

\textit{Judicial determination: Judge needs generating fact/act to determine possession quality and timing}

\textbf{Intermediate Representation:}

\begin{quote}
If causa generadora revelada is obligated, then (possession as owner is obligated and (good faith exercise is obligated and not (bad faith exercise is obligated)))
\end{quote}

\textbf{Witness Clause:}

\begin{lstlisting}[style=witnessstyle, language=Witness]
clause tesis216431_determinacion_judicial = oblig(causa_generadora_revelada) IMPLIES (oblig(posesion_como_propietario) AND (oblig(ejercicio_buena_fe) AND not(oblig(ejercicio_mala_fe))));
\end{lstlisting}

\subsubsection{Clause: tesis216431\_determinacion\_temporal}

\textit{Timing requirement: Must establish when prescription period begins}

\textbf{Intermediate Representation:}

\begin{quote}
If causa generadora revelada is obligated, then (posesion tiempo 5 anos is obligated or posesion tiempo 10 anos is obligated)
\end{quote}

\textbf{Witness Clause:}

\begin{lstlisting}[style=witnessstyle, language=Witness]
clause tesis216431_determinacion_temporal = oblig(causa_generadora_revelada) IMPLIES (oblig(posesion_tiempo_5_anos) OR oblig(posesion_tiempo_10_anos));
\end{lstlisting}

\newpage

\section{Thesis 219823: GOOD FAITH AND BAD FAITH POSSESSOR DEFINITIONS}
\label{thesis:219823}

\subsection{Original Thesis Text}

\begin{quote}
                                                   Semanario Judicial de la Federación

                                                  Tesis

 Registro digital: 219823

 Instancia: Tribunales            Octava Época                         Materia(s): Civil
 Colegiados de Circuito

                                  Fuente: Semanario Judicial de la   Tipo: Aislada
                                  Federación.
                                  Tomo IX, Abril de 1992, página 573


PRESCRIPCION ADQUISITIVA. BUENA O MALA FE DEL POSEEDOR (LEGISLACION DEL
ESTADO DE COAHUILA).


De acuerdo con el artículo 806, del Código Civil vigente en el Estado de Coahuila, la buena fe
consiste en la creencia de que la persona de quien recibió la cosa era dueño de ella y que podía
transmitir su dominio. O bien, dicho de otra manera, es poseedor de buena fe el que entra en la
posesión del inmueble en virtud de una causa generadora o título suficiente para darle derecho de
poseer, así como el que ignora los vicios de su título que le impiden poseer con derecho, es decir,
para que exista la buena fe es indispensable la satisfacción de los siguientes requisitos: un título
suficiente o causa generadora de posesión, ignorancia de vicios de dicho título, si es que existen, y
la creencia fundada de que la cosa le pertenece. En cambio, es poseedor de mala fe el que no tiene
título, causa generadora o modo de adquirir, o el que conoce los vicios del título que le impide
poseer con derecho.

SEGUNDO TRIBUNAL COLEGIADO DEL OCTAVO CIRCUITO.


Amparo directo 479/91. Marco Antonio Velarde Cruz. 21 de enero de 1992. Unanimidad de votos.
Ponente: Julio Ibarrola González. Secretario: José Martín Hernández Simental.




                                             Pág. 1 de 1            Fecha de impresión 01/10/2025
  https://sjf2.scjn.gob.mx/detalle/tesis/219823

\end{quote}

\subsection{Formalization}

\subsubsection{Clause: tesis219823\_principio\_fundamental}

\textit{Thesis-specific principle: Good faith requires sufficient title/generating cause, ignorance of title defects, and founded belief of ownership}

\textbf{Intermediate Representation:}

\begin{quote}
good faith exercise is obligated and not (conocimiento vicios is obligated)
\end{quote}

\textbf{Witness Clause:}

\begin{lstlisting}[style=witnessstyle, language=Witness]
clause tesis219823_principio_fundamental = oblig(ejercicio_buena_fe) IMPLIES oblig(causa_generadora_probada) AND not(oblig(conocimiento_vicios));
\end{lstlisting}

\subsubsection{Clause: tesis219823\_mala\_fe}

\textit{Supporting logic: Bad faith possessor has no title/generating cause OR knows title defects that prevent rightful possession}

\textbf{Intermediate Representation:}

\begin{quote}
If bad faith exercise is obligated, then (not (proven generating cause is obligated) or conocimiento vicios is obligated)
\end{quote}

\textbf{Witness Clause:}

\begin{lstlisting}[style=witnessstyle, language=Witness]
clause tesis219823_mala_fe = oblig(ejercicio_mala_fe) IMPLIES (not(oblig(causa_generadora_probada)) OR oblig(conocimiento_vicios));
\end{lstlisting}

\subsubsection{Clause: tesis219823\_requisitos\_buena\_fe}

\textit{Good faith requirements: Must have sufficient title and ignorance of defects}

\textbf{Intermediate Representation:}

\begin{quote}
good faith exercise is obligated
\end{quote}

\textbf{Witness Clause:}

\begin{lstlisting}[style=witnessstyle, language=Witness]
clause tesis219823_requisitos_buena_fe = oblig(ejercicio_buena_fe) IMPLIES oblig(causa_generadora_probada);
\end{lstlisting}

\subsubsection{Clause: tesis219823\_condiciones\_mala\_fe}

\textit{Bad faith conditions: No title or knowledge of defects}

\textbf{Intermediate Representation:}

\begin{quote}
(not(oblig(causa\_generadora\_probada)) OR oblig(conocimiento\_vicios)) IMPLIES oblig(ejercicio\_mala\_fe)
\end{quote}

\textbf{Witness Clause:}

\begin{lstlisting}[style=witnessstyle, language=Witness]
clause tesis219823_condiciones_mala_fe = (not(oblig(causa_generadora_probada)) OR oblig(conocimiento_vicios)) IMPLIES oblig(ejercicio_mala_fe);
\end{lstlisting}

\newpage

\section{Thesis 221460: ORIGIN AND TIMING REQUIREMENT FOR ADVERSE POSSESSION}
\label{thesis:221460}

\subsection{Original Thesis Text}

\begin{quote}
                                                   Semanario Judicial de la Federación

                                                  Tesis

 Registro digital: 221460

 Instancia: Tribunales            Octava Época                         Materia(s): Civil
 Colegiados de Circuito

                                  Fuente: Semanario Judicial de la     Tipo: Aislada
                                  Federación.
                                  Tomo VIII, Noviembre de 1991,
                                  página 269


PRESCRIPCION POSITIVA. PARA QUE OPERE ES NECESARIO SEÑALAR EL HECHO O EL
ACTO QUE LA ORIGINO Y CUANDO OCURRIO.


Para usucapir es imprescindible señalar el hecho o el acto que la originó y cuándo ocurrió, a fin de
que sea posible determinar la naturaleza de dicha posesión, o sea si es originaria o derivada, de
buena o mala fe, y cuál es el tiempo de su duración para que se consume la prescripción positiva.

TRIBUNAL COLEGIADO DEL VIGESIMO CIRCUITO.


Amparo directo 568/90. Noé Ramírez Ovalle. 4 de abril de 1991. Unanimidad de votos. Ponente:
Angel Suárez Torres. Secretario: Casto Ambrosio Domínguez Bermúdez.




                                             Pág. 1 de 1             Fecha de impresión 01/10/2025
  https://sjf2.scjn.gob.mx/detalle/tesis/221460

\end{quote}

\subsection{Formalization}

\subsubsection{Clause: tesis221460\_principio\_fundamental}

\textit{Thesis-specific principle: Must specify the fact or act that originated possession and when it occurred for adverse possession to operate}

\textbf{Intermediate Representation:}

\begin{quote}
ejercer prescripcion positiva is obligated
\end{quote}

\textbf{Witness Clause:}

\begin{lstlisting}[style=witnessstyle, language=Witness]
clause tesis221460_principio_fundamental = oblig(ejercer_prescripcion_positiva) IMPLIES oblig(causa_generadora_revelada);
\end{lstlisting}

\subsubsection{Clause: tesis221460\_determinacion\_naturaleza}

\textit{Supporting logic: Judge needs this information to determine possession nature (original or derived, good or bad faith)}

\textbf{Intermediate Representation:}

\begin{quote}
If causa generadora revelada is obligated, then (possession as owner is obligated and not (temporal possession is obligated)) and (good faith exercise is obligated and not (bad faith exercise is obligated))
\end{quote}

\textbf{Witness Clause:}

\begin{lstlisting}[style=witnessstyle, language=Witness]
clause tesis221460_determinacion_naturaleza = oblig(causa_generadora_revelada) IMPLIES (oblig(posesion_como_propietario) AND not(oblig(posesion_temporal))) AND (oblig(ejercicio_buena_fe) AND not(oblig(ejercicio_mala_fe)));
\end{lstlisting}

\subsubsection{Clause: tesis221460\_determinacion\_temporal}

\textit{Timing determination: Must establish duration of possession for prescription to be consummated}

\textbf{Intermediate Representation:}

\begin{quote}
If causa generadora revelada is obligated, then (posesion tiempo 5 anos is obligated or posesion tiempo 10 anos is obligated)
\end{quote}

\textbf{Witness Clause:}

\begin{lstlisting}[style=witnessstyle, language=Witness]
clause tesis221460_determinacion_temporal = oblig(causa_generadora_revelada) IMPLIES (oblig(posesion_tiempo_5_anos) OR oblig(posesion_tiempo_10_anos));
\end{lstlisting}

\subsubsection{Clause: tesis221460\_requisito\_usucapion}

\textit{Usucapion requirement: Essential for adverse possession to operate}

\textbf{Intermediate Representation:}

\begin{quote}
ejercer prescripcion positiva is obligated
\end{quote}

\textbf{Witness Clause:}

\begin{lstlisting}[style=witnessstyle, language=Witness]
clause tesis221460_requisito_usucapion = oblig(ejercer_prescripcion_positiva) IMPLIES oblig(causa_generadora_revelada);
\end{lstlisting}

\newpage

\section{Thesis 221758: REVEALING AND PROVING POSSESSION ORIGIN REQUIREMENT}
\label{thesis:221758}

\subsection{Original Thesis Text}

\begin{quote}
                                                   Semanario Judicial de la Federación

                                                  Tesis

 Registro digital: 221758

 Instancia: Tribunales            Octava Época                          Materia(s): Civil
 Colegiados de Circuito

                                  Fuente: Semanario Judicial de la   Tipo: Aislada
                                  Federación.
                                  Tomo VIII, Octubre de 1991, página
                                  236


PRESCRIPCION POSITIVA. ES NECESARIO REVELAR Y ACREDITAR EL ORIGEN DE LA
POSESION. (LEGISLACION DEL ESTADO DE BAJA CALIFORNIA).


La posesión en concepto de dueño admite doctrinalmente tres formas de poseer: a) Poseer con
justo título objetivamente válido, b) Poseer con justo título subjetivamente válido y, c) Poseer sin
título, pero con animus domini; y en la especie, la frase que emplea el código sustantivo local
poseer en concepto de propietario, abarca el título subjetivamente válido, que se funda en el título o
causa que se cree suficiente para adquirir el dominio por prescripción, de donde se infiere que es
necesario revelar y acreditar el origen de la posesión, para poder determinar si es en calidad de
propietario, originaria o derivada, de buena o mala fe, y así precisar el momento en que debe
empezar a contar el plazo de la prescripción y, en el caso concreto, tales circunstancias no se
acreditaron, pues la frase "me la dio", no es suficiente para acreditar la causa generadora de su
posesión en concepto de propietario, siendo irrelevante que los testigos hayan sido contestes y
uniformes en manifestar que el actor siempre se ha conducido respecto del inmueble como dueño y
señor, toda vez que los mismos no pueden subsanar las omisiones del actor, en cuanto al origen de
su posesión.

PRIMER TRIBUNAL COLEGIADO DEL DECIMO QUINTO CIRCUITO.


Amparo directo 381/87. Urbano Bahena Sales. 3 de mayo de 1988. Unanimidad de votos. Ponente:
Miguel Angel Morales Hernández. Secretario: Héctor Gómez Hernández.




                                             Pág. 1 de 1             Fecha de impresión 01/10/2025
  https://sjf2.scjn.gob.mx/detalle/tesis/221758

\end{quote}

\subsection{Formalization}

\subsubsection{Clause: tesis221758\_principio\_fundamental}

\textit{Thesis-specific principle: Must reveal and prove the origin of possession to determine possession quality and timing}

\textbf{Intermediate Representation:}

\begin{quote}
ejercer prescripcion positiva is obligated and proven generating cause is obligated
\end{quote}

\textbf{Witness Clause:}

\begin{lstlisting}[style=witnessstyle, language=Witness]
clause tesis221758_principio_fundamental = oblig(ejercer_prescripcion_positiva) IMPLIES oblig(causa_generadora_revelada) AND oblig(causa_generadora_probada);
\end{lstlisting}

\subsubsection{Clause: tesis221758\_determinacion\_posesion}

\textit{Supporting logic: Judge needs this information to determine if possession is as owner, original or derived, good or bad faith}

\textbf{Intermediate Representation:}

\begin{quote}
If causa generadora revelada is obligated, then (possession as owner is obligated and (good faith exercise is obligated and not (bad faith exercise is obligated)))
\end{quote}

\textbf{Witness Clause:}

\begin{lstlisting}[style=witnessstyle, language=Witness]
clause tesis221758_determinacion_posesion = oblig(causa_generadora_revelada) IMPLIES (oblig(posesion_como_propietario) AND (oblig(ejercicio_buena_fe) AND not(oblig(ejercicio_mala_fe))));
\end{lstlisting}

\subsubsection{Clause: tesis221758\_determinacion\_temporal}

\textit{Timing determination: Must establish when prescription period should start}

\textbf{Intermediate Representation:}

\begin{quote}
If causa generadora revelada is obligated, then (posesion tiempo 5 anos is obligated or posesion tiempo 10 anos is obligated)
\end{quote}

\textbf{Witness Clause:}

\begin{lstlisting}[style=witnessstyle, language=Witness]
clause tesis221758_determinacion_temporal = oblig(causa_generadora_revelada) IMPLIES (oblig(posesion_tiempo_5_anos) OR oblig(posesion_tiempo_10_anos));
\end{lstlisting}

\subsubsection{Clause: tesis221758\_requisito\_prueba}

\textit{Proof requirement: Insufficient proof like "me la dio" cannot establish generating cause}

\textbf{Intermediate Representation:}

\begin{quote}
proven generating cause is obligated
\end{quote}

\textbf{Witness Clause:}

\begin{lstlisting}[style=witnessstyle, language=Witness]
clause tesis221758_requisito_prueba = oblig(causa_generadora_probada) IMPLIES oblig(causa_generadora_revelada);
\end{lstlisting}

\newpage

\section{Thesis 222153: JUST TITLE REQUIREMENT FOR GOOD FAITH POSSESSION}
\label{thesis:222153}

\subsection{Original Thesis Text}

\begin{quote}
                                                   Semanario Judicial de la Federación

                                                  Tesis

 Registro digital: 222153

 Instancia: Tribunales            Octava Época                           Materia(s): Civil
 Colegiados de Circuito

                                  Fuente: Semanario Judicial de la       Tipo: Aislada
                                  Federación.
                                  Tomo VIII, Agosto de 1991, página
                                  204


PRESCRIPCION ADQUISITIVA. JUSTO TITULO (LEGISLACION DEL ESTADO DE JALISCO).


El justo título, aun cuando no es absolutamente necesario para prescribir, no ha sido desterrado del
Código Civil del Estado de Jalisco, pues a él corresponden las nociones de título objetiva y
subjetivamente válido a que se hace referencia en el artículo 849, en la medida en que previene, en
lo que interesa, que es poseedor de buena fe el que entra en la posesión en virtud de un título
suficiente para darle derecho a poseer y el que ignora los vicios de su título que le impide poseer
con derecho.

PRIMER TRIBUNAL COLEGIADO EN MATERIA CIVIL DEL TERCER CIRCUITO.


Amparo directo 227/91. Jorge Pagua Montaño y Carmen Servín Pagua. 28 de junio de 1991.
Unanimidad de votos. Ponente: Carlos Arturo González Zárate. Secretario: Juan Bonilla Pizano.

Amparo directo 840/88. Carlota Reynoso Castillo. 3 de febrero de 1989. Unanimidad de votos.
Ponente: Carlos Arturo González Zárate. Secretario: Juan Bonilla Pizano.

Nota: Este criterio ha integrado la jurisprudencia III.1o.C. J/6, publicada en el Semanario Judicial de
la Federación y su Gaceta, Novena Época, Tomo II, septiembre de 1995, página 475, de rubro:
"PRESCRIPCION ADQUISITIVA, JUSTO TITULO."




                                             Pág. 1 de 1              Fecha de impresión 01/10/2025
  https://sjf2.scjn.gob.mx/detalle/tesis/222153

\end{quote}

\subsection{Formalization}

\subsubsection{Clause: tesis222153\_principio\_fundamental}

\textit{Thesis-specific principle: Just title is not absolutely necessary for adverse possession, but good faith requires sufficient title and ignorance of defects}

\textbf{Intermediate Representation:}

\begin{quote}
good faith exercise is obligated and not (conocimiento vicios is obligated)
\end{quote}

\textbf{Witness Clause:}

\begin{lstlisting}[style=witnessstyle, language=Witness]
clause tesis222153_principio_fundamental = oblig(ejercicio_buena_fe) IMPLIES oblig(causa_generadora_probada) AND not(oblig(conocimiento_vicios));
\end{lstlisting}

\subsubsection{Clause: tesis222153\_titulo\_suficiente}

\textit{Supporting logic: Good faith possessor enters possession by virtue of sufficient title to give right to possess}

\textbf{Intermediate Representation:}

\begin{quote}
good faith exercise is obligated
\end{quote}

\textbf{Witness Clause:}

\begin{lstlisting}[style=witnessstyle, language=Witness]
clause tesis222153_titulo_suficiente = oblig(ejercicio_buena_fe) IMPLIES oblig(causa_generadora_probada);
\end{lstlisting}

\subsubsection{Clause: tesis222153\_ignorancia\_vicios}

\textit{Ignorance requirement: Must ignore title defects that prevent rightful possession}

\textbf{Intermediate Representation:}

\begin{quote}
good faith exercise is obligated
\end{quote}

\textbf{Witness Clause:}

\begin{lstlisting}[style=witnessstyle, language=Witness]
clause tesis222153_ignorancia_vicios = oblig(ejercicio_buena_fe) IMPLIES not(oblig(conocimiento_vicios));
\end{lstlisting}

\subsubsection{Clause: tesis222153\_flexibilidad\_titulo}

\textit{Just title flexibility: Not absolutely necessary but relevant for good faith determination}

\textbf{Intermediate Representation:}

\begin{quote}
If ejercer prescripcion positiva is obligated, then (good faith exercise is obligated or bad faith exercise is obligated)
\end{quote}

\textbf{Witness Clause:}

\begin{lstlisting}[style=witnessstyle, language=Witness]
clause tesis222153_flexibilidad_titulo = oblig(ejercer_prescripcion_positiva) IMPLIES (oblig(ejercicio_buena_fe) OR oblig(ejercicio_mala_fe));
\end{lstlisting}

\newpage

\section{Thesis 222185: IMPERFECT GENERATING ACT FOR ADVERSE POSSESSION}
\label{thesis:222185}

\subsection{Original Thesis Text}

\begin{quote}
                                                   Semanario Judicial de la Federación

                                                  Tesis

 Registro digital: 222185

 Instancia: Tribunales            Octava Época                        Materia(s): Civil
 Colegiados de Circuito

                                  Fuente: Semanario Judicial de la    Tipo: Aislada
                                  Federación.
                                  Tomo VIII, Agosto de 1991, página
                                  229


USUCAPION. EL ACTO GENERADOR DE LA POSESION NO REQUIERE SER PERFECTO
CONFORME A LA LEY.


La circunstancia de que algún acto generador de la posesión no sea, conforme a la ley, suficiente
para transmitir el dominio pleno de una cosa, no implica, que tal acto, aun cuando adolezca de
vicios, no pueda invocarse como causa generadora de la posesión, para efectos de la usucapión,
pues tratándose de dicha acción, y para acreditar la buena fe del poseedor, sólo se requiere que
éste hubiera entrado en posesión de la cosa, mediante un acto que, en apariencia, reúne los
requisitos legales, y que produzcan la convicción al adquirente de que el acto es bastante para
transferirle los derechos del dominio.

SEGUNDO TRIBUNAL COLEGIADO DEL SEXTO CIRCUITO.


Amparo directo 409/89. Matilde Juárez Salazar. 23 de mayo de 1990. Unanimidad de votos.
Ponente: Arnoldo Nájera Virgen. Secretario: Guillermo Báez Pérez.




                                             Pág. 1 de 1          Fecha de impresión 01/10/2025
  https://sjf2.scjn.gob.mx/detalle/tesis/222185

\end{quote}

\subsection{Formalization}

\subsubsection{Clause: tesis222185\_principio\_fundamental}

\textit{Thesis-specific principle: Generating act of possession does not need to be perfect according to law, even acts with defects can serve as generating cause}

\textbf{Intermediate Representation:}

\begin{quote}
ejercer prescripcion positiva is obligated
\end{quote}

\textbf{Witness Clause:}

\begin{lstlisting}[style=witnessstyle, language=Witness]
clause tesis222185_principio_fundamental = oblig(ejercer_prescripcion_positiva) IMPLIES oblig(causa_generadora_probada);
\end{lstlisting}

\subsubsection{Clause: tesis222185\_buena\_fe\_aparente}

\textit{Supporting logic: For good faith, only requires act that appears to meet legal requirements and produces conviction in acquirer}

\textbf{Intermediate Representation:}

\begin{quote}
good faith exercise is obligated
\end{quote}

\textbf{Witness Clause:}

\begin{lstlisting}[style=witnessstyle, language=Witness]
clause tesis222185_buena_fe_aparente = oblig(ejercicio_buena_fe) IMPLIES oblig(causa_generadora_probada);
\end{lstlisting}

\subsubsection{Clause: tesis222185\_actos\_viciosos}

\textit{Defective acts validity: Even acts with defects can be invoked as generating cause for usucapion}

\textbf{Intermediate Representation:}

\begin{quote}
conocimiento vicios is obligated and proven generating cause is obligated
\end{quote}

\textbf{Witness Clause:}

\begin{lstlisting}[style=witnessstyle, language=Witness]
clause tesis222185_actos_viciosos = oblig(conocimiento_vicios) AND oblig(causa_generadora_probada) IMPLIES oblig(ejercer_prescripcion_positiva);
\end{lstlisting}

\subsubsection{Clause: tesis222185\_validez\_aparente}

\textit{Apparent validity requirement: Act must appear to meet legal requirements to produce conviction}

\textbf{Intermediate Representation:}

\begin{quote}
proven generating cause is obligated
\end{quote}

\textbf{Witness Clause:}

\begin{lstlisting}[style=witnessstyle, language=Witness]
clause tesis222185_validez_aparente = oblig(causa_generadora_probada) IMPLIES oblig(ejercicio_buena_fe);
\end{lstlisting}

\newpage

\section{Thesis 223431: GOOD/BAD FAITH AS CONDITIONS NOT QUALITIES OF POSSESSION}
\label{thesis:223431}

\subsection{Original Thesis Text}

\begin{quote}
                                                   Semanario Judicial de la Federación

                                                  Tesis

 Registro digital: 223431

 Instancia: Tribunales              Octava Época                            Materia(s): Civil
 Colegiados de Circuito

                                    Fuente: Semanario Judicial de la        Tipo: Aislada
                                    Federación.
                                    Tomo VII, Marzo de 1991, página
                                    194


PRESCRIPCION ADQUISITIVA. LA BUENA O MALA FE NO FORMAN PARTE DE LAS
CUALIDADES DE LA POSESION YA QUE SON CONDICIONES DE LA MISMA. (LEGISLACION
DEL ESTADO DE JALISCO).


El artículo 1180 del Código Civil del Estado de Jalisco, establece las cualidades que debe tener la
posesión a fin de que sea apta para prescribir, puesto que proviene: "La posesión necesaria para
prescribir debe ser: I.- En concepto de propietario; II.- Pacífica; III.- Continua; IV.- Pública; (esto es,
no alude a la buena o a la mala fe como un requisito que deba tener la posesión a fin de que sea
apta para prescribir); en cambio el artículo siguiente (el 1181) se refiere que los bienes inmuebles
prescriben en cinco años cuando, además de ser en concepto de dueño, pacífica, continua y
públicamente, la posesión es de buena fe, y en diez cuando es de mala fe. Una correcta
interpretación de los artículos citados permite concluir que la buena o la mala fe no forman parte de
las cualidades que debe reunir la posesión a fin de que sea apta para usucapir, sino que por ser
útiles únicamente para aumentar o disminuir el tiempo que la ley señala como suficiente para lograr
la prescripción, en realidad constituyen condiciones (no cualidades) de dicha posesión.
Consiguientemente, se viola el principio de congruencia en que descansa toda sentencia si al
pronunciarse ésta no se hace el estudio de la buena y de la mala fe del hecho posesorio, aunque
los actores hayan manifestado que han sido poseedores de buena fe, porque estando satisfechas
las cualidades de la posesión (en concepto de propietario, pacífica, continua y pública) debe
analizarse, en vía de consecuencia, su duración.

TERCER TRIBUNAL COLEGIADO EN MATERIA CIVIL DEL TERCER CIRCUITO.


Amparo directo 806/90. Antonio Zaragoza Romero. 15 de diciembre de 1990. Unanimidad de votos.
Ponente: Jorge Figueroa Cacho. Secretario: Luis Rubén Baltazar Aceves.

Nota: Esta tesis contendió en la contradicción 18/2005-PS resuelta por la Primera Sala, de la que
derivó la tesis 1a./J. 200/2005, que aparece publicada en el Semanario Judicial de la Federación y
su Gaceta, Novena Época, Tomo XXIII, febrero de 2006, página 441, con el rubro:
"PRESCRIPCIÓN POSITIVA. SI LA ACCIÓN SE EJERCE CON BASE EN LA POSESIÓN DE
BUENA FE, EL JUZGADOR SE ENCUENTRA IMPEDIDO PARA ANALIZAR DE OFICIO LA
POSESIÓN DE MALA FE."




                                               Pág. 1 de 2              Fecha de impresión 01/10/2025
  https://sjf2.scjn.gob.mx/detalle/tesis/223431
                                                Semanario Judicial de la Federación




                                           Pág. 2 de 2       Fecha de impresión 01/10/2025
https://sjf2.scjn.gob.mx/detalle/tesis/223431

\end{quote}

\subsection{Formalization}

\subsubsection{Clause: tesis223431\_principio\_fundamental}

\textit{Thesis-specific principle: Good/bad faith are conditions (not qualities) of possession, used only to determine time periods for prescription}

\textbf{Intermediate Representation:}

\begin{quote}
possession as owner is obligated and (oblig(ejercicio\_buena\_fe) AND not(oblig(ejercicio\_mala\_fe))) IMPLIES oblig(ejercer\_prescripcion\_positiva)
\end{quote}

\textbf{Witness Clause:}

\begin{lstlisting}[style=witnessstyle, language=Witness]
clause tesis223431_principio_fundamental = oblig(posesion_como_propietario) AND (oblig(ejercicio_buena_fe) AND not(oblig(ejercicio_mala_fe))) IMPLIES oblig(ejercer_prescripcion_positiva);
\end{lstlisting}

\subsubsection{Clause: tesis223431\_analisis\_judicial}

\textit{Supporting logic: Judge must analyze good/bad faith even if actors claim good faith, to determine appropriate time period}

\textbf{Intermediate Representation:}

\begin{quote}
If ejercer prescripcion positiva is obligated, then (good faith exercise is obligated and not (bad faith exercise is obligated))
\end{quote}

\textbf{Witness Clause:}

\begin{lstlisting}[style=witnessstyle, language=Witness]
clause tesis223431_analisis_judicial = oblig(ejercer_prescripcion_positiva) IMPLIES (oblig(ejercicio_buena_fe) AND not(oblig(ejercicio_mala_fe)));
\end{lstlisting}

\subsubsection{Clause: tesis223431\_plazos\_temporales}

\textit{Time determination: Good faith requires 5 years, bad faith requires 10 years}

\textbf{Intermediate Representation:}

\begin{quote}
good faith exercise is obligated
\end{quote}

\textbf{Witness Clause:}

\begin{lstlisting}[style=witnessstyle, language=Witness]
clause tesis223431_plazos_temporales = oblig(ejercicio_buena_fe) IMPLIES oblig(posesion_tiempo_5_anos);
\end{lstlisting}

\subsubsection{Clause: tesis223431\_plazos\_mala\_fe}

\textbf{Intermediate Representation:}

\begin{quote}
bad faith exercise is obligated
\end{quote}

\textbf{Witness Clause:}

\begin{lstlisting}[style=witnessstyle, language=Witness]
clause tesis223431_plazos_mala_fe = oblig(ejercicio_mala_fe) IMPLIES oblig(posesion_tiempo_10_anos);
\end{lstlisting}

\subsubsection{Clause: tesis223431\_congruencia}

\textit{Congruence requirement: Judge must analyze good/bad faith to avoid violating congruence principle}

\textbf{Intermediate Representation:}

\begin{quote}
ejercer prescripcion positiva is obligated
\end{quote}

\textbf{Witness Clause:}

\begin{lstlisting}[style=witnessstyle, language=Witness]
clause tesis223431_congruencia = oblig(ejercer_prescripcion_positiva) IMPLIES oblig(congruencia_procesal);
\end{lstlisting}

\newpage

\section{Thesis 224038: PROOF REQUIREMENT FOR GENERATING CAUSE OF POSSESSION}
\label{thesis:224038}

\subsection{Original Thesis Text}

\begin{quote}
                                                   Semanario Judicial de la Federación

                                                  Tesis

 Registro digital: 224038

 Instancia: Tribunales            Octava Época                         Materia(s): Civil
 Colegiados de Circuito

                                  Fuente: Semanario Judicial de la     Tipo: Aislada
                                  Federación.
                                  Tomo VII, Enero de 1991, página
                                  357


PRESCRIPCION ADQUISITIVA. LA CAUSA GENERADORA DE LA POSESION DEBE
PROBARSE EN EL JUICIO.


En el juicio sobre prescripción adquisitiva no basta con que el actor revele la causa que generó su
posesión, ni que se considere asimismo propietario, sino que debe justificar además la existencia de
aquella causa con prueba directa, pues si no lo hace, tampoco se está en condiciones de
determinar si su posesión es de buena o mala fe y en qué momentos empezó y se consumó la
usucapión.

TRIBUNAL COLEGIADO DEL DECIMO CIRCUITO.


Amparo directo 78/90. María Esperanza Pérez Gómez. 4 de mayo de 1990. Unanimidad de votos.
Ponente: Moisés Duarte Aguíñiga. Secretario: Juan García Orozco.




                                             Pág. 1 de 1             Fecha de impresión 01/10/2025
  https://sjf2.scjn.gob.mx/detalle/tesis/224038

\end{quote}

\subsection{Formalization}

\subsubsection{Clause: tesis224038\_principio\_fundamental}

\textit{Thesis-specific principle: Generating cause of possession must be proven in trial, not just revealed or claimed}

\textbf{Intermediate Representation:}

\begin{quote}
ejercer prescripcion positiva is obligated
\end{quote}

\textbf{Witness Clause:}

\begin{lstlisting}[style=witnessstyle, language=Witness]
clause tesis224038_principio_fundamental = oblig(ejercer_prescripcion_positiva) IMPLIES oblig(causa_generadora_probada);
\end{lstlisting}

\subsubsection{Clause: tesis224038\_prueba\_directa}

\textit{Supporting logic: Must justify existence of generating cause with direct proof, not just revelation}

\textbf{Intermediate Representation:}

\begin{quote}
proven generating cause is obligated
\end{quote}

\textbf{Witness Clause:}

\begin{lstlisting}[style=witnessstyle, language=Witness]
clause tesis224038_prueba_directa = oblig(causa_generadora_probada) IMPLIES oblig(causa_generadora_revelada);
\end{lstlisting}

\subsubsection{Clause: tesis224038\_determinacion\_fe}

\textit{Judicial determination: Without proof, cannot determine if possession is good or bad faith}

\textbf{Intermediate Representation:}

\begin{quote}
If proven generating cause is obligated, then (good faith exercise is obligated and not (bad faith exercise is obligated))
\end{quote}

\textbf{Witness Clause:}

\begin{lstlisting}[style=witnessstyle, language=Witness]
clause tesis224038_determinacion_fe = oblig(causa_generadora_probada) IMPLIES (oblig(ejercicio_buena_fe) AND not(oblig(ejercicio_mala_fe)));
\end{lstlisting}

\subsubsection{Clause: tesis224038\_determinacion\_temporal}

\textit{Timing determination: Must establish when usucapion began and was consummated}

\textbf{Intermediate Representation:}

\begin{quote}
If proven generating cause is obligated, then (posesion tiempo 5 anos is obligated or posesion tiempo 10 anos is obligated)
\end{quote}

\textbf{Witness Clause:}

\begin{lstlisting}[style=witnessstyle, language=Witness]
clause tesis224038_determinacion_temporal = oblig(causa_generadora_probada) IMPLIES (oblig(posesion_tiempo_5_anos) OR oblig(posesion_tiempo_10_anos));
\end{lstlisting}

\newpage

\section{Thesis 224737: REVEALING AND PROVING GENERATING CAUSE FOR USUCAPION}
\label{thesis:224737}

\subsection{Original Thesis Text}

\begin{quote}
                                                   Semanario Judicial de la Federación

                                                  Tesis

 Registro digital: 224737

 Instancia: Tribunales            Octava Época                         Materia(s): Civil
 Colegiados de Circuito

                                  Fuente: Semanario Judicial de la     Tipo: Aislada
                                  Federación.
                                  Tomo VI, Segunda Parte-1, Julio-
                                  Diciembre de 1990, página 304


USUCAPION, NECESIDAD DE REVELAR Y PROBAR LA CAUSA GENERADORA DE LA
POSESION PARA QUE PROSPERE LA. (LEGISLACION DEL ESTADO DE JALISCO).


Aun cuando el Código Civil del Estado de Jalisco no exige que la posesión necesaria para adquirir
por prescripción se base en justo título, sí es menester que se revele la causa que la originó, para
poder determinar la calidad con que se ejerce, desde cuándo y si es de buena o mala fe, por lo que,
además, no basta con afirmar los hechos que constituyen esa causa, sino que es necesario
probarlos.

SEGUNDO TRIBUNAL COLEGIADO EN MATERIA CIVIL DEL TERCER CIRCUITO.


Amparo directo 418/90. Susana Dodero Ceceña. 6 de septiembre de 1990. Unanimidad de votos.
Ponente: Gilda Rincón Orta. Secretario: Enrique Gómez Mendoza.




                                             Pág. 1 de 1             Fecha de impresión 01/10/2025
  https://sjf2.scjn.gob.mx/detalle/tesis/224737

\end{quote}

\subsection{Formalization}

\subsubsection{Clause: tesis224737\_principio\_fundamental}

\textit{Thesis-specific principle: Must reveal and prove the generating cause of possession for usucapion to prosper, even though just title not required}

\textbf{Intermediate Representation:}

\begin{quote}
ejercer prescripcion positiva is obligated and proven generating cause is obligated
\end{quote}

\textbf{Witness Clause:}

\begin{lstlisting}[style=witnessstyle, language=Witness]
clause tesis224737_principio_fundamental = oblig(ejercer_prescripcion_positiva) IMPLIES oblig(causa_generadora_revelada) AND oblig(causa_generadora_probada);
\end{lstlisting}

\subsubsection{Clause: tesis224737\_determinacion\_posesion}

\textit{Supporting logic: Must reveal cause to determine possession quality, timing, and good/bad faith}

\textbf{Intermediate Representation:}

\begin{quote}
If causa generadora revelada is obligated, then (possession as owner is obligated and (good faith exercise is obligated and not (bad faith exercise is obligated)))
\end{quote}

\textbf{Witness Clause:}

\begin{lstlisting}[style=witnessstyle, language=Witness]
clause tesis224737_determinacion_posesion = oblig(causa_generadora_revelada) IMPLIES (oblig(posesion_como_propietario) AND (oblig(ejercicio_buena_fe) AND not(oblig(ejercicio_mala_fe))));
\end{lstlisting}

\subsubsection{Clause: tesis224737\_requisito\_prueba}

\textit{Proof requirement: Not sufficient to just affirm facts, must prove the generating cause}

\textbf{Intermediate Representation:}

\begin{quote}
proven generating cause is obligated
\end{quote}

\textbf{Witness Clause:}

\begin{lstlisting}[style=witnessstyle, language=Witness]
clause tesis224737_requisito_prueba = oblig(causa_generadora_probada) IMPLIES oblig(causa_generadora_revelada);
\end{lstlisting}

\subsubsection{Clause: tesis224737\_determinacion\_temporal}

\textit{Timing determination: Must establish when possession began and if good or bad faith}

\textbf{Intermediate Representation:}

\begin{quote}
If causa generadora revelada is obligated, then (posesion tiempo 5 anos is obligated or posesion tiempo 10 anos is obligated)
\end{quote}

\textbf{Witness Clause:}

\begin{lstlisting}[style=witnessstyle, language=Witness]
clause tesis224737_determinacion_temporal = oblig(causa_generadora_revelada) IMPLIES (oblig(posesion_tiempo_5_anos) OR oblig(posesion_tiempo_10_anos));
\end{lstlisting}

\newpage

\section{Thesis 224820: POSSESSION IN CONCEPT OF OWNER REQUIREMENTS FOR ADVERSE POSSESSION}
\label{thesis:224820}

\subsection{Original Thesis Text}

\begin{quote}
                                                   Semanario Judicial de la Federación

                                                  Tesis

 Registro digital: 224820

 Instancia: Tribunales             Octava Época                           Materia(s): Civil
 Colegiados de Circuito

 Tesis: I.4o.C. J/30               Fuente: Semanario Judicial de la       Tipo: Jurisprudencia
                                   Federación.
                                   Tomo VI, Segunda Parte-1, Julio-
                                   Diciembre de 1990, página 385


PRESCRIPCION ADQUISITIVA. HECHOS SUSCEPTIBLES DE GENERAR LA POSESION APTA
PARA LA.


Conforme a los artículos 1151 y 1152 del Código Civil para el Distrito Federal, la posesión necesaria
para prescribir debe ser en concepto de propietario, pacífica, continua, pública y por el tiempo que
señala el segundo de esos preceptos, según se trate de posesión de buena o de mala fe, o de la
que hubiera sido inscrita en el Registro Público de la Propiedad y del Comercio. Esta institución,
como medio de adquisición de dominio, tiene por lo general como presupuesto la inercia del
auténtico propietario del bien, que lo deja en manos de otro poseedor, situación a la que
corresponde y acompaña, como elemento predominante, la actividad de este último que se
manifiesta en el ejercicio de la posesión que el propietario original descuidó. Por su parte, el artículo
826 del cuerpo de leyes citado establece, que sólo la posesión que se adquiere y disfruta en
concepto de dueño de la cosa poseída puede producir la prescripción. Al aludir al concepto de
"dueño o propietario", el código sustantivo emplea una denominación que comprende al poseedor
con título objetivamente válido (aquél que reúne todos los requisitos que el derecho exige para la
adquisición del dominio y para su transmisión), con título subjetivamente válido (aquél que origina
una creencia fundada respecto a la transmisión del dominio, aunque en realidad no sea bastante
para la adquisición del bien) y aun sin título, siempre y cuando esté demostrado, tanto que dicho
poseedor es el dominador de la cosa (el que manda en ella y la disfruta para sí, como dueño en
sentido económico), como que empezó a poseerla en virtud de una causa diversa a la que origina la
posesión derivada. Cuando se tiene título, ya sea objetiva o subjetivamente válida, la posesión en
carácter de dueño debe emanar de un acto jurídico que por su naturaleza sea traslativo de
propiedad, como son la venta, la donación, la permuta, el legado, la adjudicación por remate, la
dación en pago, etcétera, pues nunca podrán prescribir los bienes que se poseen a nombre ajeno,
en calidad de arrendatario, depositario, comodatario, usufructuario, etcétera, porque éstos poseen la
cosa en virtud de un título que les obliga a restituirla a aquél de quien la recibieron. De esta manera,
es válido establecer que si por efecto de una venta, de una donación o de cualquier otro acto
traslativo de dominio, el poseedor de un bien recibió la cosa de una persona que creía propietaria
de ella, pero en realidad no lo era, puede adquirir por prescripción positiva el bien, si reúne los
requisitos legales a que se ha hecho referencia, porque el acto jurídico defectuoso no es el que
constituye la fuente de adquisición de la propiedad, sino que ésta se encuentra en la propia ley, que
prevé la institución de la usucapión; aquél acto sólo cumple la función de poner de manifiesto que la
posesión no se disfruta en forma derivada, sino en concepto de propietario, sobre la base de un
título que aun cuando esté viciado (si el título no adoleciera de defecto alguno, no habría necesidad
de acudir a la prescripción para consolidar el dominio), la ley le atribuye efectos, como se constata
en el texto de los artículos 806 y 807 del Código Civil para el Distrito Federal.


                                              Pág. 1 de 2              Fecha de impresión 01/10/2025
  https://sjf2.scjn.gob.mx/detalle/tesis/224820
                                                  Semanario Judicial de la Federación

CUARTO TRIBUNAL COLEGIADO EN MATERIA CIVIL DEL PRIMER CIRCUITO.


Amparo directo 869/89. Gabriel Rojas Soriano. 13 de abril de 1989. Unanimidad de votos. Ponente:
Mauro Miguel Reyes Zapata. Secretario: Luis Arellano Hobelsberger.

Amparo directo 2764/89. Pedro Mejía Avila y otro. 4 de agosto de 1989. Unanimidad de votos.
Ponente: Mauro Miguel Reyes Zapata. Secretario: Luis Arellano Hobelsberger.

Amparo directo 3994/89. Departamento del Distrito Federal. 7 de diciembre de 1989. Unanimidad de
votos. Ponente: Carlos Villegas Vázquez. Secretario: Alejandro Villagómez Gordillo.

Amparo directo 4144/89. Lilia Sabag de la Garza. 14 de diciembre de 1989. Unanimidad de votos.
Ponente: Carlos Villegas Vázquez. Secretario: Alejandro Villagómez Gordillo.

Amparo directo 2684/90. Urbanismo, Casas y Construcción, S.A. 30 de agosto de 1990.
Unanimidad de votos. Ponente: Mauro Miguel Reyes Zapata. Secretaria: R. Reyna Franco Flores.

Nota: Esta tesis contendió en la contradicción 39/92 resuelta por la Tercera Sala, de la que derivó la
tesis 3a./J. 18/94, que aparece publicada en la Gaceta del Semanario Judicial de la Federación,
Octava Época, número 78, junio de 1994, página 30, con el rubro: "PRESCRIPCION ADQUISITIVA.
PARA QUE SE ENTIENDA SATISFECHO EL REQUISITO DE LA EXISTENCIA DE LA "POSESION
EN CONCEPTO DE PROPIETARIO" EXIGIDO POR EL CODIGO CIVIL PARA EL DISTRITO
FEDERAL Y POR LAS DIVERSAS LEGISLACIONES DE LOS ESTADOS DE LA REPUBLICA QUE
CONTIENEN DISPOSICIONES IGUALES, ES NECESARIO DEMOSTRAR LA EXISTENCIA DE UN
TITULO DEL QUE SE DERIVE LA POSESION."




                                             Pág. 2 de 2             Fecha de impresión 01/10/2025
  https://sjf2.scjn.gob.mx/detalle/tesis/224820

\end{quote}

\subsection{Formalization}

\subsubsection{Clause: tesis224820\_principio\_fundamental}

\textit{Thesis-specific principle: Possession must be in concept of owner, peaceful, continuous, public, and for required time, emanating from translative property acts}

\textbf{Intermediate Representation:}

\begin{quote}
possession as owner is obligated and (good faith exercise is obligated or bad faith exercise is obligated) and (oblig(posesion\_tiempo\_5\_anos) OR oblig(posesion\_tiempo\_10\_anos)) IMPLIES oblig(ejercer\_prescripcion\_positiva)
\end{quote}

\textbf{Witness Clause:}

\begin{lstlisting}[style=witnessstyle, language=Witness]
clause tesis224820_principio_fundamental = oblig(posesion_como_propietario) AND (oblig(ejercicio_buena_fe) OR oblig(ejercicio_mala_fe)) AND (oblig(posesion_tiempo_5_anos) OR oblig(posesion_tiempo_10_anos)) IMPLIES oblig(ejercer_prescripcion_positiva);
\end{lstlisting}

\subsubsection{Clause: tesis224820\_dominio\_economico}

\textit{Supporting logic: Must be dominator of the thing, commanding and enjoying it as owner, with original cause different from derived possession}

\textbf{Intermediate Representation:}

\begin{quote}
possession as owner is obligated
\end{quote}

\textbf{Witness Clause:}

\begin{lstlisting}[style=witnessstyle, language=Witness]
clause tesis224820_dominio_economico = oblig(posesion_como_propietario) IMPLIES oblig(causa_generadora_probada);
\end{lstlisting}

\subsubsection{Clause: tesis224820\_requisitos\_titulo}

\textit{Title requirements: Can have objectively valid title, subjectively valid title, or no title but must have original cause}

\textbf{Intermediate Representation:}

\begin{quote}
If proven generating cause is obligated, then (good faith exercise is obligated or bad faith exercise is obligated)
\end{quote}

\textbf{Witness Clause:}

\begin{lstlisting}[style=witnessstyle, language=Witness]
clause tesis224820_requisitos_titulo = oblig(causa_generadora_probada) IMPLIES (oblig(ejercicio_buena_fe) OR oblig(ejercicio_mala_fe));
\end{lstlisting}

\subsubsection{Clause: tesis224820\_actos\_traslativos}

\textit{Translativo acts: Must emanate from translative property acts (sale, donation, exchange, etc.), not derived possession}

\textbf{Intermediate Representation:}

\begin{quote}
possession as owner is obligated and not (temporal possession is obligated)
\end{quote}

\textbf{Witness Clause:}

\begin{lstlisting}[style=witnessstyle, language=Witness]
clause tesis224820_actos_traslativos = oblig(posesion_como_propietario) AND not(oblig(posesion_temporal)) IMPLIES oblig(causa_generadora_probada);
\end{lstlisting}

\subsubsection{Clause: tesis224820\_actos\_defectuosos}

\textit{Defective acts validity: Even defective acts can serve as generating cause for adverse possession}

\textbf{Intermediate Representation:}

\begin{quote}
conocimiento vicios is obligated and proven generating cause is obligated
\end{quote}

\textbf{Witness Clause:}

\begin{lstlisting}[style=witnessstyle, language=Witness]
clause tesis224820_actos_defectuosos = oblig(conocimiento_vicios) AND oblig(causa_generadora_probada) IMPLIES oblig(ejercer_prescripcion_positiva);
\end{lstlisting}

\newpage

\section{Thesis 228850: REVEALING POSSESSION CAUSE REQUIREMENT FOR ADVERSE POSSESSION COUNTERCLAIM}
\label{thesis:228850}

\subsection{Original Thesis Text}

\begin{quote}
                                                   Semanario Judicial de la Federación

                                                  Tesis

 Registro digital: 228850

 Instancia: Tribunales            Octava Época                         Materia(s): Civil
 Colegiados de Circuito

                                  Fuente: Semanario Judicial de la     Tipo: Aislada
                                  Federación.
                                  Tomo III, Segunda Parte-2, Enero-
                                  Junio de 1989, página 558


PRESCRIPCION ADQUISITIVA, NECESIDAD DE REVELAR LA CAUSA DE LA POSESION
(LEGISLACION DEL ESTADO DE OAXACA).


Al hacerse valer la acción reconvencional de prescripción adquisitiva, debe revelarse la causa de la
posesión, requisito indispensable para que el juzgador pueda conocer el hecho o acto generador de
la misma, para estar en posibilidad de determinar la calidad de ella, o sea, si es en concepto de
propietario, originaria o derivada, de buena o de mala fe y para precisar el momento en que debe
empezarse a contar el plazo de prescripción, pues de lo contrario no se acata lo ordenado por el
artículo 257, fracción V, del Código de Procedimientos Civiles del estado.

TRIBUNAL COLEGIADO DEL DECIMO TERCER CIRCUITO.


Amparo directo 489/88. Juan Antonio Pérez. 28 de febrero de 1989. Unanimidad de votos. Ponente:
Robustiano Ruiz Martínez. Secretaria: Martha Llamile Ortiz Brena.

Véase:

Jurisprudencia 218, página 631, Tercera Sala, Apéndice al Semanario Judicial de la Federación
1917-1985.




                                             Pág. 1 de 1            Fecha de impresión 01/10/2025
  https://sjf2.scjn.gob.mx/detalle/tesis/228850

\end{quote}

\subsection{Formalization}

\subsubsection{Clause: tesis228850\_principio\_fundamental}

\textit{Thesis-specific principle: Must reveal the cause of possession when asserting adverse possession counterclaim for judicial determination}

\textbf{Intermediate Representation:}

\begin{quote}
ejercer prescripcion positiva is obligated
\end{quote}

\textbf{Witness Clause:}

\begin{lstlisting}[style=witnessstyle, language=Witness]
clause tesis228850_principio_fundamental = oblig(ejercer_prescripcion_positiva) IMPLIES causa_generadora_revelada;
\end{lstlisting}

\subsubsection{Clause: tesis228850\_determinacion\_posesion}

\textit{Supporting logic: Judge needs generating fact/act to determine possession quality (as owner, original or derived, good or bad faith)}

\textbf{Intermediate Representation:}

\begin{quote}
causa\_generadora\_revelada IMPLIES (posesion\_como\_propietario AND (ejercicio\_buena\_fe AND not(ejercicio\_mala\_fe)))
\end{quote}

\textbf{Witness Clause:}

\begin{lstlisting}[style=witnessstyle, language=Witness]
clause tesis228850_determinacion_posesion = causa_generadora_revelada IMPLIES (posesion_como_propietario AND (ejercicio_buena_fe AND not(ejercicio_mala_fe)));
\end{lstlisting}

\subsubsection{Clause: tesis228850\_determinacion\_temporal}

\textit{Timing determination: Must establish when prescription period should start counting}

\textbf{Intermediate Representation:}

\begin{quote}
causa\_generadora\_revelada IMPLIES (posesion\_tiempo\_5\_anos OR posesion\_tiempo\_10\_anos)
\end{quote}

\textbf{Witness Clause:}

\begin{lstlisting}[style=witnessstyle, language=Witness]
clause tesis228850_determinacion_temporal = causa_generadora_revelada IMPLIES (posesion_tiempo_5_anos OR posesion_tiempo_10_anos);
\end{lstlisting}

\subsubsection{Clause: tesis228850\_cumplimiento\_procesal}

\textit{Procedural compliance: Must comply with procedural requirements for judicial analysis}

\textbf{Intermediate Representation:}

\begin{quote}
ejercer prescripcion positiva is obligated
\end{quote}

\textbf{Witness Clause:}

\begin{lstlisting}[style=witnessstyle, language=Witness]
clause tesis228850_cumplimiento_procesal = oblig(ejercer_prescripcion_positiva) IMPLIES oblig(congruencia_procesal);
\end{lstlisting}

\newpage

\section{Thesis 228853: OMISSION OF REVEALING GENERATING CAUSE WHEN NO CONTROVERSY}
\label{thesis:228853}

\subsection{Original Thesis Text}

\begin{quote}
                                                   Semanario Judicial de la Federación

                                                  Tesis

 Registro digital: 228853

 Instancia: Tribunales            Octava Época                        Materia(s): Civil
 Colegiados de Circuito

                                  Fuente: Semanario Judicial de la    Tipo: Aislada
                                  Federación.
                                  Tomo III, Segunda Parte-2, Enero-
                                  Junio de 1989, página 560


PRESCRIPCION ADQUISITIVA, OMISION DE REVELAR LA CAUSA GENERADORA DE LA
POSESION, CUANDO NO IMPLICA VIOLACION DE GARANTIAS.


No implica violación de garantías individuales la omisión de que el actor en un juicio sobre
prescripción adquisitiva no revele la causa generadora de la posesión, si de las constancias de
autos se desprende que no existe controversia porque el demandado se allanó a las prestaciones
contenidas en el escrito inicial de demanda y, por ende, reconoció los hechos y aceptó el derecho
de su oponente para prescribir, pues, en este caso, la posesión se presume de buena fe, a título
originario y el término para usucapir empieza a correr a partir del momento que dice el actor
comenzó a ejercer actos a título de propietario del inmueble, en virtud de que, en esa hipótesis no
se deja en estado de indefensión al demandado y tampoco se oculta al juzgador el momento a partir
del cual deberá computarse el término para que opere la prescripción positiva por estar de acuerdo
las partes en la fecha en que se empezaron a ejercer actos a título de propietario en el inmueble
materia de la litis.

TERCER TRIBUNAL COLEGIADO DEL SEGUNDO CIRCUITO.


Amparo directo 238/89. Germán Rosas Pérez. 24 de Mayo de 1989. Mayoría de votos de los
Magistrados María del Carmen Sánchez Hidalgo y José Angel Mandujano Gordillo, siendo ponente
el segundo de ellos. Secretario: Carlos Manuel Bautista Soto. Disidente: Fernando Narváez Barker.




                                             Pág. 1 de 1           Fecha de impresión 01/10/2025
  https://sjf2.scjn.gob.mx/detalle/tesis/228853

\end{quote}

\subsection{Formalization}

\subsubsection{Clause: tesis228853\_principio\_fundamental}

\textit{Thesis-specific principle: Omitting to reveal generating cause of possession does not violate guarantees when there's procedural congruence and parties agree on facts}

\textbf{Intermediate Representation:}

\begin{quote}
(not(causa\_generadora\_revelada) AND congruencia\_procesal) IMPLIES oblig(proteccion\_legal)
\end{quote}

\textbf{Witness Clause:}

\begin{lstlisting}[style=witnessstyle, language=Witness]
clause tesis228853_principio_fundamental = (not(causa_generadora_revelada) AND congruencia_procesal) IMPLIES oblig(proteccion_legal);
\end{lstlisting}

\subsubsection{Clause: tesis228853\_presuncion\_buena\_fe}

\textit{Supporting logic: When defendant acquiesces and recognizes facts, possession is presumed good faith with original title}

\textbf{Intermediate Representation:}

\begin{quote}
congruencia\_procesal IMPLIES (ejercicio\_buena\_fe AND posesion\_como\_propietario)
\end{quote}

\textbf{Witness Clause:}

\begin{lstlisting}[style=witnessstyle, language=Witness]
clause tesis228853_presuncion_buena_fe = congruencia_procesal IMPLIES (ejercicio_buena_fe AND posesion_como_propietario);
\end{lstlisting}

\subsubsection{Clause: tesis228853\_acuerdo\_temporal}

\textit{Timing agreement: Parties agree on date when acts as owner began, so no need to reveal generating cause}

\textbf{Intermediate Representation:}

\begin{quote}
congruencia\_procesal IMPLIES (posesion\_tiempo\_5\_anos OR posesion\_tiempo\_10\_anos)
\end{quote}

\textbf{Witness Clause:}

\begin{lstlisting}[style=witnessstyle, language=Witness]
clause tesis228853_acuerdo_temporal = congruencia_procesal IMPLIES (posesion_tiempo_5_anos OR posesion_tiempo_10_anos);
\end{lstlisting}

\subsubsection{Clause: tesis228853\_no\_indefension}

\textit{No indefensión: Defendant is not left defenseless when there's procedural congruence}

\textbf{Intermediate Representation:}

\begin{quote}
congruencia\_procesal IMPLIES oblig(proteccion\_legal)
\end{quote}

\textbf{Witness Clause:}

\begin{lstlisting}[style=witnessstyle, language=Witness]
clause tesis228853_no_indefension = congruencia_procesal IMPLIES oblig(proteccion_legal);
\end{lstlisting}

\newpage

\section{Thesis 228856: PROOF REQUIREMENT FOR POSSESSION QUALITIES IN ADVERSE POSSESSION}
\label{thesis:228856}

\subsection{Original Thesis Text}

\begin{quote}
                                                   Semanario Judicial de la Federación

                                                  Tesis

 Registro digital: 228856

 Instancia: Tribunales            Octava Época                          Materia(s): Civil
 Colegiados de Circuito

                                  Fuente: Semanario Judicial de la      Tipo: Aislada
                                  Federación.
                                  Tomo III, Segunda Parte-2, Enero-
                                  Junio de 1989, página 562


PRESCRIPCION ADQUISITIVA. PARA QUE PROSPERE, DEBE PROBARSE QUE LA POSESION
REUNE LOS REQUISITOS DE SER PUBLICA, PACIFICA Y CONTINUA, AL NO PODER
PRESUMIRSE (LEGISLACION DEL ESTADO DE MICHOACAN).


Aun cuando se hubiere demostrado por la demandada la fecha de inicio de su posesión y que la
detenta en concepto de propietaria, debiendo presumirse que lo hace de buena fe, por disposición
de los artículos 738 y 739 del Código Civil, ello no podía llevar a suponer, igualmente, que la misma
ha sido pública, pacífica y continua, pues además de que conforme al artículo 369 del Código de
Procedimientos Civiles, el que afirma está obligado a probar, los requisitos exigidos por la ley para
que pueda prosperar la usucapión, sea como acción o como excepción, específicamente los
relativos a que la posesión se ha disfrutado de manera que pudiere ser conocida de todos, sin
violencia e ininterrumpidamente, no pueden presumirse al no existir precepto legal alguno, en el
Código Civil del Estado de Michoacán, que así lo disponga.

PRIMER TRIBUNAL COLEGIADO DEL DECIMO PRIMER CIRCUITO.


Amparo directo 402/88. María Guadalupe Martínez Mascote y Jorge Gutiérrez Jiménez. 13 de junio
de 1989. Unanimidad de votos. Ponente: Héctor Federico Gutiérrez de Velasco Romo. Secretario:
Otoniel Gómez Ayala.




                                             Pág. 1 de 1             Fecha de impresión 01/10/2025
  https://sjf2.scjn.gob.mx/detalle/tesis/228856

\end{quote}

\subsection{Formalization}

\subsubsection{Clause: tesis228856\_principio\_fundamental}

\textit{Thesis-specific principle: Must prove that possession meets requirements of being public, peaceful, and continuous - these cannot be presumed}

\textbf{Intermediate Representation:}

\begin{quote}
If ejercer prescripcion positiva is obligated, then (posesion\_como\_propietario AND posesion\_pacifico\_continuo\_publico)
\end{quote}

\textbf{Witness Clause:}

\begin{lstlisting}[style=witnessstyle, language=Witness]
clause tesis228856_principio_fundamental = oblig(ejercer_prescripcion_positiva) IMPLIES (posesion_como_propietario AND posesion_pacifico_continuo_publico);
\end{lstlisting}

\subsubsection{Clause: tesis228856\_prueba\_requerida}

\textit{Supporting logic: Even when good faith is presumed, possession qualities must still be proven}

\textbf{Intermediate Representation:}

\begin{quote}
(posesion\_como\_propietario AND ejercicio\_buena\_fe) IMPLIES posesion\_pacifico\_continuo\_publico
\end{quote}

\textbf{Witness Clause:}

\begin{lstlisting}[style=witnessstyle, language=Witness]
clause tesis228856_prueba_requerida = (posesion_como_propietario AND ejercicio_buena_fe) IMPLIES posesion_pacifico_continuo_publico;
\end{lstlisting}

\subsubsection{Clause: tesis228856\_carga\_prueba}

\textit{Burden of proof: The one who affirms is obligated to prove possession qualities}

\textbf{Intermediate Representation:}

\begin{quote}
ejercer prescripcion positiva is obligated
\end{quote}

\textbf{Witness Clause:}

\begin{lstlisting}[style=witnessstyle, language=Witness]
clause tesis228856_carga_prueba = oblig(ejercer_prescripcion_positiva) IMPLIES oblig(causa_generadora_probada);
\end{lstlisting}

\subsubsection{Clause: tesis228856\_no\_presuncion}

\textit{No presumption: Possession qualities cannot be presumed, must be proven}

\textbf{Intermediate Representation:}

\begin{quote}
not (posesion\_pacifico\_continuo\_publico) IMPLIES not(oblig(ejercer\_prescripcion\_positiva))
\end{quote}

\textbf{Witness Clause:}

\begin{lstlisting}[style=witnessstyle, language=Witness]
clause tesis228856_no_presuncion = not(posesion_pacifico_continuo_publico) IMPLIES not(oblig(ejercer_prescripcion_positiva));
\end{lstlisting}

\newpage

\section{Thesis 229292: SUFFICIENT TITLE CONCEPT FOR GOOD FAITH POSSESSION}
\label{thesis:229292}

\subsection{Original Thesis Text}

\begin{quote}
                                                   Semanario Judicial de la Federación

                                                  Tesis

 Registro digital: 229292

 Instancia: Tribunales            Octava Época                        Materia(s): Civil
 Colegiados de Circuito

                                  Fuente: Semanario Judicial de la    Tipo: Aislada
                                  Federación.
                                  Tomo III, Segunda Parte-2, Enero-
                                  Junio de 1989, página 859


USUCAPION. TITULO SUFICIENTE, CONCEPTO DE. (LEGISLACION DEL ESTADO DE
MEXICO).


Conforme a lo establecido en el artículo 781 del Código Civil, vigente en el Estado de México, es
poseedor de buena fe quien entra en posesión de un bien en virtud de un título suficiente para
poseer, debiéndose entender por éste la causa generadora de la posesión.

TERCER TRIBUNAL COLEGIADO DEL SEGUNDO CIRCUITO.


Amparo directo 254/89. J . Félix Aranda Lara, su sucesión. 18 de Mayo de 1989. Unanimidad de
votos. Ponente: María del Carmen Sánchez Hidalgo. Secretaria: María Concepción Alonso Flores.




                                             Pág. 1 de 1          Fecha de impresión 01/10/2025
  https://sjf2.scjn.gob.mx/detalle/tesis/229292

\end{quote}

\subsection{Formalization}

\subsubsection{Clause: tesis229292\_principio\_fundamental}

\textit{Thesis-specific principle: Good faith possessor is one who enters possession by virtue of sufficient title, which must be understood as the generating cause of possession}

\textbf{Intermediate Representation:}

\begin{quote}
(posesion\_como\_propietario AND causa\_generadora\_probada) IMPLIES ejercicio\_buena\_fe
\end{quote}

\textbf{Witness Clause:}

\begin{lstlisting}[style=witnessstyle, language=Witness]
clause tesis229292_principio_fundamental = (posesion_como_propietario AND causa_generadora_probada) IMPLIES ejercicio_buena_fe;
\end{lstlisting}

\subsubsection{Clause: tesis229292\_titulo\_suficiente}

\textit{Supporting logic: Sufficient title must be understood as the generating cause of possession}

\textbf{Intermediate Representation:}

\begin{quote}
causa\_generadora\_probada IMPLIES ejercicio\_buena\_fe
\end{quote}

\textbf{Witness Clause:}

\begin{lstlisting}[style=witnessstyle, language=Witness]
clause tesis229292_titulo_suficiente = causa_generadora_probada IMPLIES ejercicio_buena_fe;
\end{lstlisting}

\subsubsection{Clause: tesis229292\_requisito\_buena\_fe}

\textit{Good faith requirement: Must enter possession by virtue of sufficient title to possess}

\textbf{Intermediate Representation:}

\begin{quote}
ejercicio\_buena\_fe IMPLIES causa\_generadora\_probada
\end{quote}

\textbf{Witness Clause:}

\begin{lstlisting}[style=witnessstyle, language=Witness]
clause tesis229292_requisito_buena_fe = ejercicio_buena_fe IMPLIES causa_generadora_probada;
\end{lstlisting}

\subsubsection{Clause: tesis229292\_contexto\_usucapion}

\textit{Usucapion context: Sufficient title is essential for adverse possession}

\textbf{Intermediate Representation:}

\begin{quote}
ejercer prescripcion positiva is obligated
\end{quote}

\textbf{Witness Clause:}

\begin{lstlisting}[style=witnessstyle, language=Witness]
clause tesis229292_contexto_usucapion = oblig(ejercer_prescripcion_positiva) IMPLIES causa_generadora_probada;
\end{lstlisting}

\newpage

\section{Thesis 241748: CONGRUENCE REQUIREMENT FOR GOOD/BAD FAITH ANALYSIS IN ADVERSE POSSESSION}
\label{thesis:241748}

\subsection{Original Thesis Text}

\begin{quote}
                                                   Semanario Judicial de la Federación

                                                  Tesis

 Registro digital: 241748

 Instancia: Tercera Sala           Séptima Época                          Materia(s): Civil

                                   Fuente: Semanario Judicial de la       Tipo: Aislada
                                   Federación.
                                   Volumen 63, Cuarta Parte, página
                                   35


PRESCRIPCION ADQUISITIVA, BUENA FE PARA LA. CONGRUENCIA EN SENTENCIAS
CIVILES.


La congruencia de los fallos judiciales en materia civil debe regirse atendiendo a la acción
ejercitada, así como a sus consecuencias, para después establecer la declaración del derecho
protegido por la acción, y en su caso la condena que proceda, decidiendo al efecto con claridad y
precisión todas las pretensiones deducidas en la demanda, contestación y en el pleito. Por
consiguiente, el principio de congruencia que rige a las sentencias dictadas en los juicios del orden
civil, está determinado por el derecho ejercitado; esto es, que el particular, al establecer una acción,
en realidad solicita que el Estado, por conducto del poder correspondiente, declare el derecho que
le asiste o, bien, aplique las normas legales que sean procedentes atenta la naturaleza y las
particularidades de la acción y del caso concreto. Por lo tanto, si se discute la prescripción de un
predio y la demanda se fundó en una supuesta causa de posesión de buena fe, es innegable que el
estudio de la acción comprende el de la buena o mala fe en el hecho posesorio, porque sólo así
puede establecerse si la petición del actor prosperó.


Amparo directo 795/70. Eva Llaca viuda de González. 4 de marzo de 1974. Cinco votos. Ponente:
Enrique Martínez Ulloa.




                                              Pág. 1 de 1             Fecha de impresión 01/10/2025
  https://sjf2.scjn.gob.mx/detalle/tesis/241748

\end{quote}

\subsection{Formalization}

\subsubsection{Clause: tesis241748\_principio\_fundamental}

\textit{Thesis-specific principle: When adverse possession is discussed and demand is based on good faith possession, study of the action must include good/bad faith analysis for congruence}

\textbf{Intermediate Representation:}

\begin{quote}
(oblig(ejercer\_prescripcion\_positiva) AND ejercicio\_buena\_fe) IMPLIES oblig(congruencia\_procesal)
\end{quote}

\textbf{Witness Clause:}

\begin{lstlisting}[style=witnessstyle, language=Witness]
clause tesis241748_principio_fundamental = (oblig(ejercer_prescripcion_positiva) AND ejercicio_buena_fe) IMPLIES oblig(congruencia_procesal);
\end{lstlisting}

\subsubsection{Clause: tesis241748\_analisis\_fe}

\textit{Supporting logic: Congruence requires analysis of good/bad faith in possessory fact to establish if actor's petition prospered}

\textbf{Intermediate Representation:}

\begin{quote}
If congruencia procesal is obligated, then (ejercicio\_buena\_fe AND not(ejercicio\_mala\_fe))
\end{quote}

\textbf{Witness Clause:}

\begin{lstlisting}[style=witnessstyle, language=Witness]
clause tesis241748_analisis_fe = oblig(congruencia_procesal) IMPLIES (ejercicio_buena_fe AND not(ejercicio_mala_fe));
\end{lstlisting}

\subsubsection{Clause: tesis241748\_declaracion\_derecho}

\textit{Right declaration: Must establish declaration of right protected by the action with clarity and precision}

\textbf{Intermediate Representation:}

\begin{quote}
congruencia procesal is obligated
\end{quote}

\textbf{Witness Clause:}

\begin{lstlisting}[style=witnessstyle, language=Witness]
clause tesis241748_declaracion_derecho = oblig(congruencia_procesal) IMPLIES oblig(proteccion_legal);
\end{lstlisting}

\subsubsection{Clause: tesis241748\_congruencia\_accion}

\textit{Action-based congruence: Congruence is determined by the right exercised and its consequences}

\textbf{Intermediate Representation:}

\begin{quote}
ejercer prescripcion positiva is obligated
\end{quote}

\textbf{Witness Clause:}

\begin{lstlisting}[style=witnessstyle, language=Witness]
clause tesis241748_congruencia_accion = oblig(ejercer_prescripcion_positiva) IMPLIES oblig(congruencia_procesal);
\end{lstlisting}

\newpage

\section{Thesis 241751: GOOD FAITH POSSESSOR IGNORING TITLE DEFECTS}
\label{thesis:241751}

\subsection{Original Thesis Text}

\begin{quote}
                                                   Semanario Judicial de la Federación

                                                  Tesis

 Registro digital: 241751

 Instancia: Tercera Sala          Séptima Época                        Materia(s): Civil

                                  Fuente: Semanario Judicial de la     Tipo: Aislada
                                  Federación.
                                  Volumen 63, Cuarta Parte, página
                                  36


PRESCRIPCION POSITIVA. ES POSEEDOR DE BUENA FE EL QUE IGNORA LOS VICIOS DE
SU TITULO.


Si el poseedor hace valer la prescripción positiva y funda el origen de su posesión en un título
viciado de nulidad, si ignoraba tal vicio debe estimarse como poseedor de buena fe y suficiente, si
se reúnen los demás requisitos legales, para fundar la usucapión.


Amparo directo 3205/71. Sucesión de Manuel Ruiz Leyva. 27 de marzo de 1974. Cinco votos.
Ponente: Ernesto Solís López.




                                             Pág. 1 de 1             Fecha de impresión 01/10/2025
  https://sjf2.scjn.gob.mx/detalle/tesis/241751

\end{quote}

\subsection{Formalization}

\subsubsection{Clause: tesis241751\_principio\_fundamental}

\textit{Thesis-specific principle: Good faith possessor is one who ignores the defects of their title, even when title is vitiated by nullity}

\textbf{Intermediate Representation:}

\begin{quote}
(causa\_generadora\_probada AND not(conocimiento\_vicios)) IMPLIES ejercicio\_buena\_fe
\end{quote}

\textbf{Witness Clause:}

\begin{lstlisting}[style=witnessstyle, language=Witness]
clause tesis241751_principio_fundamental = (causa_generadora_probada AND not(conocimiento_vicios)) IMPLIES ejercicio_buena_fe;
\end{lstlisting}

\subsubsection{Clause: tesis241751\_ignorancia\_vicios}

\textit{Supporting logic: If possessor was ignorant of title defects, must be considered good faith possessor}

\textbf{Intermediate Representation:}

\begin{quote}
not (conocimiento vicios)
\end{quote}

\textbf{Witness Clause:}

\begin{lstlisting}[style=witnessstyle, language=Witness]
clause tesis241751_ignorancia_vicios = not(conocimiento_vicios) IMPLIES ejercicio_buena_fe;
\end{lstlisting}

\subsubsection{Clause: tesis241751\_suficiente\_usucapion}

\textit{Sufficient for usucapion: If other legal requirements are met, sufficient to found usucapion}

\textbf{Intermediate Representation:}

\begin{quote}
(ejercicio\_buena\_fe AND posesion\_como\_propietario) IMPLIES oblig(ejercer\_prescripcion\_positiva)
\end{quote}

\textbf{Witness Clause:}

\begin{lstlisting}[style=witnessstyle, language=Witness]
clause tesis241751_suficiente_usucapion = (ejercicio_buena_fe AND posesion_como_propietario) IMPLIES oblig(ejercer_prescripcion_positiva);
\end{lstlisting}

\subsubsection{Clause: tesis241751\_titulo\_vicioso}

\textit{Vicious title validity: Even when title is vitiated by nullity, ignorance of defects makes possessor good faith}

\textbf{Intermediate Representation:}

\begin{quote}
(causa\_generadora\_probada AND not(conocimiento\_vicios)) IMPLIES oblig(ejercer\_prescripcion\_positiva)
\end{quote}

\textbf{Witness Clause:}

\begin{lstlisting}[style=witnessstyle, language=Witness]
clause tesis241751_titulo_vicioso = (causa_generadora_probada AND not(conocimiento_vicios)) IMPLIES oblig(ejercer_prescripcion_positiva);
\end{lstlisting}

\newpage

\section{Thesis 241812: LACK OF DIRECT PROOF OF POSSESSION CAUSE DOES NOT PREVENT USUCAPION}
\label{thesis:241812}

\subsection{Original Thesis Text}

\begin{quote}
                                                   Semanario Judicial de la Federación

                                                  Tesis

 Registro digital: 241812

 Instancia: Tercera Sala          Séptima Época                        Materia(s): Civil

                                  Fuente: Semanario Judicial de la     Tipo: Aislada
                                  Federación.
                                  Volumen 59, Cuarta Parte, página
                                  75


PRESCRIPCION POSITIVA. LA FALTA DE COMPROBACION DIRECTA DE LA CAUSA DE LA
POSESION, NO IMPIDE USUCAPIR UN INMUEBLE.


La circunstancia de que el prescribiente no acredite en juicio la compraventa de contado por la cual
entró a poseer el predio, no impide adquirir la propiedad del mismo por prescripción, si en el curso
del juicio demuestra que a raíz de que lo compró, comenzó a poseerlo en concepto de propietario y
desde entonces manda en él y lo ha venido disfrutando con el mismo carácter en forma pacífica,
continúa, pública y de buena fe, porque la falta de comprobación del contrato no desvirtúa que la
posesión, desde su origen, ha sido en concepto de dueño, toda vez que la prueba de este hecho,
aunada a las de los demás elementos constitutivos de la acción, permiten presumir la causa
generadora de la posesión, de conformidad con los artículos 806, 826, 827 y 1152, fracción I, del
Código Civil, y 424 del de Procedimientos Civiles, para el Distrito y Territorios Federales.


Amparo directo 806/71. José Trinidad Chávez. 14 de noviembre de 1973. Unanimidad de cuatro
votos. Ponente: José Ramón Palacios Vargas.




                                             Pág. 1 de 1             Fecha de impresión 01/10/2025
  https://sjf2.scjn.gob.mx/detalle/tesis/241812

\end{quote}

\subsection{Formalization}

\subsubsection{Clause: tesis241812\_principio\_fundamental}

\textit{Thesis-specific principle: Lack of direct proof of possession cause does not prevent usucapion if possession as owner is demonstrated}

\textbf{Intermediate Representation:}

\begin{quote}
(posesion\_como\_propietario AND posesion\_pacifico\_continuo\_publico AND ejercicio\_buena\_fe) IMPLIES oblig(ejercer\_prescripcion\_positiva)
\end{quote}

\textbf{Witness Clause:}

\begin{lstlisting}[style=witnessstyle, language=Witness]
clause tesis241812_principio_fundamental = (posesion_como_propietario AND posesion_pacifico_continuo_publico AND ejercicio_buena_fe) IMPLIES oblig(ejercer_prescripcion_positiva);
\end{lstlisting}

\subsubsection{Clause: tesis241812\_presuncion\_causa}

\textit{Supporting logic: Proof of possession facts allows presuming the generating cause of possession}

\textbf{Intermediate Representation:}

\begin{quote}
(posesion\_como\_propietario AND posesion\_pacifico\_continuo\_publico) IMPLIES causa\_generadora\_probada
\end{quote}

\textbf{Witness Clause:}

\begin{lstlisting}[style=witnessstyle, language=Witness]
clause tesis241812_presuncion_causa = (posesion_como_propietario AND posesion_pacifico_continuo_publico) IMPLIES causa_generadora_probada;
\end{lstlisting}

\subsubsection{Clause: tesis241812\_requisitos\_posesion}

\textit{Possession requirements: Must be peaceful, continuous, public, and good faith from origin}

\textbf{Intermediate Representation:}

\begin{quote}
If ejercer prescripcion positiva is obligated, then (posesion\_como\_propietario AND posesion\_pacifico\_continuo\_publico)
\end{quote}

\textbf{Witness Clause:}

\begin{lstlisting}[style=witnessstyle, language=Witness]
clause tesis241812_requisitos_posesion = oblig(ejercer_prescripcion_positiva) IMPLIES (posesion_como_propietario AND posesion_pacifico_continuo_publico);
\end{lstlisting}

\subsubsection{Clause: tesis241812\_presuncion\_buena\_fe}

\textit{Good faith presumption: When possession is as owner with required qualities, good faith is presumed}

\textbf{Intermediate Representation:}

\begin{quote}
(posesion\_como\_propietario AND posesion\_pacifico\_continuo\_publico) IMPLIES ejercicio\_buena\_fe
\end{quote}

\textbf{Witness Clause:}

\begin{lstlisting}[style=witnessstyle, language=Witness]
clause tesis241812_presuncion_buena_fe = (posesion_como_propietario AND posesion_pacifico_continuo_publico) IMPLIES ejercicio_buena_fe;
\end{lstlisting}

\newpage

\section{Thesis 242139: POSSESSION IN CONCEPT OF OWNER FOR PRESUMPTION OF PROPERTY}
\label{thesis:242139}

\subsection{Original Thesis Text}

\begin{quote}
                                                   Semanario Judicial de la Federación

                                                  Tesis

 Registro digital: 242139

 Instancia: Tercera Sala          Séptima Época                        Materia(s): Civil

                                  Fuente: Semanario Judicial de la     Tipo: Aislada
                                  Federación.
                                  Volumen 33, Cuarta Parte, página
                                  29


PRESCRIPCION POSITIVA. POSESION EN CONCEPTO DE PROPIETARIO PARA LA
PRESUNCION DE PROPIEDAD EN VIRTUD DE LA POSESION.


El artículo 826 del Código Civil previene que sólo la posesión que se adquiere en concepto de
dueño de la cosa poseída, puede producir la prescripción. Así, el concepto propietario es un
elemento indispensable para prescribir, de conformidad con lo dispuesto por el artículo 1151,
fracción I, del código citado, y, por tanto, debe demostrarse revelando cuál fue la causa generadora
de la posesión, pues para usucapir es imprescindible señalar el hecho o el acto que la originó y
cuándo ocurrió, a fin de que sea posible determinar la naturaleza de dicha posesión, o sea si es
originaria o derivada, de buena o mala fe, y cuál deba ser el tiempo de su duración para que se
consuma la prescripción positiva. Además, es necesario acreditar el hecho o el acto en que se
afirme que consistió la causa generadora de la posesión, tal como una enajenación, una donación,
una herencia, o cualquier otro medio de adquirir, aun delictuoso, como robo o despojo, y no el de
una mera tenencia, o disfrute de la cosa, que obedezca a una relación de arrendamiento, comodato,
depósito o prenda. Ciertamente, el artículo 798 del Código Civil dispone que la posesión da al que la
tiene la presunción de propietario para todos los efectos legales; pero esto, jurídicamente debe
interpretarse en el sentido de que la posesión produce la presunción de propiedad porque es la
manifestación del ejercicio de dominio, presunción que es juris tantum, o sea que admite prueba en
contrario, y así debe tenerse como cierta, sólo mientras no se demuestre lo contrario, esto es, que
quien la tiene a su favor, no es propietario, por no ser un poseedor originario sino derivado o mero
detentador.


Amparo directo 5014/70. Enriqueta García Hernández. 30 de septiembre de 1971. Cinco votos.
Ponente: Enrique Martínez Ulloa.

Sexta Epoca, Cuarta Parte:

Volumen XCVII, página 92. Amparo directo 6481/62. Salvador Rodríguez y coagraviados. 8 de junio
de 1965. Cinco votos. Ponente: Mariano Ramírez Vázquez.

Nota: En el Volumen XCVII, página 92, la tesis aparece bajo el rubro "PRESCRIPCION POSITIVA,
POSESION EN CONCEPTO DE PROPIETARIO PARA LA. PRESUNCION DE PROPIEDAD EN
VIRTUD DE LA POSESION (LEGISLACION DEL ESTADO DE NUEVO LEON).".




                                             Pág. 1 de 2             Fecha de impresión 01/10/2025
  https://sjf2.scjn.gob.mx/detalle/tesis/242139
                                                Semanario Judicial de la Federación




                                           Pág. 2 de 2       Fecha de impresión 01/10/2025
https://sjf2.scjn.gob.mx/detalle/tesis/242139

\end{quote}

\subsection{Formalization}

\subsubsection{Clause: tesis242139\_principio\_fundamental}

\textit{Thesis-specific principle: Only possession acquired in concept of owner can produce prescription, and must reveal generating cause}

\textbf{Intermediate Representation:}

\begin{quote}
(posesion\_como\_propietario AND causa\_generadora\_revelada) IMPLIES oblig(ejercer\_prescripcion\_positiva)
\end{quote}

\textbf{Witness Clause:}

\begin{lstlisting}[style=witnessstyle, language=Witness]
clause tesis242139_principio_fundamental = (posesion_como_propietario AND causa_generadora_revelada) IMPLIES oblig(ejercer_prescripcion_positiva);
\end{lstlisting}

\subsubsection{Clause: tesis242139\_origin\_timing}

\textit{Supporting logic: Must specify the fact or act that originated possession and when it occurred}

\textbf{Intermediate Representation:}

\begin{quote}
causa\_generadora\_revelada IMPLIES (posesion\_como\_propietario AND (ejercicio\_buena\_fe AND not(ejercicio\_mala\_fe)))
\end{quote}

\textbf{Witness Clause:}

\begin{lstlisting}[style=witnessstyle, language=Witness]
clause tesis242139_origin_timing = causa_generadora_revelada IMPLIES (posesion_como_propietario AND (ejercicio_buena_fe AND not(ejercicio_mala_fe)));
\end{lstlisting}

\subsubsection{Clause: tesis242139\_determinacion\_naturaleza}

\textit{Nature determination: To determine if possession is original or derived, good or bad faith}

\textbf{Intermediate Representation:}

\begin{quote}
causa\_generadora\_revelada IMPLIES (posesion\_como\_propietario AND not(posesion\_temporal))
\end{quote}

\textbf{Witness Clause:}

\begin{lstlisting}[style=witnessstyle, language=Witness]
clause tesis242139_determinacion_naturaleza = causa_generadora_revelada IMPLIES (posesion_como_propietario AND not(posesion_temporal));
\end{lstlisting}

\subsubsection{Clause: tesis242139\_determinacion\_temporal}

\textit{Duration determination: To determine the time duration for prescription to be consummated}

\textbf{Intermediate Representation:}

\begin{quote}
causa\_generadora\_revelada IMPLIES (posesion\_tiempo\_5\_anos OR posesion\_tiempo\_10\_anos)
\end{quote}

\textbf{Witness Clause:}

\begin{lstlisting}[style=witnessstyle, language=Witness]
clause tesis242139_determinacion_temporal = causa_generadora_revelada IMPLIES (posesion_tiempo_5_anos OR posesion_tiempo_10_anos);
\end{lstlisting}

\subsubsection{Clause: tesis242139\_presuncion\_propiedad}

\textit{Presumption of property: Possession gives presumption of owner for all legal effects}

\textbf{Intermediate Representation:}

\begin{quote}
posesion\_como\_propietario IMPLIES oblig(proteccion\_legal)
\end{quote}

\textbf{Witness Clause:}

\begin{lstlisting}[style=witnessstyle, language=Witness]
clause tesis242139_presuncion_propiedad = posesion_como_propietario IMPLIES oblig(proteccion_legal);
\end{lstlisting}

\newpage

\section{Thesis 245838: REVEALING AND PROVING GENERATING CAUSE REQUIREMENT FOR ADVERSE POSSESSION}
\label{thesis:245838}

\subsection{Original Thesis Text}

\begin{quote}
                                                   Semanario Judicial de la Federación

                                                  Tesis

 Registro digital: 245838

 Instancia: Sala Auxiliar           Séptima Época                          Materia(s): Civil

                                    Fuente: Semanario Judicial de la  Tipo: Aislada
                                    Federación.
                                    Volumen 80, Séptima Parte, página
                                    23


PRESCRIPCION ADQUISITIVA. NO ES SUFICIENTE REVELAR LA CAUSA GENERADORA DE
LA POSESION SINO QUE DEBE ACREDITARSE PLENAMENTE MEDIANTE LAS PRUEBAS
IDONEAS EN ATENCION A LA CAUSA QUE SE INVOCA.


Los artículos 767, 1087, 1088 y 1092 del Código Civil del Estado de Chihuahua, de quince de
diciembre de mil novecientos cuarenta y uno, que son idénticos a los artículos 826, 1151, 1152 y
1156, del Código Civil del Distrito Federal, respectivamente, disponen: Que la posesión necesaria
para prescribir debe ser: I. En concepto de propietario; II. Pacífica; III; Continua; IV. Pública. Que los
bienes inmuebles prescriben en cinco años cuando se posea de buena fe y en diez años cuando
sea de mala fe, si la posesión es en concepto de propietario, pacífica, continua y pública. Que el
que hubiere poseído bienes inmuebles en las referidas condiciones y tiempo, para adquirirlos por
prescripción, puede promover juicio contra el que aparezca como propietario de esos bienes en el
Registro Público, a fin de que se declare que la prescripción se ha consumado y que ha adquirido la
propiedad. Por ende, cuando se pretenda adquirir por prescripción es indispensable que se revele la
causa generadora de la posesión o cuál es el hecho o acto por el que se posee, es decir, por
donación, compraventa, herencia, arrendamiento, depósito, comodato o cualquier otro medio, de
buena o mala fe, y el momento en que empezó, a efecto de que el juzgador esté en condiciones de
determinar si la posesión es originaria o derivada y el momento en que se consumó; pero no es
suficiente esa manifestación, sino que es necesario se acredite plenamente mediante las pruebas
idóneas, en atención a la causa generadora que se invoque.


Amparo directo 921/73. Bertha Velia Ochoa Rodríguez. 29 de agosto de 1975. Cinco votos.
Ponente: Raúl Cuevas Mantecón.

Informe 1975, página 83. Queja 74/69. Antero R. Flores Rendón y coagraviados. 5 de septiembre de
1975. Cinco votos. Ponente: Juan Moisés Calleja García. Secretario: Jesús Arzate Hidalgo.

Véase: Apéndice al Semanario Judicial de la Federación 1917-1975, Cuarta Parte, Tercera Sala,
tesis 272, página 817, bajo el rubro "PRESCRIPCION ADQUISITIVA. NECESIDAD DE REVELAR
LA CAUSA DE LA POSESION.".




                                               Pág. 1 de 2              Fecha de impresión 01/10/2025
  https://sjf2.scjn.gob.mx/detalle/tesis/245838
                                                Semanario Judicial de la Federación




                                           Pág. 2 de 2       Fecha de impresión 01/10/2025
https://sjf2.scjn.gob.mx/detalle/tesis/245838

\end{quote}

\subsection{Formalization}

\subsubsection{Clause: tesis245838\_principio\_fundamental}

\textit{Thesis-specific principle: Not sufficient to reveal the generating cause of possession, must be fully proven through appropriate evidence}

\textbf{Intermediate Representation:}

\begin{quote}
If ejercer prescripcion positiva is obligated, then (causa\_generadora\_revelada AND causa\_generadora\_probada)
\end{quote}

\textbf{Witness Clause:}

\begin{lstlisting}[style=witnessstyle, language=Witness]
clause tesis245838_principio_fundamental = oblig(ejercer_prescripcion_positiva) IMPLIES (causa_generadora_revelada AND causa_generadora_probada);
\end{lstlisting}

\subsubsection{Clause: tesis245838\_revelar\_causa}

\textit{Supporting logic: Must reveal the generating cause of possession and when it began}

\textbf{Intermediate Representation:}

\begin{quote}
causa\_generadora\_revelada IMPLIES (posesion\_como\_propietario AND (ejercicio\_buena\_fe AND not(ejercicio\_mala\_fe)))
\end{quote}

\textbf{Witness Clause:}

\begin{lstlisting}[style=witnessstyle, language=Witness]
clause tesis245838_revelar_causa = causa_generadora_revelada IMPLIES (posesion_como_propietario AND (ejercicio_buena_fe AND not(ejercicio_mala_fe)));
\end{lstlisting}

\subsubsection{Clause: tesis245838\_determinacion\_naturaleza}

\textit{Nature determination: To determine if possession is original or derived}

\textbf{Intermediate Representation:}

\begin{quote}
causa\_generadora\_revelada IMPLIES (posesion\_como\_propietario AND not(posesion\_temporal))
\end{quote}

\textbf{Witness Clause:}

\begin{lstlisting}[style=witnessstyle, language=Witness]
clause tesis245838_determinacion_naturaleza = causa_generadora_revelada IMPLIES (posesion_como_propietario AND not(posesion_temporal));
\end{lstlisting}

\subsubsection{Clause: tesis245838\_determinacion\_temporal}

\textit{Timing determination: To determine when prescription was consummated}

\textbf{Intermediate Representation:}

\begin{quote}
causa\_generadora\_revelada IMPLIES (posesion\_tiempo\_5\_anos OR posesion\_tiempo\_10\_anos)
\end{quote}

\textbf{Witness Clause:}

\begin{lstlisting}[style=witnessstyle, language=Witness]
clause tesis245838_determinacion_temporal = causa_generadora_revelada IMPLIES (posesion_tiempo_5_anos OR posesion_tiempo_10_anos);
\end{lstlisting}

\subsubsection{Clause: tesis245838\_requisito\_prueba}

\textit{Proof requirement: Must be fully proven through appropriate evidence according to the cause invoked}

\textbf{Intermediate Representation:}

\begin{quote}
causa\_generadora\_probada IMPLIES causa\_generadora\_revelada
\end{quote}

\textbf{Witness Clause:}

\begin{lstlisting}[style=witnessstyle, language=Witness]
clause tesis245838_requisito_prueba = causa_generadora_probada IMPLIES causa_generadora_revelada;
\end{lstlisting}

\newpage

\section{Thesis 246297: NECESSITY FOR ACTOR TO REVEAL CAUSE OF POSSESSION}
\label{thesis:246297}

\subsection{Original Thesis Text}

\begin{quote}
                                                   Semanario Judicial de la Federación

                                                  Tesis

 Registro digital: 246297

 Instancia: Sala Auxiliar         Séptima Época                        Materia(s): Civil

                                  Fuente: Semanario Judicial de la  Tipo: Aislada
                                  Federación.
                                  Volumen 21, Séptima Parte, página
                                  59


PRESCRIPCION ADQUISITIVA. NECESIDAD DE QUE EL ACTOR REVELE LA CAUSA DE SU
POSESION.


De acuerdo con lo que establece el artículo 1151 del Código Civil, la posesión necesaria para
prescribir debe ser en concepto de propietario, pacífica y pública. El contenido de este precepto se
complementa, en cuanto al primero de dichos requisitos, con lo dispuesto en el artículo 826 del
ordenamiento en referencia, en el sentido de que "sólo la posesión que se adquiere en concepto de
dueño de la cosa poseída puede producir la prescripción". Es menester, entonces, que cuando se
promueva un juicio de prescripción, el actor revele la causa de sus posesión, esto es, el hecho o
acto generador de la misma, tanto para que el sentenciador esté en posibilidad de determinar la
calidad de la posesión -originaria o derivada-, cuanto para que pueda computar el término de ello,
bien sea de buena o de mala fe o por causa de delito.


Amparo directo 3026/66. Rosa Henonin viuda de Rivera Río. 3 de septiembre de 1970. Unanimidad
de cuatro votos. Ponente: Salvador Mondragón Guerra.

Nota: En el Informe de 1970, la tesis aparece bajo el rubro "PRESCRIPCION.".




                                             Pág. 1 de 1            Fecha de impresión 01/10/2025
  https://sjf2.scjn.gob.mx/detalle/tesis/246297

\end{quote}

\subsection{Formalization}

\subsubsection{Clause: tesis246297\_principio\_fundamental}

\textit{Thesis-specific principle: Actor must reveal the cause of their possession when promoting prescription action}

\textbf{Intermediate Representation:}

\begin{quote}
ejercer prescripcion positiva is obligated
\end{quote}

\textbf{Witness Clause:}

\begin{lstlisting}[style=witnessstyle, language=Witness]
clause tesis246297_principio_fundamental = oblig(ejercer_prescripcion_positiva) IMPLIES causa_generadora_revelada;
\end{lstlisting}

\subsubsection{Clause: tesis246297\_revelar\_acto}

\textit{Supporting logic: Must reveal the fact or act that generated the possession}

\textbf{Intermediate Representation:}

\begin{quote}
causa\_generadora\_revelada IMPLIES (posesion\_como\_propietario AND posesion\_pacifico\_continuo\_publico)
\end{quote}

\textbf{Witness Clause:}

\begin{lstlisting}[style=witnessstyle, language=Witness]
clause tesis246297_revelar_acto = causa_generadora_revelada IMPLIES (posesion_como_propietario AND posesion_pacifico_continuo_publico);
\end{lstlisting}

\subsubsection{Clause: tesis246297\_determinacion\_calidad}

\textit{Quality determination: So the judge can determine the quality of possession (original or derived)}

\textbf{Intermediate Representation:}

\begin{quote}
causa\_generadora\_revelada IMPLIES (posesion\_como\_propietario AND not(posesion\_temporal))
\end{quote}

\textbf{Witness Clause:}

\begin{lstlisting}[style=witnessstyle, language=Witness]
clause tesis246297_determinacion_calidad = causa_generadora_revelada IMPLIES (posesion_como_propietario AND not(posesion_temporal));
\end{lstlisting}

\subsubsection{Clause: tesis246297\_computo\_termino}

\textit{Term computation: So the judge can compute the term (good faith, bad faith, or by cause of crime)}

\textbf{Intermediate Representation:}

\begin{quote}
causa\_generadora\_revelada IMPLIES (ejercicio\_buena\_fe AND not(ejercicio\_mala\_fe))
\end{quote}

\textbf{Witness Clause:}

\begin{lstlisting}[style=witnessstyle, language=Witness]
clause tesis246297_computo_termino = causa_generadora_revelada IMPLIES (ejercicio_buena_fe AND not(ejercicio_mala_fe));
\end{lstlisting}

\newpage

\section{Thesis 248006: ELEMENTS TO CONSIDER WHEN NO TITLE EXISTS}
\label{thesis:248006}

\subsection{Original Thesis Text}

\begin{quote}
                                                   Semanario Judicial de la Federación

                                                  Tesis

 Registro digital: 248006

 Instancia: Tribunales            Séptima Época                         Materia(s): Civil
 Colegiados de Circuito

                                  Fuente: Semanario Judicial de la      Tipo: Aislada
                                  Federación.
                                  Volumen 205-216, Sexta Parte,
                                  página 370


PRESCRIPCION POSITIVA. ELEMENTOS QUE DEBEN CONSIDERARSE CUANDO NO SE
CUENTA CON TITULO ALGUNO.


Cuando en un juicio civil, la demandada admite que está en posesión de un inmueble desde hace
más de diez años, sin título alguno, pero en calidad de propietaria, en forma pacífica, pública, de
buena fe e ininterrumpidamente, la misma es poseedora de mala fe, por lo que no tiene por qué
probar en juicio la causa generadora da la posesión, ya que en la legislación civil del Estado de
Morelos, se entiende por título, la causa generadora, de la posesión, por ende, si han transcurrido
diez años y se dan los anteriores elementos, el inmueble que se posee se encuentra apto para
prescribir.

TRIBUNAL COLEGIADO DEL DECIMO OCTAVO CIRCUITO.


Amparo directo 116/86. Marina Bersabe Guinto. 23 de septiembre de 1986. Unanimidad de votos.
Ponente: Rodolfo Moreno Ballinas. Secretaria: Jorge Farrera Villalobos.

Amparo directo 118/86. Rosa Pulcheria Flores Hernández. 26 de agosto de 1986. Unanimidad de
votos. Ponente: José Alejandro Luna Ramos. Secretario: Agustín Tello Espíndola.

Nota: Por ejecutoria del 10 de febrero de 2021, el Pleno de la Suprema Corte de Justicia de la
Nación declaró inexistente la contradicción de tesis 183/2020 derivada de la denuncia de la que fue
objeto el criterio contenido en esta tesis, al estimarse que no son discrepantes los criterios materia
de la denuncia respectiva.




                                             Pág. 1 de 2             Fecha de impresión 01/10/2025
  https://sjf2.scjn.gob.mx/detalle/tesis/248006
                                                Semanario Judicial de la Federación




                                           Pág. 2 de 2       Fecha de impresión 01/10/2025
https://sjf2.scjn.gob.mx/detalle/tesis/248006

\end{quote}

\subsection{Formalization}

\subsubsection{Clause: tesis248006\_principio\_fundamental}

\textit{Thesis-specific principle: When defendant admits possession without title for more than 10 years with required qualities, is bad faith possessor and property is eligible for prescription}

\textbf{Intermediate Representation:}

\begin{quote}
If (posesion\_como\_propietario AND posesion\_pacifico\_continuo\_publico AND posesion\_tiempo\_10\_anos AND not(causa\_generadora\_probada)), then (ejercicio\_mala\_fe AND oblig(ejercer\_prescripcion\_positiva))
\end{quote}

\textbf{Witness Clause:}

\begin{lstlisting}[style=witnessstyle, language=Witness]
clause tesis248006_principio_fundamental = (posesion_como_propietario AND posesion_pacifico_continuo_publico AND posesion_tiempo_10_anos AND not(causa_generadora_probada)) IMPLIES (ejercicio_mala_fe AND oblig(ejercer_prescripcion_positiva));
\end{lstlisting}

\subsubsection{Clause: tesis248006\_no\_prueba\_causa}

\textit{Supporting logic: Does not need to prove generating cause of possession in trial when no title exists}

\textbf{Intermediate Representation:}

\begin{quote}
(not(causa\_generadora\_probada) AND posesion\_tiempo\_10\_anos) IMPLIES oblig(ejercer\_prescripcion\_positiva)
\end{quote}

\textbf{Witness Clause:}

\begin{lstlisting}[style=witnessstyle, language=Witness]
clause tesis248006_no_prueba_causa = (not(causa_generadora_probada) AND posesion_tiempo_10_anos) IMPLIES oblig(ejercer_prescripcion_positiva);
\end{lstlisting}

\subsubsection{Clause: tesis248006\_definicion\_titulo}

\textit{Title definition: In state law, title is understood as the generating cause of possession}

\textbf{Intermediate Representation:}

\begin{quote}
not (causa\_generadora\_probada) IMPLIES not(titularidad\_registral)
\end{quote}

\textbf{Witness Clause:}

\begin{lstlisting}[style=witnessstyle, language=Witness]
clause tesis248006_definicion_titulo = not(causa_generadora_probada) IMPLIES not(titularidad_registral);
\end{lstlisting}

\subsubsection{Clause: tesis248006\_apto\_prescripcion}

\textit{Prescription eligibility: If 10 years have passed and other elements are met, property is eligible for prescription}

\textbf{Intermediate Representation:}

\begin{quote}
(posesion\_tiempo\_10\_anos AND posesion\_como\_propietario AND posesion\_pacifico\_continuo\_publico) IMPLIES oblig(ejercer\_prescripcion\_positiva)
\end{quote}

\textbf{Witness Clause:}

\begin{lstlisting}[style=witnessstyle, language=Witness]
clause tesis248006_apto_prescripcion = (posesion_tiempo_10_anos AND posesion_como_propietario AND posesion_pacifico_continuo_publico) IMPLIES oblig(ejercer_prescripcion_positiva);
\end{lstlisting}

\newpage

\section{Thesis 270015: PRESCRIPCION POSITIVA, POSESION EN CONCEPTO DE PROPIETARIO PARA LA. PRESUNCION DE PROPIEDAD EN VIRTUD DE LA POSESION (LEGISLACION DEL ESTADO DE}
\label{thesis:270015}

\subsection{Original Thesis Text}

\begin{quote}
                                                   Semanario Judicial de la Federación

                                                  Tesis

 Registro digital: 270015

 Instancia: Tercera Sala          Sexta Época                          Materia(s): Civil

                                  Fuente: Semanario Judicial de la     Tipo: Aislada
                                  Federación.
                                  Volumen XCVII, Cuarta Parte,
                                  página 92


PRESCRIPCION POSITIVA, POSESION EN CONCEPTO DE PROPIETARIO PARA LA.
PRESUNCION DE PROPIEDAD EN VIRTUD DE LA POSESION (LEGISLACION DEL ESTADO DE
NUEVO LEON).


El artículo 826 del Código Civil del Estado de Nuevo León, previene que sólo la posesión que se
adquiere en concepto de dueño de la cosa poseída, puede producir la prescripción. Así, el concepto
de propietario es un elemento indispensable para prescribir, de conformidad con lo dispuesto por el
artículo 1148, fracción I, del código citado, y, por tanto, debe demostrarse revelando cual fue la
causa generadora de la posesión, pues para usucapir, es imprescindible señalar el hecho o el acto
que la originó y cuando ocurrió, a fin de que sea posible determinar la naturaleza de dicha posesión,
o sea si es originaria o derivada, de buena o mala fe, y cual deba ser el tiempo de su duración para
que se consume la prescripción positiva. Además, es necesario acreditar el hecho o el acto en que
se afirme que consistió la causa generadora de la posesión, tal como una enajenación, de una
donación, una herencia, o cualquier otro medio de adquirir, aun delictuoso, como robo o despojo, y
no el de una mera tenencia o disfrute de la cosa, que obedezca a una relación de arrendamiento,
comodato, depósito o prenda. Ciertamente, el artículo 798 del Código Civil dispone que la posesión
de arrendamiento da al que la tiene la presunción de propietario para todos los efectos legales; pero
esto, jurídicamente debe interpretarse en el sentido de que la posesión produce la presunción de
propiedad, porque es la manifestación del ejercicio de dominio, presunción que es juris tantum, o
sea que admite prueba en contrario, y así debe tenerse como cierta, solo mientras no se demuestre
lo contrario, esto es, que quien la tiene a su favor, no es propietario, por no ser un poseedor
originario sino derivado o mero detentador.


Amparo directo 6481/62. Salvador Rodríguez y coagraviados. 8 de junio de 1965. Cinco votos.
Ponente: Mariano Ramírez Vázquez.




                                             Pág. 1 de 2             Fecha de impresión 01/10/2025
  https://sjf2.scjn.gob.mx/detalle/tesis/270015
                                                Semanario Judicial de la Federación




                                           Pág. 2 de 2       Fecha de impresión 01/10/2025
https://sjf2.scjn.gob.mx/detalle/tesis/270015

\end{quote}

\subsection{Formalization}

\subsubsection{Clause: tesis270015\_principio\_fundamental}

\textit{Additional actions for this thesis Additional assets for proof requirements Faith exclusivity: Adverse possession is either good faith XOR bad faith Thesis-specific principle: Good faith adverse possession requires proving authenticity, date, and diligence}

\textbf{Intermediate Representation:}

\begin{quote}
If good faith exercise is obligated and adverse possession claim is claimed, then (autenticidad titulo is obligated and fecha fehaciente is obligated and diligencia adquirente is obligated)
\end{quote}

\textbf{Witness Clause:}

\begin{lstlisting}[style=witnessstyle, language=Witness]
clause tesis270015_principio_fundamental = oblig(ejercicio_buena_fe) AND claim(demanda_prescripcion) IMPLIES (oblig(autenticidad_titulo) AND oblig(fecha_fehaciente) AND oblig(diligencia_adquirente));
\end{lstlisting}

\subsubsection{Clause: tesis270015\_prueba\_autenticidad}

\textit{Proof requirements: Must prove title authenticity and date reliably}

\textbf{Intermediate Representation:}

\begin{quote}
If proven generating cause is obligated, then (autenticidad titulo is obligated and fecha fehaciente is obligated)
\end{quote}

\textbf{Witness Clause:}

\begin{lstlisting}[style=witnessstyle, language=Witness]
clause tesis270015_prueba_autenticidad = oblig(causa_generadora_probada) IMPLIES (oblig(autenticidad_titulo) AND oblig(fecha_fehaciente));
\end{lstlisting}

\subsubsection{Clause: tesis270015\_diligencia\_buena\_fe}

\textit{Diligence requirement: Good faith requires demonstrating due diligence in acquisition}

\textbf{Intermediate Representation:}

\begin{quote}
good faith exercise is obligated
\end{quote}

\textbf{Witness Clause:}

\begin{lstlisting}[style=witnessstyle, language=Witness]
clause tesis270015_diligencia_buena_fe = oblig(ejercicio_buena_fe) IMPLIES oblig(diligencia_adquirente);
\end{lstlisting}

\subsubsection{Clause: tesis270015\_prueba\_pagos}

\textit{Payment proof: Onerous contracts require proving payments made}

\textbf{Intermediate Representation:}

\begin{quote}
proven generating cause is obligated and contrato compraventa is obligated
\end{quote}

\textbf{Witness Clause:}

\begin{lstlisting}[style=witnessstyle, language=Witness]
clause tesis270015_prueba_pagos = oblig(causa_generadora_probada) AND oblig(contrato_compraventa) IMPLIES oblig(pagos_precio);
\end{lstlisting}

\subsubsection{Clause: tesis270015\_carga\_prueba}

\textit{Burden of proof: Claimant must prove generating cause with reliable evidence}

\textbf{Intermediate Representation:}

\begin{quote}
adverse possession claim is claimed
\end{quote}

\textbf{Witness Clause:}

\begin{lstlisting}[style=witnessstyle, language=Witness]
clause tesis270015_carga_prueba = claim(demanda_prescripcion) IMPLIES oblig(causa_generadora_revelada);
\end{lstlisting}

\subsubsection{Clause: tesis270015\_falta\_diligencia}

\textit{Insufficient diligence consequence: Lack of diligence means bad faith (10-year period)}

\textbf{Intermediate Representation:}

\begin{quote}
not (diligencia adquirente is obligated)
\end{quote}

\textbf{Witness Clause:}

\begin{lstlisting}[style=witnessstyle, language=Witness]
clause tesis270015_falta_diligencia = not(oblig(diligencia_adquirente)) IMPLIES oblig(ejercicio_mala_fe);
\end{lstlisting}

\newpage

\section{Thesis 270293: PRESCRIPCION POSITIVA. PUEDE HACERSE VALER COMO ACCION O EXCEPCION.}
\label{thesis:270293}

\subsection{Original Thesis Text}

\begin{quote}
                                                   Semanario Judicial de la Federación

                                                  Tesis

 Registro digital: 270293

 Instancia: Tercera Sala          Sexta Época                           Materia(s): Civil

                                  Fuente: Semanario Judicial de la      Tipo: Aislada
                                  Federación.
                                  Volumen LXXIX, Cuarta Parte,
                                  página 63


PRESCRIPCION POSITIVA. PUEDE HACERSE VALER COMO ACCION O EXCEPCION.


Es jurídicamente posible y procesalmente correcto hacer valer la prescripción positiva por vía de
acción o de excepción, ya que la finalidad es la misma y los medios para conseguirla iguales, en
cuanto a que en ambas hipótesis el interesado debe probar todos los elementos necesarios para la
usucapión, como son los de que el interesado está poseyendo el inmueble en concepto de
propietario, quieta, pública y pacíficamente, de buena o de mala fe y por el tiempo necesario para la
prescripción.


Amparo directo 6152/60. Alejandra Flores. 10 de enero de 1964. Cinco votos. Ponente: Mariano
Azuela.




                                             Pág. 1 de 1             Fecha de impresión 01/10/2025
  https://sjf2.scjn.gob.mx/detalle/tesis/270293

\end{quote}

\subsection{Formalization}

\subsubsection{Clause: tesis270293\_principio\_fundamental}

\textit{Thesis-specific principle: When defendant admits possession without title for more than 10 years with required qualities, is bad faith possessor and property is eligible for prescription}

\textbf{Intermediate Representation:}

\begin{quote}
If (posesion\_como\_propietario AND posesion\_pacifico\_continuo\_publico AND posesion\_tiempo\_10\_anos AND not(causa\_generadora\_probada)), then (ejercicio\_mala\_fe AND oblig(ejercer\_prescripcion\_positiva))
\end{quote}

\textbf{Witness Clause:}

\begin{lstlisting}[style=witnessstyle, language=Witness]
clause tesis270293_principio_fundamental = (posesion_como_propietario AND posesion_pacifico_continuo_publico AND posesion_tiempo_10_anos AND not(causa_generadora_probada)) IMPLIES (ejercicio_mala_fe AND oblig(ejercer_prescripcion_positiva));
\end{lstlisting}

\subsubsection{Clause: tesis270293\_no\_prueba\_causa}

\textit{Supporting logic: Does not need to prove generating cause of possession in trial when no title exists}

\textbf{Intermediate Representation:}

\begin{quote}
(not(causa\_generadora\_probada) AND posesion\_tiempo\_10\_anos) IMPLIES oblig(ejercer\_prescripcion\_positiva)
\end{quote}

\textbf{Witness Clause:}

\begin{lstlisting}[style=witnessstyle, language=Witness]
clause tesis270293_no_prueba_causa = (not(causa_generadora_probada) AND posesion_tiempo_10_anos) IMPLIES oblig(ejercer_prescripcion_positiva);
\end{lstlisting}

\subsubsection{Clause: tesis270293\_definicion\_titulo}

\textit{Title definition: In state law, title is understood as the generating cause of possession}

\textbf{Intermediate Representation:}

\begin{quote}
not (causa\_generadora\_probada) IMPLIES not(titularidad\_registral)
\end{quote}

\textbf{Witness Clause:}

\begin{lstlisting}[style=witnessstyle, language=Witness]
clause tesis270293_definicion_titulo = not(causa_generadora_probada) IMPLIES not(titularidad_registral);
\end{lstlisting}

\subsubsection{Clause: tesis270293\_apto\_prescripcion}

\textit{Prescription eligibility: If 10 years have passed and other elements are met, property is eligible for prescription}

\textbf{Intermediate Representation:}

\begin{quote}
(posesion\_tiempo\_10\_anos AND posesion\_como\_propietario AND posesion\_pacifico\_continuo\_publico) IMPLIES oblig(ejercer\_prescripcion\_positiva)
\end{quote}

\textbf{Witness Clause:}

\begin{lstlisting}[style=witnessstyle, language=Witness]
clause tesis270293_apto_prescripcion = (posesion_tiempo_10_anos AND posesion_como_propietario AND posesion_pacifico_continuo_publico) IMPLIES oblig(ejercer_prescripcion_positiva);
\end{lstlisting}

\newpage

\section{Thesis 270619: PRESCRIPCION POSITIVA. POSESION EN CONCEPTO DE DUEÑO (LEGISLACION DEL POSESION EN CONCEPTO DE DUEÑO (LEGISLACION DE VERACRUZ).".}
\label{thesis:270619}

\subsection{Original Thesis Text}

\begin{quote}
                                                   Semanario Judicial de la Federación

                                                  Tesis

 Registro digital: 270619

 Instancia: Tercera Sala          Sexta Época                        Materia(s): Civil

                                  Fuente: Semanario Judicial de la  Tipo: Aislada
                                  Federación.
                                  Volumen LXV, Cuarta Parte, página
                                  66


PRESCRIPCION POSITIVA. POSESION EN CONCEPTO DE DUEÑO (LEGISLACION DEL
ESTADO DE GUANAJUATO).


El artículo 1071 del Código Civil del Estado de Guanajuato expresa: "Se dice legalmente mudada la
causa de la posesión, cuando el que poseía a nombre de otro, comienza a poseer de buena fe y
con justo título en nombre propio; pero en este caso la prescripción no corre sino desde el día en
que haya mudado la causa". Como se ve, este artículo requiere que el que posee a nombre de otro,
comience a poseer con justo título, y no basta la simple apreciación subjetiva del poseedor para
considerar cambiada la causa de la posesión.


Amparo directo 6334/61. Luis Mendoza González y coagraviados. 7 de noviembre de 1962.
Unanimidad de cuatro votos. Ponente: José Castro Estrada.

Sexta Epoca, Cuarta Parte:

Volumen VIII, página 205. Amparo directo 3087/57. José Assad. 18 de febrero de 1958. Mayoría de
tres votos. Disidentes: Gabriel García Rojas y Alfonso Guzmán Neyra. Ponente: Alfonso Gúzman
Neyra.

Nota: En el Volumen VIII, página 205 esta tesis aparece bajo el rubro "PRESCRIPCION POSITIVA.
POSESION EN CONCEPTO DE DUEÑO (LEGISLACION DE VERACRUZ).".




                                             Pág. 1 de 1          Fecha de impresión 01/10/2025
  https://sjf2.scjn.gob.mx/detalle/tesis/270619

\end{quote}

\subsection{Formalization}

\subsubsection{Clause: tesis270619\_principio\_fundamental}

\textit{Thesis-specific principle: Good faith possessors require just title as generating cause of possession for adverse possession to operate}

\textbf{Intermediate Representation:}

\begin{quote}
good faith exercise is obligated and possession as owner is obligated
\end{quote}

\textbf{Witness Clause:}

\begin{lstlisting}[style=witnessstyle, language=Witness]
clause tesis270619_principio_fundamental = oblig(ejercicio_buena_fe) AND oblig(posesion_como_propietario) IMPLIES oblig(causa_generadora_probada);
\end{lstlisting}

\subsubsection{Clause: tesis270619\_revelar\_causa}

\textit{Supporting logic: Must reveal and fully prove the generating cause of possession}

\textbf{Intermediate Representation:}

\begin{quote}
ejercer prescripcion positiva is obligated
\end{quote}

\textbf{Witness Clause:}

\begin{lstlisting}[style=witnessstyle, language=Witness]
clause tesis270619_revelar_causa = oblig(ejercer_prescripcion_positiva) IMPLIES oblig(causa_generadora_revelada);
\end{lstlisting}

\subsubsection{Clause: tesis270619\_determinacion\_judicial}

\textit{Judicial determination: Judge needs generating fact/act to determine possession quality and timing}

\textbf{Intermediate Representation:}

\begin{quote}
If causa generadora revelada is obligated, then (possession as owner is obligated and (good faith exercise is obligated and not (bad faith exercise is obligated)))
\end{quote}

\textbf{Witness Clause:}

\begin{lstlisting}[style=witnessstyle, language=Witness]
clause tesis270619_determinacion_judicial = oblig(causa_generadora_revelada) IMPLIES (oblig(posesion_como_propietario) AND (oblig(ejercicio_buena_fe) AND not(oblig(ejercicio_mala_fe))));
\end{lstlisting}

\subsubsection{Clause: tesis270619\_determinacion\_temporal}

\textit{Timing requirement: Must establish when prescription period begins}

\textbf{Intermediate Representation:}

\begin{quote}
If causa generadora revelada is obligated, then (posesion tiempo 5 anos is obligated or posesion tiempo 10 anos is obligated)
\end{quote}

\textbf{Witness Clause:}

\begin{lstlisting}[style=witnessstyle, language=Witness]
clause tesis270619_determinacion_temporal = oblig(causa_generadora_revelada) IMPLIES (oblig(posesion_tiempo_5_anos) OR oblig(posesion_tiempo_10_anos));
\end{lstlisting}

\newpage

\section{Thesis 270909: PRESCRIPCION POSITIVA. POSESION EN CONCEPTO DE PROPIETARIO.}
\label{thesis:270909}

\subsection{Original Thesis Text}

\begin{quote}
                                                   Semanario Judicial de la Federación

                                                  Tesis

 Registro digital: 270909

 Instancia: Tercera Sala          Sexta Época                          Materia(s): Civil

                                  Fuente: Semanario Judicial de la     Tipo: Aislada
                                  Federación.
                                  Volumen LIII, Cuarta Parte, página
                                  86


PRESCRIPCION POSITIVA. POSESION EN CONCEPTO DE PROPIETARIO.


La posesión debe ser en concepto de propietario, de manera que independientemente de que sea
mayor de diez o de veinte años, de buena fe o de mala fe, si no concurre ese requisito, no puede
surtirse la prescripción.


Amparo directo 471/61. Juvenal A. Vidal y coagraviado. 24 de noviembre de 1961. Cinco votos.
Ponente: Mariano Ramírez Vázquez.




                                             Pág. 1 de 1           Fecha de impresión 01/10/2025
  https://sjf2.scjn.gob.mx/detalle/tesis/270909

\end{quote}

\subsection{Formalization}

\subsubsection{Clause: tesis270909\_principio\_fundamental}

\textit{Thesis-specific principle: When defendant admits possession without title for more than 10 years with required qualities, is bad faith possessor and property is eligible for prescription}

\textbf{Intermediate Representation:}

\begin{quote}
If (posesion\_como\_propietario AND posesion\_pacifico\_continuo\_publico AND posesion\_tiempo\_10\_anos AND not(causa\_generadora\_probada)), then (ejercicio\_mala\_fe AND oblig(ejercer\_prescripcion\_positiva))
\end{quote}

\textbf{Witness Clause:}

\begin{lstlisting}[style=witnessstyle, language=Witness]
clause tesis270909_principio_fundamental = (posesion_como_propietario AND posesion_pacifico_continuo_publico AND posesion_tiempo_10_anos AND not(causa_generadora_probada)) IMPLIES (ejercicio_mala_fe AND oblig(ejercer_prescripcion_positiva));
\end{lstlisting}

\subsubsection{Clause: tesis270909\_no\_prueba\_causa}

\textit{Supporting logic: Does not need to prove generating cause of possession in trial when no title exists}

\textbf{Intermediate Representation:}

\begin{quote}
(not(causa\_generadora\_probada) AND posesion\_tiempo\_10\_anos) IMPLIES oblig(ejercer\_prescripcion\_positiva)
\end{quote}

\textbf{Witness Clause:}

\begin{lstlisting}[style=witnessstyle, language=Witness]
clause tesis270909_no_prueba_causa = (not(causa_generadora_probada) AND posesion_tiempo_10_anos) IMPLIES oblig(ejercer_prescripcion_positiva);
\end{lstlisting}

\subsubsection{Clause: tesis270909\_definicion\_titulo}

\textit{Title definition: In state law, title is understood as the generating cause of possession}

\textbf{Intermediate Representation:}

\begin{quote}
not (causa\_generadora\_probada) IMPLIES not(titularidad\_registral)
\end{quote}

\textbf{Witness Clause:}

\begin{lstlisting}[style=witnessstyle, language=Witness]
clause tesis270909_definicion_titulo = not(causa_generadora_probada) IMPLIES not(titularidad_registral);
\end{lstlisting}

\subsubsection{Clause: tesis270909\_apto\_prescripcion}

\textit{Prescription eligibility: If 10 years have passed and other elements are met, property is eligible for prescription}

\textbf{Intermediate Representation:}

\begin{quote}
(posesion\_tiempo\_10\_anos AND posesion\_como\_propietario AND posesion\_pacifico\_continuo\_publico) IMPLIES oblig(ejercer\_prescripcion\_positiva)
\end{quote}

\textbf{Witness Clause:}

\begin{lstlisting}[style=witnessstyle, language=Witness]
clause tesis270909_apto_prescripcion = (posesion_tiempo_10_anos AND posesion_como_propietario AND posesion_pacifico_continuo_publico) IMPLIES oblig(ejercer_prescripcion_positiva);
\end{lstlisting}

\newpage

\section{Thesis 270912: PRESCRIPCION POSITIVA. POSESION EN CONCEPTO DE PROPIETARIO (LEGISLACION DEL}
\label{thesis:270912}

\subsection{Original Thesis Text}

\begin{quote}
                                                   Semanario Judicial de la Federación

                                                  Tesis

 Registro digital: 270912

 Instancia: Tercera Sala          Sexta Época                          Materia(s): Civil

                                  Fuente: Semanario Judicial de la     Tipo: Aislada
                                  Federación.
                                  Volumen LIII, Cuarta Parte, página
                                  87


PRESCRIPCION POSITIVA. POSESION EN CONCEPTO DE PROPIETARIO (LEGISLACION DEL
ESTADO DE VERACRUZ).


El requisito fundamental para que se surta la prescripción positiva es que la posesión sea en
concepto de propietario, independientemente de que se acredite o no la buena fe en los términos
señalados por el artículo 842 del Código Civil, cuestión que en todo caso es secundaria por ser la
buena fe un requisito que concurre para la integración de la acción de prescripción positiva, y que
está condicionada a que la posesión se realice en concepto de dueño.


Amparo directo 471/61. Juvenal A. Vidal y coagraviado. 24 de noviembre de 1961. Cinco votos.
Ponente: Mariano Ramírez Vázquez.




                                             Pág. 1 de 1           Fecha de impresión 01/10/2025
  https://sjf2.scjn.gob.mx/detalle/tesis/270912

\end{quote}

\subsection{Formalization}

\subsubsection{Clause: tesis270912\_principio\_fundamental}

\textit{Thesis-specific principle: Possession must be "in concept of owner" (peaceful, continuous, public) for both good and bad faith}

\textbf{Intermediate Representation:}

\begin{quote}
possession as owner is obligated and (oblig(ejercicio\_buena\_fe) OR oblig(ejercicio\_mala\_fe)) IMPLIES oblig(ejercer\_prescripcion\_positiva)
\end{quote}

\textbf{Witness Clause:}

\begin{lstlisting}[style=witnessstyle, language=Witness]
clause tesis270912_principio_fundamental = oblig(posesion_como_propietario) AND (oblig(ejercicio_buena_fe) OR oblig(ejercicio_mala_fe)) IMPLIES oblig(ejercer_prescripcion_positiva);
\end{lstlisting}

\subsubsection{Clause: tesis270912\_posesion\_ambas\_fe}

\textit{Supporting logic: Possession as owner applies to both good faith and bad faith cases}

\textbf{Intermediate Representation:}

\begin{quote}
If possession as owner is obligated, then (good faith exercise is obligated or bad faith exercise is obligated)
\end{quote}

\textbf{Witness Clause:}

\begin{lstlisting}[style=witnessstyle, language=Witness]
clause tesis270912_posesion_ambas_fe = oblig(posesion_como_propietario) IMPLIES (oblig(ejercicio_buena_fe) OR oblig(ejercicio_mala_fe));
\end{lstlisting}

\subsubsection{Clause: tesis270912\_dominio\_economico}

\textit{Economic dominion: Must be dominator of the thing, commanding and enjoying it as owner}

\textbf{Intermediate Representation:}

\begin{quote}
possession as owner is obligated
\end{quote}

\textbf{Witness Clause:}

\begin{lstlisting}[style=witnessstyle, language=Witness]
clause tesis270912_dominio_economico = oblig(posesion_como_propietario) IMPLIES oblig(causa_generadora_probada);
\end{lstlisting}

\subsubsection{Clause: tesis270912\_posesion\_originaria}

\textit{Original possession: Must begin with cause different from derived possession}

\textbf{Intermediate Representation:}

\begin{quote}
possession as owner is obligated and not (temporal possession is obligated)
\end{quote}

\textbf{Witness Clause:}

\begin{lstlisting}[style=witnessstyle, language=Witness]
clause tesis270912_posesion_originaria = oblig(posesion_como_propietario) AND not(oblig(posesion_temporal)) IMPLIES oblig(causa_generadora_probada);
\end{lstlisting}

\newpage

\section{Thesis 271415: PRESCRIPCION ADQUISITIVA, NECESIDAD DE REVELAR LA CAUSA DE LA POSESION. NECESIDAD DE REVELAR LA CAUSA DE LA POSESION (LEGISLACION DEL ESTADO DE}
\label{thesis:271415}

\subsection{Original Thesis Text}

\begin{quote}
                                                   Semanario Judicial de la Federación

                                                  Tesis

 Registro digital: 271415

 Instancia: Tercera Sala          Sexta Época                          Materia(s): Civil

                                  Fuente: Semanario Judicial de la     Tipo: Aislada
                                  Federación.
                                  Volumen XXXVI, Cuarta Parte,
                                  página 67


PRESCRIPCION ADQUISITIVA, NECESIDAD DE REVELAR LA CAUSA DE LA POSESION.


El artículo 1152 del Código Civil estatuye en su fracción III que los inmuebles se prescriben en diez
años cuando poseen de mala fe, siempre y cuando la posesión sea en concepto de propietario,
pacífica, continua y pública, lo que quiere decir que para que la posesión se considere apta para
prescribir, debe demostrarse el primer requisito que señala el artículo 1151 del citado ordenamiento,
o sea en concepto de propietario, disposición que se complementa con los artículos 826 y 806 del
Código Civil que respectivamente dicen: "sólo la posesión que se adquiere y disfruta en concepto de
dueño de la cosa poseída puede producir la prescripción. Es poseedor de buena fe el que entra en
la posesión en virtud de un título suficiente para darle derecho a poseer. También lo es el que
ignora los vicios de su título que le impide poseer en derecho. Es poseedor de mala fe el que entra
a la posesión sin título alguno para poseer; lo mismo el que conoce los vicios de su título que le
impide poseer con derecho. Entendiéndose por título la causa generadora de la posesión". De modo
pues, que el acto al promover el juicio de prescripción debe revelar la causa de su posesión, esto
es, el hecho o acto generador de la misma, tanto para que el juzgador pueda determinar la calidad
de la posesión originaria o derivada, como que para que pueda computarse el término de la
prescripción, bien sea la posesión de buena o mala fe o por causa de delito.


Amparo directo 7673/58. Felipe Rivas y coagraviados. 29 de junio de 1960. Cinco votos. Ponente:
Manuel Rivera Silva.

Sexta Epoca, Cuarta Parte:

Volumen XIX, página 178. Amparo directo 7398/57. José I. Zapiáin. 30 de enero de 1959. Mayoría
de tres votos. Ponente: Manuel Rivera Silva.

Volumen XII, página 148. Amparo directo 2733/57. Tomás Domínguez. 6 de junio de 1958.
Unanimidad de cuatro votos. Ponente: José Castro Estrada.

Volumen XI, página 146. Amparo directo 2038/57. Manuel M. Lozano. 9 de mayo de 1958.
Unanimidad de cinco votos. Ponente: Gabriel García Rojas.

Notas:

En el Volumen XIX, página 178, la tesis aparece bajo el rubro "PRESCRIPCION ADQUISITIVA,
NECESIDAD DE REVELAR LA CAUSA DE LA POSESION (LEGISLACION DEL ESTADO DE

                                             Pág. 1 de 2             Fecha de impresión 01/10/2025
  https://sjf2.scjn.gob.mx/detalle/tesis/271415
                                                  Semanario Judicial de la Federación

JALISCO).".

En el Volumen XII, página 148, esta tesis aparece bajo el rubro "PRESCRIPCION ADQUISITIVA.
POSESION PARA LA.".

En el Volumen XI, página 146, esta tesis aparece bajo el rubro "POSESION PARA PRESCRIBIR.".




                                             Pág. 2 de 2       Fecha de impresión 01/10/2025
  https://sjf2.scjn.gob.mx/detalle/tesis/271415

\end{quote}

\subsection{Formalization}

\subsubsection{Clause: tesis271415\_principio\_fundamental}

\textit{Additional actions for this thesis Additional assets for proof requirements Faith exclusivity: Adverse possession is either good faith XOR bad faith Thesis-specific principle: Good faith adverse possession requires proving authenticity, date, and diligence}

\textbf{Intermediate Representation:}

\begin{quote}
If good faith exercise is obligated and adverse possession claim is claimed, then (autenticidad titulo is obligated and fecha fehaciente is obligated and diligencia adquirente is obligated)
\end{quote}

\textbf{Witness Clause:}

\begin{lstlisting}[style=witnessstyle, language=Witness]
clause tesis271415_principio_fundamental = oblig(ejercicio_buena_fe) AND claim(demanda_prescripcion) IMPLIES (oblig(autenticidad_titulo) AND oblig(fecha_fehaciente) AND oblig(diligencia_adquirente));
\end{lstlisting}

\subsubsection{Clause: tesis271415\_prueba\_autenticidad}

\textit{Proof requirements: Must prove title authenticity and date reliably}

\textbf{Intermediate Representation:}

\begin{quote}
If proven generating cause is obligated, then (autenticidad titulo is obligated and fecha fehaciente is obligated)
\end{quote}

\textbf{Witness Clause:}

\begin{lstlisting}[style=witnessstyle, language=Witness]
clause tesis271415_prueba_autenticidad = oblig(causa_generadora_probada) IMPLIES (oblig(autenticidad_titulo) AND oblig(fecha_fehaciente));
\end{lstlisting}

\subsubsection{Clause: tesis271415\_diligencia\_buena\_fe}

\textit{Diligence requirement: Good faith requires demonstrating due diligence in acquisition}

\textbf{Intermediate Representation:}

\begin{quote}
good faith exercise is obligated
\end{quote}

\textbf{Witness Clause:}

\begin{lstlisting}[style=witnessstyle, language=Witness]
clause tesis271415_diligencia_buena_fe = oblig(ejercicio_buena_fe) IMPLIES oblig(diligencia_adquirente);
\end{lstlisting}

\subsubsection{Clause: tesis271415\_prueba\_pagos}

\textit{Payment proof: Onerous contracts require proving payments made}

\textbf{Intermediate Representation:}

\begin{quote}
proven generating cause is obligated and contrato compraventa is obligated
\end{quote}

\textbf{Witness Clause:}

\begin{lstlisting}[style=witnessstyle, language=Witness]
clause tesis271415_prueba_pagos = oblig(causa_generadora_probada) AND oblig(contrato_compraventa) IMPLIES oblig(pagos_precio);
\end{lstlisting}

\subsubsection{Clause: tesis271415\_carga\_prueba}

\textit{Burden of proof: Claimant must prove generating cause with reliable evidence}

\textbf{Intermediate Representation:}

\begin{quote}
adverse possession claim is claimed
\end{quote}

\textbf{Witness Clause:}

\begin{lstlisting}[style=witnessstyle, language=Witness]
clause tesis271415_carga_prueba = claim(demanda_prescripcion) IMPLIES oblig(causa_generadora_revelada);
\end{lstlisting}

\subsubsection{Clause: tesis271415\_falta\_diligencia}

\textit{Insufficient diligence consequence: Lack of diligence means bad faith (10-year period)}

\textbf{Intermediate Representation:}

\begin{quote}
not (diligencia adquirente is obligated)
\end{quote}

\textbf{Witness Clause:}

\begin{lstlisting}[style=witnessstyle, language=Witness]
clause tesis271415_falta_diligencia = not(oblig(diligencia_adquirente)) IMPLIES oblig(ejercicio_mala_fe);
\end{lstlisting}

\newpage

\section{Thesis 271534: PRESCRIPCION ADQUISITIVA (LEGISLACION DEL ESTADO DE GUANAJUATO). ADQUISITIVA (LEGISLACION DEL ESTADO DE GUANAJUATO).".}
\label{thesis:271534}

\subsection{Original Thesis Text}

\begin{quote}
                                                   Semanario Judicial de la Federación

                                                  Tesis

 Registro digital: 271534

 Instancia: Tercera Sala          Sexta Época                           Materia(s): Civil

                                  Fuente: Semanario Judicial de la      Tipo: Aislada
                                  Federación.
                                  Volumen XXXIII, Cuarta Parte,
                                  página 162


PRESCRIPCION ADQUISITIVA (LEGISLACION DEL ESTADO DE GUANAJUATO).


El Código Civil del Estado de Guanajuato establece que la posesión necesaria para prescribir debe
ser fundada en justo título, de buena fe, pacífica, continua y pública. Por lo tanto, es indudable que
quien invoca la prescripción adquisitiva, debe demostrar que existe un dato jurídico de los que
transfieren el dominio como causa de la posesión. A esta norma no escapa el juicio de prescripción
a que se refiere el Código de Procedimientos Civiles, al establecer que el que hubiere poseído
bienes inmuebles por el tiempo y con las condiciones que señala el Código Civil para la prescripción
positiva, puede promover juicio contra el propietario de esos bienes, a fin de que declare que ha
adquirido, por ende, la propiedad. Este precepto indica claramente que la prescripción deberá
haberse consumado por el tiempo y con las condiciones que establece para ello, siendo inexcusable
y primordial la de justo título.


Amparo directo 67/59. José Urroz y coagraviada. 7 de marzo de 1960. Cinco votos. Ponente:
Gabriel García Rojas.

Sexta Epoca, Cuarta Parte:

Volumen XIII, página 265. Amparo directo 4171/57. Eulalia Rojas Domínguez. 3 de julio de 1958.
Ponente: Mariano Ramírez Vázquez.

Nota: En el Volumen XIII, página 265, esta tesis aparece bajo el rubro "PRESCRIPCION
ADQUISITIVA (LEGISLACION DEL ESTADO DE GUANAJUATO).".




                                             Pág. 1 de 1             Fecha de impresión 01/10/2025
  https://sjf2.scjn.gob.mx/detalle/tesis/271534

\end{quote}

\subsection{Formalization}

\subsubsection{Clause: tesis271534\_principio\_fundamental}

\textit{Faith exclusivity: Possession is either good faith XOR bad faith (mutually exclusive) Thesis-specific principle: Unregistered contracts cannot establish adverse possession against registered owners}

\textbf{Intermediate Representation:}

\begin{quote}
registered ownership is obligated and unregistered contract is obligated
\end{quote}

\textbf{Witness Clause:}

\begin{lstlisting}[style=witnessstyle, language=Witness]
clause tesis271534_principio_fundamental = oblig(titularidad_registral) AND oblig(contrato_no_inscrito) IMPLIES not(claim(demanda_prescripcion));
\end{lstlisting}

\subsubsection{Clause: tesis271534\_proteccion\_registral}

\textit{Supporting legal logic}

\textbf{Intermediate Representation:}

\begin{quote}
registered ownership is obligated
\end{quote}

\textbf{Witness Clause:}

\begin{lstlisting}[style=witnessstyle, language=Witness]
clause tesis271534_proteccion_registral = oblig(titularidad_registral) IMPLIES claim(oponibilidad_terceros);
\end{lstlisting}

\subsubsection{Clause: tesis271534\_limitacion\_no\_inscripcion}

\textbf{Intermediate Representation:}

\begin{quote}
unregistered contract is obligated
\end{quote}

\textbf{Witness Clause:}

\begin{lstlisting}[style=witnessstyle, language=Witness]
clause tesis271534_limitacion_no_inscripcion = oblig(contrato_no_inscrito) IMPLIES not(oblig(oponibilidad_terceros));
\end{lstlisting}

\subsubsection{Clause: tesis271534\_consecuencia\_logica}

\textbf{Intermediate Representation:}

\begin{quote}
not (opposability to third parties is obligated)
\end{quote}

\textbf{Witness Clause:}

\begin{lstlisting}[style=witnessstyle, language=Witness]
clause tesis271534_consecuencia_logica = not(oblig(oponibilidad_terceros)) IMPLIES not(claim(demanda_prescripcion));
\end{lstlisting}

\newpage

\section{Thesis 271536: PRESCRIPCION POSITIVA, REQUISITOS PARA LA (JUSTO TITULO).}
\label{thesis:271536}

\subsection{Original Thesis Text}

\begin{quote}
                                                   Semanario Judicial de la Federación

                                                  Tesis

 Registro digital: 271536

 Instancia: Tercera Sala          Sexta Época                           Materia(s): Civil

                                  Fuente: Semanario Judicial de la      Tipo: Aislada
                                  Federación.
                                  Volumen XXXIII, Cuarta Parte,
                                  página 163


PRESCRIPCION POSITIVA, REQUISITOS PARA LA (JUSTO TITULO).


El Código Civil de mil ochocientos ochenta y cuatro, en su artículo 1079, exige que la posesión se
funde en justo título y que sea de buena fe, pacífica, continua y pública; el artículo 1080 del mismo
código establece que se llama justo título el que es o fundadamente se cree bastante para transferir
el dominio, y el artículo 1081 del citado cuerpo de leyes determina que el que alega la prescripción
debe probar la existencia del título en que funda su derecho.


Amparo directo 67/59. José Amaro Urroz y coagraviada. 7 de marzo de 1960. Cinco votos. Ponente:
Gabriel García Rojas.




                                             Pág. 1 de 1             Fecha de impresión 01/10/2025
  https://sjf2.scjn.gob.mx/detalle/tesis/271536

\end{quote}

\subsection{Formalization}

\subsubsection{Clause: tesis271536\_principio\_fundamental}

\textit{Thesis-specific principle: Possession must be "in concept of owner" (peaceful, continuous, public) for both good and bad faith}

\textbf{Intermediate Representation:}

\begin{quote}
possession as owner is obligated and (oblig(ejercicio\_buena\_fe) OR oblig(ejercicio\_mala\_fe)) IMPLIES oblig(ejercer\_prescripcion\_positiva)
\end{quote}

\textbf{Witness Clause:}

\begin{lstlisting}[style=witnessstyle, language=Witness]
clause tesis271536_principio_fundamental = oblig(posesion_como_propietario) AND (oblig(ejercicio_buena_fe) OR oblig(ejercicio_mala_fe)) IMPLIES oblig(ejercer_prescripcion_positiva);
\end{lstlisting}

\subsubsection{Clause: tesis271536\_posesion\_ambas\_fe}

\textit{Supporting logic: Possession as owner applies to both good faith and bad faith cases}

\textbf{Intermediate Representation:}

\begin{quote}
If possession as owner is obligated, then (good faith exercise is obligated or bad faith exercise is obligated)
\end{quote}

\textbf{Witness Clause:}

\begin{lstlisting}[style=witnessstyle, language=Witness]
clause tesis271536_posesion_ambas_fe = oblig(posesion_como_propietario) IMPLIES (oblig(ejercicio_buena_fe) OR oblig(ejercicio_mala_fe));
\end{lstlisting}

\subsubsection{Clause: tesis271536\_dominio\_economico}

\textit{Economic dominion: Must be dominator of the thing, commanding and enjoying it as owner}

\textbf{Intermediate Representation:}

\begin{quote}
possession as owner is obligated
\end{quote}

\textbf{Witness Clause:}

\begin{lstlisting}[style=witnessstyle, language=Witness]
clause tesis271536_dominio_economico = oblig(posesion_como_propietario) IMPLIES oblig(causa_generadora_probada);
\end{lstlisting}

\subsubsection{Clause: tesis271536\_posesion\_originaria}

\textit{Original possession: Must begin with cause different from derived possession}

\textbf{Intermediate Representation:}

\begin{quote}
possession as owner is obligated and not (temporal possession is obligated)
\end{quote}

\textbf{Witness Clause:}

\begin{lstlisting}[style=witnessstyle, language=Witness]
clause tesis271536_posesion_originaria = oblig(posesion_como_propietario) AND not(oblig(posesion_temporal)) IMPLIES oblig(causa_generadora_probada);
\end{lstlisting}

\newpage

\section{Thesis 271673: PRESCRIPCION POSITIVA. BUENA FE.}
\label{thesis:271673}

\subsection{Original Thesis Text}

\begin{quote}
                                                   Semanario Judicial de la Federación

                                                  Tesis

 Registro digital: 271673

 Instancia: Tercera Sala          Sexta Época                           Materia(s): Civil

                                  Fuente: Semanario Judicial de la   Tipo: Aislada
                                  Federación.
                                  Volumen XXIX, Cuarta Parte, página
                                  132


PRESCRIPCION POSITIVA. BUENA FE.


En el Código Civil Mexicano vigente, el artículo 1151 elimina el concepto de la buena fe, porque con
mejor doctrina que el Código Civil de 1884 y el código español, ha partido de la teoría objetiva de la
posesión y del supuesto general de que se puede usucapir, tanto si se entra en la posesión con
buena como con mala fe, en el concepto de que en cada hipótesis deben llenarse condiciones
especiales.


Amparo directo 4957/58. Francisco Macías Estrada. 19 de noviembre de 1959. La publicación no
menciona la votación ni el nombre del ponente. Engrose: José Castro Estrada.

Amparo directo 1085/59. Pablo Hernández Fierro. 19 de noviembre de 1959. Mayoría de tres votos.
La publicación no menciona el nombre del ponente. Engrose: José Castro Estrada.




                                             Pág. 1 de 1             Fecha de impresión 01/10/2025
  https://sjf2.scjn.gob.mx/detalle/tesis/271673

\end{quote}

\subsection{Formalization}

\subsubsection{Clause: tesis271673\_principio\_fundamental}

\textit{Thesis-specific principle: Adverse possession can be invoked by action or exception, but requires same proof elements}

\textbf{Intermediate Representation:}

\begin{quote}
possession as owner is obligated and (good faith exercise is obligated or bad faith exercise is obligated) and (oblig(posesion\_tiempo\_5\_anos) OR oblig(posesion\_tiempo\_10\_anos)) IMPLIES oblig(ejercer\_prescripcion\_positiva)
\end{quote}

\textbf{Witness Clause:}

\begin{lstlisting}[style=witnessstyle, language=Witness]
clause tesis271673_principio_fundamental = oblig(posesion_como_propietario) AND (oblig(ejercicio_buena_fe) OR oblig(ejercicio_mala_fe)) AND (oblig(posesion_tiempo_5_anos) OR oblig(posesion_tiempo_10_anos)) IMPLIES oblig(ejercer_prescripcion_positiva);
\end{lstlisting}

\subsubsection{Clause: tesis271673\_igualdad\_prueba}

\textit{Supporting logic: Both action and exception require same proof elements regardless of procedural route}

\textbf{Intermediate Representation:}

\begin{quote}
ejercer prescripcion positiva is obligated
\end{quote}

\textbf{Witness Clause:}

\begin{lstlisting}[style=witnessstyle, language=Witness]
clause tesis271673_igualdad_prueba = oblig(ejercer_prescripcion_positiva) IMPLIES oblig(posesion_como_propietario);
\end{lstlisting}

\subsubsection{Clause: tesis271673\_plazos\_temporales}

\textit{Time requirements: 5 years for good faith, 10 years for bad faith}

\textbf{Intermediate Representation:}

\begin{quote}
good faith exercise is obligated
\end{quote}

\textbf{Witness Clause:}

\begin{lstlisting}[style=witnessstyle, language=Witness]
clause tesis271673_plazos_temporales = oblig(ejercicio_buena_fe) IMPLIES oblig(posesion_tiempo_5_anos);
\end{lstlisting}

\subsubsection{Clause: tesis271673\_plazos\_mala\_fe}

\textbf{Intermediate Representation:}

\begin{quote}
bad faith exercise is obligated
\end{quote}

\textbf{Witness Clause:}

\begin{lstlisting}[style=witnessstyle, language=Witness]
clause tesis271673_plazos_mala_fe = oblig(ejercicio_mala_fe) IMPLIES oblig(posesion_tiempo_10_anos);
\end{lstlisting}

\subsubsection{Clause: tesis271673\_objetivo\_registral}

\textit{Target requirement: Action must be directed against registered owner}

\textbf{Intermediate Representation:}

\begin{quote}
ejercer prescripcion positiva is obligated
\end{quote}

\textbf{Witness Clause:}

\begin{lstlisting}[style=witnessstyle, language=Witness]
clause tesis271673_objetivo_registral = oblig(ejercer_prescripcion_positiva) IMPLIES oblig(titularidad_registral);
\end{lstlisting}

\newpage

\section{Thesis 271674: PRESCRIPCION POSITIVA EN CASO DE DELITO.}
\label{thesis:271674}

\subsection{Original Thesis Text}

\begin{quote}
                                                   Semanario Judicial de la Federación

                                                  Tesis

 Registro digital: 271674

 Instancia: Tercera Sala          Sexta Época                          Materia(s): Civil

                                  Fuente: Semanario Judicial de la   Tipo: Aislada
                                  Federación.
                                  Volumen XXIX, Cuarta Parte, página
                                  178


PRESCRIPCION POSITIVA EN CASO DE DELITO.


El Código Civil actual ha superado el sistema en la regla general del artículo 1151 y ha tomado en
cuenta la teoría objetiva de la posesión para permitir que aun en los casos de delito o de violencia
(artículos 1154 y 1155) se pueda usucapir a partir del momento en que cesen los efectos de la
ilicitud, siempre que se entre a la posesión y disfrute de la misma en concepto de dueño, es decir,
como propietario de la cosa, ejercitando actos de señorío sobre ella; requisitos sin los cuales ni en
el caso de la buena fe ni en el de la mala fe se puede producir usucapión.


Amparo directo 1085/59. Pablo Hernández Fierro. 19 de noviembre de 1959. Cinco votos. Ponente:
Gabriel García Rojas.




                                             Pág. 1 de 1            Fecha de impresión 01/10/2025
  https://sjf2.scjn.gob.mx/detalle/tesis/271674

\end{quote}

\subsection{Formalization}

\subsubsection{Clause: tesis271674\_principio\_fundamental}

\textit{Thesis-specific principle: Just title not required for adverse possession, but must reveal origin of possession and affirm possession as owner}

\textbf{Intermediate Representation:}

\begin{quote}
ejercer prescripcion positiva is obligated and possession as owner is obligated
\end{quote}

\textbf{Witness Clause:}

\begin{lstlisting}[style=witnessstyle, language=Witness]
clause tesis271674_principio_fundamental = oblig(ejercer_prescripcion_positiva) IMPLIES oblig(causa_generadora_revelada) AND oblig(posesion_como_propietario);
\end{lstlisting}

\subsubsection{Clause: tesis271674\_determinacion\_judicial}

\textit{Supporting logic: Judge needs generating fact/act to determine possession quality and timing}

\textbf{Intermediate Representation:}

\begin{quote}
If causa generadora revelada is obligated, then (possession as owner is obligated and (good faith exercise is obligated and not (bad faith exercise is obligated)))
\end{quote}

\textbf{Witness Clause:}

\begin{lstlisting}[style=witnessstyle, language=Witness]
clause tesis271674_determinacion_judicial = oblig(causa_generadora_revelada) IMPLIES (oblig(posesion_como_propietario) AND (oblig(ejercicio_buena_fe) AND not(oblig(ejercicio_mala_fe))));
\end{lstlisting}

\subsubsection{Clause: tesis271674\_tipo\_posesion}

\textit{Possession type determination: Can determine if original or derived possession}

\textbf{Intermediate Representation:}

\begin{quote}
If causa generadora revelada is obligated, then (possession as owner is obligated and not (temporal possession is obligated))
\end{quote}

\textbf{Witness Clause:}

\begin{lstlisting}[style=witnessstyle, language=Witness]
clause tesis271674_tipo_posesion = oblig(causa_generadora_revelada) IMPLIES (oblig(posesion_como_propietario) AND not(oblig(posesion_temporal)));
\end{lstlisting}

\subsubsection{Clause: tesis271674\_determinacion\_temporal}

\textit{Timing requirement: Must establish when prescription period starts counting}

\textbf{Intermediate Representation:}

\begin{quote}
If causa generadora revelada is obligated, then (posesion tiempo 5 anos is obligated or posesion tiempo 10 anos is obligated)
\end{quote}

\textbf{Witness Clause:}

\begin{lstlisting}[style=witnessstyle, language=Witness]
clause tesis271674_determinacion_temporal = oblig(causa_generadora_revelada) IMPLIES (oblig(posesion_tiempo_5_anos) OR oblig(posesion_tiempo_10_anos));
\end{lstlisting}

\newpage

\section{Thesis 271675: PRESCRIPCION POSITIVA. POSESION SIN TITULO, REGLAMENTACION DE LA.}
\label{thesis:271675}

\subsection{Original Thesis Text}

\begin{quote}
                                                   Semanario Judicial de la Federación

                                                  Tesis

 Registro digital: 271675

 Instancia: Tercera Sala          Sexta Época                           Materia(s): Civil

                                  Fuente: Semanario Judicial de la   Tipo: Aislada
                                  Federación.
                                  Volumen XXIX, Cuarta Parte, página
                                  179


PRESCRIPCION POSITIVA. POSESION SIN TITULO, REGLAMENTACION DE LA.


El espíritu del Código Civil vigente fue el de implantar la teoría objetiva de la posesión, llevándola
mas lejos de donde habían llegado los códigos alemán y suizo: para conceptuar poseedora a una
persona no se exigió el "animus domini" de la escuela clásica, ni siquiera el "animus posidendi" de
la escuela de la transmisión, sino que basta, para adquirir la posesión que se ejerza un verdadero
poder de hecho sobre la cosa en provecho de quien la tiene y sin perjudicar la colectividad. La
posesión es la consagración que el derecho hace de una situación de hecho que no necesita
averiguación del criterio subjetivo del beneficiario, y que sólo varía por el reconocimiento que del
derecho de pro posesión independizándola del derecho de propiedad y de cualquier otro acto que le
sirviera de título, es decir, se reglamento la posesión sin título, como el poder de hecho que se
adquiere sobre una cosa, independientemente de toda autorización de su dueño. Por ello se cambió
el criterio subjetivo para analizar la posesión de buena fe, por elemento objetivo, estimándose que la
creencia de tener título bastante para transferir el dominio, no constituye en realidad un verdadero
título sino una posesión sin título.


Amparo directo 1085/59. Pablo Hernández Fierro. 19 de noviembre de 1959. Mayoría de tres votos.
La publicación no menciona el nombre del ponente. Engrose: José Castro Estrada.




                                             Pág. 1 de 1             Fecha de impresión 01/10/2025
  https://sjf2.scjn.gob.mx/detalle/tesis/271675

\end{quote}

\subsection{Formalization}

\subsubsection{Clause: tesis271675\_principio\_fundamental}

\textit{Thesis-specific principle: Good faith possessors require just title as generating cause of possession for adverse possession to operate}

\textbf{Intermediate Representation:}

\begin{quote}
good faith exercise is obligated and possession as owner is obligated
\end{quote}

\textbf{Witness Clause:}

\begin{lstlisting}[style=witnessstyle, language=Witness]
clause tesis271675_principio_fundamental = oblig(ejercicio_buena_fe) AND oblig(posesion_como_propietario) IMPLIES oblig(causa_generadora_probada);
\end{lstlisting}

\subsubsection{Clause: tesis271675\_revelar\_causa}

\textit{Supporting logic: Must reveal and fully prove the generating cause of possession}

\textbf{Intermediate Representation:}

\begin{quote}
ejercer prescripcion positiva is obligated
\end{quote}

\textbf{Witness Clause:}

\begin{lstlisting}[style=witnessstyle, language=Witness]
clause tesis271675_revelar_causa = oblig(ejercer_prescripcion_positiva) IMPLIES oblig(causa_generadora_revelada);
\end{lstlisting}

\subsubsection{Clause: tesis271675\_determinacion\_judicial}

\textit{Judicial determination: Judge needs generating fact/act to determine possession quality and timing}

\textbf{Intermediate Representation:}

\begin{quote}
If causa generadora revelada is obligated, then (possession as owner is obligated and (good faith exercise is obligated and not (bad faith exercise is obligated)))
\end{quote}

\textbf{Witness Clause:}

\begin{lstlisting}[style=witnessstyle, language=Witness]
clause tesis271675_determinacion_judicial = oblig(causa_generadora_revelada) IMPLIES (oblig(posesion_como_propietario) AND (oblig(ejercicio_buena_fe) AND not(oblig(ejercicio_mala_fe))));
\end{lstlisting}

\subsubsection{Clause: tesis271675\_determinacion\_temporal}

\textit{Timing requirement: Must establish when prescription period begins}

\textbf{Intermediate Representation:}

\begin{quote}
If causa generadora revelada is obligated, then (posesion tiempo 5 anos is obligated or posesion tiempo 10 anos is obligated)
\end{quote}

\textbf{Witness Clause:}

\begin{lstlisting}[style=witnessstyle, language=Witness]
clause tesis271675_determinacion_temporal = oblig(causa_generadora_revelada) IMPLIES (oblig(posesion_tiempo_5_anos) OR oblig(posesion_tiempo_10_anos));
\end{lstlisting}

\newpage

\section{Thesis 272051: USUCAPION. NO ES NECESARIO EL JUSTO TITULO PARA FUNDARLA.}
\label{thesis:272051}

\subsection{Original Thesis Text}

\begin{quote}
                                                   Semanario Judicial de la Federación

                                                  Tesis

 Registro digital: 272051

 Instancia: Tercera Sala           Sexta Época                            Materia(s): Civil

                                   Fuente: Semanario Judicial de la   Tipo: Aislada
                                   Federación.
                                   Volumen XXII, Cuarta Parte, página
                                   376


USUCAPION. NO ES NECESARIO EL JUSTO TITULO PARA FUNDARLA.


El artículo 806 del Código Civil, al establecer que es poseedor de buena fe el que entra en la
posesión en virtud de un título suficiente para darle derecho de poseer y que también lo es el que
ignora los vicios de su título que le impiden poseer con derecho y que es poseedor de mala fe el
que entra a la posesión sin título alguno para poseer, lo mismo que el que conoce los vicios de su
título que le impiden poseer con derecho, pone de manifiesto que el propósito del legislador fue
cambiar el sistema del Código de 1884, que exigía en su artículo 1079, fracción I, que la posesión
ad usucapionem debería fundarse el justo título. El código actual exige en la fracción I del artículo
1151, que la posesión necesaria para prescribir debe ser en concepto de propietario, y no exige ya
el justo título sino que adopta un sistema objetivo sobre la materia de la posesión tendiente a
facilitar la solución de los problemas que en este punto planteaba la legislación anterior; pero si bien
no exige la ley el justo título, es necesario probar el origen de la posesión, no como acto traslativo
de dominio sino como hecho jurídico que produce consecuencias de derecho, para conocer la fecha
cierta a partir de la cual ha de computarse el término legal de la prescripción, pero a condición de
que el poseedor se comporte como propietario, esto es, que se conduzca ostensiblemente y de
manera objetiva, susceptible de apreciarse por los sentidos, mediante actos que revelen que el
poseedor es el dominador de la cosa, el señor de ella, el que manda en la misma, como dueño en
sentido económico, para hacer suya la cosa desde el punto de vista de los hechos, aun cuando
carezca de título jurídicamente hablando.


Amparo directo 7279/56. Francisco Cuevas Cansino y coagraviados. 27 de abril de 1959. Mayoría
de tres votos. Ponente: Gabriel García Rojas.




                                              Pág. 1 de 1              Fecha de impresión 01/10/2025
  https://sjf2.scjn.gob.mx/detalle/tesis/272051

\end{quote}

\subsection{Formalization}

\subsubsection{Clause: tesis272051\_principio\_fundamental}

\textit{Thesis-specific principle: Good faith possessors require just title as generating cause of possession for adverse possession to operate}

\textbf{Intermediate Representation:}

\begin{quote}
good faith exercise is obligated and possession as owner is obligated
\end{quote}

\textbf{Witness Clause:}

\begin{lstlisting}[style=witnessstyle, language=Witness]
clause tesis272051_principio_fundamental = oblig(ejercicio_buena_fe) AND oblig(posesion_como_propietario) IMPLIES oblig(causa_generadora_probada);
\end{lstlisting}

\subsubsection{Clause: tesis272051\_revelar\_causa}

\textit{Supporting logic: Must reveal and fully prove the generating cause of possession}

\textbf{Intermediate Representation:}

\begin{quote}
ejercer prescripcion positiva is obligated
\end{quote}

\textbf{Witness Clause:}

\begin{lstlisting}[style=witnessstyle, language=Witness]
clause tesis272051_revelar_causa = oblig(ejercer_prescripcion_positiva) IMPLIES oblig(causa_generadora_revelada);
\end{lstlisting}

\subsubsection{Clause: tesis272051\_determinacion\_judicial}

\textit{Judicial determination: Judge needs generating fact/act to determine possession quality and timing}

\textbf{Intermediate Representation:}

\begin{quote}
If causa generadora revelada is obligated, then (possession as owner is obligated and (good faith exercise is obligated and not (bad faith exercise is obligated)))
\end{quote}

\textbf{Witness Clause:}

\begin{lstlisting}[style=witnessstyle, language=Witness]
clause tesis272051_determinacion_judicial = oblig(causa_generadora_revelada) IMPLIES (oblig(posesion_como_propietario) AND (oblig(ejercicio_buena_fe) AND not(oblig(ejercicio_mala_fe))));
\end{lstlisting}

\subsubsection{Clause: tesis272051\_determinacion\_temporal}

\textit{Timing requirement: Must establish when prescription period begins}

\textbf{Intermediate Representation:}

\begin{quote}
If causa generadora revelada is obligated, then (posesion tiempo 5 anos is obligated or posesion tiempo 10 anos is obligated)
\end{quote}

\textbf{Witness Clause:}

\begin{lstlisting}[style=witnessstyle, language=Witness]
clause tesis272051_determinacion_temporal = oblig(causa_generadora_revelada) IMPLIES (oblig(posesion_tiempo_5_anos) OR oblig(posesion_tiempo_10_anos));
\end{lstlisting}

\newpage

\section{Thesis 272139: PRESCRIPCION POSITIVA. POSESION SIN TITULO DE DOMINIO}
\label{thesis:272139}

\subsection{Original Thesis Text}

\begin{quote}
                                                   Semanario Judicial de la Federación

                                                  Tesis

 Registro digital: 272139

 Instancia: Tercera Sala          Sexta Época                           Materia(s): Civil

                                  Fuente: Semanario Judicial de la      Tipo: Aislada
                                  Federación.
                                  Volumen XX, Cuarta Parte, página
                                  193


PRESCRIPCION POSITIVA. POSESION SIN TITULO DE DOMINIO


Es necesario que el actor, al promover el juicio de prescripción revele la causa de su posesión, esto
es, el hecho o acto generador de la misma, tanto para que el juzgador pueda determinar la calidad
de la posesión, originaria o derivada, como para que pueda computarse el término de la
prescripción, bien sea la posesión de buena o mala fe o por causa de delito. En consecuencia, si el
actor no alude a la causa de su posesión, ni en el curso del juicio la demuestra, no puede estimarse
que haya acreditado que su posesión reúne el requisito a que se contrae la fracción I del artículo
1151.


Amparo directo 4215/57. Inversiones de México, S. A. 4 de febrero de 1959. Cinco votos. Ponente:
José Castro Estrada.




                                             Pág. 1 de 1             Fecha de impresión 01/10/2025
  https://sjf2.scjn.gob.mx/detalle/tesis/272139

\end{quote}

\subsection{Formalization}

\subsubsection{Clause: tesis272139\_principio\_fundamental}

\textit{Thesis-specific principle: Good faith possessors require just title as generating cause of possession for adverse possession to operate}

\textbf{Intermediate Representation:}

\begin{quote}
good faith exercise is obligated and possession as owner is obligated
\end{quote}

\textbf{Witness Clause:}

\begin{lstlisting}[style=witnessstyle, language=Witness]
clause tesis272139_principio_fundamental = oblig(ejercicio_buena_fe) AND oblig(posesion_como_propietario) IMPLIES oblig(causa_generadora_probada);
\end{lstlisting}

\subsubsection{Clause: tesis272139\_revelar\_causa}

\textit{Supporting logic: Must reveal and fully prove the generating cause of possession}

\textbf{Intermediate Representation:}

\begin{quote}
ejercer prescripcion positiva is obligated
\end{quote}

\textbf{Witness Clause:}

\begin{lstlisting}[style=witnessstyle, language=Witness]
clause tesis272139_revelar_causa = oblig(ejercer_prescripcion_positiva) IMPLIES oblig(causa_generadora_revelada);
\end{lstlisting}

\subsubsection{Clause: tesis272139\_determinacion\_judicial}

\textit{Judicial determination: Judge needs generating fact/act to determine possession quality and timing}

\textbf{Intermediate Representation:}

\begin{quote}
If causa generadora revelada is obligated, then (possession as owner is obligated and (good faith exercise is obligated and not (bad faith exercise is obligated)))
\end{quote}

\textbf{Witness Clause:}

\begin{lstlisting}[style=witnessstyle, language=Witness]
clause tesis272139_determinacion_judicial = oblig(causa_generadora_revelada) IMPLIES (oblig(posesion_como_propietario) AND (oblig(ejercicio_buena_fe) AND not(oblig(ejercicio_mala_fe))));
\end{lstlisting}

\subsubsection{Clause: tesis272139\_determinacion\_temporal}

\textit{Timing requirement: Must establish when prescription period begins}

\textbf{Intermediate Representation:}

\begin{quote}
If causa generadora revelada is obligated, then (posesion tiempo 5 anos is obligated or posesion tiempo 10 anos is obligated)
\end{quote}

\textbf{Witness Clause:}

\begin{lstlisting}[style=witnessstyle, language=Witness]
clause tesis272139_determinacion_temporal = oblig(causa_generadora_revelada) IMPLIES (oblig(posesion_tiempo_5_anos) OR oblig(posesion_tiempo_10_anos));
\end{lstlisting}

\newpage

\section{Thesis 272536: PRESCRIPCION ADQUISITIVA (LEGISLACION DE PUEBLA).}
\label{thesis:272536}

\subsection{Original Thesis Text}

\begin{quote}
                                                   Semanario Judicial de la Federación

                                                  Tesis

 Registro digital: 272536

 Instancia: Tercera Sala          Sexta Época                          Materia(s): Civil

                                  Fuente: Semanario Judicial de la     Tipo: Aislada
                                  Federación.
                                  Volumen XIII, Cuarta Parte, página
                                  265


PRESCRIPCION ADQUISITIVA (LEGISLACION DE PUEBLA).


El artículo 1069 del Código Civil establece que la posesión necesaria para prescribir debe ser
fundada en justo título, de buena fe, pacífica, contínua y pública. Por lo tanto, es indudable que
quien invoca la prescripción adquisitiva, debe demostrar que existe un acto jurídico de los que
transfieren el dominio, como causa de la posesión. A esta norma no escapa el juicio de prescripción.
El artículo 567 del Código de Procedimientos Civiles establece: El que hubiere poseído bienes
inmuebles por el tiempo y con las condiciones que señala el Código Civil para la prescripción
positiva, puede promover juicio contra el propietario de esos bienes, a fin de que se declare que ha
adquirido, por ende, la propiedad. Este precepto indica claramente que la prescripción deberá
haberse consumado por el tiempo y con las condiciones que establece para ello, siendo inexcusable
y primordial, la del justo título.


Amparo directo 4171/57. Eulalia Rojas Domínguez. 3 de julio de 1958. Cinco votos. Ponente:
Mariano Ramírez Vázquez.




                                             Pág. 1 de 1            Fecha de impresión 01/10/2025
  https://sjf2.scjn.gob.mx/detalle/tesis/272536

\end{quote}

\subsection{Formalization}

\subsubsection{Clause: tesis272536\_principio\_fundamental}

\textit{Thesis-specific principle: Possession must be "in concept of owner" (peaceful, continuous, public) for both good and bad faith}

\textbf{Intermediate Representation:}

\begin{quote}
possession as owner is obligated and (oblig(ejercicio\_buena\_fe) OR oblig(ejercicio\_mala\_fe)) IMPLIES oblig(ejercer\_prescripcion\_positiva)
\end{quote}

\textbf{Witness Clause:}

\begin{lstlisting}[style=witnessstyle, language=Witness]
clause tesis272536_principio_fundamental = oblig(posesion_como_propietario) AND (oblig(ejercicio_buena_fe) OR oblig(ejercicio_mala_fe)) IMPLIES oblig(ejercer_prescripcion_positiva);
\end{lstlisting}

\subsubsection{Clause: tesis272536\_posesion\_ambas\_fe}

\textit{Supporting logic: Possession as owner applies to both good faith and bad faith cases}

\textbf{Intermediate Representation:}

\begin{quote}
If possession as owner is obligated, then (good faith exercise is obligated or bad faith exercise is obligated)
\end{quote}

\textbf{Witness Clause:}

\begin{lstlisting}[style=witnessstyle, language=Witness]
clause tesis272536_posesion_ambas_fe = oblig(posesion_como_propietario) IMPLIES (oblig(ejercicio_buena_fe) OR oblig(ejercicio_mala_fe));
\end{lstlisting}

\subsubsection{Clause: tesis272536\_dominio\_economico}

\textit{Economic dominion: Must be dominator of the thing, commanding and enjoying it as owner}

\textbf{Intermediate Representation:}

\begin{quote}
possession as owner is obligated
\end{quote}

\textbf{Witness Clause:}

\begin{lstlisting}[style=witnessstyle, language=Witness]
clause tesis272536_dominio_economico = oblig(posesion_como_propietario) IMPLIES oblig(causa_generadora_probada);
\end{lstlisting}

\subsubsection{Clause: tesis272536\_posesion\_originaria}

\textit{Original possession: Must begin with cause different from derived possession}

\textbf{Intermediate Representation:}

\begin{quote}
possession as owner is obligated and not (temporal possession is obligated)
\end{quote}

\textbf{Witness Clause:}

\begin{lstlisting}[style=witnessstyle, language=Witness]
clause tesis272536_posesion_originaria = oblig(posesion_como_propietario) AND not(oblig(posesion_temporal)) IMPLIES oblig(causa_generadora_probada);
\end{lstlisting}

\newpage

\section{Thesis 273071: PRESCRIPCION POSITIVA, POSESION PARA LA.}
\label{thesis:273071}

\subsection{Original Thesis Text}

\begin{quote}
                                                   Semanario Judicial de la Federación

                                                  Tesis

 Registro digital: 273071

 Instancia: Tercera Sala          Sexta Época                           Materia(s): Civil

                                  Fuente: Semanario Judicial de la      Tipo: Aislada
                                  Federación.
                                  Volumen II, Cuarta Parte, página
                                  125


PRESCRIPCION POSITIVA, POSESION PARA LA.


Es cierto que el primer párrafo del artículo 798 del Código Civil establece que la posesión da al que
la tiene la presunción de propietario para todos los efectos legales, pero aquí la ley se refiere al
poseedor original, y para que el juzgador pueda determinar la calidad de la posesión del actor en un
juicio de prescripción positiva, es indispensable que éste dé a conocer el hecho o acto generador de
la posesión, incluso para que pueda computarse el término de la prescripción, bien sea de buena o
mala fe, o por causa de delito.


Amparo directo 2261/57. Carlos Espíritu Nava. 28 de agosto de 1957. Cinco votos. Ponente: José
Castro Estrada.




                                             Pág. 1 de 1             Fecha de impresión 01/10/2025
  https://sjf2.scjn.gob.mx/detalle/tesis/273071

\end{quote}

\subsection{Formalization}

\subsubsection{Clause: tesis273071\_principio\_fundamental}

\textit{Additional actions for this thesis Additional assets for proof requirements Faith exclusivity: Adverse possession is either good faith XOR bad faith Thesis-specific principle: Good faith adverse possession requires proving authenticity, date, and diligence}

\textbf{Intermediate Representation:}

\begin{quote}
If good faith exercise is obligated and adverse possession claim is claimed, then (autenticidad titulo is obligated and fecha fehaciente is obligated and diligencia adquirente is obligated)
\end{quote}

\textbf{Witness Clause:}

\begin{lstlisting}[style=witnessstyle, language=Witness]
clause tesis273071_principio_fundamental = oblig(ejercicio_buena_fe) AND claim(demanda_prescripcion) IMPLIES (oblig(autenticidad_titulo) AND oblig(fecha_fehaciente) AND oblig(diligencia_adquirente));
\end{lstlisting}

\subsubsection{Clause: tesis273071\_prueba\_autenticidad}

\textit{Proof requirements: Must prove title authenticity and date reliably}

\textbf{Intermediate Representation:}

\begin{quote}
If proven generating cause is obligated, then (autenticidad titulo is obligated and fecha fehaciente is obligated)
\end{quote}

\textbf{Witness Clause:}

\begin{lstlisting}[style=witnessstyle, language=Witness]
clause tesis273071_prueba_autenticidad = oblig(causa_generadora_probada) IMPLIES (oblig(autenticidad_titulo) AND oblig(fecha_fehaciente));
\end{lstlisting}

\subsubsection{Clause: tesis273071\_diligencia\_buena\_fe}

\textit{Diligence requirement: Good faith requires demonstrating due diligence in acquisition}

\textbf{Intermediate Representation:}

\begin{quote}
good faith exercise is obligated
\end{quote}

\textbf{Witness Clause:}

\begin{lstlisting}[style=witnessstyle, language=Witness]
clause tesis273071_diligencia_buena_fe = oblig(ejercicio_buena_fe) IMPLIES oblig(diligencia_adquirente);
\end{lstlisting}

\subsubsection{Clause: tesis273071\_prueba\_pagos}

\textit{Payment proof: Onerous contracts require proving payments made}

\textbf{Intermediate Representation:}

\begin{quote}
proven generating cause is obligated and contrato compraventa is obligated
\end{quote}

\textbf{Witness Clause:}

\begin{lstlisting}[style=witnessstyle, language=Witness]
clause tesis273071_prueba_pagos = oblig(causa_generadora_probada) AND oblig(contrato_compraventa) IMPLIES oblig(pagos_precio);
\end{lstlisting}

\subsubsection{Clause: tesis273071\_carga\_prueba}

\textit{Burden of proof: Claimant must prove generating cause with reliable evidence}

\textbf{Intermediate Representation:}

\begin{quote}
adverse possession claim is claimed
\end{quote}

\textbf{Witness Clause:}

\begin{lstlisting}[style=witnessstyle, language=Witness]
clause tesis273071_carga_prueba = claim(demanda_prescripcion) IMPLIES oblig(causa_generadora_revelada);
\end{lstlisting}

\subsubsection{Clause: tesis273071\_falta\_diligencia}

\textit{Insufficient diligence consequence: Lack of diligence means bad faith (10-year period)}

\textbf{Intermediate Representation:}

\begin{quote}
not (diligencia adquirente is obligated)
\end{quote}

\textbf{Witness Clause:}

\begin{lstlisting}[style=witnessstyle, language=Witness]
clause tesis273071_falta_diligencia = not(oblig(diligencia_adquirente)) IMPLIES oblig(ejercicio_mala_fe);
\end{lstlisting}

\newpage

\section{Thesis 273127: PRESCRIPCION POSITIVA. POSESION DE BUENA FE.}
\label{thesis:273127}

\subsection{Original Thesis Text}

\begin{quote}
                                                   Semanario Judicial de la Federación

                                                  Tesis

 Registro digital: 273127

 Instancia: Tercera Sala          Sexta Época                         Materia(s): Civil

                                  Fuente: Semanario Judicial de la    Tipo: Aislada
                                  Federación.
                                  Volumen I, Cuarta Parte, página 134


PRESCRIPCION POSITIVA. POSESION DE BUENA FE.


No puede admitirse que el comprador que confiesa haber comprado a una persona un bien que es
de otra, posea de buena fe, ya que conoce el vicio fundamental de que adolece su título.


Amparo directo 6179/56. Teodomiro González y coagraviado. 12 de julio de 1957. Cinco votos.
Ponente: Mariano Ramírez Vázquez.




                                             Pág. 1 de 1           Fecha de impresión 01/10/2025
  https://sjf2.scjn.gob.mx/detalle/tesis/273127

\end{quote}

\subsection{Formalization}

\subsubsection{Clause: tesis273127\_principio\_fundamental}

\textit{Additional actions for this thesis Additional assets for proof requirements Faith exclusivity: Adverse possession is either good faith XOR bad faith Thesis-specific principle: Good faith adverse possession requires proving authenticity, date, and diligence}

\textbf{Intermediate Representation:}

\begin{quote}
If good faith exercise is obligated and adverse possession claim is claimed, then (autenticidad titulo is obligated and fecha fehaciente is obligated and diligencia adquirente is obligated)
\end{quote}

\textbf{Witness Clause:}

\begin{lstlisting}[style=witnessstyle, language=Witness]
clause tesis273127_principio_fundamental = oblig(ejercicio_buena_fe) AND claim(demanda_prescripcion) IMPLIES (oblig(autenticidad_titulo) AND oblig(fecha_fehaciente) AND oblig(diligencia_adquirente));
\end{lstlisting}

\subsubsection{Clause: tesis273127\_prueba\_autenticidad}

\textit{Proof requirements: Must prove title authenticity and date reliably}

\textbf{Intermediate Representation:}

\begin{quote}
If proven generating cause is obligated, then (autenticidad titulo is obligated and fecha fehaciente is obligated)
\end{quote}

\textbf{Witness Clause:}

\begin{lstlisting}[style=witnessstyle, language=Witness]
clause tesis273127_prueba_autenticidad = oblig(causa_generadora_probada) IMPLIES (oblig(autenticidad_titulo) AND oblig(fecha_fehaciente));
\end{lstlisting}

\subsubsection{Clause: tesis273127\_diligencia\_buena\_fe}

\textit{Diligence requirement: Good faith requires demonstrating due diligence in acquisition}

\textbf{Intermediate Representation:}

\begin{quote}
good faith exercise is obligated
\end{quote}

\textbf{Witness Clause:}

\begin{lstlisting}[style=witnessstyle, language=Witness]
clause tesis273127_diligencia_buena_fe = oblig(ejercicio_buena_fe) IMPLIES oblig(diligencia_adquirente);
\end{lstlisting}

\subsubsection{Clause: tesis273127\_prueba\_pagos}

\textit{Payment proof: Onerous contracts require proving payments made}

\textbf{Intermediate Representation:}

\begin{quote}
proven generating cause is obligated and contrato compraventa is obligated
\end{quote}

\textbf{Witness Clause:}

\begin{lstlisting}[style=witnessstyle, language=Witness]
clause tesis273127_prueba_pagos = oblig(causa_generadora_probada) AND oblig(contrato_compraventa) IMPLIES oblig(pagos_precio);
\end{lstlisting}

\subsubsection{Clause: tesis273127\_carga\_prueba}

\textit{Burden of proof: Claimant must prove generating cause with reliable evidence}

\textbf{Intermediate Representation:}

\begin{quote}
adverse possession claim is claimed
\end{quote}

\textbf{Witness Clause:}

\begin{lstlisting}[style=witnessstyle, language=Witness]
clause tesis273127_carga_prueba = claim(demanda_prescripcion) IMPLIES oblig(causa_generadora_revelada);
\end{lstlisting}

\subsubsection{Clause: tesis273127\_falta\_diligencia}

\textit{Insufficient diligence consequence: Lack of diligence means bad faith (10-year period)}

\textbf{Intermediate Representation:}

\begin{quote}
not (diligencia adquirente is obligated)
\end{quote}

\textbf{Witness Clause:}

\begin{lstlisting}[style=witnessstyle, language=Witness]
clause tesis273127_falta_diligencia = not(oblig(diligencia_adquirente)) IMPLIES oblig(ejercicio_mala_fe);
\end{lstlisting}

\newpage

\section{Thesis 339015: PRESCRIPCION ADQUISITIVA, EN CASO DE COMPRAVENTA DEL BIEN.}
\label{thesis:339015}

\subsection{Original Thesis Text}

\begin{quote}
                                                   Semanario Judicial de la Federación

                                                  Tesis

 Registro digital: 339015

 Instancia: Tercera Sala          Quinta Época                           Materia(s): Civil

                                  Fuente: Semanario Judicial de la       Tipo: Aislada
                                  Federación.
                                  Tomo CXXXI, página 689


PRESCRIPCION ADQUISITIVA, EN CASO DE COMPRAVENTA DEL BIEN.


No siendo necesaria la inscripción en el registro de la posesión del actor de un juicio de prescripción
para que éste adquiriese por prescripción, tiene que aceptarse que al mismo no puede perjudicar,
como medio para adquirir la propiedad por prescripción, la inscripción de los actos registrables, ni la
buena fe, por tanto, del comprador del bien usucapido, compra hecha a quien tenía el bien
registrado a su nombre.


Amparo directo 5745/54. Praxedis Rosales. 20 de marzo de 1957. Unanimidad de cinco votos.
Ponente: Gabriel García Rojas.




                                             Pág. 1 de 1              Fecha de impresión 01/10/2025
  https://sjf2.scjn.gob.mx/detalle/tesis/339015

\end{quote}

\subsection{Formalization}

\subsubsection{Clause: tesis339015\_principio\_fundamental}

\textit{Additional actions for this thesis Additional assets for proof requirements Faith exclusivity: Adverse possession is either good faith XOR bad faith Thesis-specific principle: Good faith adverse possession requires proving authenticity, date, and diligence}

\textbf{Intermediate Representation:}

\begin{quote}
If good faith exercise is obligated and adverse possession claim is claimed, then (autenticidad titulo is obligated and fecha fehaciente is obligated and diligencia adquirente is obligated)
\end{quote}

\textbf{Witness Clause:}

\begin{lstlisting}[style=witnessstyle, language=Witness]
clause tesis339015_principio_fundamental = oblig(ejercicio_buena_fe) AND claim(demanda_prescripcion) IMPLIES (oblig(autenticidad_titulo) AND oblig(fecha_fehaciente) AND oblig(diligencia_adquirente));
\end{lstlisting}

\subsubsection{Clause: tesis339015\_prueba\_autenticidad}

\textit{Proof requirements: Must prove title authenticity and date reliably}

\textbf{Intermediate Representation:}

\begin{quote}
If proven generating cause is obligated, then (autenticidad titulo is obligated and fecha fehaciente is obligated)
\end{quote}

\textbf{Witness Clause:}

\begin{lstlisting}[style=witnessstyle, language=Witness]
clause tesis339015_prueba_autenticidad = oblig(causa_generadora_probada) IMPLIES (oblig(autenticidad_titulo) AND oblig(fecha_fehaciente));
\end{lstlisting}

\subsubsection{Clause: tesis339015\_diligencia\_buena\_fe}

\textit{Diligence requirement: Good faith requires demonstrating due diligence in acquisition}

\textbf{Intermediate Representation:}

\begin{quote}
good faith exercise is obligated
\end{quote}

\textbf{Witness Clause:}

\begin{lstlisting}[style=witnessstyle, language=Witness]
clause tesis339015_diligencia_buena_fe = oblig(ejercicio_buena_fe) IMPLIES oblig(diligencia_adquirente);
\end{lstlisting}

\subsubsection{Clause: tesis339015\_prueba\_pagos}

\textit{Payment proof: Onerous contracts require proving payments made}

\textbf{Intermediate Representation:}

\begin{quote}
proven generating cause is obligated and contrato compraventa is obligated
\end{quote}

\textbf{Witness Clause:}

\begin{lstlisting}[style=witnessstyle, language=Witness]
clause tesis339015_prueba_pagos = oblig(causa_generadora_probada) AND oblig(contrato_compraventa) IMPLIES oblig(pagos_precio);
\end{lstlisting}

\subsubsection{Clause: tesis339015\_carga\_prueba}

\textit{Burden of proof: Claimant must prove generating cause with reliable evidence}

\textbf{Intermediate Representation:}

\begin{quote}
adverse possession claim is claimed
\end{quote}

\textbf{Witness Clause:}

\begin{lstlisting}[style=witnessstyle, language=Witness]
clause tesis339015_carga_prueba = claim(demanda_prescripcion) IMPLIES oblig(causa_generadora_revelada);
\end{lstlisting}

\subsubsection{Clause: tesis339015\_falta\_diligencia}

\textit{Insufficient diligence consequence: Lack of diligence means bad faith (10-year period)}

\textbf{Intermediate Representation:}

\begin{quote}
not (diligencia adquirente is obligated)
\end{quote}

\textbf{Witness Clause:}

\begin{lstlisting}[style=witnessstyle, language=Witness]
clause tesis339015_falta_diligencia = not(oblig(diligencia_adquirente)) IMPLIES oblig(ejercicio_mala_fe);
\end{lstlisting}

\newpage

\section{Thesis 340211: PRESCRIPCION POSITIVA (LEGISLACIONES DE VERACRUZ Y DEL DISTRITO FEDERAL).}
\label{thesis:340211}

\subsection{Original Thesis Text}

\begin{quote}
                                                   Semanario Judicial de la Federación

                                                  Tesis

 Registro digital: 340211

 Instancia: Tercera Sala          Quinta Época                           Materia(s): Civil

                                  Fuente: Semanario Judicial de la       Tipo: Aislada
                                  Federación.
                                  Tomo CXXIII, página 313


PRESCRIPCION POSITIVA (LEGISLACIONES DE VERACRUZ Y DEL DISTRITO FEDERAL).


El derecho francés distingue los vicios de forma respecto de la adquisición de los frutos y esos
mismos vicios en materia de usucapión, y establece, por una parte, que los primeros son
inoperantes, si los ignora el poseedor, respecto de la adquisición de tales frutos, y por otra, que "el
título nulo por defecto de forma no puede servir de base a la prescripción de diez y veinte años
(artículo 2267 del Código de Napoleón). Ahora bien, en la legislación mexicana no hay esa
distinción, ya que un título nulo por vicio de forma puede servir de base a la posesión de buena fe,
tanto para la adquisición de los frutos como para la usucapión, con tal de que se cumplan las
condiciones para que por el comportamiento de las partes se subsane la nulidad por defecto de
forma (artículo 2167 del Código Civil de Veracruz, igual al 2234 del código de la misma materia del
Distrito Federal), aun cuando sin dejar de reconocer que hay vicios de forma que no pueden ser
subsanados por dicho comportamiento de las partes, como son la falta de formalidad para la venta
de bienes de menores, formalidad que establecen el artículo 365 del código de Veracruz, igual al
436 del código del Distrito Federal, para los sujetos a la patria potestad; el 491 y siguientes del
primero, igual al 561 y siguiente del segundo, respecto de los tutores, y el 105 de Veracruz, que
corresponde al 174 del Distrito, para la mujer que contrate con su marido; pero en ninguno de esos
casos ni en otro alguno de excepción está el de quien trata de prescribir con un título nulo, por el
contrario, dicho caso queda comprendido dentro de la regla general de que el título nulo por vicio de
forma puede servir de base para la prescripción adquisitiva.


Amparo Civil Directo 4824/53. Díaz viuda de Carmona Catalina y coagraviados. 19 de enero de
1955. Unanimidad de cuatro votos. Ponente Gabriel García Rojas.




                                             Pág. 1 de 1              Fecha de impresión 01/10/2025
  https://sjf2.scjn.gob.mx/detalle/tesis/340211

\end{quote}

\subsection{Formalization}

\subsubsection{Clause: tesis340211\_principio\_fundamental}

\textit{Thesis-specific principle: Good faith possessors require just title as generating cause of possession for adverse possession to operate}

\textbf{Intermediate Representation:}

\begin{quote}
good faith exercise is obligated and possession as owner is obligated
\end{quote}

\textbf{Witness Clause:}

\begin{lstlisting}[style=witnessstyle, language=Witness]
clause tesis340211_principio_fundamental = oblig(ejercicio_buena_fe) AND oblig(posesion_como_propietario) IMPLIES oblig(causa_generadora_probada);
\end{lstlisting}

\subsubsection{Clause: tesis340211\_revelar\_causa}

\textit{Supporting logic: Must reveal and fully prove the generating cause of possession}

\textbf{Intermediate Representation:}

\begin{quote}
ejercer prescripcion positiva is obligated
\end{quote}

\textbf{Witness Clause:}

\begin{lstlisting}[style=witnessstyle, language=Witness]
clause tesis340211_revelar_causa = oblig(ejercer_prescripcion_positiva) IMPLIES oblig(causa_generadora_revelada);
\end{lstlisting}

\subsubsection{Clause: tesis340211\_determinacion\_judicial}

\textit{Judicial determination: Judge needs generating fact/act to determine possession quality and timing}

\textbf{Intermediate Representation:}

\begin{quote}
If causa generadora revelada is obligated, then (possession as owner is obligated and (good faith exercise is obligated and not (bad faith exercise is obligated)))
\end{quote}

\textbf{Witness Clause:}

\begin{lstlisting}[style=witnessstyle, language=Witness]
clause tesis340211_determinacion_judicial = oblig(causa_generadora_revelada) IMPLIES (oblig(posesion_como_propietario) AND (oblig(ejercicio_buena_fe) AND not(oblig(ejercicio_mala_fe))));
\end{lstlisting}

\subsubsection{Clause: tesis340211\_determinacion\_temporal}

\textit{Timing requirement: Must establish when prescription period begins}

\textbf{Intermediate Representation:}

\begin{quote}
If causa generadora revelada is obligated, then (posesion tiempo 5 anos is obligated or posesion tiempo 10 anos is obligated)
\end{quote}

\textbf{Witness Clause:}

\begin{lstlisting}[style=witnessstyle, language=Witness]
clause tesis340211_determinacion_temporal = oblig(causa_generadora_revelada) IMPLIES (oblig(posesion_tiempo_5_anos) OR oblig(posesion_tiempo_10_anos));
\end{lstlisting}

\newpage

\section{Thesis 341704: PRESCRIPCION POSITIVA, REQUISITOS PARA LA (JUSTO TITULO).}
\label{thesis:341704}

\subsection{Original Thesis Text}

\begin{quote}
                                                   Semanario Judicial de la Federación

                                                  Tesis

 Registro digital: 341704

 Instancia: Sala Auxiliar         Quinta Época                         Materia(s): Civil

                                  Fuente: Semanario Judicial de la     Tipo: Aislada
                                  Federación.
                                  Tomo CXVII, página 1131


PRESCRIPCION POSITIVA, REQUISITOS PARA LA (JUSTO TITULO).


El Código Civil de mil ochocientos ochenta y cuatro, en su artículo 1079, exige que la posesión se
funde en justo título y que sea de buena fe, pacífica, continua y pública; el artículo 1080 del mismo
código establece que se llama justo título el que es o fundadamente se cree bastante para transferir
el dominio, y el artículo 1081 del citado cuerpo de leyes determina que el que alega la prescripción
debe probar la existencia del título en que funda su derecho. El Código Civil vigente, en su artículo
1151, exige que la posesión necesaria para prescribir debe ser en concepto de propietario, pacífica,
continua y pública; por lo que, en consecuencia, es requisito esencial, de acuerdo con los preceptos
citados, que la posesión provenga de un justo título y sea en concepto de propietario, para que
pueda dar lugar a la prescripción adquisitiva.


Amparo civil directo 4679/50. Benavides R. Jesús y coagraviado. 10 de octubre de 1952. Mayoría
de tres votos. Disidentes: Rafael Matos Escobedo y Felipe Tena Ramírez. La publicación no
menciona el nombre del ponente.




                                             Pág. 1 de 1             Fecha de impresión 01/10/2025
  https://sjf2.scjn.gob.mx/detalle/tesis/341704

\end{quote}

\subsection{Formalization}

\subsubsection{Clause: tesis341704\_principio\_fundamental}

\textit{Thesis-specific principle: Possession must be "in concept of owner" (peaceful, continuous, public) for both good and bad faith}

\textbf{Intermediate Representation:}

\begin{quote}
possession as owner is obligated and (oblig(ejercicio\_buena\_fe) OR oblig(ejercicio\_mala\_fe)) IMPLIES oblig(ejercer\_prescripcion\_positiva)
\end{quote}

\textbf{Witness Clause:}

\begin{lstlisting}[style=witnessstyle, language=Witness]
clause tesis341704_principio_fundamental = oblig(posesion_como_propietario) AND (oblig(ejercicio_buena_fe) OR oblig(ejercicio_mala_fe)) IMPLIES oblig(ejercer_prescripcion_positiva);
\end{lstlisting}

\subsubsection{Clause: tesis341704\_posesion\_ambas\_fe}

\textit{Supporting logic: Possession as owner applies to both good faith and bad faith cases}

\textbf{Intermediate Representation:}

\begin{quote}
If possession as owner is obligated, then (good faith exercise is obligated or bad faith exercise is obligated)
\end{quote}

\textbf{Witness Clause:}

\begin{lstlisting}[style=witnessstyle, language=Witness]
clause tesis341704_posesion_ambas_fe = oblig(posesion_como_propietario) IMPLIES (oblig(ejercicio_buena_fe) OR oblig(ejercicio_mala_fe));
\end{lstlisting}

\subsubsection{Clause: tesis341704\_dominio\_economico}

\textit{Economic dominion: Must be dominator of the thing, commanding and enjoying it as owner}

\textbf{Intermediate Representation:}

\begin{quote}
possession as owner is obligated
\end{quote}

\textbf{Witness Clause:}

\begin{lstlisting}[style=witnessstyle, language=Witness]
clause tesis341704_dominio_economico = oblig(posesion_como_propietario) IMPLIES oblig(causa_generadora_probada);
\end{lstlisting}

\subsubsection{Clause: tesis341704\_posesion\_originaria}

\textit{Original possession: Must begin with cause different from derived possession}

\textbf{Intermediate Representation:}

\begin{quote}
possession as owner is obligated and not (temporal possession is obligated)
\end{quote}

\textbf{Witness Clause:}

\begin{lstlisting}[style=witnessstyle, language=Witness]
clause tesis341704_posesion_originaria = oblig(posesion_como_propietario) AND not(oblig(posesion_temporal)) IMPLIES oblig(causa_generadora_probada);
\end{lstlisting}

\newpage

\section{Thesis 343163: PRESCRIPCION POSITIVA, EXCEPCION DE (LEGISLACION DE ZACATECAS).}
\label{thesis:343163}

\subsection{Original Thesis Text}

\begin{quote}
                                                   Semanario Judicial de la Federación

                                                  Tesis

 Registro digital: 343163

 Instancia: Tercera Sala          Quinta Época                         Materia(s): Civil

                                  Fuente: Semanario Judicial de la     Tipo: Aislada
                                  Federación.
                                  Tomo CVII, página 2024


PRESCRIPCION POSITIVA, EXCEPCION DE (LEGISLACION DE ZACATECAS).


Si la autoridad responsable reconoció que la parte demandada en un juicio reivindicatorio está en
posesión, desde hace muchos años, de los bienes reclamados, no obstante lo cual, nada resolvió
respecto a la excepción de prescripción opuesta por la propia demandada, al contestar la demanda,
alegando una posesión quieta, pública, continua y de buena fe, con apoyo, entre otros preceptos, en
los artículos 1086 y 1087 del Código Civil; en tales condiciones, debe concederse el amparo para el
efecto de que dicha autoridad estudie y resuelva sobre la excepción de prescripción opuesta.


Amparo civil directo 2427/40. Torres Vda. de Valdés Petra. 26 de marzo de 1951. Mayoría de tres
votos. Ponente: Hilario Medina.




                                             Pág. 1 de 1             Fecha de impresión 01/10/2025
  https://sjf2.scjn.gob.mx/detalle/tesis/343163

\end{quote}

\subsection{Formalization}

\subsubsection{Clause: tesis343163\_principio\_fundamental}

\textit{Thesis-specific principle: Good faith possessors require just title as generating cause of possession for adverse possession to operate}

\textbf{Intermediate Representation:}

\begin{quote}
good faith exercise is obligated and possession as owner is obligated
\end{quote}

\textbf{Witness Clause:}

\begin{lstlisting}[style=witnessstyle, language=Witness]
clause tesis343163_principio_fundamental = oblig(ejercicio_buena_fe) AND oblig(posesion_como_propietario) IMPLIES oblig(causa_generadora_probada);
\end{lstlisting}

\subsubsection{Clause: tesis343163\_revelar\_causa}

\textit{Supporting logic: Must reveal and fully prove the generating cause of possession}

\textbf{Intermediate Representation:}

\begin{quote}
ejercer prescripcion positiva is obligated
\end{quote}

\textbf{Witness Clause:}

\begin{lstlisting}[style=witnessstyle, language=Witness]
clause tesis343163_revelar_causa = oblig(ejercer_prescripcion_positiva) IMPLIES oblig(causa_generadora_revelada);
\end{lstlisting}

\subsubsection{Clause: tesis343163\_determinacion\_judicial}

\textit{Judicial determination: Judge needs generating fact/act to determine possession quality and timing}

\textbf{Intermediate Representation:}

\begin{quote}
If causa generadora revelada is obligated, then (possession as owner is obligated and (good faith exercise is obligated and not (bad faith exercise is obligated)))
\end{quote}

\textbf{Witness Clause:}

\begin{lstlisting}[style=witnessstyle, language=Witness]
clause tesis343163_determinacion_judicial = oblig(causa_generadora_revelada) IMPLIES (oblig(posesion_como_propietario) AND (oblig(ejercicio_buena_fe) AND not(oblig(ejercicio_mala_fe))));
\end{lstlisting}

\subsubsection{Clause: tesis343163\_determinacion\_temporal}

\textit{Timing requirement: Must establish when prescription period begins}

\textbf{Intermediate Representation:}

\begin{quote}
If causa generadora revelada is obligated, then (posesion tiempo 5 anos is obligated or posesion tiempo 10 anos is obligated)
\end{quote}

\textbf{Witness Clause:}

\begin{lstlisting}[style=witnessstyle, language=Witness]
clause tesis343163_determinacion_temporal = oblig(causa_generadora_revelada) IMPLIES (oblig(posesion_tiempo_5_anos) OR oblig(posesion_tiempo_10_anos));
\end{lstlisting}

\newpage

\section{Thesis 343922: PRESCRIPCION POSITIVA, POSESION DE BUENA FE Y EN CONCEPTO DE PROPIETARIO,}
\label{thesis:343922}

\subsection{Original Thesis Text}

\begin{quote}
                                                   Semanario Judicial de la Federación

                                                  Tesis

 Registro digital: 343922

 Instancia: Tercera Sala          Quinta Época                         Materia(s): Civil

                                  Fuente: Semanario Judicial de la     Tipo: Aislada
                                  Federación.
                                  Tomo CIV, página 930


PRESCRIPCION POSITIVA, POSESION DE BUENA FE Y EN CONCEPTO DE PROPIETARIO,
PARA LOS EFECTOS DE LA.


La buena fe y el concepto de propietario a que alude la ley, para los efectos de la prescripción
positiva, son esencialmente subjetivos. En consecuencia, la mala fe del causante o el hecho de que
éste no hubiera tenido tal concepto, no perjudica al causahabiente a título particular, pues por ser
tales situaciones estrictamente personales, pueden faltar en el causante y existir en el
causahabiente, y la mala fe del primero no implica mala fe en el segundo.


Amparo civil directo 9095/49. Garza Alfonso y coagraviados. 26 de abril de 1950. Mayoría de cuatro
votos. Disidente: Hilario Medina. Ponente: Carlos I. Meléndez.




                                             Pág. 1 de 1             Fecha de impresión 01/10/2025
  https://sjf2.scjn.gob.mx/detalle/tesis/343922

\end{quote}

\subsection{Formalization}

\subsubsection{Clause: tesis343922\_principio\_fundamental}

\textit{Thesis-specific principle: Just title not required for adverse possession, but must reveal origin of possession and affirm possession as owner}

\textbf{Intermediate Representation:}

\begin{quote}
ejercer prescripcion positiva is obligated and possession as owner is obligated
\end{quote}

\textbf{Witness Clause:}

\begin{lstlisting}[style=witnessstyle, language=Witness]
clause tesis343922_principio_fundamental = oblig(ejercer_prescripcion_positiva) IMPLIES oblig(causa_generadora_revelada) AND oblig(posesion_como_propietario);
\end{lstlisting}

\subsubsection{Clause: tesis343922\_determinacion\_judicial}

\textit{Supporting logic: Judge needs generating fact/act to determine possession quality and timing}

\textbf{Intermediate Representation:}

\begin{quote}
If causa generadora revelada is obligated, then (possession as owner is obligated and (good faith exercise is obligated and not (bad faith exercise is obligated)))
\end{quote}

\textbf{Witness Clause:}

\begin{lstlisting}[style=witnessstyle, language=Witness]
clause tesis343922_determinacion_judicial = oblig(causa_generadora_revelada) IMPLIES (oblig(posesion_como_propietario) AND (oblig(ejercicio_buena_fe) AND not(oblig(ejercicio_mala_fe))));
\end{lstlisting}

\subsubsection{Clause: tesis343922\_tipo\_posesion}

\textit{Possession type determination: Can determine if original or derived possession}

\textbf{Intermediate Representation:}

\begin{quote}
If causa generadora revelada is obligated, then (possession as owner is obligated and not (temporal possession is obligated))
\end{quote}

\textbf{Witness Clause:}

\begin{lstlisting}[style=witnessstyle, language=Witness]
clause tesis343922_tipo_posesion = oblig(causa_generadora_revelada) IMPLIES (oblig(posesion_como_propietario) AND not(oblig(posesion_temporal)));
\end{lstlisting}

\subsubsection{Clause: tesis343922\_determinacion\_temporal}

\textit{Timing requirement: Must establish when prescription period starts counting}

\textbf{Intermediate Representation:}

\begin{quote}
If causa generadora revelada is obligated, then (posesion tiempo 5 anos is obligated or posesion tiempo 10 anos is obligated)
\end{quote}

\textbf{Witness Clause:}

\begin{lstlisting}[style=witnessstyle, language=Witness]
clause tesis343922_determinacion_temporal = oblig(causa_generadora_revelada) IMPLIES (oblig(posesion_tiempo_5_anos) OR oblig(posesion_tiempo_10_anos));
\end{lstlisting}

\newpage

\section{Thesis 345885: PRESCRIPCION POSITIVA FUNDADA EN POSESION SIN JUSTO TITULO (LEGISLACION DE}
\label{thesis:345885}

\subsection{Original Thesis Text}

\begin{quote}
                                                   Semanario Judicial de la Federación

                                                  Tesis

 Registro digital: 345885

 Instancia: Tercera Sala          Quinta Época                         Materia(s): Civil

                                  Fuente: Semanario Judicial de la     Tipo: Aislada
                                  Federación.
                                  Tomo XCVI, página 1588


PRESCRIPCION POSITIVA FUNDADA EN POSESION SIN JUSTO TITULO (LEGISLACION DE
VERACRUZ).


El artículo 1185 del Código Civil del Estado de Veracruz establece, en su fracción I, la prescripción
de diez años, para quienes posean bienes inmuebles en concepto de propietario, con justo título,
con buena fe, pacífica, continua y públicamente; y, en su fracción III, establece la prescripción de
veinte años, para quienes los posean de mala fe, si la posesión es en concepto de propietario,
pacífica, continua y públicamente. De lo dispuesto por estos preceptos, en relación con lo que
previene el artículo 842 del código citado, se advierte que la posesión por veinte años, con las
cualidades legales, origina la prescripción, aun sin justo título.


Amparo civil directo 5333/47. Márquez Pioquinto Salvador. 10 de junio de 1948. Unanimidad de
cuatro votos. Ausente: Carlos I. Meléndez. La publicación no menciona el nombre del ponente.




                                             Pág. 1 de 1             Fecha de impresión 01/10/2025
  https://sjf2.scjn.gob.mx/detalle/tesis/345885

\end{quote}

\subsection{Formalization}

\subsubsection{Clause: tesis345885\_principio\_fundamental}

\textit{Thesis-specific principle: When defendant admits possession without title for more than 10 years with required qualities, is bad faith possessor and property is eligible for prescription}

\textbf{Intermediate Representation:}

\begin{quote}
If (posesion\_como\_propietario AND posesion\_pacifico\_continuo\_publico AND posesion\_tiempo\_10\_anos AND not(causa\_generadora\_probada)), then (ejercicio\_mala\_fe AND oblig(ejercer\_prescripcion\_positiva))
\end{quote}

\textbf{Witness Clause:}

\begin{lstlisting}[style=witnessstyle, language=Witness]
clause tesis345885_principio_fundamental = (posesion_como_propietario AND posesion_pacifico_continuo_publico AND posesion_tiempo_10_anos AND not(causa_generadora_probada)) IMPLIES (ejercicio_mala_fe AND oblig(ejercer_prescripcion_positiva));
\end{lstlisting}

\subsubsection{Clause: tesis345885\_no\_prueba\_causa}

\textit{Supporting logic: Does not need to prove generating cause of possession in trial when no title exists}

\textbf{Intermediate Representation:}

\begin{quote}
(not(causa\_generadora\_probada) AND posesion\_tiempo\_10\_anos) IMPLIES oblig(ejercer\_prescripcion\_positiva)
\end{quote}

\textbf{Witness Clause:}

\begin{lstlisting}[style=witnessstyle, language=Witness]
clause tesis345885_no_prueba_causa = (not(causa_generadora_probada) AND posesion_tiempo_10_anos) IMPLIES oblig(ejercer_prescripcion_positiva);
\end{lstlisting}

\subsubsection{Clause: tesis345885\_definicion\_titulo}

\textit{Title definition: In state law, title is understood as the generating cause of possession}

\textbf{Intermediate Representation:}

\begin{quote}
not (causa\_generadora\_probada) IMPLIES not(titularidad\_registral)
\end{quote}

\textbf{Witness Clause:}

\begin{lstlisting}[style=witnessstyle, language=Witness]
clause tesis345885_definicion_titulo = not(causa_generadora_probada) IMPLIES not(titularidad_registral);
\end{lstlisting}

\subsubsection{Clause: tesis345885\_apto\_prescripcion}

\textit{Prescription eligibility: If 10 years have passed and other elements are met, property is eligible for prescription}

\textbf{Intermediate Representation:}

\begin{quote}
(posesion\_tiempo\_10\_anos AND posesion\_como\_propietario AND posesion\_pacifico\_continuo\_publico) IMPLIES oblig(ejercer\_prescripcion\_positiva)
\end{quote}

\textbf{Witness Clause:}

\begin{lstlisting}[style=witnessstyle, language=Witness]
clause tesis345885_apto_prescripcion = (posesion_tiempo_10_anos AND posesion_como_propietario AND posesion_pacifico_continuo_publico) IMPLIES oblig(ejercer_prescripcion_positiva);
\end{lstlisting}

\newpage

\section{Thesis 347316: PRESCRIPCION POSITIVA, BUENA O MALA FE EN LA POSESION EN QUE SE FUNDA (LEGISLACION DEL ESTADO DE MEXICO).}
\label{thesis:347316}

\subsection{Original Thesis Text}

\begin{quote}
                                                   Semanario Judicial de la Federación

                                                  Tesis

 Registro digital: 347316

 Instancia: Tercera Sala          Quinta Época                          Materia(s): Civil

                                  Fuente: Semanario Judicial de la      Tipo: Aislada
                                  Federación.
                                  Tomo XCI, página 2512


PRESCRIPCION POSITIVA, BUENA O MALA FE EN LA POSESION EN QUE SE FUNDA
(LEGISLACION DEL ESTADO DE MEXICO).


El decreto de nueve de agosto de mil novecientos treinta y siete, que declaró vigente en el Estado
de México el Código Civil para el Distrito y Territorios Federales de mil novecientos treinta y dos,
modificó el artículo 1152 en el sentido de que, de conformidad con el artículo 200 de la Constitución
del Estado, los bienes inmuebles se prescribían en veinte años; pero el decreto de veintiuno de
diciembre de mil novecientos treinta y ocho, seguramente por virtud de la reforma al precepto
constitucional, derogó el citado artículo del Código Civil, tal como había sido reformado por el
decreto de nueve de agosto de mil novecientos treinta y siete, quedando en definitiva, como rige en
el Distrito Federal. Ahora bien, tanto por esta razón de que el artículo 1152 obliga actualmente en el
Estado de México, tal como se halla redactado en el Código Civil vigente en el Distrito Federal,
como porque la reforma del artículo 200 de la Constitución local sólo fija término máximo para la
prescripción de la propiedad raíz, es completamente inexacto que por disposición expresa del
mencionado artículo constitucional, no se requiera en la posesión el elemento de buena o de mala
fe, a que se contraen las fracciones I y III del artículo 1152 de la ley civil y que solamente debe
exigirse que la posesión sea en concepto de propietario, pacífica, contínua y pública; debiendo
agregarse que son precisamente esos elementos de buena o mala fe, los que constituyen la base
para computar el término en que se realiza la prescripción, que es de cinco años cuando concurre el
primero y diez cuando se está en presencia del último.


Amparo civil directo 3240/45. Terrón Dolores F. y coagraviado. 19 de marzo de 1947. Unanimidad
de cinco votos. La publicación no menciona el nombre del ponente.




                                             Pág. 1 de 1             Fecha de impresión 01/10/2025
  https://sjf2.scjn.gob.mx/detalle/tesis/347316

\end{quote}

\subsection{Formalization}

\subsubsection{Clause: tesis347316\_principio\_fundamental}

\textit{Thesis-specific principle: Good faith requires sufficient title/generating cause, ignorance of title defects, and founded belief of ownership}

\textbf{Intermediate Representation:}

\begin{quote}
good faith exercise is obligated and not (conocimiento vicios is obligated)
\end{quote}

\textbf{Witness Clause:}

\begin{lstlisting}[style=witnessstyle, language=Witness]
clause tesis347316_principio_fundamental = oblig(ejercicio_buena_fe) IMPLIES oblig(causa_generadora_probada) AND not(oblig(conocimiento_vicios));
\end{lstlisting}

\subsubsection{Clause: tesis347316\_mala\_fe}

\textit{Supporting logic: Bad faith possessor has no title/generating cause OR knows title defects that prevent rightful possession}

\textbf{Intermediate Representation:}

\begin{quote}
If bad faith exercise is obligated, then (not (proven generating cause is obligated) or conocimiento vicios is obligated)
\end{quote}

\textbf{Witness Clause:}

\begin{lstlisting}[style=witnessstyle, language=Witness]
clause tesis347316_mala_fe = oblig(ejercicio_mala_fe) IMPLIES (not(oblig(causa_generadora_probada)) OR oblig(conocimiento_vicios));
\end{lstlisting}

\subsubsection{Clause: tesis347316\_requisitos\_buena\_fe}

\textit{Good faith requirements: Must have sufficient title and ignorance of defects}

\textbf{Intermediate Representation:}

\begin{quote}
good faith exercise is obligated
\end{quote}

\textbf{Witness Clause:}

\begin{lstlisting}[style=witnessstyle, language=Witness]
clause tesis347316_requisitos_buena_fe = oblig(ejercicio_buena_fe) IMPLIES oblig(causa_generadora_probada);
\end{lstlisting}

\subsubsection{Clause: tesis347316\_condiciones\_mala\_fe}

\textit{Bad faith conditions: No title or knowledge of defects}

\textbf{Intermediate Representation:}

\begin{quote}
(not(oblig(causa\_generadora\_probada)) OR oblig(conocimiento\_vicios)) IMPLIES oblig(ejercicio\_mala\_fe)
\end{quote}

\textbf{Witness Clause:}

\begin{lstlisting}[style=witnessstyle, language=Witness]
clause tesis347316_condiciones_mala_fe = (not(oblig(causa_generadora_probada)) OR oblig(conocimiento_vicios)) IMPLIES oblig(ejercicio_mala_fe);
\end{lstlisting}

\newpage

\section{Thesis 347475: PRESCRIPCION POSITIVA, JUSTO TITULO EN LA POSESION, PARA LOS EFECTOS DE LA (LEGISLACION DE SINALOA).}
\label{thesis:347475}

\subsection{Original Thesis Text}

\begin{quote}
                                                   Semanario Judicial de la Federación

                                                  Tesis

 Registro digital: 347475

 Instancia: Tercera Sala          Quinta Época                          Materia(s): Civil

                                  Fuente: Semanario Judicial de la      Tipo: Aislada
                                  Federación.
                                  Tomo XC, página 1211


PRESCRIPCION POSITIVA, JUSTO TITULO EN LA POSESION, PARA LOS EFECTOS DE LA
(LEGISLACION DE SINALOA).


Conforme el artículo 1079 del Código Civil del Estado de Sinaloa, anterior al vigente, la posesión
necesaria para prescribir debía ser fundada en justo título y de buena fe. Ahora bien, debe
estimarse que posee con justo título quien adquirió la posesión mediante un contrato que, por su
naturaleza, transmite el dominio, como es el de compraventa, sin que importe que éste haya sido
declarado nulo, pues esto sólo significa que carecía de eficacia para transmitir la propiedad, pero no
lo afecta como causa generadora de la posesión, que es lo que la ley requiere como título para
considerar que se posee de buena fe.


Amparo civil directo 4957/45. Espinosa Nicolás. 28 de octubre de 1946. Unanimidad de cuatro
votos. Ausente: Hilario Medina. La publicación no menciona el nombre del ponente.

Quinta Epoca:

Tomo XC, página 2981. Indice Alfabético. Amparo directo 4511/45. Hernández Venancio. 28 de
octubre de 1946. Unanimidad de cuatro votos. Ausente: Hilario Medina. Ponente: Agustín Mercado
Alarcón.




                                             Pág. 1 de 1             Fecha de impresión 01/10/2025
  https://sjf2.scjn.gob.mx/detalle/tesis/347475

\end{quote}

\subsection{Formalization}

\subsubsection{Clause: tesis347475\_principio\_fundamental}

\textit{Thesis-specific principle: Good faith possessors require just title as generating cause of possession for adverse possession to operate}

\textbf{Intermediate Representation:}

\begin{quote}
good faith exercise is obligated and possession as owner is obligated
\end{quote}

\textbf{Witness Clause:}

\begin{lstlisting}[style=witnessstyle, language=Witness]
clause tesis347475_principio_fundamental = oblig(ejercicio_buena_fe) AND oblig(posesion_como_propietario) IMPLIES oblig(causa_generadora_probada);
\end{lstlisting}

\subsubsection{Clause: tesis347475\_revelar\_causa}

\textit{Supporting logic: Must reveal and fully prove the generating cause of possession}

\textbf{Intermediate Representation:}

\begin{quote}
ejercer prescripcion positiva is obligated
\end{quote}

\textbf{Witness Clause:}

\begin{lstlisting}[style=witnessstyle, language=Witness]
clause tesis347475_revelar_causa = oblig(ejercer_prescripcion_positiva) IMPLIES oblig(causa_generadora_revelada);
\end{lstlisting}

\subsubsection{Clause: tesis347475\_determinacion\_judicial}

\textit{Judicial determination: Judge needs generating fact/act to determine possession quality and timing}

\textbf{Intermediate Representation:}

\begin{quote}
If causa generadora revelada is obligated, then (possession as owner is obligated and (good faith exercise is obligated and not (bad faith exercise is obligated)))
\end{quote}

\textbf{Witness Clause:}

\begin{lstlisting}[style=witnessstyle, language=Witness]
clause tesis347475_determinacion_judicial = oblig(causa_generadora_revelada) IMPLIES (oblig(posesion_como_propietario) AND (oblig(ejercicio_buena_fe) AND not(oblig(ejercicio_mala_fe))));
\end{lstlisting}

\subsubsection{Clause: tesis347475\_determinacion\_temporal}

\textit{Timing requirement: Must establish when prescription period begins}

\textbf{Intermediate Representation:}

\begin{quote}
If causa generadora revelada is obligated, then (posesion tiempo 5 anos is obligated or posesion tiempo 10 anos is obligated)
\end{quote}

\textbf{Witness Clause:}

\begin{lstlisting}[style=witnessstyle, language=Witness]
clause tesis347475_determinacion_temporal = oblig(causa_generadora_revelada) IMPLIES (oblig(posesion_tiempo_5_anos) OR oblig(posesion_tiempo_10_anos));
\end{lstlisting}

\newpage

\section{Thesis 347751: PRESCRIPCION POSITIVA FUNDADA EN POSESION DE BUENA FE (LEGISLACION DE SAN}
\label{thesis:347751}

\subsection{Original Thesis Text}

\begin{quote}
                                                   Semanario Judicial de la Federación

                                                  Tesis

 Registro digital: 347751

 Instancia: Tercera Sala          Quinta Época                          Materia(s): Civil

                                  Fuente: Semanario Judicial de la      Tipo: Aislada
                                  Federación.
                                  Tomo LXXXIX, página 1948


PRESCRIPCION POSITIVA FUNDADA EN POSESION DE BUENA FE (LEGISLACION DE SAN
LUIS POTOSI).


El artículo 830 del Código Civil del Estado de San Luis Potosí, define al poseedor de buena fe como
aquel que tiene o fundadamente cree tener título bastante para transferir el dominio; y no puede
estimarse que haya tenido esa creencia, quien admite que en la operación de compraventa del
predio que pretende haber adquirido por prescripción, recibió sólo la posesión y no así una titulación
perfecta, porque el vendedor también carecía de ella.


Amparo civil directo 1922/45. Céspedes Lauro. 21 de agosto de 1946. Unanimidad de cuatro votos.
Ausente: Vicente Santos Guajardo. Ponente: Agustín Mercado Alarcón.




                                             Pág. 1 de 1             Fecha de impresión 01/10/2025
  https://sjf2.scjn.gob.mx/detalle/tesis/347751

\end{quote}

\subsection{Formalization}

\subsubsection{Clause: tesis347751\_principio\_fundamental}

\textit{Additional actions for this thesis Additional assets for proof requirements Faith exclusivity: Adverse possession is either good faith XOR bad faith Thesis-specific principle: Good faith adverse possession requires proving authenticity, date, and diligence}

\textbf{Intermediate Representation:}

\begin{quote}
If good faith exercise is obligated and adverse possession claim is claimed, then (autenticidad titulo is obligated and fecha fehaciente is obligated and diligencia adquirente is obligated)
\end{quote}

\textbf{Witness Clause:}

\begin{lstlisting}[style=witnessstyle, language=Witness]
clause tesis347751_principio_fundamental = oblig(ejercicio_buena_fe) AND claim(demanda_prescripcion) IMPLIES (oblig(autenticidad_titulo) AND oblig(fecha_fehaciente) AND oblig(diligencia_adquirente));
\end{lstlisting}

\subsubsection{Clause: tesis347751\_prueba\_autenticidad}

\textit{Proof requirements: Must prove title authenticity and date reliably}

\textbf{Intermediate Representation:}

\begin{quote}
If proven generating cause is obligated, then (autenticidad titulo is obligated and fecha fehaciente is obligated)
\end{quote}

\textbf{Witness Clause:}

\begin{lstlisting}[style=witnessstyle, language=Witness]
clause tesis347751_prueba_autenticidad = oblig(causa_generadora_probada) IMPLIES (oblig(autenticidad_titulo) AND oblig(fecha_fehaciente));
\end{lstlisting}

\subsubsection{Clause: tesis347751\_diligencia\_buena\_fe}

\textit{Diligence requirement: Good faith requires demonstrating due diligence in acquisition}

\textbf{Intermediate Representation:}

\begin{quote}
good faith exercise is obligated
\end{quote}

\textbf{Witness Clause:}

\begin{lstlisting}[style=witnessstyle, language=Witness]
clause tesis347751_diligencia_buena_fe = oblig(ejercicio_buena_fe) IMPLIES oblig(diligencia_adquirente);
\end{lstlisting}

\subsubsection{Clause: tesis347751\_prueba\_pagos}

\textit{Payment proof: Onerous contracts require proving payments made}

\textbf{Intermediate Representation:}

\begin{quote}
proven generating cause is obligated and contrato compraventa is obligated
\end{quote}

\textbf{Witness Clause:}

\begin{lstlisting}[style=witnessstyle, language=Witness]
clause tesis347751_prueba_pagos = oblig(causa_generadora_probada) AND oblig(contrato_compraventa) IMPLIES oblig(pagos_precio);
\end{lstlisting}

\subsubsection{Clause: tesis347751\_carga\_prueba}

\textit{Burden of proof: Claimant must prove generating cause with reliable evidence}

\textbf{Intermediate Representation:}

\begin{quote}
adverse possession claim is claimed
\end{quote}

\textbf{Witness Clause:}

\begin{lstlisting}[style=witnessstyle, language=Witness]
clause tesis347751_carga_prueba = claim(demanda_prescripcion) IMPLIES oblig(causa_generadora_revelada);
\end{lstlisting}

\subsubsection{Clause: tesis347751\_falta\_diligencia}

\textit{Insufficient diligence consequence: Lack of diligence means bad faith (10-year period)}

\textbf{Intermediate Representation:}

\begin{quote}
not (diligencia adquirente is obligated)
\end{quote}

\textbf{Witness Clause:}

\begin{lstlisting}[style=witnessstyle, language=Witness]
clause tesis347751_falta_diligencia = not(oblig(diligencia_adquirente)) IMPLIES oblig(ejercicio_mala_fe);
\end{lstlisting}

\newpage

\section{Thesis 348113: THESIS 348113}
\label{thesis:348113}

\subsection{Original Thesis Text}

\begin{quote}
                                                   Semanario Judicial de la Federación

                                                  Tesis

 Registro digital: 348113

 Instancia: Tercera Sala            Quinta Época                            Materia(s): Civil

                                    Fuente: Semanario Judicial de la        Tipo: Aislada
                                    Federación.
                                    Tomo LXXXVIII, página 2541


PRESCRIPCION POSITIVA.


Si la causa de la posesión, por el actor, en un juicio reivindicatorio, del predio, objeto de dicho juicio,
independientemente de su eficacia actual como título de propiedad perfecto, procedente del dueño
de la cosa, es, en derecho, de los que bastan para trasmitir el dominio, como la compraventa, la
permuta, etcétera, y además, existe por expresa disposición de la ley la presunción de buena fe, y la
demandada no contradijo ni desvirtuó semejante presunción, que admite prueba en contrario, debe
concluirse que la posesión por el tiempo requerido por la ley y no interrumpida, fue suficiente para
que el actor hubiera adquirido por usucapión, el predio de referencia.


Amparo civil directo 5227/42. Montes Isidro. 19 de junio de 1946. Unanimidad de cuatro votos.
Ausente: Carlos I. Meléndez. La publicación no menciona el nombre del ponente.




                                               Pág. 1 de 1              Fecha de impresión 01/10/2025
  https://sjf2.scjn.gob.mx/detalle/tesis/348113

\end{quote}

\subsection{Formalization}

\subsubsection{Clause: tesis348113\_principio\_fundamental}

\textit{Thesis-specific principle: Adverse possession can be invoked by action or exception, but requires same proof elements}

\textbf{Intermediate Representation:}

\begin{quote}
possession as owner is obligated and (good faith exercise is obligated or bad faith exercise is obligated) and (oblig(posesion\_tiempo\_5\_anos) OR oblig(posesion\_tiempo\_10\_anos)) IMPLIES oblig(ejercer\_prescripcion\_positiva)
\end{quote}

\textbf{Witness Clause:}

\begin{lstlisting}[style=witnessstyle, language=Witness]
clause tesis348113_principio_fundamental = oblig(posesion_como_propietario) AND (oblig(ejercicio_buena_fe) OR oblig(ejercicio_mala_fe)) AND (oblig(posesion_tiempo_5_anos) OR oblig(posesion_tiempo_10_anos)) IMPLIES oblig(ejercer_prescripcion_positiva);
\end{lstlisting}

\subsubsection{Clause: tesis348113\_igualdad\_prueba}

\textit{Supporting logic: Both action and exception require same proof elements regardless of procedural route}

\textbf{Intermediate Representation:}

\begin{quote}
ejercer prescripcion positiva is obligated
\end{quote}

\textbf{Witness Clause:}

\begin{lstlisting}[style=witnessstyle, language=Witness]
clause tesis348113_igualdad_prueba = oblig(ejercer_prescripcion_positiva) IMPLIES oblig(posesion_como_propietario);
\end{lstlisting}

\subsubsection{Clause: tesis348113\_plazos\_temporales}

\textit{Time requirements: 5 years for good faith, 10 years for bad faith}

\textbf{Intermediate Representation:}

\begin{quote}
good faith exercise is obligated
\end{quote}

\textbf{Witness Clause:}

\begin{lstlisting}[style=witnessstyle, language=Witness]
clause tesis348113_plazos_temporales = oblig(ejercicio_buena_fe) IMPLIES oblig(posesion_tiempo_5_anos);
\end{lstlisting}

\subsubsection{Clause: tesis348113\_plazos\_mala\_fe}

\textbf{Intermediate Representation:}

\begin{quote}
bad faith exercise is obligated
\end{quote}

\textbf{Witness Clause:}

\begin{lstlisting}[style=witnessstyle, language=Witness]
clause tesis348113_plazos_mala_fe = oblig(ejercicio_mala_fe) IMPLIES oblig(posesion_tiempo_10_anos);
\end{lstlisting}

\subsubsection{Clause: tesis348113\_objetivo\_registral}

\textit{Target requirement: Action must be directed against registered owner}

\textbf{Intermediate Representation:}

\begin{quote}
ejercer prescripcion positiva is obligated
\end{quote}

\textbf{Witness Clause:}

\begin{lstlisting}[style=witnessstyle, language=Witness]
clause tesis348113_objetivo_registral = oblig(ejercer_prescripcion_positiva) IMPLIES oblig(titularidad_registral);
\end{lstlisting}

\newpage

\section{Thesis 348467: PRESCRIPCION POSITIVA, RESPECTO DE INMUEBLES ADQUIRIDOS ANTERIORMENTE POR PERSONAS INTERPOSITAS DEL CLERO.}
\label{thesis:348467}

\subsection{Original Thesis Text}

\begin{quote}
                                                   Semanario Judicial de la Federación

                                                  Tesis

 Registro digital: 348467

 Instancia: Tercera Sala           Quinta Época                          Materia(s): Civil

                                   Fuente: Semanario Judicial de la      Tipo: Aislada
                                   Federación.
                                   Tomo LXXXVI, página 727


PRESCRIPCION POSITIVA, RESPECTO DE INMUEBLES ADQUIRIDOS ANTERIORMENTE POR
PERSONAS INTERPOSITAS DEL CLERO.


La posesión en concepto de propietario constituye un hecho jurídico y no un acto jurídico. El título
de posesión que fundamenta, se considera bastante para adquirir y transmitir el dominio, es
subjetivamente válido, y esta situación subjetiva, que presume la ley y que el opositor debe destruir
con pruebas directas, no es un efecto jurídico que se derive de la validez objetiva del título de
propiedad por virtud del cual adquiere un poseedor determinado, sino que es una simple cuestión
de hecho, esencialmente interna. De considerarse que sólo aquellos que poseen con título
objetivamente válido, es decir, perfecto con arreglo a derecho, pueden adquirir por prescripción,
ésta resultaría absolutamente inútil, ya que se partiría de la hipótesis de que el título es perfecto.
Por consiguientes, la posesión en concepto de propietario, puede existir como situación netamente
subjetiva, para adquirir por prescripción el dominio de un inmueble, aun cuando uno de los títulos
anteriores o de los actos jurídicos que le sirvieron de antecedentes, hayan sido nulos o inexistentes,
por haber intervenido en ellos personas interpósitas de la iglesia Católica. Por otra parte, siendo la
buena fue y el concepto de propietario o animus domini a que el alude la ley, para los efectos de la
prescripción positiva, esencialmente subjetivos, la mala fe del causante o el hecho de que no
hubiera tenido tal animus, no perjudica al causahabiente a título particular; de manera que la
circunstancia de que quien era interpósita persona del clero, no haya tenido el animus domini, no
impide ni puede impedir que un posterior adquirente lo haya tenido por virtud de la compra del
inmueble que hizo a tercera persona, y del mismo modo, la mala fe del causante, no implica mala fe
en el causahabiente a título particular.


Amparo civil directo 4648/35. Estrella viuda de Fregoso Plácida. 24 de octubre de 1945. Unanimidad
de cuatro votos. El Ministro Hilario Medina, no intervino en este juicio por encontrarse disfrutando de
licencia. Ponente: Vicente Santos Guajardo.




                                              Pág. 1 de 2             Fecha de impresión 01/10/2025
  https://sjf2.scjn.gob.mx/detalle/tesis/348467
                                                Semanario Judicial de la Federación




                                           Pág. 2 de 2       Fecha de impresión 01/10/2025
https://sjf2.scjn.gob.mx/detalle/tesis/348467

\end{quote}

\subsection{Formalization}

\subsubsection{Clause: tesis348467\_principio\_fundamental}

\textit{Thesis-specific principle: Good faith possessors require just title as generating cause of possession for adverse possession to operate}

\textbf{Intermediate Representation:}

\begin{quote}
good faith exercise is obligated and possession as owner is obligated
\end{quote}

\textbf{Witness Clause:}

\begin{lstlisting}[style=witnessstyle, language=Witness]
clause tesis348467_principio_fundamental = oblig(ejercicio_buena_fe) AND oblig(posesion_como_propietario) IMPLIES oblig(causa_generadora_probada);
\end{lstlisting}

\subsubsection{Clause: tesis348467\_revelar\_causa}

\textit{Supporting logic: Must reveal and fully prove the generating cause of possession}

\textbf{Intermediate Representation:}

\begin{quote}
ejercer prescripcion positiva is obligated
\end{quote}

\textbf{Witness Clause:}

\begin{lstlisting}[style=witnessstyle, language=Witness]
clause tesis348467_revelar_causa = oblig(ejercer_prescripcion_positiva) IMPLIES oblig(causa_generadora_revelada);
\end{lstlisting}

\subsubsection{Clause: tesis348467\_determinacion\_judicial}

\textit{Judicial determination: Judge needs generating fact/act to determine possession quality and timing}

\textbf{Intermediate Representation:}

\begin{quote}
If causa generadora revelada is obligated, then (possession as owner is obligated and (good faith exercise is obligated and not (bad faith exercise is obligated)))
\end{quote}

\textbf{Witness Clause:}

\begin{lstlisting}[style=witnessstyle, language=Witness]
clause tesis348467_determinacion_judicial = oblig(causa_generadora_revelada) IMPLIES (oblig(posesion_como_propietario) AND (oblig(ejercicio_buena_fe) AND not(oblig(ejercicio_mala_fe))));
\end{lstlisting}

\subsubsection{Clause: tesis348467\_determinacion\_temporal}

\textit{Timing requirement: Must establish when prescription period begins}

\textbf{Intermediate Representation:}

\begin{quote}
If causa generadora revelada is obligated, then (posesion tiempo 5 anos is obligated or posesion tiempo 10 anos is obligated)
\end{quote}

\textbf{Witness Clause:}

\begin{lstlisting}[style=witnessstyle, language=Witness]
clause tesis348467_determinacion_temporal = oblig(causa_generadora_revelada) IMPLIES (oblig(posesion_tiempo_5_anos) OR oblig(posesion_tiempo_10_anos));
\end{lstlisting}

\newpage

\section{Thesis 348821: THESIS 348821}
\label{thesis:348821}

\subsection{Original Thesis Text}

\begin{quote}
                                                   Semanario Judicial de la Federación

                                                  Tesis

 Registro digital: 348821

 Instancia: Tercera Sala          Quinta Época                         Materia(s): Civil

                                  Fuente: Semanario Judicial de la     Tipo: Aislada
                                  Federación.
                                  Tomo LXXXIV, página 45


PRESCRIPCION POSITIVA.


Tanto la prescripción de buena fe, como la de mala fe, requieren que el prescribiente haya poseído
en concepto de propietario; de manera que si esta condición no quedó probada con la prueba
testimonial rendida para acreditar la posesión, y por tal motivo la responsable omitió el estudio de
los demás requisitos de la prescripción adquisitiva y absolvió a la parte demandada, no puede
considerarse que haya incurrido en violación de garantías individuales.


Amparo civil directo 1754/41. Ponce Adela. 2 de abril de 1945. Unanimidad de cuatro votos.
Ausente: Hilario Medina. La publicación no menciona el nombre del ponente.




                                             Pág. 1 de 1             Fecha de impresión 01/10/2025
  https://sjf2.scjn.gob.mx/detalle/tesis/348821

\end{quote}

\subsection{Formalization}

\subsubsection{Clause: tesis348821\_principio\_fundamental}

\textit{Thesis-specific principle: When defendant admits possession without title for more than 10 years with required qualities, is bad faith possessor and property is eligible for prescription}

\textbf{Intermediate Representation:}

\begin{quote}
If (posesion\_como\_propietario AND posesion\_pacifico\_continuo\_publico AND posesion\_tiempo\_10\_anos AND not(causa\_generadora\_probada)), then (ejercicio\_mala\_fe AND oblig(ejercer\_prescripcion\_positiva))
\end{quote}

\textbf{Witness Clause:}

\begin{lstlisting}[style=witnessstyle, language=Witness]
clause tesis348821_principio_fundamental = (posesion_como_propietario AND posesion_pacifico_continuo_publico AND posesion_tiempo_10_anos AND not(causa_generadora_probada)) IMPLIES (ejercicio_mala_fe AND oblig(ejercer_prescripcion_positiva));
\end{lstlisting}

\subsubsection{Clause: tesis348821\_no\_prueba\_causa}

\textit{Supporting logic: Does not need to prove generating cause of possession in trial when no title exists}

\textbf{Intermediate Representation:}

\begin{quote}
(not(causa\_generadora\_probada) AND posesion\_tiempo\_10\_anos) IMPLIES oblig(ejercer\_prescripcion\_positiva)
\end{quote}

\textbf{Witness Clause:}

\begin{lstlisting}[style=witnessstyle, language=Witness]
clause tesis348821_no_prueba_causa = (not(causa_generadora_probada) AND posesion_tiempo_10_anos) IMPLIES oblig(ejercer_prescripcion_positiva);
\end{lstlisting}

\subsubsection{Clause: tesis348821\_definicion\_titulo}

\textit{Title definition: In state law, title is understood as the generating cause of possession}

\textbf{Intermediate Representation:}

\begin{quote}
not (causa\_generadora\_probada) IMPLIES not(titularidad\_registral)
\end{quote}

\textbf{Witness Clause:}

\begin{lstlisting}[style=witnessstyle, language=Witness]
clause tesis348821_definicion_titulo = not(causa_generadora_probada) IMPLIES not(titularidad_registral);
\end{lstlisting}

\subsubsection{Clause: tesis348821\_apto\_prescripcion}

\textit{Prescription eligibility: If 10 years have passed and other elements are met, property is eligible for prescription}

\textbf{Intermediate Representation:}

\begin{quote}
(posesion\_tiempo\_10\_anos AND posesion\_como\_propietario AND posesion\_pacifico\_continuo\_publico) IMPLIES oblig(ejercer\_prescripcion\_positiva)
\end{quote}

\textbf{Witness Clause:}

\begin{lstlisting}[style=witnessstyle, language=Witness]
clause tesis348821_apto_prescripcion = (posesion_tiempo_10_anos AND posesion_como_propietario AND posesion_pacifico_continuo_publico) IMPLIES oblig(ejercer_prescripcion_positiva);
\end{lstlisting}

\newpage

\section{Thesis 348860: PRESCRIPCION POSITIVA (LEGISLACION DEL ESTADO DE MEXICO).}
\label{thesis:348860}

\subsection{Original Thesis Text}

\begin{quote}
                                                   Semanario Judicial de la Federación

                                                  Tesis

 Registro digital: 348860

 Instancia: Tercera Sala          Quinta Época                           Materia(s): Civil

                                  Fuente: Semanario Judicial de la       Tipo: Aislada
                                  Federación.
                                  Tomo LXXXIV, página 554


PRESCRIPCION POSITIVA (LEGISLACION DEL ESTADO DE MEXICO).


El artículo 200 de la Constitución Política del Estado de México, en su texto original disponía: "En el
Estado, todas las propiedades raíces, o urbanas, prescribirán por la sola posesión de veinte años,
con título o sin él y con buena fe o sin ella.". Este precepto fue reformado por decreto de veintitrés
de octubre de mil novecientos treinta y cinco, en los siguientes términos: "Se fija como término
máximo para la prescripción de la propiedad raíz en el Estado, veinte años.". La modificación
sustancial del precepto consiste en que, mientras el primero requería el transcurso de veinte años
para que se operara la prescripción, el segundo sólo fijó el término máximo, para que se tuviera por
efectuada, dejando así que en la legislación común se señalaran los términos para la prescripción,
sin traspasar el máximo aludido, y en la cantidad citada, fue adoptado el Código Civil para el Distrito
y Territorios Federales, de treinta de agosto de mil novecientos veintiocho, por decreto de nueve de
agosto de mil novecientos treinta y siete, con las adiciones y modificaciones que se especifican en
el mismo decreto, entre los cuales se hace referencia al artículo 1152, que está concebido en los
siguientes términos: "De conformidad con lo prescrito por el artículo 200 de la Constitución particular
del Estado, los bienes inmuebles en el Estado se prescriben en veinte años.".


Amparo civil directo 5323/42. Cruz Pablo. 12 de abril de 1945. Unanimidad de cuatro votos.
Ausente: Emilio Pardo Aspe. Relator: Hilario Medina.




                                             Pág. 1 de 1              Fecha de impresión 01/10/2025
  https://sjf2.scjn.gob.mx/detalle/tesis/348860

\end{quote}

\subsection{Formalization}

\subsubsection{Clause: tesis348860\_principio\_fundamental}

\textit{Thesis-specific principle: Good faith requires sufficient title/generating cause, ignorance of title defects, and founded belief of ownership}

\textbf{Intermediate Representation:}

\begin{quote}
good faith exercise is obligated and not (conocimiento vicios is obligated)
\end{quote}

\textbf{Witness Clause:}

\begin{lstlisting}[style=witnessstyle, language=Witness]
clause tesis348860_principio_fundamental = oblig(ejercicio_buena_fe) IMPLIES oblig(causa_generadora_probada) AND not(oblig(conocimiento_vicios));
\end{lstlisting}

\subsubsection{Clause: tesis348860\_mala\_fe}

\textit{Supporting logic: Bad faith possessor has no title/generating cause OR knows title defects that prevent rightful possession}

\textbf{Intermediate Representation:}

\begin{quote}
If bad faith exercise is obligated, then (not (proven generating cause is obligated) or conocimiento vicios is obligated)
\end{quote}

\textbf{Witness Clause:}

\begin{lstlisting}[style=witnessstyle, language=Witness]
clause tesis348860_mala_fe = oblig(ejercicio_mala_fe) IMPLIES (not(oblig(causa_generadora_probada)) OR oblig(conocimiento_vicios));
\end{lstlisting}

\subsubsection{Clause: tesis348860\_requisitos\_buena\_fe}

\textit{Good faith requirements: Must have sufficient title and ignorance of defects}

\textbf{Intermediate Representation:}

\begin{quote}
good faith exercise is obligated
\end{quote}

\textbf{Witness Clause:}

\begin{lstlisting}[style=witnessstyle, language=Witness]
clause tesis348860_requisitos_buena_fe = oblig(ejercicio_buena_fe) IMPLIES oblig(causa_generadora_probada);
\end{lstlisting}

\subsubsection{Clause: tesis348860\_condiciones\_mala\_fe}

\textit{Bad faith conditions: No title or knowledge of defects}

\textbf{Intermediate Representation:}

\begin{quote}
(not(oblig(causa\_generadora\_probada)) OR oblig(conocimiento\_vicios)) IMPLIES oblig(ejercicio\_mala\_fe)
\end{quote}

\textbf{Witness Clause:}

\begin{lstlisting}[style=witnessstyle, language=Witness]
clause tesis348860_condiciones_mala_fe = (not(oblig(causa_generadora_probada)) OR oblig(conocimiento_vicios)) IMPLIES oblig(ejercicio_mala_fe);
\end{lstlisting}

\newpage

\section{Thesis 349370: PRESCRIPCION POSITIVA FUNDADA EN POSESION CON TITULO (LEGISLACION DE}
\label{thesis:349370}

\subsection{Original Thesis Text}

\begin{quote}
                                                   Semanario Judicial de la Federación

                                                  Tesis

 Registro digital: 349370

 Instancia: Tercera Sala          Quinta Época                         Materia(s): Civil

                                  Fuente: Semanario Judicial de la     Tipo: Aislada
                                  Federación.
                                  Tomo CII, página 146


PRESCRIPCION POSITIVA FUNDADA EN POSESION CON TITULO (LEGISLACION DE
VERACRUZ).


Si el actor en un juicio reivindicatorio, reconoce que el demandado tenía un título cuya nulidad se
declaró por sentencia, debe estimarse que para los efectos de la prescripción, dicho título es
bastante para darle al poseedor el concepto de propietario, en los términos que regulan este
elemento esencial los artículo 827, 834, 862 y 1184 del Código Civil. No es exacto que el título de
dueño o concepto de propietario que se necesita para prescribir, deba ser un título de dominio
perfecto, pues de estimarse así, sería inútil la prescripción, ya que el propietario con título de
dominio indiscutible no tendría que recurrir a otro medio para adquirir la propiedad que ya tiene.
Ahora bien, la circunstancia de que se haya declarado la nulidad del título del demandado, no trae
como consecuencia que se le deba considerar poseedor de mala fe, pues en términos generales, la
buena fe se presume y toca a quien invoca la mal fe, demostrar que el poseedor carecía de título o
conocía los vicios del mismo, de acuerdo con lo que dispone los artículos 842 a 844 del Código
Civil. En consecuencia, si el demandado tenía un título, como la ley presume la buena fe, debe
estimarse que si el acreedor no probó que aquél conocía los vicios de su título, su buena fe se
mantuvo, así como la posesión adquirida en ese concepto, la cual sólo podía perder dicho carácter,
hasta el momento en que el poseedor conociera tales vicios.


Amparo civil directo 598/48. García Astolfo M. 6 de octubre de 1949. Mayoría de tres votos.
Disidentes: Agustín Mercado Alarcón e Hilario Medina. La publicación no menciona el nombre del
ponente.




                                             Pág. 1 de 1             Fecha de impresión 01/10/2025
  https://sjf2.scjn.gob.mx/detalle/tesis/349370

\end{quote}

\subsection{Formalization}

\subsubsection{Clause: tesis349370\_principio\_fundamental}

\textit{Thesis-specific principle: Just title not required for adverse possession, but must reveal origin of possession and affirm possession as owner}

\textbf{Intermediate Representation:}

\begin{quote}
ejercer prescripcion positiva is obligated and possession as owner is obligated
\end{quote}

\textbf{Witness Clause:}

\begin{lstlisting}[style=witnessstyle, language=Witness]
clause tesis349370_principio_fundamental = oblig(ejercer_prescripcion_positiva) IMPLIES oblig(causa_generadora_revelada) AND oblig(posesion_como_propietario);
\end{lstlisting}

\subsubsection{Clause: tesis349370\_determinacion\_judicial}

\textit{Supporting logic: Judge needs generating fact/act to determine possession quality and timing}

\textbf{Intermediate Representation:}

\begin{quote}
If causa generadora revelada is obligated, then (possession as owner is obligated and (good faith exercise is obligated and not (bad faith exercise is obligated)))
\end{quote}

\textbf{Witness Clause:}

\begin{lstlisting}[style=witnessstyle, language=Witness]
clause tesis349370_determinacion_judicial = oblig(causa_generadora_revelada) IMPLIES (oblig(posesion_como_propietario) AND (oblig(ejercicio_buena_fe) AND not(oblig(ejercicio_mala_fe))));
\end{lstlisting}

\subsubsection{Clause: tesis349370\_tipo\_posesion}

\textit{Possession type determination: Can determine if original or derived possession}

\textbf{Intermediate Representation:}

\begin{quote}
If causa generadora revelada is obligated, then (possession as owner is obligated and not (temporal possession is obligated))
\end{quote}

\textbf{Witness Clause:}

\begin{lstlisting}[style=witnessstyle, language=Witness]
clause tesis349370_tipo_posesion = oblig(causa_generadora_revelada) IMPLIES (oblig(posesion_como_propietario) AND not(oblig(posesion_temporal)));
\end{lstlisting}

\subsubsection{Clause: tesis349370\_determinacion\_temporal}

\textit{Timing requirement: Must establish when prescription period starts counting}

\textbf{Intermediate Representation:}

\begin{quote}
If causa generadora revelada is obligated, then (posesion tiempo 5 anos is obligated or posesion tiempo 10 anos is obligated)
\end{quote}

\textbf{Witness Clause:}

\begin{lstlisting}[style=witnessstyle, language=Witness]
clause tesis349370_determinacion_temporal = oblig(causa_generadora_revelada) IMPLIES (oblig(posesion_tiempo_5_anos) OR oblig(posesion_tiempo_10_anos));
\end{lstlisting}

\newpage

\section{Thesis 349788: PRESCRIPCION POSITIVA, TITULO PARA LA (LEGISLACION DE NUEVO LEON).}
\label{thesis:349788}

\subsection{Original Thesis Text}

\begin{quote}
                                                   Semanario Judicial de la Federación

                                                  Tesis

 Registro digital: 349788

 Instancia: Tercera Sala           Quinta Época                            Materia(s): Civil

                                   Fuente: Semanario Judicial de la        Tipo: Aislada
                                   Federación.
                                   Tomo LXXXII, página 2789


PRESCRIPCION POSITIVA, TITULO PARA LA (LEGISLACION DE NUEVO LEON).


De acuerdo con el Código Civil vigente en el año de mil novecientos veintiocho, en el Estado de
Nuevo León, era indispensable para adquirir en propiedad un inmueble por medio de la
prescripción, que ésta se fincará en un título suficiente para transferir el dominio, si se alegaba una
posesión de buena fe. Ahora bien, si el albacea de una sucesión comprometió en venta un inmueble
de ésta, firmando una minuta en la que se obligó a obtener el consentimiento de los herederos para
elevar a escritura pública la operación, y el comprador, después de esperar algún tiempo, intentó un
juicio en el que demandó el perfeccionamiento de la compraventa, exigiendo el consenso de los
herederos y el otorgamiento de la escritura pública respectiva, juicio que se falló desfavorablemente
a sus pretensiones, tanto en primera como en segunda instancia, por lo que el actor solicitó amparo,
el cual le fue negado, por estimarse que en el caso había un título sin validez alguna como contrato
de compraventa, por carecer de un requisito esencialísimo, que era el consentimiento de los
herederos dueños del inmueble comprometido, en tales condiciones, debe estimarse que el
comprador no pudo oponer la prescripción de buena fe, en el juicio iniciado en su contra por la
sucesión, sobre reivindicación del inmueble, porque nunca tuvo título suficiente que le hubiera
transmitido el dominio de éste o le hubiera servido de base para una posesión legal, como lo exige
el Código Civil citado, pues si bien el título suficiente para transferir el dominio requerido por la ley,
como base de la prescripción de buena fe, no es un título irreprochable, porque entonces bastaría
por sí sólo para demostrar la propiedad, sin la necesidad de invocar la prescripción, tampoco puede
ser cualquier documento y menos aún el que carezca de un elemento esencial, como es el
consentimiento del dueño, para enajenar el inmueble. La ley se refiere a los títulos que tienen todas
las apariencias de legalidad, de fondo y de forma, pero en los que luego aparece un vicio en la
propiedad del que vendió, que puede nulificar la enajenación hecha por éste a un tercero que
compró de buena fe. La única defensa de este último para conservar la propiedad del inmueble, es
alegar la prescripción, cuando ha transcurrido el tiempo necesario señalado por la ley, demostrando
que tiene un título constante en escritura pública, debidamente registrado, que reúne todos los
elementos esenciales para transferir el dominio, como son: consentimiento del vendedor, capacidad,
precio y objeto lícito; y si después aparece que el vendedor adquirió el bien indebidamente, por
causa de error, dolo, violencia, incapacidad, etcétera, surge entonces la lesión o deficiencia en la
propiedad a que se refieren los autores, y que da lugar a que el último adquiriente del bien, se
defienda alegando la prescripción de buena fe, contra la sentencia que haya declarado nula la
adquisición del que la vendió.


Amparo civil directo 7425/41. Pérez Jesús C. y coagraviada, sucesiones acumuladas. 7 de
noviembre de 1944. Mayoría de tres votos. Disidentes: Agustín Mercado Alarcón y Emilio Pardo
Aspe. La publicación no menciona el nombre del ponente. Engrose: Vicente Santos Guajardo.

                                               Pág. 1 de 2              Fecha de impresión 01/10/2025
  https://sjf2.scjn.gob.mx/detalle/tesis/349788
                                                Semanario Judicial de la Federación




                                           Pág. 2 de 2       Fecha de impresión 01/10/2025
https://sjf2.scjn.gob.mx/detalle/tesis/349788

\end{quote}

\subsection{Formalization}

\subsubsection{Clause: tesis349788\_principio\_fundamental}

\textit{Thesis-specific principle: Good faith requires sufficient title/generating cause, ignorance of title defects, and founded belief of ownership}

\textbf{Intermediate Representation:}

\begin{quote}
good faith exercise is obligated and not (conocimiento vicios is obligated)
\end{quote}

\textbf{Witness Clause:}

\begin{lstlisting}[style=witnessstyle, language=Witness]
clause tesis349788_principio_fundamental = oblig(ejercicio_buena_fe) IMPLIES oblig(causa_generadora_probada) AND not(oblig(conocimiento_vicios));
\end{lstlisting}

\subsubsection{Clause: tesis349788\_mala\_fe}

\textit{Supporting logic: Bad faith possessor has no title/generating cause OR knows title defects that prevent rightful possession}

\textbf{Intermediate Representation:}

\begin{quote}
If bad faith exercise is obligated, then (not (proven generating cause is obligated) or conocimiento vicios is obligated)
\end{quote}

\textbf{Witness Clause:}

\begin{lstlisting}[style=witnessstyle, language=Witness]
clause tesis349788_mala_fe = oblig(ejercicio_mala_fe) IMPLIES (not(oblig(causa_generadora_probada)) OR oblig(conocimiento_vicios));
\end{lstlisting}

\subsubsection{Clause: tesis349788\_requisitos\_buena\_fe}

\textit{Good faith requirements: Must have sufficient title and ignorance of defects}

\textbf{Intermediate Representation:}

\begin{quote}
good faith exercise is obligated
\end{quote}

\textbf{Witness Clause:}

\begin{lstlisting}[style=witnessstyle, language=Witness]
clause tesis349788_requisitos_buena_fe = oblig(ejercicio_buena_fe) IMPLIES oblig(causa_generadora_probada);
\end{lstlisting}

\subsubsection{Clause: tesis349788\_condiciones\_mala\_fe}

\textit{Bad faith conditions: No title or knowledge of defects}

\textbf{Intermediate Representation:}

\begin{quote}
(not(oblig(causa\_generadora\_probada)) OR oblig(conocimiento\_vicios)) IMPLIES oblig(ejercicio\_mala\_fe)
\end{quote}

\textbf{Witness Clause:}

\begin{lstlisting}[style=witnessstyle, language=Witness]
clause tesis349788_condiciones_mala_fe = (not(oblig(causa_generadora_probada)) OR oblig(conocimiento_vicios)) IMPLIES oblig(ejercicio_mala_fe);
\end{lstlisting}

\newpage

\section{Thesis 351465: PRESCRIPCION POSITIVA, LEY APLICABLE PARA COMPUTARLA (LEGISLACION DE}
\label{thesis:351465}

\subsection{Original Thesis Text}

\begin{quote}
                                                   Semanario Judicial de la Federación

                                                  Tesis

 Registro digital: 351465

 Instancia: Tercera Sala           Quinta Época                           Materia(s): Civil

                                   Fuente: Semanario Judicial de la       Tipo: Aislada
                                   Federación.
                                   Tomo LXXV, página 6901


PRESCRIPCION POSITIVA, LEY APLICABLE PARA COMPUTARLA (LEGISLACION DE
TAMAULIPAS).


El artículo 1o., transitorio, del Código Civil vigente en Tamaulipas, establece: "Los actos anteriores a
la vigencia del presente código, se regirán por sus disposiciones si con la aplicación de éstas no se
violan derechos adquiridos"; y el artículo 3o., transitorio, del mismo código, previene: "Los plazos
que estén corriendo para prescribir o para cualquier otro efecto, se computarán conforme al
presente código, sin perjuicio de lo dispuesto por el artículo primero transitorio". Ahora bien, estos
preceptos no contienen un criterio exacto para saber cuándo la aplicación de una nueva norma viola
o no derechos adquiridos, tratándose de prescripciones que comenzaron a correr durante la
vigencia del código anterior. Este problema se resuelve de manera clara y precisa, en los artículos
2o. y 6o., transitorios, del Código Civil del Distrito Federal; y para la solución buscada, tratándose de
la legislación de Tamaulipas, es dable invocar, si bien sólo como doctrina, el criterio en que se
apoyaron estas disposiciones legales. Por tanto, para computar el término de la prescripción
adquisitiva, que comenzó a correr antes de la vigencia del actual Código Civil de Tamaulipas, debe
aplicarse este código que establece el término de cinco años para la prescripción de buena fe, pero
el tiempo transcurrido hasta la fecha en que entró en vigor, debe reducirse a la mitad, pues en la
misma proporción redujo dicho código el término de la prescripción.


Amparo civil directo 3968/42. Ibarra Carlos. 17 de marzo de 1943. Mayoría de tres votos.
Disidentes: Nicéforo Guerrero e Hilario Medina. La publicación no menciona el nombre del ponente.




                                              Pág. 1 de 1              Fecha de impresión 01/10/2025
  https://sjf2.scjn.gob.mx/detalle/tesis/351465

\end{quote}

\subsection{Formalization}

\subsubsection{Clause: tesis351465\_principio\_fundamental}

\textit{Thesis-specific principle: Adverse possession can be invoked by action or exception, but requires same proof elements}

\textbf{Intermediate Representation:}

\begin{quote}
possession as owner is obligated and (good faith exercise is obligated or bad faith exercise is obligated) and (oblig(posesion\_tiempo\_5\_anos) OR oblig(posesion\_tiempo\_10\_anos)) IMPLIES oblig(ejercer\_prescripcion\_positiva)
\end{quote}

\textbf{Witness Clause:}

\begin{lstlisting}[style=witnessstyle, language=Witness]
clause tesis351465_principio_fundamental = oblig(posesion_como_propietario) AND (oblig(ejercicio_buena_fe) OR oblig(ejercicio_mala_fe)) AND (oblig(posesion_tiempo_5_anos) OR oblig(posesion_tiempo_10_anos)) IMPLIES oblig(ejercer_prescripcion_positiva);
\end{lstlisting}

\subsubsection{Clause: tesis351465\_igualdad\_prueba}

\textit{Supporting logic: Both action and exception require same proof elements regardless of procedural route}

\textbf{Intermediate Representation:}

\begin{quote}
ejercer prescripcion positiva is obligated
\end{quote}

\textbf{Witness Clause:}

\begin{lstlisting}[style=witnessstyle, language=Witness]
clause tesis351465_igualdad_prueba = oblig(ejercer_prescripcion_positiva) IMPLIES oblig(posesion_como_propietario);
\end{lstlisting}

\subsubsection{Clause: tesis351465\_plazos\_temporales}

\textit{Time requirements: 5 years for good faith, 10 years for bad faith}

\textbf{Intermediate Representation:}

\begin{quote}
good faith exercise is obligated
\end{quote}

\textbf{Witness Clause:}

\begin{lstlisting}[style=witnessstyle, language=Witness]
clause tesis351465_plazos_temporales = oblig(ejercicio_buena_fe) IMPLIES oblig(posesion_tiempo_5_anos);
\end{lstlisting}

\subsubsection{Clause: tesis351465\_plazos\_mala\_fe}

\textbf{Intermediate Representation:}

\begin{quote}
bad faith exercise is obligated
\end{quote}

\textbf{Witness Clause:}

\begin{lstlisting}[style=witnessstyle, language=Witness]
clause tesis351465_plazos_mala_fe = oblig(ejercicio_mala_fe) IMPLIES oblig(posesion_tiempo_10_anos);
\end{lstlisting}

\subsubsection{Clause: tesis351465\_objetivo\_registral}

\textit{Target requirement: Action must be directed against registered owner}

\textbf{Intermediate Representation:}

\begin{quote}
ejercer prescripcion positiva is obligated
\end{quote}

\textbf{Witness Clause:}

\begin{lstlisting}[style=witnessstyle, language=Witness]
clause tesis351465_objetivo_registral = oblig(ejercer_prescripcion_positiva) IMPLIES oblig(titularidad_registral);
\end{lstlisting}

\newpage

\section{Thesis 351663: PRESCRIPCION POSITIVA, POR VIRTUD DE LA POSESION (LEGISLACION DE VERACRUZ).}
\label{thesis:351663}

\subsection{Original Thesis Text}

\begin{quote}
                                                   Semanario Judicial de la Federación

                                                  Tesis

 Registro digital: 351663

 Instancia: Tercera Sala          Quinta Época                          Materia(s): Civil

                                  Fuente: Semanario Judicial de la      Tipo: Aislada
                                  Federación.
                                  Tomo LXXIV, página 3162


PRESCRIPCION POSITIVA, POR VIRTUD DE LA POSESION (LEGISLACION DE VERACRUZ).


Por título en la posesión, debe entenderse la causa que la genera, y según el artículo 1074 del
Código Civil de Veracruz, se llama justo título, el que es o fundadamente se cree bastante para
transferir el dominio. Ahora bien, si una persona dirigió a otra una carta, en la cual convino en
venderle una casa, y le acusó recibo de parte del precio, obligándose a extender la escritura pública
tan pronto como recibiera el total del importe de la venta concertada, debe estimarse que dicha
carta constituye un título para la posesión que de la casa haya tenido el comprador, y si la
mencionada posesión fue de buena fe y por más de diez años, debe concluirse que se operó la
prescripción, conforme a lo dispuesto por el artículo 1080 del código citado. No obsta en contrario
que la venta no se hubiere celebrado en escritura pública, requerida por la ley, de acuerdo con la
cuantía de la operación, pues si así hubiere sido, se habría transmitido el dominio de la cosa, desde
el momento del otorgamiento de la escritura y no sería necesaria la prueba de la posesión, para
adquirir el inmueble por prescripción positiva.


Amparo civil directo 3461/42. Guarneros Ambrosio. 5 de noviembre de 1942. Mayoría de tres votos.
Disidentes: Felipe de J. Tena y Nicéforo Guerrero. La publicación no menciona el nombre del
ponente.




                                             Pág. 1 de 1             Fecha de impresión 01/10/2025
  https://sjf2.scjn.gob.mx/detalle/tesis/351663

\end{quote}

\subsection{Formalization}

\subsubsection{Clause: tesis351663\_principio\_fundamental}

\textit{Additional actions for this thesis Additional assets for proof requirements Faith exclusivity: Adverse possession is either good faith XOR bad faith Thesis-specific principle: Good faith adverse possession requires proving authenticity, date, and diligence}

\textbf{Intermediate Representation:}

\begin{quote}
If good faith exercise is obligated and adverse possession claim is claimed, then (autenticidad titulo is obligated and fecha fehaciente is obligated and diligencia adquirente is obligated)
\end{quote}

\textbf{Witness Clause:}

\begin{lstlisting}[style=witnessstyle, language=Witness]
clause tesis351663_principio_fundamental = oblig(ejercicio_buena_fe) AND claim(demanda_prescripcion) IMPLIES (oblig(autenticidad_titulo) AND oblig(fecha_fehaciente) AND oblig(diligencia_adquirente));
\end{lstlisting}

\subsubsection{Clause: tesis351663\_prueba\_autenticidad}

\textit{Proof requirements: Must prove title authenticity and date reliably}

\textbf{Intermediate Representation:}

\begin{quote}
If proven generating cause is obligated, then (autenticidad titulo is obligated and fecha fehaciente is obligated)
\end{quote}

\textbf{Witness Clause:}

\begin{lstlisting}[style=witnessstyle, language=Witness]
clause tesis351663_prueba_autenticidad = oblig(causa_generadora_probada) IMPLIES (oblig(autenticidad_titulo) AND oblig(fecha_fehaciente));
\end{lstlisting}

\subsubsection{Clause: tesis351663\_diligencia\_buena\_fe}

\textit{Diligence requirement: Good faith requires demonstrating due diligence in acquisition}

\textbf{Intermediate Representation:}

\begin{quote}
good faith exercise is obligated
\end{quote}

\textbf{Witness Clause:}

\begin{lstlisting}[style=witnessstyle, language=Witness]
clause tesis351663_diligencia_buena_fe = oblig(ejercicio_buena_fe) IMPLIES oblig(diligencia_adquirente);
\end{lstlisting}

\subsubsection{Clause: tesis351663\_prueba\_pagos}

\textit{Payment proof: Onerous contracts require proving payments made}

\textbf{Intermediate Representation:}

\begin{quote}
proven generating cause is obligated and contrato compraventa is obligated
\end{quote}

\textbf{Witness Clause:}

\begin{lstlisting}[style=witnessstyle, language=Witness]
clause tesis351663_prueba_pagos = oblig(causa_generadora_probada) AND oblig(contrato_compraventa) IMPLIES oblig(pagos_precio);
\end{lstlisting}

\subsubsection{Clause: tesis351663\_carga\_prueba}

\textit{Burden of proof: Claimant must prove generating cause with reliable evidence}

\textbf{Intermediate Representation:}

\begin{quote}
adverse possession claim is claimed
\end{quote}

\textbf{Witness Clause:}

\begin{lstlisting}[style=witnessstyle, language=Witness]
clause tesis351663_carga_prueba = claim(demanda_prescripcion) IMPLIES oblig(causa_generadora_revelada);
\end{lstlisting}

\subsubsection{Clause: tesis351663\_falta\_diligencia}

\textit{Insufficient diligence consequence: Lack of diligence means bad faith (10-year period)}

\textbf{Intermediate Representation:}

\begin{quote}
not (diligencia adquirente is obligated)
\end{quote}

\textbf{Witness Clause:}

\begin{lstlisting}[style=witnessstyle, language=Witness]
clause tesis351663_falta_diligencia = not(oblig(diligencia_adquirente)) IMPLIES oblig(ejercicio_mala_fe);
\end{lstlisting}

\newpage

\section{Thesis 352533: PRESCRIPCION ADQUISITIVA (LEGISLACION DE HIDALGO).}
\label{thesis:352533}

\subsection{Original Thesis Text}

\begin{quote}
                                                   Semanario Judicial de la Federación

                                                  Tesis

 Registro digital: 352533

 Instancia: Tercera Sala          Quinta Época                         Materia(s): Civil

                                  Fuente: Semanario Judicial de la     Tipo: Aislada
                                  Federación.
                                  Tomo LXXII, página 4948


PRESCRIPCION ADQUISITIVA (LEGISLACION DE HIDALGO).


Si el demandado en un juicio sobre nulidad de un contrato de compraventa, relativo a un predio,
opuso la excepción de prescripción y demostró en ese juicio, que ha estado en posesión pacífica,
continua y pública, del predio de que se trata, durante más de veinticinco años, a partir de la fecha
del contrato, y más de veintiuno, desde que el mismo fue registrado, debe decirse que ya sea que
esa posesión la hubiera tenido de buena o mala fe, de todos modos transcurrieron los términos
señalados en el artículo 1045 del Código Civil del Estado de Hidalgo, y por lo mismo, es de
considerarse que el demandado adquirió en propiedad el inmueble, por prescripción adquisitiva, y
que la acción ejercitada en su contra, sobre nulidad del contrato de compraventa, no pudo ser
procedente.


Amparo civil directo 5362/39. Ramírez Modesto. 11 de junio de 1942. Mayoría de cuatro votos.
Disidente: Carlos I. Meléndez. La publicación no menciona el nombre del ponente.




                                             Pág. 1 de 1             Fecha de impresión 01/10/2025
  https://sjf2.scjn.gob.mx/detalle/tesis/352533

\end{quote}

\subsection{Formalization}

\subsubsection{Clause: tesis352533\_principio\_fundamental}

\textit{Thesis-specific principle: When defendant admits possession without title for more than 10 years with required qualities, is bad faith possessor and property is eligible for prescription}

\textbf{Intermediate Representation:}

\begin{quote}
If (posesion\_como\_propietario AND posesion\_pacifico\_continuo\_publico AND posesion\_tiempo\_10\_anos AND not(causa\_generadora\_probada)), then (ejercicio\_mala\_fe AND oblig(ejercer\_prescripcion\_positiva))
\end{quote}

\textbf{Witness Clause:}

\begin{lstlisting}[style=witnessstyle, language=Witness]
clause tesis352533_principio_fundamental = (posesion_como_propietario AND posesion_pacifico_continuo_publico AND posesion_tiempo_10_anos AND not(causa_generadora_probada)) IMPLIES (ejercicio_mala_fe AND oblig(ejercer_prescripcion_positiva));
\end{lstlisting}

\subsubsection{Clause: tesis352533\_no\_prueba\_causa}

\textit{Supporting logic: Does not need to prove generating cause of possession in trial when no title exists}

\textbf{Intermediate Representation:}

\begin{quote}
(not(causa\_generadora\_probada) AND posesion\_tiempo\_10\_anos) IMPLIES oblig(ejercer\_prescripcion\_positiva)
\end{quote}

\textbf{Witness Clause:}

\begin{lstlisting}[style=witnessstyle, language=Witness]
clause tesis352533_no_prueba_causa = (not(causa_generadora_probada) AND posesion_tiempo_10_anos) IMPLIES oblig(ejercer_prescripcion_positiva);
\end{lstlisting}

\subsubsection{Clause: tesis352533\_definicion\_titulo}

\textit{Title definition: In state law, title is understood as the generating cause of possession}

\textbf{Intermediate Representation:}

\begin{quote}
not (causa\_generadora\_probada) IMPLIES not(titularidad\_registral)
\end{quote}

\textbf{Witness Clause:}

\begin{lstlisting}[style=witnessstyle, language=Witness]
clause tesis352533_definicion_titulo = not(causa_generadora_probada) IMPLIES not(titularidad_registral);
\end{lstlisting}

\subsubsection{Clause: tesis352533\_apto\_prescripcion}

\textit{Prescription eligibility: If 10 years have passed and other elements are met, property is eligible for prescription}

\textbf{Intermediate Representation:}

\begin{quote}
(posesion\_tiempo\_10\_anos AND posesion\_como\_propietario AND posesion\_pacifico\_continuo\_publico) IMPLIES oblig(ejercer\_prescripcion\_positiva)
\end{quote}

\textbf{Witness Clause:}

\begin{lstlisting}[style=witnessstyle, language=Witness]
clause tesis352533_apto_prescripcion = (posesion_tiempo_10_anos AND posesion_como_propietario AND posesion_pacifico_continuo_publico) IMPLIES oblig(ejercer_prescripcion_positiva);
\end{lstlisting}

\newpage

\section{Thesis 352678: PRESCRIPCION ADQUISITIVA, JUSTO TITULO EN LA POSESION, COMO REQUISITO DE LA (LEGISLACION DE MICHOACAN).}
\label{thesis:352678}

\subsection{Original Thesis Text}

\begin{quote}
                                                   Semanario Judicial de la Federación

                                                  Tesis

 Registro digital: 352678

 Instancia: Tercera Sala          Quinta Época                         Materia(s): Civil

                                  Fuente: Semanario Judicial de la     Tipo: Aislada
                                  Federación.
                                  Tomo LXXI, página 834


PRESCRIPCION ADQUISITIVA, JUSTO TITULO EN LA POSESION, COMO REQUISITO DE LA
(LEGISLACION DE MICHOACAN).


Conforme al artículo 1080 del Código Civil del Estado de Michoacán, se entiende por justo título, el
que es bastante para transferir el dominio o fundadamente se cree así. Por tanto, un contrato de
compraventa celebrado respecto de una casa que es objeto de un juicio de reivindicación, debe
considerarse como justo título de la posesión que de ella haya tenido el heredero del comprador,
que hubiere sido demandado en dicho juicio, aun cuando ese contrato adolezca de deficiencias, si
dicho demandado pudo legalmente, considerar que, mediante el mismo, había adquirido el de cujus
la propiedad de la finca en cuestión, y por lo mismo, si el demandado en el juicio reivindicatorio
opuso la excepción de prescripción, fundada en el artículo 1085 del código citado, y rindió además
prueba testimonial, con la que acreditó que, de buena fe, ha venido poseyendo la finca por más de
20 años, debe tenerse por comprobada suficientemente la excepción mencionada, y la sentencia
que en tal sentido se dicte, no es violatoria de garantías.


Amparo civil directo 2680/36. Molina Dionisio. 19 de enero de 1942. Unanimidad de cuatros votos.
Ausente: Felipe de J. Tena. Relator: Hilario Medina.




                                             Pág. 1 de 1             Fecha de impresión 01/10/2025
  https://sjf2.scjn.gob.mx/detalle/tesis/352678

\end{quote}

\subsection{Formalization}

\subsubsection{Clause: tesis352678\_principio\_fundamental}

\textit{Faith exclusivity: Possession is either good faith XOR bad faith (mutually exclusive) Thesis-specific principle: Unregistered contracts cannot establish adverse possession against registered owners}

\textbf{Intermediate Representation:}

\begin{quote}
registered ownership is obligated and unregistered contract is obligated
\end{quote}

\textbf{Witness Clause:}

\begin{lstlisting}[style=witnessstyle, language=Witness]
clause tesis352678_principio_fundamental = oblig(titularidad_registral) AND oblig(contrato_no_inscrito) IMPLIES not(claim(demanda_prescripcion));
\end{lstlisting}

\subsubsection{Clause: tesis352678\_proteccion\_registral}

\textit{Supporting legal logic}

\textbf{Intermediate Representation:}

\begin{quote}
registered ownership is obligated
\end{quote}

\textbf{Witness Clause:}

\begin{lstlisting}[style=witnessstyle, language=Witness]
clause tesis352678_proteccion_registral = oblig(titularidad_registral) IMPLIES claim(oponibilidad_terceros);
\end{lstlisting}

\subsubsection{Clause: tesis352678\_limitacion\_no\_inscripcion}

\textbf{Intermediate Representation:}

\begin{quote}
unregistered contract is obligated
\end{quote}

\textbf{Witness Clause:}

\begin{lstlisting}[style=witnessstyle, language=Witness]
clause tesis352678_limitacion_no_inscripcion = oblig(contrato_no_inscrito) IMPLIES not(oblig(oponibilidad_terceros));
\end{lstlisting}

\subsubsection{Clause: tesis352678\_consecuencia\_logica}

\textbf{Intermediate Representation:}

\begin{quote}
not (opposability to third parties is obligated)
\end{quote}

\textbf{Witness Clause:}

\begin{lstlisting}[style=witnessstyle, language=Witness]
clause tesis352678_consecuencia_logica = not(oblig(oponibilidad_terceros)) IMPLIES not(claim(demanda_prescripcion));
\end{lstlisting}

\newpage

\section{Thesis 353700: PRESCRIPCION POSITIVA, SOLO TIENE LUGAR RESPECTO DE LOS POSEEDORES DE BUENA FE (LEGISLACION DEL PUEBLA).}
\label{thesis:353700}

\subsection{Original Thesis Text}

\begin{quote}
                                                   Semanario Judicial de la Federación

                                                  Tesis

 Registro digital: 353700

 Instancia: Tercera Sala          Quinta Época                          Materia(s): Civil

                                  Fuente: Semanario Judicial de la      Tipo: Aislada
                                  Federación.
                                  Tomo LXIX, página 4968


PRESCRIPCION POSITIVA, SOLO TIENE LUGAR RESPECTO DE LOS POSEEDORES DE
BUENA FE (LEGISLACION DEL PUEBLA).


Si el quejoso estuvo poseyendo por más de veinte años un bien inmueble; pero esa posesión no fue
a nombre propio, sino a nombre del dueño del bien, es de estimarse, conforme a los artículos 816 y
817 del Código Civil del Estado de Puebla, que no ha sido poseedor en derecho, del inmueble de
referencia, por lo que, de acuerdo con lo dispuesto por el artículo 1060 del mismo ordenamiento, no
pudo adquirirlo por prescripción; tanto más si el quejoso declaró que no tiene ningún título a su
favor, con relación al mismo bien, pues entonces es inconcuso que no reúne los requisitos exigidos
por los artículos 1069 y 1071 del código citado, para poder adquirir por prescripción, el inmueble de
que se trata.


Amparo civil directo 2016/38. Quiroz Filiberta y coagraviada. 29 de septiembre de 1941. Unanimidad
de cinco votos. Relator: Hilario Medina.




                                             Pág. 1 de 1             Fecha de impresión 01/10/2025
  https://sjf2.scjn.gob.mx/detalle/tesis/353700

\end{quote}

\subsection{Formalization}

\subsubsection{Clause: tesis353700\_principio\_fundamental}

\textit{Thesis-specific principle: Good faith requires sufficient title/generating cause, ignorance of title defects, and founded belief of ownership}

\textbf{Intermediate Representation:}

\begin{quote}
good faith exercise is obligated and not (conocimiento vicios is obligated)
\end{quote}

\textbf{Witness Clause:}

\begin{lstlisting}[style=witnessstyle, language=Witness]
clause tesis353700_principio_fundamental = oblig(ejercicio_buena_fe) IMPLIES oblig(causa_generadora_probada) AND not(oblig(conocimiento_vicios));
\end{lstlisting}

\subsubsection{Clause: tesis353700\_mala\_fe}

\textit{Supporting logic: Bad faith possessor has no title/generating cause OR knows title defects that prevent rightful possession}

\textbf{Intermediate Representation:}

\begin{quote}
If bad faith exercise is obligated, then (not (proven generating cause is obligated) or conocimiento vicios is obligated)
\end{quote}

\textbf{Witness Clause:}

\begin{lstlisting}[style=witnessstyle, language=Witness]
clause tesis353700_mala_fe = oblig(ejercicio_mala_fe) IMPLIES (not(oblig(causa_generadora_probada)) OR oblig(conocimiento_vicios));
\end{lstlisting}

\subsubsection{Clause: tesis353700\_requisitos\_buena\_fe}

\textit{Good faith requirements: Must have sufficient title and ignorance of defects}

\textbf{Intermediate Representation:}

\begin{quote}
good faith exercise is obligated
\end{quote}

\textbf{Witness Clause:}

\begin{lstlisting}[style=witnessstyle, language=Witness]
clause tesis353700_requisitos_buena_fe = oblig(ejercicio_buena_fe) IMPLIES oblig(causa_generadora_probada);
\end{lstlisting}

\subsubsection{Clause: tesis353700\_condiciones\_mala\_fe}

\textit{Bad faith conditions: No title or knowledge of defects}

\textbf{Intermediate Representation:}

\begin{quote}
(not(oblig(causa\_generadora\_probada)) OR oblig(conocimiento\_vicios)) IMPLIES oblig(ejercicio\_mala\_fe)
\end{quote}

\textbf{Witness Clause:}

\begin{lstlisting}[style=witnessstyle, language=Witness]
clause tesis353700_condiciones_mala_fe = (not(oblig(causa_generadora_probada)) OR oblig(conocimiento_vicios)) IMPLIES oblig(ejercicio_mala_fe);
\end{lstlisting}

\newpage

\section{Thesis 353892: POSESION, PRESCRIPCION ADQUISITIVA DERIVADA DE LA.}
\label{thesis:353892}

\subsection{Original Thesis Text}

\begin{quote}
                                                   Semanario Judicial de la Federación

                                                  Tesis

 Registro digital: 353892

 Instancia: Tercera Sala          Quinta Época                           Materia(s): Civil

                                  Fuente: Semanario Judicial de la       Tipo: Aislada
                                  Federación.
                                  Tomo LXVIII, página 3029


POSESION, PRESCRIPCION ADQUISITIVA DERIVADA DE LA.


El Código Civil de 1884 para el Distrito Federal y Territorios de Baja California, sólo exigía el justo
título de la posesión, para la prescripción de buena fe, pues con mala fe, que es la del que posee,
sabiendo que no tiene título, también se adquiere por prescripción, si bien en un plazo mucho más
largo.


Amparo civil directo 5937/38. Jiménez Roque. 27 de junio de 1941. Unanimidad de cuatro votos. El
Ministro Francisco Barba no estuvo presente en este asunto, por el motivo que se expresa en el
acta del día. La publicación no menciona el nombre del ponente.




                                             Pág. 1 de 1              Fecha de impresión 01/10/2025
  https://sjf2.scjn.gob.mx/detalle/tesis/353892

\end{quote}

\subsection{Formalization}

\subsubsection{Clause: tesis353892\_principio\_fundamental}

\textit{Thesis-specific principle: Good faith requires sufficient title/generating cause, ignorance of title defects, and founded belief of ownership}

\textbf{Intermediate Representation:}

\begin{quote}
good faith exercise is obligated and not (conocimiento vicios is obligated)
\end{quote}

\textbf{Witness Clause:}

\begin{lstlisting}[style=witnessstyle, language=Witness]
clause tesis353892_principio_fundamental = oblig(ejercicio_buena_fe) IMPLIES oblig(causa_generadora_probada) AND not(oblig(conocimiento_vicios));
\end{lstlisting}

\subsubsection{Clause: tesis353892\_mala\_fe}

\textit{Supporting logic: Bad faith possessor has no title/generating cause OR knows title defects that prevent rightful possession}

\textbf{Intermediate Representation:}

\begin{quote}
If bad faith exercise is obligated, then (not (proven generating cause is obligated) or conocimiento vicios is obligated)
\end{quote}

\textbf{Witness Clause:}

\begin{lstlisting}[style=witnessstyle, language=Witness]
clause tesis353892_mala_fe = oblig(ejercicio_mala_fe) IMPLIES (not(oblig(causa_generadora_probada)) OR oblig(conocimiento_vicios));
\end{lstlisting}

\subsubsection{Clause: tesis353892\_requisitos\_buena\_fe}

\textit{Good faith requirements: Must have sufficient title and ignorance of defects}

\textbf{Intermediate Representation:}

\begin{quote}
good faith exercise is obligated
\end{quote}

\textbf{Witness Clause:}

\begin{lstlisting}[style=witnessstyle, language=Witness]
clause tesis353892_requisitos_buena_fe = oblig(ejercicio_buena_fe) IMPLIES oblig(causa_generadora_probada);
\end{lstlisting}

\subsubsection{Clause: tesis353892\_condiciones\_mala\_fe}

\textit{Bad faith conditions: No title or knowledge of defects}

\textbf{Intermediate Representation:}

\begin{quote}
(not(oblig(causa\_generadora\_probada)) OR oblig(conocimiento\_vicios)) IMPLIES oblig(ejercicio\_mala\_fe)
\end{quote}

\textbf{Witness Clause:}

\begin{lstlisting}[style=witnessstyle, language=Witness]
clause tesis353892_condiciones_mala_fe = (not(oblig(causa_generadora_probada)) OR oblig(conocimiento_vicios)) IMPLIES oblig(ejercicio_mala_fe);
\end{lstlisting}

\newpage

\section{Thesis 354119: POSESION, PRESCRIPCION ADQUISITIVA DE LA (LEGISLACION DEL ESTADO DE MEXICO).}
\label{thesis:354119}

\subsection{Original Thesis Text}

\begin{quote}
                                                   Semanario Judicial de la Federación

                                                  Tesis

 Registro digital: 354119

 Instancia: Tercera Sala           Quinta Época                          Materia(s): Civil

                                   Fuente: Semanario Judicial de la      Tipo: Aislada
                                   Federación.
                                   Tomo LXVII, página 3599


POSESION, PRESCRIPCION ADQUISITIVA DE LA (LEGISLACION DEL ESTADO DE MEXICO).


De los términos del artículo 200 de la Constitución Política del Estado de México, que establece que
todas las propiedades raíces rústicas y urbanas, prescriben por la sola posesión de veinte años, con
título o sin él, o con buena fe o sin ella, se desprende que para que opere la prescripción adquisitiva
de un inmueble, no es necesario que la posesión reúna todos los elementos a que se contrae la
teoría.


Amparo civil directo 1440/38. Morán Mariana. 26 de marzo de 1941. Unanimidad de cinco votos. La
publicación no menciona el nombre del ponente.




                                              Pág. 1 de 1             Fecha de impresión 01/10/2025
  https://sjf2.scjn.gob.mx/detalle/tesis/354119

\end{quote}

\subsection{Formalization}

\subsubsection{Clause: tesis354119\_principio\_fundamental}

\textit{Additional actions for this thesis Additional assets for proof requirements Faith exclusivity: Adverse possession is either good faith XOR bad faith Thesis-specific principle: Good faith adverse possession requires proving authenticity, date, and diligence}

\textbf{Intermediate Representation:}

\begin{quote}
If good faith exercise is obligated and adverse possession claim is claimed, then (autenticidad titulo is obligated and fecha fehaciente is obligated and diligencia adquirente is obligated)
\end{quote}

\textbf{Witness Clause:}

\begin{lstlisting}[style=witnessstyle, language=Witness]
clause tesis354119_principio_fundamental = oblig(ejercicio_buena_fe) AND claim(demanda_prescripcion) IMPLIES (oblig(autenticidad_titulo) AND oblig(fecha_fehaciente) AND oblig(diligencia_adquirente));
\end{lstlisting}

\subsubsection{Clause: tesis354119\_prueba\_autenticidad}

\textit{Proof requirements: Must prove title authenticity and date reliably}

\textbf{Intermediate Representation:}

\begin{quote}
If proven generating cause is obligated, then (autenticidad titulo is obligated and fecha fehaciente is obligated)
\end{quote}

\textbf{Witness Clause:}

\begin{lstlisting}[style=witnessstyle, language=Witness]
clause tesis354119_prueba_autenticidad = oblig(causa_generadora_probada) IMPLIES (oblig(autenticidad_titulo) AND oblig(fecha_fehaciente));
\end{lstlisting}

\subsubsection{Clause: tesis354119\_diligencia\_buena\_fe}

\textit{Diligence requirement: Good faith requires demonstrating due diligence in acquisition}

\textbf{Intermediate Representation:}

\begin{quote}
good faith exercise is obligated
\end{quote}

\textbf{Witness Clause:}

\begin{lstlisting}[style=witnessstyle, language=Witness]
clause tesis354119_diligencia_buena_fe = oblig(ejercicio_buena_fe) IMPLIES oblig(diligencia_adquirente);
\end{lstlisting}

\subsubsection{Clause: tesis354119\_prueba\_pagos}

\textit{Payment proof: Onerous contracts require proving payments made}

\textbf{Intermediate Representation:}

\begin{quote}
proven generating cause is obligated and contrato compraventa is obligated
\end{quote}

\textbf{Witness Clause:}

\begin{lstlisting}[style=witnessstyle, language=Witness]
clause tesis354119_prueba_pagos = oblig(causa_generadora_probada) AND oblig(contrato_compraventa) IMPLIES oblig(pagos_precio);
\end{lstlisting}

\subsubsection{Clause: tesis354119\_carga\_prueba}

\textit{Burden of proof: Claimant must prove generating cause with reliable evidence}

\textbf{Intermediate Representation:}

\begin{quote}
adverse possession claim is claimed
\end{quote}

\textbf{Witness Clause:}

\begin{lstlisting}[style=witnessstyle, language=Witness]
clause tesis354119_carga_prueba = claim(demanda_prescripcion) IMPLIES oblig(causa_generadora_revelada);
\end{lstlisting}

\subsubsection{Clause: tesis354119\_falta\_diligencia}

\textit{Insufficient diligence consequence: Lack of diligence means bad faith (10-year period)}

\textbf{Intermediate Representation:}

\begin{quote}
not (diligencia adquirente is obligated)
\end{quote}

\textbf{Witness Clause:}

\begin{lstlisting}[style=witnessstyle, language=Witness]
clause tesis354119_falta_diligencia = not(oblig(diligencia_adquirente)) IMPLIES oblig(ejercicio_mala_fe);
\end{lstlisting}

\newpage

\section{Thesis 358628: PRESCRIPCION POSITIVA, NATURALEZA DE LA.}
\label{thesis:358628}

\subsection{Original Thesis Text}

\begin{quote}
                                                   Semanario Judicial de la Federación

                                                  Tesis

 Registro digital: 358628

 Instancia: Tercera Sala          Quinta Época                          Materia(s): Civil

                                  Fuente: Semanario Judicial de la      Tipo: Aislada
                                  Federación.
                                  Tomo XLVIII, página 1891


PRESCRIPCION POSITIVA, NATURALEZA DE LA.


Cuando el artículo 1086 del Código Civil de 1884, previene que la prescripción adquisitiva de los
inmuebles, se consuma, con buena fe, en diez años, y con mala en veinte, y se limita la prevención
con lo dispuesto en el artículo 1070 del propio ordenamiento, el cual requiere que la posesión ad
usucapione, se realice cuando la posesión es a nombre propio, se fija una condición a la cual se
sujeta la eficacia de la prescripción, porque la propia ley, en su artículo 826, previene que el que
posee en nombre de otro, no es poseedor en derecho, de lo que se colige que es condición
indispensable para que se opere la prescripción adquisitiva, que la posesión de donde la misma se
hace derivar, se haya desarrollado jurídicamente, esto es, con el ánimo de apropiación, pues
cuando la posesión es precaria o derivada, no da materia a la adquisición, por prescripción, del bien
poseído.


Amparo civil directo 3904/33. González viuda de Romero Mariana. 4 de mayo de 1936. Unanimidad
de cinco votos. La publicación no menciona el nombre del ponente.




                                             Pág. 1 de 1             Fecha de impresión 01/10/2025
  https://sjf2.scjn.gob.mx/detalle/tesis/358628

\end{quote}

\subsection{Formalization}

\subsubsection{Clause: tesis358628\_principio\_fundamental}

\textit{Thesis-specific principle: Good faith requires sufficient title/generating cause, ignorance of title defects, and founded belief of ownership}

\textbf{Intermediate Representation:}

\begin{quote}
good faith exercise is obligated and not (conocimiento vicios is obligated)
\end{quote}

\textbf{Witness Clause:}

\begin{lstlisting}[style=witnessstyle, language=Witness]
clause tesis358628_principio_fundamental = oblig(ejercicio_buena_fe) IMPLIES oblig(causa_generadora_probada) AND not(oblig(conocimiento_vicios));
\end{lstlisting}

\subsubsection{Clause: tesis358628\_mala\_fe}

\textit{Supporting logic: Bad faith possessor has no title/generating cause OR knows title defects that prevent rightful possession}

\textbf{Intermediate Representation:}

\begin{quote}
If bad faith exercise is obligated, then (not (proven generating cause is obligated) or conocimiento vicios is obligated)
\end{quote}

\textbf{Witness Clause:}

\begin{lstlisting}[style=witnessstyle, language=Witness]
clause tesis358628_mala_fe = oblig(ejercicio_mala_fe) IMPLIES (not(oblig(causa_generadora_probada)) OR oblig(conocimiento_vicios));
\end{lstlisting}

\subsubsection{Clause: tesis358628\_requisitos\_buena\_fe}

\textit{Good faith requirements: Must have sufficient title and ignorance of defects}

\textbf{Intermediate Representation:}

\begin{quote}
good faith exercise is obligated
\end{quote}

\textbf{Witness Clause:}

\begin{lstlisting}[style=witnessstyle, language=Witness]
clause tesis358628_requisitos_buena_fe = oblig(ejercicio_buena_fe) IMPLIES oblig(causa_generadora_probada);
\end{lstlisting}

\subsubsection{Clause: tesis358628\_condiciones\_mala\_fe}

\textit{Bad faith conditions: No title or knowledge of defects}

\textbf{Intermediate Representation:}

\begin{quote}
(not(oblig(causa\_generadora\_probada)) OR oblig(conocimiento\_vicios)) IMPLIES oblig(ejercicio\_mala\_fe)
\end{quote}

\textbf{Witness Clause:}

\begin{lstlisting}[style=witnessstyle, language=Witness]
clause tesis358628_condiciones_mala_fe = (not(oblig(causa_generadora_probada)) OR oblig(conocimiento_vicios)) IMPLIES oblig(ejercicio_mala_fe);
\end{lstlisting}

\newpage

\section{Thesis 359660: PRESCRIPCION POSITIVA (LEGISLACION DE NUEVO LEON).}
\label{thesis:359660}

\subsection{Original Thesis Text}

\begin{quote}
                                                   Semanario Judicial de la Federación

                                                  Tesis

 Registro digital: 359660

 Instancia: Tercera Sala          Quinta Época                         Materia(s): Civil

                                  Fuente: Semanario Judicial de la     Tipo: Aislada
                                  Federación.
                                  Tomo XLV, página 3386


PRESCRIPCION POSITIVA (LEGISLACION DE NUEVO LEON).


No es exacto que para la prescripción, en todas sus categorías, sea necesario un documento o un
título independientemente de la posesión misma, pues la exigencia de la fracción I del artículo 1029
del Código Civil del Estado de Nuevo León, se refiere a que la posesión necesaria para prescribir
sea "a título", es decir, con ánimo o intención de dueño; porque de otro modo, no existiría la
clasificación que la ley hace de posesiones de mala y buena fe, y nunca de mala fe sería eficaz para
la usucapión, lo que sería antijurídico, puesto que existe disposición legal que establece, respecto
de los inmuebles, la prescripción de veinte años para el poseedor de mala fe.


Amparo civil directo 14762/32. Arámbula Leonides y coagraviado. 23 de agosto de 1935. Mayoría
de tres votos. El Ministro Abenamar Eboli Paniagua no votó en este asunto por las razones que
constan en el acta del día. Disidente: Luis Bazdresch. La publicación no menciona el nombre del
ponente.




                                             Pág. 1 de 1             Fecha de impresión 01/10/2025
  https://sjf2.scjn.gob.mx/detalle/tesis/359660

\end{quote}

\subsection{Formalization}

\subsubsection{Clause: tesis359660\_principio\_fundamental}

\textit{Thesis-specific principle: Good faith requires sufficient title/generating cause, ignorance of title defects, and founded belief of ownership}

\textbf{Intermediate Representation:}

\begin{quote}
good faith exercise is obligated and not (conocimiento vicios is obligated)
\end{quote}

\textbf{Witness Clause:}

\begin{lstlisting}[style=witnessstyle, language=Witness]
clause tesis359660_principio_fundamental = oblig(ejercicio_buena_fe) IMPLIES oblig(causa_generadora_probada) AND not(oblig(conocimiento_vicios));
\end{lstlisting}

\subsubsection{Clause: tesis359660\_mala\_fe}

\textit{Supporting logic: Bad faith possessor has no title/generating cause OR knows title defects that prevent rightful possession}

\textbf{Intermediate Representation:}

\begin{quote}
If bad faith exercise is obligated, then (not (proven generating cause is obligated) or conocimiento vicios is obligated)
\end{quote}

\textbf{Witness Clause:}

\begin{lstlisting}[style=witnessstyle, language=Witness]
clause tesis359660_mala_fe = oblig(ejercicio_mala_fe) IMPLIES (not(oblig(causa_generadora_probada)) OR oblig(conocimiento_vicios));
\end{lstlisting}

\subsubsection{Clause: tesis359660\_requisitos\_buena\_fe}

\textit{Good faith requirements: Must have sufficient title and ignorance of defects}

\textbf{Intermediate Representation:}

\begin{quote}
good faith exercise is obligated
\end{quote}

\textbf{Witness Clause:}

\begin{lstlisting}[style=witnessstyle, language=Witness]
clause tesis359660_requisitos_buena_fe = oblig(ejercicio_buena_fe) IMPLIES oblig(causa_generadora_probada);
\end{lstlisting}

\subsubsection{Clause: tesis359660\_condiciones\_mala\_fe}

\textit{Bad faith conditions: No title or knowledge of defects}

\textbf{Intermediate Representation:}

\begin{quote}
(not(oblig(causa\_generadora\_probada)) OR oblig(conocimiento\_vicios)) IMPLIES oblig(ejercicio\_mala\_fe)
\end{quote}

\textbf{Witness Clause:}

\begin{lstlisting}[style=witnessstyle, language=Witness]
clause tesis359660_condiciones_mala_fe = (not(oblig(causa_generadora_probada)) OR oblig(conocimiento_vicios)) IMPLIES oblig(ejercicio_mala_fe);
\end{lstlisting}

\newpage

\section{Thesis 360341: PRESCRIPCION ADQUISITIVA EN VERACRUZ.}
\label{thesis:360341}

\subsection{Original Thesis Text}

\begin{quote}
                                                   Semanario Judicial de la Federación

                                                  Tesis

 Registro digital: 360341

 Instancia: Tercera Sala          Quinta Época                         Materia(s): Civil

                                  Fuente: Semanario Judicial de la     Tipo: Aislada
                                  Federación.
                                  Tomo XLIII, página 2649


PRESCRIPCION ADQUISITIVA EN VERACRUZ.


Con arreglo a lo dispuesto por los artículos 1073, 1074 y 1086 del Código Civil de Veracruz, la
posesión necesaria para prescribir, debe ser fundada en justo título; éste lo constituye el que
fundadamente es o se cree bastante para transferir el dominio, y los inmuebles se prescriben, con
buena fe, en diez años, y con mala fe, en veinte.


Amparo civil directo 525/30. Riveros Balbina. 16 de marzo de 1935. Unanimidad de cinco votos. La
publicación no menciona el nombre del ponente.




                                             Pág. 1 de 1             Fecha de impresión 01/10/2025
  https://sjf2.scjn.gob.mx/detalle/tesis/360341

\end{quote}

\subsection{Formalization}

\subsubsection{Clause: tesis360341\_principio\_fundamental}

\textit{Thesis-specific principle: When defendant admits possession without title for more than 10 years with required qualities, is bad faith possessor and property is eligible for prescription}

\textbf{Intermediate Representation:}

\begin{quote}
If (posesion\_como\_propietario AND posesion\_pacifico\_continuo\_publico AND posesion\_tiempo\_10\_anos AND not(causa\_generadora\_probada)), then (ejercicio\_mala\_fe AND oblig(ejercer\_prescripcion\_positiva))
\end{quote}

\textbf{Witness Clause:}

\begin{lstlisting}[style=witnessstyle, language=Witness]
clause tesis360341_principio_fundamental = (posesion_como_propietario AND posesion_pacifico_continuo_publico AND posesion_tiempo_10_anos AND not(causa_generadora_probada)) IMPLIES (ejercicio_mala_fe AND oblig(ejercer_prescripcion_positiva));
\end{lstlisting}

\subsubsection{Clause: tesis360341\_no\_prueba\_causa}

\textit{Supporting logic: Does not need to prove generating cause of possession in trial when no title exists}

\textbf{Intermediate Representation:}

\begin{quote}
(not(causa\_generadora\_probada) AND posesion\_tiempo\_10\_anos) IMPLIES oblig(ejercer\_prescripcion\_positiva)
\end{quote}

\textbf{Witness Clause:}

\begin{lstlisting}[style=witnessstyle, language=Witness]
clause tesis360341_no_prueba_causa = (not(causa_generadora_probada) AND posesion_tiempo_10_anos) IMPLIES oblig(ejercer_prescripcion_positiva);
\end{lstlisting}

\subsubsection{Clause: tesis360341\_definicion\_titulo}

\textit{Title definition: In state law, title is understood as the generating cause of possession}

\textbf{Intermediate Representation:}

\begin{quote}
not (causa\_generadora\_probada) IMPLIES not(titularidad\_registral)
\end{quote}

\textbf{Witness Clause:}

\begin{lstlisting}[style=witnessstyle, language=Witness]
clause tesis360341_definicion_titulo = not(causa_generadora_probada) IMPLIES not(titularidad_registral);
\end{lstlisting}

\subsubsection{Clause: tesis360341\_apto\_prescripcion}

\textit{Prescription eligibility: If 10 years have passed and other elements are met, property is eligible for prescription}

\textbf{Intermediate Representation:}

\begin{quote}
(posesion\_tiempo\_10\_anos AND posesion\_como\_propietario AND posesion\_pacifico\_continuo\_publico) IMPLIES oblig(ejercer\_prescripcion\_positiva)
\end{quote}

\textbf{Witness Clause:}

\begin{lstlisting}[style=witnessstyle, language=Witness]
clause tesis360341_apto_prescripcion = (posesion_tiempo_10_anos AND posesion_como_propietario AND posesion_pacifico_continuo_publico) IMPLIES oblig(ejercer_prescripcion_positiva);
\end{lstlisting}

\newpage

\section{Thesis 360740: PRESCRIPCION ADQUISITIVA EN LA LEGISLACION DE SINALOA.}
\label{thesis:360740}

\subsection{Original Thesis Text}

\begin{quote}
                                                   Semanario Judicial de la Federación

                                                  Tesis

 Registro digital: 360740

 Instancia: Tercera Sala          Quinta Época                         Materia(s): Civil

                                  Fuente: Semanario Judicial de la     Tipo: Aislada
                                  Federación.
                                  Tomo XLII, página 3767


PRESCRIPCION ADQUISITIVA EN LA LEGISLACION DE SINALOA.


Son condiciones indispensables para que la prescripción adquisitiva opere, que la posesión se
funde en justo título, se realice de buena fe y que concurran las tres circunstancias específicas,
enumeradas en el artículo 1079 del Código Civil de Sinaloa.


Amparo civil directo 195/34. Imaña Federico y coagraviados. 13 de diciembre de 1934. Unanimidad
de cinco votos. La publicación no menciona el nombre del ponente.




                                             Pág. 1 de 1             Fecha de impresión 01/10/2025
  https://sjf2.scjn.gob.mx/detalle/tesis/360740

\end{quote}

\subsection{Formalization}

\subsubsection{Clause: tesis360740\_principio\_fundamental}

\textit{Thesis-specific principle: Good faith requires sufficient title/generating cause, ignorance of title defects, and founded belief of ownership}

\textbf{Intermediate Representation:}

\begin{quote}
good faith exercise is obligated and not (conocimiento vicios is obligated)
\end{quote}

\textbf{Witness Clause:}

\begin{lstlisting}[style=witnessstyle, language=Witness]
clause tesis360740_principio_fundamental = oblig(ejercicio_buena_fe) IMPLIES oblig(causa_generadora_probada) AND not(oblig(conocimiento_vicios));
\end{lstlisting}

\subsubsection{Clause: tesis360740\_mala\_fe}

\textit{Supporting logic: Bad faith possessor has no title/generating cause OR knows title defects that prevent rightful possession}

\textbf{Intermediate Representation:}

\begin{quote}
If bad faith exercise is obligated, then (not (proven generating cause is obligated) or conocimiento vicios is obligated)
\end{quote}

\textbf{Witness Clause:}

\begin{lstlisting}[style=witnessstyle, language=Witness]
clause tesis360740_mala_fe = oblig(ejercicio_mala_fe) IMPLIES (not(oblig(causa_generadora_probada)) OR oblig(conocimiento_vicios));
\end{lstlisting}

\subsubsection{Clause: tesis360740\_requisitos\_buena\_fe}

\textit{Good faith requirements: Must have sufficient title and ignorance of defects}

\textbf{Intermediate Representation:}

\begin{quote}
good faith exercise is obligated
\end{quote}

\textbf{Witness Clause:}

\begin{lstlisting}[style=witnessstyle, language=Witness]
clause tesis360740_requisitos_buena_fe = oblig(ejercicio_buena_fe) IMPLIES oblig(causa_generadora_probada);
\end{lstlisting}

\subsubsection{Clause: tesis360740\_condiciones\_mala\_fe}

\textit{Bad faith conditions: No title or knowledge of defects}

\textbf{Intermediate Representation:}

\begin{quote}
(not(oblig(causa\_generadora\_probada)) OR oblig(conocimiento\_vicios)) IMPLIES oblig(ejercicio\_mala\_fe)
\end{quote}

\textbf{Witness Clause:}

\begin{lstlisting}[style=witnessstyle, language=Witness]
clause tesis360740_condiciones_mala_fe = (not(oblig(causa_generadora_probada)) OR oblig(conocimiento_vicios)) IMPLIES oblig(ejercicio_mala_fe);
\end{lstlisting}

\newpage

\section{Thesis 360742: PRESCRIPCION ADQUISITIVA, BUENA FE PARA LA.}
\label{thesis:360742}

\subsection{Original Thesis Text}

\begin{quote}
                                                   Semanario Judicial de la Federación

                                                  Tesis

 Registro digital: 360742

 Instancia: Tercera Sala          Quinta Época                         Materia(s): Civil

                                  Fuente: Semanario Judicial de la     Tipo: Aislada
                                  Federación.
                                  Tomo XLII, página 3768


PRESCRIPCION ADQUISITIVA, BUENA FE PARA LA.


La buena fe requiere que el poseedor tenga o crea tener su propio título bastante para transferir el
dominio o ignore los vicios de título. Solamente el poseedor que reúna las anteriores condiciones,
puede ser considerado como de buena fe; la presunción de buena fe ha sido establecida en favor
del poseedor.


Amparo civil directo 195/34. Imaña Federico y coagraviados. 13 de diciembre de 1934. Unanimidad
de cinco votos. La publicación no menciona el nombre del ponente.




                                             Pág. 1 de 1             Fecha de impresión 01/10/2025
  https://sjf2.scjn.gob.mx/detalle/tesis/360742

\end{quote}

\subsection{Formalization}

\subsubsection{Clause: tesis360742\_principio\_fundamental}

\textit{Faith exclusivity: Possession is either good faith XOR bad faith (mutually exclusive) Thesis-specific principle: Unregistered contracts cannot establish adverse possession against registered owners}

\textbf{Intermediate Representation:}

\begin{quote}
registered ownership is obligated and unregistered contract is obligated
\end{quote}

\textbf{Witness Clause:}

\begin{lstlisting}[style=witnessstyle, language=Witness]
clause tesis360742_principio_fundamental = oblig(titularidad_registral) AND oblig(contrato_no_inscrito) IMPLIES not(claim(demanda_prescripcion));
\end{lstlisting}

\subsubsection{Clause: tesis360742\_proteccion\_registral}

\textit{Supporting legal logic}

\textbf{Intermediate Representation:}

\begin{quote}
registered ownership is obligated
\end{quote}

\textbf{Witness Clause:}

\begin{lstlisting}[style=witnessstyle, language=Witness]
clause tesis360742_proteccion_registral = oblig(titularidad_registral) IMPLIES claim(oponibilidad_terceros);
\end{lstlisting}

\subsubsection{Clause: tesis360742\_limitacion\_no\_inscripcion}

\textbf{Intermediate Representation:}

\begin{quote}
unregistered contract is obligated
\end{quote}

\textbf{Witness Clause:}

\begin{lstlisting}[style=witnessstyle, language=Witness]
clause tesis360742_limitacion_no_inscripcion = oblig(contrato_no_inscrito) IMPLIES not(oblig(oponibilidad_terceros));
\end{lstlisting}

\subsubsection{Clause: tesis360742\_consecuencia\_logica}

\textbf{Intermediate Representation:}

\begin{quote}
not (opposability to third parties is obligated)
\end{quote}

\textbf{Witness Clause:}

\begin{lstlisting}[style=witnessstyle, language=Witness]
clause tesis360742_consecuencia_logica = not(oblig(oponibilidad_terceros)) IMPLIES not(claim(demanda_prescripcion));
\end{lstlisting}

\newpage

\section{Thesis 361460: THESIS 361460}
\label{thesis:361460}

\subsection{Original Thesis Text}

\begin{quote}
                                                   Semanario Judicial de la Federación

                                                  Tesis

 Registro digital: 361460

 Instancia: Tercera Sala          Quinta Época                         Materia(s): Civil

                                  Fuente: Semanario Judicial de la     Tipo: Aislada
                                  Federación.
                                  Tomo XL, página 2937


PRESCRIPCION POSITIVA.


La excepción de prescripción opuesta en un juicio reivindicatorio, es ineficaz cuando no han
transcurrido veinte años entre la fecha del remate y la en que se presentó la demanda
reivindicatoria, tratándose de los dos primeros adquirentes, ni diez años, tratándose de los
sucesivos, de acuerdo con el artículo 1086 del Código Civil de 1884, para el Distrito Federal, para
que la prescripción positiva opere respecto de bienes inmuebles, cuando son poseídos de mala o de
buena fe.


Amparo civil directo 3937/28. Compañía de Hipotecas y Préstamos, S. C. de R. L. 2 de abril de
1934. Unanimidad de cuatro votos. Ausente: Francisco Díaz Lombardo. La publicación no menciona
el nombre del ponente.




                                             Pág. 1 de 1             Fecha de impresión 01/10/2025
  https://sjf2.scjn.gob.mx/detalle/tesis/361460

\end{quote}

\subsection{Formalization}

\subsubsection{Clause: tesis361460\_principio\_fundamental}

\textit{Thesis-specific principle: Good faith requires sufficient title/generating cause, ignorance of title defects, and founded belief of ownership}

\textbf{Intermediate Representation:}

\begin{quote}
good faith exercise is obligated and not (conocimiento vicios is obligated)
\end{quote}

\textbf{Witness Clause:}

\begin{lstlisting}[style=witnessstyle, language=Witness]
clause tesis361460_principio_fundamental = oblig(ejercicio_buena_fe) IMPLIES oblig(causa_generadora_probada) AND not(oblig(conocimiento_vicios));
\end{lstlisting}

\subsubsection{Clause: tesis361460\_mala\_fe}

\textit{Supporting logic: Bad faith possessor has no title/generating cause OR knows title defects that prevent rightful possession}

\textbf{Intermediate Representation:}

\begin{quote}
If bad faith exercise is obligated, then (not (proven generating cause is obligated) or conocimiento vicios is obligated)
\end{quote}

\textbf{Witness Clause:}

\begin{lstlisting}[style=witnessstyle, language=Witness]
clause tesis361460_mala_fe = oblig(ejercicio_mala_fe) IMPLIES (not(oblig(causa_generadora_probada)) OR oblig(conocimiento_vicios));
\end{lstlisting}

\subsubsection{Clause: tesis361460\_requisitos\_buena\_fe}

\textit{Good faith requirements: Must have sufficient title and ignorance of defects}

\textbf{Intermediate Representation:}

\begin{quote}
good faith exercise is obligated
\end{quote}

\textbf{Witness Clause:}

\begin{lstlisting}[style=witnessstyle, language=Witness]
clause tesis361460_requisitos_buena_fe = oblig(ejercicio_buena_fe) IMPLIES oblig(causa_generadora_probada);
\end{lstlisting}

\subsubsection{Clause: tesis361460\_condiciones\_mala\_fe}

\textit{Bad faith conditions: No title or knowledge of defects}

\textbf{Intermediate Representation:}

\begin{quote}
(not(oblig(causa\_generadora\_probada)) OR oblig(conocimiento\_vicios)) IMPLIES oblig(ejercicio\_mala\_fe)
\end{quote}

\textbf{Witness Clause:}

\begin{lstlisting}[style=witnessstyle, language=Witness]
clause tesis361460_condiciones_mala_fe = (not(oblig(causa_generadora_probada)) OR oblig(conocimiento_vicios)) IMPLIES oblig(ejercicio_mala_fe);
\end{lstlisting}

\newpage

\section{Thesis 385312: REIVINDICACION ACCION DE, FUNDADA EN LA PRESCRIPCION POSITIVA.}
\label{thesis:385312}

\subsection{Original Thesis Text}

\begin{quote}
                                                   Semanario Judicial de la Federación

                                                  Tesis

 Registro digital: 385312

 Instancia: Sala Auxiliar         Quinta Época                         Materia(s): Civil

                                  Fuente: Semanario Judicial de la     Tipo: Aislada
                                  Federación.
                                  Tomo CXVII, página 588


REIVINDICACION ACCION DE, FUNDADA EN LA PRESCRIPCION POSITIVA.


Aun teniendo en cuenta la tesis conforme a la cual la reivindicación se detiene ante la buena fe del
tercer adquirente, la usucapión se sustrae a tal tesis, porque el que posee, por el tiempo y con las
condiciones exigidas por la ley, adquiere por prescripción precisamente contra quien aparece como
propietario del inmueble en el Registro Público de la Propiedad. Por tanto, si el adjudicatario
adquiere el dominio del inmueble afectado con la prescripción, y si ésta se ha consumado contra el
ejecutado, sólo adquiere en apariencia; y en tal virtud, su título sólo sería válido a condición de
proceder de un causante que fuese a su vez propietario.


Amparo civil directo 2961/50. Cerda José. 6 de agosto de 1953. Unanimidad de cuatro votos.
Ausente: Angel González de la Vega. La publicación no menciona el nombre del ponente.




                                             Pág. 1 de 1             Fecha de impresión 01/10/2025
  https://sjf2.scjn.gob.mx/detalle/tesis/385312

\end{quote}

\subsection{Formalization}

\subsubsection{Clause: tesis385312\_principio\_fundamental}

\textit{Thesis-specific principle: Possession must be "in concept of owner" (peaceful, continuous, public) for both good and bad faith}

\textbf{Intermediate Representation:}

\begin{quote}
possession as owner is obligated and (oblig(ejercicio\_buena\_fe) OR oblig(ejercicio\_mala\_fe)) IMPLIES oblig(ejercer\_prescripcion\_positiva)
\end{quote}

\textbf{Witness Clause:}

\begin{lstlisting}[style=witnessstyle, language=Witness]
clause tesis385312_principio_fundamental = oblig(posesion_como_propietario) AND (oblig(ejercicio_buena_fe) OR oblig(ejercicio_mala_fe)) IMPLIES oblig(ejercer_prescripcion_positiva);
\end{lstlisting}

\subsubsection{Clause: tesis385312\_posesion\_ambas\_fe}

\textit{Supporting logic: Possession as owner applies to both good faith and bad faith cases}

\textbf{Intermediate Representation:}

\begin{quote}
If possession as owner is obligated, then (good faith exercise is obligated or bad faith exercise is obligated)
\end{quote}

\textbf{Witness Clause:}

\begin{lstlisting}[style=witnessstyle, language=Witness]
clause tesis385312_posesion_ambas_fe = oblig(posesion_como_propietario) IMPLIES (oblig(ejercicio_buena_fe) OR oblig(ejercicio_mala_fe));
\end{lstlisting}

\subsubsection{Clause: tesis385312\_dominio\_economico}

\textit{Economic dominion: Must be dominator of the thing, commanding and enjoying it as owner}

\textbf{Intermediate Representation:}

\begin{quote}
possession as owner is obligated
\end{quote}

\textbf{Witness Clause:}

\begin{lstlisting}[style=witnessstyle, language=Witness]
clause tesis385312_dominio_economico = oblig(posesion_como_propietario) IMPLIES oblig(causa_generadora_probada);
\end{lstlisting}

\subsubsection{Clause: tesis385312\_posesion\_originaria}

\textit{Original possession: Must begin with cause different from derived possession}

\textbf{Intermediate Representation:}

\begin{quote}
possession as owner is obligated and not (temporal possession is obligated)
\end{quote}

\textbf{Witness Clause:}

\begin{lstlisting}[style=witnessstyle, language=Witness]
clause tesis385312_posesion_originaria = oblig(posesion_como_propietario) AND not(oblig(posesion_temporal)) IMPLIES oblig(causa_generadora_probada);
\end{lstlisting}

\newpage

\section{Thesis 385657: PRESCRIPCION POSITIVA, INTERRUPCION DE LA EN CASO DE (LEGISLACION DE JALISCO).}
\label{thesis:385657}

\subsection{Original Thesis Text}

\begin{quote}
                                                   Semanario Judicial de la Federación

                                                  Tesis

 Registro digital: 385657

 Instancia: Sala Auxiliar         Quinta Época                         Materia(s): Civil

                                  Fuente: Semanario Judicial de la     Tipo: Aislada
                                  Federación.
                                  Tomo CXIII, página 178


PRESCRIPCION POSITIVA, INTERRUPCION DE LA EN CASO DE (LEGISLACION DE JALISCO).


Si la demanda de reivindicación se notifica al poseedor después de transcurridos los cinco años que
para los poseedores de buena fe y aun los diez que para los de mala fe señala el Código Civil, es
indudable que no pudo interrumpirse la prescripción, pues ésta ya se había consumado.


Amparo civil directo 6712/49. Orozco Camarena Jovita y coagraviada. 9 de julio de 1952.
Unanimidad de cinco votos por lo que respecta al punto resolutivo de esta ejecutoria y por mayoría
de tres votos respecto de la parte considerativa. Disidentes: Angel González de la Vega y Gabriel
García Rojas. Ponente: Felipe Tena Ramírez.




                                             Pág. 1 de 1             Fecha de impresión 01/10/2025
  https://sjf2.scjn.gob.mx/detalle/tesis/385657

\end{quote}

\subsection{Formalization}

\subsubsection{Clause: tesis385657\_principio\_fundamental}

\textit{Thesis-specific principle: Adverse possession can be invoked by action or exception, but requires same proof elements}

\textbf{Intermediate Representation:}

\begin{quote}
possession as owner is obligated and (good faith exercise is obligated or bad faith exercise is obligated) and (oblig(posesion\_tiempo\_5\_anos) OR oblig(posesion\_tiempo\_10\_anos)) IMPLIES oblig(ejercer\_prescripcion\_positiva)
\end{quote}

\textbf{Witness Clause:}

\begin{lstlisting}[style=witnessstyle, language=Witness]
clause tesis385657_principio_fundamental = oblig(posesion_como_propietario) AND (oblig(ejercicio_buena_fe) OR oblig(ejercicio_mala_fe)) AND (oblig(posesion_tiempo_5_anos) OR oblig(posesion_tiempo_10_anos)) IMPLIES oblig(ejercer_prescripcion_positiva);
\end{lstlisting}

\subsubsection{Clause: tesis385657\_igualdad\_prueba}

\textit{Supporting logic: Both action and exception require same proof elements regardless of procedural route}

\textbf{Intermediate Representation:}

\begin{quote}
ejercer prescripcion positiva is obligated
\end{quote}

\textbf{Witness Clause:}

\begin{lstlisting}[style=witnessstyle, language=Witness]
clause tesis385657_igualdad_prueba = oblig(ejercer_prescripcion_positiva) IMPLIES oblig(posesion_como_propietario);
\end{lstlisting}

\subsubsection{Clause: tesis385657\_plazos\_temporales}

\textit{Time requirements: 5 years for good faith, 10 years for bad faith}

\textbf{Intermediate Representation:}

\begin{quote}
good faith exercise is obligated
\end{quote}

\textbf{Witness Clause:}

\begin{lstlisting}[style=witnessstyle, language=Witness]
clause tesis385657_plazos_temporales = oblig(ejercicio_buena_fe) IMPLIES oblig(posesion_tiempo_5_anos);
\end{lstlisting}

\subsubsection{Clause: tesis385657\_plazos\_mala\_fe}

\textbf{Intermediate Representation:}

\begin{quote}
bad faith exercise is obligated
\end{quote}

\textbf{Witness Clause:}

\begin{lstlisting}[style=witnessstyle, language=Witness]
clause tesis385657_plazos_mala_fe = oblig(ejercicio_mala_fe) IMPLIES oblig(posesion_tiempo_10_anos);
\end{lstlisting}

\subsubsection{Clause: tesis385657\_objetivo\_registral}

\textit{Target requirement: Action must be directed against registered owner}

\textbf{Intermediate Representation:}

\begin{quote}
ejercer prescripcion positiva is obligated
\end{quote}

\textbf{Witness Clause:}

\begin{lstlisting}[style=witnessstyle, language=Witness]
clause tesis385657_objetivo_registral = oblig(ejercer_prescripcion_positiva) IMPLIES oblig(titularidad_registral);
\end{lstlisting}

\newpage

\section{Thesis 800677: PRESCRIPCION ADQUISITIVA. NECESIDAD DE REVELAR LA CAUSA DE LA POSESION.}
\label{thesis:800677}

\subsection{Original Thesis Text}

\begin{quote}
                                                   Semanario Judicial de la Federación

                                                  Tesis

 Registro digital: 800677

 Instancia: Tercera Sala          Séptima Época                        Materia(s): Civil

                                  Fuente: Semanario Judicial de la     Tipo: Aislada
                                  Federación.
                                  Volumen 62, Cuarta Parte, página
                                  51


PRESCRIPCION ADQUISITIVA. NECESIDAD DE REVELAR LA CAUSA DE LA POSESION.


El actor en un juicio de prescripción positiva, debe revelar la causa de su posesión, aun en el caso
del poseedor de mala fe, porque es necesario que el juzgador conozca el hecho o acto generador
de la misma, para poder determinar la calidad de la posesión, si es en concepto de propietario,
originaria o derivada, de buena o mala fe, y para precisar el momento en que debe empezar a
contar el plazo de la prescripción.


Amparo directo 1617/72. Sergio Téllez Durazo. 15 de febrero de 1974. Cinco votos. Ponente: David
Franco Rodríguez.




                                             Pág. 1 de 1             Fecha de impresión 01/10/2025
  https://sjf2.scjn.gob.mx/detalle/tesis/800677

\end{quote}

\subsection{Formalization}

\subsubsection{Clause: tesis800677\_principio\_fundamental}

\textit{Thesis-specific principle: When defendant admits possession without title for more than 10 years with required qualities, is bad faith possessor and property is eligible for prescription}

\textbf{Intermediate Representation:}

\begin{quote}
If (posesion\_como\_propietario AND posesion\_pacifico\_continuo\_publico AND posesion\_tiempo\_10\_anos AND not(causa\_generadora\_probada)), then (ejercicio\_mala\_fe AND oblig(ejercer\_prescripcion\_positiva))
\end{quote}

\textbf{Witness Clause:}

\begin{lstlisting}[style=witnessstyle, language=Witness]
clause tesis800677_principio_fundamental = (posesion_como_propietario AND posesion_pacifico_continuo_publico AND posesion_tiempo_10_anos AND not(causa_generadora_probada)) IMPLIES (ejercicio_mala_fe AND oblig(ejercer_prescripcion_positiva));
\end{lstlisting}

\subsubsection{Clause: tesis800677\_no\_prueba\_causa}

\textit{Supporting logic: Does not need to prove generating cause of possession in trial when no title exists}

\textbf{Intermediate Representation:}

\begin{quote}
(not(causa\_generadora\_probada) AND posesion\_tiempo\_10\_anos) IMPLIES oblig(ejercer\_prescripcion\_positiva)
\end{quote}

\textbf{Witness Clause:}

\begin{lstlisting}[style=witnessstyle, language=Witness]
clause tesis800677_no_prueba_causa = (not(causa_generadora_probada) AND posesion_tiempo_10_anos) IMPLIES oblig(ejercer_prescripcion_positiva);
\end{lstlisting}

\subsubsection{Clause: tesis800677\_definicion\_titulo}

\textit{Title definition: In state law, title is understood as the generating cause of possession}

\textbf{Intermediate Representation:}

\begin{quote}
not (causa\_generadora\_probada) IMPLIES not(titularidad\_registral)
\end{quote}

\textbf{Witness Clause:}

\begin{lstlisting}[style=witnessstyle, language=Witness]
clause tesis800677_definicion_titulo = not(causa_generadora_probada) IMPLIES not(titularidad_registral);
\end{lstlisting}

\subsubsection{Clause: tesis800677\_apto\_prescripcion}

\textit{Prescription eligibility: If 10 years have passed and other elements are met, property is eligible for prescription}

\textbf{Intermediate Representation:}

\begin{quote}
(posesion\_tiempo\_10\_anos AND posesion\_como\_propietario AND posesion\_pacifico\_continuo\_publico) IMPLIES oblig(ejercer\_prescripcion\_positiva)
\end{quote}

\textbf{Witness Clause:}

\begin{lstlisting}[style=witnessstyle, language=Witness]
clause tesis800677_apto_prescripcion = (posesion_tiempo_10_anos AND posesion_como_propietario AND posesion_pacifico_continuo_publico) IMPLIES oblig(ejercer_prescripcion_positiva);
\end{lstlisting}

\newpage

\section{Thesis 800681: PRESCRIPCION ADQUISITIVA. POSESION EN CONCEPTO DE PROPIETARIO. tesis 220, página 637, bajo el rubro "PRESCRIPCION ADQUISITIVA. POSESION EN CONCEPTO}
\label{thesis:800681}

\subsection{Original Thesis Text}

\begin{quote}
                                                   Semanario Judicial de la Federación

                                                  Tesis

 Registro digital: 800681

 Instancia: Tercera Sala          Séptima Época                        Materia(s): Civil

                                  Fuente: Semanario Judicial de la     Tipo: Aislada
                                  Federación.
                                  Volumen 78, Cuarta Parte, página
                                  38


PRESCRIPCION ADQUISITIVA. POSESION EN CONCEPTO DE PROPIETARIO.


La exigencia del Código Civil para el Distrito y Territorios Federales y las legislaciones de los
Estados de la República que contienen disposiciones iguales, de poseer en concepto de propietario
para poder adquirir por prescripción, comprende no sólo los casos de buena fe, sino también el caso
de la posesión de mala fe, por lo que no basta la simple intención de poseer como dueño, sino que
es necesario probar la ejecución de actos o hechos susceptibles de ser apreciados por los sentidos,
que de manera indiscutible y objetiva demuestren que el poseedor es el dominador de la cosa, el
que manda en ella y la disfruta para sí, como dueño, en sentido económico, aun cuando carezca de
un título legítimo, frente a todo el mundo, y siempre que haya comenzado a poseer en virtud de una
causa diversa de la que origina la posesión derivada.


Amparo directo 3865/73. Rosaura Chanes de Dounce. 23 de junio de 1975. Cinco votos. Ponente:
Rafael Rojina Villegas. Secretario: Jaime M. Marroquín Zaleta

Séptima Epoca, Cuarta Parte:

Volumen 62, página 51. Amparo directo 1617/72. Sergio Téllez Durazo. 15 de febrero de 1974.
Cinco votos. Ponente: David Franco Rodríguez.

Véase: Apéndice al Semanario Judicial de la Federación 1917-1985, Cuarta Parte, Tercera Sala,
tesis 220, página 637, bajo el rubro "PRESCRIPCION ADQUISITIVA. POSESION EN CONCEPTO
DE PROPIETARIO.".




                                             Pág. 1 de 1             Fecha de impresión 01/10/2025
  https://sjf2.scjn.gob.mx/detalle/tesis/800681

\end{quote}

\subsection{Formalization}

\subsubsection{Clause: tesis800681\_principio\_fundamental}

\textit{Thesis-specific principle: Adverse possession can be invoked by action or exception, but requires same proof elements}

\textbf{Intermediate Representation:}

\begin{quote}
possession as owner is obligated and (good faith exercise is obligated or bad faith exercise is obligated) and (oblig(posesion\_tiempo\_5\_anos) OR oblig(posesion\_tiempo\_10\_anos)) IMPLIES oblig(ejercer\_prescripcion\_positiva)
\end{quote}

\textbf{Witness Clause:}

\begin{lstlisting}[style=witnessstyle, language=Witness]
clause tesis800681_principio_fundamental = oblig(posesion_como_propietario) AND (oblig(ejercicio_buena_fe) OR oblig(ejercicio_mala_fe)) AND (oblig(posesion_tiempo_5_anos) OR oblig(posesion_tiempo_10_anos)) IMPLIES oblig(ejercer_prescripcion_positiva);
\end{lstlisting}

\subsubsection{Clause: tesis800681\_igualdad\_prueba}

\textit{Supporting logic: Both action and exception require same proof elements regardless of procedural route}

\textbf{Intermediate Representation:}

\begin{quote}
ejercer prescripcion positiva is obligated
\end{quote}

\textbf{Witness Clause:}

\begin{lstlisting}[style=witnessstyle, language=Witness]
clause tesis800681_igualdad_prueba = oblig(ejercer_prescripcion_positiva) IMPLIES oblig(posesion_como_propietario);
\end{lstlisting}

\subsubsection{Clause: tesis800681\_plazos\_temporales}

\textit{Time requirements: 5 years for good faith, 10 years for bad faith}

\textbf{Intermediate Representation:}

\begin{quote}
good faith exercise is obligated
\end{quote}

\textbf{Witness Clause:}

\begin{lstlisting}[style=witnessstyle, language=Witness]
clause tesis800681_plazos_temporales = oblig(ejercicio_buena_fe) IMPLIES oblig(posesion_tiempo_5_anos);
\end{lstlisting}

\subsubsection{Clause: tesis800681\_plazos\_mala\_fe}

\textbf{Intermediate Representation:}

\begin{quote}
bad faith exercise is obligated
\end{quote}

\textbf{Witness Clause:}

\begin{lstlisting}[style=witnessstyle, language=Witness]
clause tesis800681_plazos_mala_fe = oblig(ejercicio_mala_fe) IMPLIES oblig(posesion_tiempo_10_anos);
\end{lstlisting}

\subsubsection{Clause: tesis800681\_objetivo\_registral}

\textit{Target requirement: Action must be directed against registered owner}

\textbf{Intermediate Representation:}

\begin{quote}
ejercer prescripcion positiva is obligated
\end{quote}

\textbf{Witness Clause:}

\begin{lstlisting}[style=witnessstyle, language=Witness]
clause tesis800681_objetivo_registral = oblig(ejercer_prescripcion_positiva) IMPLIES oblig(titularidad_registral);
\end{lstlisting}

\newpage

\section{Thesis 801244: PRESCRIPCION ADQUISITIVA. POSESION PARA LA.}
\label{thesis:801244}

\subsection{Original Thesis Text}

\begin{quote}
                                                   Semanario Judicial de la Federación

                                                  Tesis

 Registro digital: 801244

 Instancia: Tercera Sala          Sexta Época                           Materia(s): Civil

                                  Fuente: Semanario Judicial de la      Tipo: Aislada
                                  Federación.
                                  Volumen XII, Cuarta Parte, página
                                  148


PRESCRIPCION ADQUISITIVA. POSESION PARA LA.


Quien alega la usucapión debe manifestar el hecho que generó la posesión para que el juzgador
esté en aptitud de estimar si se trata de una posesión original o derivada. Es esencial conocer el
hecho o acto que engendró la posesión, incluso para computar el término legal que corresponda al
caso de que se trate de una posesión de buena o mala fe, si se adquirió por medios violentos o de
un delito.


Amparo directo 2733/57. Tomás Domínguez. 6 de junio de 1958. Unanimidad de cuatro votos.
Ponente: José Castro Estrada.




                                             Pág. 1 de 1              Fecha de impresión 01/10/2025
  https://sjf2.scjn.gob.mx/detalle/tesis/801244

\end{quote}

\subsection{Formalization}

\subsubsection{Clause: tesis801244\_principio\_fundamental}

\textit{Thesis-specific principle: Good faith requires sufficient title/generating cause, ignorance of title defects, and founded belief of ownership}

\textbf{Intermediate Representation:}

\begin{quote}
good faith exercise is obligated and not (conocimiento vicios is obligated)
\end{quote}

\textbf{Witness Clause:}

\begin{lstlisting}[style=witnessstyle, language=Witness]
clause tesis801244_principio_fundamental = oblig(ejercicio_buena_fe) IMPLIES oblig(causa_generadora_probada) AND not(oblig(conocimiento_vicios));
\end{lstlisting}

\subsubsection{Clause: tesis801244\_mala\_fe}

\textit{Supporting logic: Bad faith possessor has no title/generating cause OR knows title defects that prevent rightful possession}

\textbf{Intermediate Representation:}

\begin{quote}
If bad faith exercise is obligated, then (not (proven generating cause is obligated) or conocimiento vicios is obligated)
\end{quote}

\textbf{Witness Clause:}

\begin{lstlisting}[style=witnessstyle, language=Witness]
clause tesis801244_mala_fe = oblig(ejercicio_mala_fe) IMPLIES (not(oblig(causa_generadora_probada)) OR oblig(conocimiento_vicios));
\end{lstlisting}

\subsubsection{Clause: tesis801244\_requisitos\_buena\_fe}

\textit{Good faith requirements: Must have sufficient title and ignorance of defects}

\textbf{Intermediate Representation:}

\begin{quote}
good faith exercise is obligated
\end{quote}

\textbf{Witness Clause:}

\begin{lstlisting}[style=witnessstyle, language=Witness]
clause tesis801244_requisitos_buena_fe = oblig(ejercicio_buena_fe) IMPLIES oblig(causa_generadora_probada);
\end{lstlisting}

\subsubsection{Clause: tesis801244\_condiciones\_mala\_fe}

\textit{Bad faith conditions: No title or knowledge of defects}

\textbf{Intermediate Representation:}

\begin{quote}
(not(oblig(causa\_generadora\_probada)) OR oblig(conocimiento\_vicios)) IMPLIES oblig(ejercicio\_mala\_fe)
\end{quote}

\textbf{Witness Clause:}

\begin{lstlisting}[style=witnessstyle, language=Witness]
clause tesis801244_condiciones_mala_fe = (not(oblig(causa_generadora_probada)) OR oblig(conocimiento_vicios)) IMPLIES oblig(ejercicio_mala_fe);
\end{lstlisting}

\newpage

\section{Thesis 801247: PRESCRIPCION ADQUISITIVA. POSESION PARA LA.}
\label{thesis:801247}

\subsection{Original Thesis Text}

\begin{quote}
                                                   Semanario Judicial de la Federación

                                                  Tesis

 Registro digital: 801247

 Instancia: Tercera Sala          Sexta Época                           Materia(s): Civil

                                  Fuente: Semanario Judicial de la      Tipo: Aislada
                                  Federación.
                                  Volumen XII, Cuarta Parte, página
                                  148


PRESCRIPCION ADQUISITIVA. POSESION PARA LA.


Como conforme al artículo 806 del Código Civil, es poseedor de buena fe el que entra en la
posesión en virtud de un título suficiente para darle el derecho de poseer, quien afirma que es
poseedor de buena fe, debe demostrar ese título, o sea la causa generadora de la posesión.


Amparo directo 2733/57. Tomás Domínguez. 6 de junio de 1958. Unanimidad de cuatro votos.
Ponente: José Castro Estrada.




                                             Pág. 1 de 1              Fecha de impresión 01/10/2025
  https://sjf2.scjn.gob.mx/detalle/tesis/801247

\end{quote}

\subsection{Formalization}

\subsubsection{Clause: tesis801247\_principio\_fundamental}

\textit{Additional actions for this thesis Additional assets for proof requirements Faith exclusivity: Adverse possession is either good faith XOR bad faith Thesis-specific principle: Good faith adverse possession requires proving authenticity, date, and diligence}

\textbf{Intermediate Representation:}

\begin{quote}
If good faith exercise is obligated and adverse possession claim is claimed, then (autenticidad titulo is obligated and fecha fehaciente is obligated and diligencia adquirente is obligated)
\end{quote}

\textbf{Witness Clause:}

\begin{lstlisting}[style=witnessstyle, language=Witness]
clause tesis801247_principio_fundamental = oblig(ejercicio_buena_fe) AND claim(demanda_prescripcion) IMPLIES (oblig(autenticidad_titulo) AND oblig(fecha_fehaciente) AND oblig(diligencia_adquirente));
\end{lstlisting}

\subsubsection{Clause: tesis801247\_prueba\_autenticidad}

\textit{Proof requirements: Must prove title authenticity and date reliably}

\textbf{Intermediate Representation:}

\begin{quote}
If proven generating cause is obligated, then (autenticidad titulo is obligated and fecha fehaciente is obligated)
\end{quote}

\textbf{Witness Clause:}

\begin{lstlisting}[style=witnessstyle, language=Witness]
clause tesis801247_prueba_autenticidad = oblig(causa_generadora_probada) IMPLIES (oblig(autenticidad_titulo) AND oblig(fecha_fehaciente));
\end{lstlisting}

\subsubsection{Clause: tesis801247\_diligencia\_buena\_fe}

\textit{Diligence requirement: Good faith requires demonstrating due diligence in acquisition}

\textbf{Intermediate Representation:}

\begin{quote}
good faith exercise is obligated
\end{quote}

\textbf{Witness Clause:}

\begin{lstlisting}[style=witnessstyle, language=Witness]
clause tesis801247_diligencia_buena_fe = oblig(ejercicio_buena_fe) IMPLIES oblig(diligencia_adquirente);
\end{lstlisting}

\subsubsection{Clause: tesis801247\_prueba\_pagos}

\textit{Payment proof: Onerous contracts require proving payments made}

\textbf{Intermediate Representation:}

\begin{quote}
proven generating cause is obligated and contrato compraventa is obligated
\end{quote}

\textbf{Witness Clause:}

\begin{lstlisting}[style=witnessstyle, language=Witness]
clause tesis801247_prueba_pagos = oblig(causa_generadora_probada) AND oblig(contrato_compraventa) IMPLIES oblig(pagos_precio);
\end{lstlisting}

\subsubsection{Clause: tesis801247\_carga\_prueba}

\textit{Burden of proof: Claimant must prove generating cause with reliable evidence}

\textbf{Intermediate Representation:}

\begin{quote}
adverse possession claim is claimed
\end{quote}

\textbf{Witness Clause:}

\begin{lstlisting}[style=witnessstyle, language=Witness]
clause tesis801247_carga_prueba = claim(demanda_prescripcion) IMPLIES oblig(causa_generadora_revelada);
\end{lstlisting}

\subsubsection{Clause: tesis801247\_falta\_diligencia}

\textit{Insufficient diligence consequence: Lack of diligence means bad faith (10-year period)}

\textbf{Intermediate Representation:}

\begin{quote}
not (diligencia adquirente is obligated)
\end{quote}

\textbf{Witness Clause:}

\begin{lstlisting}[style=witnessstyle, language=Witness]
clause tesis801247_falta_diligencia = not(oblig(diligencia_adquirente)) IMPLIES oblig(ejercicio_mala_fe);
\end{lstlisting}

\newpage

\section{Thesis 804158: PRESCRIPCION POSITIVA. POSESION DE BUENA FE.}
\label{thesis:804158}

\subsection{Original Thesis Text}

\begin{quote}
                                                   Semanario Judicial de la Federación

                                                  Tesis

 Registro digital: 804158

 Instancia: Tercera Sala          Quinta Época                         Materia(s): Civil

                                  Fuente: Semanario Judicial de la     Tipo: Aislada
                                  Federación.
                                  Tomo CXXIV, página 1107


PRESCRIPCION POSITIVA. POSESION DE BUENA FE.


Si no se prueba que el adquirente de un bien lo haya adquirido conociendo los vicios de origen de
los que aparecían como anteriores propietarios, debe considerársele como poseedor de buena fe
para los efectos de la prescripción adquisitiva, y con justo título.


Amparo civil directo 1837/54. Evia viuda de Valencia Manuela. 22 de junio de 1955. Mayoría de tres
votos. Disidentes: Gabriel García Rojas y Mariano Ramírez Vázquez. Ponente: Gilberto Valenzuela.




                                             Pág. 1 de 1             Fecha de impresión 01/10/2025
  https://sjf2.scjn.gob.mx/detalle/tesis/804158

\end{quote}

\subsection{Formalization}

\subsubsection{Clause: tesis804158\_principio\_fundamental}

\textit{Faith exclusivity: Possession is either good faith XOR bad faith (mutually exclusive) Thesis-specific principle: Unregistered contracts cannot establish adverse possession against registered owners}

\textbf{Intermediate Representation:}

\begin{quote}
registered ownership is obligated and unregistered contract is obligated
\end{quote}

\textbf{Witness Clause:}

\begin{lstlisting}[style=witnessstyle, language=Witness]
clause tesis804158_principio_fundamental = oblig(titularidad_registral) AND oblig(contrato_no_inscrito) IMPLIES not(claim(demanda_prescripcion));
\end{lstlisting}

\subsubsection{Clause: tesis804158\_proteccion\_registral}

\textit{Supporting legal logic}

\textbf{Intermediate Representation:}

\begin{quote}
registered ownership is obligated
\end{quote}

\textbf{Witness Clause:}

\begin{lstlisting}[style=witnessstyle, language=Witness]
clause tesis804158_proteccion_registral = oblig(titularidad_registral) IMPLIES claim(oponibilidad_terceros);
\end{lstlisting}

\subsubsection{Clause: tesis804158\_limitacion\_no\_inscripcion}

\textbf{Intermediate Representation:}

\begin{quote}
unregistered contract is obligated
\end{quote}

\textbf{Witness Clause:}

\begin{lstlisting}[style=witnessstyle, language=Witness]
clause tesis804158_limitacion_no_inscripcion = oblig(contrato_no_inscrito) IMPLIES not(oblig(oponibilidad_terceros));
\end{lstlisting}

\subsubsection{Clause: tesis804158\_consecuencia\_logica}

\textbf{Intermediate Representation:}

\begin{quote}
not (opposability to third parties is obligated)
\end{quote}

\textbf{Witness Clause:}

\begin{lstlisting}[style=witnessstyle, language=Witness]
clause tesis804158_consecuencia_logica = not(oblig(oponibilidad_terceros)) IMPLIES not(claim(demanda_prescripcion));
\end{lstlisting}

\newpage

\section{Thesis 804320: PRESCRIPCION POSITIVA, BUENA FE EN LA POSESION PARA ADQUIRIR POR (LEGISLACION DE JALISCO).}
\label{thesis:804320}

\subsection{Original Thesis Text}

\begin{quote}
                                                   Semanario Judicial de la Federación

                                                  Tesis

 Registro digital: 804320

 Instancia: Tercera Sala          Quinta Época                         Materia(s): Civil

                                  Fuente: Semanario Judicial de la     Tipo: Aislada
                                  Federación.
                                  Tomo CXXII, página 1227


PRESCRIPCION POSITIVA, BUENA FE EN LA POSESION PARA ADQUIRIR POR
(LEGISLACION DE JALISCO).


De conformidad con el artículo 1082 del Código Civil, aun cuando a los pocos días de celebrarse un
contrato de compraventa, el comprador se percate de que el vendedor vendió una cosa que no era
de su propiedad, basta que en el momento de la celebración del contrato, el adquirente haya
actuado de buena fe, para que toda su posesión tenga que estimarse, para los efectos de la
usucapión, de la misma clase, y toca a la otra parte demostrar que se contrató de mala fe desde el
principio, esto es, que desde el momento mismo de la celebración del contrato el comprador supo
que el vendedor disponía de bienes que no eran de su propiedad.


Amparo civil directo 3667/53. Compañía Jalisciense Urbanizadora, S.A. 18 de noviembre de 1954.
Mayoría de cuatro votos. Disidente: Gabriel García Rojas. Ponente: Vicente Santos Guajardo.




                                             Pág. 1 de 1             Fecha de impresión 01/10/2025
  https://sjf2.scjn.gob.mx/detalle/tesis/804320

\end{quote}

\subsection{Formalization}

\subsubsection{Clause: tesis804320\_principio\_fundamental}

\textit{Thesis-specific principle: Possession must be "in concept of owner" (peaceful, continuous, public) for both good and bad faith}

\textbf{Intermediate Representation:}

\begin{quote}
possession as owner is obligated and (oblig(ejercicio\_buena\_fe) OR oblig(ejercicio\_mala\_fe)) IMPLIES oblig(ejercer\_prescripcion\_positiva)
\end{quote}

\textbf{Witness Clause:}

\begin{lstlisting}[style=witnessstyle, language=Witness]
clause tesis804320_principio_fundamental = oblig(posesion_como_propietario) AND (oblig(ejercicio_buena_fe) OR oblig(ejercicio_mala_fe)) IMPLIES oblig(ejercer_prescripcion_positiva);
\end{lstlisting}

\subsubsection{Clause: tesis804320\_posesion\_ambas\_fe}

\textit{Supporting logic: Possession as owner applies to both good faith and bad faith cases}

\textbf{Intermediate Representation:}

\begin{quote}
If possession as owner is obligated, then (good faith exercise is obligated or bad faith exercise is obligated)
\end{quote}

\textbf{Witness Clause:}

\begin{lstlisting}[style=witnessstyle, language=Witness]
clause tesis804320_posesion_ambas_fe = oblig(posesion_como_propietario) IMPLIES (oblig(ejercicio_buena_fe) OR oblig(ejercicio_mala_fe));
\end{lstlisting}

\subsubsection{Clause: tesis804320\_dominio\_economico}

\textit{Economic dominion: Must be dominator of the thing, commanding and enjoying it as owner}

\textbf{Intermediate Representation:}

\begin{quote}
possession as owner is obligated
\end{quote}

\textbf{Witness Clause:}

\begin{lstlisting}[style=witnessstyle, language=Witness]
clause tesis804320_dominio_economico = oblig(posesion_como_propietario) IMPLIES oblig(causa_generadora_probada);
\end{lstlisting}

\subsubsection{Clause: tesis804320\_posesion\_originaria}

\textit{Original possession: Must begin with cause different from derived possession}

\textbf{Intermediate Representation:}

\begin{quote}
possession as owner is obligated and not (temporal possession is obligated)
\end{quote}

\textbf{Witness Clause:}

\begin{lstlisting}[style=witnessstyle, language=Witness]
clause tesis804320_posesion_originaria = oblig(posesion_como_propietario) AND not(oblig(posesion_temporal)) IMPLIES oblig(causa_generadora_probada);
\end{lstlisting}

\newpage

\section{Thesis 804441: PRESCRIPCION POSITIVA (LEGISLACION DE JALISCO).}
\label{thesis:804441}

\subsection{Original Thesis Text}

\begin{quote}
                                                   Semanario Judicial de la Federación

                                                  Tesis

 Registro digital: 804441

 Instancia: Sala Auxiliar         Quinta Época                         Materia(s): Civil

                                  Fuente: Semanario Judicial de la     Tipo: Aislada
                                  Federación.
                                  Tomo CXVI, página 829


PRESCRIPCION POSITIVA (LEGISLACION DE JALISCO).


Si la quejosa no demostró que fuera poseedora del inmueble, de buena fe, a virtud de un título
suficiente para darle derecho a poseer, como lo determina el artículo 849 del Código Civil, sino que
sólo demostró que se encontraba unida en matrimonio eclesiástico con el autor de la sucesión, por
el cual motivo debe estimarse que el inmueble de referencia ha venido estando en su poder en
virtud de la situación de dependencia en que se encontraba con respecto del propietario del mismo,
atentos los términos del artículo 836 del ordenamiento citado, no puede considerarse a la quejosa
como poseedora del mencionado inmueble y menos que la tenencia del mismo pueda constituir una
posesión apta para prescribir.


Amparo civil directo 443/52. Velasco Cecilia R. 14 de octubre de 1952. Unanimidad de cinco votos.
La publicación no menciona el nombre del ponente.




                                             Pág. 1 de 1             Fecha de impresión 01/10/2025
  https://sjf2.scjn.gob.mx/detalle/tesis/804441

\end{quote}

\subsection{Formalization}

\subsubsection{Clause: tesis804441\_principio\_fundamental}

\textit{Thesis-specific principle: Adverse possession can be invoked by action or exception, but requires same proof elements}

\textbf{Intermediate Representation:}

\begin{quote}
possession as owner is obligated and (good faith exercise is obligated or bad faith exercise is obligated) and (oblig(posesion\_tiempo\_5\_anos) OR oblig(posesion\_tiempo\_10\_anos)) IMPLIES oblig(ejercer\_prescripcion\_positiva)
\end{quote}

\textbf{Witness Clause:}

\begin{lstlisting}[style=witnessstyle, language=Witness]
clause tesis804441_principio_fundamental = oblig(posesion_como_propietario) AND (oblig(ejercicio_buena_fe) OR oblig(ejercicio_mala_fe)) AND (oblig(posesion_tiempo_5_anos) OR oblig(posesion_tiempo_10_anos)) IMPLIES oblig(ejercer_prescripcion_positiva);
\end{lstlisting}

\subsubsection{Clause: tesis804441\_igualdad\_prueba}

\textit{Supporting logic: Both action and exception require same proof elements regardless of procedural route}

\textbf{Intermediate Representation:}

\begin{quote}
ejercer prescripcion positiva is obligated
\end{quote}

\textbf{Witness Clause:}

\begin{lstlisting}[style=witnessstyle, language=Witness]
clause tesis804441_igualdad_prueba = oblig(ejercer_prescripcion_positiva) IMPLIES oblig(posesion_como_propietario);
\end{lstlisting}

\subsubsection{Clause: tesis804441\_plazos\_temporales}

\textit{Time requirements: 5 years for good faith, 10 years for bad faith}

\textbf{Intermediate Representation:}

\begin{quote}
good faith exercise is obligated
\end{quote}

\textbf{Witness Clause:}

\begin{lstlisting}[style=witnessstyle, language=Witness]
clause tesis804441_plazos_temporales = oblig(ejercicio_buena_fe) IMPLIES oblig(posesion_tiempo_5_anos);
\end{lstlisting}

\subsubsection{Clause: tesis804441\_plazos\_mala\_fe}

\textbf{Intermediate Representation:}

\begin{quote}
bad faith exercise is obligated
\end{quote}

\textbf{Witness Clause:}

\begin{lstlisting}[style=witnessstyle, language=Witness]
clause tesis804441_plazos_mala_fe = oblig(ejercicio_mala_fe) IMPLIES oblig(posesion_tiempo_10_anos);
\end{lstlisting}

\subsubsection{Clause: tesis804441\_objetivo\_registral}

\textit{Target requirement: Action must be directed against registered owner}

\textbf{Intermediate Representation:}

\begin{quote}
ejercer prescripcion positiva is obligated
\end{quote}

\textbf{Witness Clause:}

\begin{lstlisting}[style=witnessstyle, language=Witness]
clause tesis804441_objetivo_registral = oblig(ejercer_prescripcion_positiva) IMPLIES oblig(titularidad_registral);
\end{lstlisting}

\newpage

\section{Thesis 806970: PRESCRIPCION ADQUISITIVA.}
\label{thesis:806970}

\subsection{Original Thesis Text}

\begin{quote}
                                                   Semanario Judicial de la Federación

                                                  Tesis

 Registro digital: 806970

 Instancia: Segunda Sala           Quinta Época                           Materia(s): Civil

                                   Fuente: Semanario Judicial de la       Tipo: Aislada
                                   Federación.
                                   Tomo LXXX, página 1542


PRESCRIPCION ADQUISITIVA.


Es infundado el agravio que se refiere a la falta de aplicación de los artículos 5o., 73, fracciones IX,
XVIII, 78 y 154 de la Ley de Amparo y 1149 del Código Civil, porque según el recurrente estaba
favorecido por la prescripción adquisitiva, por el sólo transcurso de cinco años; porque el inferior no
tenía motivo legal para hacer en el juicio de garantías, declaración alguna acerca de si se operó o
no, en favor de tercero perjudicado, la prescripción adquisitiva, pues ésta es materia de acción o de
excepción entre un juicio contradictorio entre las partes, carácter que no tiene el juicio de garantías,
en que se encausan actos de autoridad. Esta clase de prescripción se opera con una posesión con
justo título, quieta y pacífica, por determinado lapso pero la prueba de sus elementos debe
acreditarse en un juicio contradictorio, que tenga por objeto la discusión del dominio de la cosa y
que culmine con un fallo que la eleve a la categoría de título, para quien la invoca. En el juicio de
garantías no es jurídico contravenir ese dominio entre reivindicador y reivindicante, porque la
materia de este juicio está limitada por la ley constitucional, a resolver en cada caso concreto, sobre
la legalidad de un acto de autoridad, para destruirlo como medio eficaz de mantener al quejoso en el
goce de la garantía violada. No tiene aplicación en el caso, la ejecutoria publicada en la página 112
del Tomo XXIX, del Semanario Judicial de la Federación, que invoca el recurrente y que se refiere al
plazo que la ley concede para que la prescripción se verifique en favor del que posee de buena fe,
porque en ese asunto, el amparo fue solicitado por la misma deudora que fue parte en el juicio en
que se realizó el remate y después quiso destruir, y en el presente asunto es un tercero extraño al
procedimiento administrativo en el cual no fue oído. Además no se trata de desposeimiento alguno
que se hubiera llevado a efecto con perjuicio del recurrente, sino de la falta de citación en el juicio
economicocoactivo y sus consecuencias, especialmente, la cancelación del crédito hipotecario, sin
audiencia alguna, rédito constituido a favor de la quejosa, todo esto sin tomar en cuenta que no se
ha allegado la prescripción y no pudo haber sido motivo de su estudio por parte del inferior.


Amparo administrativo en revisión 10052/44. "Las Filipinas", S. A. 25 de abril de 1944. Unanimidad
de cuatro votos. Ausente: Franco Carreño. La publicación no menciona el nombre del ponente.




                                              Pág. 1 de 2             Fecha de impresión 01/10/2025
  https://sjf2.scjn.gob.mx/detalle/tesis/806970
                                                Semanario Judicial de la Federación




                                           Pág. 2 de 2       Fecha de impresión 01/10/2025
https://sjf2.scjn.gob.mx/detalle/tesis/806970

\end{quote}

\subsection{Formalization}

\subsubsection{Clause: tesis806970\_principio\_fundamental}

\textit{Thesis-specific principle: Just title not required for adverse possession, but must reveal origin of possession and affirm possession as owner}

\textbf{Intermediate Representation:}

\begin{quote}
ejercer prescripcion positiva is obligated and possession as owner is obligated
\end{quote}

\textbf{Witness Clause:}

\begin{lstlisting}[style=witnessstyle, language=Witness]
clause tesis806970_principio_fundamental = oblig(ejercer_prescripcion_positiva) IMPLIES oblig(causa_generadora_revelada) AND oblig(posesion_como_propietario);
\end{lstlisting}

\subsubsection{Clause: tesis806970\_determinacion\_judicial}

\textit{Supporting logic: Judge needs generating fact/act to determine possession quality and timing}

\textbf{Intermediate Representation:}

\begin{quote}
If causa generadora revelada is obligated, then (possession as owner is obligated and (good faith exercise is obligated and not (bad faith exercise is obligated)))
\end{quote}

\textbf{Witness Clause:}

\begin{lstlisting}[style=witnessstyle, language=Witness]
clause tesis806970_determinacion_judicial = oblig(causa_generadora_revelada) IMPLIES (oblig(posesion_como_propietario) AND (oblig(ejercicio_buena_fe) AND not(oblig(ejercicio_mala_fe))));
\end{lstlisting}

\subsubsection{Clause: tesis806970\_tipo\_posesion}

\textit{Possession type determination: Can determine if original or derived possession}

\textbf{Intermediate Representation:}

\begin{quote}
If causa generadora revelada is obligated, then (possession as owner is obligated and not (temporal possession is obligated))
\end{quote}

\textbf{Witness Clause:}

\begin{lstlisting}[style=witnessstyle, language=Witness]
clause tesis806970_tipo_posesion = oblig(causa_generadora_revelada) IMPLIES (oblig(posesion_como_propietario) AND not(oblig(posesion_temporal)));
\end{lstlisting}

\subsubsection{Clause: tesis806970\_determinacion\_temporal}

\textit{Timing requirement: Must establish when prescription period starts counting}

\textbf{Intermediate Representation:}

\begin{quote}
If causa generadora revelada is obligated, then (posesion tiempo 5 anos is obligated or posesion tiempo 10 anos is obligated)
\end{quote}

\textbf{Witness Clause:}

\begin{lstlisting}[style=witnessstyle, language=Witness]
clause tesis806970_determinacion_temporal = oblig(causa_generadora_revelada) IMPLIES (oblig(posesion_tiempo_5_anos) OR oblig(posesion_tiempo_10_anos));
\end{lstlisting}

\newpage

\section{Thesis 811551: THESIS 811551}
\label{thesis:811551}

\subsection{Original Thesis Text}

\begin{quote}
                                                   Semanario Judicial de la Federación

                                                  Tesis

 Registro digital: 811551

 Instancia: Pleno                 Quinta Época                         Materia(s): Civil

                                  Fuente: Semanario Judicial de la     Tipo: Aislada
                                  Federación.
                                  Tomo II, página 449


PRESCRIPCION POSITIVA.


Para que sea eficaz, es indispensable que concurran estos elementos: que la posesión sea pacífica,
continua, pública, de buena fe y basada en justo título. Estos dos últimos requisitos son
innecesarios, cuando la posesión data de más de treinta años.


Amparo civil interpuesto directamente ante la Corte. González Angel. 11 de febrero de 1918.
Unanimidad de once votos. La publicación no menciona el nombre del ponente.




                                             Pág. 1 de 1             Fecha de impresión 01/10/2025
  https://sjf2.scjn.gob.mx/detalle/tesis/811551

\end{quote}

\subsection{Formalization}

\subsubsection{Clause: tesis811551\_principio\_fundamental}

\textit{Additional actions for this thesis Additional assets for proof requirements Faith exclusivity: Adverse possession is either good faith XOR bad faith Thesis-specific principle: Good faith adverse possession requires proving authenticity, date, and diligence}

\textbf{Intermediate Representation:}

\begin{quote}
If good faith exercise is obligated and adverse possession claim is claimed, then (autenticidad titulo is obligated and fecha fehaciente is obligated and diligencia adquirente is obligated)
\end{quote}

\textbf{Witness Clause:}

\begin{lstlisting}[style=witnessstyle, language=Witness]
clause tesis811551_principio_fundamental = oblig(ejercicio_buena_fe) AND claim(demanda_prescripcion) IMPLIES (oblig(autenticidad_titulo) AND oblig(fecha_fehaciente) AND oblig(diligencia_adquirente));
\end{lstlisting}

\subsubsection{Clause: tesis811551\_prueba\_autenticidad}

\textit{Proof requirements: Must prove title authenticity and date reliably}

\textbf{Intermediate Representation:}

\begin{quote}
If proven generating cause is obligated, then (autenticidad titulo is obligated and fecha fehaciente is obligated)
\end{quote}

\textbf{Witness Clause:}

\begin{lstlisting}[style=witnessstyle, language=Witness]
clause tesis811551_prueba_autenticidad = oblig(causa_generadora_probada) IMPLIES (oblig(autenticidad_titulo) AND oblig(fecha_fehaciente));
\end{lstlisting}

\subsubsection{Clause: tesis811551\_diligencia\_buena\_fe}

\textit{Diligence requirement: Good faith requires demonstrating due diligence in acquisition}

\textbf{Intermediate Representation:}

\begin{quote}
good faith exercise is obligated
\end{quote}

\textbf{Witness Clause:}

\begin{lstlisting}[style=witnessstyle, language=Witness]
clause tesis811551_diligencia_buena_fe = oblig(ejercicio_buena_fe) IMPLIES oblig(diligencia_adquirente);
\end{lstlisting}

\subsubsection{Clause: tesis811551\_prueba\_pagos}

\textit{Payment proof: Onerous contracts require proving payments made}

\textbf{Intermediate Representation:}

\begin{quote}
proven generating cause is obligated and contrato compraventa is obligated
\end{quote}

\textbf{Witness Clause:}

\begin{lstlisting}[style=witnessstyle, language=Witness]
clause tesis811551_prueba_pagos = oblig(causa_generadora_probada) AND oblig(contrato_compraventa) IMPLIES oblig(pagos_precio);
\end{lstlisting}

\subsubsection{Clause: tesis811551\_carga\_prueba}

\textit{Burden of proof: Claimant must prove generating cause with reliable evidence}

\textbf{Intermediate Representation:}

\begin{quote}
adverse possession claim is claimed
\end{quote}

\textbf{Witness Clause:}

\begin{lstlisting}[style=witnessstyle, language=Witness]
clause tesis811551_carga_prueba = claim(demanda_prescripcion) IMPLIES oblig(causa_generadora_revelada);
\end{lstlisting}

\subsubsection{Clause: tesis811551\_falta\_diligencia}

\textit{Insufficient diligence consequence: Lack of diligence means bad faith (10-year period)}

\textbf{Intermediate Representation:}

\begin{quote}
not (diligencia adquirente is obligated)
\end{quote}

\textbf{Witness Clause:}

\begin{lstlisting}[style=witnessstyle, language=Witness]
clause tesis811551_falta_diligencia = not(oblig(diligencia_adquirente)) IMPLIES oblig(ejercicio_mala_fe);
\end{lstlisting}

\newpage

\section{Thesis 812166: PRESCRIPCION POSITIVA, LA PRUEBA DE LA CAUSA GENERADORA DE LA POSESION ES ELEMENTO DE LA PRETENSION, AUN CUANDO LA LEY NO EXIJA JUSTO TITULO.}
\label{thesis:812166}

\subsection{Original Thesis Text}

\begin{quote}
                                                   Semanario Judicial de la Federación

                                                  Tesis

 Registro digital: 812166

 Instancia: Tribunales            Octava Época                          Materia(s): Civil
 Colegiados de Circuito

 Tesis: 16                        Fuente: Informes.                     Tipo: Aislada
                                  Informe 1989, Parte III, página 272


PRESCRIPCION POSITIVA, LA PRUEBA DE LA CAUSA GENERADORA DE LA POSESION ES
ELEMENTO DE LA PRETENSION, AUN CUANDO LA LEY NO EXIJA JUSTO TITULO.


Aun cuando se pretende probar que la posesión de que se goza sobre el inmueble disputado es en
concepto de dueño o propietario, tratando al efecto de acreditar que durante el tiempo que se ha
ocupado se han realizado actos o hechos que demuestran que se es dominador de la cosa, el que
manda en ella y la disfruta para sí, como dueño en sentido económico, al no acreditarse la causa
que dio origen a la posesión, no se puede admitir que se satisfacen los requisitos para obtener la
propiedad por prescripción, debiéndose aclarar que la causa generadora de la posesión que debe
revelarse y probarse no es concepto jurídico que se confunda con el justo título, (entendiendo por tal
el acto que es bastante para adquirir el dominio), pues aun cuando ambas figuras se pueden
presentar en algunos supuestos en que se posea un bien susceptible de adquirirse por prescripción
(cuando la posesión es de buena fe) no en todas las hipótesis pueden actualizarse (como lo es la
posesión de mala fe), pudiendo así decir que toda posesión susceptible de crear la prescripción
supone una causa generadora, pero no toda posesión susceptible de crear la prescripción supone
un justo título.

TERCER TRIBUNAL COLEGIADO EN MATERIA CIVIL DEL PRIMER CIRCUITO.


Amparo directo 1068/89. Francisca Rodríguez Méndez. 6 de abril de 1989. Ponente: Manuel
Ernesto Saloma Vera.

Amparo directo 3328/88. Fernando Cubero Flores. 24 de noviembre de 1988. Unanimidad de votos.
Ponente: José Rojas Aja. Secretario Francisco Sánchez Planells.




                                             Pág. 1 de 2             Fecha de impresión 01/10/2025
  https://sjf2.scjn.gob.mx/detalle/tesis/812166
                                                Semanario Judicial de la Federación




                                           Pág. 2 de 2       Fecha de impresión 01/10/2025
https://sjf2.scjn.gob.mx/detalle/tesis/812166

\end{quote}

\subsection{Formalization}

\subsubsection{Clause: tesis812166\_principio\_fundamental}

\textit{Faith exclusivity: Possession is either good faith XOR bad faith (mutually exclusive) Thesis-specific principle: Unregistered contracts cannot establish adverse possession against registered owners}

\textbf{Intermediate Representation:}

\begin{quote}
registered ownership is obligated and unregistered contract is obligated
\end{quote}

\textbf{Witness Clause:}

\begin{lstlisting}[style=witnessstyle, language=Witness]
clause tesis812166_principio_fundamental = oblig(titularidad_registral) AND oblig(contrato_no_inscrito) IMPLIES not(claim(demanda_prescripcion));
\end{lstlisting}

\subsubsection{Clause: tesis812166\_proteccion\_registral}

\textit{Supporting legal logic}

\textbf{Intermediate Representation:}

\begin{quote}
registered ownership is obligated
\end{quote}

\textbf{Witness Clause:}

\begin{lstlisting}[style=witnessstyle, language=Witness]
clause tesis812166_proteccion_registral = oblig(titularidad_registral) IMPLIES claim(oponibilidad_terceros);
\end{lstlisting}

\subsubsection{Clause: tesis812166\_limitacion\_no\_inscripcion}

\textbf{Intermediate Representation:}

\begin{quote}
unregistered contract is obligated
\end{quote}

\textbf{Witness Clause:}

\begin{lstlisting}[style=witnessstyle, language=Witness]
clause tesis812166_limitacion_no_inscripcion = oblig(contrato_no_inscrito) IMPLIES not(oblig(oponibilidad_terceros));
\end{lstlisting}

\subsubsection{Clause: tesis812166\_consecuencia\_logica}

\textbf{Intermediate Representation:}

\begin{quote}
not (opposability to third parties is obligated)
\end{quote}

\textbf{Witness Clause:}

\begin{lstlisting}[style=witnessstyle, language=Witness]
clause tesis812166_consecuencia_logica = not(oblig(oponibilidad_terceros)) IMPLIES not(claim(demanda_prescripcion));
\end{lstlisting}

\newpage

\section{Thesis 815363: PRESCRIPCION ADQUISITIVA, NECESIDAD DE REVELAR LA CAUSA DE LA POSESION. REVELAR LA CAUSA DE LA POSESION".}
\label{thesis:815363}

\subsection{Original Thesis Text}

\begin{quote}
                                                   Semanario Judicial de la Federación

                                                  Tesis

 Registro digital: 815363

 Instancia: Tercera Sala          Sexta Época                          Materia(s): Civil

                                  Fuente: Informes.                    Tipo: Aislada
                                  Informe 1964, página 49


PRESCRIPCION ADQUISITIVA, NECESIDAD DE REVELAR LA CAUSA DE LA POSESION.


El actor en un juicio de prescripción positiva debe revelar la causa de su posesión, o sea el hecho o
acto generador de la misma tanto para que el juzgador pueda determinar la calidad de la posesión,
originaria o derivada, de buena o mala fe, cuanto para que esté en aptitud de contar el plazo de
prescripción. Esta tesis confirma el criterio de esta Tercera Sala contenido en las ejecutorias
pronunciadas en los siguientes directos: 7673/58; 7523/58; 2733/57 y 7385/56.



Amparo directo 2070/61. Magaña Méndez Heliodoro. 17 de junio de 1964. Unanimidad de cinco
votos. Relator: Rafael Rojina Villegas.

Véase:

Apéndice al Semanario Judicial de la Federación 1917-1995, Sexta Epoca, Tomo IV, parte SCJN,
tesis 316, página 213, Jurisprudencia de rubro "PRESCRIPCION ADQUISITIVA. NECESIDAD DE
REVELAR LA CAUSA DE LA POSESION".




                                             Pág. 1 de 1            Fecha de impresión 01/10/2025
  https://sjf2.scjn.gob.mx/detalle/tesis/815363

\end{quote}

\subsection{Formalization}

\subsubsection{Clause: tesis815363\_principio\_fundamental}

\textit{Thesis-specific principle: Good faith possessors require just title as generating cause of possession for adverse possession to operate}

\textbf{Intermediate Representation:}

\begin{quote}
good faith exercise is obligated and possession as owner is obligated
\end{quote}

\textbf{Witness Clause:}

\begin{lstlisting}[style=witnessstyle, language=Witness]
clause tesis815363_principio_fundamental = oblig(ejercicio_buena_fe) AND oblig(posesion_como_propietario) IMPLIES oblig(causa_generadora_probada);
\end{lstlisting}

\subsubsection{Clause: tesis815363\_revelar\_causa}

\textit{Supporting logic: Must reveal and fully prove the generating cause of possession}

\textbf{Intermediate Representation:}

\begin{quote}
ejercer prescripcion positiva is obligated
\end{quote}

\textbf{Witness Clause:}

\begin{lstlisting}[style=witnessstyle, language=Witness]
clause tesis815363_revelar_causa = oblig(ejercer_prescripcion_positiva) IMPLIES oblig(causa_generadora_revelada);
\end{lstlisting}

\subsubsection{Clause: tesis815363\_determinacion\_judicial}

\textit{Judicial determination: Judge needs generating fact/act to determine possession quality and timing}

\textbf{Intermediate Representation:}

\begin{quote}
If causa generadora revelada is obligated, then (possession as owner is obligated and (good faith exercise is obligated and not (bad faith exercise is obligated)))
\end{quote}

\textbf{Witness Clause:}

\begin{lstlisting}[style=witnessstyle, language=Witness]
clause tesis815363_determinacion_judicial = oblig(causa_generadora_revelada) IMPLIES (oblig(posesion_como_propietario) AND (oblig(ejercicio_buena_fe) AND not(oblig(ejercicio_mala_fe))));
\end{lstlisting}

\subsubsection{Clause: tesis815363\_determinacion\_temporal}

\textit{Timing requirement: Must establish when prescription period begins}

\textbf{Intermediate Representation:}

\begin{quote}
If causa generadora revelada is obligated, then (posesion tiempo 5 anos is obligated or posesion tiempo 10 anos is obligated)
\end{quote}

\textbf{Witness Clause:}

\begin{lstlisting}[style=witnessstyle, language=Witness]
clause tesis815363_determinacion_temporal = oblig(causa_generadora_revelada) IMPLIES (oblig(posesion_tiempo_5_anos) OR oblig(posesion_tiempo_10_anos));
\end{lstlisting}

\newpage

\section{Thesis 2000627: PRESCRIPCIÓN POSITIVA. EL TÍTULO EXHIBIDO COMO CAUSA GENERADORA DE LA POSESIÓN DE BUENA FE DEBE EXPRESAR LA TRANSMISIÓN SUCESIVA DEL DOMINIO DEL}
\label{thesis:2000627}

\subsection{Original Thesis Text}

\begin{quote}
                                                   Semanario Judicial de la Federación

                                                 Tesis

 Registro digital: 2000627

 Instancia: Tribunales              Décima Época                           Materia(s): Civil
 Colegiados de Circuito

 Tesis: I.3o.C.9 C (10a.)           Fuente: Semanario Judicial de la       Tipo: Aislada
                                    Federación y su Gaceta.
                                    Libro VII, Abril de 2012, Tomo 2,
                                    página 1836


PRESCRIPCIÓN POSITIVA. EL TÍTULO EXHIBIDO COMO CAUSA GENERADORA DE LA
POSESIÓN DE BUENA FE DEBE EXPRESAR LA TRANSMISIÓN SUCESIVA DEL DOMINIO DEL
INMUEBLE, PARTIENDO DE QUIEN ESTÉ INSCRITO COMO PROPIETARIO EN EL REGISTRO
PÚBLICO DE LA PROPIEDAD.


La inscripción registral surte efecto contra terceros porque da publicidad al acto inscrito. Por tanto, si
ante el Registro Público de la Propiedad se encuentra la inscripción de la propiedad en favor de una
persona determinada, y sobre ese inmueble el actor promueve la acción de prescripción positiva de
buena fe, el título traslativo de dominio exhibido como causa generadora de su posesión debe
probar la relación de los actos jurídicos que de manera sucesiva hayan servido para efectuar esa
transmisión de dominio partiendo del acto inscrito hasta llegar al actor, porque de lo contrario, es
decir, que el actor exhiba un documento sin vinculación alguna con el propietario inscrito, haría
nugatorios los efectos frente a terceros que tienen las inscripciones registrales.

TERCER TRIBUNAL COLEGIADO EN MATERIA CIVIL DEL PRIMER CIRCUITO.


Amparo directo 56/2012. Roberto Buendía Álvarez. 23 de febrero de 2012. Unanimidad de votos.
Ponente: Benito Alva Zenteno. Secretario: Ricardo Núñez Ayala.

Nota: Por ejecutoria del 22 de enero de 2014, la Primera Sala declaró inexistente la contradicción de
tesis 187/2013 derivada de la denuncia de la que fue objeto el criterio contenido en esta tesis, al
estimarse que no son discrepantes los criterios materia de la denuncia respectiva.




                                               Pág. 1 de 2              Fecha de impresión 01/10/2025
  https://sjf2.scjn.gob.mx/detalle/tesis/2000627
                                                 Semanario Judicial de la Federación




                                            Pág. 2 de 2       Fecha de impresión 01/10/2025
https://sjf2.scjn.gob.mx/detalle/tesis/2000627

\end{quote}

\subsection{Formalization}

\subsubsection{Clause: tesis2000627\_principio\_fundamental}

\textit{Thesis-specific principle: Good faith possessors require just title as generating cause of possession for adverse possession to operate}

\textbf{Intermediate Representation:}

\begin{quote}
good faith exercise is obligated and possession as owner is obligated
\end{quote}

\textbf{Witness Clause:}

\begin{lstlisting}[style=witnessstyle, language=Witness]
clause tesis2000627_principio_fundamental = oblig(ejercicio_buena_fe) AND oblig(posesion_como_propietario) IMPLIES oblig(causa_generadora_probada);
\end{lstlisting}

\subsubsection{Clause: tesis2000627\_revelar\_causa}

\textit{Supporting logic: Must reveal and fully prove the generating cause of possession}

\textbf{Intermediate Representation:}

\begin{quote}
ejercer prescripcion positiva is obligated
\end{quote}

\textbf{Witness Clause:}

\begin{lstlisting}[style=witnessstyle, language=Witness]
clause tesis2000627_revelar_causa = oblig(ejercer_prescripcion_positiva) IMPLIES oblig(causa_generadora_revelada);
\end{lstlisting}

\subsubsection{Clause: tesis2000627\_determinacion\_judicial}

\textit{Judicial determination: Judge needs generating fact/act to determine possession quality and timing}

\textbf{Intermediate Representation:}

\begin{quote}
If causa generadora revelada is obligated, then (possession as owner is obligated and (good faith exercise is obligated and not (bad faith exercise is obligated)))
\end{quote}

\textbf{Witness Clause:}

\begin{lstlisting}[style=witnessstyle, language=Witness]
clause tesis2000627_determinacion_judicial = oblig(causa_generadora_revelada) IMPLIES (oblig(posesion_como_propietario) AND (oblig(ejercicio_buena_fe) AND not(oblig(ejercicio_mala_fe))));
\end{lstlisting}

\subsubsection{Clause: tesis2000627\_determinacion\_temporal}

\textit{Timing requirement: Must establish when prescription period begins}

\textbf{Intermediate Representation:}

\begin{quote}
If causa generadora revelada is obligated, then (posesion tiempo 5 anos is obligated or posesion tiempo 10 anos is obligated)
\end{quote}

\textbf{Witness Clause:}

\begin{lstlisting}[style=witnessstyle, language=Witness]
clause tesis2000627_determinacion_temporal = oblig(causa_generadora_revelada) IMPLIES (oblig(posesion_tiempo_5_anos) OR oblig(posesion_tiempo_10_anos));
\end{lstlisting}

\newpage

\section{Thesis 2003880: PRESCRIPCIÓN POSITIVA O USUCAPIÓN FUNDADA EN INMATRICULACIÓN. CORRESPONDE AL REGISTRO PÚBLICO DE LA PROPIEDAD Y DE COMERCIO ORDENAR LA INSCRIPCIÓN DE}
\label{thesis:2003880}

\subsection{Original Thesis Text}

\begin{quote}
                                                   Semanario Judicial de la Federación

                                                Tesis

 Registro digital: 2003880

 Instancia: Tribunales             Décima Época                           Materia(s): Civil
 Colegiados de Circuito

 Tesis: I.3o.C.104 C (10a.)        Fuente: Semanario Judicial de la       Tipo: Aislada
                                   Federación y su Gaceta.
                                   Libro XXI, Junio de 2013, Tomo 2,
                                   página 1288


PRESCRIPCIÓN POSITIVA O USUCAPIÓN FUNDADA EN INMATRICULACIÓN. CORRESPONDE
AL REGISTRO PÚBLICO DE LA PROPIEDAD Y DE COMERCIO ORDENAR LA INSCRIPCIÓN DE
LA PROPIEDAD (ARTÍCULOS 1151, 1152, 3055 Y 3056 DEL CÓDIGO CIVIL PARA EL DISTRITO
FEDERAL).


Conforme a los preceptos citados con antelación, la prescripción positiva o usucapión debe reunir
los requisitos siguientes: a) Debe tenerse en concepto de propietario; b) Debe ser pacífica; c)
Continua; y, d) Pública. Los bienes inmuebles se prescriben en cinco años, cuando se poseen
reuniendo los requisitos anteriores; en cinco años cuando los inmuebles hayan sido objeto de una
inscripción de posesión, o en diez años cuando se poseen de mala fe, si la posesión es en concepto
de propietario, pacífica, continua y pública. Por tanto, la prescripción positiva o usucapión constituye
un derecho que el propietario de un bien inmueble puede adquirir, siempre y cuando haya obtenido
mediante resolución judicial o administrativa la inscripción de la posesión de un bien inmueble, una
vez que hayan transcurrido cinco años, siempre y cuando se reúnan los requisitos para la
prescripción. Si la posesión es de buena fe, podrá acudir el interesado ante el Registro Público de la
Propiedad y de Comercio para que éste ordene la inscripción de la propiedad adquirida por
prescripción positiva, quien la ordenará siempre y cuando el interesado acredite fehacientemente
haber continuado en la posesión del bien inmueble, sin que exista asiento alguno que contradiga la
posesión inscrita.

TERCER TRIBUNAL COLEGIADO EN MATERIA CIVIL DEL PRIMER CIRCUITO.


Amparo directo 665/2011. Francisco Benítez Soriano. 26 de enero de 2012. Unanimidad de votos.
Ponente: Neófito López Ramos. Secretaria: Ana Lilia Osorno Arroyo.




                                              Pág. 1 de 2              Fecha de impresión 01/10/2025
  https://sjf2.scjn.gob.mx/detalle/tesis/2003880
                                                 Semanario Judicial de la Federación




                                            Pág. 2 de 2       Fecha de impresión 01/10/2025
https://sjf2.scjn.gob.mx/detalle/tesis/2003880

\end{quote}

\subsection{Formalization}

\subsubsection{Clause: tesis2003880\_principio\_fundamental}

\textit{Thesis-specific principle: Possession must be "in concept of owner" (peaceful, continuous, public) for both good and bad faith}

\textbf{Intermediate Representation:}

\begin{quote}
possession as owner is obligated and (oblig(ejercicio\_buena\_fe) OR oblig(ejercicio\_mala\_fe)) IMPLIES oblig(ejercer\_prescripcion\_positiva)
\end{quote}

\textbf{Witness Clause:}

\begin{lstlisting}[style=witnessstyle, language=Witness]
clause tesis2003880_principio_fundamental = oblig(posesion_como_propietario) AND (oblig(ejercicio_buena_fe) OR oblig(ejercicio_mala_fe)) IMPLIES oblig(ejercer_prescripcion_positiva);
\end{lstlisting}

\subsubsection{Clause: tesis2003880\_posesion\_ambas\_fe}

\textit{Supporting logic: Possession as owner applies to both good faith and bad faith cases}

\textbf{Intermediate Representation:}

\begin{quote}
If possession as owner is obligated, then (good faith exercise is obligated or bad faith exercise is obligated)
\end{quote}

\textbf{Witness Clause:}

\begin{lstlisting}[style=witnessstyle, language=Witness]
clause tesis2003880_posesion_ambas_fe = oblig(posesion_como_propietario) IMPLIES (oblig(ejercicio_buena_fe) OR oblig(ejercicio_mala_fe));
\end{lstlisting}

\subsubsection{Clause: tesis2003880\_dominio\_economico}

\textit{Economic dominion: Must be dominator of the thing, commanding and enjoying it as owner}

\textbf{Intermediate Representation:}

\begin{quote}
possession as owner is obligated
\end{quote}

\textbf{Witness Clause:}

\begin{lstlisting}[style=witnessstyle, language=Witness]
clause tesis2003880_dominio_economico = oblig(posesion_como_propietario) IMPLIES oblig(causa_generadora_probada);
\end{lstlisting}

\subsubsection{Clause: tesis2003880\_posesion\_originaria}

\textit{Original possession: Must begin with cause different from derived possession}

\textbf{Intermediate Representation:}

\begin{quote}
possession as owner is obligated and not (temporal possession is obligated)
\end{quote}

\textbf{Witness Clause:}

\begin{lstlisting}[style=witnessstyle, language=Witness]
clause tesis2003880_posesion_originaria = oblig(posesion_como_propietario) AND not(oblig(posesion_temporal)) IMPLIES oblig(causa_generadora_probada);
\end{lstlisting}

\newpage

\section{Thesis 2008083: GOOD FAITH ADVERSE POSSESSION - PROOF REQUIREMENTS FOR JUST TITLE}
\label{thesis:2008083}

\subsection{Original Thesis Text}

\begin{quote}
                                                   Semanario Judicial de la Federación

                                                Tesis

 Registro digital: 2008083

 Instancia: Primera Sala           Décima Época                           Materia(s): Civil

 Tesis: 1a./J. 82/2014 (10a.)      Fuente: Semanario Judicial de la       Tipo: Jurisprudencia
                                   Federación.


PRESCRIPCIÓN ADQUISITIVA. AUNQUE LA LEGISLACIÓN APLICABLE NO EXIJA QUE EL
JUSTO TÍTULO O ACTO TRASLATIVO DE DOMINIO QUE CONSTITUYE LA CAUSA
GENERADORA DE LA POSESIÓN DE BUENA FE, SEA DE FECHA CIERTA, LA CERTEZA DE LA
FECHA DEL ACTO JURÍDICO DEBE PROBARSE EN FORMA FEHACIENTE POR SER UN
ELEMENTO DEL JUSTO TÍTULO (INTERRUPCIÓN DE LA JURISPRUDENCIA 1a./J. 9/2008).


Esta Primera Sala de la Suprema Corte de Justicia de la Nación, en la jurisprudencia citada,
estableció que para la procedencia de la acción de prescripción positiva de buena fe es
indispensable que el documento privado que se exhiba como causa generadora de la posesión sea
de fecha cierta, porque: a) se inscribió en el Registro Público de la Propiedad; b) fue presentado
ante algún funcionario por razón de su oficio; o, c) alguno de sus firmantes falleció. Ahora bien, una
nueva reflexión sobre el tema lleva a apartarse de ese criterio y, por ende, a interrumpir dicha
jurisprudencia, ya que, tanto la certeza de la fecha como la celebración misma del acto jurídico
traslativo de dominio, incluyendo la autenticidad del documento, pueden acreditarse con diversos
medios de prueba que deben quedar a la valoración del juzgador, además de que el cumplimiento
con alguno de los tres requisitos señalados no es óptimo para acreditar el "justo título". En efecto, el
justo título es un acto traslativo de dominio "imperfecto", que quien pretende usucapir el bien a su
favor cree fundadamente bastante para transferirle el dominio, lo que implica que esa creencia debe
ser seria y descansar en un error que, en concepto del juzgador, sea fundado, al tratarse de uno
que "en cualquier persona" pueda provocar una creencia respecto de la validez del título. Por tanto,
para probar su justo título, el promovente debe aportar al juicio de usucapión las pruebas necesarias
para acreditar: 1) que el acto traslativo de dominio que constituye su justo título tuvo lugar, lo cual
debe acompañarse de pruebas que demuestren que objetivamente existían bases suficientes para
creer fundadamente que el enajenante podía disponer del bien, lo cual prueba cierta diligencia e
interés en el adquirente en conocer el origen del título que aduce tener su enajenante; 2) si el acto
traslativo de dominio de que se trata es oneroso, que se hicieron pagos a cuenta del precio pactado;
en caso contrario, tendrá que probar que la transmisión del bien se le hizo en forma gratuita; y, 3) la
fecha de celebración del acto jurídico traslativo de dominio, la cual deberá acreditarse en forma
fehaciente, pues constituye el punto de partida para el cómputo del plazo necesario para que opere
la prescripción adquisitiva de buena fe; además de probar que ha poseído en concepto de
propietario con su justo título, de forma pacífica, pública y continua durante cinco años, como lo
establecen los Códigos Civiles de los Estados de México, de Nuevo León y de Jalisco. De manera
que todo aquel que no pueda demostrar un nivel mínimo de diligencia, podrá prescribir, pero en el
plazo más largo de diez años, previsto en los códigos citados, ya que, de otra forma, se estará
ampliando injustificadamente el régimen especial que el legislador creó para aquellas personas que
puedan probar que su creencia en la validez de su título es fundada, con base en circunstancias
objetivas, y no apreciaciones meramente subjetivas ajenas a la realidad. Así, la procedencia de la
prescripción adquisitiva que ejerce un poseedor que aduce ser de buena fe, tendrá que cimentarse
en la convicción que adquiera el juzgador de la autenticidad del propio título y de la fecha a partir de


                                              Pág. 1 de 3             Fecha de impresión 01/10/2025
  https://sjf2.scjn.gob.mx/detalle/tesis/2008083
                                                   Semanario Judicial de la Federación

la cual se inició la posesión en concepto de propietario, con base en la valoración de los diversos
medios de convicción que ofrezca la parte actora para demostrar que es fundada su creencia en la
validez de su título, debiendo precisar que la carga de la prueba recae en la parte actora.

PRIMERA SALA


Contradicción de tesis 204/2014. Entre las sustentadas por el Quinto Tribunal Colegiado en Materia
Civil del Tercer Circuito y el Segundo Tribunal Colegiado en Materia Civil del Segundo Circuito. 5 de
noviembre de 2014. La votación se dividió en dos partes: mayoría de cuatro votos por la
competencia. Disidente: José Ramón Cossío Díaz. Unanimidad de cinco votos de los Ministros
Arturo Zaldívar Lelo de Larrea, José Ramón Cossío Díaz, Jorge Mario Pardo Rebolledo, Olga
Sánchez Cordero de García Villegas y Alfredo Gutiérrez Ortiz Mena, en cuanto al fondo. Ponente:
Jorge Mario Pardo Rebolledo. Secretaria: Rosa María Rojas Vértiz Contreras.

Tesis y/o criterios contendientes:

El Segundo Tribunal Colegiado en Materia Civil del Segundo Circuito al resolver los juicios de
amparo directo 9/2010, 74/2010, 622/2010, 899/2010 y 860/2010 que dieron origen a la tesis
jurisprudencial II.2o.C J/31, de rubro: "ACCIÓN DE USUCAPIÓN. NO LE ES APLICABLE LA
FIGURA DE LA FECHA CIERTA PARA ACREDITAR LA CAUSA GENERADORA DE LA
POSESIÓN (LEGISLACIÓN DEL ESTADO DE MÉXICO ABROGADA).", publicada en el Semanario
Judicial de la Federación y su Gaceta, Novena Época, Tomo XXXIII, mayo de 2011, página 833, con
número de registro digital: 162244; el Quinto Tribunal Colegiado en Materia Civil del Tercer Circuito,
al resolver el juicio de amparo directo 253/2014, en el que consideró fundamentalmente que si bien
la legislación del Estado de Jalisco no exige que la posesión necesaria para usucapir deba
apoyarse en un "justo título", ello no significa que la actora quede exenta de revelar y justificar la
causa generadora de su ocupación, debiendo demostrar que el documento en que sustenta el
motivo de su posesión sea de fecha cierta, no como acto traslativo de dominio perfecto, sino como
hecho jurídico para conocer la fecha a partir de la que ha de computarse el término legal de la
prescripción.

Nota: La presente tesis interrumpe el criterio sostenido en la diversa 1a./J. 9/2008, de rubro:
"PRESCRIPCIÓN ADQUISITIVA. EL CONTRATO PRIVADO DE COMPRAVENTA QUE SE
EXHIBE PARA ACREDITAR EL JUSTO TÍTULO O LA CAUSA GENERADORA DE LA POSESIÓN,
DEBE SER DE FECHA CIERTA (LEGISLACIÓN DEL ESTADO DE NUEVO LEÓN).", que aparece
publicada en el Semanario Judicial de la Federación y su Gaceta, Novena Época, Tomo XXVII, abril
de 2008, página 315.

Tesis de jurisprudencia 82/2014 (10a.). Aprobada por la Primera Sala de este Alto Tribunal, en
sesión de fecha diecinueve de noviembre de dos mil catorce.


Esta tesis se publicó el viernes 05 de diciembre de 2014 a las 10:05 horas en el Semanario Judicial
de la Federación y, por ende, se considera de aplicación obligatoria a partir del lunes 08 de
diciembre de 2014, para los efectos previstos en el punto séptimo del Acuerdo General Plenario
19/2013.




                                              Pág. 2 de 3            Fecha de impresión 01/10/2025
  https://sjf2.scjn.gob.mx/detalle/tesis/2008083
                                                 Semanario Judicial de la Federación




                                            Pág. 3 de 3       Fecha de impresión 01/10/2025
https://sjf2.scjn.gob.mx/detalle/tesis/2008083

\end{quote}

\subsection{Formalization}

\subsubsection{Clause: tesis2008083\_principio\_fundamental}

\textit{Additional actions for this thesis Additional assets for proof requirements Faith exclusivity: Adverse possession is either good faith XOR bad faith Thesis-specific principle: Good faith adverse possession requires proving authenticity, date, and diligence}

\textbf{Intermediate Representation:}

\begin{quote}
If good faith exercise is obligated and adverse possession claim is claimed, then (autenticidad titulo is obligated and fecha fehaciente is obligated and diligencia adquirente is obligated)
\end{quote}

\textbf{Witness Clause:}

\begin{lstlisting}[style=witnessstyle, language=Witness]
clause tesis2008083_principio_fundamental = oblig(ejercicio_buena_fe) AND claim(demanda_prescripcion) IMPLIES (oblig(autenticidad_titulo) AND oblig(fecha_fehaciente) AND oblig(diligencia_adquirente));
\end{lstlisting}

\subsubsection{Clause: tesis2008083\_prueba\_autenticidad}

\textit{Proof requirements: Must prove title authenticity and date reliably}

\textbf{Intermediate Representation:}

\begin{quote}
If proven generating cause is obligated, then (autenticidad titulo is obligated and fecha fehaciente is obligated)
\end{quote}

\textbf{Witness Clause:}

\begin{lstlisting}[style=witnessstyle, language=Witness]
clause tesis2008083_prueba_autenticidad = oblig(causa_generadora_probada) IMPLIES (oblig(autenticidad_titulo) AND oblig(fecha_fehaciente));
\end{lstlisting}

\subsubsection{Clause: tesis2008083\_diligencia\_buena\_fe}

\textit{Diligence requirement: Good faith requires demonstrating due diligence in acquisition}

\textbf{Intermediate Representation:}

\begin{quote}
good faith exercise is obligated
\end{quote}

\textbf{Witness Clause:}

\begin{lstlisting}[style=witnessstyle, language=Witness]
clause tesis2008083_diligencia_buena_fe = oblig(ejercicio_buena_fe) IMPLIES oblig(diligencia_adquirente);
\end{lstlisting}

\subsubsection{Clause: tesis2008083\_prueba\_pagos}

\textit{Payment proof: Onerous contracts require proving payments made}

\textbf{Intermediate Representation:}

\begin{quote}
proven generating cause is obligated and contrato compraventa is obligated
\end{quote}

\textbf{Witness Clause:}

\begin{lstlisting}[style=witnessstyle, language=Witness]
clause tesis2008083_prueba_pagos = oblig(causa_generadora_probada) AND oblig(contrato_compraventa) IMPLIES oblig(pagos_precio);
\end{lstlisting}

\subsubsection{Clause: tesis2008083\_carga\_prueba}

\textit{Burden of proof: Claimant must prove generating cause with reliable evidence}

\textbf{Intermediate Representation:}

\begin{quote}
adverse possession claim is claimed
\end{quote}

\textbf{Witness Clause:}

\begin{lstlisting}[style=witnessstyle, language=Witness]
clause tesis2008083_carga_prueba = claim(demanda_prescripcion) IMPLIES oblig(causa_generadora_revelada);
\end{lstlisting}

\subsubsection{Clause: tesis2008083\_falta\_diligencia}

\textit{Insufficient diligence consequence: Lack of diligence means bad faith (10-year period)}

\textbf{Intermediate Representation:}

\begin{quote}
not (diligencia adquirente is obligated)
\end{quote}

\textbf{Witness Clause:}

\begin{lstlisting}[style=witnessstyle, language=Witness]
clause tesis2008083_falta_diligencia = not(oblig(diligencia_adquirente)) IMPLIES oblig(ejercicio_mala_fe);
\end{lstlisting}

\newpage

\section{Thesis 2009239: BAD FAITH ADVERSE POSSESSION WITHOUT JUST TITLE}
\label{thesis:2009239}

\subsection{Original Thesis Text}

\begin{quote}
                                                   Semanario Judicial de la Federación

                                               Tesis

 Registro digital: 2009239

 Instancia: Tribunales             Décima Época                          Materia(s): Civil
 Colegiados de Circuito

 Tesis: XII.C.1 C (10a.)           Fuente: Semanario Judicial de la      Tipo: Aislada
                                   Federación.


PRESCRIPCIÓN POSITIVA. POSESIÓN DE MALA FE SIN JUSTO TÍTULO ES APTA PARA QUE
OPERE (LEGISLACIÓN DEL ESTADO DE SINALOA).


El artículo 1150, fracciones I y III, del Código Civil para el Estado de Sinaloa establece que los
bienes inmuebles se prescriben en cinco años cuando se poseen en concepto de propietario, con
buena fe, pacífica, continua y públicamente, y en diez años, cuando se poseen de mala fe, si la
posesión es en concepto de propietario, pacífica, continua y pública, en tanto que el diverso numeral
807 del propio código, precisa que es poseedor de buena fe, el que entra en posesión en virtud de
un título suficiente para darle derecho de poseer, también el que ignora los vicios de su título que le
impiden poseer con derecho, y que es poseedor de mala fe, el que entra a la posesión sin título
alguno para hacerlo, lo mismo que el que conoce los vicios de su título que le impiden poseer con
derecho. Por tanto, si bien es verdad que el artículo 827 del código mencionado establece que sólo
la posesión que se adquiere y disfruta en concepto de dueño de la cosa puede producir la
prescripción, ello no significa que dicha disposición establezca que para prescribir en todos los
casos necesariamente se requiere acreditar la existencia de un acto traslativo de dominio, aun
cuando se trate de posesión de mala fe, pues es evidente que se refiere a que esa posesión en
concepto de propietario se puede adquirir de buena o mala fe, con justo título o sin él, pues de otra
manera no se hubiera realizado la distinción de referencia. De ahí que la interpretación que debe
darse al citado numeral, en concordancia con los artículos 807 y 1150, fracciones I y III, de dicho
código, consiste en que se deben distinguir dos formas para adquirir la propiedad por prescripción
positiva: la primera, cuando la posesión es de buena fe, la cual debe ser en concepto de propietario,
pacífica, continua y pública, por más de cinco años donde se requiere justo título, que puede ser
objetiva o subjetivamente válido; y, la segunda, cuando la posesión es de mala fe, en la que no se
requiere justo título, sino que únicamente debe revelarse y probarse la causa generadora de la
posesión que, además, sea en concepto de propietario, pacífica, continua y pública, y por más de
diez años.

TRIBUNAL COLEGIADO EN MATERIA CIVIL DEL DÉCIMO SEGUNDO CIRCUITO.


Amparo directo 203/2014. Ariel Eduardo López Gaxiola. 5 de marzo de 2015. Unanimidad de votos.
Ponente: Ramona Manuela Campos Sauceda. Secretario: Roberto Cisneros Delgado.

Nota: Por ejecutoria del 8 de diciembre de 2021, la Primera Sala declaró improcedente la
contradicción de tesis 119/2021, derivada de la denuncia de la que fue objeto el criterio contenido
en esta tesis.




                                              Pág. 1 de 2             Fecha de impresión 01/10/2025
  https://sjf2.scjn.gob.mx/detalle/tesis/2009239
                                                  Semanario Judicial de la Federación


Esta tesis se publicó el viernes 22 de mayo de 2015 a las 09:30 horas en el Semanario Judicial de
la Federación.




                                             Pág. 2 de 2          Fecha de impresión 01/10/2025
 https://sjf2.scjn.gob.mx/detalle/tesis/2009239

\end{quote}

\subsection{Formalization}

\subsubsection{Clause: tesis2009239\_principio\_fundamental}

\textit{Faith exclusivity: Adverse possession is either good faith XOR bad faith Thesis-specific principle: Bad faith adverse possession does not require just title, only revealing generating cause}

\textbf{Intermediate Representation:}

\begin{quote}
If bad faith exercise is obligated and adverse possession claim is claimed, then (causa generadora revelada is obligated and not (proven generating cause is obligated))
\end{quote}

\textbf{Witness Clause:}

\begin{lstlisting}[style=witnessstyle, language=Witness]
clause tesis2009239_principio_fundamental = oblig(ejercicio_mala_fe) AND claim(demanda_prescripcion) IMPLIES (oblig(causa_generadora_revelada) AND not(oblig(causa_generadora_probada)));
\end{lstlisting}

\subsubsection{Clause: tesis2009239\_mala\_fe\_sin\_titulo}

\textit{Bad faith distinction: Bad faith possession requires only revealing cause, not proving just title}

\textbf{Intermediate Representation:}

\begin{quote}
bad faith exercise is obligated
\end{quote}

\textbf{Witness Clause:}

\begin{lstlisting}[style=witnessstyle, language=Witness]
clause tesis2009239_mala_fe_sin_titulo = oblig(ejercicio_mala_fe) IMPLIES not(oblig(causa_generadora_probada));
\end{lstlisting}

\subsubsection{Clause: tesis2009239\_requisitos\_mala\_fe}

\textit{Possession requirements: Bad faith requires possession as owner, peaceful, continuous, public for 10 years}

\textbf{Intermediate Representation:}

\begin{quote}
If bad faith exercise is obligated and adverse possession claim is claimed, then (possession as owner is obligated and temporal possession is obligated)
\end{quote}

\textbf{Witness Clause:}

\begin{lstlisting}[style=witnessstyle, language=Witness]
clause tesis2009239_requisitos_mala_fe = oblig(ejercicio_mala_fe) AND claim(demanda_prescripcion) IMPLIES (oblig(posesion_como_propietario) AND oblig(posesion_temporal));
\end{lstlisting}

\subsubsection{Clause: tesis2009239\_distincion\_fe}

\textit{Good faith vs bad faith distinction: Good faith requires just title, bad faith only requires revealing cause}

\textbf{Intermediate Representation:}

\begin{quote}
good faith exercise is obligated
\end{quote}

\textbf{Witness Clause:}

\begin{lstlisting}[style=witnessstyle, language=Witness]
clause tesis2009239_distincion_fe = oblig(ejercicio_buena_fe) IMPLIES oblig(causa_generadora_probada);
\end{lstlisting}

\subsubsection{Clause: tesis2009239\_revelacion\_causa}

\textit{Generating cause revelation: Even in bad faith, must reveal how possession was acquired}

\textbf{Intermediate Representation:}

\begin{quote}
adverse possession claim is claimed
\end{quote}

\textbf{Witness Clause:}

\begin{lstlisting}[style=witnessstyle, language=Witness]
clause tesis2009239_revelacion_causa = claim(demanda_prescripcion) IMPLIES oblig(causa_generadora_revelada);
\end{lstlisting}

\newpage

\section{Thesis 2011927: AGRARIAN ADVERSE POSSESSION - PROCEDURAL CONGRUENCE PRINCIPLE}
\label{thesis:2011927}

\subsection{Original Thesis Text}

\begin{quote}
                                                   Semanario Judicial de la Federación

                                                Tesis

 Registro digital: 2011927

 Instancia: Tribunales             Décima Época                           Materia(s): Administrativa
 Colegiados de Circuito

 Tesis: II.1o.21 A (10a.)          Fuente: Semanario Judicial de la       Tipo: Aislada
                                   Federación.


PRESCRIPCIÓN ADQUISITIVA EN MATERIA AGRARIA. SI LA ACCIÓN SE EJERCE EN LA
MODALIDAD DE POSESIÓN DE BUENA FE, EL TRIBUNAL NO DEBE VARIAR LA LITIS Y
ANALIZAR EL SUPUESTO DE MALA FE.


Si en el juicio agrario se ejerce la prescripción adquisitiva y la causa de pedir se sustenta en hechos
relativos a la modalidad de buena fe y, con los mismos, así como con lo alegado por la demandada,
se integra la litis en la audiencia de ley prevista en el artículo 185 de la Ley Agraria, que sirve como
límite para cualquier sentencia de fondo que resuelva la controversia efectivamente planteada,
porque ésta debe sujetarse exclusivamente a las pretensiones de las partes, con las que se
construye el objeto del proceso, y no puede decidir sobre cuestiones distintas a éste, a fin de
respetar el artículo 17 de la Constitución Política de los Estados Unidos Mexicanos, que contiene el
derecho a la impartición de justicia, que debe atender al principio de congruencia, por lo que las
pretensiones y hechos de las partes no pueden ser rebasados en virtud de que el demandado hará
valer sus defensas respecto de las peticiones y hechos de la demanda, aun cuando no se acredite
esa prescripción de buena fe, el tribunal de la materia no puede analizar la usucapión de mala fe,
aunque ambas estén en el artículo 48 de la Ley Agraria. Estimar lo contrario equivaldría a variar la
litis y violar el derecho de defensa del enjuiciado, sin que esto pueda soslayarse por lo indicado en
el artículo 189 del ordenamiento referido, que permite a las autoridades de la materia resolver a
verdad sabida, porque ese principio opera en la justipreciación probatoria, lo que tampoco puede
ser destruido por la suplencia de la queja en materia agraria, pues ésta no puede hacer procedentes
acciones no planteadas por las partes, lo que se refuerza al tener presente el artículo 349 del
Código Federal de Procedimientos Civiles, de aplicación supletoria, el cual establece que las
sentencias sólo se ocuparán de las personas, cosas, acciones y excepciones que hayan sido
materia del juicio.

PRIMER TRIBUNAL COLEGIADO DEL SEGUNDO CIRCUITO CON RESIDENCIA EN CIUDAD
NEZAHUALCÓYOTL, ESTADO DE MÉXICO.


Amparo directo 640/2015. Silvestre Cadena Ramírez. 28 de enero de 2016. Unanimidad de votos.
Ponente: Miguel Enrique Sánchez Frías. Secretaria: Erika Yazmín Zárate Villa.


Esta tesis se publicó el viernes 17 de junio de 2016 a las 10:17 horas en el Semanario Judicial de la
Federación.




                                              Pág. 1 de 2             Fecha de impresión 01/10/2025
  https://sjf2.scjn.gob.mx/detalle/tesis/2011927
                                                 Semanario Judicial de la Federación




                                            Pág. 2 de 2       Fecha de impresión 01/10/2025
https://sjf2.scjn.gob.mx/detalle/tesis/2011927

\end{quote}

\subsection{Formalization}

\subsubsection{Clause: tesis2011927\_principio\_fundamental}

\textit{Additional actions for this thesis No additional assets needed - using universal assets Faith exclusivity: Adverse possession is either good faith XOR bad faith Thesis-specific principle: If action is filed as good faith, court cannot analyze bad faith variant}

\textbf{Intermediate Representation:}

\begin{quote}
good faith exercise is obligated and adverse possession claim is claimed
\end{quote}

\textbf{Witness Clause:}

\begin{lstlisting}[style=witnessstyle, language=Witness]
clause tesis2011927_principio_fundamental = oblig(ejercicio_buena_fe) AND claim(demanda_prescripcion) IMPLIES not(oblig(variacion_litis));
\end{lstlisting}

\subsubsection{Clause: tesis2011927\_congruencia\_procesal}

\textit{Procedural congruence: Court must respect the procedural framework established by parties}

\textbf{Intermediate Representation:}

\begin{quote}
adverse possession claim is claimed
\end{quote}

\textbf{Witness Clause:}

\begin{lstlisting}[style=witnessstyle, language=Witness]
clause tesis2011927_congruencia_procesal = claim(demanda_prescripcion) IMPLIES oblig(congruencia_procesal);
\end{lstlisting}

\subsubsection{Clause: tesis2011927\_prohibicion\_variacion}

\textit{Prohibition of litis variation: Cannot change from good faith to bad faith analysis}

\textbf{Intermediate Representation:}

\begin{quote}
congruencia procesal is obligated
\end{quote}

\textbf{Witness Clause:}

\begin{lstlisting}[style=witnessstyle, language=Witness]
clause tesis2011927_prohibicion_variacion = oblig(congruencia_procesal) IMPLIES not(oblig(variacion_litis));
\end{lstlisting}

\subsubsection{Clause: tesis2011927\_proteccion\_defensa}

\textit{Defense rights protection: Changing legal theory violates defendant's defense rights}

\textbf{Intermediate Representation:}

\begin{quote}
variacion litis is obligated
\end{quote}

\textbf{Witness Clause:}

\begin{lstlisting}[style=witnessstyle, language=Witness]
clause tesis2011927_proteccion_defensa = oblig(variacion_litis) IMPLIES not(oblig(proteccion_registral));
\end{lstlisting}

\newpage

\section{Thesis 2014858: THIRD PARTY EXCLUSION - UNREGISTERED DOCUMENTS VS GOOD FAITH ACQUIRERS}
\label{thesis:2014858}

\subsection{Original Thesis Text}

\begin{quote}
                                                   Semanario Judicial de la Federación

                                               Tesis

 Registro digital: 2014858

 Instancia: Tribunales             Décima Época                          Materia(s): Civil
 Colegiados de Circuito

 Tesis: XX.2o.P.C.4 C (10a.)       Fuente: Gaceta del Semanario       Tipo: Aislada
                                   Judicial de la Federación.
                                   Libro 45, Agosto de 2017, Tomo IV,
                                   página 3175


TERCERÍA EXCLUYENTE DE DOMINIO. LA FALTA DE INSCRIPCIÓN REGISTRAL DEL
DOCUMENTO CON EL QUE UN TERCERISTA PRETENDA FUNDAR ESA ACCIÓN, CONLLEVA
QUE NO PUEDA OPONERSE A UN TERCERO ADQUIRENTE DE BUENA FE, DE UN DERECHO
REAL SOBRE ESE MISMO BIEN, POR ADJUDICACIÓN EN UN JUICIO EJECUTIVO
MERCANTIL.


La otrora Tercera Sala de la Suprema Corte de Justicia de la Nación, sostuvo en la jurisprudencia
3a./J. 22/94, de rubro: "EMBARGO, ES ILEGAL EL TRABADO EN BIENES SALIDOS DEL
DOMINIO DEL DEUDOR, AUN CUANDO NO SE ENCUENTREN INSCRITOS EN EL REGISTRO
PÚBLICO DE LA PROPIEDAD A NOMBRE DEL NUEVO ADQUIRENTE. (LEGISLACIÓN DE
DURANGO SIMILAR A LA DEL DISTRITO FEDERAL).", que el acreedor quirografario no tiene un
derecho real, ni poder directo e inmediato sobre la cosa embargada, porque el embargo realizado
en un juicio ejecutivo mercantil aun cuando se encuentre registrado no puede ser oponible a
quienes adquirieron con anterioridad la propiedad del bien. Por otro lado, la Primera Sala de la
Suprema Corte de Justicia de la Nación, al resolver la contradicción de tesis 333/2012, el dieciséis
de enero de dos mil trece, sostuvo que el derecho real de propiedad del cónyuge que no fue oído en
un juicio ejecutivo mercantil, respecto de un inmueble adquirido con motivo de la sociedad conyugal
que no se encuentra inscrita en el Registro Público de la Propiedad, no es oponible mediante el
juicio de amparo contra el derecho de quien adquirió por remate ese inmueble embargado, esto es,
contra el adjudicatario de buena fe en el juicio ejecutivo mercantil pues, a partir de la venta, ésta
constituye un acto traslativo de dominio que importa la concurrencia de un derecho real. Sin que
obste el origen remoto del derecho de propiedad, consistente en el derecho personal que se ejerció
judicialmente pues, una vez adjudicado el bien, el derecho real es totalmente independiente de su
origen. Las anteriores consideraciones permiten afirmar, que mediante la tramitación y culminación
del procedimiento de remate y adjudicación, derivado de un juicio ejecutivo mercantil, el postor o el
adjudicatario de buena fe que haya cubierto el precio fijado para la venta, adquiere un derecho de
propiedad sobre el inmueble rematado en el procedimiento de venta judicial, lo que implica que
adquiere coetáneamente un auténtico derecho real sobre ese bien raíz, en forma análoga a como
sucedería en el caso de una compraventa privada. En esas condiciones, por identidad jurídica
sustancial, la falta de inscripción registral del documento con el que un tercerista pretende fundar su
acción de tercería excluyente de dominio, conlleva que no pueda oponerse a un tercero adquirente
de buena fe, de un derecho real sobre ese mismo bien, por adjudicación en un juicio ejecutivo
mercantil, pues este último se encuentra inmune frente al derecho real que pretende hacer valer el
tercerista, que carece de inscripción en el Registro Público de la Propiedad.

SEGUNDO TRIBUNAL COLEGIADO EN MATERIAS PENAL Y CIVIL DEL VIGÉSIMO CIRCUITO.


                                              Pág. 1 de 2             Fecha de impresión 18/08/2025
  https://sjf2.scjn.gob.mx/detalle/tesis/2014858
                                                  Semanario Judicial de la Federación



Amparo directo 797/2016. Esaú Nucamendi Molina y otros. 9 de febrero de 2017. Unanimidad de
votos. Ponente: Susana Teresa Sánchez González. Secretario: José Luis Martínez Villarreal.

Nota: La tesis de jurisprudencia 3a./J. 22/94 citada, aparece publicada en la Gaceta del Semanario
Judicial de la Federación, Octava Época, Número 80, agosto de 1994, página 21.

La parte conducente de la ejecutoria relativa a la contradicción de tesis 333/2012 citada, aparece
publicada en el Semanario Judicial de la Federación y su Gaceta, Décima Época, Libro XXI, Tomo
1, junio de 2013, página 553.

Por ejecutoria del 18 de abril de 2018, la Primera Sala declaró improcedente la contradicción de
tesis 371/2017 derivada de la denuncia de la que fue objeto el criterio contenido en esta tesis, al
estimarse que uno de los criterios en contradicción solamente constituye la aplicación de una
jurisprudencia de la Suprema Corte de Justicia de la Nación.


Esta tesis se publicó el viernes 04 de agosto de 2017 a las 10:12 horas en el Semanario Judicial de
la Federación.




                                             Pág. 2 de 2           Fecha de impresión 18/08/2025
 https://sjf2.scjn.gob.mx/detalle/tesis/2014858

\end{quote}

\subsection{Formalization}

\subsubsection{Clause: tesis2014858\_principio\_fundamental}

\textit{No additional subjects or actions needed - using universal assets Faith exclusivity: Adverse possession is either good faith XOR bad faith Thesis-specific principle: Unregistered documents cannot be opposed to good faith acquirers through judicial auction}

\textbf{Intermediate Representation:}

\begin{quote}
unregistered contract is obligated and adquisicion remate is obligated and good faith exercise is obligated
\end{quote}

\textbf{Witness Clause:}

\begin{lstlisting}[style=witnessstyle, language=Witness]
clause tesis2014858_principio_fundamental = oblig(contrato_no_inscrito) AND oblig(adquisicion_remate) AND oblig(ejercicio_buena_fe) IMPLIES not(oblig(terceria_excluyente));
\end{lstlisting}

\subsubsection{Clause: tesis2014858\_proteccion\_registral}

\textit{Registry protection: Registration is necessary for opposability to good faith third parties}

\textbf{Intermediate Representation:}

\begin{quote}
good faith exercise is obligated and adquisicion remate is obligated
\end{quote}

\textbf{Witness Clause:}

\begin{lstlisting}[style=witnessstyle, language=Witness]
clause tesis2014858_proteccion_registral = oblig(ejercicio_buena_fe) AND oblig(adquisicion_remate) IMPLIES oblig(proteccion_registral);
\end{lstlisting}

\subsubsection{Clause: tesis2014858\_adquisicion\_buena\_fe}

\textit{Good faith acquisition: Judicial auction creates real property rights for good faith acquirers}

\textbf{Intermediate Representation:}

\begin{quote}
adquisicion remate is obligated and good faith exercise is obligated
\end{quote}

\textbf{Witness Clause:}

\begin{lstlisting}[style=witnessstyle, language=Witness]
clause tesis2014858_adquisicion_buena_fe = oblig(adquisicion_remate) AND oblig(ejercicio_buena_fe) IMPLIES oblig(ejercer_derecho_legal);
\end{lstlisting}

\subsubsection{Clause: tesis2014858\_limitacion\_oponibilidad}

\textit{Opposability limitation: Unregistered documents lack opposability against registered good faith acquirers}

\textbf{Intermediate Representation:}

\begin{quote}
unregistered contract is obligated
\end{quote}

\textbf{Witness Clause:}

\begin{lstlisting}[style=witnessstyle, language=Witness]
clause tesis2014858_limitacion_oponibilidad = oblig(contrato_no_inscrito) IMPLIES not(oblig(oponibilidad_terceros));
\end{lstlisting}

\newpage

\section{Thesis 2014940: RES JUDICATA - GOOD FAITH VS BAD FAITH AS DIFFERENT CAUSES}
\label{thesis:2014940}

\subsection{Original Thesis Text}

\begin{quote}
                                                   Semanario Judicial de la Federación

                                                Tesis

 Registro digital: 2014940

 Instancia: Tribunales             Décima Época                           Materia(s): Civil
 Colegiados de Circuito

 Tesis: III.5o.C.42 C (10a.)       Fuente: Semanario Judicial de la       Tipo: Aislada
                                   Federación.


COSA JUZGADA. NO SE ACTUALIZA DICHA EXCEPCIÓN SI EN EL PRIMER JUICIO SE
DEMANDÓ LA PRESCRIPCIÓN POSITIVA APOYADA EN UNA POSESIÓN DE BUENA FE Y EN
EL SEGUNDO SE EJERCE LA MISMA ACCIÓN, CONTRA EL MISMO DEMANDADO, PERO
SUSTENTADA EN UNA POSESIÓN DE MALA FE (LEGISLACIÓN DEL ESTADO DE COLIMA).


De los artículos 35, fracción IX, 260, 421, 425 y 469 del Código de Procedimientos Civiles para el
Estado de Colima, se advierte que entre las excepciones procesales se prevé la de cosa juzgada, la
cual, para que surta efectos en otro juicio, es necesario que entre el caso resuelto por la sentencia
ejecutoria y aquel en que ésta sea invocada, concurra identidad en las cosas, las causas, las
personas de los litigantes y la calidad con que lo fueren. Por lo que ve al requisito de la identidad en
la causa, la Primera Sala de la Suprema Corte de Justicia de la Nación, en la jurisprudencia 1a./J.
161/2007, de rubro: "COSA JUZGADA. PRESUPUESTOS PARA SU EXISTENCIA.",(1) estableció
que aquélla debe entenderse como el hecho generador que las partes hacen valer como
fundamento de las pretensiones que reclaman, por lo que es requisito indispensable para que exista
cosa juzgada, que se atienda no únicamente a la causa próxima (consecuencia directa e inmediata
de la realización del acto jurídico) sino, además, a la causa remota (causal supeditada a
acontecimientos supervenientes para su consumación), pues sólo si existe esa identidad podría
afirmarse que las cuestiones propuestas en el segundo procedimiento, ya fueron materia de análisis
en el primero y que, por ello, deba declararse procedente la excepción con la finalidad de no dar
pauta a posibles sentencias contradictorias. Por tanto, cuando en el primer juicio, donde se dictó
sentencia ejecutoria, se demandó la prescripción positiva sustentada en una posesión de buena fe
y, en el segundo, se ejerce la misma acción contra el mismo demandado, pero en una posesión de
mala fe, en realidad no se actualiza la excepción de cosa juzgada, porque la causa remota (calidad
de esa posesión) es diversa, lo que permite considerar que en el primer procedimiento no se
analizaron los mismos hechos que son materia del segundo juicio. Sin que esto impida que en la
sentencia definitiva que, en su caso, se dicte, se analice el alcance de lo resuelto en un caso
anterior, especialmente aquellos elementos que ya se encuentren decididos, lo anterior, a fin de
evitar la emisión de sentencias contradictorias.

QUINTO TRIBUNAL COLEGIADO EN MATERIA CIVIL DEL TERCER CIRCUITO.


Amparo en revisión 339/2016. 27 de abril de 2017. Mayoría de votos. Disidente: Alicia Guadalupe
Cabral Parra. Ponente: Rodolfo Castro León. Secretaria: Lizette Arroyo Delgadillo.


Esta tesis se publicó el viernes 18 de agosto de 2017 a las 10:26 horas en el Semanario Judicial de


                                              Pág. 1 de 2             Fecha de impresión 01/10/2025
  https://sjf2.scjn.gob.mx/detalle/tesis/2014940
                                                  Semanario Judicial de la Federación

la Federación.




                                             Pág. 2 de 2       Fecha de impresión 01/10/2025
 https://sjf2.scjn.gob.mx/detalle/tesis/2014940

\end{quote}

\subsection{Formalization}

\subsubsection{Clause: tesis2014940\_principio\_fundamental}

\textit{No additional actions or assets needed - using universal assets Faith exclusivity: Adverse possession is either good faith XOR bad faith Thesis-specific principle: Good faith and bad faith adverse possession have different remote causes, so res judicata does not apply}

\textbf{Intermediate Representation:}

\begin{quote}
good faith exercise is obligated and bad faith exercise is obligated
\end{quote}

\textbf{Witness Clause:}

\begin{lstlisting}[style=witnessstyle, language=Witness]
clause tesis2014940_principio_fundamental = oblig(ejercicio_buena_fe) AND oblig(ejercicio_mala_fe) IMPLIES oblig(causa_remota_diversa);
\end{lstlisting}

\subsubsection{Clause: tesis2014940\_causas\_diferentes}

\textit{Different legal causes: Good faith vs bad faith constitute different remote causes}

\textbf{Intermediate Representation:}

\begin{quote}
causa remota diversa is obligated
\end{quote}

\textbf{Witness Clause:}

\begin{lstlisting}[style=witnessstyle, language=Witness]
clause tesis2014940_causas_diferentes = oblig(causa_remota_diversa) IMPLIES not(oblig(cosa_juzgada));
\end{lstlisting}

\subsubsection{Clause: tesis2014940\_identidad\_causas}

\textit{Res judicata limitation: Res judicata requires identity in remote and proximate causes}

\textbf{Intermediate Representation:}

\begin{quote}
cosa juzgada is obligated
\end{quote}

\textbf{Witness Clause:}

\begin{lstlisting}[style=witnessstyle, language=Witness]
clause tesis2014940_identidad_causas = oblig(cosa_juzgada) IMPLIES not(oblig(causa_remota_diversa));
\end{lstlisting}

\subsubsection{Clause: tesis2014940\_procedimientos\_separados}

\textit{Sequential analysis allowed: Different possession qualities allow separate proceedings}

\textbf{Intermediate Representation:}

\begin{quote}
causa remota diversa is obligated
\end{quote}

\textbf{Witness Clause:}

\begin{lstlisting}[style=witnessstyle, language=Witness]
clause tesis2014940_procedimientos_separados = oblig(causa_remota_diversa) IMPLIES oblig(cumplir_requisitos_legales);
\end{lstlisting}

\newpage

\section{Thesis 2016571: GOOD FAITH ADVERSE POSSESSION - PROHIBITION AGAINST VARYING CLAIMS DURING TRIAL}
\label{thesis:2016571}

\subsection{Original Thesis Text}

\begin{quote}
                                                   Semanario Judicial de la Federación

                                               Tesis

 Registro digital: 2016571

 Instancia: Tribunales             Décima Época                          Materia(s): Civil
 Colegiados de Circuito

 Tesis: I.3o.C.320 C (10a.)        Fuente: Semanario Judicial de la      Tipo: Aislada
                                   Federación.


PRESCRIPCIÓN POSITIVA DE BUENA FE. QUIEN LA EJERZA NO PUEDE, DURANTE EL
JUICIO, VARIAR SUS PRETENSIONES PARA AJUSTARSE A OTRO DE LOS SUPUESTOS
PREVISTOS EN EL ARTÍCULO 1152 DEL CÓDIGO CIVIL PARA EL DISTRITO FEDERAL,
APLICABLE PARA LA CIUDAD DE MÉXICO, A EFECTO DE QUE PROCEDA AQUÉLLA, PERO
DE MALA FE.


De acuerdo con el artículo 1152 del Código Civil para el Distrito Federal, aplicable para la Ciudad de
México, los inmuebles prescriben en: (I) cinco años, cuando se poseen en concepto de propietario,
con buena fe, pacífica, continua y públicamente; (II) cinco años, cuando los inmuebles hayan sido
objeto de una inscripción de posesión; (III) diez años, cuando se poseen de mala fe, si la posesión
es en concepto de propietario, pacífica, continua y pública; y, (IV) aumentando en una tercera parte
el tiempo señalado en las fracciones I y III del artículo citado, si se demuestra, por quien tenga
interés jurídico en ello, que el poseedor de la finca rústica no la ha cultivado durante la mayor parte
del tiempo que la ha poseído, o que por no haber hecho el poseedor de la finca urbana las
reparaciones necesarias, ésta ha permanecido deshabitada la mayor parte del tiempo que ha
estado en poder de aquél. Ahora bien, el artículo 806 del propio ordenamiento, determina que es
poseedor de buena fe el que entra en la posesión en virtud de un título suficiente para darle derecho
a poseer, o bien, el que ignora los vicios de su título que le impiden poseer con derecho. Por otro
lado, también describe a la posesión de mala fe como la de aquel que la ostenta sin título alguno
para poseer, o bien, el que conoce los vicios de su título que le impiden poseer con derecho. Con
base en lo anterior, se concluye que el que ejerza la prescripción positiva de un inmueble de buena
fe, de acuerdo con la fracción I del artículo 1152 citado, no puede, durante el juicio, variar sus
pretensiones para ajustarse a otro de sus supuestos, como lo es aquélla pero de mala fe, contenida
en su fracción III. Lo anterior es así, puesto que el Juez no puede enderezar la litis para ajustarse a
un supuesto no pedido, porque invariablemente eso significaría dejar en estado de indefensión a la
contraparte.

TERCER TRIBUNAL COLEGIADO EN MATERIA CIVIL DEL PRIMER CIRCUITO.


Amparo directo 40/2018. Alma Rosa Guzmán Esquivel. 7 de febrero de 2018. Unanimidad de votos.
Ponente: Paula María García Villegas Sánchez Cordero. Secretaria: Montserrat Cesarina Camberos
Funes.


Esta tesis se publicó el viernes 06 de abril de 2018 a las 10:10 horas en el Semanario Judicial de la
Federación.


                                              Pág. 1 de 2             Fecha de impresión 01/10/2025
  https://sjf2.scjn.gob.mx/detalle/tesis/2016571
                                                 Semanario Judicial de la Federación




                                            Pág. 2 de 2       Fecha de impresión 01/10/2025
https://sjf2.scjn.gob.mx/detalle/tesis/2016571

\end{quote}

\subsection{Formalization}

\subsubsection{Clause: tesis2016571\_principio\_fundamental}

\textit{Additional actions for this thesis (reusing most existing assets) Additional assets for due process protection (reusing most existing assets) Faith exclusivity: Possession is either good faith XOR bad faith (mutually exclusive) Thesis-specific principle: Good faith adverse possession claimant cannot vary claims to bad faith during trial}

\textbf{Intermediate Representation:}

\begin{quote}
good faith exercise is obligated and adverse possession claim is claimed
\end{quote}

\textbf{Witness Clause:}

\begin{lstlisting}[style=witnessstyle, language=Witness]
clause tesis2016571_principio_fundamental = oblig(ejercicio_buena_fe) AND claim(demanda_prescripcion) IMPLIES not(oblig(variacion_litis));
\end{lstlisting}

\subsubsection{Clause: tesis2016571\_prohibicion\_variacion}

\textit{Prohibition of claim variation: Cannot change from good faith (5 years) to bad faith (10 years) theory}

\textbf{Intermediate Representation:}

\begin{quote}
good faith exercise is obligated and adverse possession claim is claimed
\end{quote}

\textbf{Witness Clause:}

\begin{lstlisting}[style=witnessstyle, language=Witness]
clause tesis2016571_prohibicion_variacion = oblig(ejercicio_buena_fe) AND claim(demanda_prescripcion) IMPLIES not(oblig(ejercicio_mala_fe));
\end{lstlisting}

\subsubsection{Clause: tesis2016571\_congruencia\_procesal}

\textit{Procedural congruence requirement: Court must respect the legal theory chosen by plaintiff}

\textbf{Intermediate Representation:}

\begin{quote}
adverse possession claim is claimed
\end{quote}

\textbf{Witness Clause:}

\begin{lstlisting}[style=witnessstyle, language=Witness]
clause tesis2016571_congruencia_procesal = claim(demanda_prescripcion) IMPLIES oblig(congruencia_procesal);
\end{lstlisting}

\subsubsection{Clause: tesis2016571\_limitacion\_judicial}

\textit{Judicial limitation: Judge cannot fix lawsuit to adjust to different legal theory not requested}

\textbf{Intermediate Representation:}

\begin{quote}
congruencia procesal is obligated
\end{quote}

\textbf{Witness Clause:}

\begin{lstlisting}[style=witnessstyle, language=Witness]
clause tesis2016571_limitacion_judicial = oblig(congruencia_procesal) IMPLIES not(oblig(enderezamiento_litis));
\end{lstlisting}

\subsubsection{Clause: tesis2016571\_proteccion\_debido\_proceso}

\textit{Due process protection: Varying claims would leave defendant defenseless}

\textbf{Intermediate Representation:}

\begin{quote}
variacion litis is obligated
\end{quote}

\textbf{Witness Clause:}

\begin{lstlisting}[style=witnessstyle, language=Witness]
clause tesis2016571_proteccion_debido_proceso = oblig(variacion_litis) IMPLIES not(oblig(proteccion_defensa));
\end{lstlisting}

\subsubsection{Clause: tesis2016571\_preservacion\_defensa}

\textit{Defense rights preservation: Procedural congruence protects defendant's defense strategy}

\textbf{Intermediate Representation:}

\begin{quote}
congruencia procesal is obligated
\end{quote}

\textbf{Witness Clause:}

\begin{lstlisting}[style=witnessstyle, language=Witness]
clause tesis2016571_preservacion_defensa = oblig(congruencia_procesal) IMPLIES oblig(proteccion_defensa);
\end{lstlisting}

\subsubsection{Clause: tesis2016571\_consistencia\_temporal}

\textit{Time period consistency: Must maintain consistent time period based on chosen faith theory}

\textbf{Intermediate Representation:}

\begin{quote}
good faith exercise is obligated and adverse possession claim is claimed
\end{quote}

\textbf{Witness Clause:}

\begin{lstlisting}[style=witnessstyle, language=Witness]
clause tesis2016571_consistencia_temporal = oblig(ejercicio_buena_fe) AND claim(demanda_prescripcion) IMPLIES oblig(posesion_temporal);
\end{lstlisting}

\subsubsection{Clause: tesis2016571\_requisitos\_buena\_fe}

\textit{Good faith requirements: Must prove sufficient title or ignorance of title defects}

\textbf{Intermediate Representation:}

\begin{quote}
If good faith exercise is obligated, then (proven generating cause is obligated or conocimiento vicios is obligated)
\end{quote}

\textbf{Witness Clause:}

\begin{lstlisting}[style=witnessstyle, language=Witness]
clause tesis2016571_requisitos_buena_fe = oblig(ejercicio_buena_fe) IMPLIES (oblig(causa_generadora_probada) OR oblig(conocimiento_vicios));
\end{lstlisting}

\subsubsection{Clause: tesis2016571\_compromiso\_teoria}

\textit{Legal theory commitment: Once good faith theory is chosen, must prove all its elements}

\textbf{Intermediate Representation:}

\begin{quote}
If good faith exercise is obligated and adverse possession claim is claimed, then (proven generating cause is obligated and possession as owner is obligated and temporal possession is obligated)
\end{quote}

\textbf{Witness Clause:}

\begin{lstlisting}[style=witnessstyle, language=Witness]
clause tesis2016571_compromiso_teoria = oblig(ejercicio_buena_fe) AND claim(demanda_prescripcion) IMPLIES (oblig(causa_generadora_probada) AND oblig(posesion_como_propietario) AND oblig(posesion_temporal));
\end{lstlisting}

\newpage

\section{Thesis 2017366: ADVERSE POSSESSION CLASSIFICATION BY QUALITY OF POSSESSION}
\label{thesis:2017366}

\subsection{Original Thesis Text}

\begin{quote}
                                                   Semanario Judicial de la Federación

                                               Tesis

 Registro digital: 2017366

 Instancia: Tribunales             Décima Época                          Materia(s): Civil
 Colegiados de Circuito

 Tesis: I.3o.C.330 C (10a.)        Fuente: Semanario Judicial de la      Tipo: Aislada
                                   Federación.


PRESCRIPCIÓN ADQUISITIVA. SU CLASIFICACIÓN ATENDIENDO A LA CALIDAD DE LA
POSESIÓN (LEGISLACIÓN APLICABLE PARA LA CIUDAD DE MÉXICO).


La acción de prescripción adquisitiva atiende a la calidad de la posesión, toda vez que la finalidad
perseguida por quien pretende usucapir es convertirse en propietario a partir de su posesión
calificada. Así, la detentación directa de la cosa que se tiene físicamente, si bien se posee por ser
propietario, lo cierto es que puede suceder que éste transmita contractualmente la posesión,
conservando la posesión originaria. Asimismo, también puede poseerse sin tener derecho, como en
el caso del apoderamiento en forma ilícita; y en otras ocasiones, hay quien está poseyendo por
encargo del propietario y quien posee porque cree que le transmitieron el derecho, pero quien lo
hizo no estaba legitimado para ello. El mero detentador es aquel quien posee a nombre de otro o
por encargo, como el depositario judicial. Si se parte de estas premisas, la posesión puede ser de
buena y de mala fe, la primera cuando el poseedor estima tener derecho a la posesión de que
disfruta pero se basa en un error, es decir, el que ignora que en un título o modo de adquirir exista
un vicio que lo invalide; mientras que la segunda se refiere al poseedor que conoce o sabe que
posee indebidamente. Luego, los poseedores de buena fe pueden ser de dos tipos: a) los que
tienen un título suficiente o bien uno viciado y lo ignoran; y, b) los que no tienen título pero
fundadamente creen tenerlo. En cuanto a los poseedores de mala fe, también pueden ser de dos
tipos: a) los que poseen con título viciado y lo saben y b) los que poseen sin título y lo saben. Pero,
además, para los poseedores de mala fe, establece una subclasificación: 1) poseedores que
poseen en forma delictuosa y 2) poseedores que no poseen en forma delictuosa. Así, encontramos
que la ley regula dichas situaciones en las cuales, sin que alguien tenga la cosa físicamente en su
poder, se le considera poseedor de ella y, por el contrario, hay quien, de hecho la tiene, pero carece
del carácter de poseedor, de acuerdo con la ley. Lo anterior, se basa en los artículos 806 y 807,
ambos del Código Civil para el Distrito Federal, aplicable para la Ciudad de México. En ese orden
de ideas, se advierte que el legislador estableció en favor del adquirente que ignora los vicios de su
título que le impiden poseer con derecho, una presunción iuris tantum, pues lo considera poseedor
de buena fe, hasta el momento en que existan actos que acrediten que adquiere el conocimiento de
que posee la cosa indebidamente. En consecuencia, cuando esa presunción se destruye, al
acreditarse que quien la tiene en su favor conoce los vicios de su título, la posesión, aun adquirida
inicialmente de buena fe, bajo el supuesto de ignorancia, deja de tener esa calidad, para convertirse
en posesión de mala fe. Esto deja claro que la posesión considerada de buena fe, por actualizarse
la presunción derivada de la ignorancia de los vicios de un título, no confiere de manera indefinida al
poseedor que se beneficia de ella, un reconocimiento legal irrevocable de buena fe, pues su
permanencia está condicionada, precisamente, a que persista la de ignorancia sobre los vicios de
su título.

TERCER TRIBUNAL COLEGIADO EN MATERIA CIVIL DEL PRIMER CIRCUITO.


                                              Pág. 1 de 2             Fecha de impresión 01/10/2025
  https://sjf2.scjn.gob.mx/detalle/tesis/2017366
                                                   Semanario Judicial de la Federación



Amparo directo 486/2017. 12 de julio de 2017. Unanimidad de votos. Ponente: Paula María García
Villegas Sánchez Cordero. Secretaria: Cinthia Monserrat Ortega Mondragón.


Esta tesis se publicó el viernes 06 de julio de 2018 a las 10:13 horas en el Semanario Judicial de la
Federación.




                                              Pág. 2 de 2            Fecha de impresión 01/10/2025
  https://sjf2.scjn.gob.mx/detalle/tesis/2017366

\end{quote}

\subsection{Formalization}

\subsubsection{Clause: tesis2017366\_principio\_fundamental}

\textit{Additional actions for this thesis (most assets now universal) Additional assets for possession classification (most assets now universal) Faith exclusivity: Possession is either good faith XOR bad faith (mutually exclusive) Thesis-specific principle: Possession classification depends on knowledge of title defects and legal basis}

\textbf{Intermediate Representation:}

\begin{quote}
If adverse possession claim is claimed, then (good faith exercise is obligated or conocimiento vicios is obligated)
\end{quote}

\textbf{Witness Clause:}

\begin{lstlisting}[style=witnessstyle, language=Witness]
clause tesis2017366_principio_fundamental = claim(demanda_prescripcion) IMPLIES (oblig(ejercicio_buena_fe) OR oblig(conocimiento_vicios));
\end{lstlisting}

\subsubsection{Clause: tesis2017366\_buena\_fe\_tipos}

\textit{Good faith types: With sufficient/defective title (unknowing) or without title but reasonable belief}

\textbf{Intermediate Representation:}

\begin{quote}
If good faith exercise is obligated, then (diligencia adquirente is obligated or not (documento con vicios is obligated))
\end{quote}

\textbf{Witness Clause:}

\begin{lstlisting}[style=witnessstyle, language=Witness]
clause tesis2017366_buena_fe_tipos = oblig(ejercicio_buena_fe) IMPLIES (oblig(diligencia_adquirente) OR not(oblig(documento_con_vicios)));
\end{lstlisting}

\subsubsection{Clause: tesis2017366\_mala\_fe\_tipos}

\textit{Bad faith types: With defective title (knowing) or without title (knowing)}

\textbf{Intermediate Representation:}

\begin{quote}
If bad faith exercise is obligated, then (conocimiento vicios is obligated or not (proven generating cause is obligated))
\end{quote}

\textbf{Witness Clause:}

\begin{lstlisting}[style=witnessstyle, language=Witness]
clause tesis2017366_mala_fe_tipos = oblig(ejercicio_mala_fe) IMPLIES (oblig(conocimiento_vicios) OR not(oblig(causa_generadora_probada)));
\end{lstlisting}

\subsubsection{Clause: tesis2017366\_presuncion\_legal}

\textit{Presumption of good faith: Legal presumption favors good faith until proven otherwise (integrated into universal good faith)}

\textbf{Intermediate Representation:}

\begin{quote}
adverse possession claim is claimed
\end{quote}

\textbf{Witness Clause:}

\begin{lstlisting}[style=witnessstyle, language=Witness]
clause tesis2017366_presuncion_legal = claim(demanda_prescripcion) IMPLIES oblig(ejercicio_buena_fe);
\end{lstlisting}

\subsubsection{Clause: tesis2017366\_conversion\_fe}

\textit{Conversion mechanism: Good faith converts to bad faith upon knowledge of title defects}

\textbf{Intermediate Representation:}

\begin{quote}
conocimiento vicios is obligated and documento con vicios is obligated
\end{quote}

\textbf{Witness Clause:}

\begin{lstlisting}[style=witnessstyle, language=Witness]
clause tesis2017366_conversion_fe = oblig(conocimiento_vicios) AND oblig(documento_con_vicios) IMPLIES oblig(conversion_fe);
\end{lstlisting}

\subsubsection{Clause: tesis2017366\_naturaleza\_condicional}

\textit{Conditional nature: Good faith status depends on continued diligence and lack of knowledge of defects}

\textbf{Intermediate Representation:}

\begin{quote}
good faith exercise is obligated
\end{quote}

\textbf{Witness Clause:}

\begin{lstlisting}[style=witnessstyle, language=Witness]
clause tesis2017366_naturaleza_condicional = oblig(ejercicio_buena_fe) IMPLIES oblig(diligencia_adquirente);
\end{lstlisting}

\subsubsection{Clause: tesis2017366\_posesion\_delictuosa}

\textit{Delictual possession: Criminal possession is always bad faith}

\textbf{Intermediate Representation:}

\begin{quote}
posesion delictuosa is obligated
\end{quote}

\textbf{Witness Clause:}

\begin{lstlisting}[style=witnessstyle, language=Witness]
clause tesis2017366_posesion_delictuosa = oblig(posesion_delictuosa) IMPLIES oblig(ejercicio_mala_fe);
\end{lstlisting}

\newpage

\section{Thesis 2017539: DOMAIN FORFEITURE - INDEPENDENT RESIDENCES AND GOOD FAITH KNOWLEDGE}
\label{thesis:2017539}

\subsection{Original Thesis Text}

\begin{quote}
                                                  Semanario Judicial de la Federación

                                              Tesis

 Registro digital: 2017539

 Instancia: Tribunales            Décima Época                        Materia(s): Civil
 Colegiados de Circuito

 Tesis: I.12o.C.44 C (10a.)       Fuente: Gaceta del Semanario        Tipo: Aislada
                                  Judicial de la Federación.
                                  Libro 57, Agosto de 2018, Tomo III,
                                  página 2805


EXTINCIÓN DE DOMINIO. CUANDO EXISTEN DOMICILIOS INDEPENDIENTES DENTRO DE UN
MISMO INMUEBLE, CORRESPONDE AL ÓRGANO JURISDICCIONAL ESTABLECER SI EL
TERCERO AFECTADO QUE ALEGA BUENA FE, SABÍA QUE SU INMUEBLE SE UTILIZABA
PARA COMETER UN HECHO ILÍCITO.


Es un hecho notorio por los datos de la "Encuesta Intercensal 2015" del Instituto Nacional de
Estadística y Geografía que señalan que veintiocho de cada cien hogares en México son ampliados
[un hogar nuclear más otros parientes (tíos, primos, hermanos, suegros, etcétera) y uno de cada
cien hogares son compuestos (constituido por un hogar nuclear o ampliado, más personas sin
parentesco con el jefe del hogar)], por lo que diversas familias pueden convivir en un mismo
inmueble pero en viviendas separadas con accesos independientes que les permiten delimitar cierto
espacio de intimidad. En este sentido, puede considerarse que residen en domicilios diversos, sin
que obste para llegar a esta conclusión que la numeración oficial no se haya definido de esa forma,
ni se haya constituido el régimen de propiedad en condominio, pues la ausencia de esas
formalidades no puede ocultar que en la realidad se ha generalizado la existencia de dichas formas
de coexistencia propias de algún grupo social del país. Ya sea que tengan "forma de vecindad", o
que consistan en la construcción de cuartos con accesos propios, es evidente que pueden
considerarse domicilios independientes para efectos legales. Así, la existencia de viviendas
construidas, incluso, sin una autorización de construcción, genera conjuntos habitacionales de
hecho, que tienen consecuencias en el mundo del derecho; sin embargo, para la acción de extinción
de dominio lo relevante es si físicamente era posible o no, que el dueño tuviera o debiera tener
conocimiento de que su inmueble se utilizaba para un hecho ilícito; para lo cual la presunción
humana es idónea. Por tanto, corresponde al órgano jurisdiccional establecer si el tercero afectado
que alega buena fe habitaba en el mismo inmueble y si por su superficie y la de las viviendas, así
como su intercomunicación o total independencia, era realmente muy difícil o improbable que
advirtiera que su inmueble se utilizaba para cometer un hecho ilícito.

DÉCIMO SEGUNDO TRIBUNAL COLEGIADO EN MATERIA CIVIL DEL PRIMER CIRCUITO.


Amparo directo 523/2017. Gobierno de la Ciudad de México. 6 de abril de 2018. Unanimidad de
votos. Ponente: Neófito López Ramos. Secretario: Hugo Alfonso Carreón Muñoz.


Esta tesis se publicó el viernes 10 de agosto de 2018 a las 10:18 horas en el Semanario Judicial


                                             Pág. 1 de 2           Fecha de impresión 18/08/2025
 https://sjf2.scjn.gob.mx/detalle/tesis/2017539
                                                  Semanario Judicial de la Federación

de la Federación.




                                             Pág. 2 de 2       Fecha de impresión 18/08/2025
 https://sjf2.scjn.gob.mx/detalle/tesis/2017539

\end{quote}

\subsection{Formalization}

\subsubsection{Clause: tesis2017539\_principio\_fundamental}

\textit{Additional assets for domain forfeiture (all actions now universal) Faith exclusivity: Good faith vs knowledge of illegal activities (mutually exclusive) Thesis-specific principle: Courts must determine if third party with good faith claim could reasonably know about illegal use}

\textbf{Intermediate Representation:}

\begin{quote}
good faith exercise is obligated and accion extincion dominio is claimed
\end{quote}

\textbf{Witness Clause:}

\begin{lstlisting}[style=witnessstyle, language=Witness]
clause tesis2017539_principio_fundamental = oblig(ejercicio_buena_fe) AND claim(accion_extincion_dominio) IMPLIES oblig(posibilidad_fisica_conocimiento);
\end{lstlisting}

\subsubsection{Clause: tesis2017539\_proteccion\_independencia}

\textit{Independent occupation protection: Separate spaces within same property can support good faith defense}

\textbf{Intermediate Representation:}

\begin{quote}
ocupacion independiente is obligated
\end{quote}

\textbf{Witness Clause:}

\begin{lstlisting}[style=witnessstyle, language=Witness]
clause tesis2017539_proteccion_independencia = oblig(ocupacion_independiente) IMPLIES oblig(ejercicio_buena_fe);
\end{lstlisting}

\subsubsection{Clause: tesis2017539\_estandar\_conocimiento}

\textit{Knowledge standard: Physical possibility determines whether owner should have known about illegal activities}

\textbf{Intermediate Representation:}

\begin{quote}
If posibilidad fisica conocimiento is obligated, then (bad faith exercise is obligated or good faith exercise is obligated)
\end{quote}

\textbf{Witness Clause:}

\begin{lstlisting}[style=witnessstyle, language=Witness]
clause tesis2017539_estandar_conocimiento = oblig(posibilidad_fisica_conocimiento) IMPLIES (oblig(ejercicio_mala_fe) OR oblig(ejercicio_buena_fe));
\end{lstlisting}

\subsubsection{Clause: tesis2017539\_presuncion\_aplicable}

\textit{Human presumption use: Courts can use reasonable inferences about what people would likely know}

\textbf{Intermediate Representation:}

\begin{quote}
accion extincion dominio is claimed
\end{quote}

\textbf{Witness Clause:}

\begin{lstlisting}[style=witnessstyle, language=Witness]
clause tesis2017539_presuncion_aplicable = claim(accion_extincion_dominio) IMPLIES oblig(presuncion_humana);
\end{lstlisting}

\subsubsection{Clause: tesis2017539\_analisis\_factual}

\textit{Factual analysis requirement: Courts must examine property layout and intercommunication}

\textbf{Intermediate Representation:}

\begin{quote}
presuncion humana is obligated
\end{quote}

\textbf{Witness Clause:}

\begin{lstlisting}[style=witnessstyle, language=Witness]
clause tesis2017539_analisis_factual = oblig(presuncion_humana) IMPLIES oblig(posibilidad_fisica_conocimiento);
\end{lstlisting}

\subsubsection{Clause: tesis2017539\_defensa\_buena\_fe}

\textit{Good faith defense: Third parties can defend against forfeiture if they lacked knowledge of illegal use}

\textbf{Intermediate Representation:}

\begin{quote}
good faith exercise is obligated and ocupacion independiente is obligated
\end{quote}

\textbf{Witness Clause:}

\begin{lstlisting}[style=witnessstyle, language=Witness]
clause tesis2017539_defensa_buena_fe = oblig(ejercicio_buena_fe) AND oblig(ocupacion_independiente) IMPLIES not(claim(accion_extincion_dominio));
\end{lstlisting}

\subsubsection{Clause: tesis2017539\_carga\_probatoria}

\textit{Burden of proof: Government must prove third party knew or should have known about illegal activities}

\textbf{Intermediate Representation:}

\begin{quote}
If accion extincion dominio is claimed, then (bad faith exercise is obligated or posibilidad fisica conocimiento is obligated)
\end{quote}

\textbf{Witness Clause:}

\begin{lstlisting}[style=witnessstyle, language=Witness]
clause tesis2017539_carga_probatoria = claim(accion_extincion_dominio) IMPLIES (oblig(ejercicio_mala_fe) OR oblig(posibilidad_fisica_conocimiento));
\end{lstlisting}

\newpage

\section{Thesis 2017540: DOMAIN FORFEITURE - DYNAMIC BURDEN OF PROOF AND BAD FAITH STANDARD}
\label{thesis:2017540}

\subsection{Original Thesis Text}

\begin{quote}
                                                   Semanario Judicial de la Federación

                                               Tesis

 Registro digital: 2017540

 Instancia: Tribunales             Décima Época                          Materia(s): Civil
 Colegiados de Circuito

 Tesis: I.12o.C.45 C (10a.)        Fuente: Gaceta del Semanario        Tipo: Aislada
                                   Judicial de la Federación.
                                   Libro 57, Agosto de 2018, Tomo III,
                                   página 2806


EXTINCIÓN DE DOMINIO. ESTA ACCIÓN PROCEDE CONTRA EL TERCERO AFECTADO SI SE
ACREDITA QUE EXISTIÓ MALA FE, ES DECIR, SI TENÍA CONOCIMIENTO DE QUE SUS
BIENES ERAN UTILIZADOS PARA LA COMISIÓN DE DELITOS Y NO LO NOTIFICÓ A LA
AUTORIDAD O HIZO ALGO PARA IMPEDIRLO (CARGA DINÁMICA O SUCESIVA DE LA
PRUEBA).


Cuando los bienes objeto de la acción de extinción de dominio son utilizados para la comisión de
delitos por un tercero, resulta necesario el estudio del concepto de buena y mala fe. En este
supuesto, la litis del juicio de extinción de dominio consiste en definir si existió mala fe del
propietario del inmueble. Esto es necesario, porque dicho tercero resulta afectado al tener un
derecho real sobre los bienes materia de la acción de extinción de dominio y no existen evidencias
de que haya participado en los hechos delictivos, por lo que únicamente puede proceder la acción
de extinción de dominio en su contra si se acredita que existió mala fe, es decir, si tuvo
conocimiento de que sus bienes eran utilizados para la comisión de delitos por un tercero y no lo
notificó a la autoridad o hizo algo para impedirlo, o bien, porque debió tener conocimiento de los
ilícitos; por tanto, el Ministerio Público debe aportar datos que razonablemente permitan sostener
que los delitos se realizaron con conocimiento del propietario de los bienes. Ahora bien, la
legislación en materia de extinción de dominio sostiene que corresponde al tercero afectado probar
su buena fe, al tratarse de un principio inherente a la conducta de la persona que se presume es
honesta, diligente y se conduce generalmente con respeto a la ley y a las buenas costumbres; de
ahí que es necesaria la existencia de elementos que acrediten que los hechos ilícitos se realizaron
con su conocimiento o debieron ser de su conocimiento y así arrojar sobre la persona la carga en
juicio de aportar elementos de prueba que abonen sobre la buena fe y que no pudo tener
conocimiento de los hechos ilícitos. Así, la carga probatoria de que se trata se actualiza hasta que el
representante social ha ofrecido elementos o datos que permitan presumir su conocimiento sobre
los actos delictivos; por tanto, la oportunidad de acreditar la buena fe se presenta cuando se da al
afectado la posibilidad de desvirtuar los datos o elementos que haya aportado previamente el
Ministerio Público. Esta forma de comprender las cargas probatorias en el juicio de extinción de
dominio lleva a reconocer la existencia de una carga dinámica o sucesiva de la prueba. El concepto
de mala fe se traduce en que el afectado tuviera o debiera tener conocimiento de los hechos ilícitos
que ocurrían dentro del inmueble de su propiedad y no en un mero actuar negligente. Es pertinente
aclarar que cuando la jurisprudencia expresa que el tercero "deba" tener conocimiento del hecho
ilícito, no se refiere a un deber abstracto inherente al derecho de propiedad, sino a una probabilidad
fáctica que haga suponer razonablemente que podía enterarse de los hechos ilícitos.

DÉCIMO SEGUNDO TRIBUNAL COLEGIADO EN MATERIA CIVIL DEL PRIMER CIRCUITO.


                                              Pág. 1 de 2             Fecha de impresión 18/08/2025
  https://sjf2.scjn.gob.mx/detalle/tesis/2017540
                                                  Semanario Judicial de la Federación



Amparo directo 523/2017. Gobierno de la Ciudad de México. 6 de abril de 2018. Unanimidad de
votos. Ponente: Neófito López Ramos. Secretario: Hugo Alfonso Carreón Muñoz.


Esta tesis se publicó el viernes 10 de agosto de 2018 a las 10:18 horas en el Semanario Judicial
de la Federación.




                                             Pág. 2 de 2         Fecha de impresión 18/08/2025
 https://sjf2.scjn.gob.mx/detalle/tesis/2017540

\end{quote}

\subsection{Formalization}

\subsubsection{Clause: tesis2017540\_principio\_fundamental}

\textit{All actions now universal - no thesis-specific actions needed Faith exclusivity: Good faith vs bad faith (mutually exclusive) Thesis-specific principle: Domain forfeiture against third parties requires proving bad faith (knowledge + failure to act)}

\textbf{Intermediate Representation:}

\begin{quote}
If accion extincion dominio is claimed, then (bad faith exercise is obligated and cumplimiento deber legal is obligated)
\end{quote}

\textbf{Witness Clause:}

\begin{lstlisting}[style=witnessstyle, language=Witness]
clause tesis2017540_principio_fundamental = claim(accion_extincion_dominio) IMPLIES (oblig(ejercicio_mala_fe) AND oblig(cumplimiento_deber_legal));
\end{lstlisting}

\subsubsection{Clause: tesis2017540\_estandar\_mala\_fe}

\textit{Knowledge standard: Bad faith requires actual or constructive knowledge of illegal activities}

\textbf{Intermediate Representation:}

\begin{quote}
If bad faith exercise is obligated, then (conocimiento problemas legales is obligated or posibilidad fisica conocimiento is obligated)
\end{quote}

\textbf{Witness Clause:}

\begin{lstlisting}[style=witnessstyle, language=Witness]
clause tesis2017540_estandar_mala_fe = oblig(ejercicio_mala_fe) IMPLIES (oblig(conocimiento_problemas_legales) OR oblig(posibilidad_fisica_conocimiento));
\end{lstlisting}

\subsubsection{Clause: tesis2017540\_deber\_actuar}

\textit{Duty to act: Knowledge creates obligation to notify authorities or prevent illegal activities}

\textbf{Intermediate Representation:}

\begin{quote}
conocimiento problemas legales is obligated
\end{quote}

\textbf{Witness Clause:}

\begin{lstlisting}[style=witnessstyle, language=Witness]
clause tesis2017540_deber_actuar = oblig(conocimiento_problemas_legales) IMPLIES oblig(cumplimiento_deber_legal);
\end{lstlisting}

\subsubsection{Clause: tesis2017540\_carga\_inicial}

\textit{Dynamic burden of proof: Initial burden on prosecutor to show evidence of knowledge}

\textbf{Intermediate Representation:}

\begin{quote}
accion extincion dominio is claimed
\end{quote}

\textbf{Witness Clause:}

\begin{lstlisting}[style=witnessstyle, language=Witness]
clause tesis2017540_carga_inicial = claim(accion_extincion_dominio) IMPLIES oblig(elementos_probatorios);
\end{lstlisting}

\subsubsection{Clause: tesis2017540\_carga\_dinamica}

\textit{Burden shift: Once prosecutor shows evidence, third party must prove good faith}

\textbf{Intermediate Representation:}

\begin{quote}
elementos probatorios is obligated
\end{quote}

\textbf{Witness Clause:}

\begin{lstlisting}[style=witnessstyle, language=Witness]
clause tesis2017540_carga_dinamica = oblig(elementos_probatorios) IMPLIES oblig(carga_dinamica_prueba);
\end{lstlisting}

\subsubsection{Clause: tesis2017540\_presuncion\_buena\_fe}

\textit{Presumption of good faith: Legal presumption favors honest, diligent conduct}

\textbf{Intermediate Representation:}

\begin{quote}
carga dinamica prueba is obligated
\end{quote}

\textbf{Witness Clause:}

\begin{lstlisting}[style=witnessstyle, language=Witness]
clause tesis2017540_presuncion_buena_fe = oblig(carga_dinamica_prueba) IMPLIES oblig(ejercicio_buena_fe);
\end{lstlisting}

\subsubsection{Clause: tesis2017540\_probabilidad\_factica}

\textit{Factual probability standard: "Should have known" refers to factual probability, not abstract duty}

\textbf{Intermediate Representation:}

\begin{quote}
posibilidad fisica conocimiento is obligated
\end{quote}

\textbf{Witness Clause:}

\begin{lstlisting}[style=witnessstyle, language=Witness]
clause tesis2017540_probabilidad_factica = oblig(posibilidad_fisica_conocimiento) IMPLIES oblig(presuncion_humana);
\end{lstlisting}

\subsubsection{Clause: tesis2017540\_defensa\_exitosa}

\textit{Good faith defense success: Third party can refute prosecutor's evidence}

\textbf{Intermediate Representation:}

\begin{quote}
good faith exercise is obligated and carga dinamica prueba is obligated
\end{quote}

\textbf{Witness Clause:}

\begin{lstlisting}[style=witnessstyle, language=Witness]
clause tesis2017540_defensa_exitosa = oblig(ejercicio_buena_fe) AND oblig(carga_dinamica_prueba) IMPLIES not(claim(accion_extincion_dominio));
\end{lstlisting}

\subsubsection{Clause: tesis2017540\_requisito\_conocimiento}

\textit{Burden of proof completion: Prosecutor must establish knowledge before burden shifts}

\textbf{Intermediate Representation:}

\begin{quote}
carga dinamica prueba is obligated
\end{quote}

\textbf{Witness Clause:}

\begin{lstlisting}[style=witnessstyle, language=Witness]
clause tesis2017540_requisito_conocimiento = oblig(carga_dinamica_prueba) IMPLIES oblig(elementos_probatorios);
\end{lstlisting}

\newpage

\section{Thesis 2018505: ADVERSE POSSESSION - ALLANAMIENTO CANNOT DEMONSTRATE POSSESSION ATTRIBUTES}
\label{thesis:2018505}

\subsection{Original Thesis Text}

\begin{quote}
                                                   Semanario Judicial de la Federación

                                               Tesis

 Registro digital: 2018505

 Instancia: Plenos de Circuito     Décima Época                         Materia(s): Civil

 Tesis: PC.XVII. J/17 C (10a.)     Fuente: Semanario Judicial de la     Tipo: Jurisprudencia
                                   Federación.


PRESCRIPCIÓN POSITIVA. EL ALLANAMIENTO A LA DEMANDA NO ES APTO PARA
DEMOSTRAR LOS "ATRIBUTOS DE LA POSESIÓN" (LEGISLACIÓN DEL ESTADO DE
CHIHUAHUA).


La prescripción positiva es una institución del derecho civil de orden público, que dota de seguridad
jurídica a los poseedores de un bien que acrediten los "atributos de la posesión", en términos de los
artículos 1153 y 1154 del Código Civil del Estado de Chihuahua, esto es, de manera pública,
pacífica, continua de buena fe y en un lapso suficiente; de modo que, si al contestar la demanda el
enjuiciado se allana a las pretensiones del actor, ese reconocimiento sólo produce el acreditamiento
de la causa generadora de la posesión a título de dueño, pero no es apto para demostrar los
atributos de la posesión, pues las cualidades de ésta no son hechos propios del demandado, por lo
que no se releva al actor de probar los hechos intrínsecos y fundatorios de su pretensión; de ahí
que le corresponda probar los demás elementos constitutivos de su acción, para no afectar
derechos de terceros.

PLENO DEL DECIMOSÉPTIMO CIRCUITO.


Contradicción de tesis 4/2017. Entre las sustentadas por el Primer Tribunal Colegiado del Décimo
Séptimo Circuito, actualmente Primer Tribunal Colegiado en Materias Civil y de Trabajo del Décimo
Séptimo Circuito, y el Segundo Tribunal Colegiado en Materias Civil y de Trabajo del Décimo
Séptimo Circuito. 9 de octubre de 2018. Mayoría de seis votos de los Magistrados María Teresa
Zambrano Calero, Refugio Noel Montoya Moreno, Juan Carlos Zamora Tejeda, José Raymundo
Cornejo Olvera, Presidente del Pleno de Circuito, José de Jesús González Ruiz y Rogelio Alberto
Montoya Rodríguez. Disidente: Magistrada María del Carmen Cordero Martínez, quien formuló voto
particular. Ponente: María Teresa Zambrano Calero. Secretaria: Emma Margarita Aréchiga
Rodríguez.

Tesis y criterios contendientes:

Tesis XVII.1o.30 C, de rubro: "PRESCRIPCIÓN POSITIVA. EL ALLANAMIENTO A LA DEMANDA
ES APTO PARA DEMOSTAR LA POSESIÓN POR PARTE DEL ACTOR Y SUS DEMÁS
ATRIBUTOS (LEGISLACIÓN DEL ESTADO DE CHIHUAHUA).", aprobada por el Primer Tribunal
Colegiado del Décimo Séptimo Circuito y publicada en el Semanario Judicial de la Federación y su
Gaceta, Novena Época, Tomo XV, abril de 2002, página 1313, y

El sustentado por el Primer Tribunal Colegiado en Materias Civil y de Trabajo del Décimo Séptimo
Circuito, al resolver los amparos directos 208/2017 y 265/2017, y el diverso sustentado por el
Segundo Tribunal Colegiado en Materias Civil y de Trabajo del Décimo Séptimo Circuito, al resolver

                                              Pág. 1 de 2             Fecha de impresión 01/10/2025
  https://sjf2.scjn.gob.mx/detalle/tesis/2018505
                                                   Semanario Judicial de la Federación

el amparo directo civil 160/2017.


Esta tesis se publicó el viernes 30 de noviembre de 2018 a las 10:41 horas en el Semanario
Judicial de la Federación y, por ende, se considera de aplicación obligatoria a partir del lunes 03 de
diciembre de 2018, para los efectos previstos en el punto séptimo del Acuerdo General Plenario
19/2013.




                                              Pág. 2 de 2            Fecha de impresión 01/10/2025
  https://sjf2.scjn.gob.mx/detalle/tesis/2018505

\end{quote}

\subsection{Formalization}

\subsubsection{Clause: tesis2018505\_principio\_fundamental}

\textit{Additional actions for this thesis (reusing most existing assets) Additional assets for public order and legal security (reusing most existing assets) Faith exclusivity: Possession is either good faith XOR bad faith (mutually exclusive) Thesis-specific principle: Allanamiento can prove generating cause but NOT possession attributes (public, peaceful, continuous)}

\textbf{Intermediate Representation:}

\begin{quote}
defendant acquiescence is obligated and not (peaceful, continuous, and public possession is obligated)
\end{quote}

\textbf{Witness Clause:}

\begin{lstlisting}[style=witnessstyle, language=Witness]
clause tesis2018505_principio_fundamental = oblig(allanamiento_demandado) IMPLIES oblig(causa_generadora_probada) AND not(oblig(posesion_pacifico_continuo_publico));
\end{lstlisting}

\subsubsection{Clause: tesis2018505\_orden\_publico}

\textit{Public order institution: Adverse possession provides legal security for possessors who prove possession attributes}

\textbf{Intermediate Representation:}

\begin{quote}
If adverse possession claim is claimed, then (seguridad juridica is obligated and peaceful, continuous, and public possession is obligated)
\end{quote}

\textbf{Witness Clause:}

\begin{lstlisting}[style=witnessstyle, language=Witness]
clause tesis2018505_orden_publico = claim(demanda_prescripcion) IMPLIES (oblig(seguridad_juridica) AND oblig(posesion_pacifico_continuo_publico));
\end{lstlisting}

\subsubsection{Clause: tesis2018505\_prueba\_independiente}

\textit{Independent proof requirement: Plaintiff must prove possession attributes independently despite defendant's admission}

\textbf{Intermediate Representation:}

\begin{quote}
defendant acquiescence is obligated
\end{quote}

\textbf{Witness Clause:}

\begin{lstlisting}[style=witnessstyle, language=Witness]
clause tesis2018505_prueba_independiente = oblig(allanamiento_demandado) IMPLIES oblig(prueba_independiente);
\end{lstlisting}

\subsubsection{Clause: tesis2018505\_hechos\_intrinsecos}

\textit{Intrinsic facts limitation: Possession qualities are intrinsic to plaintiff, defendant cannot admit them}

\textbf{Intermediate Representation:}

\begin{quote}
peaceful, continuous, and public possession is obligated
\end{quote}

\textbf{Witness Clause:}

\begin{lstlisting}[style=witnessstyle, language=Witness]
clause tesis2018505_hechos_intrinsecos = oblig(posesion_pacifico_continuo_publico) IMPLIES oblig(prueba_independiente);
\end{lstlisting}

\subsubsection{Clause: tesis2018505\_proteccion\_terceros}

\textit{Third party protection: Independent proof required to protect third party rights}

\textbf{Intermediate Representation:}

\begin{quote}
third party protection is obligated
\end{quote}

\textbf{Witness Clause:}

\begin{lstlisting}[style=witnessstyle, language=Witness]
clause tesis2018505_proteccion_terceros = oblig(proteccion_terceros) IMPLIES oblig(prueba_independiente);
\end{lstlisting}

\subsubsection{Clause: tesis2018505\_carga\_probatoria}

\textit{Burden of proof maintenance: Allanamiento does not relieve plaintiff of proving foundational facts}

\textbf{Intermediate Representation:}

\begin{quote}
If adverse possession claim is claimed, then (proven generating cause is obligated and peaceful, continuous, and public possession is obligated)
\end{quote}

\textbf{Witness Clause:}

\begin{lstlisting}[style=witnessstyle, language=Witness]
clause tesis2018505_carga_probatoria = claim(demanda_prescripcion) IMPLIES (oblig(causa_generadora_probada) AND oblig(posesion_pacifico_continuo_publico));
\end{lstlisting}

\subsubsection{Clause: tesis2018505\_proposito\_seguridad}

\textit{Legal security purpose: System protects possessors who meet all requirements}

\textbf{Intermediate Representation:}

\begin{quote}
peaceful, continuous, and public possession is obligated and good faith exercise is obligated and temporal possession is obligated
\end{quote}

\textbf{Witness Clause:}

\begin{lstlisting}[style=witnessstyle, language=Witness]
clause tesis2018505_proposito_seguridad = oblig(posesion_pacifico_continuo_publico) AND oblig(ejercicio_buena_fe) AND oblig(posesion_temporal) IMPLIES oblig(seguridad_juridica);
\end{lstlisting}

\subsubsection{Clause: tesis2018505\_requisito\_temporal}

\textit{Temporal requirement: Must prove sufficient time period along with other attributes}

\textbf{Intermediate Representation:}

\begin{quote}
If adverse possession claim is claimed, then (temporal possession is obligated and peaceful, continuous, and public possession is obligated)
\end{quote}

\textbf{Witness Clause:}

\begin{lstlisting}[style=witnessstyle, language=Witness]
clause tesis2018505_requisito_temporal = claim(demanda_prescripcion) IMPLIES (oblig(posesion_temporal) AND oblig(posesion_pacifico_continuo_publico));
\end{lstlisting}

\newpage

\section{Thesis 2020418: ADVERSE POSSESSION AND UNREGISTERED CONTRACTS}
\label{thesis:2020418}

\subsection{Original Thesis Text}

\begin{quote}
                                                   Semanario Judicial de la Federación

                                                Tesis

 Registro digital: 2020418

 Instancia: Tribunales             Décima Época                            Materia(s): Civil
 Colegiados de Circuito

 Tesis: VII.2o.C.195 C (10a.)      Fuente: Semanario Judicial de la        Tipo: Aislada
                                   Federación.


PRESCRIPCIÓN POSITIVA. EL CONTRATO DE COMPRAVENTA OTORGADO EN ESCRITURA
PÚBLICA Y NO INSCRITO EN EL REGISTRO PÚBLICO DE LA PROPIEDAD Y DE COMERCIO,
NO ES OPONIBLE AL TERCERO QUE APARECE COMO TITULAR REGISTRAL NI ES APTO
PARA ACREDITAR LA CAUSA GENERADORA DE LA POSESIÓN DE AQUELLA ACCIÓN
(LEGISLACIÓN DEL ESTADO DE VERACRUZ).


Conforme a los artículos 862, 1168, 1184, 1185 y 1189 del Código Civil para el Estado de Veracruz,
cuando la acción de prescripción positiva se ejerce bajo la hipótesis de que se es poseedor de
buena fe del inmueble, al ser uno de sus requisitos que la posesión se tenga con carácter de dueño,
en términos del artículo 228 del Código de Procedimientos Civiles para el Estado de Veracruz, al
accionante le corresponde demostrar la causa generadora de su posesión a título de propietario, es
decir, el justo título en virtud del cual entró en posesión del inmueble, en el entendido de que éste se
trata de un acto traslativo de dominio imperfecto, pero suficiente para transferirle el dominio, como
puede ser, entre otros, un contrato de compraventa. Ahora bien, el sujeto pasivo de la acción de
prescripción positiva de un inmueble lo es quien aparece como su propietario ante el Registro
Público de la Propiedad y de Comercio. En tales condiciones, cuando el actor en un juicio de
prescripción positiva exhibe un contrato de compraventa celebrado con un tercero y no por quien
aparece como propietario en el Registro Público de la Propiedad y de Comercio, otorgado conforme
al artículo 2255 del Código Civil en escritura pública, el cual debió inscribirse, éste no es apto para
tener por acreditada la causa generadora de la posesión, porque ante la falta de inscripción, no
resulta oponible contra el titular registral que no participó en la celebración del acto jurídico. Ello es
así, porque conforme a los artículos 2935, fracción I, 2936 y 2944, fracción I, del Código Civil citado,
en el Registro Público de la Propiedad deben inscribirse los documentos públicos que contengan los
títulos relativos a la posesión o dominio de inmuebles; y la sanción en caso de que no se haga, es
que dichos actos sólo surtirán efectos entre quienes los otorgaron, mas no podrán producir
perjuicios contra terceros ajenos a su celebración, quienes sí podrán aprovecharlos en cuanto les
resulten favorables. Por tanto, la falta de observancia a dichas disposiciones, sólo puede repercutir
al titular de los derechos de dominio que debieron inscribirse, mas no a terceros. Por otro lado, el
contrato de compraventa otorgado en escritura pública no inscrito en el Registro Público de la
Propiedad y de Comercio, además de no ser oponible al tercero que aparece como titular registral
del inmueble que se pretende prescribir, tampoco es apto para acreditar la causa generadora de la
posesión, porque al haberse otorgado con las formalidades de ley, se trata de un título perfecto que
transmite el dominio de forma plena y es inscribible en el Registro Público de la Propiedad y de
Comercio, por lo que la prescripción positiva como forma no convencional de adquirir el dominio
desaparece. Por tanto, el derecho de propiedad que pudiera tener quien pretendió usucapir en
relación con un inmueble inscrito a favor de un tercero deberá dilucidarse a través de la forma y vía
que correspondan, pero no por la acción de prescripción positiva, máxime que uno de los efectos de
ésta es la obtención de un título susceptible de inscribirse en el Registro Público de la Propiedad y


                                               Pág. 1 de 2              Fecha de impresión 01/10/2025
  https://sjf2.scjn.gob.mx/detalle/tesis/2020418
                                                  Semanario Judicial de la Federación

de Comercio, por lo que no puede tener por objeto subsanar la falta de observancia de las
disposiciones registrales, en perjuicio de un tercero que sí actuó con la debida diligencia al haber
inscrito su derecho de propiedad.

SEGUNDO TRIBUNAL COLEGIADO EN MATERIA CIVIL DEL SÉPTIMO CIRCUITO.


Amparo directo 660/2018. 9 de mayo de 2019. Mayoría de votos. Disidente: José Manuel De Alba
De Alba. Ponente: Isidro Pedro Alcántara Valdés. Secretario: Flavio Bernardo Galván Zilli.


Esta tesis se publicó el viernes 16 de agosto de 2019 a las 10:24 horas en el Semanario Judicial
de la Federación.




                                             Pág. 2 de 2            Fecha de impresión 01/10/2025
 https://sjf2.scjn.gob.mx/detalle/tesis/2020418

\end{quote}

\subsection{Formalization}

\subsubsection{Clause: tesis2020418\_principio\_fundamental}

\textit{Faith exclusivity: Possession is either good faith XOR bad faith (mutually exclusive) Thesis-specific principle: Unregistered contracts cannot establish adverse possession against registered owners}

\textbf{Intermediate Representation:}

\begin{quote}
registered ownership is obligated and unregistered contract is obligated
\end{quote}

\textbf{Witness Clause:}

\begin{lstlisting}[style=witnessstyle, language=Witness]
clause tesis2020418_principio_fundamental = oblig(titularidad_registral) AND oblig(contrato_no_inscrito) IMPLIES not(claim(demanda_prescripcion));
\end{lstlisting}

\subsubsection{Clause: tesis2020418\_proteccion\_registral}

\textit{Supporting legal logic}

\textbf{Intermediate Representation:}

\begin{quote}
registered ownership is obligated
\end{quote}

\textbf{Witness Clause:}

\begin{lstlisting}[style=witnessstyle, language=Witness]
clause tesis2020418_proteccion_registral = oblig(titularidad_registral) IMPLIES claim(oponibilidad_terceros);
\end{lstlisting}

\subsubsection{Clause: tesis2020418\_limitacion\_no\_inscripcion}

\textbf{Intermediate Representation:}

\begin{quote}
unregistered contract is obligated
\end{quote}

\textbf{Witness Clause:}

\begin{lstlisting}[style=witnessstyle, language=Witness]
clause tesis2020418_limitacion_no_inscripcion = oblig(contrato_no_inscrito) IMPLIES not(oblig(oponibilidad_terceros));
\end{lstlisting}

\subsubsection{Clause: tesis2020418\_consecuencia\_logica}

\textbf{Intermediate Representation:}

\begin{quote}
not (opposability to third parties is obligated)
\end{quote}

\textbf{Witness Clause:}

\begin{lstlisting}[style=witnessstyle, language=Witness]
clause tesis2020418_consecuencia_logica = not(oblig(oponibilidad_terceros)) IMPLIES not(claim(demanda_prescripcion));
\end{lstlisting}

\newpage

\section{Thesis 2020599: PARCELAMIENTO ECONÓMICO O "DE HECHO" SOBRE TIERRAS DE USO COMÚN. ES NECESARIO QUE EL REGISTRO AGRARIO NACIONAL LO VALIDE Y EXPIDA LOS}
\label{thesis:2020599}

\subsection{Original Thesis Text}

\begin{quote}
                                                   Semanario Judicial de la Federación

                                                Tesis

 Registro digital: 2020599

 Instancia: Tribunales             Décima Época                           Materia(s): Administrativa
 Colegiados de Circuito

 Tesis: XVI.1o.A.193 A (10a.)      Fuente: Semanario Judicial de la       Tipo: Aislada
                                   Federación.


PARCELAMIENTO ECONÓMICO O "DE HECHO" SOBRE TIERRAS DE USO COMÚN. ES
NECESARIO QUE EL REGISTRO AGRARIO NACIONAL LO VALIDE Y EXPIDA LOS
CERTIFICADOS CORRESPONDIENTES, PARA QUE PUEDA RECONOCERSE A UN
POSESIONARIO COMO TITULAR DE LA PARCELA CON MOTIVO DE LA PRESCRIPCIÓN
ADQUISITIVA.


La prescripción a que alude el artículo 48 de la Ley Agraria se limita a las tierras parceladas por la
asamblea general de ejidatarios, cuando se cumplen los presupuestos siguientes: a) se trate de un
posesionario; b) la posesión se ejerza "en concepto de titular de derechos de ejidatario"; y, c) la
posesión sea pacífica, continua y pública durante cinco años, si es de buena fe, o de diez si fuera
de mala fe. Ahora bien, de conformidad con el artículo 23, fracción VIII, del mismo ordenamiento, la
asamblea de ejidatarios está facultada para reconocer el parcelamiento económico o "de hecho"
realizado al interior del ejido, el cual surte efectos jurídicos una vez que el Registro Agrario Nacional
lo valide y expida los certificados correspondientes, los que se retrotraen al momento en que aquél
se realizó. Así, el lapso entre la determinación de la asamblea de realizar el parcelamiento y la
oficialización por parte del Registro Agrario Nacional, constituye un estado transitorio en el cual la
asignación de tierras sigue siendo "de hecho". Por ende, los posesionarios de tierras de uso común
parceladas económicamente tienen la expectativa de derecho para obtener el certificado
correspondiente, pero no es jurídicamente factible reconocerlos como titulares de la parcela con
motivo de la prescripción adquisitiva, antes de que ocurra el parcelamiento formal; de ahí que pueda
computarse el plazo para que opere la prescripción durante el parcelamiento económico, pero
solamente si existe una validación por parte de la autoridad agraria mencionada, porque puede
darse el caso de que ésta no lo reconozca, anulándose todos los efectos jurídicos que puede tener
la repartición efectuada por la asamblea.

PRIMER TRIBUNAL COLEGIADO EN MATERIA ADMINISTRATIVA DEL DÉCIMO SEXTO
CIRCUITO.


Amparo directo 687/2018. Ricardo Saldaña Gómez. 22 de marzo de 2019. Unanimidad de votos.
Ponente: Ariel Alberto Rojas Caballero. Secretario: Javier Cruz Vázquez.


Esta tesis se publicó el viernes 13 de septiembre de 2019 a las 10:22 horas en el Semanario
Judicial de la Federación.




                                              Pág. 1 de 2              Fecha de impresión 01/10/2025
  https://sjf2.scjn.gob.mx/detalle/tesis/2020599
                                                 Semanario Judicial de la Federación




                                            Pág. 2 de 2       Fecha de impresión 01/10/2025
https://sjf2.scjn.gob.mx/detalle/tesis/2020599

\end{quote}

\subsection{Formalization}

\subsubsection{Clause: tesis2020599\_principio\_fundamental}

\textit{Additional actions for this thesis Additional assets for proof requirements Faith exclusivity: Adverse possession is either good faith XOR bad faith Thesis-specific principle: Good faith adverse possession requires proving authenticity, date, and diligence}

\textbf{Intermediate Representation:}

\begin{quote}
If good faith exercise is obligated and adverse possession claim is claimed, then (autenticidad titulo is obligated and fecha fehaciente is obligated and diligencia adquirente is obligated)
\end{quote}

\textbf{Witness Clause:}

\begin{lstlisting}[style=witnessstyle, language=Witness]
clause tesis2020599_principio_fundamental = oblig(ejercicio_buena_fe) AND claim(demanda_prescripcion) IMPLIES (oblig(autenticidad_titulo) AND oblig(fecha_fehaciente) AND oblig(diligencia_adquirente));
\end{lstlisting}

\subsubsection{Clause: tesis2020599\_prueba\_autenticidad}

\textit{Proof requirements: Must prove title authenticity and date reliably}

\textbf{Intermediate Representation:}

\begin{quote}
If proven generating cause is obligated, then (autenticidad titulo is obligated and fecha fehaciente is obligated)
\end{quote}

\textbf{Witness Clause:}

\begin{lstlisting}[style=witnessstyle, language=Witness]
clause tesis2020599_prueba_autenticidad = oblig(causa_generadora_probada) IMPLIES (oblig(autenticidad_titulo) AND oblig(fecha_fehaciente));
\end{lstlisting}

\subsubsection{Clause: tesis2020599\_diligencia\_buena\_fe}

\textit{Diligence requirement: Good faith requires demonstrating due diligence in acquisition}

\textbf{Intermediate Representation:}

\begin{quote}
good faith exercise is obligated
\end{quote}

\textbf{Witness Clause:}

\begin{lstlisting}[style=witnessstyle, language=Witness]
clause tesis2020599_diligencia_buena_fe = oblig(ejercicio_buena_fe) IMPLIES oblig(diligencia_adquirente);
\end{lstlisting}

\subsubsection{Clause: tesis2020599\_prueba\_pagos}

\textit{Payment proof: Onerous contracts require proving payments made}

\textbf{Intermediate Representation:}

\begin{quote}
proven generating cause is obligated and contrato compraventa is obligated
\end{quote}

\textbf{Witness Clause:}

\begin{lstlisting}[style=witnessstyle, language=Witness]
clause tesis2020599_prueba_pagos = oblig(causa_generadora_probada) AND oblig(contrato_compraventa) IMPLIES oblig(pagos_precio);
\end{lstlisting}

\subsubsection{Clause: tesis2020599\_carga\_prueba}

\textit{Burden of proof: Claimant must prove generating cause with reliable evidence}

\textbf{Intermediate Representation:}

\begin{quote}
adverse possession claim is claimed
\end{quote}

\textbf{Witness Clause:}

\begin{lstlisting}[style=witnessstyle, language=Witness]
clause tesis2020599_carga_prueba = claim(demanda_prescripcion) IMPLIES oblig(causa_generadora_revelada);
\end{lstlisting}

\subsubsection{Clause: tesis2020599\_falta\_diligencia}

\textit{Insufficient diligence consequence: Lack of diligence means bad faith (10-year period)}

\textbf{Intermediate Representation:}

\begin{quote}
not (diligencia adquirente is obligated)
\end{quote}

\textbf{Witness Clause:}

\begin{lstlisting}[style=witnessstyle, language=Witness]
clause tesis2020599_falta_diligencia = not(oblig(diligencia_adquirente)) IMPLIES oblig(ejercicio_mala_fe);
\end{lstlisting}

\newpage

\section{Thesis 2020615: USUCAPIÓN. EL ALLANAMIENTO DE LOS ENJUICIADOS A LA DEMANDA, ES INSUFICIENTE PARA EVIDENCIAR LOS ATRIBUTOS O CUALIDADES DE LA POSESIÓN REQUERIDOS PARA}
\label{thesis:2020615}

\subsection{Original Thesis Text}

\begin{quote}
                                                  Semanario Judicial de la Federación

                                              Tesis

 Registro digital: 2020615

 Instancia: Tribunales            Décima Época                         Materia(s): Civil
 Colegiados de Circuito

 Tesis: II.4o.C.30 C (10a.)       Fuente: Semanario Judicial de la     Tipo: Aislada
                                  Federación.


USUCAPIÓN. EL ALLANAMIENTO DE LOS ENJUICIADOS A LA DEMANDA, ES INSUFICIENTE
PARA EVIDENCIAR LOS ATRIBUTOS O CUALIDADES DE LA POSESIÓN REQUERIDOS PARA
LA PROCEDENCIA DE ESA ACCIÓN (LEGISLACIÓN DEL ESTADO DE MÉXICO ABROGADA).


La usucapión es un medio de adquirir la propiedad de los bienes mediante la posesión de los
mismos, durante el tiempo y con las condiciones establecidas en la ley; no obstante, la posesión
requerida para efectos de la procedencia de dicha acción, en términos de los artículos 910 y 911 del
Código Civil del Estado de México abrogado, requiere de cualidades específicas y concretas, como
es que se tenga en concepto de propietario, de manera pacífica, continua, pública y de buena fe. Si
se trata de bienes inmuebles, la posesión en concepto de propietario debe de ser por un lapso de
cinco años. En esa virtud, el allanamiento de los enjuiciados a la demanda de usucapión, es
insuficiente para evidenciar los atributos o cualidades de la posesión requeridos en esa acción pues,
en su caso, con éste se robustece la causa generadora de la posesión a título de dueño, pero no
resulta apta para demostrar las cualidades de la posesión, puesto que el allanamiento se traduce en
una confesión, la que para ser apta y trascender procesalmente debe ser de hechos propios, en
términos del artículo 1.271, fracción III, del Código de Procedimientos Civiles local, no así sobre
cuestiones ajenas, como es lo relativo a que se tenga una posesión pacífica, continua, pública y de
buena fe; por ende, pese a existir el allanamiento, debe abrirse el juicio a prueba, a fin de que la
actora esté en condiciones de demostrar las cualidades de la posesión ya referidas, en caso
contrario, la acción resultará improcedente.

CUARTO TRIBUNAL COLEGIADO EN MATERIA CIVIL DEL SEGUNDO CIRCUITO.


Amparo directo 890/2018. Agueda Hernández Rebollar. 11 de abril de 2019. Unanimidad de votos.
Ponente: Fernando Sánchez Calderón. Secretario: Antonio Salazar López.


Esta tesis se publicó el viernes 13 de septiembre de 2019 a las 10:22 horas en el Semanario
Judicial de la Federación.




                                             Pág. 1 de 2             Fecha de impresión 01/10/2025
 https://sjf2.scjn.gob.mx/detalle/tesis/2020615
                                                 Semanario Judicial de la Federación




                                            Pág. 2 de 2       Fecha de impresión 01/10/2025
https://sjf2.scjn.gob.mx/detalle/tesis/2020615

\end{quote}

\subsection{Formalization}

\subsubsection{Clause: tesis2020615\_principio\_fundamental}

\textit{Additional actions for this thesis Additional assets for proof requirements Faith exclusivity: Adverse possession is either good faith XOR bad faith Thesis-specific principle: Good faith adverse possession requires proving authenticity, date, and diligence}

\textbf{Intermediate Representation:}

\begin{quote}
If good faith exercise is obligated and adverse possession claim is claimed, then (autenticidad titulo is obligated and fecha fehaciente is obligated and diligencia adquirente is obligated)
\end{quote}

\textbf{Witness Clause:}

\begin{lstlisting}[style=witnessstyle, language=Witness]
clause tesis2020615_principio_fundamental = oblig(ejercicio_buena_fe) AND claim(demanda_prescripcion) IMPLIES (oblig(autenticidad_titulo) AND oblig(fecha_fehaciente) AND oblig(diligencia_adquirente));
\end{lstlisting}

\subsubsection{Clause: tesis2020615\_prueba\_autenticidad}

\textit{Proof requirements: Must prove title authenticity and date reliably}

\textbf{Intermediate Representation:}

\begin{quote}
If proven generating cause is obligated, then (autenticidad titulo is obligated and fecha fehaciente is obligated)
\end{quote}

\textbf{Witness Clause:}

\begin{lstlisting}[style=witnessstyle, language=Witness]
clause tesis2020615_prueba_autenticidad = oblig(causa_generadora_probada) IMPLIES (oblig(autenticidad_titulo) AND oblig(fecha_fehaciente));
\end{lstlisting}

\subsubsection{Clause: tesis2020615\_diligencia\_buena\_fe}

\textit{Diligence requirement: Good faith requires demonstrating due diligence in acquisition}

\textbf{Intermediate Representation:}

\begin{quote}
good faith exercise is obligated
\end{quote}

\textbf{Witness Clause:}

\begin{lstlisting}[style=witnessstyle, language=Witness]
clause tesis2020615_diligencia_buena_fe = oblig(ejercicio_buena_fe) IMPLIES oblig(diligencia_adquirente);
\end{lstlisting}

\subsubsection{Clause: tesis2020615\_prueba\_pagos}

\textit{Payment proof: Onerous contracts require proving payments made}

\textbf{Intermediate Representation:}

\begin{quote}
proven generating cause is obligated and contrato compraventa is obligated
\end{quote}

\textbf{Witness Clause:}

\begin{lstlisting}[style=witnessstyle, language=Witness]
clause tesis2020615_prueba_pagos = oblig(causa_generadora_probada) AND oblig(contrato_compraventa) IMPLIES oblig(pagos_precio);
\end{lstlisting}

\subsubsection{Clause: tesis2020615\_carga\_prueba}

\textit{Burden of proof: Claimant must prove generating cause with reliable evidence}

\textbf{Intermediate Representation:}

\begin{quote}
adverse possession claim is claimed
\end{quote}

\textbf{Witness Clause:}

\begin{lstlisting}[style=witnessstyle, language=Witness]
clause tesis2020615_carga_prueba = claim(demanda_prescripcion) IMPLIES oblig(causa_generadora_revelada);
\end{lstlisting}

\subsubsection{Clause: tesis2020615\_falta\_diligencia}

\textit{Insufficient diligence consequence: Lack of diligence means bad faith (10-year period)}

\textbf{Intermediate Representation:}

\begin{quote}
not (diligencia adquirente is obligated)
\end{quote}

\textbf{Witness Clause:}

\begin{lstlisting}[style=witnessstyle, language=Witness]
clause tesis2020615_falta_diligencia = not(oblig(diligencia_adquirente)) IMPLIES oblig(ejercicio_mala_fe);
\end{lstlisting}

\newpage

\section{Thesis 2021246: Bad Faith Adverse Possession - Proof of Generating Cause}
\label{thesis:2021246}

\subsection{Original Thesis Text}

\begin{quote}
                                                   Semanario Judicial de la Federación

                                                Tesis

 Registro digital: 2021246

 Instancia: Tribunales             Décima Época                           Materia(s): Civil
 Colegiados de Circuito

 Tesis: I.12o.C.148 C (10a.)       Fuente: Semanario Judicial de la       Tipo: Aislada
                                   Federación.


PRESCRIPCIÓN ADQUISITIVA. AUNQUE AL POSEEDOR DE MALA FE NO LE ES EXIGIBLE
QUE DEMUESTRE EL JUSTO TÍTULO COMO BASE DE SU PRETENSIÓN, ES NECESARIO QUE
ACREDITE LA CAUSA GENERADORA DE LA POSESIÓN (LEGISLACIÓN APLICABLE PARA LA
CIUDAD DE MÉXICO).


Conforme a los artículos 1135 y 1136, en relación con los diversos 1151 y 1152, todos del Código
Civil para el Distrito Federal, aplicable para la Ciudad de México, la prescripción positiva es el medio
para adquirir bienes, por el transcurso de cierto tiempo y bajo las condiciones establecidas en la ley,
en el caso de inmuebles, mediante la posesión por cinco años si ésta es de buena fe o por diez
años cuando es de mala fe, y en ambos supuestos dicha posesión debe ser en concepto de
propietario, pacífica, continua y pública. Ahora bien, el artículo 806 del código citado establece que
es poseedor de buena fe el que entra en la posesión en virtud de un título suficiente para darle
derecho a poseer, así como el que ignora los vicios de su título que le impiden poseer con derecho,
en tanto que lo es de mala fe el que entra a la posesión sin título alguno, al igual que el que conoce
los vicios de su título que le impiden poseer con derecho, y dispone que se entiende por título, la
causa generadora de la posesión. Por otra parte, el concepto de propietario comprende al poseedor
con un título objetiva o subjetivamente válido e, incluso, sin título, siempre que demuestre ser el que
tiene el dominio de la cosa y que empezó a poseerla en virtud de una causa que le permite
ostentarse como dueño, la cual es ajena a la buena o mala fe, pues no proviene del fuero interno
del poseedor, sino que la tiene quien entró a poseer mediante un acto o hecho que le permite
ostentarse como dueño, con exclusión de los demás, y resulta relevante porque legalmente no es
apta para usucapir la posesión derivada (a nombre de otro), sino sólo la originaria (en concepto de
dueño) sea jurídica o de hecho, por lo que, además de probar el tiempo por el que
ininterrumpidamente ha poseído (cinco o diez años, según sea el caso), el actor debe demostrar
siempre la causa generadora de la posesión, si es de buena fe precisa acreditar el justo título o el
hecho generador en que basa su pretensión, en tanto que si ésta la sustenta en la posesión por diez
años en calidad de poseedor originario, de hecho y de mala fe, debe probar el hecho generador de
la posesión a título de dueño, esto es, cualquier acto que fundadamente considere bastante para
transferir al poseedor el dominio sobre el bien de que se trate. En ese sentido, si bien no puede
exigirse la acreditación de un justo título cuando la acción relativa se apoye en la posesión de mala
fe, lo cierto es que resulta necesario que el promovente justifique la causa generadora de la
posesión, debido a que la voluntad del legislador, al establecer la usucapión, no fue incentivar el
incumplimiento de las obligaciones o el apoderamiento de bienes ajenos, sino formalizar una
cuestión de hecho, pero sólo cuando sea evidente que el titular del derecho de propiedad no tuvo
interés en conservarlo durante el plazo en que se consumó la prescripción; por tanto, en el caso de
la posesión de mala fe, debe exigirse un estándar probatorio elevado, a fin de que el accionante
revele y acredite en forma fehaciente dicha causa generadora y las calidades de la posesión que
exige la ley, por más de diez años, pues de no ser así, el juzgador estaría imposibilitado para


                                              Pág. 1 de 2             Fecha de impresión 01/10/2025
  https://sjf2.scjn.gob.mx/detalle/tesis/2021246
                                                  Semanario Judicial de la Federación

determinar si la posesión aducida es originaria o derivada, de buena o de mala fe y a partir de qué
momento debe computarse el plazo para prescribir.

DÉCIMO SEGUNDO TRIBUNAL COLEGIADO EN MATERIA CIVIL DEL PRIMER CIRCUITO.


Amparo directo 761/2018. Jorge Manuel Suárez Díaz y otro. 22 de marzo de 2019. Unanimidad de
votos. Ponente: Rómulo Amadeo Figueroa Salmorán. Secretaria: Nancy Michelle Álvarez Díaz
Barriga.


Esta tesis se publicó el viernes 06 de diciembre de 2019 a las 10:18 horas en el Semanario Judicial
de la Federación.




                                             Pág. 2 de 2           Fecha de impresión 01/10/2025
 https://sjf2.scjn.gob.mx/detalle/tesis/2021246

\end{quote}

\subsection{Formalization}

\subsubsection{Clause: tesis2021246\_principio\_fundamental}

\textit{Faith exclusivity: Possession is either good faith XOR bad faith (mutually exclusive) Thesis-specific principle: Bad faith possessors must prove generating cause even without just title}

\textbf{Intermediate Representation:}

\begin{quote}
bad faith exercise is obligated and adverse possession claim is claimed
\end{quote}

\textbf{Witness Clause:}

\begin{lstlisting}[style=witnessstyle, language=Witness]
clause tesis2021246_principio_fundamental = oblig(ejercicio_mala_fe) AND claim(demanda_prescripcion) IMPLIES oblig(causa_generadora_probada);
\end{lstlisting}

\subsubsection{Clause: tesis2021246\_causa\_generadora}

\textit{Supporting legal logic: Bad faith possession requires demonstrating ownership character}

\textbf{Intermediate Representation:}

\begin{quote}
proven generating cause is obligated
\end{quote}

\textbf{Witness Clause:}

\begin{lstlisting}[style=witnessstyle, language=Witness]
clause tesis2021246_causa_generadora = oblig(causa_generadora_probada) IMPLIES oblig(posesion_como_propietario);
\end{lstlisting}

\subsubsection{Clause: tesis2021246\_plazo\_mala\_fe}

\textit{Temporal requirement: Bad faith possession requires extended time period (10 years)}

\textbf{Intermediate Representation:}

\begin{quote}
bad faith exercise is obligated and adverse possession claim is claimed
\end{quote}

\textbf{Witness Clause:}

\begin{lstlisting}[style=witnessstyle, language=Witness]
clause tesis2021246_plazo_mala_fe = oblig(ejercicio_mala_fe) AND claim(demanda_prescripcion) IMPLIES oblig(posesion_temporal);
\end{lstlisting}

\subsubsection{Clause: tesis2021246\_estandar\_probatorio}

\textit{Elevated proof standard: Bad faith requires higher burden of proof than good faith}

\textbf{Intermediate Representation:}

\begin{quote}
bad faith exercise is obligated
\end{quote}

\textbf{Witness Clause:}

\begin{lstlisting}[style=witnessstyle, language=Witness]
clause tesis2021246_estandar_probatorio = oblig(ejercicio_mala_fe) IMPLIES oblig(causa_generadora_probada);
\end{lstlisting}

\newpage

\section{Thesis 2021398: Usucapion - Limits of Allanamiento in Proving Possession Qualities}
\label{thesis:2021398}

\subsection{Original Thesis Text}

\begin{quote}
                                                   Semanario Judicial de la Federación

                                               Tesis

 Registro digital: 2021398

 Instancia: Tribunales             Décima Época                         Materia(s): Civil
 Colegiados de Circuito

 Tesis: II.2o.C.26 C (10a.)        Fuente: Semanario Judicial de la     Tipo: Aislada
                                   Federación.


USUCAPIÓN. LA RATIFICACIÓN DEL ALLANAMIENTO A LAS PRESTACIONES DEL ACTOR NO
TIENE COMO ALCANCE UN RECONOCIMIENTO PROPIO ACERCA DE QUE LA POSESIÓN SE
HAYA DESARROLLADO EN FORMA PACÍFICA, CONTINUA Y PÚBLICA –INTERRUPCIÓN DEL
CRITERIO SOSTENIDO EN LA TESIS AISLADA II.2o.C.258 C– (LEGISLACIÓN DEL ESTADO DE
MÉXICO).


La interpretación sistemática de los artículos 5.59, 5.60, 5.61, 5.127, 5.128, 5.129 y 5.130 del
Código Civil del Estado de México, permite establecer que la usucapión es un medio de adquirir la
propiedad de un bien por el transcurso del tiempo, con las condiciones previstas por el legislador, y
cuando la acción se promueve de buena fe respecto del inmueble, el actor debe acreditar que
cuenta con justo título para poseerlo en concepto de propietario, en forma pacífica, continua y
pública, por más de cinco años. Ahora bien, el allanamiento constituye la actitud autocompositiva
que implica el sometimiento incondicional por quien resiste en el proceso a las pretensiones del
accionante, respecto a derechos renunciables, lo cual lleva implícita la admisión de la exactitud de
los hechos, que puede hacerse en cualquier estado del juicio. Por ello, la ratificación del
allanamiento por admisión de los hechos expresados en la demanda resultaría apto para probar, en
su caso, la transmisión de la posesión en concepto de propietario y de buena fe, al tenor del justo
título requerido, cuando se reúna en el demandado la calidad de titular registral del inmueble objeto
del juicio, pues dicha parte está en posibilidad de renunciar a la transmisión de la propiedad aducida
como causa generadora de la posesión, así como a la inscripción de propiedad que obrare en su
favor en la institución registral, porque la buena fe se presume siempre en favor de quien la invoca;
sin embargo, la ratificación del allanamiento a las pretensiones del actor no puede tener como
alcance un reconocimiento propio acerca de que esa posesión se haya desarrollado en forma
pacífica, continua y pública, pues aun cuando quedara acreditado que el actor adquirió sin violencia
la posesión, por haber sido directamente otorgada por quien aparece como propietario en la
inscripción registral, la continuidad en esa posesión no dependería del sometimiento del
demandado, pues confluye en un hecho que pudiera incidir en derechos de terceros o con
afectación al interés público, condición exigida para poder realizar esa renuncia en el precepto 1.3
de dicha ley, pues el numeral 5.139 del propio código establece diversas hipótesis de interrupción
en el plazo para usucapir. Asimismo, el actor debe probar que ha poseído ininterrumpidamente ese
bien por el tiempo exigido, por lo cual, resolver con base en ese mero allanamiento, sería tanto
como considerar la imposibilidad de terceros para hacer valer acciones en defensa de derechos de
propiedad o posesión; además, si la posesión pública es la que se disfruta de manera que pueda
ser conocida por todos, el allanamiento a la demanda no puede representar renuncia de algún
derecho del que pudiera disponer el demandado, pues atiende al conocimiento de más de una
persona sobre el disfrute del bien respectivo ante la colectividad, incluso, el plazo de más de cinco
años previsto para prescribir, tampoco podría justificarse con dicho allanamiento, debido a las
causas que pueden dar lugar a la interrupción de la posesión, las cuales resultan ajenas al


                                              Pág. 1 de 2             Fecha de impresión 01/10/2025
  https://sjf2.scjn.gob.mx/detalle/tesis/2021398
                                                  Semanario Judicial de la Federación

demandado; consecuentemente, en una nueva reflexión, este Tribunal Colegiado de Circuito
interrumpe el criterio sostenido en la tesis II.2o.C.258 C, de rubro: "USUCAPIÓN. EL
ALLANAMIENTO A LA DEMANDA ES APTO PARA DEMOSTRAR LA POSESIÓN POR PARTE
DEL ACTOR Y SUS DEMÁS ATRIBUTOS (LEGISLACIÓN DEL ESTADO DE MÉXICO).", publicada
en el Semanario Judicial de la Federación y su Gaceta, Novena Época, Tomo XIII, enero de 2001,
página 1810, registro digital: 190437.

SEGUNDO TRIBUNAL COLEGIADO EN MATERIA CIVIL DEL SEGUNDO CIRCUITO.


Amparo directo 552/2015. Ma. del Rocío Loza Díaz y otro. 19 de octubre de 2015. Unanimidad de
votos. Ponente: José Antonio Rodríguez Rodríguez. Secretaria: Claudia Valeria Dávila Montero.

Nota: Esta tesis interrumpe el criterio sostenido por el propio tribunal en la diversa II.2o.C.258 C,
publicada en el Semanario Judicial de la Federación y su Gaceta, Novena Época, Tomo XIII, enero
de 2001, página 1810, con número de registro digital: 190437.


Esta tesis se publicó el viernes 10 de enero de 2020 a las 10:11 horas en el Semanario Judicial de
la Federación.




                                             Pág. 2 de 2            Fecha de impresión 01/10/2025
 https://sjf2.scjn.gob.mx/detalle/tesis/2021398

\end{quote}

\subsection{Formalization}

\subsubsection{Clause: tesis2021398\_principio\_fundamental}

\textit{Additional actions for this thesis Additional assets for allanamiento limitations Faith exclusivity: Possession is either good faith XOR bad faith (mutually exclusive) Thesis-specific principle: Allanamiento can prove good faith and transmission, but NOT peaceful/continuous/public possession}

\textbf{Intermediate Representation:}

\begin{quote}
If defendant acquiescence is obligated, then (good faith exercise is obligated and proven generating cause is obligated) and not (peaceful, continuous, and public possession is obligated)
\end{quote}

\textbf{Witness Clause:}

\begin{lstlisting}[style=witnessstyle, language=Witness]
clause tesis2021398_principio_fundamental = oblig(allanamiento_demandado) IMPLIES (oblig(ejercicio_buena_fe) AND oblig(causa_generadora_probada)) AND not(oblig(posesion_pacifico_continuo_publico));
\end{lstlisting}

\subsubsection{Clause: tesis2021398\_prueba\_independiente}

\textit{Limitation rule: Peaceful, continuous, public possession must be proven independently}

\textbf{Intermediate Representation:}

\begin{quote}
adverse possession claim is claimed
\end{quote}

\textbf{Witness Clause:}

\begin{lstlisting}[style=witnessstyle, language=Witness]
clause tesis2021398_prueba_independiente = claim(demanda_prescripcion) IMPLIES oblig(prueba_independiente);
\end{lstlisting}

\subsubsection{Clause: tesis2021398\_proteccion\_terceros}

\textit{Protection rule: Qualities affecting third parties cannot be established by allanamiento alone}

\textbf{Intermediate Representation:}

\begin{quote}
third party protection is obligated
\end{quote}

\textbf{Witness Clause:}

\begin{lstlisting}[style=witnessstyle, language=Witness]
clause tesis2021398_proteccion_terceros = oblig(proteccion_terceros) IMPLIES not(oblig(allanamiento_demandado));
\end{lstlisting}

\subsubsection{Clause: tesis2021398\_requisito\_temporal}

\textit{Temporal requirement: Good faith usucapion requires 5+ years with independent proof}

\textbf{Intermediate Representation:}

\begin{quote}
good faith exercise is obligated and adverse possession claim is claimed
\end{quote}

\textbf{Witness Clause:}

\begin{lstlisting}[style=witnessstyle, language=Witness]
clause tesis2021398_requisito_temporal = oblig(ejercicio_buena_fe) AND claim(demanda_prescripcion) IMPLIES oblig(posesion_temporal);
\end{lstlisting}

\newpage

\section{Thesis 2021537: Bad Faith Adverse Possession - Defective Documents Are Less Significant}
\label{thesis:2021537}

\subsection{Original Thesis Text}

\begin{quote}
                                                   Semanario Judicial de la Federación

                                                Tesis

 Registro digital: 2021537

 Instancia: Tribunales             Décima Época                           Materia(s): Civil
 Colegiados de Circuito

 Tesis: VII.2o.C.222 C (10a.)      Fuente: Semanario Judicial de la       Tipo: Aislada
                                   Federación.


PRESCRIPCIÓN POSITIVA DE MALA FE. LOS VICIOS CONTENIDOS EN EL DOCUMENTO QUE
ACREDITA LA CAUSA GENERADORA DE LA POSESIÓN SON MENOS TRASCENDENTES
PARA LA PROCEDENCIA DE DICHA ACCIÓN (LEGISLACIÓN DEL ESTADO DE VERACRUZ).


El concepto de propietario constituye el eje rector bajo el cual se analiza la procedencia o
improcedencia de la acción adquisitiva por prescripción, pues únicamente la posesión originaria es
apta para usucapir; por tanto, para la actualización de ese supuesto es indispensable que el
poseedor del bien cuente con el derecho de disposición (ius abutendi), el derecho de apropiarse de
los frutos del bien (ius fruendi) y el derecho de usar el bien (ius utendi); es decir, que el propietario
de la cosa se conduzca como el dueño y cuente con todos los derechos inherentes a ella. De ahí
que si la posesión ejercida por una persona es originaria, a título de dueño y por el tiempo
suficiente, de forma continua, pública y pacífica, es procedente la prescripción adquisitiva de mala fe
del bien a su favor. En ese sentido, sólo debe acreditar la causa generadora de la posesión y que
ésta sea apta para prescribir, por lo que necesariamente debe aportar una prueba objetiva con la
cual revele el origen de ésta. Sin embargo, en caso de que dicha prueba consista en un contrato de
compraventa, los vicios contenidos en dicho documento son menos trascendentes si se ejerce la
prescripción adquisitiva por mala fe, esto, ya que al exigir la ley un lapso mayor para su procedencia
(veinte años en términos del artículo 1185, fracción III, del Código Civil para el Estado de Veracruz),
implica un desinterés prolongado y continuo del legítimo propietario para recuperar el bien. Por lo
cual, la diferencia entre este tipo de prescripción adquisitiva y la ejercida de buena fe radica en el
grado de perfección del justo título, toda vez que, en el segundo supuesto, al establecerse un lapso
menor para su procedencia forzosamente requiere que la causa generadora sea lo más perfecta
posible. En este contexto, si una persona suscribió un contrato de compraventa de un inmueble, se
estima que dicha documental resulta apta para demostrar la posesión originaria del bien, porque
implica su adquisición tanto jurídica como material, es decir, que el inmueble es adquirido de forma
íntegra y completa con todos sus derechos y obligaciones; de ahí que si bien el carácter del
vendedor del inmueble puede afectar la validez del contrato, eso no demerita la calidad de la
posesión del comprador, ya que para ello debe anularse el título de mérito por la ausencia de las
facultades del vendedor para transmitir el dominio del bien, vicio que acarrea su nulidad. No
obstante, si eso no es demostrado en el juicio y el legítimo dueño del inmueble no hizo las gestiones
necesarias para su recuperación en el tiempo que exige la ley para usucapir, entonces, procede la
prescripción adquisitiva del bien, aun cuando la causa generadora de la posesión presente vicios,
como lo es la falta de facultades del vendedor para transmitir la propiedad del inmueble.

SEGUNDO TRIBUNAL COLEGIADO EN MATERIA CIVIL DEL SÉPTIMO CIRCUITO.


Amparo directo 335/2019. Rufina Sanchez Badillo. 22 de noviembre de 2019. Unanimidad de votos.

                                              Pág. 1 de 2              Fecha de impresión 01/10/2025
  https://sjf2.scjn.gob.mx/detalle/tesis/2021537
                                                   Semanario Judicial de la Federación

Ponente: Lucio Huesca Ballesteros, secretario de tribunal autorizado por la Comisión de Carrera
Judicial del Consejo de la Judicatura Federal para desempeñar las funciones de Magistrado, en
términos de los artículos 26, párrafo segundo y 81, fracción XXII, de la Ley Orgánica del Poder
Judicial de la Federación, en relación con el diverso 40, fracción V, del Acuerdo General del Pleno
del Consejo de la Judicatura Federal, que reglamenta la organización y funcionamiento del propio
Consejo; y reforma y deroga disposiciones de otros acuerdos generales. Secretaria: Rubí Sindirely
Aguilar Lasserre.

Nota: Esta tesis fue objeto de la denuncia relativa a la contradicción de criterios 378/2023 del índice
de la Suprema Corte de Justicia de la Nación, desechada por notoriamente improcedente, mediante
acuerdo de presidencia de 4 de diciembre de 2023.


Esta tesis se publicó el viernes 31 de enero de 2020 a las 10:32 horas en el Semanario Judicial de
la Federación.




                                              Pág. 2 de 2             Fecha de impresión 01/10/2025
  https://sjf2.scjn.gob.mx/detalle/tesis/2021537

\end{quote}

\subsection{Formalization}

\subsubsection{Clause: tesis2021537\_principio\_fundamental}

\textit{Additional actions for this thesis Additional assets for defective document tolerance Faith exclusivity: Adverse possession is either good faith XOR bad faith Thesis-specific principle: Bad faith adverse possession tolerates defective documents due to extended time period}

\textbf{Intermediate Representation:}

\begin{quote}
bad faith exercise is obligated and documento con vicios is obligated and temporal possession is obligated
\end{quote}

\textbf{Witness Clause:}

\begin{lstlisting}[style=witnessstyle, language=Witness]
clause tesis2021537_principio_fundamental = oblig(ejercicio_mala_fe) AND oblig(documento_con_vicios) AND oblig(posesion_temporal) IMPLIES claim(demanda_prescripcion);
\end{lstlisting}

\subsubsection{Clause: tesis2021537\_compensacion\_temporal}

\textit{Temporal compensation: Longer time period compensates for document defects}

\textbf{Intermediate Representation:}

\begin{quote}
documento con vicios is obligated and bad faith exercise is obligated
\end{quote}

\textbf{Witness Clause:}

\begin{lstlisting}[style=witnessstyle, language=Witness]
clause tesis2021537_compensacion_temporal = oblig(documento_con_vicios) AND oblig(ejercicio_mala_fe) IMPLIES oblig(posesion_temporal);
\end{lstlisting}

\subsubsection{Clause: tesis2021537\_desinteres\_propietario}

\textit{Disinterest demonstration: Extended possession shows legitimate owner's disinterest}

\textbf{Intermediate Representation:}

\begin{quote}
temporal possession is obligated
\end{quote}

\textbf{Witness Clause:}

\begin{lstlisting}[style=witnessstyle, language=Witness]
clause tesis2021537_desinteres_propietario = oblig(posesion_temporal) IMPLIES oblig(desinteres_propietario);
\end{lstlisting}

\subsubsection{Clause: tesis2021537\_ejercicio\_dominio}

\textit{Ownership exercise: Must demonstrate full ownership rights regardless of document defects}

\textbf{Intermediate Representation:}

\begin{quote}
adverse possession claim is claimed
\end{quote}

\textbf{Witness Clause:}

\begin{lstlisting}[style=witnessstyle, language=Witness]
clause tesis2021537_ejercicio_dominio = claim(demanda_prescripcion) IMPLIES oblig(posesion_como_propietario);
\end{lstlisting}

\subsubsection{Clause: tesis2021537\_contraste\_buena\_fe}

\textit{Contrast with good faith: Good faith requires more perfect documentation}

\textbf{Intermediate Representation:}

\begin{quote}
good faith exercise is obligated
\end{quote}

\textbf{Witness Clause:}

\begin{lstlisting}[style=witnessstyle, language=Witness]
clause tesis2021537_contraste_buena_fe = oblig(ejercicio_buena_fe) IMPLIES not(oblig(documento_con_vicios));
\end{lstlisting}

\newpage

\section{Thesis 2022776: DERECHO DEL TANTO EN MATERIA AGRARIA. LOS HIJOS DEL EJIDATARIO QUE CEDIÓ LOS DERECHOS PARCELARIOS, AL SER TITULARES DE ESA PRERROGATIVA, TIENEN INTERÉS}
\label{thesis:2022776}

\subsection{Original Thesis Text}

\begin{quote}
                                                   Semanario Judicial de la Federación

                                                Tesis

 Registro digital: 2022776

 Instancia: Tribunales             Décima Época                           Materia(s): Común
 Colegiados de Circuito

 Tesis: XXVIII.1o.3 A (10a.)       Fuente: Semanario Judicial de la       Tipo: Aislada
                                   Federación.


DERECHO DEL TANTO EN MATERIA AGRARIA. LOS HIJOS DEL EJIDATARIO QUE CEDIÓ LOS
DERECHOS PARCELARIOS, AL SER TITULARES DE ESA PRERROGATIVA, TIENEN INTERÉS
JURÍDICO EN EL JUICIO DE AMPARO PARA RECLAMAR LA OMISIÓN DE LLAMARLOS AL
JUICIO DONDE SE EJERCIÓ LA PRESCRIPCIÓN ADQUISITIVA DE BUENA FE DE AQUÉLLOS.


De los artículos 80, 83 y 84 de la Ley Agraria se advierte que tratándose de operaciones onerosas,
el legislador dejó al ejidatario en libertad para disponer de sus bienes, adoptando las formas que
considere más adecuadas, permitiéndole celebrar cualquier contrato que diversifique riesgos e
incremente sus ingresos, con la única limitante de que, en caso de enajenación de parcelas a
terceros no ejidatarios, debe concederse el derecho del tanto, entre otros, a sus hijos, so pena de
nulidad. Así, la traslación de dominio en materia agraria se encuentra limitada a que se respete ese
derecho, lo cual se cumple cuando se le comunica al beneficiario para que lo pueda ejercer en el
plazo de treinta días naturales a partir de su notificación o, en su caso, éste renuncie de manera
expresa o lo pierda por caducidad, con la finalidad de que los derechos parcelarios se conserven en
el núcleo familiar, a pesar de la existencia de una probable transmisión, lo cual le otorga
legitimación para salir en defensa de esa prerrogativa preferencial. Por tanto, los hijos del ejidatario
que cedió los derechos parcelarios tienen interés jurídico en el juicio de amparo para reclamar la
omisión de llamarlos al juicio agrario donde se ejerció la prescripción adquisitiva de buena fe de
éstos, al ser titulares de un derecho subjetivo; además, la tramitación del juicio y la sentencia que
declaró procedente la acción les origina una afectación personal y directa, pues trastocó su derecho
del tanto, ya que no se les otorgó el derecho de audiencia para ejercer su defensa, pues estuvieron
en imposibilidad jurídica y material para cuestionar el contrato de cesión que constituyó el
documento base de la acción.

PRIMER TRIBUNAL COLEGIADO DEL VIGÉSIMO OCTAVO CIRCUITO.


Amparo en revisión 141/2020. Rocío Espejel Lazcano. 19 de noviembre de 2020. Unanimidad de
votos. Ponente: Jesús Eduardo Hernández Fonseca. Secretario: Arturo Fonseca Mendoza.


Esta tesis se publicó el viernes 05 de marzo de 2021 a las 10:08 horas en el Semanario Judicial de
la Federación.




                                              Pág. 1 de 2             Fecha de impresión 01/10/2025
  https://sjf2.scjn.gob.mx/detalle/tesis/2022776
                                                 Semanario Judicial de la Federación




                                            Pág. 2 de 2       Fecha de impresión 01/10/2025
https://sjf2.scjn.gob.mx/detalle/tesis/2022776

\end{quote}

\subsection{Formalization}

\subsubsection{Clause: tesis2022776\_principio\_fundamental}

\textit{Thesis-specific principle: Possession must be "in concept of owner" (peaceful, continuous, public) for both good and bad faith}

\textbf{Intermediate Representation:}

\begin{quote}
possession as owner is obligated and (oblig(ejercicio\_buena\_fe) OR oblig(ejercicio\_mala\_fe)) IMPLIES oblig(ejercer\_prescripcion\_positiva)
\end{quote}

\textbf{Witness Clause:}

\begin{lstlisting}[style=witnessstyle, language=Witness]
clause tesis2022776_principio_fundamental = oblig(posesion_como_propietario) AND (oblig(ejercicio_buena_fe) OR oblig(ejercicio_mala_fe)) IMPLIES oblig(ejercer_prescripcion_positiva);
\end{lstlisting}

\subsubsection{Clause: tesis2022776\_posesion\_ambas\_fe}

\textit{Supporting logic: Possession as owner applies to both good faith and bad faith cases}

\textbf{Intermediate Representation:}

\begin{quote}
If possession as owner is obligated, then (good faith exercise is obligated or bad faith exercise is obligated)
\end{quote}

\textbf{Witness Clause:}

\begin{lstlisting}[style=witnessstyle, language=Witness]
clause tesis2022776_posesion_ambas_fe = oblig(posesion_como_propietario) IMPLIES (oblig(ejercicio_buena_fe) OR oblig(ejercicio_mala_fe));
\end{lstlisting}

\subsubsection{Clause: tesis2022776\_dominio\_economico}

\textit{Economic dominion: Must be dominator of the thing, commanding and enjoying it as owner}

\textbf{Intermediate Representation:}

\begin{quote}
possession as owner is obligated
\end{quote}

\textbf{Witness Clause:}

\begin{lstlisting}[style=witnessstyle, language=Witness]
clause tesis2022776_dominio_economico = oblig(posesion_como_propietario) IMPLIES oblig(causa_generadora_probada);
\end{lstlisting}

\subsubsection{Clause: tesis2022776\_posesion\_originaria}

\textit{Original possession: Must begin with cause different from derived possession}

\textbf{Intermediate Representation:}

\begin{quote}
possession as owner is obligated and not (temporal possession is obligated)
\end{quote}

\textbf{Witness Clause:}

\begin{lstlisting}[style=witnessstyle, language=Witness]
clause tesis2022776_posesion_originaria = oblig(posesion_como_propietario) AND not(oblig(posesion_temporal)) IMPLIES oblig(causa_generadora_probada);
\end{lstlisting}

\newpage

\section{Thesis 2023487: Good Faith Adverse Possession - Donation Contract for Different Property Cannot Constitute Just Title}
\label{thesis:2023487}

\subsection{Original Thesis Text}

\begin{quote}
                                                   Semanario Judicial de la Federación

                                                Tesis

 Registro digital: 2023487

 Instancia: Tribunales             Undécima Época                         Materia(s): Civil
 Colegiados de Circuito

 Tesis: XXII.1o.A.C.8 C (10a.)     Fuente: Semanario Judicial de la       Tipo: Aislada
                                   Federación.


PRESCRIPCIÓN ADQUISITIVA DE BUENA FE. EL JUSTO TÍTULO, COMO REQUISITO PARA SU
PROCEDENCIA, NO LO CONSTITUYE EL CONTRATO DE DONACIÓN CELEBRADO ENTRE
LAS PARTES RESPECTO DE UN INMUEBLE DIVERSO AL QUE SE PRETENDE PRESCRIBIR
(LEGISLACIÓN DEL ESTADO DE QUERÉTARO).


De los artículos 799, 802, 803, 804, 806, 822, 823, 1131, 1143 y 1144, fracciones I y III, del Código
Civil del Estado de Querétaro, cuyo texto, además, es similar al que fue materia de interpretación
por la Primera Sala de la Suprema Corte de Justicia de la Nación, al resolver la contradicción de
tesis 204/2014, de la que derivó la tesis de jurisprudencia 1a./J. 82/2014 (10a.), de título y subtítulo:
"PRESCRIPCIÓN ADQUISITIVA. AUNQUE LA LEGISLACIÓN APLICABLE NO EXIJA QUE EL
JUSTO TÍTULO O ACTO TRASLATIVO DE DOMINIO QUE CONSTITUYE LA CAUSA
GENERADORA DE LA POSESIÓN DE BUENA FE, SEA DE FECHA CIERTA, LA CERTEZA DE LA
FECHA DEL ACTO JURÍDICO DEBE PROBARSE EN FORMA FEHACIENTE POR SER UN
ELEMENTO DEL JUSTO TÍTULO (INTERRUPCIÓN DE LA JURISPRUDENCIA 1a./J. 9/2008).", se
advierte que la posesión legalmente exigida para la procedencia de la prescripción adquisitiva de
buena fe debe ser originaria, es decir, en concepto de dueño o propietario y por este último debe
entenderse al que detenta con "justo título", que equivale a un acto jurídico "imperfecto", per se
suficiente para poseer en concepto de propietario, pero no para transferirle el dominio y cuyo titular
cree, en forma seria, que es bastante para ello, a más de que descansa en un error que el juzgador
estima fundado, porque en cualquier persona pudo haber motivado creencia sobre su validez. Bajo
esas premisas, el contrato de donación celebrado entre las partes materiales sobre un inmueble
diverso al que se pretende usucapir de buena fe y que, con posterioridad, materialmente se ubicó
equivocadamente por un tercero, no constituye el justo título que exige la ley sustantiva civil del
Estado de Querétaro para la procedencia de dicha prescripción, ya que: 1) dicho acto jurídico debe
entenderse en los términos en que fue pactado por las partes, es decir, respecto del bien raíz que
ahí se precisó, de manera que no es apto para transmitir la posesión de un inmueble diverso al que
se pretende prescribir, ni podía subjetivamente hacer creer al usucapista que ello era factible, pues
el donante no externó su voluntad de transmitir ese dominio, lo que no se vincula a un requisito de
validez, sino que trasciende a la inexistencia del acto jurídico, que es insubsanable, en términos de
los artículos 1672, 1673, 1675, 1682 y 2110 del ordenamiento en cita; 2) los vicios de validez del
justo título deben emanar del propio acto jurídico imperfecto, no así de un error humano
independiente, ajeno y posterior a la celebración del propio acto traslativo, como lo es la incorrecta
ubicación material del predio afecto; 3) no considerarlo así, podría dar lugar a operaciones
fraudulentas en perjuicio de terceros, contra un principio de seguridad jurídica en las operaciones
realizadas, al estimar que un acto jurídico celebrado entre particulares sobre un determinado objeto
puede constituir justo título para demandar la usucapión de un predio diverso; y, 4) conforme a un
principio de equidad entre los contendientes, debe razonarse que: 4.1 es insoslayable que los
donantes actuaron voluntariamente y de buena fe, al transmitir gratuitamente el bien raíz; 4.2 el


                                              Pág. 1 de 2              Fecha de impresión 01/10/2025
  https://sjf2.scjn.gob.mx/detalle/tesis/2023487
                                                   Semanario Judicial de la Federación

yerro en la ubicación material del predio no es imputable a los donantes; 4.3 no puede aceptarse
que en perjuicio de estos últimos, se permitiera al donatario ostentar no sólo el predio que real y
efectivamente se donó, sino también el diverso cuya prescripción adquisitiva reclama; y, 4.4 de
manera incongruente existiría la posibilidad de que con base en un mismo acto jurídico de donación
se reconozcan dos derechos reales, uno, el de propiedad derivado del contrato de donación y dos,
el derivado de la prescripción adquisitiva a favor del donatario sobre diferente bien raíz. Por otro
lado, el contrato de donación sobre un predio diverso al que se pretende prescribir, tampoco es apto
para detonar el inicio del plazo legal de cinco años de posesión continua requerida para prescribir
de buena fe pues, se reitera, ese acto jurídico sólo transmite la posesión real, material y jurídica del
predio ahí pactado y no de uno diverso, al no ser objeto de la donación.

PRIMER TRIBUNAL COLEGIADO EN MATERIAS ADMINISTRATIVA Y CIVIL DEL VIGÉSIMO
SEGUNDO CIRCUITO.


Amparo directo 323/2019. 23 de enero de 2020. Unanimidad de votos. Ponente: Ramiro Rodríguez
Pérez. Secretario: Omar Alejandro Elizalde Herrera.

Nota: La parte conducente de la sentencia relativa a la contradicción de tesis 204/2014 y la tesis de
jurisprudencia 1a./J. 82/2014 (10a.) citadas, aparecen publicadas en el Semanario Judicial de la
Federación del viernes 5 de diciembre de 2014 a las 10:05 horas y en la Gaceta del Semanario
Judicial de la Federación, Décima Época, Libros 14, Tomo I, enero de 2015, página 688 y 13, Tomo
I, diciembre de 2014, página 200, con números de registro digital: 25449 y 2008083,
respectivamente.


Esta tesis se publicó el viernes 27 de agosto de 2021 a las 10:35 horas en el Semanario Judicial de
la Federación.




                                              Pág. 2 de 2             Fecha de impresión 01/10/2025
  https://sjf2.scjn.gob.mx/detalle/tesis/2023487

\end{quote}

\subsection{Formalization}

\subsubsection{Clause: tesis2023487\_principio\_fundamental}

\textit{Additional actions for this thesis Additional assets for property specificity requirement Faith exclusivity: Adverse possession is either good faith XOR bad faith Thesis-specific principle: Donation contract for different property cannot constitute just title for good faith adverse possession}

\textbf{Intermediate Representation:}

\begin{quote}
good faith exercise is obligated and contrato inmueble diverso is obligated
\end{quote}

\textbf{Witness Clause:}

\begin{lstlisting}[style=witnessstyle, language=Witness]
clause tesis2023487_principio_fundamental = oblig(ejercicio_buena_fe) AND oblig(contrato_inmueble_diverso) IMPLIES not(oblig(causa_generadora_probada));
\end{lstlisting}

\subsubsection{Clause: tesis2023487\_especificidad\_titulo}

\textit{Specificity requirement: Just title must relate to the exact property being claimed}

\textbf{Intermediate Representation:}

\begin{quote}
good faith exercise is obligated and adverse possession claim is claimed
\end{quote}

\textbf{Witness Clause:}

\begin{lstlisting}[style=witnessstyle, language=Witness]
clause tesis2023487_especificidad_titulo = oblig(ejercicio_buena_fe) AND claim(demanda_prescripcion) IMPLIES not(oblig(contrato_inmueble_diverso));
\end{lstlisting}

\subsubsection{Clause: tesis2023487\_error\_material}

\textit{Material error limitation: Physical mistakes don't cure legal defects}

\textbf{Intermediate Representation:}

\begin{quote}
error ubicacion material is obligated
\end{quote}

\textbf{Witness Clause:}

\begin{lstlisting}[style=witnessstyle, language=Witness]
clause tesis2023487_error_material = oblig(error_ubicacion_material) IMPLIES not(oblig(causa_generadora_probada));
\end{lstlisting}

\subsubsection{Clause: tesis2023487\_prevencion\_fraude}

\textit{Fraud prevention: Protects against using one contract for multiple properties}

\textbf{Intermediate Representation:}

\begin{quote}
contrato inmueble diverso is obligated
\end{quote}

\textbf{Witness Clause:}

\begin{lstlisting}[style=witnessstyle, language=Witness]
clause tesis2023487_prevencion_fraude = oblig(contrato_inmueble_diverso) IMPLIES oblig(proteccion_terceros);
\end{lstlisting}

\subsubsection{Clause: tesis2023487\_proteccion\_equidad}

\textit{Equity protection: Protects original parties from unintended consequences}

\textbf{Intermediate Representation:}

\begin{quote}
contrato inmueble diverso is obligated
\end{quote}

\textbf{Witness Clause:}

\begin{lstlisting}[style=witnessstyle, language=Witness]
clause tesis2023487_proteccion_equidad = oblig(contrato_inmueble_diverso) IMPLIES oblig(proteccion_registral);
\end{lstlisting}

\subsubsection{Clause: tesis2023487\_requisito\_buena\_fe}

\textit{Good faith requirement: Must have valid just title for good faith adverse possession}

\textbf{Intermediate Representation:}

\begin{quote}
good faith exercise is obligated and adverse possession claim is claimed
\end{quote}

\textbf{Witness Clause:}

\begin{lstlisting}[style=witnessstyle, language=Witness]
clause tesis2023487_requisito_buena_fe = oblig(ejercicio_buena_fe) AND claim(demanda_prescripcion) IMPLIES oblig(causa_generadora_probada);
\end{lstlisting}

\newpage

\section{Thesis 2024088: BAD FAITH ADVERSE POSSESSION WITHOUT TITLE}
\label{thesis:2024088}

\subsection{Original Thesis Text}

\begin{quote}
                                                   Semanario Judicial de la Federación

                                               Tesis

 Registro digital: 2024088

 Instancia: Primera Sala           Undécima Época                       Materia(s): Civil

 Tesis: 1a./J. 2/2022 (11a.)       Fuente: Semanario Judicial de la     Tipo: Jurisprudencia
                                   Federación.


PRESCRIPCIÓN POSITIVA DE MALA FE Y SIN TÍTULO. PARA QUE OPERE, DEBE
ACREDITARSE FEHACIENTEMENTE LA CAUSA GENERADORA DE LA POSESIÓN.


Hechos: En un juicio de amparo directo se reclamó una resolución dictada en apelación en la que se
consideró que operó la figura de la prescripción positiva sobre un bien inmueble, en beneficio de
diversa persona. El Tribunal Colegiado de Circuito concedió el amparo a la parte quejosa principal al
considerar que en el caso no se había acreditado la causa generadora de la posesión en concepto
de propietario. Inconforme con el fallo anterior, el tercero interesado y quejoso adherente interpuso
recurso de revisión.

Criterio jurídico: La Primera Sala de la Suprema Corte de Justicia de la Nación confirma la decisión
del Tribunal Colegiado, al estimar que el requisito de acreditar fehacientemente la causa generadora
de la posesión, en un juicio de prescripción positiva obtenida de mala fe y sin título, no resulta
desproporcional al grado de impedir el ejercicio del derecho a la propiedad, sino que constituye una
adecuada garantía procesal de seguridad jurídica.

Justificación: Para poder actualizar la prescripción positiva es necesario demostrar la causa
generadora de la posesión, sin importar si se trata de buena o mala fe, a fin de establecer si se
encuentra ante una posesión originaria o derivada. Esto, se erige como un requisito procesal
esencial de la acción, que brinda certeza a los titulares originales del bien sobre su derecho de
propiedad, sin que esto implique en forma alguna que el accionante de la usucapión se vea
impedido para instar la función jurisdiccional en aras de reivindicar el derecho que pretende. Por lo
que es menester cumplir con los requisitos que exige la ley, mismos que fueron establecidos por el
legislador en uso de su libertad configurativa para crear certeza en la propiedad y en la posesión,
materializando así la protección que, a través del debido proceso, otorga la Constitución General.
De ahí que, de ninguna manera la posesión derivada puede ser apta para prescribir ya que se trata
de una posesión ejercida sin ánimo de apropiación, es decir, enfocada exclusivamente en el uso y
disfrute temporal de un bien, sin que pueda colegirse, sin previa demostración, que se tuvo desde
un principio la posesión con ánimo de dueño. Por lo anterior, el requisito de acreditar
fehacientemente la causa generadora de la posesión no resulta desproporcionado ni hace nugatorio
el derecho al debido proceso en su vertiente de acceso a la justicia o los derechos a la posesión o a
la propiedad, por el contrario, constituye una adecuada garantía procesal de seguridad jurídica en
aras de la tutela de todos estos derechos.

PRIMERA SALA.


Amparo directo en revisión 2173/2020. María de Lourdes Gerarda Georgina González Silva. 16 de
junio de 2021. Cinco votos de las Ministras Norma Lucía Piña Hernández y Ana Margarita Ríos

                                              Pág. 1 de 2             Fecha de impresión 01/10/2025
  https://sjf2.scjn.gob.mx/detalle/tesis/2024088
                                                  Semanario Judicial de la Federación

Farjat, quien reservó su derecho para formular voto concurrente, y los Ministros Juan Luis González
Alcántara Carrancá, quien formuló voto aclaratorio, Jorge Mario Pardo Rebolledo y Alfredo Gutiérrez
Ortiz Mena, quien reservó su derecho para formular voto concurrente. Ponente: Ministro Jorge Mario
Pardo Rebolledo. Secretario: Alejandro Castañón Ramírez.

Tesis de jurisprudencia 2/2022 (11a.). Aprobada por la Primera Sala de este Alto Tribunal, en sesión
privada de doce de enero de dos mil veintidós.


Esta tesis se publicó el viernes 21 de enero de 2022 a las 10:22 horas en el Semanario Judicial de
la Federación y, por ende, se considera de aplicación obligatoria a partir del lunes 24 de enero de
2022, para los efectos previstos en el punto noveno del Acuerdo General Plenario 1/2021.




                                             Pág. 2 de 2            Fecha de impresión 01/10/2025
 https://sjf2.scjn.gob.mx/detalle/tesis/2024088

\end{quote}

\subsection{Formalization}

\subsubsection{Clause: tesis2024088\_principio\_fundamental}

\textit{Thesis-specific principle: Bad faith adverse possession without title requires proving the generating cause of possession}

\textbf{Intermediate Representation:}

\begin{quote}
bad faith exercise is obligated and not (registered ownership is obligated)
\end{quote}

\textbf{Witness Clause:}

\begin{lstlisting}[style=witnessstyle, language=Witness]
clause tesis2024088_principio_fundamental = oblig(ejercicio_mala_fe) AND not(oblig(titularidad_registral)) IMPLIES oblig(causa_generadora_probada);
\end{lstlisting}

\subsubsection{Clause: tesis2024088\_requisito\_causa\_generadora}

\textit{Supporting legal logic: Proving generating cause is essential for both good and bad faith possession}

\textbf{Intermediate Representation:}

\begin{quote}
adverse possession claim is obligated
\end{quote}

\textbf{Witness Clause:}

\begin{lstlisting}[style=witnessstyle, language=Witness]
clause tesis2024088_requisito_causa_generadora = oblig(demanda_prescripcion) IMPLIES oblig(causa_generadora_probada);
\end{lstlisting}

\subsubsection{Clause: tesis2024088\_limitacion\_posesion\_derivada}

\textit{Derived possession cannot prescribe because it lacks ownership intent}

\textbf{Intermediate Representation:}

\begin{quote}
temporal possession is obligated
\end{quote}

\textbf{Witness Clause:}

\begin{lstlisting}[style=witnessstyle, language=Witness]
clause tesis2024088_limitacion_posesion_derivada = oblig(posesion_temporal) IMPLIES not(oblig(posesion_como_propietario));
\end{lstlisting}

\subsubsection{Clause: tesis2024088\_seguridad\_juridica}

\textit{This requirement provides legal security without violating due process}

\textbf{Intermediate Representation:}

\begin{quote}
proven generating cause is obligated
\end{quote}

\textbf{Witness Clause:}

\begin{lstlisting}[style=witnessstyle, language=Witness]
clause tesis2024088_seguridad_juridica = oblig(causa_generadora_probada) IMPLIES oblig(proteccion_legal);
\end{lstlisting}

\newpage

\section{Thesis 2025237: Agrarian Adverse Possession - Third Party Requirement for Interruption}
\label{thesis:2025237}

\subsection{Original Thesis Text}

\begin{quote}
                                                   Semanario Judicial de la Federación

                                               Tesis

 Registro digital: 2025237

 Instancia: Tribunales             Undécima Época                        Materia(s): Administrativa
 Colegiados de Circuito

 Tesis: XVIII.2o.P.A.3 A (11a.)    Fuente: Semanario Judicial de la      Tipo: Aislada
                                   Federación.


PRESCRIPCIÓN ADQUISITIVA EN MATERIA AGRARIA. PARA QUE UNA DEMANDA
INTERRUMPA EL PLAZO RELATIVO CONFORME AL TERCER PÁRRAFO DEL ARTÍCULO 48
DE LA LEY AGRARIA, DEBE SER PRESENTADA POR UN TERCERO AJENO AL POSEEDOR.


Hechos: El quejoso promovió procedimiento de jurisdicción voluntaria ante el Tribunal Unitario
Agrario con la pretensión de que se declarara la prescripción adquisitiva en su favor de unas
parcelas, el cual se declaró improcedente al sostenerse que la posesión ejercida no se ajusta al
término previsto en el artículo 48 de la Ley Agraria, pues la promoción del juicio por cualquier
interesado interrumpe el plazo legal previsto en dicho precepto para tal efecto.

Criterio jurídico: Este Tribunal Colegiado de Circuito determina que si bien es cierto que el tercer
párrafo del artículo 48 de la Ley Agraria prevé que la demanda presentada por cualquier interesado
ante el tribunal agrario o la denuncia ante el Ministerio Público por despojo interrumpirá el plazo de
la prescripción adquisitiva establecido en su primer párrafo, hasta que se dicte resolución definitiva,
también lo es que ello no debe interpretarse en un sentido genérico, sino que quien promueva tenga
intereses contrarios a quien tiene la posesión y, en su caso, acude a ejercer dicha acción, esto es,
un tercero ajeno al poseedor y accionante y quien, derivado de dicha acción, incida o pretenda
interrumpir dicha posesión.

Justificación: Lo anterior, en virtud de que de dicho precepto deriva que aquella persona que
hubiere poseído tierras ejidales, en concepto de titular de derechos de ejidatario, que no sean las
destinadas al asentamiento humano ni se trate de bosques o selvas, de manera pacífica, continua y
pública durante un periodo de cinco años si la posesión es de buena fe, o de diez si fuera de mala
fe, adquirirá sobre dichas tierras los mismos derechos que cualquier ejidatario sobre su parcela y
que el poseedor podrá acudir ante el tribunal agrario para que, previa audiencia de los interesados,
del comisariado ejidal y de los colindantes, en la vía de jurisdicción voluntaria o mediante el
desahogo del juicio correspondiente, emita resolución sobre la adquisición de los derechos sobre la
parcela o tierras de que se trate, lo que comunicará al Registro Agrario Nacional para que expida de
inmediato el certificado correspondiente. Lo anterior conforme al criterio sustentado por la Segunda
Sala de la Suprema Corte de Justicia de la Nación, al resolver la contradicción de tesis 320/2016, de
la que derivó la tesis de jurisprudencia 2a./J. 86/2017 (10a.), de título y subtítulo: "JURISDICCIÓN
VOLUNTARIA SOBRE PRESCRIPCIÓN ADQUISITIVA EN MATERIA AGRARIA. LA RESOLUCIÓN
CON LA QUE CONCLUYE EL PROCEDIMIENTO RELATIVO CONSTITUYE UN ACTO PRIVATIVO
RESPECTO DEL CUAL ES NECESARIO RESPETAR EL DERECHO DE AUDIENCIA PREVIA.".
De esta manera, el hecho de que el aludido precepto legal refiera que interrumpirá el plazo la
demanda presentada por cualquier interesado, no debe entenderse en un sentido genérico, sino
que el que promueva tenga intereses contrarios al que tenga la posesión, esto es, debe ser un
tercero ajeno al poseedor, que derivado de dicha acción incida o pretenda interrumpir dicha


                                              Pág. 1 de 2             Fecha de impresión 01/10/2025
  https://sjf2.scjn.gob.mx/detalle/tesis/2025237
                                                  Semanario Judicial de la Federación

posesión.

SEGUNDO TRIBUNAL COLEGIADO EN MATERIAS PENAL Y ADMINISTRATIVA DEL DÉCIMO
OCTAVO CIRCUITO.


Amparo directo 263/2020. Vicencio Marín Castillo. 14 de mayo de 2021. Unanimidad de votos.
Ponente: Juan José Franco Luna. Secretario: Héctor Toledo Bárcenas.

Nota: La tesis de jurisprudencia 2a./J. 86/2017 (10a.) y la parte conducente de la sentencia relativa
a la contradicción de tesis 320/2016 citadas, aparecen publicadas en el Semanario Judicial de la
Federación del viernes 11 de agosto de 2017 a las 10:19 horas y en la Gaceta del Semanario
Judicial de la Federación, Décima Época, Libro 45, Tomo II, agosto de 2017, páginas 1005 y 957,
con números de registro digital: 2014866 y 27272, respectivamente.


Esta tesis se publicó el viernes 09 de septiembre de 2022 a las 10:18 horas en el Semanario
Judicial de la Federación.




                                             Pág. 2 de 2            Fecha de impresión 01/10/2025
 https://sjf2.scjn.gob.mx/detalle/tesis/2025237

\end{quote}

\subsection{Formalization}

\subsubsection{Clause: tesis2025237\_principio\_fundamental}

\textit{Additional subjects for agrarian context No additional actions needed - using universal assets Faith exclusivity: Adverse possession is either good faith XOR bad faith Thesis-specific principle: Only third parties with contrary interests can interrupt agrarian adverse possession}

\textbf{Intermediate Representation:}

\begin{quote}
If interrupcion prescripcion is obligated, then (third party protection is obligated and ejercicio derecho legal is obligated)
\end{quote}

\textbf{Witness Clause:}

\begin{lstlisting}[style=witnessstyle, language=Witness]
clause tesis2025237_principio_fundamental = oblig(interrupcion_prescripcion) IMPLIES (oblig(proteccion_terceros) AND oblig(ejercicio_derecho_legal));
\end{lstlisting}

\subsubsection{Clause: tesis2025237\_requisito\_tercero}

\textit{Third party requirement: Interruptor must be separate from possessor}

\textbf{Intermediate Representation:}

\begin{quote}
interrupcion prescripcion is claimed
\end{quote}

\textbf{Witness Clause:}

\begin{lstlisting}[style=witnessstyle, language=Witness]
clause tesis2025237_requisito_tercero = claim(interrupcion_prescripcion) IMPLIES oblig(proteccion_terceros);
\end{lstlisting}

\subsubsection{Clause: tesis2025237\_intereses\_contrarios}

\textit{Contrary interest requirement: Interruptor must have opposing interests}

\textbf{Intermediate Representation:}

\begin{quote}
interrupcion prescripcion is obligated
\end{quote}

\textbf{Witness Clause:}

\begin{lstlisting}[style=witnessstyle, language=Witness]
clause tesis2025237_intereses_contrarios = oblig(interrupcion_prescripcion) IMPLIES oblig(ejercicio_derecho_legal);
\end{lstlisting}

\subsubsection{Clause: tesis2025237\_contexto\_agrario}

\textit{Agrarian context: Applies specifically to ejidal land possession}

\textbf{Intermediate Representation:}

\begin{quote}
adverse possession claim is claimed
\end{quote}

\textbf{Witness Clause:}

\begin{lstlisting}[style=witnessstyle, language=Witness]
clause tesis2025237_contexto_agrario = claim(demanda_prescripcion) IMPLIES oblig(posesion_como_propietario);
\end{lstlisting}

\subsubsection{Clause: tesis2025237\_plazo\_agrario}

\textit{Temporal framework: Agrarian adverse possession follows same time periods}

\textbf{Intermediate Representation:}

\begin{quote}
possession as owner is obligated
\end{quote}

\textbf{Witness Clause:}

\begin{lstlisting}[style=witnessstyle, language=Witness]
clause tesis2025237_plazo_agrario = oblig(posesion_como_propietario) IMPLIES oblig(posesion_temporal);
\end{lstlisting}

\newpage

\section{Thesis 2025433: Good Faith Adverse Possession - Knowledge of INFONAVIT Mortgage Does Not Vitiate Just Title}
\label{thesis:2025433}

\subsection{Original Thesis Text}

\begin{quote}
                                                   Semanario Judicial de la Federación

                                                 Tesis

 Registro digital: 2025433

 Instancia: Tribunales              Undécima Época                         Materia(s): Civil
 Colegiados de Circuito

 Tesis: VII.2o.C.13 C (11a.)        Fuente: Semanario Judicial de la       Tipo: Aislada
                                    Federación.


PRESCRIPCIÓN ADQUISITIVA DE BUENA FE. CUANDO SE SUSTENTA EN EL HECHO DE
HABER CELEBRADO COMPRAVENTA VERBAL, EL CONOCIMIENTO DE QUE EL INMUEBLE
SE ENCUENTRA HIPOTECADO EN FAVOR DEL INSTITUTO DEL FONDO NACIONAL DE LA
VIVIENDA PARA LOS TRABAJADORES (INFONAVIT), EN TÉRMINOS DEL ARTÍCULO 49,
PÁRRAFO PRIMERO, DE LA LEY QUE LO RIGE, NO VICIA EL JUSTO TÍTULO PARA LA
PROCEDENCIA DE LA ACCIÓN (LEGISLACIÓN DEL ESTADO DE VERACRUZ DE IGNACIO DE
LA LLAVE).


Hechos: Una persona demandó la prescripción adquisitiva de buena fe de un terreno y de la casa
en él construida bajo el dicho de haber celebrado un contrato de compraventa verbal, en el que se
pactó el pago parcial de la operación y el resto sería entregado en cuanto se expidiera el mandato
para realizar las gestiones necesarias ante el Instituto del Fondo Nacional de la Vivienda para los
Trabajadores; sin embargo, la acción se declaró improcedente al estimarse que carecía de justo
título, en tanto el comprador había tenido conocimiento de que el inmueble se encontraba
hipotecado ante dicha institución, con la que se había pactado la no enajenación en favor de
terceros.

Criterio jurídico: Este Tribunal Colegiado de Circuito determina que el conocimiento de que un
inmueble se encuentra hipotecado en favor del Instituto del Fondo Nacional de la Vivienda para los
Trabajadores (Infonavit), en términos del artículo 49, párrafo primero, de la ley que lo rige, no vicia el
justo título para la procedencia de la acción de prescripción adquisitiva de buena fe.

Justificación: Lo anterior, porque de conformidad con el artículo 2181 del Código Civil para el Estado
de Veracruz de Ignacio de la Llave, en virtud del contrato de compraventa se adquiere la propiedad
de una cosa o de un derecho; por su parte, el diverso artículo 867 establece "El que tiene la
propiedad de una cosa puede gozar y disponer de ella; pero tal aprovechamiento sólo es lícito en
cuanto se hace de acuerdo con los intereses de la sociedad.". De este modo, el derecho de
propiedad implica un poder jurídico directo sobre la cosa para aprovecharla totalmente; es el
derecho real que otorga la mayor potestad jurídica en relación con un bien, pues los otros derechos
reales sólo comprenden formas de aprovechamiento parcial. Así, en virtud de la propiedad se
adquiere el derecho de disposición (ius abutendi), el de apropiarse de los frutos del bien (ius fruendi)
y el de usar el bien (ius utendi). El ius abutendi confiere el dominio jurídico pleno de la cosa; en
virtud de éste el propietario tiene la potestad de decidir sobre todos los actos que afecten la
situación jurídica de la cosa, incluyendo la traslación del inmueble en propiedad a otras personas y
no se pierde por obligarse a no ejecutar ciertos actos frente a terceros, sino sólo en el caso en que
se transfiera a un tercero esa potestad. Por tanto, cuando una persona adquiere en compraventa la
propiedad de un inmueble sin limitación alguna, en virtud del ius abutendi, a él le corresponde
decidir sobre el estatus jurídico de la cosa, por lo que válidamente puede obligarse en la medida y


                                               Pág. 1 de 2              Fecha de impresión 01/10/2025
  https://sjf2.scjn.gob.mx/detalle/tesis/2025433
                                                  Semanario Judicial de la Federación

grado que conforme a sus intereses convengan; sin embargo, si conforme a la potestad de domino
que le corresponde, el propietario se obliga a un no hacer, como ocurre en el caso a no enajenar
(entre otras) el inmueble sin consentimiento del acreedor hipotecario bajo la sanción de tener por
vencido anticipadamente un crédito y rescindida la relación contractual, ello no elimina el ius
abutendi con el que cuenta el propietario; por ello, el que una persona pacte en un convenio o
contrato la enajenación del inmueble en favor de terceros, siguiendo las directrices y términos
prescritos por el artículo 49, párrafo primero, de la Ley del Instituto del Fondo Nacional de la
Vivienda para los Trabajadores, no desaparece o anula su potestad de dominio jurídico sobre la
cosa para enajenarla, de modo que pase del propietario a la institución acreditante, porque entender
con ese alcance a las cláusulas contractuales acordadas al amparo de dicha disposición normativa,
iría en contra de lo previsto en el artículo 2234 del Código Civil para el Estado de Veracruz de
Ignacio de la Llave que prohíbe convenir la no enajenación del inmueble a tercera persona y,
además, haría nudo propietario al Infonavit y usufructuario al comprador del inmueble, aspecto que
es contrario al artículo 3o., fracción II, inciso a), de la ley de dicho organismo, que dispone como
objetivo la operación de un sistema de financiamiento que permita a los trabajadores obtener
créditos para la adquisición de habitaciones. Por tanto, si la cláusula pactada conforme al artículo
49, párrafo primero, referido, no deslegitima al comprador del inmueble para disponer de él y, desde
luego, para enajenarlo; entonces, el que un ulterior adquiriente del inmueble cuente con el
conocimiento del contrato de crédito con pacto de no enajenar, no puede fungir como un vicio que
anule la buena fe del título imperfecto en la acción de prescripción adquisitiva.

SEGUNDO TRIBUNAL COLEGIADO EN MATERIA CIVIL DEL SÉPTIMO CIRCUITO.


Amparo directo 624/2021. 8 de septiembre de 2022. Unanimidad de votos. Ponente: Lucio Huesca
Ballesteros, secretario de tribunal autorizado por la Comisión de Carrera Judicial del Consejo de la
Judicatura Federal para desempeñar las funciones de Magistrado. Secretario: Alan Iván Torres
Hinojosa.


Esta tesis se publicó el viernes 04 de noviembre de 2022 a las 10:15 horas en el Semanario Judicial
de la Federación.




                                             Pág. 2 de 2            Fecha de impresión 01/10/2025
 https://sjf2.scjn.gob.mx/detalle/tesis/2025433

\end{quote}

\subsection{Formalization}

\subsubsection{Clause: tesis2025433\_principio\_fundamental}

\textit{Additional subjects for mortgage context Additional actions for this thesis Additional assets for mortgage knowledge and ownership rights Faith exclusivity: Adverse possession is either good faith XOR bad faith Thesis-specific principle: Knowledge of INFONAVIT mortgage does not vitiate just title for good faith adverse possession}

\textbf{Intermediate Representation:}

\begin{quote}
good faith exercise is obligated and conocimiento hipoteca infonavit is obligated
\end{quote}

\textbf{Witness Clause:}

\begin{lstlisting}[style=witnessstyle, language=Witness]
clause tesis2025433_principio_fundamental = oblig(ejercicio_buena_fe) AND oblig(conocimiento_hipoteca_infonavit) IMPLIES oblig(causa_generadora_probada);
\end{lstlisting}

\subsubsection{Clause: tesis2025433\_conservacion\_dominio}

\textit{Ius abutendi preservation: Contractual obligations don't eliminate ownership powers}

\textbf{Intermediate Representation:}

\begin{quote}
cumplir requisitos legales is obligated
\end{quote}

\textbf{Witness Clause:}

\begin{lstlisting}[style=witnessstyle, language=Witness]
clause tesis2025433_conservacion_dominio = oblig(cumplir_requisitos_legales) IMPLIES oblig(posesion_como_propietario);
\end{lstlisting}

\subsubsection{Clause: tesis2025433\_proteccion\_buena\_fe}

\textit{Good faith protection: Knowledge of restrictions doesn't destroy good faith}

\textbf{Intermediate Representation:}

\begin{quote}
conocimiento hipoteca infonavit is obligated
\end{quote}

\textbf{Witness Clause:}

\begin{lstlisting}[style=witnessstyle, language=Witness]
clause tesis2025433_proteccion_buena_fe = oblig(conocimiento_hipoteca_infonavit) IMPLIES not(not(oblig(ejercicio_buena_fe)));
\end{lstlisting}

\subsubsection{Clause: tesis2025433\_validez\_verbal}

\textit{Verbal contract validity: Verbal sale can constitute just title}

\textbf{Intermediate Representation:}

\begin{quote}
unregistered contract is obligated and good faith exercise is obligated
\end{quote}

\textbf{Witness Clause:}

\begin{lstlisting}[style=witnessstyle, language=Witness]
clause tesis2025433_validez_verbal = oblig(contrato_no_inscrito) AND oblig(ejercicio_buena_fe) IMPLIES oblig(causa_generadora_probada);
\end{lstlisting}

\subsubsection{Clause: tesis2025433\_distincion\_obligacion}

\textit{Ownership vs. obligation distinction: Obligations don't transfer ownership powers}

\textbf{Intermediate Representation:}

\begin{quote}
cumplir requisitos legales is obligated
\end{quote}

\textbf{Witness Clause:}

\begin{lstlisting}[style=witnessstyle, language=Witness]
clause tesis2025433_distincion_obligacion = oblig(cumplir_requisitos_legales) IMPLIES not(not(oblig(posesion_como_propietario)));
\end{lstlisting}

\newpage

\section{Thesis 2029251: Good Faith Adverse Possession Interruption - Constitutional Validity of Recognition-Based Interruption}
\label{thesis:2029251}

\subsection{Original Thesis Text}

\begin{quote}
                                                   Semanario Judicial de la Federación

                                                Tesis

 Registro digital: 2029251

 Instancia: Tribunales             Undécima Época                         Materia(s): Constitucional
 Colegiados de Circuito

 Tesis: XXVII.1o.1 C (11a.)        Fuente: Semanario Judicial de la       Tipo: Aislada
                                   Federación.


INTERRUPCIÓN DE LA PRESCRIPCIÓN ADQUISITIVA DE BUENA FE. EL ARTÍCULO 1846,
FRACCIÓN III, DEL CÓDIGO CIVIL PARA EL ESTADO DE QUINTANA ROO QUE LA PREVÉ, NO
VIOLA EL DERECHO HUMANO A LA PROPIEDAD PRIVADA.


Hechos: Una persona demandó en la vía civil la prescripción adquisitiva de buena fe de un
inmueble, con el argumento de haber recibido un ofrecimiento de venta y, posteriormente, celebrado
un convenio de promesa de compraventa; sin embargo, la acción se declaró improcedente, ya que
los documentos ofrecidos como justo título no fueron suscritos por quien al momento de su
celebración era propietario de los predios referidos; además, porque después de la fecha en que se
celebraron los documentos fundatorios de la acción, el actor reconoció como copropietarios del
citado inmueble a los demandados, circunstancia que interrumpió la prescripción adquisitiva de
buena fe en términos del artículo 1846, fracción III, del Código Civil para el Estado de Quintana Roo.
En el amparo directo el actor solicitó se realizara un control de convencionalidad ex officio sobre el
mencionado precepto legal.

Criterio jurídico: Este Tribunal Colegiado de Circuito determina que el artículo 1846, fracción III, del
Código Civil para el Estado de Quintana Roo, que establece la interrupción de la usucapión, no viola
el derecho humano a la propiedad privada.

Justificación: La interrupción de la prescripción adquisitiva de buena fe prevista en dicho artículo,
porque la persona a cuyo favor transcurre reconozca expresamente, de palabra, por escrito o
tácitamente por hechos indubitables, el derecho de la persona contra quien prescribe, no priva a
aquélla de adquirir el derecho de propiedad sobre un bien, ya sea mueble o inmueble, pues la
interrupción sólo tiene como efecto inutilizar todo el tiempo transcurrido antes de que se
materializara el reconocimiento expreso o los hechos indubitables que permitan acreditar el
reconocimiento tácito del derecho de la persona en contra de la cual se prescribe. La racionalidad
de la norma parte de la circunstancia de que la posesión de buena fe, necesaria para la
actualización de la usucapión, debe ser permanente por el plazo legal respectivo, y no únicamente
presentarse durante el momento de la adquisición, de modo que desde que se conozcan los vicios
del justo título se interrumpe el tiempo de posesión de buena fe y vuelve a iniciar el plazo para que
se actualice la usucapión. Por tanto, la interrupción en estudio no viola el artículo 21 de la
Convención Americana sobre Derechos Humanos, pues tiene como único objeto brindar seguridad
jurídica a los titulares de un justo título.

PRIMER TRIBUNAL COLEGIADO DEL VIGÉSIMO SÉPTIMO CIRCUITO.


Amparo directo 511/2022. Claudio Sergio de la Fuente Burton y/o Sergio de la Fuente Burton, su

                                              Pág. 1 de 2              Fecha de impresión 01/10/2025
  https://sjf2.scjn.gob.mx/detalle/tesis/2029251
                                                  Semanario Judicial de la Federación

sucesión. 11 de marzo de 2024. Unanimidad de votos. Ponente: José Antonio Belda Rodríguez.
Secretario: Carlos Hernández García.


Esta tesis se publicó el viernes 09 de agosto de 2024 a las 10:17 horas en el Semanario Judicial de
la Federación.




                                             Pág. 2 de 2           Fecha de impresión 01/10/2025
 https://sjf2.scjn.gob.mx/detalle/tesis/2029251

\end{quote}

\subsection{Formalization}

\subsubsection{Clause: tesis2029251\_principio\_fundamental}

\textit{Additional actions for this thesis Additional assets for recognition-based interruption Faith exclusivity: Adverse possession is either good faith XOR bad faith Thesis-specific principle: Recognition-based interruption is constitutionally valid and doesn't violate property rights}

\textbf{Intermediate Representation:}

\begin{quote}
If ejercicio derecho legal is obligated, then (interrupcion prescripcion is obligated and proteccion registral is obligated)
\end{quote}

\textbf{Witness Clause:}

\begin{lstlisting}[style=witnessstyle, language=Witness]
clause tesis2029251_principio_fundamental = oblig(ejercicio_derecho_legal) IMPLIES (oblig(interrupcion_prescripcion) AND oblig(proteccion_registral));
\end{lstlisting}

\subsubsection{Clause: tesis2029251\_efecto\_interrupcion}

\textit{Interruption effect: Recognition nullifies elapsed time and restarts prescription period}

\textbf{Intermediate Representation:}

\begin{quote}
ejercicio derecho legal is obligated
\end{quote}

\textbf{Witness Clause:}

\begin{lstlisting}[style=witnessstyle, language=Witness]
clause tesis2029251_efecto_interrupcion = oblig(ejercicio_derecho_legal) IMPLIES oblig(inutilizacion_tiempo);
\end{lstlisting}

\subsubsection{Clause: tesis2029251\_permanencia\_buena\_fe}

\textit{Good faith permanence: Good faith must be continuous throughout prescription period}

\textbf{Intermediate Representation:}

\begin{quote}
good faith exercise is obligated and ejercicio derecho legal is obligated
\end{quote}

\textbf{Witness Clause:}

\begin{lstlisting}[style=witnessstyle, language=Witness]
clause tesis2029251_permanencia_buena_fe = oblig(ejercicio_buena_fe) AND oblig(ejercicio_derecho_legal) IMPLIES oblig(ejercicio_mala_fe);
\end{lstlisting}

\subsubsection{Clause: tesis2029251\_validez\_constitucional}

\textit{Constitutional validity: Interruption mechanism provides legal security without violating rights}

\textbf{Intermediate Representation:}

\begin{quote}
interrupcion prescripcion is obligated
\end{quote}

\textbf{Witness Clause:}

\begin{lstlisting}[style=witnessstyle, language=Witness]
clause tesis2029251_validez_constitucional = oblig(interrupcion_prescripcion) IMPLIES oblig(proteccion_registral);
\end{lstlisting}

\subsubsection{Clause: tesis2029251\_proteccion\_titulares}

\textit{Legal security purpose: Rights exercise leads to protection}

\textbf{Intermediate Representation:}

\begin{quote}
ejercicio derecho legal is obligated
\end{quote}

\textbf{Witness Clause:}

\begin{lstlisting}[style=witnessstyle, language=Witness]
clause tesis2029251_proteccion_titulares = oblig(ejercicio_derecho_legal) IMPLIES oblig(proteccion_registral);
\end{lstlisting}

\newpage

\section{Thesis 2030499: Agrarian Adverse Possession - Supplementary Application of Civil Law with Agrarian Adaptations}
\label{thesis:2030499}

\subsection{Original Thesis Text}

\begin{quote}
                                                   Semanario Judicial de la Federación

                                                Tesis

 Registro digital: 2030499

 Instancia: Tribunales             Undécima Época                         Materia(s): Administrativa
 Colegiados de Circuito

 Tesis: II.2o.A.56 A (11a.)        Fuente: Semanario Judicial de la       Tipo: Aislada
                                   Federación.


PRESCRIPCIÓN ADQUISITIVA EN MATERIA AGRARIA. LA LEGISLACIÓN CIVIL PUEDE
APLICARSE SUPLETORIAMENTE CON LOS MATICES PROPIOS DE AQUELLA MATERIA.


Hechos: Una persona demandó la prescripción positiva de diversas parcelas que adquirió mediante
contratos de cesión de derechos. El Tribunal Unitario Agrario consideró que no acreditó que los
contratos se celebraron en fecha cierta para efecto del cómputo de la prescripción, en términos del
Código Federal de Procedimientos Civiles, de aplicación supletoria en materia agraria. Contra esa
determinación la actora promovió amparo directo, en el que argumentó que para que prospere su
pretensión no es necesario contar con un justo título, sino que sólo se debe acreditar la causa
generadora de la posesión.

Criterio jurídico: Este Tribunal Colegiado de Circuito determina que la legislación civil que regula la
prescripción adquisitiva puede aplicarse supletoriamente en los juicios agrarios en los que se ejerza
esa acción, con los matices propios de la materia agraria.

Justificación: En la jurisprudencia 2a./J. 207/2004, la Segunda Sala de la Suprema Corte de Justicia
de la Nación sostuvo que en los juicios agrarios no se requiere de un justo título para que prospere
la prescripción positiva, sino que sólo es indispensable acreditar la causa generadora de la posesión
para conocer la calidad con la que se ejerce. Esto es, si es originaria o derivada, si es de buena o
mala fe, y la fecha a partir de la cual se computará el término legal de la prescripción. En los juicios
agrarios en los que se ejerza la acción de prescripción y se señale como causa generadora de la
posesión un contrato de cesión de derechos, el órgano jurisdiccional debe valorar el documento que
contenga el contrato, junto con los restantes medios de convicción allegados al juicio agrario, con la
finalidad de determinar si son suficientes o no para crear plena convicción sobre una fecha a partir
de la cual puede comenzar a computarse el plazo de la prescripción. En ese contexto, no puede
exigirse que los contratos de cesión de derechos agrarios sean de fecha cierta con el rigor de la
materia civil, pues es una práctica común entre las personas ejidatarias, comuneras o aspirantes a
adquirir esas calidades, celebrarlos sin cumplir con todas las formalidades de la ley, pues
generalmente enajenan sus derechos conforme a sus prácticas o costumbres, e incluso llegan a
celebrar actos bajo el principio de buena fe y la confianza que implica "la palabra".

SEGUNDO TRIBUNAL COLEGIADO EN MATERIA ADMINISTRATIVA DEL SEGUNDO CIRCUITO.


Amparo directo 624/2023. 20 de junio de 2024. Unanimidad de votos. Ponente: Mónica Alejandra
Soto Bueno. Secretario: Edgar Iván Jiménez Sánchez.

Nota: La tesis de jurisprudencia 2a./J. 207/2004 citada, aparece publicada con el rubro:

                                              Pág. 1 de 2              Fecha de impresión 01/10/2025
  https://sjf2.scjn.gob.mx/detalle/tesis/2030499
                                                   Semanario Judicial de la Federación

"PRESCRIPCIÓN ADQUISITIVA EN MATERIA AGRARIA. PARA SU PROCEDENCIA NO SE
REQUIERE DE JUSTO TÍTULO.", en el Semanario Judicial de la Federación y su Gaceta, Novena
Época, Tomo XXI, enero de 2005, página 575, con número de registro digital: 179504.


Esta tesis se publicó el viernes 06 de junio de 2025 a las 10:12 horas en el Semanario Judicial de la
Federación.




                                              Pág. 2 de 2            Fecha de impresión 01/10/2025
  https://sjf2.scjn.gob.mx/detalle/tesis/2030499

\end{quote}

\subsection{Formalization}

\subsubsection{Clause: tesis2030499\_principio\_fundamental}

\textit{No additional actions needed - using universal assets Faith exclusivity: Adverse possession is either good faith XOR bad faith Thesis-specific principle: Civil law applies supplementarily in agrarian matters with appropriate adaptations}

\textbf{Intermediate Representation:}

\begin{quote}
If adverse possession claim is claimed, then (aplicacion supletoria is obligated and ejercicio derecho legal is obligated)
\end{quote}

\textbf{Witness Clause:}

\begin{lstlisting}[style=witnessstyle, language=Witness]
clause tesis2030499_principio_fundamental = claim(demanda_prescripcion) IMPLIES (oblig(aplicacion_supletoria) AND oblig(ejercicio_derecho_legal));
\end{lstlisting}

\subsubsection{Clause: tesis2030499\_sin\_justo\_titulo}

\textit{No just title requirement: Only generating cause needed in agrarian adverse possession}

\textbf{Intermediate Representation:}

\begin{quote}
aplicacion supletoria is obligated
\end{quote}

\textbf{Witness Clause:}

\begin{lstlisting}[style=witnessstyle, language=Witness]
clause tesis2030499_sin_justo_titulo = oblig(aplicacion_supletoria) IMPLIES oblig(causa_generadora_probada);
\end{lstlisting}

\subsubsection{Clause: tesis2030499\_adaptacion\_cultural}

\textit{Cultural adaptation: Formal requirements relaxed for agrarian customs}

\textbf{Intermediate Representation:}

\begin{quote}
ejercicio derecho legal is obligated
\end{quote}

\textbf{Witness Clause:}

\begin{lstlisting}[style=witnessstyle, language=Witness]
clause tesis2030499_adaptacion_cultural = oblig(ejercicio_derecho_legal) IMPLIES oblig(ejercicio_buena_fe);
\end{lstlisting}

\subsubsection{Clause: tesis2030499\_cesion\_valida}

\textit{Rights assignment validity: Agrarian rights transfers follow custom-based practices}

\textbf{Intermediate Representation:}

\begin{quote}
unregistered contract is obligated and good faith exercise is obligated
\end{quote}

\textbf{Witness Clause:}

\begin{lstlisting}[style=witnessstyle, language=Witness]
clause tesis2030499_cesion_valida = oblig(contrato_no_inscrito) AND oblig(ejercicio_buena_fe) IMPLIES oblig(causa_generadora_probada);
\end{lstlisting}

\subsubsection{Clause: tesis2030499\_alcance\_supletorio}

\textit{Supplementary application scope: Civil law fills gaps but respects agrarian specificity}

\textbf{Intermediate Representation:}

\begin{quote}
aplicacion supletoria is obligated
\end{quote}

\textbf{Witness Clause:}

\begin{lstlisting}[style=witnessstyle, language=Witness]
clause tesis2030499_alcance_supletorio = oblig(aplicacion_supletoria) IMPLIES oblig(posesion_como_propietario);
\end{lstlisting}

\newpage

\section{Thesis 2030500: Adverse Possession - Judicially Nullified Title Cannot Establish Generating Cause}
\label{thesis:2030500}

\subsection{Original Thesis Text}

\begin{quote}
                                                   Semanario Judicial de la Federación

                                                Tesis

 Registro digital: 2030500

 Instancia: Tribunales             Undécima Época                         Materia(s): Civil
 Colegiados de Circuito

 Tesis: I.11o.C.41 C (11a.)        Fuente: Semanario Judicial de la       Tipo: Aislada
                                   Federación.


PRESCRIPCIÓN POSITIVA. SI EL TÍTULO EN EL QUE SE SUSTENTA LA POSESIÓN DEL BIEN
INMUEBLE SE DECLARA JUDICIALMENTE NULO, NO ES APTO PARA DEMOSTRAR LA CAUSA
QUE DIO ORIGEN A LA POSESIÓN (LEGISLACIÓN APLICABLE PARA LA CIUDAD DE MÉXICO).


Hechos: En un juicio ordinario civil se ejerció la acción reivindicatoria y se reconvino la prescripción
positiva; se dictó sentencia definitiva firme que declaró fundada la acción principal y se absolvió
respecto de la reconvención. En segunda instancia se confirmó el fallo apelado. Inconforme, el
demandado en el juicio de origen promovió amparo directo en el que adujo que sí fundó su acción
de prescripción con el título de propiedad que exhibió.

Criterio jurídico: Este Tribunal Colegiado de Circuito determina que si la acción de prescripción
positiva se sustenta en un título declarado judicialmente nulo, éste no es apto para demostrar la
causa que dio origen a la posesión del bien inmueble.

Justificación: Si la nulidad se declara en virtud de que la parte compradora no compareció ante el
notario ni firmó la escritura, implica que no existió el acto traslativo de dominio en que se apoyó
quien accionó la prescripción positiva. Al ser inexistente el acto traslativo del inmueble respectivo,
significa que nunca produjo efecto alguno, en términos del artículo 2224 del Código Civil para el
Distrito Federal, aplicable para la Ciudad de México, sin que sean aplicables los diversos 2225 y
2226, en cuanto a que los actos puedan producir provisionalmente sus efectos, pues se refieren a la
ilicitud en el objeto, y no respecto a la inexistencia por falta de consentimiento. De ahí que al no
acreditarse la existencia del acto traslativo de dominio que dio origen a la posesión, no se genera la
convicción de que la parte accionante posee en concepto de dueña o propietaria, pues debe
tomarse en cuenta que la usucapión es una institución de orden público que pretende generar
certidumbre en la propiedad de los bienes materiales, por lo que es necesario cumplir la integridad
de los requisitos legales para poder actualizar esta figura, en este caso, los requisitos previstos en
los artículos 826, 1151, fracción I y 1152, fracción I, del código citado, los cuales exigen que la
posesión debe ejercerse en concepto de propietario. Además, esa nulidad genera que no se
acredite la posesión de buena fe por carecer precisamente de un título que sea traslativo de dominio
suficiente para otorgar el derecho a poseer.

DÉCIMO PRIMER TRIBUNAL COLEGIADO EN MATERIA CIVIL DEL PRIMER CIRCUITO.


Amparo directo 480/2021. José Guadalupe Naranjo Sánchez. 24 de noviembre de 2022.
Unanimidad de votos. Ponente: Fernando Rangel Ramírez. Secretario: Octavio Rosales Rivera.




                                              Pág. 1 de 2              Fecha de impresión 01/10/2025
  https://sjf2.scjn.gob.mx/detalle/tesis/2030500
                                                   Semanario Judicial de la Federación


Esta tesis se publicó el viernes 06 de junio de 2025 a las 10:12 horas en el Semanario Judicial de la
Federación.




                                              Pág. 2 de 2            Fecha de impresión 01/10/2025
  https://sjf2.scjn.gob.mx/detalle/tesis/2030500

\end{quote}

\subsection{Formalization}

\subsubsection{Clause: tesis2030500\_principio\_fundamental}

\textit{Additional assets for judicial nullification effects Faith exclusivity: Adverse possession is either good faith XOR bad faith Thesis-specific principle: Judicially nullified title cannot establish generating cause for adverse possession}

\textbf{Intermediate Representation:}

\begin{quote}
nulidad judicial titulo is obligated
\end{quote}

\textbf{Witness Clause:}

\begin{lstlisting}[style=witnessstyle, language=Witness]
clause tesis2030500_principio_fundamental = oblig(nulidad_judicial_titulo) IMPLIES not(oblig(causa_generadora_probada));
\end{lstlisting}

\subsubsection{Clause: tesis2030500\_efecto\_nulidad}

\textit{Nullification effect: Judicially nullified acts lack legal effect}

\textbf{Intermediate Representation:}

\begin{quote}
nulidad judicial titulo is obligated
\end{quote}

\textbf{Witness Clause:}

\begin{lstlisting}[style=witnessstyle, language=Witness]
clause tesis2030500_efecto_nulidad = oblig(nulidad_judicial_titulo) IMPLIES oblig(ausencia_efecto_juridico);
\end{lstlisting}

\subsubsection{Clause: tesis2030500\_acto\_inexistente}

\textit{Inexistent act consequence: Nullified acts are treated as legally inexistent}

\textbf{Intermediate Representation:}

\begin{quote}
ausencia efecto juridico is obligated
\end{quote}

\textbf{Witness Clause:}

\begin{lstlisting}[style=witnessstyle, language=Witness]
clause tesis2030500_acto_inexistente = oblig(ausencia_efecto_juridico) IMPLIES oblig(acto_inexistente);
\end{lstlisting}

\subsubsection{Clause: tesis2030500\_imposibilidad\_propietario}

\textit{Ownership character impossibility: Nullified title cannot establish ownership possession}

\textbf{Intermediate Representation:}

\begin{quote}
acto inexistente is obligated
\end{quote}

\textbf{Witness Clause:}

\begin{lstlisting}[style=witnessstyle, language=Witness]
clause tesis2030500_imposibilidad_propietario = oblig(acto_inexistente) IMPLIES not(oblig(posesion_como_propietario));
\end{lstlisting}

\subsubsection{Clause: tesis2030500\_imposibilidad\_buena\_fe}

\textit{Good faith impossibility: Nullified title prevents good faith establishment}

\textbf{Intermediate Representation:}

\begin{quote}
nulidad judicial titulo is obligated and adverse possession claim is claimed
\end{quote}

\textbf{Witness Clause:}

\begin{lstlisting}[style=witnessstyle, language=Witness]
clause tesis2030500_imposibilidad_buena_fe = oblig(nulidad_judicial_titulo) AND claim(demanda_prescripcion) IMPLIES oblig(ejercicio_mala_fe);
\end{lstlisting}

\subsubsection{Clause: tesis2030500\_orden\_publico}

\textit{Public order purpose: Adverse possession serves property certainty}

\textbf{Intermediate Representation:}

\begin{quote}
adverse possession claim is claimed
\end{quote}

\textbf{Witness Clause:}

\begin{lstlisting}[style=witnessstyle, language=Witness]
clause tesis2030500_orden_publico = claim(demanda_prescripcion) IMPLIES oblig(proteccion_registral);
\end{lstlisting}

\newpage

\section{Thesis 2030717: Adverse Possession - Improcedence Against Municipal Public Domain Properties}
\label{thesis:2030717}

\subsection{Original Thesis Text}

\begin{quote}
                                                  Semanario Judicial de la Federación

                                              Tesis

 Registro digital: 2030717

 Instancia: Tribunales            Undécima Época                       Materia(s): Civil
 Colegiados de Circuito

 Tesis: XX.1o.P.C.3 C (11a.)      Fuente: Semanario Judicial de la     Tipo: Aislada
                                  Federación.


PRESCRIPCIÓN POSITIVA. ES IMPROCEDENTE RESPECTO DE BIENES INMUEBLES DE
DOMINIO PÚBLICO PROPIEDAD DE LOS AYUNTAMIENTOS MUNICIPALES DESTINADOS A LA
RESERVA ECOLÓGICA, AL ESTAR FUERA DEL COMERCIO (LEGISLACIÓN DEL ESTADO DE
CHIAPAS).


Hechos: Una persona demandó al Ayuntamiento de Tuxtla Gutiérrez, Chiapas, en la vía ordinaria
civil, la prescripción positiva de buena fe de un bien inmueble, con el argumento de que contaba con
la posesión material y jurídica de éste, porque adquirió la propiedad mediante un contrato de
compraventa. En primera instancia se absolvió a la parte demandada, ya que el propietario del bien
es un Ayuntamiento municipal, por lo que resolvió que el inmueble en litigio es imprescriptible.
Determinación que fue confirmada en apelación.

Criterio jurídico: Este Tribunal Colegiado de Circuito determina que los bienes inmuebles propiedad
de los Ayuntamientos municipales, que tienen el carácter de dominio público y están fuera del
comercio, al ser destinados a la reserva ecológica, no pueden ser objeto de prescripción adquisitiva,
dada la prohibición expresa prevista en el Código Civil para el Estado de Chiapas y en la Ley de
Desarrollo Constitucional en Materia de Gobierno y Administración Municipal de la misma entidad
federativa.

Justificación: De la interpretación sistemática de los artículos 736 a 738, 753 a 759, 1123 a 1125,
1136 y 1138 del citado código; 127 y 128 de la mencionada ley, se advierte un hecho relevante
consistente en que los Ayuntamientos se considerarán como particulares "para la prescripción de
sus bienes" que "sean susceptibles de propiedad privada", entonces el Ayuntamiento debe
considerarse como particular para efectos de la prescripción de bienes; sin embargo, ello sólo es
posible para aquellos que son susceptibles de propiedad privada, es decir, los que no son de
dominio público y no estén excluidos del comercio; empero, cuando se trata de inmuebles que sí
son de dominio público que están destinados a la conservación ecológica y están fuera del
comercio, éstos no pueden ser prescriptibles en términos de las legislaciones civil y administrativa
citadas; de manera que la posesión de un inmueble de esa naturaleza que detente una persona no
es apta para la usucapión.

PRIMER TRIBUNAL COLEGIADO EN MATERIAS PENAL Y CIVIL DEL VIGÉSIMO CIRCUITO.


Amparo directo 260/2024. 20 de febrero de 2025. Unanimidad de votos. Ponente: Refugio Noel
Montoya Moreno. Secretario: Salomón Zenteno Urbina.




                                             Pág. 1 de 2             Fecha de impresión 01/10/2025
 https://sjf2.scjn.gob.mx/detalle/tesis/2030717
                                                   Semanario Judicial de la Federación


Esta tesis se publicó el viernes 04 de julio de 2025 a las 10:11 horas en el Semanario Judicial de la
Federación.




                                              Pág. 2 de 2            Fecha de impresión 01/10/2025
  https://sjf2.scjn.gob.mx/detalle/tesis/2030717

\end{quote}

\subsection{Formalization}

\subsubsection{Clause: tesis2030717\_principio\_fundamental}

\textit{Additional subjects for municipal context Additional assets for municipal public domain (using universal actions) Faith exclusivity: Adverse possession is either good faith XOR bad faith Thesis-specific principle: Municipal public domain properties designated for ecological reserves are imprescriptible}

\textbf{Intermediate Representation:}

\begin{quote}
dominio publico is obligated and reserva ecologica is obligated
\end{quote}

\textbf{Witness Clause:}

\begin{lstlisting}[style=witnessstyle, language=Witness]
clause tesis2030717_principio_fundamental = oblig(dominio_publico) AND oblig(reserva_ecologica) IMPLIES oblig(caracter_imprescriptible);
\end{lstlisting}

\subsubsection{Clause: tesis2030717\_exclusion\_comercial}

\textit{Commercial exclusion: Properties outside commerce cannot be prescribed}

\textbf{Intermediate Representation:}

\begin{quote}
exclusion comercial is obligated
\end{quote}

\textbf{Witness Clause:}

\begin{lstlisting}[style=witnessstyle, language=Witness]
clause tesis2030717_exclusion_comercial = oblig(exclusion_comercial) IMPLIES oblig(caracter_imprescriptible);
\end{lstlisting}

\subsubsection{Clause: tesis2030717\_proteccion\_ecologica}

\textit{Ecological protection: Ecological reserves are protected from prescription}

\textbf{Intermediate Representation:}

\begin{quote}
reserva ecologica is obligated
\end{quote}

\textbf{Witness Clause:}

\begin{lstlisting}[style=witnessstyle, language=Witness]
clause tesis2030717_proteccion_ecologica = oblig(reserva_ecologica) IMPLIES oblig(exclusion_comercial);
\end{lstlisting}

\subsubsection{Clause: tesis2030717\_estatus\_municipal}

\textit{Municipal ownership effect: Municipal public domain creates special legal status}

\textbf{Intermediate Representation:}

\begin{quote}
propiedad municipal is obligated and dominio publico is obligated
\end{quote}

\textbf{Witness Clause:}

\begin{lstlisting}[style=witnessstyle, language=Witness]
clause tesis2030717_estatus_municipal = oblig(propiedad_municipal) AND oblig(dominio_publico) IMPLIES oblig(proteccion_registral);
\end{lstlisting}

\subsubsection{Clause: tesis2030717\_improcedencia}

\textit{Imprescriptibility consequence: Imprescriptible properties cannot be subject to adverse possession}

\textbf{Intermediate Representation:}

\begin{quote}
caracter imprescriptible is obligated
\end{quote}

\textbf{Witness Clause:}

\begin{lstlisting}[style=witnessstyle, language=Witness]
clause tesis2030717_improcedencia = oblig(caracter_imprescriptible) IMPLIES not(claim(demanda_prescripcion));
\end{lstlisting}

\newpage


\end{document}
